% AUTOMATICALLY GENERATED LIBRARY INDEX

\chapter{Canonical Modules}

\section{InfoGeometry.Canonical.AQFTOperatorInterface}

\noindent\textbf{Buildable}: yes

\subsection*{AQFTOperatorRealization}

\begin{theorem}[blueprint="infogeometry:canonical:aqftoperatorinterface:aqftoperatorrealization"]
  \inputleannode{InfoGeometry.Canonical.AQFTOperatorInterface.AQFTOperatorRealization}
\end{theorem}

\subsection*{ComplexAQFTOperatorRealization}

\begin{theorem}[blueprint="infogeometry:canonical:aqftoperatorinterface:complexaqftoperatorrealization"]
  \inputleannode{InfoGeometry.Canonical.AQFTOperatorInterface.ComplexAQFTOperatorRealization}
\end{theorem}

\subsection*{ComplexHilbertObs}

\begin{theorem}[blueprint="infogeometry:canonical:aqftoperatorinterface:complexhilbertobs"]
  \inputleannode{InfoGeometry.Canonical.AQFTOperatorInterface.ComplexHilbertObs}
\end{theorem}

\subsection*{ComplexHilbertOp}

\begin{theorem}[blueprint="infogeometry:canonical:aqftoperatorinterface:complexhilbertop"]
  \inputleannode{InfoGeometry.Canonical.AQFTOperatorInterface.ComplexHilbertOp}
\end{theorem}

\subsection*{IsCStarReady}

\begin{theorem}[blueprint="infogeometry:canonical:aqftoperatorinterface:iscstarready"]
  \inputleannode{InfoGeometry.Canonical.AQFTOperatorInterface.IsCStarReady}
\end{theorem}

\subsection*{IsVonNeumannReady}

\begin{theorem}[blueprint="infogeometry:canonical:aqftoperatorinterface:isvonneumannready"]
  \inputleannode{InfoGeometry.Canonical.AQFTOperatorInterface.IsVonNeumannReady}
\end{theorem}

\subsection*{RealHilbertObs}

\begin{theorem}[blueprint="infogeometry:canonical:aqftoperatorinterface:realhilbertobs"]
  \inputleannode{InfoGeometry.Canonical.AQFTOperatorInterface.RealHilbertObs}
\end{theorem}

\subsection*{aqft_interface_package_realHilbert}

\begin{theorem}[blueprint="infogeometry:canonical:aqftoperatorinterface:aqft_interface_package_realhilbert"]
  \inputleannode{InfoGeometry.Canonical.AQFTOperatorInterface.aqft_interface_package_realHilbert}
\end{theorem}

\subsection*{aqft_interface_package_with_realizations}

\begin{theorem}[blueprint="infogeometry:canonical:aqftoperatorinterface:aqft_interface_package_with_realizations"]
  \inputleannode{InfoGeometry.Canonical.AQFTOperatorInterface.aqft_interface_package_with_realizations}
\end{theorem}

\subsection*{complexFirstLegCompression}

\begin{theorem}[blueprint="infogeometry:canonical:aqftoperatorinterface:complexfirstlegcompression"]
  \inputleannode{InfoGeometry.Canonical.AQFTOperatorInterface.complexFirstLegCompression}
\end{theorem}

\subsection*{complexFirstLegEmbedding}

\begin{theorem}[blueprint="infogeometry:canonical:aqftoperatorinterface:complexfirstlegembedding"]
  \inputleannode{InfoGeometry.Canonical.AQFTOperatorInterface.complexFirstLegEmbedding}
\end{theorem}

\subsection*{complexFirstLegProjection}

\begin{theorem}[blueprint="infogeometry:canonical:aqftoperatorinterface:complexfirstlegprojection"]
  \inputleannode{InfoGeometry.Canonical.AQFTOperatorInterface.complexFirstLegProjection}
\end{theorem}

\subsection*{complexHilbertCompressionRealization}

\begin{theorem}[blueprint="infogeometry:canonical:aqftoperatorinterface:complexhilbertcompressionrealization"]
  \inputleannode{InfoGeometry.Canonical.AQFTOperatorInterface.complexHilbertCompressionRealization}
\end{theorem}

\subsection*{complexHilbertOp_cstarReady}

\begin{theorem}[blueprint="infogeometry:canonical:aqftoperatorinterface:complexhilbertop_cstarready"]
  \inputleannode{InfoGeometry.Canonical.AQFTOperatorInterface.complexHilbertOp_cstarReady}
\end{theorem}

\subsection*{complexHilbertOp_vonNeumannReady}

\begin{theorem}[blueprint="infogeometry:canonical:aqftoperatorinterface:complexhilbertop_vonneumannready"]
  \inputleannode{InfoGeometry.Canonical.AQFTOperatorInterface.complexHilbertOp_vonNeumannReady}
\end{theorem}

\subsection*{cstarReady_of_instance}

\begin{theorem}[blueprint="infogeometry:canonical:aqftoperatorinterface:cstarready_of_instance"]
  \inputleannode{InfoGeometry.Canonical.AQFTOperatorInterface.cstarReady_of_instance}
\end{theorem}

\subsection*{firstLegCompression}

\begin{theorem}[blueprint="infogeometry:canonical:aqftoperatorinterface:firstlegcompression"]
  \inputleannode{InfoGeometry.Canonical.AQFTOperatorInterface.firstLegCompression}
\end{theorem}

\subsection*{firstLegEmbedding}

\begin{theorem}[blueprint="infogeometry:canonical:aqftoperatorinterface:firstlegembedding"]
  \inputleannode{InfoGeometry.Canonical.AQFTOperatorInterface.firstLegEmbedding}
\end{theorem}

\subsection*{firstLegProjection}

\begin{theorem}[blueprint="infogeometry:canonical:aqftoperatorinterface:firstlegprojection"]
  \inputleannode{InfoGeometry.Canonical.AQFTOperatorInterface.firstLegProjection}
\end{theorem}

\subsection*{grandCanonicalFockEulerStep_with_vonNeumannRealization}

\begin{theorem}[blueprint="infogeometry:canonical:aqftoperatorinterface:grandcanonicalfockeulerstep_with_vonneumannrealization"]
  \inputleannode{InfoGeometry.Canonical.AQFTOperatorInterface.grandCanonicalFockEulerStep_with_vonNeumannRealization}
\end{theorem}

\subsection*{realHilbertCompressionRealization}

\begin{theorem}[blueprint="infogeometry:canonical:aqftoperatorinterface:realhilbertcompressionrealization"]
  \inputleannode{InfoGeometry.Canonical.AQFTOperatorInterface.realHilbertCompressionRealization}
\end{theorem}

\subsection*{realHilbertOp_cstarReady}

\begin{theorem}[blueprint="infogeometry:canonical:aqftoperatorinterface:realhilbertop_cstarready"]
  \inputleannode{InfoGeometry.Canonical.AQFTOperatorInterface.realHilbertOp_cstarReady}
\end{theorem}

\subsection*{realHilbertOp_vonNeumannReady}

\begin{theorem}[blueprint="infogeometry:canonical:aqftoperatorinterface:realhilbertop_vonneumannready"]
  \inputleannode{InfoGeometry.Canonical.AQFTOperatorInterface.realHilbertOp_vonNeumannReady}
\end{theorem}

\subsection*{sinkhorn_kmsClosure_with_cstarRealization}

\begin{theorem}[blueprint="infogeometry:canonical:aqftoperatorinterface:sinkhorn_kmsclosure_with_cstarrealization"]
  \inputleannode{InfoGeometry.Canonical.AQFTOperatorInterface.sinkhorn_kmsClosure_with_cstarRealization}
\end{theorem}

\subsection*{vonNeumannReady_of_instance}

\begin{theorem}[blueprint="infogeometry:canonical:aqftoperatorinterface:vonneumannready_of_instance"]
  \inputleannode{InfoGeometry.Canonical.AQFTOperatorInterface.vonNeumannReady_of_instance}
\end{theorem}

\section{InfoGeometry.Canonical.AQFTOperatorInterface.AQFTOperatorRealization}

\subsection*{casesOn}

\begin{theorem}[blueprint="infogeometry:canonical:aqftoperatorinterface:aqftoperatorrealization:caseson"]
  \inputleannode{InfoGeometry.Canonical.AQFTOperatorInterface.AQFTOperatorRealization.casesOn}
\end{theorem}

\subsection*{ctorIdx}

\begin{theorem}[blueprint="infogeometry:canonical:aqftoperatorinterface:aqftoperatorrealization:ctoridx"]
  \inputleannode{InfoGeometry.Canonical.AQFTOperatorInterface.AQFTOperatorRealization.ctorIdx}
\end{theorem}

\subsection*{mk}

\begin{theorem}[blueprint="infogeometry:canonical:aqftoperatorinterface:aqftoperatorrealization:mk"]
  \inputleannode{InfoGeometry.Canonical.AQFTOperatorInterface.AQFTOperatorRealization.mk}
\end{theorem}

\subsection*{noConfusion}

\begin{theorem}[blueprint="infogeometry:canonical:aqftoperatorinterface:aqftoperatorrealization:noconfusion"]
  \inputleannode{InfoGeometry.Canonical.AQFTOperatorInterface.AQFTOperatorRealization.noConfusion}
\end{theorem}

\subsection*{noConfusionType}

\begin{theorem}[blueprint="infogeometry:canonical:aqftoperatorinterface:aqftoperatorrealization:noconfusiontype"]
  \inputleannode{InfoGeometry.Canonical.AQFTOperatorInterface.AQFTOperatorRealization.noConfusionType}
\end{theorem}

\subsection*{realize}

\begin{theorem}[blueprint="infogeometry:canonical:aqftoperatorinterface:aqftoperatorrealization:realize"]
  \inputleannode{InfoGeometry.Canonical.AQFTOperatorInterface.AQFTOperatorRealization.realize}
\end{theorem}

\subsection*{rec}

\begin{theorem}[blueprint="infogeometry:canonical:aqftoperatorinterface:aqftoperatorrealization:rec"]
  \inputleannode{InfoGeometry.Canonical.AQFTOperatorInterface.AQFTOperatorRealization.rec}
\end{theorem}

\subsection*{recOn}

\begin{theorem}[blueprint="infogeometry:canonical:aqftoperatorinterface:aqftoperatorrealization:recon"]
  \inputleannode{InfoGeometry.Canonical.AQFTOperatorInterface.AQFTOperatorRealization.recOn}
\end{theorem}

\section{InfoGeometry.Canonical.AQFTOperatorInterface.ComplexAQFTOperatorRealization}

\subsection*{casesOn}

\begin{theorem}[blueprint="infogeometry:canonical:aqftoperatorinterface:complexaqftoperatorrealization:caseson"]
  \inputleannode{InfoGeometry.Canonical.AQFTOperatorInterface.ComplexAQFTOperatorRealization.casesOn}
\end{theorem}

\subsection*{ctorIdx}

\begin{theorem}[blueprint="infogeometry:canonical:aqftoperatorinterface:complexaqftoperatorrealization:ctoridx"]
  \inputleannode{InfoGeometry.Canonical.AQFTOperatorInterface.ComplexAQFTOperatorRealization.ctorIdx}
\end{theorem}

\subsection*{mk}

\begin{theorem}[blueprint="infogeometry:canonical:aqftoperatorinterface:complexaqftoperatorrealization:mk"]
  \inputleannode{InfoGeometry.Canonical.AQFTOperatorInterface.ComplexAQFTOperatorRealization.mk}
\end{theorem}

\subsection*{noConfusion}

\begin{theorem}[blueprint="infogeometry:canonical:aqftoperatorinterface:complexaqftoperatorrealization:noconfusion"]
  \inputleannode{InfoGeometry.Canonical.AQFTOperatorInterface.ComplexAQFTOperatorRealization.noConfusion}
\end{theorem}

\subsection*{noConfusionType}

\begin{theorem}[blueprint="infogeometry:canonical:aqftoperatorinterface:complexaqftoperatorrealization:noconfusiontype"]
  \inputleannode{InfoGeometry.Canonical.AQFTOperatorInterface.ComplexAQFTOperatorRealization.noConfusionType}
\end{theorem}

\subsection*{realize}

\begin{theorem}[blueprint="infogeometry:canonical:aqftoperatorinterface:complexaqftoperatorrealization:realize"]
  \inputleannode{InfoGeometry.Canonical.AQFTOperatorInterface.ComplexAQFTOperatorRealization.realize}
\end{theorem}

\subsection*{rec}

\begin{theorem}[blueprint="infogeometry:canonical:aqftoperatorinterface:complexaqftoperatorrealization:rec"]
  \inputleannode{InfoGeometry.Canonical.AQFTOperatorInterface.ComplexAQFTOperatorRealization.rec}
\end{theorem}

\subsection*{recOn}

\begin{theorem}[blueprint="infogeometry:canonical:aqftoperatorinterface:complexaqftoperatorrealization:recon"]
  \inputleannode{InfoGeometry.Canonical.AQFTOperatorInterface.ComplexAQFTOperatorRealization.recOn}
\end{theorem}

\section{InfoGeometry.Canonical.Algebra}

\noindent\textbf{Buildable}: yes

\noindent\emph{No exported declarations in the current declaration graph.}

\section{InfoGeometry.Canonical.All}

\noindent\textbf{Buildable}: yes

\noindent\emph{No exported declarations in the current declaration graph.}

\section{InfoGeometry.Canonical.AnalysisLegacy}

\noindent\textbf{Buildable}: yes

\noindent\emph{No exported declarations in the current declaration graph.}

\section{InfoGeometry.Canonical.AnalyticalIndex}

\noindent\textbf{Buildable}: yes

\subsection*{Cl11BottLaplacianZero}

\begin{theorem}[blueprint="infogeometry:canonical:analyticalindex:cl11bottlaplacianzero"]
  \inputleannode{InfoGeometry.Canonical.AnalyticalIndex.Cl11BottLaplacianZero}
\end{theorem}

\subsection*{ConcreteCl11BottFlow}

\begin{theorem}[blueprint="infogeometry:canonical:analyticalindex:concretecl11bottflow"]
  \inputleannode{InfoGeometry.Canonical.AnalyticalIndex.ConcreteCl11BottFlow}
\end{theorem}

\subsection*{FullThermoGeoIndexCapstone}

\begin{theorem}[blueprint="infogeometry:canonical:analyticalindex:fullthermogeoindexcapstone"]
  \inputleannode{InfoGeometry.Canonical.AnalyticalIndex.FullThermoGeoIndexCapstone}
\end{theorem}

\subsection*{GeometricAlgebraicState}

\begin{theorem}[blueprint="infogeometry:canonical:analyticalindex:geometricalgebraicstate"]
  \inputleannode{InfoGeometry.Canonical.AnalyticalIndex.GeometricAlgebraicState}
\end{theorem}

\subsection*{IndexInvariantAlong}

\begin{theorem}[blueprint="infogeometry:canonical:analyticalindex:indexinvariantalong"]
  \inputleannode{InfoGeometry.Canonical.AnalyticalIndex.IndexInvariantAlong}
\end{theorem}

\subsection*{SinkhornKMSCapstone}

\begin{theorem}[blueprint="infogeometry:canonical:analyticalindex:sinkhornkmscapstone"]
  \inputleannode{InfoGeometry.Canonical.AnalyticalIndex.SinkhornKMSCapstone}
\end{theorem}

\subsection*{SinkhornRicciIndexInvariant}

\begin{theorem}[blueprint="infogeometry:canonical:analyticalindex:sinkhornricciindexinvariant"]
  \inputleannode{InfoGeometry.Canonical.AnalyticalIndex.SinkhornRicciIndexInvariant}
\end{theorem}

\subsection*{ThermoGeoIndexCapstoneState}

\begin{theorem}[blueprint="infogeometry:canonical:analyticalindex:thermogeoindexcapstonestate"]
  \inputleannode{InfoGeometry.Canonical.AnalyticalIndex.ThermoGeoIndexCapstoneState}
\end{theorem}

\subsection*{ThermodynamicKMSState}

\begin{theorem}[blueprint="infogeometry:canonical:analyticalindex:thermodynamickmsstate"]
  \inputleannode{InfoGeometry.Canonical.AnalyticalIndex.ThermodynamicKMSState}
\end{theorem}

\subsection*{analyticalIndex}

\begin{theorem}[blueprint="infogeometry:canonical:analyticalindex:analyticalindex"]
  \inputleannode{InfoGeometry.Canonical.AnalyticalIndex.analyticalIndex}
\end{theorem}

\subsection*{analyticalIndex_eq_of_chiralParts_eq}

\begin{theorem}[blueprint="infogeometry:canonical:analyticalindex:analyticalindex_eq_of_chiralparts_eq"]
  \inputleannode{InfoGeometry.Canonical.AnalyticalIndex.analyticalIndex_eq_of_chiralParts_eq}
\end{theorem}

\subsection*{bottAnalyticalIndex}

\begin{theorem}[blueprint="infogeometry:canonical:analyticalindex:bottanalyticalindex"]
  \inputleannode{InfoGeometry.Canonical.AnalyticalIndex.bottAnalyticalIndex}
\end{theorem}

\subsection*{chiralPartMinus}

\begin{theorem}[blueprint="infogeometry:canonical:analyticalindex:chiralpartminus"]
  \inputleannode{InfoGeometry.Canonical.AnalyticalIndex.chiralPartMinus}
\end{theorem}

\subsection*{chiralPartPlus}

\begin{theorem}[blueprint="infogeometry:canonical:analyticalindex:chiralpartplus"]
  \inputleannode{InfoGeometry.Canonical.AnalyticalIndex.chiralPartPlus}
\end{theorem}

\subsection*{chiralProjectorMinus}

\begin{theorem}[blueprint="infogeometry:canonical:analyticalindex:chiralprojectorminus"]
  \inputleannode{InfoGeometry.Canonical.AnalyticalIndex.chiralProjectorMinus}
\end{theorem}

\subsection*{chiralProjectorPlus}

\begin{theorem}[blueprint="infogeometry:canonical:analyticalindex:chiralprojectorplus"]
  \inputleannode{InfoGeometry.Canonical.AnalyticalIndex.chiralProjectorPlus}
\end{theorem}

\subsection*{cl11BottAnalyticalIndex}

\begin{theorem}[blueprint="infogeometry:canonical:analyticalindex:cl11bottanalyticalindex"]
  \inputleannode{InfoGeometry.Canonical.AnalyticalIndex.cl11BottAnalyticalIndex}
\end{theorem}

\subsection*{cl11GlobalGrading}

\begin{theorem}[blueprint="infogeometry:canonical:analyticalindex:cl11globalgrading"]
  \inputleannode{InfoGeometry.Canonical.AnalyticalIndex.cl11GlobalGrading}
\end{theorem}

\subsection*{cl11_bottDirac_sq_eq_zero_of_laplacian_zero}

\begin{theorem}[blueprint="infogeometry:canonical:analyticalindex:cl11_bottdirac_sq_eq_zero_of_laplacian_zero"]
  \inputleannode{InfoGeometry.Canonical.AnalyticalIndex.cl11_bottDirac_sq_eq_zero_of_laplacian_zero}
\end{theorem}

\subsection*{concreteChiralPartMinus_const}

\begin{theorem}[blueprint="infogeometry:canonical:analyticalindex:concretechiralpartminus_const"]
  \inputleannode{InfoGeometry.Canonical.AnalyticalIndex.concreteChiralPartMinus_const}
\end{theorem}

\subsection*{concreteChiralPartPlus_const}

\begin{theorem}[blueprint="infogeometry:canonical:analyticalindex:concretechiralpartplus_const"]
  \inputleannode{InfoGeometry.Canonical.AnalyticalIndex.concreteChiralPartPlus_const}
\end{theorem}

\subsection*{concreteDiracFamily}

\begin{theorem}[blueprint="infogeometry:canonical:analyticalindex:concretediracfamily"]
  \inputleannode{InfoGeometry.Canonical.AnalyticalIndex.concreteDiracFamily}
\end{theorem}

\subsection*{concreteGradingFamily}

\begin{theorem}[blueprint="infogeometry:canonical:analyticalindex:concretegradingfamily"]
  \inputleannode{InfoGeometry.Canonical.AnalyticalIndex.concreteGradingFamily}
\end{theorem}

\subsection*{fullThermoGeoIndexCapstone_of_hypotheses}

\begin{theorem}[blueprint="infogeometry:canonical:analyticalindex:fullthermogeoindexcapstone_of_hypotheses"]
  \inputleannode{InfoGeometry.Canonical.AnalyticalIndex.fullThermoGeoIndexCapstone_of_hypotheses}
\end{theorem}

\subsection*{globalGrading}

\begin{theorem}[blueprint="infogeometry:canonical:analyticalindex:globalgrading"]
  \inputleannode{InfoGeometry.Canonical.AnalyticalIndex.globalGrading}
\end{theorem}

\subsection*{indexInvariantAlong_of_chiralPart_const}

\begin{theorem}[blueprint="infogeometry:canonical:analyticalindex:indexinvariantalong_of_chiralpart_const"]
  \inputleannode{InfoGeometry.Canonical.AnalyticalIndex.indexInvariantAlong_of_chiralPart_const}
\end{theorem}

\subsection*{sinkhornIterate_sinkhornKMSCapstone_of_drive}

\begin{theorem}[blueprint="infogeometry:canonical:analyticalindex:sinkhorniterate_sinkhornkmscapstone_of_drive"]
  \inputleannode{InfoGeometry.Canonical.AnalyticalIndex.sinkhornIterate_sinkhornKMSCapstone_of_drive}
\end{theorem}

\subsection*{sinkhornKMSCapstone_of_drive}

\begin{theorem}[blueprint="infogeometry:canonical:analyticalindex:sinkhornkmscapstone_of_drive"]
  \inputleannode{InfoGeometry.Canonical.AnalyticalIndex.sinkhornKMSCapstone_of_drive}
\end{theorem}

\subsection*{sinkhornKMSState_of_drive}

\begin{theorem}[blueprint="infogeometry:canonical:analyticalindex:sinkhornkmsstate_of_drive"]
  \inputleannode{InfoGeometry.Canonical.AnalyticalIndex.sinkhornKMSState_of_drive}
\end{theorem}

\subsection*{sinkhornRicciIndexInvariant_of_hypotheses}

\begin{theorem}[blueprint="infogeometry:canonical:analyticalindex:sinkhornricciindexinvariant_of_hypotheses"]
  \inputleannode{InfoGeometry.Canonical.AnalyticalIndex.sinkhornRicciIndexInvariant_of_hypotheses}
\end{theorem}

\section{InfoGeometry.Canonical.AnalyticalIndex.ConcreteCl11BottFlow}

\subsection*{Dn}

\begin{theorem}[blueprint="infogeometry:canonical:analyticalindex:concretecl11bottflow:dn"]
  \inputleannode{InfoGeometry.Canonical.AnalyticalIndex.ConcreteCl11BottFlow.Dn}
\end{theorem}

\subsection*{Dn_const}

\begin{theorem}[blueprint="infogeometry:canonical:analyticalindex:concretecl11bottflow:dn_const"]
  \inputleannode{InfoGeometry.Canonical.AnalyticalIndex.ConcreteCl11BottFlow.Dn_const}
\end{theorem}

\subsection*{casesOn}

\begin{theorem}[blueprint="infogeometry:canonical:analyticalindex:concretecl11bottflow:caseson"]
  \inputleannode{InfoGeometry.Canonical.AnalyticalIndex.ConcreteCl11BottFlow.casesOn}
\end{theorem}

\subsection*{ctorIdx}

\begin{theorem}[blueprint="infogeometry:canonical:analyticalindex:concretecl11bottflow:ctoridx"]
  \inputleannode{InfoGeometry.Canonical.AnalyticalIndex.ConcreteCl11BottFlow.ctorIdx}
\end{theorem}

\subsection*{indexInvariant}

\begin{theorem}[blueprint="infogeometry:canonical:analyticalindex:concretecl11bottflow:indexinvariant"]
  \inputleannode{InfoGeometry.Canonical.AnalyticalIndex.ConcreteCl11BottFlow.indexInvariant}
\end{theorem}

\subsection*{mk}

\begin{theorem}[blueprint="infogeometry:canonical:analyticalindex:concretecl11bottflow:mk"]
  \inputleannode{InfoGeometry.Canonical.AnalyticalIndex.ConcreteCl11BottFlow.mk}
\end{theorem}

\subsection*{noConfusion}

\begin{theorem}[blueprint="infogeometry:canonical:analyticalindex:concretecl11bottflow:noconfusion"]
  \inputleannode{InfoGeometry.Canonical.AnalyticalIndex.ConcreteCl11BottFlow.noConfusion}
\end{theorem}

\subsection*{noConfusionType}

\begin{theorem}[blueprint="infogeometry:canonical:analyticalindex:concretecl11bottflow:noconfusiontype"]
  \inputleannode{InfoGeometry.Canonical.AnalyticalIndex.ConcreteCl11BottFlow.noConfusionType}
\end{theorem}

\subsection*{rec}

\begin{theorem}[blueprint="infogeometry:canonical:analyticalindex:concretecl11bottflow:rec"]
  \inputleannode{InfoGeometry.Canonical.AnalyticalIndex.ConcreteCl11BottFlow.rec}
\end{theorem}

\subsection*{recOn}

\begin{theorem}[blueprint="infogeometry:canonical:analyticalindex:concretecl11bottflow:recon"]
  \inputleannode{InfoGeometry.Canonical.AnalyticalIndex.ConcreteCl11BottFlow.recOn}
\end{theorem}

\subsection*{Γn}

\begin{theorem}[blueprint="infogeometry:canonical:analyticalindex:concretecl11bottflow:γn"]
  \inputleannode{InfoGeometry.Canonical.AnalyticalIndex.ConcreteCl11BottFlow.Γn}
\end{theorem}

\subsection*{Γn_const}

\begin{theorem}[blueprint="infogeometry:canonical:analyticalindex:concretecl11bottflow:γn_const"]
  \inputleannode{InfoGeometry.Canonical.AnalyticalIndex.ConcreteCl11BottFlow.Γn_const}
\end{theorem}

\section{InfoGeometry.Canonical.AnalyticalIndex.FullThermoGeoIndexCapstone}

\subsection*{geometricAlgebraicState}

\begin{theorem}[blueprint="infogeometry:canonical:analyticalindex:fullthermogeoindexcapstone:geometricalgebraicstate"]
  \inputleannode{InfoGeometry.Canonical.AnalyticalIndex.FullThermoGeoIndexCapstone.geometricAlgebraicState}
\end{theorem}

\subsection*{thermodynamicKMSState}

\begin{theorem}[blueprint="infogeometry:canonical:analyticalindex:fullthermogeoindexcapstone:thermodynamickmsstate"]
  \inputleannode{InfoGeometry.Canonical.AnalyticalIndex.FullThermoGeoIndexCapstone.thermodynamicKMSState}
\end{theorem}

\section{InfoGeometry.Canonical.AnalyticalIndex.SinkhornKMSCapstone}

\subsection*{kmsState}

\begin{theorem}[blueprint="infogeometry:canonical:analyticalindex:sinkhornkmscapstone:kmsstate"]
  \inputleannode{InfoGeometry.Canonical.AnalyticalIndex.SinkhornKMSCapstone.kmsState}
\end{theorem}

\section{InfoGeometry.Canonical.AnalyticalIndex.analyticalIndex}

\subsection*{eq_1}

\begin{theorem}[blueprint="infogeometry:canonical:analyticalindex:analyticalindex:eq_1"]
  \inputleannode{InfoGeometry.Canonical.AnalyticalIndex.analyticalIndex.eq_1}
\end{theorem}

\section{InfoGeometry.Canonical.AnalyticalIndex.concreteDiracFamily}

\subsection*{eq_1}

\begin{theorem}[blueprint="infogeometry:canonical:analyticalindex:concretediracfamily:eq_1"]
  \inputleannode{InfoGeometry.Canonical.AnalyticalIndex.concreteDiracFamily.eq_1}
\end{theorem}

\section{InfoGeometry.Canonical.AnalyticalIndex.concreteGradingFamily}

\subsection*{eq_1}

\begin{theorem}[blueprint="infogeometry:canonical:analyticalindex:concretegradingfamily:eq_1"]
  \inputleannode{InfoGeometry.Canonical.AnalyticalIndex.concreteGradingFamily.eq_1}
\end{theorem}

\section{InfoGeometry.Canonical.AnomalyInflow}

\noindent\textbf{Buildable}: yes

\subsection*{AnomalyInflowClosure}

\begin{theorem}[blueprint="infogeometry:canonical:anomalyinflow:anomalyinflowclosure"]
  \inputleannode{InfoGeometry.Canonical.AnomalyInflow.AnomalyInflowClosure}
\end{theorem}

\subsection*{anomalyInflowClosure}

\begin{theorem}[blueprint="infogeometry:canonical:anomalyinflow:anomalyinflowclosure"]
  \inputleannode{InfoGeometry.Canonical.AnomalyInflow.anomalyInflowClosure}
\end{theorem}

\subsection*{anomaly_inflow_cancellation}

\begin{theorem}[blueprint="infogeometry:canonical:anomalyinflow:anomaly_inflow_cancellation"]
  \inputleannode{InfoGeometry.Canonical.AnomalyInflow.anomaly_inflow_cancellation}
\end{theorem}

\subsection*{boundaryAnomaly}

\begin{theorem}[blueprint="infogeometry:canonical:anomalyinflow:boundaryanomaly"]
  \inputleannode{InfoGeometry.Canonical.AnomalyInflow.boundaryAnomaly}
\end{theorem}

\subsection*{boundaryAnomalyDensity}

\begin{theorem}[blueprint="infogeometry:canonical:anomalyinflow:boundaryanomalydensity"]
  \inputleannode{InfoGeometry.Canonical.AnomalyInflow.boundaryAnomalyDensity}
\end{theorem}

\subsection*{bulkChernSimonsVariation}

\begin{theorem}[blueprint="infogeometry:canonical:anomalyinflow:bulkchernsimonsvariation"]
  \inputleannode{InfoGeometry.Canonical.AnomalyInflow.bulkChernSimonsVariation}
\end{theorem}

\subsection*{variationChernSimons}

\begin{theorem}[blueprint="infogeometry:canonical:anomalyinflow:variationchernsimons"]
  \inputleannode{InfoGeometry.Canonical.AnomalyInflow.variationChernSimons}
\end{theorem}

\section{InfoGeometry.Canonical.AnomalyInflow.AnomalyInflowClosure}

\subsection*{eq_1}

\begin{theorem}[blueprint="infogeometry:canonical:anomalyinflow:anomalyinflowclosure:eq_1"]
  \inputleannode{InfoGeometry.Canonical.AnomalyInflow.AnomalyInflowClosure.eq_1}
\end{theorem}

\section{InfoGeometry.Canonical.Attention}

\noindent\textbf{Buildable}: yes

\subsection*{ContextWindow}

\begin{theorem}[blueprint="infogeometry:canonical:attention:contextwindow"]
  \inputleannode{InfoGeometry.Canonical.Attention.ContextWindow}
\end{theorem}

\subsection*{Head}

\begin{theorem}[blueprint="infogeometry:canonical:attention:head"]
  \inputleannode{InfoGeometry.Canonical.Attention.Head}
\end{theorem}

\subsection*{Vec}

\begin{theorem}[blueprint="infogeometry:canonical:attention:vec"]
  \inputleannode{InfoGeometry.Canonical.Attention.Vec}
\end{theorem}

\subsection*{attentionHead}

\begin{theorem}[blueprint="infogeometry:canonical:attention:attentionhead"]
  \inputleannode{InfoGeometry.Canonical.Attention.attentionHead}
\end{theorem}

\subsection*{attentionParams}

\begin{theorem}[blueprint="infogeometry:canonical:attention:attentionparams"]
  \inputleannode{InfoGeometry.Canonical.Attention.attentionParams}
\end{theorem}

\subsection*{attentionWeights}

\begin{theorem}[blueprint="infogeometry:canonical:attention:attentionweights"]
  \inputleannode{InfoGeometry.Canonical.Attention.attentionWeights}
\end{theorem}

\subsection*{attentionWeights_sum_one}

\begin{theorem}[blueprint="infogeometry:canonical:attention:attentionweights_sum_one"]
  \inputleannode{InfoGeometry.Canonical.Attention.attentionWeights_sum_one}
\end{theorem}

\subsection*{dot}

\begin{theorem}[blueprint="infogeometry:canonical:attention:dot"]
  \inputleannode{InfoGeometry.Canonical.Attention.dot}
\end{theorem}

\subsection*{euclideanAttentionHead}

\begin{theorem}[blueprint="infogeometry:canonical:attention:euclideanattentionhead"]
  \inputleannode{InfoGeometry.Canonical.Attention.euclideanAttentionHead}
\end{theorem}

\subsection*{euclideanAttentionParams}

\begin{theorem}[blueprint="infogeometry:canonical:attention:euclideanattentionparams"]
  \inputleannode{InfoGeometry.Canonical.Attention.euclideanAttentionParams}
\end{theorem}

\subsection*{euclideanAttentionWeights}

\begin{theorem}[blueprint="infogeometry:canonical:attention:euclideanattentionweights"]
  \inputleannode{InfoGeometry.Canonical.Attention.euclideanAttentionWeights}
\end{theorem}

\subsection*{euclideanAttentionWeights_sum_one}

\begin{theorem}[blueprint="infogeometry:canonical:attention:euclideanattentionweights_sum_one"]
  \inputleannode{InfoGeometry.Canonical.Attention.euclideanAttentionWeights_sum_one}
\end{theorem}

\subsection*{euclideanMatchForm}

\begin{theorem}[blueprint="infogeometry:canonical:attention:euclideanmatchform"]
  \inputleannode{InfoGeometry.Canonical.Attention.euclideanMatchForm}
\end{theorem}

\subsection*{headOutput}

\begin{theorem}[blueprint="infogeometry:canonical:attention:headoutput"]
  \inputleannode{InfoGeometry.Canonical.Attention.headOutput}
\end{theorem}

\subsection*{headWeights}

\begin{theorem}[blueprint="infogeometry:canonical:attention:headweights"]
  \inputleannode{InfoGeometry.Canonical.Attention.headWeights}
\end{theorem}

\subsection*{interactionEnergy}

\begin{theorem}[blueprint="infogeometry:canonical:attention:interactionenergy"]
  \inputleannode{InfoGeometry.Canonical.Attention.interactionEnergy}
\end{theorem}

\subsection*{lorentzianAttentionHead}

\begin{theorem}[blueprint="infogeometry:canonical:attention:lorentzianattentionhead"]
  \inputleannode{InfoGeometry.Canonical.Attention.lorentzianAttentionHead}
\end{theorem}

\subsection*{lorentzianAttentionParams}

\begin{theorem}[blueprint="infogeometry:canonical:attention:lorentzianattentionparams"]
  \inputleannode{InfoGeometry.Canonical.Attention.lorentzianAttentionParams}
\end{theorem}

\subsection*{lorentzianAttentionWeights}

\begin{theorem}[blueprint="infogeometry:canonical:attention:lorentzianattentionweights"]
  \inputleannode{InfoGeometry.Canonical.Attention.lorentzianAttentionWeights}
\end{theorem}

\subsection*{lorentzianAttentionWeights_sum_one}

\begin{theorem}[blueprint="infogeometry:canonical:attention:lorentzianattentionweights_sum_one"]
  \inputleannode{InfoGeometry.Canonical.Attention.lorentzianAttentionWeights_sum_one}
\end{theorem}

\subsection*{partition}

\begin{theorem}[blueprint="infogeometry:canonical:attention:partition"]
  \inputleannode{InfoGeometry.Canonical.Attention.partition}
\end{theorem}

\subsection*{score}

\begin{theorem}[blueprint="infogeometry:canonical:attention:score"]
  \inputleannode{InfoGeometry.Canonical.Attention.score}
\end{theorem}

\subsection*{softmax}

\begin{theorem}[blueprint="infogeometry:canonical:attention:softmax"]
  \inputleannode{InfoGeometry.Canonical.Attention.softmax}
\end{theorem}

\subsection*{timelike_dominance}

\begin{theorem}[blueprint="infogeometry:canonical:attention:timelike_dominance"]
  \inputleannode{InfoGeometry.Canonical.Attention.timelike_dominance}
\end{theorem}

\subsection*{weightedSum}

\begin{theorem}[blueprint="infogeometry:canonical:attention:weightedsum"]
  \inputleannode{InfoGeometry.Canonical.Attention.weightedSum}
\end{theorem}

\section{InfoGeometry.Canonical.Attention.ContextWindow}

\subsection*{casesOn}

\begin{theorem}[blueprint="infogeometry:canonical:attention:contextwindow:caseson"]
  \inputleannode{InfoGeometry.Canonical.Attention.ContextWindow.casesOn}
\end{theorem}

\subsection*{ctorIdx}

\begin{theorem}[blueprint="infogeometry:canonical:attention:contextwindow:ctoridx"]
  \inputleannode{InfoGeometry.Canonical.Attention.ContextWindow.ctorIdx}
\end{theorem}

\subsection*{keys}

\begin{theorem}[blueprint="infogeometry:canonical:attention:contextwindow:keys"]
  \inputleannode{InfoGeometry.Canonical.Attention.ContextWindow.keys}
\end{theorem}

\subsection*{mk}

\begin{theorem}[blueprint="infogeometry:canonical:attention:contextwindow:mk"]
  \inputleannode{InfoGeometry.Canonical.Attention.ContextWindow.mk}
\end{theorem}

\subsection*{noConfusion}

\begin{theorem}[blueprint="infogeometry:canonical:attention:contextwindow:noconfusion"]
  \inputleannode{InfoGeometry.Canonical.Attention.ContextWindow.noConfusion}
\end{theorem}

\subsection*{noConfusionType}

\begin{theorem}[blueprint="infogeometry:canonical:attention:contextwindow:noconfusiontype"]
  \inputleannode{InfoGeometry.Canonical.Attention.ContextWindow.noConfusionType}
\end{theorem}

\subsection*{rec}

\begin{theorem}[blueprint="infogeometry:canonical:attention:contextwindow:rec"]
  \inputleannode{InfoGeometry.Canonical.Attention.ContextWindow.rec}
\end{theorem}

\subsection*{recOn}

\begin{theorem}[blueprint="infogeometry:canonical:attention:contextwindow:recon"]
  \inputleannode{InfoGeometry.Canonical.Attention.ContextWindow.recOn}
\end{theorem}

\subsection*{values}

\begin{theorem}[blueprint="infogeometry:canonical:attention:contextwindow:values"]
  \inputleannode{InfoGeometry.Canonical.Attention.ContextWindow.values}
\end{theorem}

\section{InfoGeometry.Canonical.Attention.Head}

\subsection*{WK}

\begin{theorem}[blueprint="infogeometry:canonical:attention:head:wk"]
  \inputleannode{InfoGeometry.Canonical.Attention.Head.WK}
\end{theorem}

\subsection*{WQ}

\begin{theorem}[blueprint="infogeometry:canonical:attention:head:wq"]
  \inputleannode{InfoGeometry.Canonical.Attention.Head.WQ}
\end{theorem}

\subsection*{WV}

\begin{theorem}[blueprint="infogeometry:canonical:attention:head:wv"]
  \inputleannode{InfoGeometry.Canonical.Attention.Head.WV}
\end{theorem}

\subsection*{casesOn}

\begin{theorem}[blueprint="infogeometry:canonical:attention:head:caseson"]
  \inputleannode{InfoGeometry.Canonical.Attention.Head.casesOn}
\end{theorem}

\subsection*{ctorIdx}

\begin{theorem}[blueprint="infogeometry:canonical:attention:head:ctoridx"]
  \inputleannode{InfoGeometry.Canonical.Attention.Head.ctorIdx}
\end{theorem}

\subsection*{mk}

\begin{theorem}[blueprint="infogeometry:canonical:attention:head:mk"]
  \inputleannode{InfoGeometry.Canonical.Attention.Head.mk}
\end{theorem}

\subsection*{noConfusion}

\begin{theorem}[blueprint="infogeometry:canonical:attention:head:noconfusion"]
  \inputleannode{InfoGeometry.Canonical.Attention.Head.noConfusion}
\end{theorem}

\subsection*{noConfusionType}

\begin{theorem}[blueprint="infogeometry:canonical:attention:head:noconfusiontype"]
  \inputleannode{InfoGeometry.Canonical.Attention.Head.noConfusionType}
\end{theorem}

\subsection*{rec}

\begin{theorem}[blueprint="infogeometry:canonical:attention:head:rec"]
  \inputleannode{InfoGeometry.Canonical.Attention.Head.rec}
\end{theorem}

\subsection*{recOn}

\begin{theorem}[blueprint="infogeometry:canonical:attention:head:recon"]
  \inputleannode{InfoGeometry.Canonical.Attention.Head.recOn}
\end{theorem}

\section{InfoGeometry.Canonical.Attention.attentionWeights}

\subsection*{eq_1}

\begin{theorem}[blueprint="infogeometry:canonical:attention:attentionweights:eq_1"]
  \inputleannode{InfoGeometry.Canonical.Attention.attentionWeights.eq_1}
\end{theorem}

\section{InfoGeometry.Canonical.AttentionEuclidean}

\noindent\textbf{Buildable}: yes

\noindent\emph{No exported declarations in the current declaration graph.}

\section{InfoGeometry.Canonical.AttentionSplit}

\noindent\textbf{Buildable}: yes

\noindent\emph{No exported declarations in the current declaration graph.}

\section{InfoGeometry.Canonical.BeliefAlgebra}

\noindent\textbf{Buildable}: yes

\subsection*{BeliefSystem}

\begin{theorem}[blueprint="infogeometry:canonical:beliefalgebra:beliefsystem"]
  \inputleannode{InfoGeometry.Canonical.BeliefAlgebra.BeliefSystem}
\end{theorem}

\section{InfoGeometry.Canonical.BeliefAlgebra.BeliefSystem}

\subsection*{InformationLieAlgebra}

\begin{theorem}[blueprint="infogeometry:canonical:beliefalgebra:beliefsystem:informationliealgebra"]
  \inputleannode{InfoGeometry.Canonical.BeliefAlgebra.BeliefSystem.InformationLieAlgebra}
\end{theorem}

\subsection*{Twisted}

\begin{theorem}[blueprint="infogeometry:canonical:beliefalgebra:beliefsystem:twisted"]
  \inputleannode{InfoGeometry.Canonical.BeliefAlgebra.BeliefSystem.Twisted}
\end{theorem}

\subsection*{casesOn}

\begin{theorem}[blueprint="infogeometry:canonical:beliefalgebra:beliefsystem:caseson"]
  \inputleannode{InfoGeometry.Canonical.BeliefAlgebra.BeliefSystem.casesOn}
\end{theorem}

\subsection*{ctorIdx}

\begin{theorem}[blueprint="infogeometry:canonical:beliefalgebra:beliefsystem:ctoridx"]
  \inputleannode{InfoGeometry.Canonical.BeliefAlgebra.BeliefSystem.ctorIdx}
\end{theorem}

\subsection*{mk}

\begin{theorem}[blueprint="infogeometry:canonical:beliefalgebra:beliefsystem:mk"]
  \inputleannode{InfoGeometry.Canonical.BeliefAlgebra.BeliefSystem.mk}
\end{theorem}

\subsection*{noConfusion}

\begin{theorem}[blueprint="infogeometry:canonical:beliefalgebra:beliefsystem:noconfusion"]
  \inputleannode{InfoGeometry.Canonical.BeliefAlgebra.BeliefSystem.noConfusion}
\end{theorem}

\subsection*{noConfusionType}

\begin{theorem}[blueprint="infogeometry:canonical:beliefalgebra:beliefsystem:noconfusiontype"]
  \inputleannode{InfoGeometry.Canonical.BeliefAlgebra.BeliefSystem.noConfusionType}
\end{theorem}

\subsection*{non_commutative_updates}

\begin{theorem}[blueprint="infogeometry:canonical:beliefalgebra:beliefsystem:non_commutative_updates"]
  \inputleannode{InfoGeometry.Canonical.BeliefAlgebra.BeliefSystem.non_commutative_updates}
\end{theorem}

\subsection*{rec}

\begin{theorem}[blueprint="infogeometry:canonical:beliefalgebra:beliefsystem:rec"]
  \inputleannode{InfoGeometry.Canonical.BeliefAlgebra.BeliefSystem.rec}
\end{theorem}

\subsection*{recOn}

\begin{theorem}[blueprint="infogeometry:canonical:beliefalgebra:beliefsystem:recon"]
  \inputleannode{InfoGeometry.Canonical.BeliefAlgebra.BeliefSystem.recOn}
\end{theorem}

\section{InfoGeometry.Canonical.BeliefDynamics}

\noindent\textbf{Buildable}: yes

\subsection*{exponentialTilt}

\begin{theorem}[blueprint="infogeometry:canonical:beliefdynamics:exponentialtilt"]
  \inputleannode{InfoGeometry.Canonical.BeliefDynamics.exponentialTilt}
\end{theorem}

\subsection*{parallelTransportE}

\begin{theorem}[blueprint="infogeometry:canonical:beliefdynamics:paralleltransporte"]
  \inputleannode{InfoGeometry.Canonical.BeliefDynamics.parallelTransportE}
\end{theorem}

\subsection*{parallelTransportM}

\begin{theorem}[blueprint="infogeometry:canonical:beliefdynamics:paralleltransportm"]
  \inputleannode{InfoGeometry.Canonical.BeliefDynamics.parallelTransportM}
\end{theorem}

\subsection*{quantumGeometryOp}

\begin{theorem}[blueprint="infogeometry:canonical:beliefdynamics:quantumgeometryop"]
  \inputleannode{InfoGeometry.Canonical.BeliefDynamics.quantumGeometryOp}
\end{theorem}

\subsection*{radonNikodymOp}

\begin{theorem}[blueprint="infogeometry:canonical:beliefdynamics:radonnikodymop"]
  \inputleannode{InfoGeometry.Canonical.BeliefDynamics.radonNikodymOp}
\end{theorem}

\section{InfoGeometry.Canonical.BerryPhase}

\noindent\textbf{Buildable}: yes

\subsection*{berryConnection}

\begin{theorem}[blueprint="infogeometry:canonical:berryphase:berryconnection"]
  \inputleannode{InfoGeometry.Canonical.BerryPhase.berryConnection}
\end{theorem}

\subsection*{berry_phase_vanishes_for_normal}

\begin{theorem}[blueprint="infogeometry:canonical:berryphase:berry_phase_vanishes_for_normal"]
  \inputleannode{InfoGeometry.Canonical.BerryPhase.berry_phase_vanishes_for_normal}
\end{theorem}

\subsection*{informationBerryPhase}

\begin{theorem}[blueprint="infogeometry:canonical:berryphase:informationberryphase"]
  \inputleannode{InfoGeometry.Canonical.BerryPhase.informationBerryPhase}
\end{theorem}

\section{InfoGeometry.Canonical.BerryPhase.informationBerryPhase}

\subsection*{eq_1}

\begin{theorem}[blueprint="infogeometry:canonical:berryphase:informationberryphase:eq_1"]
  \inputleannode{InfoGeometry.Canonical.BerryPhase.informationBerryPhase.eq_1}
\end{theorem}

\section{InfoGeometry.Canonical.BogoliubovFockSuper}

\noindent\textbf{Buildable}: yes

\subsection*{BogoliubovMixingParams}

\begin{theorem}[blueprint="infogeometry:canonical:bogoliubovfocksuper:bogoliubovmixingparams"]
  \inputleannode{InfoGeometry.Canonical.BogoliubovFockSuper.BogoliubovMixingParams}
\end{theorem}

\subsection*{BogoliubovParams}

\begin{theorem}[blueprint="infogeometry:canonical:bogoliubovfocksuper:bogoliubovparams"]
  \inputleannode{InfoGeometry.Canonical.BogoliubovFockSuper.BogoliubovParams}
\end{theorem}

\subsection*{CARClosure}

\begin{theorem}[blueprint="infogeometry:canonical:bogoliubovfocksuper:carclosure"]
  \inputleannode{InfoGeometry.Canonical.BogoliubovFockSuper.CARClosure}
\end{theorem}

\subsection*{CARWitness}

\begin{theorem}[blueprint="infogeometry:canonical:bogoliubovfocksuper:carwitness"]
  \inputleannode{InfoGeometry.Canonical.BogoliubovFockSuper.CARWitness}
\end{theorem}

\subsection*{CCRClosure}

\begin{theorem}[blueprint="infogeometry:canonical:bogoliubovfocksuper:ccrclosure"]
  \inputleannode{InfoGeometry.Canonical.BogoliubovFockSuper.CCRClosure}
\end{theorem}

\subsection*{CCRWitness}

\begin{theorem}[blueprint="infogeometry:canonical:bogoliubovfocksuper:ccrwitness"]
  \inputleannode{InfoGeometry.Canonical.BogoliubovFockSuper.CCRWitness}
\end{theorem}

\subsection*{FockEnd}

\begin{theorem}[blueprint="infogeometry:canonical:bogoliubovfocksuper:fockend"]
  \inputleannode{InfoGeometry.Canonical.BogoliubovFockSuper.FockEnd}
\end{theorem}

\subsection*{FockEndomorphism}

\begin{theorem}[blueprint="infogeometry:canonical:bogoliubovfocksuper:fockendomorphism"]
  \inputleannode{InfoGeometry.Canonical.BogoliubovFockSuper.FockEndomorphism}
\end{theorem}

\subsection*{SuperParity}

\begin{theorem}[blueprint="infogeometry:canonical:bogoliubovfocksuper:superparity"]
  \inputleannode{InfoGeometry.Canonical.BogoliubovFockSuper.SuperParity}
\end{theorem}

\subsection*{anticommutator}

\begin{theorem}[blueprint="infogeometry:canonical:bogoliubovfocksuper:anticommutator"]
  \inputleannode{InfoGeometry.Canonical.BogoliubovFockSuper.anticommutator}
\end{theorem}

\subsection*{anticommutator_bogoliubov_of_CAR}

\begin{theorem}[blueprint="infogeometry:canonical:bogoliubovfocksuper:anticommutator_bogoliubov_of_car"]
  \inputleannode{InfoGeometry.Canonical.BogoliubovFockSuper.anticommutator_bogoliubov_of_CAR}
\end{theorem}

\subsection*{anticommutator_symm}

\begin{theorem}[blueprint="infogeometry:canonical:bogoliubovfocksuper:anticommutator_symm"]
  \inputleannode{InfoGeometry.Canonical.BogoliubovFockSuper.anticommutator_symm}
\end{theorem}

\subsection*{bogoliubovAnnihilation}

\begin{theorem}[blueprint="infogeometry:canonical:bogoliubovfocksuper:bogoliubovannihilation"]
  \inputleannode{InfoGeometry.Canonical.BogoliubovFockSuper.bogoliubovAnnihilation}
\end{theorem}

\subsection*{bogoliubovAnnihilation_kills_vacuumVector}

\begin{theorem}[blueprint="infogeometry:canonical:bogoliubovfocksuper:bogoliubovannihilation_kills_vacuumvector"]
  \inputleannode{InfoGeometry.Canonical.BogoliubovFockSuper.bogoliubovAnnihilation_kills_vacuumVector}
\end{theorem}

\subsection*{bogoliubovCreation}

\begin{theorem}[blueprint="infogeometry:canonical:bogoliubovfocksuper:bogoliubovcreation"]
  \inputleannode{InfoGeometry.Canonical.BogoliubovFockSuper.bogoliubovCreation}
\end{theorem}

\subsection*{bogoliubovCreation_kills_vacuumVector}

\begin{theorem}[blueprint="infogeometry:canonical:bogoliubovfocksuper:bogoliubovcreation_kills_vacuumvector"]
  \inputleannode{InfoGeometry.Canonical.BogoliubovFockSuper.bogoliubovCreation_kills_vacuumVector}
\end{theorem}

\subsection*{bogoliubovNumberOperator}

\begin{theorem}[blueprint="infogeometry:canonical:bogoliubovfocksuper:bogoliubovnumberoperator"]
  \inputleannode{InfoGeometry.Canonical.BogoliubovFockSuper.bogoliubovNumberOperator}
\end{theorem}

\subsection*{commutator}

\begin{theorem}[blueprint="infogeometry:canonical:bogoliubovfocksuper:commutator"]
  \inputleannode{InfoGeometry.Canonical.BogoliubovFockSuper.commutator}
\end{theorem}

\subsection*{commutator_bogoliubov_of_CCR}

\begin{theorem}[blueprint="infogeometry:canonical:bogoliubovfocksuper:commutator_bogoliubov_of_ccr"]
  \inputleannode{InfoGeometry.Canonical.BogoliubovFockSuper.commutator_bogoliubov_of_CCR}
\end{theorem}

\subsection*{commutator_swap}

\begin{theorem}[blueprint="infogeometry:canonical:bogoliubovfocksuper:commutator_swap"]
  \inputleannode{InfoGeometry.Canonical.BogoliubovFockSuper.commutator_swap}
\end{theorem}

\subsection*{einsteinFockDeformation}

\begin{theorem}[blueprint="infogeometry:canonical:bogoliubovfocksuper:einsteinfockdeformation"]
  \inputleannode{InfoGeometry.Canonical.BogoliubovFockSuper.einsteinFockDeformation}
\end{theorem}

\subsection*{einsteinFockDeformationOperator}

\begin{theorem}[blueprint="infogeometry:canonical:bogoliubovfocksuper:einsteinfockdeformationoperator"]
  \inputleannode{InfoGeometry.Canonical.BogoliubovFockSuper.einsteinFockDeformationOperator}
\end{theorem}

\subsection*{einsteinFockDeformation_eq_zero_of_vacuumTransported}

\begin{theorem}[blueprint="infogeometry:canonical:bogoliubovfocksuper:einsteinfockdeformation_eq_zero_of_vacuumtransported"]
  \inputleannode{InfoGeometry.Canonical.BogoliubovFockSuper.einsteinFockDeformation_eq_zero_of_vacuumTransported}
\end{theorem}

\subsection*{einsteinInducedChemicalPotential}

\begin{theorem}[blueprint="infogeometry:canonical:bogoliubovfocksuper:einsteininducedchemicalpotential"]
  \inputleannode{InfoGeometry.Canonical.BogoliubovFockSuper.einsteinInducedChemicalPotential}
\end{theorem}

\subsection*{fockAnticommutator}

\begin{theorem}[blueprint="infogeometry:canonical:bogoliubovfocksuper:fockanticommutator"]
  \inputleannode{InfoGeometry.Canonical.BogoliubovFockSuper.fockAnticommutator}
\end{theorem}

\subsection*{fockAnticommutator_symm}

\begin{theorem}[blueprint="infogeometry:canonical:bogoliubovfocksuper:fockanticommutator_symm"]
  \inputleannode{InfoGeometry.Canonical.BogoliubovFockSuper.fockAnticommutator_symm}
\end{theorem}

\subsection*{fockCommutator}

\begin{theorem}[blueprint="infogeometry:canonical:bogoliubovfocksuper:fockcommutator"]
  \inputleannode{InfoGeometry.Canonical.BogoliubovFockSuper.fockCommutator}
\end{theorem}

\subsection*{fockCommutator_swap}

\begin{theorem}[blueprint="infogeometry:canonical:bogoliubovfocksuper:fockcommutator_swap"]
  \inputleannode{InfoGeometry.Canonical.BogoliubovFockSuper.fockCommutator_swap}
\end{theorem}

\subsection*{fockSuperBracket}

\begin{theorem}[blueprint="infogeometry:canonical:bogoliubovfocksuper:focksuperbracket"]
  \inputleannode{InfoGeometry.Canonical.BogoliubovFockSuper.fockSuperBracket}
\end{theorem}

\subsection*{grandCanonicalEulerStep}

\begin{theorem}[blueprint="infogeometry:canonical:bogoliubovfocksuper:grandcanonicaleulerstep"]
  \inputleannode{InfoGeometry.Canonical.BogoliubovFockSuper.grandCanonicalEulerStep}
\end{theorem}

\subsection*{grandCanonicalFockEulerStep}

\begin{theorem}[blueprint="infogeometry:canonical:bogoliubovfocksuper:grandcanonicalfockeulerstep"]
  \inputleannode{InfoGeometry.Canonical.BogoliubovFockSuper.grandCanonicalFockEulerStep}
\end{theorem}

\subsection*{grandCanonicalFockGenerator}

\begin{theorem}[blueprint="infogeometry:canonical:bogoliubovfocksuper:grandcanonicalfockgenerator"]
  \inputleannode{InfoGeometry.Canonical.BogoliubovFockSuper.grandCanonicalFockGenerator}
\end{theorem}

\subsection*{grandCanonicalFockGenerator_eq_hamiltonian_of_vacuumTransported}

\begin{theorem}[blueprint="infogeometry:canonical:bogoliubovfocksuper:grandcanonicalfockgenerator_eq_hamiltonian_of_vacuumtransported"]
  \inputleannode{InfoGeometry.Canonical.BogoliubovFockSuper.grandCanonicalFockGenerator_eq_hamiltonian_of_vacuumTransported}
\end{theorem}

\subsection*{grandCanonicalGenerator}

\begin{theorem}[blueprint="infogeometry:canonical:bogoliubovfocksuper:grandcanonicalgenerator"]
  \inputleannode{InfoGeometry.Canonical.BogoliubovFockSuper.grandCanonicalGenerator}
\end{theorem}

\subsection*{grandCanonicalGenerator_eq_hamiltonian_of_vacuumTransported}

\begin{theorem}[blueprint="infogeometry:canonical:bogoliubovfocksuper:grandcanonicalgenerator_eq_hamiltonian_of_vacuumtransported"]
  \inputleannode{InfoGeometry.Canonical.BogoliubovFockSuper.grandCanonicalGenerator_eq_hamiltonian_of_vacuumTransported}
\end{theorem}

\subsection*{inducedChemicalPotential}

\begin{theorem}[blueprint="infogeometry:canonical:bogoliubovfocksuper:inducedchemicalpotential"]
  \inputleannode{InfoGeometry.Canonical.BogoliubovFockSuper.inducedChemicalPotential}
\end{theorem}

\subsection*{inducedChemicalPotential_eq_zero_of_vacuumTransported}

\begin{theorem}[blueprint="infogeometry:canonical:bogoliubovfocksuper:inducedchemicalpotential_eq_zero_of_vacuumtransported"]
  \inputleannode{InfoGeometry.Canonical.BogoliubovFockSuper.inducedChemicalPotential_eq_zero_of_vacuumTransported}
\end{theorem}

\subsection*{instDecidableEqSuperParity}

\begin{theorem}[blueprint="infogeometry:canonical:bogoliubovfocksuper:instdecidableeqsuperparity"]
  \inputleannode{InfoGeometry.Canonical.BogoliubovFockSuper.instDecidableEqSuperParity}
\end{theorem}

\subsection*{instReprSuperParity}

\begin{theorem}[blueprint="infogeometry:canonical:bogoliubovfocksuper:instreprsuperparity"]
  \inputleannode{InfoGeometry.Canonical.BogoliubovFockSuper.instReprSuperParity}
\end{theorem}

\subsection*{numberOperator}

\begin{theorem}[blueprint="infogeometry:canonical:bogoliubovfocksuper:numberoperator"]
  \inputleannode{InfoGeometry.Canonical.BogoliubovFockSuper.numberOperator}
\end{theorem}

\subsection*{paritySign}

\begin{theorem}[blueprint="infogeometry:canonical:bogoliubovfocksuper:paritysign"]
  \inputleannode{InfoGeometry.Canonical.BogoliubovFockSuper.paritySign}
\end{theorem}

\subsection*{superBracket}

\begin{theorem}[blueprint="infogeometry:canonical:bogoliubovfocksuper:superbracket"]
  \inputleannode{InfoGeometry.Canonical.BogoliubovFockSuper.superBracket}
\end{theorem}

\subsection*{superBracket_add_left}

\begin{theorem}[blueprint="infogeometry:canonical:bogoliubovfocksuper:superbracket_add_left"]
  \inputleannode{InfoGeometry.Canonical.BogoliubovFockSuper.superBracket_add_left}
\end{theorem}

\subsection*{superBracket_add_right}

\begin{theorem}[blueprint="infogeometry:canonical:bogoliubovfocksuper:superbracket_add_right"]
  \inputleannode{InfoGeometry.Canonical.BogoliubovFockSuper.superBracket_add_right}
\end{theorem}

\subsection*{superBracket_bogoliubov_covariance}

\begin{theorem}[blueprint="infogeometry:canonical:bogoliubovfocksuper:superbracket_bogoliubov_covariance"]
  \inputleannode{InfoGeometry.Canonical.BogoliubovFockSuper.superBracket_bogoliubov_covariance}
\end{theorem}

\subsection*{superBracket_even_left}

\begin{theorem}[blueprint="infogeometry:canonical:bogoliubovfocksuper:superbracket_even_left"]
  \inputleannode{InfoGeometry.Canonical.BogoliubovFockSuper.superBracket_even_left}
\end{theorem}

\subsection*{superBracket_odd_odd}

\begin{theorem}[blueprint="infogeometry:canonical:bogoliubovfocksuper:superbracket_odd_odd"]
  \inputleannode{InfoGeometry.Canonical.BogoliubovFockSuper.superBracket_odd_odd}
\end{theorem}

\subsection*{superBracket_smul_left}

\begin{theorem}[blueprint="infogeometry:canonical:bogoliubovfocksuper:superbracket_smul_left"]
  \inputleannode{InfoGeometry.Canonical.BogoliubovFockSuper.superBracket_smul_left}
\end{theorem}

\subsection*{superBracket_smul_right}

\begin{theorem}[blueprint="infogeometry:canonical:bogoliubovfocksuper:superbracket_smul_right"]
  \inputleannode{InfoGeometry.Canonical.BogoliubovFockSuper.superBracket_smul_right}
\end{theorem}

\section{InfoGeometry.Canonical.BogoliubovFockSuper.BogoliubovParams}

\subsection*{casesOn}

\begin{theorem}[blueprint="infogeometry:canonical:bogoliubovfocksuper:bogoliubovparams:caseson"]
  \inputleannode{InfoGeometry.Canonical.BogoliubovFockSuper.BogoliubovParams.casesOn}
\end{theorem}

\subsection*{ctorIdx}

\begin{theorem}[blueprint="infogeometry:canonical:bogoliubovfocksuper:bogoliubovparams:ctoridx"]
  \inputleannode{InfoGeometry.Canonical.BogoliubovFockSuper.BogoliubovParams.ctorIdx}
\end{theorem}

\subsection*{mk}

\begin{theorem}[blueprint="infogeometry:canonical:bogoliubovfocksuper:bogoliubovparams:mk"]
  \inputleannode{InfoGeometry.Canonical.BogoliubovFockSuper.BogoliubovParams.mk}
\end{theorem}

\subsection*{noConfusion}

\begin{theorem}[blueprint="infogeometry:canonical:bogoliubovfocksuper:bogoliubovparams:noconfusion"]
  \inputleannode{InfoGeometry.Canonical.BogoliubovFockSuper.BogoliubovParams.noConfusion}
\end{theorem}

\subsection*{noConfusionType}

\begin{theorem}[blueprint="infogeometry:canonical:bogoliubovfocksuper:bogoliubovparams:noconfusiontype"]
  \inputleannode{InfoGeometry.Canonical.BogoliubovFockSuper.BogoliubovParams.noConfusionType}
\end{theorem}

\subsection*{normalization}

\begin{theorem}[blueprint="infogeometry:canonical:bogoliubovfocksuper:bogoliubovparams:normalization"]
  \inputleannode{InfoGeometry.Canonical.BogoliubovFockSuper.BogoliubovParams.normalization}
\end{theorem}

\subsection*{rec}

\begin{theorem}[blueprint="infogeometry:canonical:bogoliubovfocksuper:bogoliubovparams:rec"]
  \inputleannode{InfoGeometry.Canonical.BogoliubovFockSuper.BogoliubovParams.rec}
\end{theorem}

\subsection*{recOn}

\begin{theorem}[blueprint="infogeometry:canonical:bogoliubovfocksuper:bogoliubovparams:recon"]
  \inputleannode{InfoGeometry.Canonical.BogoliubovFockSuper.BogoliubovParams.recOn}
\end{theorem}

\subsection*{u}

\begin{theorem}[blueprint="infogeometry:canonical:bogoliubovfocksuper:bogoliubovparams:u"]
  \inputleannode{InfoGeometry.Canonical.BogoliubovFockSuper.BogoliubovParams.u}
\end{theorem}

\subsection*{v}

\begin{theorem}[blueprint="infogeometry:canonical:bogoliubovfocksuper:bogoliubovparams:v"]
  \inputleannode{InfoGeometry.Canonical.BogoliubovFockSuper.BogoliubovParams.v}
\end{theorem}

\section{InfoGeometry.Canonical.BogoliubovFockSuper.CARWitness}

\subsection*{car_annihilation}

\begin{theorem}[blueprint="infogeometry:canonical:bogoliubovfocksuper:carwitness:car_annihilation"]
  \inputleannode{InfoGeometry.Canonical.BogoliubovFockSuper.CARWitness.car_annihilation}
\end{theorem}

\subsection*{car_creation}

\begin{theorem}[blueprint="infogeometry:canonical:bogoliubovfocksuper:carwitness:car_creation"]
  \inputleannode{InfoGeometry.Canonical.BogoliubovFockSuper.CARWitness.car_creation}
\end{theorem}

\subsection*{car_mixed}

\begin{theorem}[blueprint="infogeometry:canonical:bogoliubovfocksuper:carwitness:car_mixed"]
  \inputleannode{InfoGeometry.Canonical.BogoliubovFockSuper.CARWitness.car_mixed}
\end{theorem}

\subsection*{casesOn}

\begin{theorem}[blueprint="infogeometry:canonical:bogoliubovfocksuper:carwitness:caseson"]
  \inputleannode{InfoGeometry.Canonical.BogoliubovFockSuper.CARWitness.casesOn}
\end{theorem}

\subsection*{mk}

\begin{theorem}[blueprint="infogeometry:canonical:bogoliubovfocksuper:carwitness:mk"]
  \inputleannode{InfoGeometry.Canonical.BogoliubovFockSuper.CARWitness.mk}
\end{theorem}

\subsection*{rec}

\begin{theorem}[blueprint="infogeometry:canonical:bogoliubovfocksuper:carwitness:rec"]
  \inputleannode{InfoGeometry.Canonical.BogoliubovFockSuper.CARWitness.rec}
\end{theorem}

\subsection*{recOn}

\begin{theorem}[blueprint="infogeometry:canonical:bogoliubovfocksuper:carwitness:recon"]
  \inputleannode{InfoGeometry.Canonical.BogoliubovFockSuper.CARWitness.recOn}
\end{theorem}

\section{InfoGeometry.Canonical.BogoliubovFockSuper.CCRWitness}

\subsection*{casesOn}

\begin{theorem}[blueprint="infogeometry:canonical:bogoliubovfocksuper:ccrwitness:caseson"]
  \inputleannode{InfoGeometry.Canonical.BogoliubovFockSuper.CCRWitness.casesOn}
\end{theorem}

\subsection*{ccr_annihilation}

\begin{theorem}[blueprint="infogeometry:canonical:bogoliubovfocksuper:ccrwitness:ccr_annihilation"]
  \inputleannode{InfoGeometry.Canonical.BogoliubovFockSuper.CCRWitness.ccr_annihilation}
\end{theorem}

\subsection*{ccr_creation}

\begin{theorem}[blueprint="infogeometry:canonical:bogoliubovfocksuper:ccrwitness:ccr_creation"]
  \inputleannode{InfoGeometry.Canonical.BogoliubovFockSuper.CCRWitness.ccr_creation}
\end{theorem}

\subsection*{ccr_mixed}

\begin{theorem}[blueprint="infogeometry:canonical:bogoliubovfocksuper:ccrwitness:ccr_mixed"]
  \inputleannode{InfoGeometry.Canonical.BogoliubovFockSuper.CCRWitness.ccr_mixed}
\end{theorem}

\subsection*{mk}

\begin{theorem}[blueprint="infogeometry:canonical:bogoliubovfocksuper:ccrwitness:mk"]
  \inputleannode{InfoGeometry.Canonical.BogoliubovFockSuper.CCRWitness.mk}
\end{theorem}

\subsection*{rec}

\begin{theorem}[blueprint="infogeometry:canonical:bogoliubovfocksuper:ccrwitness:rec"]
  \inputleannode{InfoGeometry.Canonical.BogoliubovFockSuper.CCRWitness.rec}
\end{theorem}

\subsection*{recOn}

\begin{theorem}[blueprint="infogeometry:canonical:bogoliubovfocksuper:ccrwitness:recon"]
  \inputleannode{InfoGeometry.Canonical.BogoliubovFockSuper.CCRWitness.recOn}
\end{theorem}

\section{InfoGeometry.Canonical.BogoliubovFockSuper.SuperParity}

\subsection*{casesOn}

\begin{theorem}[blueprint="infogeometry:canonical:bogoliubovfocksuper:superparity:caseson"]
  \inputleannode{InfoGeometry.Canonical.BogoliubovFockSuper.SuperParity.casesOn}
\end{theorem}

\subsection*{ctorElim}

\begin{theorem}[blueprint="infogeometry:canonical:bogoliubovfocksuper:superparity:ctorelim"]
  \inputleannode{InfoGeometry.Canonical.BogoliubovFockSuper.SuperParity.ctorElim}
\end{theorem}

\subsection*{ctorElimType}

\begin{theorem}[blueprint="infogeometry:canonical:bogoliubovfocksuper:superparity:ctorelimtype"]
  \inputleannode{InfoGeometry.Canonical.BogoliubovFockSuper.SuperParity.ctorElimType}
\end{theorem}

\subsection*{ctorIdx}

\begin{theorem}[blueprint="infogeometry:canonical:bogoliubovfocksuper:superparity:ctoridx"]
  \inputleannode{InfoGeometry.Canonical.BogoliubovFockSuper.SuperParity.ctorIdx}
\end{theorem}

\subsection*{even}

\begin{theorem}[blueprint="infogeometry:canonical:bogoliubovfocksuper:superparity:even"]
  \inputleannode{InfoGeometry.Canonical.BogoliubovFockSuper.SuperParity.even}
\end{theorem}

\subsection*{noConfusion}

\begin{theorem}[blueprint="infogeometry:canonical:bogoliubovfocksuper:superparity:noconfusion"]
  \inputleannode{InfoGeometry.Canonical.BogoliubovFockSuper.SuperParity.noConfusion}
\end{theorem}

\subsection*{noConfusionType}

\begin{theorem}[blueprint="infogeometry:canonical:bogoliubovfocksuper:superparity:noconfusiontype"]
  \inputleannode{InfoGeometry.Canonical.BogoliubovFockSuper.SuperParity.noConfusionType}
\end{theorem}

\subsection*{odd}

\begin{theorem}[blueprint="infogeometry:canonical:bogoliubovfocksuper:superparity:odd"]
  \inputleannode{InfoGeometry.Canonical.BogoliubovFockSuper.SuperParity.odd}
\end{theorem}

\subsection*{ofNat}

\begin{theorem}[blueprint="infogeometry:canonical:bogoliubovfocksuper:superparity:ofnat"]
  \inputleannode{InfoGeometry.Canonical.BogoliubovFockSuper.SuperParity.ofNat}
\end{theorem}

\subsection*{ofNat_ctorIdx}

\begin{theorem}[blueprint="infogeometry:canonical:bogoliubovfocksuper:superparity:ofnat_ctoridx"]
  \inputleannode{InfoGeometry.Canonical.BogoliubovFockSuper.SuperParity.ofNat_ctorIdx}
\end{theorem}

\subsection*{rec}

\begin{theorem}[blueprint="infogeometry:canonical:bogoliubovfocksuper:superparity:rec"]
  \inputleannode{InfoGeometry.Canonical.BogoliubovFockSuper.SuperParity.rec}
\end{theorem}

\subsection*{recOn}

\begin{theorem}[blueprint="infogeometry:canonical:bogoliubovfocksuper:superparity:recon"]
  \inputleannode{InfoGeometry.Canonical.BogoliubovFockSuper.SuperParity.recOn}
\end{theorem}

\subsection*{toCtorIdx}

\begin{theorem}[blueprint="infogeometry:canonical:bogoliubovfocksuper:superparity:toctoridx"]
  \inputleannode{InfoGeometry.Canonical.BogoliubovFockSuper.SuperParity.toCtorIdx}
\end{theorem}

\section{InfoGeometry.Canonical.BogoliubovFockSuper.SuperParity.even}

\subsection*{elim}

\begin{theorem}[blueprint="infogeometry:canonical:bogoliubovfocksuper:superparity:even:elim"]
  \inputleannode{InfoGeometry.Canonical.BogoliubovFockSuper.SuperParity.even.elim}
\end{theorem}

\section{InfoGeometry.Canonical.BogoliubovFockSuper.SuperParity.odd}

\subsection*{elim}

\begin{theorem}[blueprint="infogeometry:canonical:bogoliubovfocksuper:superparity:odd:elim"]
  \inputleannode{InfoGeometry.Canonical.BogoliubovFockSuper.SuperParity.odd.elim}
\end{theorem}

\section{InfoGeometry.Canonical.BogoliubovFockSuper.bogoliubovAnnihilation}

\subsection*{eq_1}

\begin{theorem}[blueprint="infogeometry:canonical:bogoliubovfocksuper:bogoliubovannihilation:eq_1"]
  \inputleannode{InfoGeometry.Canonical.BogoliubovFockSuper.bogoliubovAnnihilation.eq_1}
\end{theorem}

\section{InfoGeometry.Canonical.BogoliubovFockSuper.bogoliubovCreation}

\subsection*{eq_1}

\begin{theorem}[blueprint="infogeometry:canonical:bogoliubovfocksuper:bogoliubovcreation:eq_1"]
  \inputleannode{InfoGeometry.Canonical.BogoliubovFockSuper.bogoliubovCreation.eq_1}
\end{theorem}

\section{InfoGeometry.Canonical.BogoliubovFockSuper.instReprSuperParity}

\subsection*{repr}

\begin{theorem}[blueprint="infogeometry:canonical:bogoliubovfocksuper:instreprsuperparity:repr"]
  \inputleannode{InfoGeometry.Canonical.BogoliubovFockSuper.instReprSuperParity.repr}
\end{theorem}

\section{InfoGeometry.Canonical.BogoliubovFockSuper.paritySign}

\subsection*{eq_1}

\begin{theorem}[blueprint="infogeometry:canonical:bogoliubovfocksuper:paritysign:eq_1"]
  \inputleannode{InfoGeometry.Canonical.BogoliubovFockSuper.paritySign.eq_1}
\end{theorem}

\subsection*{eq_2}

\begin{theorem}[blueprint="infogeometry:canonical:bogoliubovfocksuper:paritysign:eq_2"]
  \inputleannode{InfoGeometry.Canonical.BogoliubovFockSuper.paritySign.eq_2}
\end{theorem}

\section{InfoGeometry.Canonical.BogoliubovFockSuper.superBracket}

\subsection*{eq_1}

\begin{theorem}[blueprint="infogeometry:canonical:bogoliubovfocksuper:superbracket:eq_1"]
  \inputleannode{InfoGeometry.Canonical.BogoliubovFockSuper.superBracket.eq_1}
\end{theorem}

\section{InfoGeometry.Canonical.BottDirac}

\noindent\textbf{Buildable}: yes

\subsection*{Endomorphism}

\begin{theorem}[blueprint="infogeometry:canonical:bottdirac:endomorphism"]
  \inputleannode{InfoGeometry.Canonical.BottDirac.Endomorphism}
\end{theorem}

\subsection*{IsChiralDirac}

\begin{theorem}[blueprint="infogeometry:canonical:bottdirac:ischiraldirac"]
  \inputleannode{InfoGeometry.Canonical.BottDirac.IsChiralDirac}
\end{theorem}

\subsection*{IsInvolutiveGrading}

\begin{theorem}[blueprint="infogeometry:canonical:bottdirac:isinvolutivegrading"]
  \inputleannode{InfoGeometry.Canonical.BottDirac.IsInvolutiveGrading}
\end{theorem}

\subsection*{bottDirac}

\begin{theorem}[blueprint="infogeometry:canonical:bottdirac:bottdirac"]
  \inputleannode{InfoGeometry.Canonical.BottDirac.bottDirac}
\end{theorem}

\subsection*{bottDirac_apply_tmul}

\begin{theorem}[blueprint="infogeometry:canonical:bottdirac:bottdirac_apply_tmul"]
  \inputleannode{InfoGeometry.Canonical.BottDirac.bottDirac_apply_tmul}
\end{theorem}

\subsection*{bottDirac_sq_eq_sum_laplacians}

\begin{theorem}[blueprint="infogeometry:canonical:bottdirac:bottdirac_sq_eq_sum_laplacians"]
  \inputleannode{InfoGeometry.Canonical.BottDirac.bottDirac_sq_eq_sum_laplacians}
\end{theorem}

\subsection*{bottDirac_sq_eq_sum_laplacians_tmul}

\begin{theorem}[blueprint="infogeometry:canonical:bottdirac:bottdirac_sq_eq_sum_laplacians_tmul"]
  \inputleannode{InfoGeometry.Canonical.BottDirac.bottDirac_sq_eq_sum_laplacians_tmul}
\end{theorem}

\subsection*{cl11BottLaplacian}

\begin{theorem}[blueprint="infogeometry:canonical:bottdirac:cl11bottlaplacian"]
  \inputleannode{InfoGeometry.Canonical.BottDirac.cl11BottLaplacian}
\end{theorem}

\subsection*{cl11DiracSeed}

\begin{theorem}[blueprint="infogeometry:canonical:bottdirac:cl11diracseed"]
  \inputleannode{InfoGeometry.Canonical.BottDirac.cl11DiracSeed}
\end{theorem}

\subsection*{cl11Grading}

\begin{theorem}[blueprint="infogeometry:canonical:bottdirac:cl11grading"]
  \inputleannode{InfoGeometry.Canonical.BottDirac.cl11Grading}
\end{theorem}

\subsection*{cl11Grading_involutive}

\begin{theorem}[blueprint="infogeometry:canonical:bottdirac:cl11grading_involutive"]
  \inputleannode{InfoGeometry.Canonical.BottDirac.cl11Grading_involutive}
\end{theorem}

\subsection*{cl11_bottDirac_sq_eq_cl11BottLaplacian}

\begin{theorem}[blueprint="infogeometry:canonical:bottdirac:cl11_bottdirac_sq_eq_cl11bottlaplacian"]
  \inputleannode{InfoGeometry.Canonical.BottDirac.cl11_bottDirac_sq_eq_cl11BottLaplacian}
\end{theorem}

\subsection*{cl11_bottDirac_sq_eq_sum_laplacians}

\begin{theorem}[blueprint="infogeometry:canonical:bottdirac:cl11_bottdirac_sq_eq_sum_laplacians"]
  \inputleannode{InfoGeometry.Canonical.BottDirac.cl11_bottDirac_sq_eq_sum_laplacians}
\end{theorem}

\subsection*{cl11_isChiralDirac}

\begin{theorem}[blueprint="infogeometry:canonical:bottdirac:cl11_ischiraldirac"]
  \inputleannode{InfoGeometry.Canonical.BottDirac.cl11_isChiralDirac}
\end{theorem}

\subsection*{spectralDiracLinear}

\begin{theorem}[blueprint="infogeometry:canonical:bottdirac:spectraldiraclinear"]
  \inputleannode{InfoGeometry.Canonical.BottDirac.spectralDiracLinear}
\end{theorem}

\subsection*{spectralDiracLinear_apply}

\begin{theorem}[blueprint="infogeometry:canonical:bottdirac:spectraldiraclinear_apply"]
  \inputleannode{InfoGeometry.Canonical.BottDirac.spectralDiracLinear_apply}
\end{theorem}

\section{InfoGeometry.Canonical.BottDirac.IsChiralDirac}

\subsection*{eq_1}

\begin{theorem}[blueprint="infogeometry:canonical:bottdirac:ischiraldirac:eq_1"]
  \inputleannode{InfoGeometry.Canonical.BottDirac.IsChiralDirac.eq_1}
\end{theorem}

\section{InfoGeometry.Canonical.BottDirac.IsInvolutiveGrading}

\subsection*{eq_1}

\begin{theorem}[blueprint="infogeometry:canonical:bottdirac:isinvolutivegrading:eq_1"]
  \inputleannode{InfoGeometry.Canonical.BottDirac.IsInvolutiveGrading.eq_1}
\end{theorem}

\section{InfoGeometry.Canonical.BottDirac.bottDirac}

\subsection*{eq_1}

\begin{theorem}[blueprint="infogeometry:canonical:bottdirac:bottdirac:eq_1"]
  \inputleannode{InfoGeometry.Canonical.BottDirac.bottDirac.eq_1}
\end{theorem}

\section{InfoGeometry.Canonical.BottDirac.cl11BottLaplacian}

\subsection*{eq_1}

\begin{theorem}[blueprint="infogeometry:canonical:bottdirac:cl11bottlaplacian:eq_1"]
  \inputleannode{InfoGeometry.Canonical.BottDirac.cl11BottLaplacian.eq_1}
\end{theorem}

\section{InfoGeometry.Canonical.BottPeriodicity}

\noindent\textbf{Buildable}: yes

\subsection*{BottTensor}

\begin{theorem}[blueprint="infogeometry:canonical:bottperiodicity:botttensor"]
  \inputleannode{InfoGeometry.Canonical.BottPeriodicity.BottTensor}
\end{theorem}

\subsection*{CompanionTensor}

\begin{theorem}[blueprint="infogeometry:canonical:bottperiodicity:companiontensor"]
  \inputleannode{InfoGeometry.Canonical.BottPeriodicity.CompanionTensor}
\end{theorem}

\subsection*{DoubledTensor}

\begin{theorem}[blueprint="infogeometry:canonical:bottperiodicity:doubledtensor"]
  \inputleannode{InfoGeometry.Canonical.BottPeriodicity.DoubledTensor}
\end{theorem}

\subsection*{bottGeneratorInjection}

\begin{theorem}[blueprint="infogeometry:canonical:bottperiodicity:bottgeneratorinjection"]
  \inputleannode{InfoGeometry.Canonical.BottPeriodicity.bottGeneratorInjection}
\end{theorem}

\subsection*{bottGeneratorInjection_apply}

\begin{theorem}[blueprint="infogeometry:canonical:bottperiodicity:bottgeneratorinjection_apply"]
  \inputleannode{InfoGeometry.Canonical.BottPeriodicity.bottGeneratorInjection_apply}
\end{theorem}

\subsection*{bottStepEquiv}

\begin{theorem}[blueprint="infogeometry:canonical:bottperiodicity:bottstepequiv"]
  \inputleannode{InfoGeometry.Canonical.BottPeriodicity.bottStepEquiv}
\end{theorem}

\subsection*{bottStepEquiv_symm_left_generator}

\begin{theorem}[blueprint="infogeometry:canonical:bottperiodicity:bottstepequiv_symm_left_generator"]
  \inputleannode{InfoGeometry.Canonical.BottPeriodicity.bottStepEquiv_symm_left_generator}
\end{theorem}

\subsection*{bottStepEquiv_symm_right_generator}

\begin{theorem}[blueprint="infogeometry:canonical:bottperiodicity:bottstepequiv_symm_right_generator"]
  \inputleannode{InfoGeometry.Canonical.BottPeriodicity.bottStepEquiv_symm_right_generator}
\end{theorem}

\subsection*{bott_step_periodicity}

\begin{theorem}[blueprint="infogeometry:canonical:bottperiodicity:bott_step_periodicity"]
  \inputleannode{InfoGeometry.Canonical.BottPeriodicity.bott_step_periodicity}
\end{theorem}

\subsection*{cl11GradeEpsilon}

\begin{theorem}[blueprint="infogeometry:canonical:bottperiodicity:cl11gradeepsilon"]
  \inputleannode{InfoGeometry.Canonical.BottPeriodicity.cl11GradeEpsilon}
\end{theorem}

\subsection*{gradedProdInjection}

\begin{theorem}[blueprint="infogeometry:canonical:bottperiodicity:gradedprodinjection"]
  \inputleannode{InfoGeometry.Canonical.BottPeriodicity.gradedProdInjection}
\end{theorem}

\subsection*{gradedProdInjection_matches_pattern_unit}

\begin{theorem}[blueprint="infogeometry:canonical:bottperiodicity:gradedprodinjection_matches_pattern_unit"]
  \inputleannode{InfoGeometry.Canonical.BottPeriodicity.gradedProdInjection_matches_pattern_unit}
\end{theorem}

\subsection*{gradedProdInjection_on_generator}

\begin{theorem}[blueprint="infogeometry:canonical:bottperiodicity:gradedprodinjection_on_generator"]
  \inputleannode{InfoGeometry.Canonical.BottPeriodicity.gradedProdInjection_on_generator}
\end{theorem}

\subsection*{newGeneratorTensorOne}

\begin{theorem}[blueprint="infogeometry:canonical:bottperiodicity:newgeneratortensorone"]
  \inputleannode{InfoGeometry.Canonical.BottPeriodicity.newGeneratorTensorOne}
\end{theorem}

\subsection*{oldGeneratorTensorEpsilon}

\begin{theorem}[blueprint="infogeometry:canonical:bottperiodicity:oldgeneratortensorepsilon"]
  \inputleannode{InfoGeometry.Canonical.BottPeriodicity.oldGeneratorTensorEpsilon}
\end{theorem}

\subsection*{tensorModularJ}

\begin{theorem}[blueprint="infogeometry:canonical:bottperiodicity:tensormodularj"]
  \inputleannode{InfoGeometry.Canonical.BottPeriodicity.tensorModularJ}
\end{theorem}

\subsection*{tensorModularJ_comp_tmul}

\begin{theorem}[blueprint="infogeometry:canonical:bottperiodicity:tensormodularj_comp_tmul"]
  \inputleannode{InfoGeometry.Canonical.BottPeriodicity.tensorModularJ_comp_tmul}
\end{theorem}

\subsection*{tensorModularJ_tmul}

\begin{theorem}[blueprint="infogeometry:canonical:bottperiodicity:tensormodularj_tmul"]
  \inputleannode{InfoGeometry.Canonical.BottPeriodicity.tensorModularJ_tmul}
\end{theorem}

\subsection*{tensorSpectralEpsilon}

\begin{theorem}[blueprint="infogeometry:canonical:bottperiodicity:tensorspectralepsilon"]
  \inputleannode{InfoGeometry.Canonical.BottPeriodicity.tensorSpectralEpsilon}
\end{theorem}

\subsection*{tensorSpectralEpsilon_comp_tmul}

\begin{theorem}[blueprint="infogeometry:canonical:bottperiodicity:tensorspectralepsilon_comp_tmul"]
  \inputleannode{InfoGeometry.Canonical.BottPeriodicity.tensorSpectralEpsilon_comp_tmul}
\end{theorem}

\subsection*{tensorSpectralEpsilon_tmul}

\begin{theorem}[blueprint="infogeometry:canonical:bottperiodicity:tensorspectralepsilon_tmul"]
  \inputleannode{InfoGeometry.Canonical.BottPeriodicity.tensorSpectralEpsilon_tmul}
\end{theorem}

\subsection*{twistedGeneratorPattern}

\begin{theorem}[blueprint="infogeometry:canonical:bottperiodicity:twistedgeneratorpattern"]
  \inputleannode{InfoGeometry.Canonical.BottPeriodicity.twistedGeneratorPattern}
\end{theorem}

\subsection*{twistedGeneratorPattern_apply}

\begin{theorem}[blueprint="infogeometry:canonical:bottperiodicity:twistedgeneratorpattern_apply"]
  \inputleannode{InfoGeometry.Canonical.BottPeriodicity.twistedGeneratorPattern_apply}
\end{theorem}

\section{InfoGeometry.Canonical.BottPeriodicity.newGeneratorTensorOne}

\subsection*{eq_1}

\begin{theorem}[blueprint="infogeometry:canonical:bottperiodicity:newgeneratortensorone:eq_1"]
  \inputleannode{InfoGeometry.Canonical.BottPeriodicity.newGeneratorTensorOne.eq_1}
\end{theorem}

\section{InfoGeometry.Canonical.BottPeriodicity.oldGeneratorTensorEpsilon}

\subsection*{eq_1}

\begin{theorem}[blueprint="infogeometry:canonical:bottperiodicity:oldgeneratortensorepsilon:eq_1"]
  \inputleannode{InfoGeometry.Canonical.BottPeriodicity.oldGeneratorTensorEpsilon.eq_1}
\end{theorem}

\section{InfoGeometry.Canonical.BottPeriodicity.tensorModularJ}

\subsection*{eq_1}

\begin{theorem}[blueprint="infogeometry:canonical:bottperiodicity:tensormodularj:eq_1"]
  \inputleannode{InfoGeometry.Canonical.BottPeriodicity.tensorModularJ.eq_1}
\end{theorem}

\section{InfoGeometry.Canonical.BottPeriodicity.tensorSpectralEpsilon}

\subsection*{eq_1}

\begin{theorem}[blueprint="infogeometry:canonical:bottperiodicity:tensorspectralepsilon:eq_1"]
  \inputleannode{InfoGeometry.Canonical.BottPeriodicity.tensorSpectralEpsilon.eq_1}
\end{theorem}

\section{InfoGeometry.Canonical.BottPeriodicity.twistedGeneratorPattern}

\subsection*{eq_1}

\begin{theorem}[blueprint="infogeometry:canonical:bottperiodicity:twistedgeneratorpattern:eq_1"]
  \inputleannode{InfoGeometry.Canonical.BottPeriodicity.twistedGeneratorPattern.eq_1}
\end{theorem}

\section{InfoGeometry.Canonical.Bregman}

\subsection*{pythagorean_law}

\begin{theorem}[blueprint="infogeometry:canonical:bregman:pythagorean_law"]
  \inputleannode{InfoGeometry.Canonical.Bregman.pythagorean_law}
\end{theorem}

\section{InfoGeometry.Canonical.BregmanTriality}

\noindent\textbf{Buildable}: yes

\subsection*{softmaxBregmanAttention}

\begin{theorem}[blueprint="infogeometry:canonical:bregmantriality:softmaxbregmanattention"]
  \inputleannode{InfoGeometry.Canonical.BregmanTriality.softmaxBregmanAttention}
\end{theorem}

\section{InfoGeometry.Canonical.CalabiYauBridge}

\noindent\textbf{Buildable}: yes

\subsection*{CalabiYauRicciState}

\begin{theorem}[blueprint="infogeometry:canonical:calabiyaubridge:calabiyauriccistate"]
  \inputleannode{InfoGeometry.Canonical.CalabiYauBridge.CalabiYauRicciState}
\end{theorem}

\subsection*{CalabiYauSpinorialState}

\begin{theorem}[blueprint="infogeometry:canonical:calabiyaubridge:calabiyauspinorialstate"]
  \inputleannode{InfoGeometry.Canonical.CalabiYauBridge.CalabiYauSpinorialState}
\end{theorem}

\subsection*{HasConstantMongeAmpereDensity}

\begin{theorem}[blueprint="infogeometry:canonical:calabiyaubridge:hasconstantmongeamperedensity"]
  \inputleannode{InfoGeometry.Canonical.CalabiYauBridge.HasConstantMongeAmpereDensity}
\end{theorem}

\subsection*{IsRicciFlat}

\begin{theorem}[blueprint="infogeometry:canonical:calabiyaubridge:isricciflat"]
  \inputleannode{InfoGeometry.Canonical.CalabiYauBridge.IsRicciFlat}
\end{theorem}

\subsection*{MongeAmpereRicciClosure}

\begin{theorem}[blueprint="infogeometry:canonical:calabiyaubridge:mongeamperericciclosure"]
  \inputleannode{InfoGeometry.Canonical.CalabiYauBridge.MongeAmpereRicciClosure}
\end{theorem}

\subsection*{MongeAmpereSpinorialClosure}

\begin{theorem}[blueprint="infogeometry:canonical:calabiyaubridge:mongeamperespinorialclosure"]
  \inputleannode{InfoGeometry.Canonical.CalabiYauBridge.MongeAmpereSpinorialClosure}
\end{theorem}

\subsection*{W_constant_of_constantMongeAmpere}

\begin{theorem}[blueprint="infogeometry:canonical:calabiyaubridge:w_constant_of_constantmongeampere"]
  \inputleannode{InfoGeometry.Canonical.CalabiYauBridge.W_constant_of_constantMongeAmpere}
\end{theorem}

\subsection*{W_constant_of_spinorialState}

\begin{theorem}[blueprint="infogeometry:canonical:calabiyaubridge:w_constant_of_spinorialstate"]
  \inputleannode{InfoGeometry.Canonical.CalabiYauBridge.W_constant_of_spinorialState}
\end{theorem}

\subsection*{constantMongeAmpereImpliesRicciFlat_of_isRicciFlat}

\begin{theorem}[blueprint="infogeometry:canonical:calabiyaubridge:constantmongeampereimpliesricciflat_of_isricciflat"]
  \inputleannode{InfoGeometry.Canonical.CalabiYauBridge.constantMongeAmpereImpliesRicciFlat_of_isRicciFlat}
\end{theorem}

\subsection*{constantMongeAmpereImpliesZeroSpinorial_of_spinorialState}

\begin{theorem}[blueprint="infogeometry:canonical:calabiyaubridge:constantmongeampereimplieszerospinorial_of_spinorialstate"]
  \inputleannode{InfoGeometry.Canonical.CalabiYauBridge.constantMongeAmpereImpliesZeroSpinorial_of_spinorialState}
\end{theorem}

\subsection*{hasConstantMongeAmpereDensity_iff}

\begin{theorem}[blueprint="infogeometry:canonical:calabiyaubridge:hasconstantmongeamperedensity_iff"]
  \inputleannode{InfoGeometry.Canonical.CalabiYauBridge.hasConstantMongeAmpereDensity_iff}
\end{theorem}

\subsection*{hasConstantMongeAmpereDensity_of_satisfiesMongeAmpere_const}

\begin{theorem}[blueprint="infogeometry:canonical:calabiyaubridge:hasconstantmongeamperedensity_of_satisfiesmongeampere_const"]
  \inputleannode{InfoGeometry.Canonical.CalabiYauBridge.hasConstantMongeAmpereDensity_of_satisfiesMongeAmpere_const}
\end{theorem}

\subsection*{isEinsteinKaehlerAtWith_zero_of_isRicciFlat}

\begin{theorem}[blueprint="infogeometry:canonical:calabiyaubridge:iseinsteinkaehleratwith_zero_of_isricciflat"]
  \inputleannode{InfoGeometry.Canonical.CalabiYauBridge.isEinsteinKaehlerAtWith_zero_of_isRicciFlat}
\end{theorem}

\subsection*{isRicciFlat_of_constantMongeAmpere}

\begin{theorem}[blueprint="infogeometry:canonical:calabiyaubridge:isricciflat_of_constantmongeampere"]
  \inputleannode{InfoGeometry.Canonical.CalabiYauBridge.isRicciFlat_of_constantMongeAmpere}
\end{theorem}

\subsection*{ricciTensor_eq_of_isRicciFlat}

\begin{theorem}[blueprint="infogeometry:canonical:calabiyaubridge:riccitensor_eq_of_isricciflat"]
  \inputleannode{InfoGeometry.Canonical.CalabiYauBridge.ricciTensor_eq_of_isRicciFlat}
\end{theorem}

\subsection*{ricciTensor_unique_of_constantMongeAmpere}

\begin{theorem}[blueprint="infogeometry:canonical:calabiyaubridge:riccitensor_unique_of_constantmongeampere"]
  \inputleannode{InfoGeometry.Canonical.CalabiYauBridge.ricciTensor_unique_of_constantMongeAmpere}
\end{theorem}

\subsection*{satisfiesMongeAmpere_const_of_hasConstantMongeAmpereDensity}

\begin{theorem}[blueprint="infogeometry:canonical:calabiyaubridge:satisfiesmongeampere_const_of_hasconstantmongeamperedensity"]
  \inputleannode{InfoGeometry.Canonical.CalabiYauBridge.satisfiesMongeAmpere_const_of_hasConstantMongeAmpereDensity}
\end{theorem}

\subsection*{spinorialScalarCurvature_eq_zero_of_constantMongeAmpere}

\begin{theorem}[blueprint="infogeometry:canonical:calabiyaubridge:spinorialscalarcurvature_eq_zero_of_constantmongeampere"]
  \inputleannode{InfoGeometry.Canonical.CalabiYauBridge.spinorialScalarCurvature_eq_zero_of_constantMongeAmpere}
\end{theorem}

\subsection*{vacuumEinsteinEquation_of_constantMongeAmpere}

\begin{theorem}[blueprint="infogeometry:canonical:calabiyaubridge:vacuumeinsteinequation_of_constantmongeampere"]
  \inputleannode{InfoGeometry.Canonical.CalabiYauBridge.vacuumEinsteinEquation_of_constantMongeAmpere}
\end{theorem}

\subsection*{vacuumEinsteinEquation_of_isRicciFlat}

\begin{theorem}[blueprint="infogeometry:canonical:calabiyaubridge:vacuumeinsteinequation_of_isricciflat"]
  \inputleannode{InfoGeometry.Canonical.CalabiYauBridge.vacuumEinsteinEquation_of_isRicciFlat}
\end{theorem}

\section{InfoGeometry.Canonical.CartanDecomposition}

\noindent\textbf{Buildable}: yes

\subsection*{InformationCartanTriple}

\begin{theorem}[blueprint="infogeometry:canonical:cartandecomposition:informationcartantriple"]
  \inputleannode{InfoGeometry.Canonical.CartanDecomposition.InformationCartanTriple}
\end{theorem}

\section{InfoGeometry.Canonical.CartanDecomposition.InformationCartanTriple}

\subsection*{A}

\begin{theorem}[blueprint="infogeometry:canonical:cartandecomposition:informationcartantriple:a"]
  \inputleannode{InfoGeometry.Canonical.CartanDecomposition.InformationCartanTriple.A}
\end{theorem}

\subsection*{A_D}

\begin{theorem}[blueprint="infogeometry:canonical:cartandecomposition:informationcartantriple:a_d"]
  \inputleannode{InfoGeometry.Canonical.CartanDecomposition.InformationCartanTriple.A_D}
\end{theorem}

\subsection*{A_MP}

\begin{theorem}[blueprint="infogeometry:canonical:cartandecomposition:informationcartantriple:a_mp"]
  \inputleannode{InfoGeometry.Canonical.CartanDecomposition.InformationCartanTriple.A_MP}
\end{theorem}

\subsection*{GammaG}

\begin{theorem}[blueprint="infogeometry:canonical:cartandecomposition:informationcartantriple:gammag"]
  \inputleannode{InfoGeometry.Canonical.CartanDecomposition.InformationCartanTriple.GammaG}
\end{theorem}

\subsection*{GammaS}

\begin{theorem}[blueprint="infogeometry:canonical:cartandecomposition:informationcartantriple:gammas"]
  \inputleannode{InfoGeometry.Canonical.CartanDecomposition.InformationCartanTriple.GammaS}
\end{theorem}

\subsection*{IsCompact}

\begin{theorem}[blueprint="infogeometry:canonical:cartandecomposition:informationcartantriple:iscompact"]
  \inputleannode{InfoGeometry.Canonical.CartanDecomposition.InformationCartanTriple.IsCompact}
\end{theorem}

\subsection*{IsNonCompact}

\begin{theorem}[blueprint="infogeometry:canonical:cartandecomposition:informationcartantriple:isnoncompact"]
  \inputleannode{InfoGeometry.Canonical.CartanDecomposition.InformationCartanTriple.IsNonCompact}
\end{theorem}

\subsection*{casesOn}

\begin{theorem}[blueprint="infogeometry:canonical:cartandecomposition:informationcartantriple:caseson"]
  \inputleannode{InfoGeometry.Canonical.CartanDecomposition.InformationCartanTriple.casesOn}
\end{theorem}

\subsection*{ctorIdx}

\begin{theorem}[blueprint="infogeometry:canonical:cartandecomposition:informationcartantriple:ctoridx"]
  \inputleannode{InfoGeometry.Canonical.CartanDecomposition.InformationCartanTriple.ctorIdx}
\end{theorem}

\subsection*{grading_commutator_anomaly}

\begin{theorem}[blueprint="infogeometry:canonical:cartandecomposition:informationcartantriple:grading_commutator_anomaly"]
  \inputleannode{InfoGeometry.Canonical.CartanDecomposition.InformationCartanTriple.grading_commutator_anomaly}
\end{theorem}

\subsection*{h_drazin}

\begin{theorem}[blueprint="infogeometry:canonical:cartandecomposition:informationcartantriple:h_drazin"]
  \inputleannode{InfoGeometry.Canonical.CartanDecomposition.InformationCartanTriple.h_drazin}
\end{theorem}

\subsection*{h_penrose}

\begin{theorem}[blueprint="infogeometry:canonical:cartandecomposition:informationcartantriple:h_penrose"]
  \inputleannode{InfoGeometry.Canonical.CartanDecomposition.InformationCartanTriple.h_penrose}
\end{theorem}

\subsection*{mk}

\begin{theorem}[blueprint="infogeometry:canonical:cartandecomposition:informationcartantriple:mk"]
  \inputleannode{InfoGeometry.Canonical.CartanDecomposition.InformationCartanTriple.mk}
\end{theorem}

\subsection*{noConfusion}

\begin{theorem}[blueprint="infogeometry:canonical:cartandecomposition:informationcartantriple:noconfusion"]
  \inputleannode{InfoGeometry.Canonical.CartanDecomposition.InformationCartanTriple.noConfusion}
\end{theorem}

\subsection*{noConfusionType}

\begin{theorem}[blueprint="infogeometry:canonical:cartandecomposition:informationcartantriple:noconfusiontype"]
  \inputleannode{InfoGeometry.Canonical.CartanDecomposition.InformationCartanTriple.noConfusionType}
\end{theorem}

\subsection*{rec}

\begin{theorem}[blueprint="infogeometry:canonical:cartandecomposition:informationcartantriple:rec"]
  \inputleannode{InfoGeometry.Canonical.CartanDecomposition.InformationCartanTriple.rec}
\end{theorem}

\subsection*{recOn}

\begin{theorem}[blueprint="infogeometry:canonical:cartandecomposition:informationcartantriple:recon"]
  \inputleannode{InfoGeometry.Canonical.CartanDecomposition.InformationCartanTriple.recOn}
\end{theorem}

\subsection*{spectral_proj_is_compact_of_normal}

\begin{theorem}[blueprint="infogeometry:canonical:cartandecomposition:informationcartantriple:spectral_proj_is_compact_of_normal"]
  \inputleannode{InfoGeometry.Canonical.CartanDecomposition.InformationCartanTriple.spectral_proj_is_compact_of_normal}
\end{theorem}

\section{InfoGeometry.Canonical.CartanDecomposition.InformationCartanTriple.IsCompact}

\subsection*{eq_1}

\begin{theorem}[blueprint="infogeometry:canonical:cartandecomposition:informationcartantriple:iscompact:eq_1"]
  \inputleannode{InfoGeometry.Canonical.CartanDecomposition.InformationCartanTriple.IsCompact.eq_1}
\end{theorem}

\section{InfoGeometry.Canonical.Cayley}

\subsection*{Bridge}

\begin{theorem}[blueprint="infogeometry:canonical:cayley:bridge"]
  \inputleannode{InfoGeometry.Canonical.Cayley.Bridge}
\end{theorem}

\subsection*{CayleyBridge}

\begin{theorem}[blueprint="infogeometry:canonical:cayley:cayleybridge"]
  \inputleannode{InfoGeometry.Canonical.Cayley.CayleyBridge}
\end{theorem}

\subsection*{CayleyCompatibleDualFlat}

\begin{theorem}[blueprint="infogeometry:canonical:cayley:cayleycompatibledualflat"]
  \inputleannode{InfoGeometry.Canonical.Cayley.CayleyCompatibleDualFlat}
\end{theorem}

\subsection*{CayleyDualFlatCompatibility}

\begin{theorem}[blueprint="infogeometry:canonical:cayley:cayleydualflatcompatibility"]
  \inputleannode{InfoGeometry.Canonical.Cayley.CayleyDualFlatCompatibility}
\end{theorem}

\subsection*{CayleyEquivalence}

\begin{theorem}[blueprint="infogeometry:canonical:cayley:cayleyequivalence"]
  \inputleannode{InfoGeometry.Canonical.Cayley.CayleyEquivalence}
\end{theorem}

\subsection*{CompatibleDualFlat}

\begin{theorem}[blueprint="infogeometry:canonical:cayley:compatibledualflat"]
  \inputleannode{InfoGeometry.Canonical.Cayley.CompatibleDualFlat}
\end{theorem}

\subsection*{cayleyIdentityBridge}

\begin{theorem}[blueprint="infogeometry:canonical:cayley:cayleyidentitybridge"]
  \inputleannode{InfoGeometry.Canonical.Cayley.cayleyIdentityBridge}
\end{theorem}

\subsection*{cayleyIdentityCompatibleGeometry}

\begin{theorem}[blueprint="infogeometry:canonical:cayley:cayleyidentitycompatiblegeometry"]
  \inputleannode{InfoGeometry.Canonical.Cayley.cayleyIdentityCompatibleGeometry}
\end{theorem}

\subsection*{cayleyPythagoreanInvariance}

\begin{theorem}[blueprint="infogeometry:canonical:cayley:cayleypythagoreaninvariance"]
  \inputleannode{InfoGeometry.Canonical.Cayley.cayleyPythagoreanInvariance}
\end{theorem}

\subsection*{cayley_pythagorean_invariance}

\begin{theorem}[blueprint="infogeometry:canonical:cayley:cayley_pythagorean_invariance"]
  \inputleannode{InfoGeometry.Canonical.Cayley.cayley_pythagorean_invariance}
\end{theorem}

\subsection*{identityBridge}

\begin{theorem}[blueprint="infogeometry:canonical:cayley:identitybridge"]
  \inputleannode{InfoGeometry.Canonical.Cayley.identityBridge}
\end{theorem}

\subsection*{identityCompatibleGeometry}

\begin{theorem}[blueprint="infogeometry:canonical:cayley:identitycompatiblegeometry"]
  \inputleannode{InfoGeometry.Canonical.Cayley.identityCompatibleGeometry}
\end{theorem}

\section{InfoGeometry.Canonical.Cayley.Bridge}

\subsection*{casesOn}

\begin{theorem}[blueprint="infogeometry:canonical:cayley:bridge:caseson"]
  \inputleannode{InfoGeometry.Canonical.Cayley.Bridge.casesOn}
\end{theorem}

\subsection*{ctorIdx}

\begin{theorem}[blueprint="infogeometry:canonical:cayley:bridge:ctoridx"]
  \inputleannode{InfoGeometry.Canonical.Cayley.Bridge.ctorIdx}
\end{theorem}

\subsection*{left_inv}

\begin{theorem}[blueprint="infogeometry:canonical:cayley:bridge:left_inv"]
  \inputleannode{InfoGeometry.Canonical.Cayley.Bridge.left_inv}
\end{theorem}

\subsection*{mk}

\begin{theorem}[blueprint="infogeometry:canonical:cayley:bridge:mk"]
  \inputleannode{InfoGeometry.Canonical.Cayley.Bridge.mk}
\end{theorem}

\subsection*{noConfusion}

\begin{theorem}[blueprint="infogeometry:canonical:cayley:bridge:noconfusion"]
  \inputleannode{InfoGeometry.Canonical.Cayley.Bridge.noConfusion}
\end{theorem}

\subsection*{noConfusionType}

\begin{theorem}[blueprint="infogeometry:canonical:cayley:bridge:noconfusiontype"]
  \inputleannode{InfoGeometry.Canonical.Cayley.Bridge.noConfusionType}
\end{theorem}

\subsection*{rec}

\begin{theorem}[blueprint="infogeometry:canonical:cayley:bridge:rec"]
  \inputleannode{InfoGeometry.Canonical.Cayley.Bridge.rec}
\end{theorem}

\subsection*{recOn}

\begin{theorem}[blueprint="infogeometry:canonical:cayley:bridge:recon"]
  \inputleannode{InfoGeometry.Canonical.Cayley.Bridge.recOn}
\end{theorem}

\subsection*{right_inv}

\begin{theorem}[blueprint="infogeometry:canonical:cayley:bridge:right_inv"]
  \inputleannode{InfoGeometry.Canonical.Cayley.Bridge.right_inv}
\end{theorem}

\subsection*{toBounded}

\begin{theorem}[blueprint="infogeometry:canonical:cayley:bridge:tobounded"]
  \inputleannode{InfoGeometry.Canonical.Cayley.Bridge.toBounded}
\end{theorem}

\subsection*{toUnbounded}

\begin{theorem}[blueprint="infogeometry:canonical:cayley:bridge:tounbounded"]
  \inputleannode{InfoGeometry.Canonical.Cayley.Bridge.toUnbounded}
\end{theorem}

\section{InfoGeometry.Canonical.Cayley.CompatibleDualFlat}

\subsection*{D_transport}

\begin{theorem}[blueprint="infogeometry:canonical:cayley:compatibledualflat:d_transport"]
  \inputleannode{InfoGeometry.Canonical.Cayley.CompatibleDualFlat.D_transport}
\end{theorem}

\subsection*{SB}

\begin{theorem}[blueprint="infogeometry:canonical:cayley:compatibledualflat:sb"]
  \inputleannode{InfoGeometry.Canonical.Cayley.CompatibleDualFlat.SB}
\end{theorem}

\subsection*{casesOn}

\begin{theorem}[blueprint="infogeometry:canonical:cayley:compatibledualflat:caseson"]
  \inputleannode{InfoGeometry.Canonical.Cayley.CompatibleDualFlat.casesOn}
\end{theorem}

\subsection*{ctorIdx}

\begin{theorem}[blueprint="infogeometry:canonical:cayley:compatibledualflat:ctoridx"]
  \inputleannode{InfoGeometry.Canonical.Cayley.CompatibleDualFlat.ctorIdx}
\end{theorem}

\subsection*{grad_transport}

\begin{theorem}[blueprint="infogeometry:canonical:cayley:compatibledualflat:grad_transport"]
  \inputleannode{InfoGeometry.Canonical.Cayley.CompatibleDualFlat.grad_transport}
\end{theorem}

\subsection*{mk}

\begin{theorem}[blueprint="infogeometry:canonical:cayley:compatibledualflat:mk"]
  \inputleannode{InfoGeometry.Canonical.Cayley.CompatibleDualFlat.mk}
\end{theorem}

\subsection*{noConfusion}

\begin{theorem}[blueprint="infogeometry:canonical:cayley:compatibledualflat:noconfusion"]
  \inputleannode{InfoGeometry.Canonical.Cayley.CompatibleDualFlat.noConfusion}
\end{theorem}

\subsection*{noConfusionType}

\begin{theorem}[blueprint="infogeometry:canonical:cayley:compatibledualflat:noconfusiontype"]
  \inputleannode{InfoGeometry.Canonical.Cayley.CompatibleDualFlat.noConfusionType}
\end{theorem}

\subsection*{rec}

\begin{theorem}[blueprint="infogeometry:canonical:cayley:compatibledualflat:rec"]
  \inputleannode{InfoGeometry.Canonical.Cayley.CompatibleDualFlat.rec}
\end{theorem}

\subsection*{recOn}

\begin{theorem}[blueprint="infogeometry:canonical:cayley:compatibledualflat:recon"]
  \inputleannode{InfoGeometry.Canonical.Cayley.CompatibleDualFlat.recOn}
\end{theorem}

\section{InfoGeometry.Canonical.CayleyBregmanBridge}

\noindent\textbf{Buildable}: yes

\noindent\emph{No exported declarations in the current declaration graph.}

\section{InfoGeometry.Canonical.ChiralAction}

\noindent\textbf{Buildable}: yes

\subsection*{chiralDirac}

\begin{theorem}[blueprint="infogeometry:canonical:chiralaction:chiraldirac"]
  \inputleannode{InfoGeometry.Canonical.ChiralAction.chiralDirac}
\end{theorem}

\subsection*{chiralInformationAction}

\begin{theorem}[blueprint="infogeometry:canonical:chiralaction:chiralinformationaction"]
  \inputleannode{InfoGeometry.Canonical.ChiralAction.chiralInformationAction}
\end{theorem}

\subsection*{chiral_action_reduces_for_normal}

\begin{theorem}[blueprint="infogeometry:canonical:chiralaction:chiral_action_reduces_for_normal"]
  \inputleannode{InfoGeometry.Canonical.ChiralAction.chiral_action_reduces_for_normal}
\end{theorem}

\section{InfoGeometry.Canonical.ChiralAnomaly}

\noindent\textbf{Buildable}: yes

\subsection*{DoublyStochasticSinkhornTrajectory}

\begin{theorem}[blueprint="infogeometry:canonical:chiralanomaly:doublystochasticsinkhorntrajectory"]
  \inputleannode{InfoGeometry.Canonical.ChiralAnomaly.DoublyStochasticSinkhornTrajectory}
\end{theorem}

\subsection*{epsilon}

\begin{theorem}[blueprint="infogeometry:canonical:chiralanomaly:epsilon"]
  \inputleannode{InfoGeometry.Canonical.ChiralAnomaly.epsilon}
\end{theorem}

\subsection*{exists_routingEpsilon_of_bistochastic}

\begin{theorem}[blueprint="infogeometry:canonical:chiralanomaly:exists_routingepsilon_of_bistochastic"]
  \inputleannode{InfoGeometry.Canonical.ChiralAnomaly.exists_routingEpsilon_of_bistochastic}
\end{theorem}

\subsection*{exists_routingEpsilon_of_mem_doublyStochastic}

\begin{theorem}[blueprint="infogeometry:canonical:chiralanomaly:exists_routingepsilon_of_mem_doublystochastic"]
  \inputleannode{InfoGeometry.Canonical.ChiralAnomaly.exists_routingEpsilon_of_mem_doublyStochastic}
\end{theorem}

\subsection*{inverseCommutator}

\begin{theorem}[blueprint="infogeometry:canonical:chiralanomaly:inversecommutator"]
  \inputleannode{InfoGeometry.Canonical.ChiralAnomaly.inverseCommutator}
\end{theorem}

\subsection*{normal_inverse_anomaly_vanishes}

\begin{theorem}[blueprint="infogeometry:canonical:chiralanomaly:normal_inverse_anomaly_vanishes"]
  \inputleannode{InfoGeometry.Canonical.ChiralAnomaly.normal_inverse_anomaly_vanishes}
\end{theorem}

\subsection*{routingChiralAnomaly}

\begin{theorem}[blueprint="infogeometry:canonical:chiralanomaly:routingchiralanomaly"]
  \inputleannode{InfoGeometry.Canonical.ChiralAnomaly.routingChiralAnomaly}
\end{theorem}

\subsection*{routingChiralAnomaly_eq_semantic_coord_gap}

\begin{theorem}[blueprint="infogeometry:canonical:chiralanomaly:routingchiralanomaly_eq_semantic_coord_gap"]
  \inputleannode{InfoGeometry.Canonical.ChiralAnomaly.routingChiralAnomaly_eq_semantic_coord_gap}
\end{theorem}

\subsection*{routingEpsilon}

\begin{theorem}[blueprint="infogeometry:canonical:chiralanomaly:routingepsilon"]
  \inputleannode{InfoGeometry.Canonical.ChiralAnomaly.routingEpsilon}
\end{theorem}

\subsection*{routingEpsilon_eq_semantic_gap_abs}

\begin{theorem}[blueprint="infogeometry:canonical:chiralanomaly:routingepsilon_eq_semantic_gap_abs"]
  \inputleannode{InfoGeometry.Canonical.ChiralAnomaly.routingEpsilon_eq_semantic_gap_abs}
\end{theorem}

\subsection*{routingEpsilon_le_one_of_simplex}

\begin{theorem}[blueprint="infogeometry:canonical:chiralanomaly:routingepsilon_le_one_of_simplex"]
  \inputleannode{InfoGeometry.Canonical.ChiralAnomaly.routingEpsilon_le_one_of_simplex}
\end{theorem}

\subsection*{routingEpsilon_nonneg}

\begin{theorem}[blueprint="infogeometry:canonical:chiralanomaly:routingepsilon_nonneg"]
  \inputleannode{InfoGeometry.Canonical.ChiralAnomaly.routingEpsilon_nonneg}
\end{theorem}

\subsection*{selectedRoutingEpsilon}

\begin{theorem}[blueprint="infogeometry:canonical:chiralanomaly:selectedroutingepsilon"]
  \inputleannode{InfoGeometry.Canonical.ChiralAnomaly.selectedRoutingEpsilon}
\end{theorem}

\subsection*{selectedRoutingEpsilon_le_one}

\begin{theorem}[blueprint="infogeometry:canonical:chiralanomaly:selectedroutingepsilon_le_one"]
  \inputleannode{InfoGeometry.Canonical.ChiralAnomaly.selectedRoutingEpsilon_le_one}
\end{theorem}

\subsection*{sinkhornIterate}

\begin{theorem}[blueprint="infogeometry:canonical:chiralanomaly:sinkhorniterate"]
  \inputleannode{InfoGeometry.Canonical.ChiralAnomaly.sinkhornIterate}
\end{theorem}

\subsection*{sinkhornIterateDoublyStochasticTrajectory}

\begin{theorem}[blueprint="infogeometry:canonical:chiralanomaly:sinkhorniteratedoublystochastictrajectory"]
  \inputleannode{InfoGeometry.Canonical.ChiralAnomaly.sinkhornIterateDoublyStochasticTrajectory}
\end{theorem}

\subsection*{sinkhornIterateTrajectory}

\begin{theorem}[blueprint="infogeometry:canonical:chiralanomaly:sinkhorniteratetrajectory"]
  \inputleannode{InfoGeometry.Canonical.ChiralAnomaly.sinkhornIterateTrajectory}
\end{theorem}

\subsection*{sinkhornIterate_generator_step_control}

\begin{theorem}[blueprint="infogeometry:canonical:chiralanomaly:sinkhorniterate_generator_step_control"]
  \inputleannode{InfoGeometry.Canonical.ChiralAnomaly.sinkhornIterate_generator_step_control}
\end{theorem}

\subsection*{sinkhornIterate_step}

\begin{theorem}[blueprint="infogeometry:canonical:chiralanomaly:sinkhorniterate_step"]
  \inputleannode{InfoGeometry.Canonical.ChiralAnomaly.sinkhornIterate_step}
\end{theorem}

\subsection*{sinkhornIterate_step_control}

\begin{theorem}[blueprint="infogeometry:canonical:chiralanomaly:sinkhorniterate_step_control"]
  \inputleannode{InfoGeometry.Canonical.ChiralAnomaly.sinkhornIterate_step_control}
\end{theorem}

\subsection*{sinkhornPermutationWeights}

\begin{theorem}[blueprint="infogeometry:canonical:chiralanomaly:sinkhornpermutationweights"]
  \inputleannode{InfoGeometry.Canonical.ChiralAnomaly.sinkhornPermutationWeights}
\end{theorem}

\subsection*{sinkhornPermutationWeights_decomposition}

\begin{theorem}[blueprint="infogeometry:canonical:chiralanomaly:sinkhornpermutationweights_decomposition"]
  \inputleannode{InfoGeometry.Canonical.ChiralAnomaly.sinkhornPermutationWeights_decomposition}
\end{theorem}

\subsection*{sinkhornPermutationWeights_nonneg}

\begin{theorem}[blueprint="infogeometry:canonical:chiralanomaly:sinkhornpermutationweights_nonneg"]
  \inputleannode{InfoGeometry.Canonical.ChiralAnomaly.sinkhornPermutationWeights_nonneg}
\end{theorem}

\subsection*{sinkhornPermutationWeights_sum_one}

\begin{theorem}[blueprint="infogeometry:canonical:chiralanomaly:sinkhornpermutationweights_sum_one"]
  \inputleannode{InfoGeometry.Canonical.ChiralAnomaly.sinkhornPermutationWeights_sum_one}
\end{theorem}

\subsection*{sinkhorn_dynamics_generator_step_control}

\begin{theorem}[blueprint="infogeometry:canonical:chiralanomaly:sinkhorn_dynamics_generator_step_control"]
  \inputleannode{InfoGeometry.Canonical.ChiralAnomaly.sinkhorn_dynamics_generator_step_control}
\end{theorem}

\subsection*{sinkhorn_dynamics_step_control}

\begin{theorem}[blueprint="infogeometry:canonical:chiralanomaly:sinkhorn_dynamics_step_control"]
  \inputleannode{InfoGeometry.Canonical.ChiralAnomaly.sinkhorn_dynamics_step_control}
\end{theorem}

\subsection*{trajectorySelectedRoutingEpsilon}

\begin{theorem}[blueprint="infogeometry:canonical:chiralanomaly:trajectoryselectedroutingepsilon"]
  \inputleannode{InfoGeometry.Canonical.ChiralAnomaly.trajectorySelectedRoutingEpsilon}
\end{theorem}

\subsection*{trajectorySelectedRoutingEpsilon_le_one}

\begin{theorem}[blueprint="infogeometry:canonical:chiralanomaly:trajectoryselectedroutingepsilon_le_one"]
  \inputleannode{InfoGeometry.Canonical.ChiralAnomaly.trajectorySelectedRoutingEpsilon_le_one}
\end{theorem}

\section{InfoGeometry.Canonical.ChiralAnomaly.DoublyStochasticSinkhornTrajectory}

\subsection*{casesOn}

\begin{theorem}[blueprint="infogeometry:canonical:chiralanomaly:doublystochasticsinkhorntrajectory:caseson"]
  \inputleannode{InfoGeometry.Canonical.ChiralAnomaly.DoublyStochasticSinkhornTrajectory.casesOn}
\end{theorem}

\subsection*{ctorIdx}

\begin{theorem}[blueprint="infogeometry:canonical:chiralanomaly:doublystochasticsinkhorntrajectory:ctoridx"]
  \inputleannode{InfoGeometry.Canonical.ChiralAnomaly.DoublyStochasticSinkhornTrajectory.ctorIdx}
\end{theorem}

\subsection*{mem_doublyStochastic}

\begin{theorem}[blueprint="infogeometry:canonical:chiralanomaly:doublystochasticsinkhorntrajectory:mem_doublystochastic"]
  \inputleannode{InfoGeometry.Canonical.ChiralAnomaly.DoublyStochasticSinkhornTrajectory.mem_doublyStochastic}
\end{theorem}

\subsection*{mk}

\begin{theorem}[blueprint="infogeometry:canonical:chiralanomaly:doublystochasticsinkhorntrajectory:mk"]
  \inputleannode{InfoGeometry.Canonical.ChiralAnomaly.DoublyStochasticSinkhornTrajectory.mk}
\end{theorem}

\subsection*{noConfusion}

\begin{theorem}[blueprint="infogeometry:canonical:chiralanomaly:doublystochasticsinkhorntrajectory:noconfusion"]
  \inputleannode{InfoGeometry.Canonical.ChiralAnomaly.DoublyStochasticSinkhornTrajectory.noConfusion}
\end{theorem}

\subsection*{noConfusionType}

\begin{theorem}[blueprint="infogeometry:canonical:chiralanomaly:doublystochasticsinkhorntrajectory:noconfusiontype"]
  \inputleannode{InfoGeometry.Canonical.ChiralAnomaly.DoublyStochasticSinkhornTrajectory.noConfusionType}
\end{theorem}

\subsection*{rec}

\begin{theorem}[blueprint="infogeometry:canonical:chiralanomaly:doublystochasticsinkhorntrajectory:rec"]
  \inputleannode{InfoGeometry.Canonical.ChiralAnomaly.DoublyStochasticSinkhornTrajectory.rec}
\end{theorem}

\subsection*{recOn}

\begin{theorem}[blueprint="infogeometry:canonical:chiralanomaly:doublystochasticsinkhorntrajectory:recon"]
  \inputleannode{InfoGeometry.Canonical.ChiralAnomaly.DoublyStochasticSinkhornTrajectory.recOn}
\end{theorem}

\subsection*{traj}

\begin{theorem}[blueprint="infogeometry:canonical:chiralanomaly:doublystochasticsinkhorntrajectory:traj"]
  \inputleannode{InfoGeometry.Canonical.ChiralAnomaly.DoublyStochasticSinkhornTrajectory.traj}
\end{theorem}

\section{InfoGeometry.Canonical.ChiralAnomaly.inverseCommutator}

\subsection*{eq_1}

\begin{theorem}[blueprint="infogeometry:canonical:chiralanomaly:inversecommutator:eq_1"]
  \inputleannode{InfoGeometry.Canonical.ChiralAnomaly.inverseCommutator.eq_1}
\end{theorem}

\section{InfoGeometry.Canonical.ChiralAnomaly.sinkhornIterate}

\subsection*{congr_simp}

\begin{theorem}[blueprint="infogeometry:canonical:chiralanomaly:sinkhorniterate:congr_simp"]
  \inputleannode{InfoGeometry.Canonical.ChiralAnomaly.sinkhornIterate.congr_simp}
\end{theorem}

\subsection*{eq_1}

\begin{theorem}[blueprint="infogeometry:canonical:chiralanomaly:sinkhorniterate:eq_1"]
  \inputleannode{InfoGeometry.Canonical.ChiralAnomaly.sinkhornIterate.eq_1}
\end{theorem}

\subsection*{eq_2}

\begin{theorem}[blueprint="infogeometry:canonical:chiralanomaly:sinkhorniterate:eq_2"]
  \inputleannode{InfoGeometry.Canonical.ChiralAnomaly.sinkhornIterate.eq_2}
\end{theorem}

\subsection*{eq_def}

\begin{theorem}[blueprint="infogeometry:canonical:chiralanomaly:sinkhorniterate:eq_def"]
  \inputleannode{InfoGeometry.Canonical.ChiralAnomaly.sinkhornIterate.eq_def}
\end{theorem}

\section{InfoGeometry.Canonical.ChiralAnomaly.sinkhornIterateDoublyStochasticTrajectory}

\subsection*{eq_1}

\begin{theorem}[blueprint="infogeometry:canonical:chiralanomaly:sinkhorniteratedoublystochastictrajectory:eq_1"]
  \inputleannode{InfoGeometry.Canonical.ChiralAnomaly.sinkhornIterateDoublyStochasticTrajectory.eq_1}
\end{theorem}

\section{InfoGeometry.Canonical.ChiralAnomaly.sinkhornIterateTrajectory}

\subsection*{eq_1}

\begin{theorem}[blueprint="infogeometry:canonical:chiralanomaly:sinkhorniteratetrajectory:eq_1"]
  \inputleannode{InfoGeometry.Canonical.ChiralAnomaly.sinkhornIterateTrajectory.eq_1}
\end{theorem}

\section{InfoGeometry.Canonical.ChiralCliffordBridge}

\noindent\textbf{Buildable}: yes

\subsection*{IsCompactBeliefUpdate}

\begin{theorem}[blueprint="infogeometry:canonical:chiralcliffordbridge:iscompactbeliefupdate"]
  \inputleannode{InfoGeometry.Canonical.ChiralCliffordBridge.IsCompactBeliefUpdate}
\end{theorem}

\subsection*{IsNonCompactBeliefUpdate}

\begin{theorem}[blueprint="infogeometry:canonical:chiralcliffordbridge:isnoncompactbeliefupdate"]
  \inputleannode{InfoGeometry.Canonical.ChiralCliffordBridge.IsNonCompactBeliefUpdate}
\end{theorem}

\subsection*{anomaly_as_structure_constant}

\begin{theorem}[blueprint="infogeometry:canonical:chiralcliffordbridge:anomaly_as_structure_constant"]
  \inputleannode{InfoGeometry.Canonical.ChiralCliffordBridge.anomaly_as_structure_constant}
\end{theorem}

\subsection*{cartan_collapse_of_normal}

\begin{theorem}[blueprint="infogeometry:canonical:chiralcliffordbridge:cartan_collapse_of_normal"]
  \inputleannode{InfoGeometry.Canonical.ChiralCliffordBridge.cartan_collapse_of_normal}
\end{theorem}

\subsection*{chiralGrading}

\begin{theorem}[blueprint="infogeometry:canonical:chiralcliffordbridge:chiralgrading"]
  \inputleannode{InfoGeometry.Canonical.ChiralCliffordBridge.chiralGrading}
\end{theorem}

\subsection*{chiralProjectorMinus}

\begin{theorem}[blueprint="infogeometry:canonical:chiralcliffordbridge:chiralprojectorminus"]
  \inputleannode{InfoGeometry.Canonical.ChiralCliffordBridge.chiralProjectorMinus}
\end{theorem}

\subsection*{chiralProjectorPlus}

\begin{theorem}[blueprint="infogeometry:canonical:chiralcliffordbridge:chiralprojectorplus"]
  \inputleannode{InfoGeometry.Canonical.ChiralCliffordBridge.chiralProjectorPlus}
\end{theorem}

\subsection*{generalizedChiralMinus}

\begin{theorem}[blueprint="infogeometry:canonical:chiralcliffordbridge:generalizedchiralminus"]
  \inputleannode{InfoGeometry.Canonical.ChiralCliffordBridge.generalizedChiralMinus}
\end{theorem}

\subsection*{generalizedChiralPlus}

\begin{theorem}[blueprint="infogeometry:canonical:chiralcliffordbridge:generalizedchiralplus"]
  \inputleannode{InfoGeometry.Canonical.ChiralCliffordBridge.generalizedChiralPlus}
\end{theorem}

\section{InfoGeometry.Canonical.ChiralEinsteinBridge}

\noindent\textbf{Buildable}: yes

\subsection*{SatisfiesAnomalyDrivenKaehlerRicciFlow}

\begin{theorem}[blueprint="infogeometry:canonical:chiraleinsteinbridge:satisfiesanomalydrivenkaehlerricciflow"]
  \inputleannode{InfoGeometry.Canonical.ChiralEinsteinBridge.SatisfiesAnomalyDrivenKaehlerRicciFlow}
\end{theorem}

\subsection*{SatisfiesAnomalyDrivenScalarRicciFlow}

\begin{theorem}[blueprint="infogeometry:canonical:chiraleinsteinbridge:satisfiesanomalydrivenscalarricciflow"]
  \inputleannode{InfoGeometry.Canonical.ChiralEinsteinBridge.SatisfiesAnomalyDrivenScalarRicciFlow}
\end{theorem}

\subsection*{anomalyDrivenScalarRicci_fixedpoint_tracks_source}

\begin{theorem}[blueprint="infogeometry:canonical:chiraleinsteinbridge:anomalydrivenscalarricci_fixedpoint_tracks_source"]
  \inputleannode{InfoGeometry.Canonical.ChiralEinsteinBridge.anomalyDrivenScalarRicci_fixedpoint_tracks_source}
\end{theorem}

\subsection*{anomalyDriven_fixedpoint_eq_inverseEpsilon}

\begin{theorem}[blueprint="infogeometry:canonical:chiraleinsteinbridge:anomalydriven_fixedpoint_eq_inverseepsilon"]
  \inputleannode{InfoGeometry.Canonical.ChiralEinsteinBridge.anomalyDriven_fixedpoint_eq_inverseEpsilon}
\end{theorem}

\subsection*{anomalyDriven_fixedpoint_tracks_source}

\begin{theorem}[blueprint="infogeometry:canonical:chiraleinsteinbridge:anomalydriven_fixedpoint_tracks_source"]
  \inputleannode{InfoGeometry.Canonical.ChiralEinsteinBridge.anomalyDriven_fixedpoint_tracks_source}
\end{theorem}

\subsection*{anomalyDriven_zeroSource_iff_normalized}

\begin{theorem}[blueprint="infogeometry:canonical:chiraleinsteinbridge:anomalydriven_zerosource_iff_normalized"]
  \inputleannode{InfoGeometry.Canonical.ChiralEinsteinBridge.anomalyDriven_zeroSource_iff_normalized}
\end{theorem}

\subsection*{anomalyStressEnergyAt}

\begin{theorem}[blueprint="infogeometry:canonical:chiraleinsteinbridge:anomalystressenergyat"]
  \inputleannode{InfoGeometry.Canonical.ChiralEinsteinBridge.anomalyStressEnergyAt}
\end{theorem}

\subsection*{anomalyStressEnergyAt_apply}

\begin{theorem}[blueprint="infogeometry:canonical:chiraleinsteinbridge:anomalystressenergyat_apply"]
  \inputleannode{InfoGeometry.Canonical.ChiralEinsteinBridge.anomalyStressEnergyAt_apply}
\end{theorem}

\subsection*{anomalyStressEnergyAt_zero}

\begin{theorem}[blueprint="infogeometry:canonical:chiraleinsteinbridge:anomalystressenergyat_zero"]
  \inputleannode{InfoGeometry.Canonical.ChiralEinsteinBridge.anomalyStressEnergyAt_zero}
\end{theorem}

\subsection*{anomalyStressEnergyModelAt}

\begin{theorem}[blueprint="infogeometry:canonical:chiraleinsteinbridge:anomalystressenergymodelat"]
  \inputleannode{InfoGeometry.Canonical.ChiralEinsteinBridge.anomalyStressEnergyModelAt}
\end{theorem}

\subsection*{einsteinEquation_of_anomaly_source}

\begin{theorem}[blueprint="infogeometry:canonical:chiraleinsteinbridge:einsteinequation_of_anomaly_source"]
  \inputleannode{InfoGeometry.Canonical.ChiralEinsteinBridge.einsteinEquation_of_anomaly_source}
\end{theorem}

\subsection*{exists_einsteinEquation_of_bistochastic_routingAnomaly}

\begin{theorem}[blueprint="infogeometry:canonical:chiraleinsteinbridge:exists_einsteinequation_of_bistochastic_routinganomaly"]
  \inputleannode{InfoGeometry.Canonical.ChiralEinsteinBridge.exists_einsteinEquation_of_bistochastic_routingAnomaly}
\end{theorem}

\subsection*{inverseEpsilonSource}

\begin{theorem}[blueprint="infogeometry:canonical:chiraleinsteinbridge:inverseepsilonsource"]
  \inputleannode{InfoGeometry.Canonical.ChiralEinsteinBridge.inverseEpsilonSource}
\end{theorem}

\section{InfoGeometry.Canonical.ChiralEinsteinBridge.SatisfiesAnomalyDrivenKaehlerRicciFlow}

\subsection*{eq_1}

\begin{theorem}[blueprint="infogeometry:canonical:chiraleinsteinbridge:satisfiesanomalydrivenkaehlerricciflow:eq_1"]
  \inputleannode{InfoGeometry.Canonical.ChiralEinsteinBridge.SatisfiesAnomalyDrivenKaehlerRicciFlow.eq_1}
\end{theorem}

\section{InfoGeometry.Canonical.ChiralEinsteinBridge.anomalyStressEnergyAt}

\subsection*{eq_1}

\begin{theorem}[blueprint="infogeometry:canonical:chiraleinsteinbridge:anomalystressenergyat:eq_1"]
  \inputleannode{InfoGeometry.Canonical.ChiralEinsteinBridge.anomalyStressEnergyAt.eq_1}
\end{theorem}

\section{InfoGeometry.Canonical.ChiralEinsteinBridge.inverseEpsilonSource}

\subsection*{eq_1}

\begin{theorem}[blueprint="infogeometry:canonical:chiraleinsteinbridge:inverseepsilonsource:eq_1"]
  \inputleannode{InfoGeometry.Canonical.ChiralEinsteinBridge.inverseEpsilonSource.eq_1}
\end{theorem}

\section{InfoGeometry.Canonical.ChiralGravity}

\noindent\textbf{Buildable}: yes

\subsection*{AnomalyCurvatureForcingStateAt}

\begin{theorem}[blueprint="infogeometry:canonical:chiralgravity:anomalycurvatureforcingstateat"]
  \inputleannode{InfoGeometry.Canonical.ChiralGravity.AnomalyCurvatureForcingStateAt}
\end{theorem}

\subsection*{anomalyEinsteinResidualAt}

\begin{theorem}[blueprint="infogeometry:canonical:chiralgravity:anomalyeinsteinresidualat"]
  \inputleannode{InfoGeometry.Canonical.ChiralGravity.anomalyEinsteinResidualAt}
\end{theorem}

\subsection*{anomalyEinsteinResidualAt_def}

\begin{theorem}[blueprint="infogeometry:canonical:chiralgravity:anomalyeinsteinresidualat_def"]
  \inputleannode{InfoGeometry.Canonical.ChiralGravity.anomalyEinsteinResidualAt_def}
\end{theorem}

\subsection*{anomalyEinsteinResidual_eq_kappa_mul_metric}

\begin{theorem}[blueprint="infogeometry:canonical:chiralgravity:anomalyeinsteinresidual_eq_kappa_mul_metric"]
  \inputleannode{InfoGeometry.Canonical.ChiralGravity.anomalyEinsteinResidual_eq_kappa_mul_metric}
\end{theorem}

\subsection*{anomaly_nonzero_excludes_vacuum}

\begin{theorem}[blueprint="infogeometry:canonical:chiralgravity:anomaly_nonzero_excludes_vacuum"]
  \inputleannode{InfoGeometry.Canonical.ChiralGravity.anomaly_nonzero_excludes_vacuum}
\end{theorem}

\subsection*{anomaly_nonzero_forces_curved_plus_component}

\begin{theorem}[blueprint="infogeometry:canonical:chiralgravity:anomaly_nonzero_forces_curved_plus_component"]
  \inputleannode{InfoGeometry.Canonical.ChiralGravity.anomaly_nonzero_forces_curved_plus_component}
\end{theorem}

\subsection*{routingAnomaly_nonzero_forces_curved_plus_component}

\begin{theorem}[blueprint="infogeometry:canonical:chiralgravity:routinganomaly_nonzero_forces_curved_plus_component"]
  \inputleannode{InfoGeometry.Canonical.ChiralGravity.routingAnomaly_nonzero_forces_curved_plus_component}
\end{theorem}

\section{InfoGeometry.Canonical.ChiralRGFlow}

\noindent\textbf{Buildable}: yes

\subsection*{ChiralAsymptoticModel}

\begin{theorem}[blueprint="infogeometry:canonical:chiralrgflow:chiralasymptoticmodel"]
  \inputleannode{InfoGeometry.Canonical.ChiralRGFlow.ChiralAsymptoticModel}
\end{theorem}

\subsection*{ChiralMassFlow}

\begin{theorem}[blueprint="infogeometry:canonical:chiralrgflow:chiralmassflow"]
  \inputleannode{InfoGeometry.Canonical.ChiralRGFlow.ChiralMassFlow}
\end{theorem}

\subsection*{IsAsymptoticallyFree}

\begin{theorem}[blueprint="infogeometry:canonical:chiralrgflow:isasymptoticallyfree"]
  \inputleannode{InfoGeometry.Canonical.ChiralRGFlow.IsAsymptoticallyFree}
\end{theorem}

\subsection*{IsInfraredConfined}

\begin{theorem}[blueprint="infogeometry:canonical:chiralrgflow:isinfraredconfined"]
  \inputleannode{InfoGeometry.Canonical.ChiralRGFlow.IsInfraredConfined}
\end{theorem}

\subsection*{asymptotic_freedom_of_negative_beta}

\begin{theorem}[blueprint="infogeometry:canonical:chiralrgflow:asymptotic_freedom_of_negative_beta"]
  \inputleannode{InfoGeometry.Canonical.ChiralRGFlow.asymptotic_freedom_of_negative_beta}
\end{theorem}

\subsection*{explicitBetaFunction}

\begin{theorem}[blueprint="infogeometry:canonical:chiralrgflow:explicitbetafunction"]
  \inputleannode{InfoGeometry.Canonical.ChiralRGFlow.explicitBetaFunction}
\end{theorem}

\subsection*{massBetaFunction}

\begin{theorem}[blueprint="infogeometry:canonical:chiralrgflow:massbetafunction"]
  \inputleannode{InfoGeometry.Canonical.ChiralRGFlow.massBetaFunction}
\end{theorem}

\section{InfoGeometry.Canonical.ChiralRGFlow.ChiralAsymptoticModel}

\subsection*{asymptotic}

\begin{theorem}[blueprint="infogeometry:canonical:chiralrgflow:chiralasymptoticmodel:asymptotic"]
  \inputleannode{InfoGeometry.Canonical.ChiralRGFlow.ChiralAsymptoticModel.asymptotic}
\end{theorem}

\subsection*{beta_eq}

\begin{theorem}[blueprint="infogeometry:canonical:chiralrgflow:chiralasymptoticmodel:beta_eq"]
  \inputleannode{InfoGeometry.Canonical.ChiralRGFlow.ChiralAsymptoticModel.beta_eq}
\end{theorem}

\subsection*{casesOn}

\begin{theorem}[blueprint="infogeometry:canonical:chiralrgflow:chiralasymptoticmodel:caseson"]
  \inputleannode{InfoGeometry.Canonical.ChiralRGFlow.ChiralAsymptoticModel.casesOn}
\end{theorem}

\subsection*{ctorIdx}

\begin{theorem}[blueprint="infogeometry:canonical:chiralrgflow:chiralasymptoticmodel:ctoridx"]
  \inputleannode{InfoGeometry.Canonical.ChiralRGFlow.ChiralAsymptoticModel.ctorIdx}
\end{theorem}

\subsection*{flow}

\begin{theorem}[blueprint="infogeometry:canonical:chiralrgflow:chiralasymptoticmodel:flow"]
  \inputleannode{InfoGeometry.Canonical.ChiralRGFlow.ChiralAsymptoticModel.flow}
\end{theorem}

\subsection*{gamma}

\begin{theorem}[blueprint="infogeometry:canonical:chiralrgflow:chiralasymptoticmodel:gamma"]
  \inputleannode{InfoGeometry.Canonical.ChiralRGFlow.ChiralAsymptoticModel.gamma}
\end{theorem}

\subsection*{gamma_pos}

\begin{theorem}[blueprint="infogeometry:canonical:chiralrgflow:chiralasymptoticmodel:gamma_pos"]
  \inputleannode{InfoGeometry.Canonical.ChiralRGFlow.ChiralAsymptoticModel.gamma_pos}
\end{theorem}

\subsection*{mk}

\begin{theorem}[blueprint="infogeometry:canonical:chiralrgflow:chiralasymptoticmodel:mk"]
  \inputleannode{InfoGeometry.Canonical.ChiralRGFlow.ChiralAsymptoticModel.mk}
\end{theorem}

\subsection*{noConfusion}

\begin{theorem}[blueprint="infogeometry:canonical:chiralrgflow:chiralasymptoticmodel:noconfusion"]
  \inputleannode{InfoGeometry.Canonical.ChiralRGFlow.ChiralAsymptoticModel.noConfusion}
\end{theorem}

\subsection*{noConfusionType}

\begin{theorem}[blueprint="infogeometry:canonical:chiralrgflow:chiralasymptoticmodel:noconfusiontype"]
  \inputleannode{InfoGeometry.Canonical.ChiralRGFlow.ChiralAsymptoticModel.noConfusionType}
\end{theorem}

\subsection*{rec}

\begin{theorem}[blueprint="infogeometry:canonical:chiralrgflow:chiralasymptoticmodel:rec"]
  \inputleannode{InfoGeometry.Canonical.ChiralRGFlow.ChiralAsymptoticModel.rec}
\end{theorem}

\subsection*{recOn}

\begin{theorem}[blueprint="infogeometry:canonical:chiralrgflow:chiralasymptoticmodel:recon"]
  \inputleannode{InfoGeometry.Canonical.ChiralRGFlow.ChiralAsymptoticModel.recOn}
\end{theorem}

\section{InfoGeometry.Canonical.ChiralTorsionBridge}

\noindent\textbf{Buildable}: yes

\subsection*{ChiralTorsionChentsovGibbsBridge}

\begin{theorem}[blueprint="infogeometry:canonical:chiraltorsionbridge:chiraltorsionchentsovgibbsbridge"]
  \inputleannode{InfoGeometry.Canonical.ChiralTorsionBridge.ChiralTorsionChentsovGibbsBridge}
\end{theorem}

\subsection*{ChiralTorsionChentsovGibbsState}

\begin{theorem}[blueprint="infogeometry:canonical:chiraltorsionbridge:chiraltorsionchentsovgibbsstate"]
  \inputleannode{InfoGeometry.Canonical.ChiralTorsionBridge.ChiralTorsionChentsovGibbsState}
\end{theorem}

\subsection*{IsVacuumApexNull}

\begin{theorem}[blueprint="infogeometry:canonical:chiraltorsionbridge:isvacuumapexnull"]
  \inputleannode{InfoGeometry.Canonical.ChiralTorsionBridge.IsVacuumApexNull}
\end{theorem}

\subsection*{UnnormalizedProjectiveState}

\begin{theorem}[blueprint="infogeometry:canonical:chiraltorsionbridge:unnormalizedprojectivestate"]
  \inputleannode{InfoGeometry.Canonical.ChiralTorsionBridge.UnnormalizedProjectiveState}
\end{theorem}

\subsection*{boltzmannRelativeVolumeEntropy}

\begin{theorem}[blueprint="infogeometry:canonical:chiraltorsionbridge:boltzmannrelativevolumeentropy"]
  \inputleannode{InfoGeometry.Canonical.ChiralTorsionBridge.boltzmannRelativeVolumeEntropy}
\end{theorem}

\subsection*{boltzmannRelativeVolumeEntropy_eq_divergence}

\begin{theorem}[blueprint="infogeometry:canonical:chiraltorsionbridge:boltzmannrelativevolumeentropy_eq_divergence"]
  \inputleannode{InfoGeometry.Canonical.ChiralTorsionBridge.boltzmannRelativeVolumeEntropy_eq_divergence}
\end{theorem}

\subsection*{chentsov_and_gibbs_of_bridge}

\begin{theorem}[blueprint="infogeometry:canonical:chiraltorsionbridge:chentsov_and_gibbs_of_bridge"]
  \inputleannode{InfoGeometry.Canonical.ChiralTorsionBridge.chentsov_and_gibbs_of_bridge}
\end{theorem}

\subsection*{chentsov_and_gibbs_of_state}

\begin{theorem}[blueprint="infogeometry:canonical:chiraltorsionbridge:chentsov_and_gibbs_of_state"]
  \inputleannode{InfoGeometry.Canonical.ChiralTorsionBridge.chentsov_and_gibbs_of_state}
\end{theorem}

\subsection*{chiralTorsionChentsovGibbsState_mk}

\begin{theorem}[blueprint="infogeometry:canonical:chiraltorsionbridge:chiraltorsionchentsovgibbsstate_mk"]
  \inputleannode{InfoGeometry.Canonical.ChiralTorsionBridge.chiralTorsionChentsovGibbsState_mk}
\end{theorem}

\subsection*{chiralTorsionChentsovGibbsState_of_twistedInference}

\begin{theorem}[blueprint="infogeometry:canonical:chiraltorsionbridge:chiraltorsionchentsovgibbsstate_of_twistedinference"]
  \inputleannode{InfoGeometry.Canonical.ChiralTorsionBridge.chiralTorsionChentsovGibbsState_of_twistedInference}
\end{theorem}

\subsection*{gibbsSmoothingOnGeneralizedKL}

\begin{theorem}[blueprint="infogeometry:canonical:chiraltorsionbridge:gibbssmoothingongeneralizedkl"]
  \inputleannode{InfoGeometry.Canonical.ChiralTorsionBridge.gibbsSmoothingOnGeneralizedKL}
\end{theorem}

\subsection*{gibbsSmoothingOnGeneralizedKL_nonneg}

\begin{theorem}[blueprint="infogeometry:canonical:chiraltorsionbridge:gibbssmoothingongeneralizedkl_nonneg"]
  \inputleannode{InfoGeometry.Canonical.ChiralTorsionBridge.gibbsSmoothingOnGeneralizedKL_nonneg}
\end{theorem}

\subsection*{gibbsSmoothingOnGeneralizedKL_pos}

\begin{theorem}[blueprint="infogeometry:canonical:chiraltorsionbridge:gibbssmoothingongeneralizedkl_pos"]
  \inputleannode{InfoGeometry.Canonical.ChiralTorsionBridge.gibbsSmoothingOnGeneralizedKL_pos}
\end{theorem}

\subsection*{relativeVacuumVolume}

\begin{theorem}[blueprint="infogeometry:canonical:chiraltorsionbridge:relativevacuumvolume"]
  \inputleannode{InfoGeometry.Canonical.ChiralTorsionBridge.relativeVacuumVolume}
\end{theorem}

\subsection*{torsion_nonzero_of_chiral}

\begin{theorem}[blueprint="infogeometry:canonical:chiraltorsionbridge:torsion_nonzero_of_chiral"]
  \inputleannode{InfoGeometry.Canonical.ChiralTorsionBridge.torsion_nonzero_of_chiral}
\end{theorem}

\subsection*{torsion_nonzero_of_state_and_chiral}

\begin{theorem}[blueprint="infogeometry:canonical:chiraltorsionbridge:torsion_nonzero_of_state_and_chiral"]
  \inputleannode{InfoGeometry.Canonical.ChiralTorsionBridge.torsion_nonzero_of_state_and_chiral}
\end{theorem}

\subsection*{twistedInference_torsion_nonzero}

\begin{theorem}[blueprint="infogeometry:canonical:chiraltorsionbridge:twistedinference_torsion_nonzero"]
  \inputleannode{InfoGeometry.Canonical.ChiralTorsionBridge.twistedInference_torsion_nonzero}
\end{theorem}

\subsection*{vacuumApexTwistor}

\begin{theorem}[blueprint="infogeometry:canonical:chiraltorsionbridge:vacuumapextwistor"]
  \inputleannode{InfoGeometry.Canonical.ChiralTorsionBridge.vacuumApexTwistor}
\end{theorem}

\section{InfoGeometry.Canonical.ChiralTorsionBridge.ChiralTorsionChentsovGibbsBridge}

\subsection*{anomaly_induces_torsion}

\begin{theorem}[blueprint="infogeometry:canonical:chiraltorsionbridge:chiraltorsionchentsovgibbsbridge:anomaly_induces_torsion"]
  \inputleannode{InfoGeometry.Canonical.ChiralTorsionBridge.ChiralTorsionChentsovGibbsBridge.anomaly_induces_torsion}
\end{theorem}

\subsection*{casesOn}

\begin{theorem}[blueprint="infogeometry:canonical:chiraltorsionbridge:chiraltorsionchentsovgibbsbridge:caseson"]
  \inputleannode{InfoGeometry.Canonical.ChiralTorsionBridge.ChiralTorsionChentsovGibbsBridge.casesOn}
\end{theorem}

\subsection*{chentsov_available}

\begin{theorem}[blueprint="infogeometry:canonical:chiraltorsionbridge:chiraltorsionchentsovgibbsbridge:chentsov_available"]
  \inputleannode{InfoGeometry.Canonical.ChiralTorsionBridge.ChiralTorsionChentsovGibbsBridge.chentsov_available}
\end{theorem}

\subsection*{gibbs_smoothing_pos}

\begin{theorem}[blueprint="infogeometry:canonical:chiraltorsionbridge:chiraltorsionchentsovgibbsbridge:gibbs_smoothing_pos"]
  \inputleannode{InfoGeometry.Canonical.ChiralTorsionBridge.ChiralTorsionChentsovGibbsBridge.gibbs_smoothing_pos}
\end{theorem}

\subsection*{mk}

\begin{theorem}[blueprint="infogeometry:canonical:chiraltorsionbridge:chiraltorsionchentsovgibbsbridge:mk"]
  \inputleannode{InfoGeometry.Canonical.ChiralTorsionBridge.ChiralTorsionChentsovGibbsBridge.mk}
\end{theorem}

\subsection*{rec}

\begin{theorem}[blueprint="infogeometry:canonical:chiraltorsionbridge:chiraltorsionchentsovgibbsbridge:rec"]
  \inputleannode{InfoGeometry.Canonical.ChiralTorsionBridge.ChiralTorsionChentsovGibbsBridge.rec}
\end{theorem}

\subsection*{recOn}

\begin{theorem}[blueprint="infogeometry:canonical:chiraltorsionbridge:chiraltorsionchentsovgibbsbridge:recon"]
  \inputleannode{InfoGeometry.Canonical.ChiralTorsionBridge.ChiralTorsionChentsovGibbsBridge.recOn}
\end{theorem}

\section{InfoGeometry.Canonical.Clifford}

\noindent\textbf{Buildable}: yes

\noindent\emph{No exported declarations in the current declaration graph.}

\section{InfoGeometry.Canonical.CliffordBridge}

\noindent\textbf{Buildable}: yes

\subsection*{B_agrees_with_Gauge_bilinear}

\begin{theorem}[blueprint="infogeometry:canonical:cliffordbridge:b_agrees_with_gauge_bilinear"]
  \inputleannode{InfoGeometry.Canonical.CliffordBridge.B_agrees_with_Gauge_bilinear}
\end{theorem}

\subsection*{q_agrees_with_Gauge_quad}

\begin{theorem}[blueprint="infogeometry:canonical:cliffordbridge:q_agrees_with_gauge_quad"]
  \inputleannode{InfoGeometry.Canonical.CliffordBridge.q_agrees_with_Gauge_quad}
\end{theorem}

\subsection*{splitBilinear_eq_gaugeBilinear}

\begin{theorem}[blueprint="infogeometry:canonical:cliffordbridge:splitbilinear_eq_gaugebilinear"]
  \inputleannode{InfoGeometry.Canonical.CliffordBridge.splitBilinear_eq_gaugeBilinear}
\end{theorem}

\subsection*{splitQuadratic_eq_gaugeQuadratic}

\begin{theorem}[blueprint="infogeometry:canonical:cliffordbridge:splitquadratic_eq_gaugequadratic"]
  \inputleannode{InfoGeometry.Canonical.CliffordBridge.splitQuadratic_eq_gaugeQuadratic}
\end{theorem}

\section{InfoGeometry.Canonical.ConformalAlgebra}

\noindent\textbf{Buildable}: yes

\subsection*{ConformalBeliefAlgebra}

\begin{theorem}[blueprint="infogeometry:canonical:conformalalgebra:conformalbeliefalgebra"]
  \inputleannode{InfoGeometry.Canonical.ConformalAlgebra.ConformalBeliefAlgebra}
\end{theorem}

\section{InfoGeometry.Canonical.ConformalAlgebra.ConformalBeliefAlgebra}

\subsection*{CI}

\begin{theorem}[blueprint="infogeometry:canonical:conformalalgebra:conformalbeliefalgebra:ci"]
  \inputleannode{InfoGeometry.Canonical.ConformalAlgebra.ConformalBeliefAlgebra.CI}
\end{theorem}

\subsection*{ConformalWeightClosure}

\begin{theorem}[blueprint="infogeometry:canonical:conformalalgebra:conformalbeliefalgebra:conformalweightclosure"]
  \inputleannode{InfoGeometry.Canonical.ConformalAlgebra.ConformalBeliefAlgebra.ConformalWeightClosure}
\end{theorem}

\subsection*{D}

\begin{theorem}[blueprint="infogeometry:canonical:conformalalgebra:conformalbeliefalgebra:d"]
  \inputleannode{InfoGeometry.Canonical.ConformalAlgebra.ConformalBeliefAlgebra.D}
\end{theorem}

\subsection*{GeneratorCartanDecomposition}

\begin{theorem}[blueprint="infogeometry:canonical:conformalalgebra:conformalbeliefalgebra:generatorcartandecomposition"]
  \inputleannode{InfoGeometry.Canonical.ConformalAlgebra.ConformalBeliefAlgebra.GeneratorCartanDecomposition}
\end{theorem}

\subsection*{M}

\begin{theorem}[blueprint="infogeometry:canonical:conformalalgebra:conformalbeliefalgebra:m"]
  \inputleannode{InfoGeometry.Canonical.ConformalAlgebra.ConformalBeliefAlgebra.M}
\end{theorem}

\subsection*{SatisfiesFlatWeights}

\begin{theorem}[blueprint="infogeometry:canonical:conformalalgebra:conformalbeliefalgebra:satisfiesflatweights"]
  \inputleannode{InfoGeometry.Canonical.ConformalAlgebra.ConformalBeliefAlgebra.SatisfiesFlatWeights}
\end{theorem}

\subsection*{SatisfiesKWeight}

\begin{theorem}[blueprint="infogeometry:canonical:conformalalgebra:conformalbeliefalgebra:satisfieskweight"]
  \inputleannode{InfoGeometry.Canonical.ConformalAlgebra.ConformalBeliefAlgebra.SatisfiesKWeight}
\end{theorem}

\subsection*{SatisfiesMasterRelation}

\begin{theorem}[blueprint="infogeometry:canonical:conformalalgebra:conformalbeliefalgebra:satisfiesmasterrelation"]
  \inputleannode{InfoGeometry.Canonical.ConformalAlgebra.ConformalBeliefAlgebra.SatisfiesMasterRelation}
\end{theorem}

\subsection*{SatisfiesPWeight}

\begin{theorem}[blueprint="infogeometry:canonical:conformalalgebra:conformalbeliefalgebra:satisfiespweight"]
  \inputleannode{InfoGeometry.Canonical.ConformalAlgebra.ConformalBeliefAlgebra.SatisfiesPWeight}
\end{theorem}

\subsection*{anomaly_breaks_weights}

\begin{theorem}[blueprint="infogeometry:canonical:conformalalgebra:conformalbeliefalgebra:anomaly_breaks_weights"]
  \inputleannode{InfoGeometry.Canonical.ConformalAlgebra.ConformalBeliefAlgebra.anomaly_breaks_weights}
\end{theorem}

\subsection*{casesOn}

\begin{theorem}[blueprint="infogeometry:canonical:conformalalgebra:conformalbeliefalgebra:caseson"]
  \inputleannode{InfoGeometry.Canonical.ConformalAlgebra.ConformalBeliefAlgebra.casesOn}
\end{theorem}

\subsection*{ctorIdx}

\begin{theorem}[blueprint="infogeometry:canonical:conformalalgebra:conformalbeliefalgebra:ctoridx"]
  \inputleannode{InfoGeometry.Canonical.ConformalAlgebra.ConformalBeliefAlgebra.ctorIdx}
\end{theorem}

\subsection*{mk}

\begin{theorem}[blueprint="infogeometry:canonical:conformalalgebra:conformalbeliefalgebra:mk"]
  \inputleannode{InfoGeometry.Canonical.ConformalAlgebra.ConformalBeliefAlgebra.mk}
\end{theorem}

\subsection*{noConfusion}

\begin{theorem}[blueprint="infogeometry:canonical:conformalalgebra:conformalbeliefalgebra:noconfusion"]
  \inputleannode{InfoGeometry.Canonical.ConformalAlgebra.ConformalBeliefAlgebra.noConfusion}
\end{theorem}

\subsection*{noConfusionType}

\begin{theorem}[blueprint="infogeometry:canonical:conformalalgebra:conformalbeliefalgebra:noconfusiontype"]
  \inputleannode{InfoGeometry.Canonical.ConformalAlgebra.ConformalBeliefAlgebra.noConfusionType}
\end{theorem}

\subsection*{rec}

\begin{theorem}[blueprint="infogeometry:canonical:conformalalgebra:conformalbeliefalgebra:rec"]
  \inputleannode{InfoGeometry.Canonical.ConformalAlgebra.ConformalBeliefAlgebra.rec}
\end{theorem}

\subsection*{recOn}

\begin{theorem}[blueprint="infogeometry:canonical:conformalalgebra:conformalbeliefalgebra:recon"]
  \inputleannode{InfoGeometry.Canonical.ConformalAlgebra.ConformalBeliefAlgebra.recOn}
\end{theorem}

\subsection*{scale_anomaly_breaks_weight_closure}

\begin{theorem}[blueprint="infogeometry:canonical:conformalalgebra:conformalbeliefalgebra:scale_anomaly_breaks_weight_closure"]
  \inputleannode{InfoGeometry.Canonical.ConformalAlgebra.ConformalBeliefAlgebra.scale_anomaly_breaks_weight_closure}
\end{theorem}

\subsection*{scale_anomaly_emergence}

\begin{theorem}[blueprint="infogeometry:canonical:conformalalgebra:conformalbeliefalgebra:scale_anomaly_emergence"]
  \inputleannode{InfoGeometry.Canonical.ConformalAlgebra.ConformalBeliefAlgebra.scale_anomaly_emergence}
\end{theorem}

\section{InfoGeometry.Canonical.ConformalAlgebra.ConformalBeliefAlgebra.SatisfiesFlatWeights}

\subsection*{eq_1}

\begin{theorem}[blueprint="infogeometry:canonical:conformalalgebra:conformalbeliefalgebra:satisfiesflatweights:eq_1"]
  \inputleannode{InfoGeometry.Canonical.ConformalAlgebra.ConformalBeliefAlgebra.SatisfiesFlatWeights.eq_1}
\end{theorem}

\section{InfoGeometry.Canonical.ConformalAlgebra.ConformalBeliefAlgebra.SatisfiesKWeight}

\subsection*{eq_1}

\begin{theorem}[blueprint="infogeometry:canonical:conformalalgebra:conformalbeliefalgebra:satisfieskweight:eq_1"]
  \inputleannode{InfoGeometry.Canonical.ConformalAlgebra.ConformalBeliefAlgebra.SatisfiesKWeight.eq_1}
\end{theorem}

\section{InfoGeometry.Canonical.ConformalAlgebra.ConformalBeliefAlgebra.SatisfiesPWeight}

\subsection*{eq_1}

\begin{theorem}[blueprint="infogeometry:canonical:conformalalgebra:conformalbeliefalgebra:satisfiespweight:eq_1"]
  \inputleannode{InfoGeometry.Canonical.ConformalAlgebra.ConformalBeliefAlgebra.SatisfiesPWeight.eq_1}
\end{theorem}

\section{InfoGeometry.Canonical.ConformalUnification}

\noindent\textbf{Buildable}: yes

\subsection*{ConformalInference}

\begin{theorem}[blueprint="infogeometry:canonical:conformalunification:conformalinference"]
  \inputleannode{InfoGeometry.Canonical.ConformalUnification.ConformalInference}
\end{theorem}

\section{InfoGeometry.Canonical.ConformalUnification.ConformalInference}

\subsection*{A}

\begin{theorem}[blueprint="infogeometry:canonical:conformalunification:conformalinference:a"]
  \inputleannode{InfoGeometry.Canonical.ConformalUnification.ConformalInference.A}
\end{theorem}

\subsection*{A_D}

\begin{theorem}[blueprint="infogeometry:canonical:conformalunification:conformalinference:a_d"]
  \inputleannode{InfoGeometry.Canonical.ConformalUnification.ConformalInference.A_D}
\end{theorem}

\subsection*{A_MP}

\begin{theorem}[blueprint="infogeometry:canonical:conformalunification:conformalinference:a_mp"]
  \inputleannode{InfoGeometry.Canonical.ConformalUnification.ConformalInference.A_MP}
\end{theorem}

\subsection*{ChiralInferenceState}

\begin{theorem}[blueprint="infogeometry:canonical:conformalunification:conformalinference:chiralinferencestate"]
  \inputleannode{InfoGeometry.Canonical.ConformalUnification.ConformalInference.ChiralInferenceState}
\end{theorem}

\subsection*{D}

\begin{theorem}[blueprint="infogeometry:canonical:conformalunification:conformalinference:d"]
  \inputleannode{InfoGeometry.Canonical.ConformalUnification.ConformalInference.D}
\end{theorem}

\subsection*{IsChiralInference}

\begin{theorem}[blueprint="infogeometry:canonical:conformalunification:conformalinference:ischiralinference"]
  \inputleannode{InfoGeometry.Canonical.ConformalUnification.ConformalInference.IsChiralInference}
\end{theorem}

\subsection*{IsNormalInference}

\begin{theorem}[blueprint="infogeometry:canonical:conformalunification:conformalinference:isnormalinference"]
  \inputleannode{InfoGeometry.Canonical.ConformalUnification.ConformalInference.IsNormalInference}
\end{theorem}

\subsection*{K}

\begin{theorem}[blueprint="infogeometry:canonical:conformalunification:conformalinference:k"]
  \inputleannode{InfoGeometry.Canonical.ConformalUnification.ConformalInference.K}
\end{theorem}

\subsection*{NormalInferenceState}

\begin{theorem}[blueprint="infogeometry:canonical:conformalunification:conformalinference:normalinferencestate"]
  \inputleannode{InfoGeometry.Canonical.ConformalUnification.ConformalInference.NormalInferenceState}
\end{theorem}

\subsection*{P}

\begin{theorem}[blueprint="infogeometry:canonical:conformalunification:conformalinference:p"]
  \inputleannode{InfoGeometry.Canonical.ConformalUnification.ConformalInference.P}
\end{theorem}

\subsection*{P_D}

\begin{theorem}[blueprint="infogeometry:canonical:conformalunification:conformalinference:p_d"]
  \inputleannode{InfoGeometry.Canonical.ConformalUnification.ConformalInference.P_D}
\end{theorem}

\subsection*{P_MP}

\begin{theorem}[blueprint="infogeometry:canonical:conformalunification:conformalinference:p_mp"]
  \inputleannode{InfoGeometry.Canonical.ConformalUnification.ConformalInference.P_MP}
\end{theorem}

\subsection*{ProjectorCommutationClosure}

\begin{theorem}[blueprint="infogeometry:canonical:conformalunification:conformalinference:projectorcommutationclosure"]
  \inputleannode{InfoGeometry.Canonical.ConformalUnification.ConformalInference.ProjectorCommutationClosure}
\end{theorem}

\subsection*{casesOn}

\begin{theorem}[blueprint="infogeometry:canonical:conformalunification:conformalinference:caseson"]
  \inputleannode{InfoGeometry.Canonical.ConformalUnification.ConformalInference.casesOn}
\end{theorem}

\subsection*{chiralAnomaly}

\begin{theorem}[blueprint="infogeometry:canonical:conformalunification:conformalinference:chiralanomaly"]
  \inputleannode{InfoGeometry.Canonical.ConformalUnification.ConformalInference.chiralAnomaly}
\end{theorem}

\subsection*{chiralAnomalyOperator}

\begin{theorem}[blueprint="infogeometry:canonical:conformalunification:conformalinference:chiralanomalyoperator"]
  \inputleannode{InfoGeometry.Canonical.ConformalUnification.ConformalInference.chiralAnomalyOperator}
\end{theorem}

\subsection*{chiralAnomaly_eq_zero_iff_projectorCommutation}

\begin{theorem}[blueprint="infogeometry:canonical:conformalunification:conformalinference:chiralanomaly_eq_zero_iff_projectorcommutation"]
  \inputleannode{InfoGeometry.Canonical.ConformalUnification.ConformalInference.chiralAnomaly_eq_zero_iff_projectorCommutation}
\end{theorem}

\subsection*{chiralScale}

\begin{theorem}[blueprint="infogeometry:canonical:conformalunification:conformalinference:chiralscale"]
  \inputleannode{InfoGeometry.Canonical.ConformalUnification.ConformalInference.chiralScale}
\end{theorem}

\subsection*{chiral_commutation_link}

\begin{theorem}[blueprint="infogeometry:canonical:conformalunification:conformalinference:chiral_commutation_link"]
  \inputleannode{InfoGeometry.Canonical.ConformalUnification.ConformalInference.chiral_commutation_link}
\end{theorem}

\subsection*{ctorIdx}

\begin{theorem}[blueprint="infogeometry:canonical:conformalunification:conformalinference:ctoridx"]
  \inputleannode{InfoGeometry.Canonical.ConformalUnification.ConformalInference.ctorIdx}
\end{theorem}

\subsection*{epsilon}

\begin{theorem}[blueprint="infogeometry:canonical:conformalunification:conformalinference:epsilon"]
  \inputleannode{InfoGeometry.Canonical.ConformalUnification.ConformalInference.epsilon}
\end{theorem}

\subsection*{h_drazin}

\begin{theorem}[blueprint="infogeometry:canonical:conformalunification:conformalinference:h_drazin"]
  \inputleannode{InfoGeometry.Canonical.ConformalUnification.ConformalInference.h_drazin}
\end{theorem}

\subsection*{h_penrose}

\begin{theorem}[blueprint="infogeometry:canonical:conformalunification:conformalinference:h_penrose"]
  \inputleannode{InfoGeometry.Canonical.ConformalUnification.ConformalInference.h_penrose}
\end{theorem}

\subsection*{isChiralInference_iff_epsilon_pos}

\begin{theorem}[blueprint="infogeometry:canonical:conformalunification:conformalinference:ischiralinference_iff_epsilon_pos"]
  \inputleannode{InfoGeometry.Canonical.ConformalUnification.ConformalInference.isChiralInference_iff_epsilon_pos}
\end{theorem}

\subsection*{isNormalInference_iff_epsilon_eq_zero}

\begin{theorem}[blueprint="infogeometry:canonical:conformalunification:conformalinference:isnormalinference_iff_epsilon_eq_zero"]
  \inputleannode{InfoGeometry.Canonical.ConformalUnification.ConformalInference.isNormalInference_iff_epsilon_eq_zero}
\end{theorem}

\subsection*{metricChiralProjector}

\begin{theorem}[blueprint="infogeometry:canonical:conformalunification:conformalinference:metricchiralprojector"]
  \inputleannode{InfoGeometry.Canonical.ConformalUnification.ConformalInference.metricChiralProjector}
\end{theorem}

\subsection*{mk}

\begin{theorem}[blueprint="infogeometry:canonical:conformalunification:conformalinference:mk"]
  \inputleannode{InfoGeometry.Canonical.ConformalUnification.ConformalInference.mk}
\end{theorem}

\subsection*{noConfusion}

\begin{theorem}[blueprint="infogeometry:canonical:conformalunification:conformalinference:noconfusion"]
  \inputleannode{InfoGeometry.Canonical.ConformalUnification.ConformalInference.noConfusion}
\end{theorem}

\subsection*{noConfusionType}

\begin{theorem}[blueprint="infogeometry:canonical:conformalunification:conformalinference:noconfusiontype"]
  \inputleannode{InfoGeometry.Canonical.ConformalUnification.ConformalInference.noConfusionType}
\end{theorem}

\subsection*{rec}

\begin{theorem}[blueprint="infogeometry:canonical:conformalunification:conformalinference:rec"]
  \inputleannode{InfoGeometry.Canonical.ConformalUnification.ConformalInference.rec}
\end{theorem}

\subsection*{recOn}

\begin{theorem}[blueprint="infogeometry:canonical:conformalunification:conformalinference:recon"]
  \inputleannode{InfoGeometry.Canonical.ConformalUnification.ConformalInference.recOn}
\end{theorem}

\subsection*{spectralChiralProjector}

\begin{theorem}[blueprint="infogeometry:canonical:conformalunification:conformalinference:spectralchiralprojector"]
  \inputleannode{InfoGeometry.Canonical.ConformalUnification.ConformalInference.spectralChiralProjector}
\end{theorem}

\section{InfoGeometry.Canonical.ConformalUnification.ConformalInference.IsChiralInference}

\subsection*{eq_1}

\begin{theorem}[blueprint="infogeometry:canonical:conformalunification:conformalinference:ischiralinference:eq_1"]
  \inputleannode{InfoGeometry.Canonical.ConformalUnification.ConformalInference.IsChiralInference.eq_1}
\end{theorem}

\section{InfoGeometry.Canonical.ConformalUnification.ConformalInference.IsNormalInference}

\subsection*{eq_1}

\begin{theorem}[blueprint="infogeometry:canonical:conformalunification:conformalinference:isnormalinference:eq_1"]
  \inputleannode{InfoGeometry.Canonical.ConformalUnification.ConformalInference.IsNormalInference.eq_1}
\end{theorem}

\section{InfoGeometry.Canonical.ConformalUnification.ConformalInference.ProjectorCommutationClosure}

\subsection*{eq_1}

\begin{theorem}[blueprint="infogeometry:canonical:conformalunification:conformalinference:projectorcommutationclosure:eq_1"]
  \inputleannode{InfoGeometry.Canonical.ConformalUnification.ConformalInference.ProjectorCommutationClosure.eq_1}
\end{theorem}

\section{InfoGeometry.Canonical.ConformalUnification.ConformalInference.chiralAnomaly}

\subsection*{eq_1}

\begin{theorem}[blueprint="infogeometry:canonical:conformalunification:conformalinference:chiralanomaly:eq_1"]
  \inputleannode{InfoGeometry.Canonical.ConformalUnification.ConformalInference.chiralAnomaly.eq_1}
\end{theorem}

\section{InfoGeometry.Canonical.ConformalUnification.ConformalInference.epsilon}

\subsection*{eq_1}

\begin{theorem}[blueprint="infogeometry:canonical:conformalunification:conformalinference:epsilon:eq_1"]
  \inputleannode{InfoGeometry.Canonical.ConformalUnification.ConformalInference.epsilon.eq_1}
\end{theorem}

\section{InfoGeometry.Canonical.ConformalWard}

\noindent\textbf{Buildable}: yes

\subsection*{conformalVariation}

\begin{theorem}[blueprint="infogeometry:canonical:conformalward:conformalvariation"]
  \inputleannode{InfoGeometry.Canonical.ConformalWard.conformalVariation}
\end{theorem}

\subsection*{conformal_ward_identity}

\begin{theorem}[blueprint="infogeometry:canonical:conformalward:conformal_ward_identity"]
  \inputleannode{InfoGeometry.Canonical.ConformalWard.conformal_ward_identity}
\end{theorem}

\section{InfoGeometry.Canonical.CoreLegacy}

\noindent\textbf{Buildable}: yes

\noindent\emph{No exported declarations in the current declaration graph.}

\section{InfoGeometry.Canonical.CountSubstrateBridge}

\noindent\textbf{Buildable}: yes

\subsection*{CountSubstrate}

\begin{theorem}[blueprint="infogeometry:canonical:countsubstratebridge:countsubstrate"]
  \inputleannode{InfoGeometry.Canonical.CountSubstrateBridge.CountSubstrate}
\end{theorem}

\subsection*{CountSubstrateNontrivial}

\begin{theorem}[blueprint="infogeometry:canonical:countsubstratebridge:countsubstratenontrivial"]
  \inputleannode{InfoGeometry.Canonical.CountSubstrateBridge.CountSubstrateNontrivial}
\end{theorem}

\subsection*{EntrywisePositive}

\begin{theorem}[blueprint="infogeometry:canonical:countsubstratebridge:entrywisepositive"]
  \inputleannode{InfoGeometry.Canonical.CountSubstrateBridge.EntrywisePositive}
\end{theorem}

\subsection*{PositiveSinkhornState}

\begin{theorem}[blueprint="infogeometry:canonical:countsubstratebridge:positivesinkhornstate"]
  \inputleannode{InfoGeometry.Canonical.CountSubstrateBridge.PositiveSinkhornState}
\end{theorem}

\subsection*{colNormalize_entrywisePositive}

\begin{theorem}[blueprint="infogeometry:canonical:countsubstratebridge:colnormalize_entrywisepositive"]
  \inputleannode{InfoGeometry.Canonical.CountSubstrateBridge.colNormalize_entrywisePositive}
\end{theorem}

\subsection*{countInducedCoupling}

\begin{theorem}[blueprint="infogeometry:canonical:countsubstratebridge:countinducedcoupling"]
  \inputleannode{InfoGeometry.Canonical.CountSubstrateBridge.countInducedCoupling}
\end{theorem}

\subsection*{countInducedCoupling_entrywisePositive}

\begin{theorem}[blueprint="infogeometry:canonical:countsubstratebridge:countinducedcoupling_entrywisepositive"]
  \inputleannode{InfoGeometry.Canonical.CountSubstrateBridge.countInducedCoupling_entrywisePositive}
\end{theorem}

\subsection*{countInducedCoupling_hasPositiveColSums}

\begin{theorem}[blueprint="infogeometry:canonical:countsubstratebridge:countinducedcoupling_haspositivecolsums"]
  \inputleannode{InfoGeometry.Canonical.CountSubstrateBridge.countInducedCoupling_hasPositiveColSums}
\end{theorem}

\subsection*{countInducedCoupling_hasPositiveRowSums}

\begin{theorem}[blueprint="infogeometry:canonical:countsubstratebridge:countinducedcoupling_haspositiverowsums"]
  \inputleannode{InfoGeometry.Canonical.CountSubstrateBridge.countInducedCoupling_hasPositiveRowSums}
\end{theorem}

\subsection*{countInducedIterate}

\begin{theorem}[blueprint="infogeometry:canonical:countsubstratebridge:countinducediterate"]
  \inputleannode{InfoGeometry.Canonical.CountSubstrateBridge.countInducedIterate}
\end{theorem}

\subsection*{countInducedIterate_entrywisePositive}

\begin{theorem}[blueprint="infogeometry:canonical:countsubstratebridge:countinducediterate_entrywisepositive"]
  \inputleannode{InfoGeometry.Canonical.CountSubstrateBridge.countInducedIterate_entrywisePositive}
\end{theorem}

\subsection*{countInducedIterate_step}

\begin{theorem}[blueprint="infogeometry:canonical:countsubstratebridge:countinducediterate_step"]
  \inputleannode{InfoGeometry.Canonical.CountSubstrateBridge.countInducedIterate_step}
\end{theorem}

\subsection*{countInducedPositiveIterate}

\begin{theorem}[blueprint="infogeometry:canonical:countsubstratebridge:countinducedpositiveiterate"]
  \inputleannode{InfoGeometry.Canonical.CountSubstrateBridge.countInducedPositiveIterate}
\end{theorem}

\subsection*{countInducedSinkhornTrajectory}

\begin{theorem}[blueprint="infogeometry:canonical:countsubstratebridge:countinducedsinkhorntrajectory"]
  \inputleannode{InfoGeometry.Canonical.CountSubstrateBridge.countInducedSinkhornTrajectory}
\end{theorem}

\subsection*{countWeight}

\begin{theorem}[blueprint="infogeometry:canonical:countsubstratebridge:countweight"]
  \inputleannode{InfoGeometry.Canonical.CountSubstrateBridge.countWeight}
\end{theorem}

\subsection*{countWeight_pos}

\begin{theorem}[blueprint="infogeometry:canonical:countsubstratebridge:countweight_pos"]
  \inputleannode{InfoGeometry.Canonical.CountSubstrateBridge.countWeight_pos}
\end{theorem}

\subsection*{countsFirst_holographicEmergence_of_countInducedTrajectory}

\begin{theorem}[blueprint="infogeometry:canonical:countsubstratebridge:countsfirst_holographicemergence_of_countinducedtrajectory"]
  \inputleannode{InfoGeometry.Canonical.CountSubstrateBridge.countsFirst_holographicEmergence_of_countInducedTrajectory}
\end{theorem}

\subsection*{countsFirst_holographicEmergence_package}

\begin{theorem}[blueprint="infogeometry:canonical:countsubstratebridge:countsfirst_holographicemergence_package"]
  \inputleannode{InfoGeometry.Canonical.CountSubstrateBridge.countsFirst_holographicEmergence_package}
\end{theorem}

\subsection*{emergentTimeFlow_countInducedSinkhornTrajectory}

\begin{theorem}[blueprint="infogeometry:canonical:countsubstratebridge:emergenttimeflow_countinducedsinkhorntrajectory"]
  \inputleannode{InfoGeometry.Canonical.CountSubstrateBridge.emergentTimeFlow_countInducedSinkhornTrajectory}
\end{theorem}

\subsection*{empiricalProbabilityState}

\begin{theorem}[blueprint="infogeometry:canonical:countsubstratebridge:empiricalprobabilitystate"]
  \inputleannode{InfoGeometry.Canonical.CountSubstrateBridge.empiricalProbabilityState}
\end{theorem}

\subsection*{empiricalProbabilityState_spec}

\begin{theorem}[blueprint="infogeometry:canonical:countsubstratebridge:empiricalprobabilitystate_spec"]
  \inputleannode{InfoGeometry.Canonical.CountSubstrateBridge.empiricalProbabilityState_spec}
\end{theorem}

\subsection*{entrywisePositive_hasPositiveColSums}

\begin{theorem}[blueprint="infogeometry:canonical:countsubstratebridge:entrywisepositive_haspositivecolsums"]
  \inputleannode{InfoGeometry.Canonical.CountSubstrateBridge.entrywisePositive_hasPositiveColSums}
\end{theorem}

\subsection*{entrywisePositive_hasPositiveRowSums}

\begin{theorem}[blueprint="infogeometry:canonical:countsubstratebridge:entrywisepositive_haspositiverowsums"]
  \inputleannode{InfoGeometry.Canonical.CountSubstrateBridge.entrywisePositive_hasPositiveRowSums}
\end{theorem}

\subsection*{exists_empiricalProbabilityState}

\begin{theorem}[blueprint="infogeometry:canonical:countsubstratebridge:exists_empiricalprobabilitystate"]
  \inputleannode{InfoGeometry.Canonical.CountSubstrateBridge.exists_empiricalProbabilityState}
\end{theorem}

\subsection*{rowNormalize_entrywisePositive}

\begin{theorem}[blueprint="infogeometry:canonical:countsubstratebridge:rownormalize_entrywisepositive"]
  \inputleannode{InfoGeometry.Canonical.CountSubstrateBridge.rowNormalize_entrywisePositive}
\end{theorem}

\section{InfoGeometry.Canonical.CountSubstrateBridge.PositiveSinkhornState}

\subsection*{M}

\begin{theorem}[blueprint="infogeometry:canonical:countsubstratebridge:positivesinkhornstate:m"]
  \inputleannode{InfoGeometry.Canonical.CountSubstrateBridge.PositiveSinkhornState.M}
\end{theorem}

\subsection*{casesOn}

\begin{theorem}[blueprint="infogeometry:canonical:countsubstratebridge:positivesinkhornstate:caseson"]
  \inputleannode{InfoGeometry.Canonical.CountSubstrateBridge.PositiveSinkhornState.casesOn}
\end{theorem}

\subsection*{ctorIdx}

\begin{theorem}[blueprint="infogeometry:canonical:countsubstratebridge:positivesinkhornstate:ctoridx"]
  \inputleannode{InfoGeometry.Canonical.CountSubstrateBridge.PositiveSinkhornState.ctorIdx}
\end{theorem}

\subsection*{mk}

\begin{theorem}[blueprint="infogeometry:canonical:countsubstratebridge:positivesinkhornstate:mk"]
  \inputleannode{InfoGeometry.Canonical.CountSubstrateBridge.PositiveSinkhornState.mk}
\end{theorem}

\subsection*{noConfusion}

\begin{theorem}[blueprint="infogeometry:canonical:countsubstratebridge:positivesinkhornstate:noconfusion"]
  \inputleannode{InfoGeometry.Canonical.CountSubstrateBridge.PositiveSinkhornState.noConfusion}
\end{theorem}

\subsection*{noConfusionType}

\begin{theorem}[blueprint="infogeometry:canonical:countsubstratebridge:positivesinkhornstate:noconfusiontype"]
  \inputleannode{InfoGeometry.Canonical.CountSubstrateBridge.PositiveSinkhornState.noConfusionType}
\end{theorem}

\subsection*{pos}

\begin{theorem}[blueprint="infogeometry:canonical:countsubstratebridge:positivesinkhornstate:pos"]
  \inputleannode{InfoGeometry.Canonical.CountSubstrateBridge.PositiveSinkhornState.pos}
\end{theorem}

\subsection*{rec}

\begin{theorem}[blueprint="infogeometry:canonical:countsubstratebridge:positivesinkhornstate:rec"]
  \inputleannode{InfoGeometry.Canonical.CountSubstrateBridge.PositiveSinkhornState.rec}
\end{theorem}

\subsection*{recOn}

\begin{theorem}[blueprint="infogeometry:canonical:countsubstratebridge:positivesinkhornstate:recon"]
  \inputleannode{InfoGeometry.Canonical.CountSubstrateBridge.PositiveSinkhornState.recOn}
\end{theorem}

\section{InfoGeometry.Canonical.CountSubstrateBridge.countInducedIterate}

\subsection*{eq_1}

\begin{theorem}[blueprint="infogeometry:canonical:countsubstratebridge:countinducediterate:eq_1"]
  \inputleannode{InfoGeometry.Canonical.CountSubstrateBridge.countInducedIterate.eq_1}
\end{theorem}

\section{InfoGeometry.Canonical.CountSubstrateBridge.countInducedPositiveIterate}

\subsection*{eq_1}

\begin{theorem}[blueprint="infogeometry:canonical:countsubstratebridge:countinducedpositiveiterate:eq_1"]
  \inputleannode{InfoGeometry.Canonical.CountSubstrateBridge.countInducedPositiveIterate.eq_1}
\end{theorem}

\subsection*{eq_2}

\begin{theorem}[blueprint="infogeometry:canonical:countsubstratebridge:countinducedpositiveiterate:eq_2"]
  \inputleannode{InfoGeometry.Canonical.CountSubstrateBridge.countInducedPositiveIterate.eq_2}
\end{theorem}

\subsection*{eq_def}

\begin{theorem}[blueprint="infogeometry:canonical:countsubstratebridge:countinducedpositiveiterate:eq_def"]
  \inputleannode{InfoGeometry.Canonical.CountSubstrateBridge.countInducedPositiveIterate.eq_def}
\end{theorem}

\section{InfoGeometry.Canonical.CurvatureRGFlow}

\noindent\textbf{Buildable}: yes

\subsection*{ScalarCurvatureFlow}

\begin{theorem}[blueprint="infogeometry:canonical:curvaturergflow:scalarcurvatureflow"]
  \inputleannode{InfoGeometry.Canonical.CurvatureRGFlow.ScalarCurvatureFlow}
\end{theorem}

\subsection*{curvatureBetaFunction}

\begin{theorem}[blueprint="infogeometry:canonical:curvaturergflow:curvaturebetafunction"]
  \inputleannode{InfoGeometry.Canonical.CurvatureRGFlow.curvatureBetaFunction}
\end{theorem}

\subsection*{curvature_invariant_at_fixed_point}

\begin{theorem}[blueprint="infogeometry:canonical:curvaturergflow:curvature_invariant_at_fixed_point"]
  \inputleannode{InfoGeometry.Canonical.CurvatureRGFlow.curvature_invariant_at_fixed_point}
\end{theorem}

\subsection*{einsteinHilbertFlow}

\begin{theorem}[blueprint="infogeometry:canonical:curvaturergflow:einsteinhilbertflow"]
  \inputleannode{InfoGeometry.Canonical.CurvatureRGFlow.einsteinHilbertFlow}
\end{theorem}

\section{InfoGeometry.Canonical.Drazin}

\noindent\textbf{Buildable}: yes

\subsection*{IsDrazinInverse}

\begin{theorem}[blueprint="infogeometry:canonical:drazin:isdrazininverse"]
  \inputleannode{InfoGeometry.Canonical.Drazin.IsDrazinInverse}
\end{theorem}

\section{InfoGeometry.Canonical.Drazin.IsDrazinInverse}

\subsection*{casesOn}

\begin{theorem}[blueprint="infogeometry:canonical:drazin:isdrazininverse:caseson"]
  \inputleannode{InfoGeometry.Canonical.Drazin.IsDrazinInverse.casesOn}
\end{theorem}

\subsection*{comm}

\begin{theorem}[blueprint="infogeometry:canonical:drazin:isdrazininverse:comm"]
  \inputleannode{InfoGeometry.Canonical.Drazin.IsDrazinInverse.comm}
\end{theorem}

\subsection*{core}

\begin{theorem}[blueprint="infogeometry:canonical:drazin:isdrazininverse:core"]
  \inputleannode{InfoGeometry.Canonical.Drazin.IsDrazinInverse.core}
\end{theorem}

\subsection*{fitting_decomposition}

\begin{theorem}[blueprint="infogeometry:canonical:drazin:isdrazininverse:fitting_decomposition"]
  \inputleannode{InfoGeometry.Canonical.Drazin.IsDrazinInverse.fitting_decomposition}
\end{theorem}

\subsection*{idempotent}

\begin{theorem}[blueprint="infogeometry:canonical:drazin:isdrazininverse:idempotent"]
  \inputleannode{InfoGeometry.Canonical.Drazin.IsDrazinInverse.idempotent}
\end{theorem}

\subsection*{inverse_eq_pow_mul_pow}

\begin{theorem}[blueprint="infogeometry:canonical:drazin:isdrazininverse:inverse_eq_pow_mul_pow"]
  \inputleannode{InfoGeometry.Canonical.Drazin.IsDrazinInverse.inverse_eq_pow_mul_pow}
\end{theorem}

\subsection*{mk}

\begin{theorem}[blueprint="infogeometry:canonical:drazin:isdrazininverse:mk"]
  \inputleannode{InfoGeometry.Canonical.Drazin.IsDrazinInverse.mk}
\end{theorem}

\subsection*{nilpotent}

\begin{theorem}[blueprint="infogeometry:canonical:drazin:isdrazininverse:nilpotent"]
  \inputleannode{InfoGeometry.Canonical.Drazin.IsDrazinInverse.nilpotent}
\end{theorem}

\subsection*{nilpotent_comm_self}

\begin{theorem}[blueprint="infogeometry:canonical:drazin:isdrazininverse:nilpotent_comm_self"]
  \inputleannode{InfoGeometry.Canonical.Drazin.IsDrazinInverse.nilpotent_comm_self}
\end{theorem}

\subsection*{power}

\begin{theorem}[blueprint="infogeometry:canonical:drazin:isdrazininverse:power"]
  \inputleannode{InfoGeometry.Canonical.Drazin.IsDrazinInverse.power}
\end{theorem}

\subsection*{power_le}

\begin{theorem}[blueprint="infogeometry:canonical:drazin:isdrazininverse:power_le"]
  \inputleannode{InfoGeometry.Canonical.Drazin.IsDrazinInverse.power_le}
\end{theorem}

\subsection*{projection}

\begin{theorem}[blueprint="infogeometry:canonical:drazin:isdrazininverse:projection"]
  \inputleannode{InfoGeometry.Canonical.Drazin.IsDrazinInverse.projection}
\end{theorem}

\subsection*{projection_comm}

\begin{theorem}[blueprint="infogeometry:canonical:drazin:isdrazininverse:projection_comm"]
  \inputleannode{InfoGeometry.Canonical.Drazin.IsDrazinInverse.projection_comm}
\end{theorem}

\subsection*{projection_is_idempotent}

\begin{theorem}[blueprint="infogeometry:canonical:drazin:isdrazininverse:projection_is_idempotent"]
  \inputleannode{InfoGeometry.Canonical.Drazin.IsDrazinInverse.projection_is_idempotent}
\end{theorem}

\subsection*{rec}

\begin{theorem}[blueprint="infogeometry:canonical:drazin:isdrazininverse:rec"]
  \inputleannode{InfoGeometry.Canonical.Drazin.IsDrazinInverse.rec}
\end{theorem}

\subsection*{recOn}

\begin{theorem}[blueprint="infogeometry:canonical:drazin:isdrazininverse:recon"]
  \inputleannode{InfoGeometry.Canonical.Drazin.IsDrazinInverse.recOn}
\end{theorem}

\section{InfoGeometry.Canonical.DualConnections}

\noindent\textbf{Buildable}: yes

\noindent\emph{No exported declarations in the current declaration graph.}

\section{InfoGeometry.Canonical.ExponentialFamilyLegacy}

\noindent\textbf{Buildable}: yes

\noindent\emph{No exported declarations in the current declaration graph.}

\section{InfoGeometry.Canonical.FormalScaffold}

\noindent\textbf{Buildable}: yes

\noindent\emph{No exported declarations in the current declaration graph.}

\section{InfoGeometry.Canonical.Foundations}

\noindent\textbf{Buildable}: yes

\noindent\emph{No exported declarations in the current declaration graph.}

\section{InfoGeometry.Canonical.FoundationsLegacy}

\noindent\textbf{Buildable}: yes

\noindent\emph{No exported declarations in the current declaration graph.}

\section{InfoGeometry.Canonical.Gauge}

\subsection*{Plane}

\begin{theorem}[blueprint="infogeometry:canonical:gauge:plane"]
  \inputleannode{InfoGeometry.Canonical.Gauge.Plane}
\end{theorem}

\subsection*{Signature}

\begin{theorem}[blueprint="infogeometry:canonical:gauge:signature"]
  \inputleannode{InfoGeometry.Canonical.Gauge.Signature}
\end{theorem}

\subsection*{act}

\begin{theorem}[blueprint="infogeometry:canonical:gauge:act"]
  \inputleannode{InfoGeometry.Canonical.Gauge.act}
\end{theorem}

\subsection*{act_preserves_bilinear}

\begin{theorem}[blueprint="infogeometry:canonical:gauge:act_preserves_bilinear"]
  \inputleannode{InfoGeometry.Canonical.Gauge.act_preserves_bilinear}
\end{theorem}

\subsection*{bilinear}

\begin{theorem}[blueprint="infogeometry:canonical:gauge:bilinear"]
  \inputleannode{InfoGeometry.Canonical.Gauge.bilinear}
\end{theorem}

\subsection*{boost}

\begin{theorem}[blueprint="infogeometry:canonical:gauge:boost"]
  \inputleannode{InfoGeometry.Canonical.Gauge.boost}
\end{theorem}

\subsection*{boost_preserves_bilinear}

\begin{theorem}[blueprint="infogeometry:canonical:gauge:boost_preserves_bilinear"]
  \inputleannode{InfoGeometry.Canonical.Gauge.boost_preserves_bilinear}
\end{theorem}

\subsection*{quad}

\begin{theorem}[blueprint="infogeometry:canonical:gauge:quad"]
  \inputleannode{InfoGeometry.Canonical.Gauge.quad}
\end{theorem}

\subsection*{rot}

\begin{theorem}[blueprint="infogeometry:canonical:gauge:rot"]
  \inputleannode{InfoGeometry.Canonical.Gauge.rot}
\end{theorem}

\subsection*{rot_preserves_bilinear}

\begin{theorem}[blueprint="infogeometry:canonical:gauge:rot_preserves_bilinear"]
  \inputleannode{InfoGeometry.Canonical.Gauge.rot_preserves_bilinear}
\end{theorem}

\section{InfoGeometry.Canonical.Gauge.Signature}

\subsection*{casesOn}

\begin{theorem}[blueprint="infogeometry:canonical:gauge:signature:caseson"]
  \inputleannode{InfoGeometry.Canonical.Gauge.Signature.casesOn}
\end{theorem}

\subsection*{ctorElim}

\begin{theorem}[blueprint="infogeometry:canonical:gauge:signature:ctorelim"]
  \inputleannode{InfoGeometry.Canonical.Gauge.Signature.ctorElim}
\end{theorem}

\subsection*{ctorElimType}

\begin{theorem}[blueprint="infogeometry:canonical:gauge:signature:ctorelimtype"]
  \inputleannode{InfoGeometry.Canonical.Gauge.Signature.ctorElimType}
\end{theorem}

\subsection*{ctorIdx}

\begin{theorem}[blueprint="infogeometry:canonical:gauge:signature:ctoridx"]
  \inputleannode{InfoGeometry.Canonical.Gauge.Signature.ctorIdx}
\end{theorem}

\subsection*{euclid}

\begin{theorem}[blueprint="infogeometry:canonical:gauge:signature:euclid"]
  \inputleannode{InfoGeometry.Canonical.Gauge.Signature.euclid}
\end{theorem}

\subsection*{noConfusion}

\begin{theorem}[blueprint="infogeometry:canonical:gauge:signature:noconfusion"]
  \inputleannode{InfoGeometry.Canonical.Gauge.Signature.noConfusion}
\end{theorem}

\subsection*{noConfusionType}

\begin{theorem}[blueprint="infogeometry:canonical:gauge:signature:noconfusiontype"]
  \inputleannode{InfoGeometry.Canonical.Gauge.Signature.noConfusionType}
\end{theorem}

\subsection*{rec}

\begin{theorem}[blueprint="infogeometry:canonical:gauge:signature:rec"]
  \inputleannode{InfoGeometry.Canonical.Gauge.Signature.rec}
\end{theorem}

\subsection*{recOn}

\begin{theorem}[blueprint="infogeometry:canonical:gauge:signature:recon"]
  \inputleannode{InfoGeometry.Canonical.Gauge.Signature.recOn}
\end{theorem}

\subsection*{split}

\begin{theorem}[blueprint="infogeometry:canonical:gauge:signature:split"]
  \inputleannode{InfoGeometry.Canonical.Gauge.Signature.split}
\end{theorem}

\subsection*{toCtorIdx}

\begin{theorem}[blueprint="infogeometry:canonical:gauge:signature:toctoridx"]
  \inputleannode{InfoGeometry.Canonical.Gauge.Signature.toCtorIdx}
\end{theorem}

\section{InfoGeometry.Canonical.Gauge.Signature.euclid}

\subsection*{elim}

\begin{theorem}[blueprint="infogeometry:canonical:gauge:signature:euclid:elim"]
  \inputleannode{InfoGeometry.Canonical.Gauge.Signature.euclid.elim}
\end{theorem}

\section{InfoGeometry.Canonical.Gauge.Signature.split}

\subsection*{elim}

\begin{theorem}[blueprint="infogeometry:canonical:gauge:signature:split:elim"]
  \inputleannode{InfoGeometry.Canonical.Gauge.Signature.split.elim}
\end{theorem}

\section{InfoGeometry.Canonical.Gauge.quad}

\subsection*{eq_1}

\begin{theorem}[blueprint="infogeometry:canonical:gauge:quad:eq_1"]
  \inputleannode{InfoGeometry.Canonical.Gauge.quad.eq_1}
\end{theorem}

\subsection*{eq_2}

\begin{theorem}[blueprint="infogeometry:canonical:gauge:quad:eq_2"]
  \inputleannode{InfoGeometry.Canonical.Gauge.quad.eq_2}
\end{theorem}

\section{InfoGeometry.Canonical.GaugeUnified}

\noindent\textbf{Buildable}: yes

\noindent\emph{No exported declarations in the current declaration graph.}

\section{InfoGeometry.Canonical.GeneralizedKL}

\noindent\textbf{Buildable}: yes

\noindent\emph{No exported declarations in the current declaration graph.}

\section{InfoGeometry.Canonical.Geometry}

\noindent\textbf{Buildable}: yes

\noindent\emph{No exported declarations in the current declaration graph.}

\section{InfoGeometry.Canonical.GrandCanonicalExperts}

\noindent\textbf{Buildable}: yes

\noindent\emph{No exported declarations in the current declaration graph.}

\section{InfoGeometry.Canonical.GrandSynthesis}

\noindent\textbf{Buildable}: yes

\subsection*{AlgebraicEquilibriumCl11}

\begin{theorem}[blueprint="infogeometry:canonical:grandsynthesis:algebraicequilibriumcl11"]
  \inputleannode{InfoGeometry.Canonical.GrandSynthesis.AlgebraicEquilibriumCl11}
\end{theorem}

\subsection*{GeometricAlgebraicFromThermodynamic}

\begin{theorem}[blueprint="infogeometry:canonical:grandsynthesis:geometricalgebraicfromthermodynamic"]
  \inputleannode{InfoGeometry.Canonical.GrandSynthesis.GeometricAlgebraicFromThermodynamic}
\end{theorem}

\subsection*{GeometricEquilibrium}

\begin{theorem}[blueprint="infogeometry:canonical:grandsynthesis:geometricequilibrium"]
  \inputleannode{InfoGeometry.Canonical.GrandSynthesis.GeometricEquilibrium}
\end{theorem}

\subsection*{LichnerowiczBalancedCl11}

\begin{theorem}[blueprint="infogeometry:canonical:grandsynthesis:lichnerowiczbalancedcl11"]
  \inputleannode{InfoGeometry.Canonical.GrandSynthesis.LichnerowiczBalancedCl11}
\end{theorem}

\subsection*{RNEntropySourcesMongeAmpere}

\begin{theorem}[blueprint="infogeometry:canonical:grandsynthesis:rnentropysourcesmongeampere"]
  \inputleannode{InfoGeometry.Canonical.GrandSynthesis.RNEntropySourcesMongeAmpere}
\end{theorem}

\subsection*{RelativeCounts}

\begin{theorem}[blueprint="infogeometry:canonical:grandsynthesis:relativecounts"]
  \inputleannode{InfoGeometry.Canonical.GrandSynthesis.RelativeCounts}
\end{theorem}

\subsection*{SinkhornEntropyMonotoneRN}

\begin{theorem}[blueprint="infogeometry:canonical:grandsynthesis:sinkhornentropymonotonern"]
  \inputleannode{InfoGeometry.Canonical.GrandSynthesis.SinkhornEntropyMonotoneRN}
\end{theorem}

\subsection*{ThermodynamicEquilibrium}

\begin{theorem}[blueprint="infogeometry:canonical:grandsynthesis:thermodynamicequilibrium"]
  \inputleannode{InfoGeometry.Canonical.GrandSynthesis.ThermodynamicEquilibrium}
\end{theorem}

\subsection*{ThermodynamicFromGeometricAlgebraic}

\begin{theorem}[blueprint="infogeometry:canonical:grandsynthesis:thermodynamicfromgeometricalgebraic"]
  \inputleannode{InfoGeometry.Canonical.GrandSynthesis.ThermodynamicFromGeometricAlgebraic}
\end{theorem}

\subsection*{bochnerWeitzenboeckBridge_of_sinkhornDrive}

\begin{theorem}[blueprint="infogeometry:canonical:grandsynthesis:bochnerweitzenboeckbridge_of_sinkhorndrive"]
  \inputleannode{InfoGeometry.Canonical.GrandSynthesis.bochnerWeitzenboeckBridge_of_sinkhornDrive}
\end{theorem}

\subsection*{bochnerWeitzenboeckBridge_of_sinkhornDrive_and_constantMongeAmpere}

\begin{theorem}[blueprint="infogeometry:canonical:grandsynthesis:bochnerweitzenboeckbridge_of_sinkhorndrive_and_constantmongeampere"]
  \inputleannode{InfoGeometry.Canonical.GrandSynthesis.bochnerWeitzenboeckBridge_of_sinkhornDrive_and_constantMongeAmpere}
\end{theorem}

\subsection*{calabiYauEntropyBridge_of_sinkhornRicciIndexHypotheses}

\begin{theorem}[blueprint="infogeometry:canonical:grandsynthesis:calabiyauentropybridge_of_sinkhornricciindexhypotheses"]
  \inputleannode{InfoGeometry.Canonical.GrandSynthesis.calabiYauEntropyBridge_of_sinkhornRicciIndexHypotheses}
\end{theorem}

\subsection*{cl11_bottDirac_sq_apply_eq_zero_of_algebraicEquilibrium}

\begin{theorem}[blueprint="infogeometry:canonical:grandsynthesis:cl11_bottdirac_sq_apply_eq_zero_of_algebraicequilibrium"]
  \inputleannode{InfoGeometry.Canonical.GrandSynthesis.cl11_bottDirac_sq_apply_eq_zero_of_algebraicEquilibrium}
\end{theorem}

\subsection*{cl11_bottDirac_sq_eq_metricOp_form}

\begin{theorem}[blueprint="infogeometry:canonical:grandsynthesis:cl11_bottdirac_sq_eq_metricop_form"]
  \inputleannode{InfoGeometry.Canonical.GrandSynthesis.cl11_bottDirac_sq_eq_metricOp_form}
\end{theorem}

\subsection*{cl11_bottDirac_sq_eq_zero_of_algebraicEquilibrium}

\begin{theorem}[blueprint="infogeometry:canonical:grandsynthesis:cl11_bottdirac_sq_eq_zero_of_algebraicequilibrium"]
  \inputleannode{InfoGeometry.Canonical.GrandSynthesis.cl11_bottDirac_sq_eq_zero_of_algebraicEquilibrium}
\end{theorem}

\subsection*{cl11_bottDirac_sq_eq_zero_of_lichnerowiczBalanced}

\begin{theorem}[blueprint="infogeometry:canonical:grandsynthesis:cl11_bottdirac_sq_eq_zero_of_lichnerowiczbalanced"]
  \inputleannode{InfoGeometry.Canonical.GrandSynthesis.cl11_bottDirac_sq_eq_zero_of_lichnerowiczBalanced}
\end{theorem}

\subsection*{doublyStochastic_sinkhornEntropyMonotoneRN}

\begin{theorem}[blueprint="infogeometry:canonical:grandsynthesis:doublystochastic_sinkhornentropymonotonern"]
  \inputleannode{InfoGeometry.Canonical.GrandSynthesis.doublyStochastic_sinkhornEntropyMonotoneRN}
\end{theorem}

\subsection*{geometricAlgebraicState_of_sinkhornRicciIndexHypotheses}

\begin{theorem}[blueprint="infogeometry:canonical:grandsynthesis:geometricalgebraicstate_of_sinkhornricciindexhypotheses"]
  \inputleannode{InfoGeometry.Canonical.GrandSynthesis.geometricAlgebraicState_of_sinkhornRicciIndexHypotheses}
\end{theorem}

\subsection*{grandSynthesis_of_axioms}

\begin{theorem}[blueprint="infogeometry:canonical:grandsynthesis:grandsynthesis_of_axioms"]
  \inputleannode{InfoGeometry.Canonical.GrandSynthesis.grandSynthesis_of_axioms}
\end{theorem}

\subsection*{grandSynthesis_step}

\begin{theorem}[blueprint="infogeometry:canonical:grandsynthesis:grandsynthesis_step"]
  \inputleannode{InfoGeometry.Canonical.GrandSynthesis.grandSynthesis_step}
\end{theorem}

\subsection*{gravity_generated_by_rnEntropy}

\begin{theorem}[blueprint="infogeometry:canonical:grandsynthesis:gravity_generated_by_rnentropy"]
  \inputleannode{InfoGeometry.Canonical.GrandSynthesis.gravity_generated_by_rnEntropy}
\end{theorem}

\subsection*{hasConstantMongeAmpereDensity_of_rnEntropySource}

\begin{theorem}[blueprint="infogeometry:canonical:grandsynthesis:hasconstantmongeamperedensity_of_rnentropysource"]
  \inputleannode{InfoGeometry.Canonical.GrandSynthesis.hasConstantMongeAmpereDensity_of_rnEntropySource}
\end{theorem}

\subsection*{information_wheeler_dewitt_equivalence}

\begin{theorem}[blueprint="infogeometry:canonical:grandsynthesis:information_wheeler_dewitt_equivalence"]
  \inputleannode{InfoGeometry.Canonical.GrandSynthesis.information_wheeler_dewitt_equivalence}
\end{theorem}

\subsection*{information_wheeler_dewitt_equivalence_of_constructive_hypotheses}

\begin{theorem}[blueprint="infogeometry:canonical:grandsynthesis:information_wheeler_dewitt_equivalence_of_constructive_hypotheses"]
  \inputleannode{InfoGeometry.Canonical.GrandSynthesis.information_wheeler_dewitt_equivalence_of_constructive_hypotheses}
\end{theorem}

\subsection*{information_wheeler_dewitt_equivalence_of_sinkhornDrive}

\begin{theorem}[blueprint="infogeometry:canonical:grandsynthesis:information_wheeler_dewitt_equivalence_of_sinkhorndrive"]
  \inputleannode{InfoGeometry.Canonical.GrandSynthesis.information_wheeler_dewitt_equivalence_of_sinkhornDrive}
\end{theorem}

\subsection*{information_wheeler_dewitt_equivalence_of_sinkhornDrive_and_calabiYau}

\begin{theorem}[blueprint="infogeometry:canonical:grandsynthesis:information_wheeler_dewitt_equivalence_of_sinkhorndrive_and_calabiyau"]
  \inputleannode{InfoGeometry.Canonical.GrandSynthesis.information_wheeler_dewitt_equivalence_of_sinkhornDrive_and_calabiYau}
\end{theorem}

\subsection*{information_wheeler_dewitt_equivalence_of_sinkhornDrive_and_indexHypotheses}

\begin{theorem}[blueprint="infogeometry:canonical:grandsynthesis:information_wheeler_dewitt_equivalence_of_sinkhorndrive_and_indexhypotheses"]
  \inputleannode{InfoGeometry.Canonical.GrandSynthesis.information_wheeler_dewitt_equivalence_of_sinkhornDrive_and_indexHypotheses}
\end{theorem}

\subsection*{isRicciFlat_of_rnEntropySource}

\begin{theorem}[blueprint="infogeometry:canonical:grandsynthesis:isricciflat_of_rnentropysource"]
  \inputleannode{InfoGeometry.Canonical.GrandSynthesis.isRicciFlat_of_rnEntropySource}
\end{theorem}

\subsection*{jacDetCLM}

\begin{theorem}[blueprint="infogeometry:canonical:grandsynthesis:jacdetclm"]
  \inputleannode{InfoGeometry.Canonical.GrandSynthesis.jacDetCLM}
\end{theorem}

\subsection*{jacDetCLM_comp}

\begin{theorem}[blueprint="infogeometry:canonical:grandsynthesis:jacdetclm_comp"]
  \inputleannode{InfoGeometry.Canonical.GrandSynthesis.jacDetCLM_comp}
\end{theorem}

\subsection*{kahlerPotentialRN}

\begin{theorem}[blueprint="infogeometry:canonical:grandsynthesis:kahlerpotentialrn"]
  \inputleannode{InfoGeometry.Canonical.GrandSynthesis.kahlerPotentialRN}
\end{theorem}

\subsection*{kahlerPotentialRN_eq_neg_logJacobian}

\begin{theorem}[blueprint="infogeometry:canonical:grandsynthesis:kahlerpotentialrn_eq_neg_logjacobian"]
  \inputleannode{InfoGeometry.Canonical.GrandSynthesis.kahlerPotentialRN_eq_neg_logJacobian}
\end{theorem}

\subsection*{logAbsDetMatrix}

\begin{theorem}[blueprint="infogeometry:canonical:grandsynthesis:logabsdetmatrix"]
  \inputleannode{InfoGeometry.Canonical.GrandSynthesis.logAbsDetMatrix}
\end{theorem}

\subsection*{logAbsDetMatrix_kronecker}

\begin{theorem}[blueprint="infogeometry:canonical:grandsynthesis:logabsdetmatrix_kronecker"]
  \inputleannode{InfoGeometry.Canonical.GrandSynthesis.logAbsDetMatrix_kronecker}
\end{theorem}

\subsection*{logAbsDetMatrix_mul}

\begin{theorem}[blueprint="infogeometry:canonical:grandsynthesis:logabsdetmatrix_mul"]
  \inputleannode{InfoGeometry.Canonical.GrandSynthesis.logAbsDetMatrix_mul}
\end{theorem}

\subsection*{logAbsDet_metricModel_eq_trace_metricOp}

\begin{theorem}[blueprint="infogeometry:canonical:grandsynthesis:logabsdet_metricmodel_eq_trace_metricop"]
  \inputleannode{InfoGeometry.Canonical.GrandSynthesis.logAbsDet_metricModel_eq_trace_metricOp}
\end{theorem}

\subsection*{logAbsDet_spectralModel_eq_spectralVolume}

\begin{theorem}[blueprint="infogeometry:canonical:grandsynthesis:logabsdet_spectralmodel_eq_spectralvolume"]
  \inputleannode{InfoGeometry.Canonical.GrandSynthesis.logAbsDet_spectralModel_eq_spectralVolume}
\end{theorem}

\subsection*{logAbsJacDetCLM}

\begin{theorem}[blueprint="infogeometry:canonical:grandsynthesis:logabsjacdetclm"]
  \inputleannode{InfoGeometry.Canonical.GrandSynthesis.logAbsJacDetCLM}
\end{theorem}

\subsection*{logAbsJacDetCLM_comp}

\begin{theorem}[blueprint="infogeometry:canonical:grandsynthesis:logabsjacdetclm_comp"]
  \inputleannode{InfoGeometry.Canonical.GrandSynthesis.logAbsJacDetCLM_comp}
\end{theorem}

\subsection*{log_mongeAmpereDensity_eq_trace_metricOp}

\begin{theorem}[blueprint="infogeometry:canonical:grandsynthesis:log_mongeamperedensity_eq_trace_metricop"]
  \inputleannode{InfoGeometry.Canonical.GrandSynthesis.log_mongeAmpereDensity_eq_trace_metricOp}
\end{theorem}

\subsection*{log_spectralMongeAmpereDensity_eq_spectralVolume}

\begin{theorem}[blueprint="infogeometry:canonical:grandsynthesis:log_spectralmongeamperedensity_eq_spectralvolume"]
  \inputleannode{InfoGeometry.Canonical.GrandSynthesis.log_spectralMongeAmpereDensity_eq_spectralVolume}
\end{theorem}

\subsection*{relativeCountDensity}

\begin{theorem}[blueprint="infogeometry:canonical:grandsynthesis:relativecountdensity"]
  \inputleannode{InfoGeometry.Canonical.GrandSynthesis.relativeCountDensity}
\end{theorem}

\subsection*{relativeCountDensity_eq_rn_lift}

\begin{theorem}[blueprint="infogeometry:canonical:grandsynthesis:relativecountdensity_eq_rn_lift"]
  \inputleannode{InfoGeometry.Canonical.GrandSynthesis.relativeCountDensity_eq_rn_lift}
\end{theorem}

\subsection*{relativeLogDensityMean}

\begin{theorem}[blueprint="infogeometry:canonical:grandsynthesis:relativelogdensitymean"]
  \inputleannode{InfoGeometry.Canonical.GrandSynthesis.relativeLogDensityMean}
\end{theorem}

\subsection*{relativeLogDensityMean_mul}

\begin{theorem}[blueprint="infogeometry:canonical:grandsynthesis:relativelogdensitymean_mul"]
  \inputleannode{InfoGeometry.Canonical.GrandSynthesis.relativeLogDensityMean_mul}
\end{theorem}

\subsection*{relativeModularHamiltonian}

\begin{theorem}[blueprint="infogeometry:canonical:grandsynthesis:relativemodularhamiltonian"]
  \inputleannode{InfoGeometry.Canonical.GrandSynthesis.relativeModularHamiltonian}
\end{theorem}

\subsection*{relativeTomitaTakesakiOp}

\begin{theorem}[blueprint="infogeometry:canonical:grandsynthesis:relativetomitatakesakiop"]
  \inputleannode{InfoGeometry.Canonical.GrandSynthesis.relativeTomitaTakesakiOp}
\end{theorem}

\subsection*{relativeTomitaTakesakiOp_apply}

\begin{theorem}[blueprint="infogeometry:canonical:grandsynthesis:relativetomitatakesakiop_apply"]
  \inputleannode{InfoGeometry.Canonical.GrandSynthesis.relativeTomitaTakesakiOp_apply}
\end{theorem}

\subsection*{relativeVolumeChangeRN}

\begin{theorem}[blueprint="infogeometry:canonical:grandsynthesis:relativevolumechangern"]
  \inputleannode{InfoGeometry.Canonical.GrandSynthesis.relativeVolumeChangeRN}
\end{theorem}

\subsection*{relativeVolumeChangeRN_eq_exp_logJacobian}

\begin{theorem}[blueprint="infogeometry:canonical:grandsynthesis:relativevolumechangern_eq_exp_logjacobian"]
  \inputleannode{InfoGeometry.Canonical.GrandSynthesis.relativeVolumeChangeRN_eq_exp_logJacobian}
\end{theorem}

\subsection*{ricci_component_constant_of_geometricEquilibrium}

\begin{theorem}[blueprint="infogeometry:canonical:grandsynthesis:ricci_component_constant_of_geometricequilibrium"]
  \inputleannode{InfoGeometry.Canonical.GrandSynthesis.ricci_component_constant_of_geometricEquilibrium}
\end{theorem}

\subsection*{sinkhornEntropyMonotoneRN}

\begin{theorem}[blueprint="infogeometry:canonical:grandsynthesis:sinkhornentropymonotonern"]
  \inputleannode{InfoGeometry.Canonical.GrandSynthesis.sinkhornEntropyMonotoneRN}
\end{theorem}

\subsection*{sinkhornStepwise_kmsResidual_le_entropyBarrier}

\begin{theorem}[blueprint="infogeometry:canonical:grandsynthesis:sinkhornstepwise_kmsresidual_le_entropybarrier"]
  \inputleannode{InfoGeometry.Canonical.GrandSynthesis.sinkhornStepwise_kmsResidual_le_entropyBarrier}
\end{theorem}

\subsection*{thermodynamicEquilibrium_of_doublyStochastic}

\begin{theorem}[blueprint="infogeometry:canonical:grandsynthesis:thermodynamicequilibrium_of_doublystochastic"]
  \inputleannode{InfoGeometry.Canonical.GrandSynthesis.thermodynamicEquilibrium_of_doublyStochastic}
\end{theorem}

\subsection*{thermodynamic_and_geometricAlgebraic_of_fullCapstone}

\begin{theorem}[blueprint="infogeometry:canonical:grandsynthesis:thermodynamic_and_geometricalgebraic_of_fullcapstone"]
  \inputleannode{InfoGeometry.Canonical.GrandSynthesis.thermodynamic_and_geometricAlgebraic_of_fullCapstone}
\end{theorem}

\subsection*{vacuumEinsteinEquation_of_rnEntropySource}

\begin{theorem}[blueprint="infogeometry:canonical:grandsynthesis:vacuumeinsteinequation_of_rnentropysource"]
  \inputleannode{InfoGeometry.Canonical.GrandSynthesis.vacuumEinsteinEquation_of_rnEntropySource}
\end{theorem}

\section{InfoGeometry.Canonical.GrandSynthesis.jacDetCLM}

\subsection*{eq_1}

\begin{theorem}[blueprint="infogeometry:canonical:grandsynthesis:jacdetclm:eq_1"]
  \inputleannode{InfoGeometry.Canonical.GrandSynthesis.jacDetCLM.eq_1}
\end{theorem}

\section{InfoGeometry.Canonical.GrandSynthesis.relativeTomitaTakesakiOp}

\subsection*{eq_1}

\begin{theorem}[blueprint="infogeometry:canonical:grandsynthesis:relativetomitatakesakiop:eq_1"]
  \inputleannode{InfoGeometry.Canonical.GrandSynthesis.relativeTomitaTakesakiOp.eq_1}
\end{theorem}

\section{InfoGeometry.Canonical.GrandUnification}

\noindent\textbf{Buildable}: yes

\subsection*{IsJordanKKTGeometry}

\begin{theorem}[blueprint="infogeometry:canonical:grandunification:isjordankktgeometry"]
  \inputleannode{InfoGeometry.Canonical.GrandUnification.IsJordanKKTGeometry}
\end{theorem}

\subsection*{JordanKKTData}

\begin{theorem}[blueprint="infogeometry:canonical:grandunification:jordankktdata"]
  \inputleannode{InfoGeometry.Canonical.GrandUnification.JordanKKTData}
\end{theorem}

\subsection*{geometry_divergence_eq_jordan_bregman}

\begin{theorem}[blueprint="infogeometry:canonical:grandunification:geometry_divergence_eq_jordan_bregman"]
  \inputleannode{InfoGeometry.Canonical.GrandUnification.geometry_divergence_eq_jordan_bregman}
\end{theorem}

\section{InfoGeometry.Canonical.GrandUnification.IsJordanKKTGeometry}

\subsection*{D_eq_DBregman}

\begin{theorem}[blueprint="infogeometry:canonical:grandunification:isjordankktgeometry:d_eq_dbregman"]
  \inputleannode{InfoGeometry.Canonical.GrandUnification.IsJordanKKTGeometry.D_eq_DBregman}
\end{theorem}

\subsection*{casesOn}

\begin{theorem}[blueprint="infogeometry:canonical:grandunification:isjordankktgeometry:caseson"]
  \inputleannode{InfoGeometry.Canonical.GrandUnification.IsJordanKKTGeometry.casesOn}
\end{theorem}

\subsection*{mk}

\begin{theorem}[blueprint="infogeometry:canonical:grandunification:isjordankktgeometry:mk"]
  \inputleannode{InfoGeometry.Canonical.GrandUnification.IsJordanKKTGeometry.mk}
\end{theorem}

\subsection*{rec}

\begin{theorem}[blueprint="infogeometry:canonical:grandunification:isjordankktgeometry:rec"]
  \inputleannode{InfoGeometry.Canonical.GrandUnification.IsJordanKKTGeometry.rec}
\end{theorem}

\subsection*{recOn}

\begin{theorem}[blueprint="infogeometry:canonical:grandunification:isjordankktgeometry:recon"]
  \inputleannode{InfoGeometry.Canonical.GrandUnification.IsJordanKKTGeometry.recOn}
\end{theorem}

\section{InfoGeometry.Canonical.GrandUnification.JordanKKTData}

\subsection*{DBregman}

\begin{theorem}[blueprint="infogeometry:canonical:grandunification:jordankktdata:dbregman"]
  \inputleannode{InfoGeometry.Canonical.GrandUnification.JordanKKTData.DBregman}
\end{theorem}

\subsection*{DBregman_def}

\begin{theorem}[blueprint="infogeometry:canonical:grandunification:jordankktdata:dbregman_def"]
  \inputleannode{InfoGeometry.Canonical.GrandUnification.JordanKKTData.DBregman_def}
\end{theorem}

\subsection*{DBregman_eq_zero_iff}

\begin{theorem}[blueprint="infogeometry:canonical:grandunification:jordankktdata:dbregman_eq_zero_iff"]
  \inputleannode{InfoGeometry.Canonical.GrandUnification.JordanKKTData.DBregman_eq_zero_iff}
\end{theorem}

\subsection*{DBregman_nonneg}

\begin{theorem}[blueprint="infogeometry:canonical:grandunification:jordankktdata:dbregman_nonneg"]
  \inputleannode{InfoGeometry.Canonical.GrandUnification.JordanKKTData.DBregman_nonneg}
\end{theorem}

\subsection*{DBregman_pos}

\begin{theorem}[blueprint="infogeometry:canonical:grandunification:jordankktdata:dbregman_pos"]
  \inputleannode{InfoGeometry.Canonical.GrandUnification.JordanKKTData.DBregman_pos}
\end{theorem}

\subsection*{DBregman_self}

\begin{theorem}[blueprint="infogeometry:canonical:grandunification:jordankktdata:dbregman_self"]
  \inputleannode{InfoGeometry.Canonical.GrandUnification.JordanKKTData.DBregman_self}
\end{theorem}

\subsection*{IsBregmanOrthogonal}

\begin{theorem}[blueprint="infogeometry:canonical:grandunification:jordankktdata:isbregmanorthogonal"]
  \inputleannode{InfoGeometry.Canonical.GrandUnification.JordanKKTData.IsBregmanOrthogonal}
\end{theorem}

\subsection*{IsBregmanProjection}

\begin{theorem}[blueprint="infogeometry:canonical:grandunification:jordankktdata:isbregmanprojection"]
  \inputleannode{InfoGeometry.Canonical.GrandUnification.JordanKKTData.IsBregmanProjection}
\end{theorem}

\subsection*{K}

\begin{theorem}[blueprint="infogeometry:canonical:grandunification:jordankktdata:k"]
  \inputleannode{InfoGeometry.Canonical.GrandUnification.JordanKKTData.K}
\end{theorem}

\subsection*{K_def}

\begin{theorem}[blueprint="infogeometry:canonical:grandunification:jordankktdata:k_def"]
  \inputleannode{InfoGeometry.Canonical.GrandUnification.JordanKKTData.K_def}
\end{theorem}

\subsection*{MetricLaws}

\begin{theorem}[blueprint="infogeometry:canonical:grandunification:jordankktdata:metriclaws"]
  \inputleannode{InfoGeometry.Canonical.GrandUnification.JordanKKTData.MetricLaws}
\end{theorem}

\subsection*{casesOn}

\begin{theorem}[blueprint="infogeometry:canonical:grandunification:jordankktdata:caseson"]
  \inputleannode{InfoGeometry.Canonical.GrandUnification.JordanKKTData.casesOn}
\end{theorem}

\subsection*{ctorIdx}

\begin{theorem}[blueprint="infogeometry:canonical:grandunification:jordankktdata:ctoridx"]
  \inputleannode{InfoGeometry.Canonical.GrandUnification.JordanKKTData.ctorIdx}
\end{theorem}

\subsection*{dbregman_local_limit}

\begin{theorem}[blueprint="infogeometry:canonical:grandunification:jordankktdata:dbregman_local_limit"]
  \inputleannode{InfoGeometry.Canonical.GrandUnification.JordanKKTData.dbregman_local_limit}
\end{theorem}

\subsection*{dbregman_pos_of_ne}

\begin{theorem}[blueprint="infogeometry:canonical:grandunification:jordankktdata:dbregman_pos_of_ne"]
  \inputleannode{InfoGeometry.Canonical.GrandUnification.JordanKKTData.dbregman_pos_of_ne}
\end{theorem}

\subsection*{detJ}

\begin{theorem}[blueprint="infogeometry:canonical:grandunification:jordankktdata:detj"]
  \inputleannode{InfoGeometry.Canonical.GrandUnification.JordanKKTData.detJ}
\end{theorem}

\subsection*{detJ_pos}

\begin{theorem}[blueprint="infogeometry:canonical:grandunification:jordankktdata:detj_pos"]
  \inputleannode{InfoGeometry.Canonical.GrandUnification.JordanKKTData.detJ_pos}
\end{theorem}

\subsection*{dualCoord}

\begin{theorem}[blueprint="infogeometry:canonical:grandunification:jordankktdata:dualcoord"]
  \inputleannode{InfoGeometry.Canonical.GrandUnification.JordanKKTData.dualCoord}
\end{theorem}

\subsection*{dualPotential}

\begin{theorem}[blueprint="infogeometry:canonical:grandunification:jordankktdata:dualpotential"]
  \inputleannode{InfoGeometry.Canonical.GrandUnification.JordanKKTData.dualPotential}
\end{theorem}

\subsection*{dualPotential_def}

\begin{theorem}[blueprint="infogeometry:canonical:grandunification:jordankktdata:dualpotential_def"]
  \inputleannode{InfoGeometry.Canonical.GrandUnification.JordanKKTData.dualPotential_def}
\end{theorem}

\subsection*{dualPotential_grad_eq}

\begin{theorem}[blueprint="infogeometry:canonical:grandunification:jordankktdata:dualpotential_grad_eq"]
  \inputleannode{InfoGeometry.Canonical.GrandUnification.JordanKKTData.dualPotential_grad_eq}
\end{theorem}

\subsection*{dual_potential_mixed_identity}

\begin{theorem}[blueprint="infogeometry:canonical:grandunification:jordankktdata:dual_potential_mixed_identity"]
  \inputleannode{InfoGeometry.Canonical.GrandUnification.JordanKKTData.dual_potential_mixed_identity}
\end{theorem}

\subsection*{fenchel_young_equality}

\begin{theorem}[blueprint="infogeometry:canonical:grandunification:jordankktdata:fenchel_young_equality"]
  \inputleannode{InfoGeometry.Canonical.GrandUnification.JordanKKTData.fenchel_young_equality}
\end{theorem}

\subsection*{g}

\begin{theorem}[blueprint="infogeometry:canonical:grandunification:jordankktdata:g"]
  \inputleannode{InfoGeometry.Canonical.GrandUnification.JordanKKTData.g}
\end{theorem}

\subsection*{generalized_pythagorean_theorem}

\begin{theorem}[blueprint="infogeometry:canonical:grandunification:jordankktdata:generalized_pythagorean_theorem"]
  \inputleannode{InfoGeometry.Canonical.GrandUnification.JordanKKTData.generalized_pythagorean_theorem}
\end{theorem}

\subsection*{gradK}

\begin{theorem}[blueprint="infogeometry:canonical:grandunification:jordankktdata:gradk"]
  \inputleannode{InfoGeometry.Canonical.GrandUnification.JordanKKTData.gradK}
\end{theorem}

\subsection*{hasFDerivAt_K}

\begin{theorem}[blueprint="infogeometry:canonical:grandunification:jordankktdata:hasfderivat_k"]
  \inputleannode{InfoGeometry.Canonical.GrandUnification.JordanKKTData.hasFDerivAt_K}
\end{theorem}

\subsection*{has_gradient}

\begin{theorem}[blueprint="infogeometry:canonical:grandunification:jordankktdata:has_gradient"]
  \inputleannode{InfoGeometry.Canonical.GrandUnification.JordanKKTData.has_gradient}
\end{theorem}

\subsection*{has_hessian}

\begin{theorem}[blueprint="infogeometry:canonical:grandunification:jordankktdata:has_hessian"]
  \inputleannode{InfoGeometry.Canonical.GrandUnification.JordanKKTData.has_hessian}
\end{theorem}

\subsection*{metric}

\begin{theorem}[blueprint="infogeometry:canonical:grandunification:jordankktdata:metric"]
  \inputleannode{InfoGeometry.Canonical.GrandUnification.JordanKKTData.metric}
\end{theorem}

\subsection*{metric_symmetry}

\begin{theorem}[blueprint="infogeometry:canonical:grandunification:jordankktdata:metric_symmetry"]
  \inputleannode{InfoGeometry.Canonical.GrandUnification.JordanKKTData.metric_symmetry}
\end{theorem}

\subsection*{mixed_bregman_identity}

\begin{theorem}[blueprint="infogeometry:canonical:grandunification:jordankktdata:mixed_bregman_identity"]
  \inputleannode{InfoGeometry.Canonical.GrandUnification.JordanKKTData.mixed_bregman_identity}
\end{theorem}

\subsection*{mk}

\begin{theorem}[blueprint="infogeometry:canonical:grandunification:jordankktdata:mk"]
  \inputleannode{InfoGeometry.Canonical.GrandUnification.JordanKKTData.mk}
\end{theorem}

\subsection*{noConfusion}

\begin{theorem}[blueprint="infogeometry:canonical:grandunification:jordankktdata:noconfusion"]
  \inputleannode{InfoGeometry.Canonical.GrandUnification.JordanKKTData.noConfusion}
\end{theorem}

\subsection*{noConfusionType}

\begin{theorem}[blueprint="infogeometry:canonical:grandunification:jordankktdata:noconfusiontype"]
  \inputleannode{InfoGeometry.Canonical.GrandUnification.JordanKKTData.noConfusionType}
\end{theorem}

\subsection*{projection_minimizes}

\begin{theorem}[blueprint="infogeometry:canonical:grandunification:jordankktdata:projection_minimizes"]
  \inputleannode{InfoGeometry.Canonical.GrandUnification.JordanKKTData.projection_minimizes}
\end{theorem}

\subsection*{projection_unique}

\begin{theorem}[blueprint="infogeometry:canonical:grandunification:jordankktdata:projection_unique"]
  \inputleannode{InfoGeometry.Canonical.GrandUnification.JordanKKTData.projection_unique}
\end{theorem}

\subsection*{pythagorean_inequality}

\begin{theorem}[blueprint="infogeometry:canonical:grandunification:jordankktdata:pythagorean_inequality"]
  \inputleannode{InfoGeometry.Canonical.GrandUnification.JordanKKTData.pythagorean_inequality}
\end{theorem}

\subsection*{rec}

\begin{theorem}[blueprint="infogeometry:canonical:grandunification:jordankktdata:rec"]
  \inputleannode{InfoGeometry.Canonical.GrandUnification.JordanKKTData.rec}
\end{theorem}

\subsection*{recOn}

\begin{theorem}[blueprint="infogeometry:canonical:grandunification:jordankktdata:recon"]
  \inputleannode{InfoGeometry.Canonical.GrandUnification.JordanKKTData.recOn}
\end{theorem}

\subsection*{three_point_law}

\begin{theorem}[blueprint="infogeometry:canonical:grandunification:jordankktdata:three_point_law"]
  \inputleannode{InfoGeometry.Canonical.GrandUnification.JordanKKTData.three_point_law}
\end{theorem}

\section{InfoGeometry.Canonical.GrandUnification.JordanKKTData.DBregman}

\subsection*{eq_1}

\begin{theorem}[blueprint="infogeometry:canonical:grandunification:jordankktdata:dbregman:eq_1"]
  \inputleannode{InfoGeometry.Canonical.GrandUnification.JordanKKTData.DBregman.eq_1}
\end{theorem}

\section{InfoGeometry.Canonical.GrandUnification.JordanKKTData.K}

\subsection*{eq_1}

\begin{theorem}[blueprint="infogeometry:canonical:grandunification:jordankktdata:k:eq_1"]
  \inputleannode{InfoGeometry.Canonical.GrandUnification.JordanKKTData.K.eq_1}
\end{theorem}

\section{InfoGeometry.Canonical.GrandUnification.JordanKKTData.MetricLaws}

\subsection*{bregman_local}

\begin{theorem}[blueprint="infogeometry:canonical:grandunification:jordankktdata:metriclaws:bregman_local"]
  \inputleannode{InfoGeometry.Canonical.GrandUnification.JordanKKTData.MetricLaws.bregman_local}
\end{theorem}

\subsection*{casesOn}

\begin{theorem}[blueprint="infogeometry:canonical:grandunification:jordankktdata:metriclaws:caseson"]
  \inputleannode{InfoGeometry.Canonical.GrandUnification.JordanKKTData.MetricLaws.casesOn}
\end{theorem}

\subsection*{metric_symm}

\begin{theorem}[blueprint="infogeometry:canonical:grandunification:jordankktdata:metriclaws:metric_symm"]
  \inputleannode{InfoGeometry.Canonical.GrandUnification.JordanKKTData.MetricLaws.metric_symm}
\end{theorem}

\subsection*{mk}

\begin{theorem}[blueprint="infogeometry:canonical:grandunification:jordankktdata:metriclaws:mk"]
  \inputleannode{InfoGeometry.Canonical.GrandUnification.JordanKKTData.MetricLaws.mk}
\end{theorem}

\subsection*{rec}

\begin{theorem}[blueprint="infogeometry:canonical:grandunification:jordankktdata:metriclaws:rec"]
  \inputleannode{InfoGeometry.Canonical.GrandUnification.JordanKKTData.MetricLaws.rec}
\end{theorem}

\subsection*{recOn}

\begin{theorem}[blueprint="infogeometry:canonical:grandunification:jordankktdata:metriclaws:recon"]
  \inputleannode{InfoGeometry.Canonical.GrandUnification.JordanKKTData.MetricLaws.recOn}
\end{theorem}

\section{InfoGeometry.Canonical.GrandUnification.JordanKKTData.dualPotential}

\subsection*{eq_1}

\begin{theorem}[blueprint="infogeometry:canonical:grandunification:jordankktdata:dualpotential:eq_1"]
  \inputleannode{InfoGeometry.Canonical.GrandUnification.JordanKKTData.dualPotential.eq_1}
\end{theorem}

\section{InfoGeometry.Canonical.GrandUnificationBlueprint}

\noindent\textbf{Buildable}: yes

\subsection*{UnificationComplete}

\begin{theorem}[blueprint="infogeometry:canonical:grandunificationblueprint:unificationcomplete"]
  \inputleannode{InfoGeometry.Canonical.GrandUnificationBlueprint.UnificationComplete}
\end{theorem}

\section{InfoGeometry.Canonical.GrandUnificationMetric}

\noindent\textbf{Buildable}: yes

\noindent\emph{No exported declarations in the current declaration graph.}

\section{InfoGeometry.Canonical.HeatKernel}

\noindent\textbf{Buildable}: yes

\subsection*{a0}

\begin{theorem}[blueprint="infogeometry:canonical:heatkernel:a0"]
  \inputleannode{InfoGeometry.Canonical.HeatKernel.a0}
\end{theorem}

\subsection*{a0_eq_spectralVolume}

\begin{theorem}[blueprint="infogeometry:canonical:heatkernel:a0_eq_spectralvolume"]
  \inputleannode{InfoGeometry.Canonical.HeatKernel.a0_eq_spectralVolume}
\end{theorem}

\subsection*{a1}

\begin{theorem}[blueprint="infogeometry:canonical:heatkernel:a1"]
  \inputleannode{InfoGeometry.Canonical.HeatKernel.a1}
\end{theorem}

\subsection*{a1_eq_neg_spectralVolume}

\begin{theorem}[blueprint="infogeometry:canonical:heatkernel:a1_eq_neg_spectralvolume"]
  \inputleannode{InfoGeometry.Canonical.HeatKernel.a1_eq_neg_spectralVolume}
\end{theorem}

\subsection*{einsteinHilbertAction}

\begin{theorem}[blueprint="infogeometry:canonical:heatkernel:einsteinhilbertaction"]
  \inputleannode{InfoGeometry.Canonical.HeatKernel.einsteinHilbertAction}
\end{theorem}

\subsection*{einsteinHilbertAction_eq_neg_six_spectralVolume}

\begin{theorem}[blueprint="infogeometry:canonical:heatkernel:einsteinhilbertaction_eq_neg_six_spectralvolume"]
  \inputleannode{InfoGeometry.Canonical.HeatKernel.einsteinHilbertAction_eq_neg_six_spectralVolume}
\end{theorem}

\subsection*{einsteinHilbertAction_eq_totalScalarCurvature}

\begin{theorem}[blueprint="infogeometry:canonical:heatkernel:einsteinhilbertaction_eq_totalscalarcurvature"]
  \inputleannode{InfoGeometry.Canonical.HeatKernel.einsteinHilbertAction_eq_totalScalarCurvature}
\end{theorem}

\subsection*{heatTrace}

\begin{theorem}[blueprint="infogeometry:canonical:heatkernel:heattrace"]
  \inputleannode{InfoGeometry.Canonical.HeatKernel.heatTrace}
\end{theorem}

\subsection*{spectralVolume}

\begin{theorem}[blueprint="infogeometry:canonical:heatkernel:spectralvolume"]
  \inputleannode{InfoGeometry.Canonical.HeatKernel.spectralVolume}
\end{theorem}

\subsection*{totalScalarCurvature}

\begin{theorem}[blueprint="infogeometry:canonical:heatkernel:totalscalarcurvature"]
  \inputleannode{InfoGeometry.Canonical.HeatKernel.totalScalarCurvature}
\end{theorem}

\subsection*{totalScalarCurvature_eq_neg_six_spectralVolume}

\begin{theorem}[blueprint="infogeometry:canonical:heatkernel:totalscalarcurvature_eq_neg_six_spectralvolume"]
  \inputleannode{InfoGeometry.Canonical.HeatKernel.totalScalarCurvature_eq_neg_six_spectralVolume}
\end{theorem}

\section{InfoGeometry.Canonical.HeatKernel.a0}

\subsection*{eq_1}

\begin{theorem}[blueprint="infogeometry:canonical:heatkernel:a0:eq_1"]
  \inputleannode{InfoGeometry.Canonical.HeatKernel.a0.eq_1}
\end{theorem}

\section{InfoGeometry.Canonical.HeatKernel.einsteinHilbertAction}

\subsection*{eq_1}

\begin{theorem}[blueprint="infogeometry:canonical:heatkernel:einsteinhilbertaction:eq_1"]
  \inputleannode{InfoGeometry.Canonical.HeatKernel.einsteinHilbertAction.eq_1}
\end{theorem}

\section{InfoGeometry.Canonical.HeatKernel.totalScalarCurvature}

\subsection*{eq_1}

\begin{theorem}[blueprint="infogeometry:canonical:heatkernel:totalscalarcurvature:eq_1"]
  \inputleannode{InfoGeometry.Canonical.HeatKernel.totalScalarCurvature.eq_1}
\end{theorem}

\section{InfoGeometry.Canonical.HolographicEmergence}

\noindent\textbf{Buildable}: yes

\subsection*{AnomalyScalePhase}

\begin{theorem}[blueprint="infogeometry:canonical:holographicemergence:anomalyscalephase"]
  \inputleannode{InfoGeometry.Canonical.HolographicEmergence.AnomalyScalePhase}
\end{theorem}

\subsection*{EmergentTimeFlow}

\begin{theorem}[blueprint="infogeometry:canonical:holographicemergence:emergenttimeflow"]
  \inputleannode{InfoGeometry.Canonical.HolographicEmergence.EmergentTimeFlow}
\end{theorem}

\subsection*{anomalyScalePhase_of_nonzeroAnomaly}

\begin{theorem}[blueprint="infogeometry:canonical:holographicemergence:anomalyscalephase_of_nonzeroanomaly"]
  \inputleannode{InfoGeometry.Canonical.HolographicEmergence.anomalyScalePhase_of_nonzeroAnomaly}
\end{theorem}

\subsection*{boundaryAnomalyCancellation}

\begin{theorem}[blueprint="infogeometry:canonical:holographicemergence:boundaryanomalycancellation"]
  \inputleannode{InfoGeometry.Canonical.HolographicEmergence.boundaryAnomalyCancellation}
\end{theorem}

\subsection*{emergentTimeFlow_of_sinkhornTrajectory}

\begin{theorem}[blueprint="infogeometry:canonical:holographicemergence:emergenttimeflow_of_sinkhorntrajectory"]
  \inputleannode{InfoGeometry.Canonical.HolographicEmergence.emergentTimeFlow_of_sinkhornTrajectory}
\end{theorem}

\subsection*{exists_gaugeOrderHysteresis_witness}

\begin{theorem}[blueprint="infogeometry:canonical:holographicemergence:exists_gaugeorderhysteresis_witness"]
  \inputleannode{InfoGeometry.Canonical.HolographicEmergence.exists_gaugeOrderHysteresis_witness}
\end{theorem}

\subsection*{holographicEmergence_package}

\begin{theorem}[blueprint="infogeometry:canonical:holographicemergence:holographicemergence_package"]
  \inputleannode{InfoGeometry.Canonical.HolographicEmergence.holographicEmergence_package}
\end{theorem}

\subsection*{pathDependence_of_twistedInference}

\begin{theorem}[blueprint="infogeometry:canonical:holographicemergence:pathdependence_of_twistedinference"]
  \inputleannode{InfoGeometry.Canonical.HolographicEmergence.pathDependence_of_twistedInference}
\end{theorem}

\subsection*{vacuumApexNull_lifts_to_twistor}

\begin{theorem}[blueprint="infogeometry:canonical:holographicemergence:vacuumapexnull_lifts_to_twistor"]
  \inputleannode{InfoGeometry.Canonical.HolographicEmergence.vacuumApexNull_lifts_to_twistor}
\end{theorem}

\section{InfoGeometry.Canonical.IB}

\noindent\textbf{Buildable}: yes

\noindent\emph{No exported declarations in the current declaration graph.}

\section{InfoGeometry.Canonical.InformationNumber}

\noindent\textbf{Buildable}: yes

\subsection*{informationNumber}

\begin{theorem}[blueprint="infogeometry:canonical:informationnumber:informationnumber"]
  \inputleannode{InfoGeometry.Canonical.InformationNumber.informationNumber}
\end{theorem}

\subsection*{informationNumber_eq_dim}

\begin{theorem}[blueprint="infogeometry:canonical:informationnumber:informationnumber_eq_dim"]
  \inputleannode{InfoGeometry.Canonical.InformationNumber.informationNumber_eq_dim}
\end{theorem}

\section{InfoGeometry.Canonical.InformationNumber.informationNumber}

\subsection*{eq_1}

\begin{theorem}[blueprint="infogeometry:canonical:informationnumber:informationnumber:eq_1"]
  \inputleannode{InfoGeometry.Canonical.InformationNumber.informationNumber.eq_1}
\end{theorem}

\section{InfoGeometry.Canonical.InformationTorsion}

\noindent\textbf{Buildable}: yes

\subsection*{Connection}

\begin{theorem}[blueprint="infogeometry:canonical:informationtorsion:connection"]
  \inputleannode{InfoGeometry.Canonical.InformationTorsion.Connection}
\end{theorem}

\subsection*{DualConnections}

\begin{theorem}[blueprint="infogeometry:canonical:informationtorsion:dualconnections"]
  \inputleannode{InfoGeometry.Canonical.InformationTorsion.DualConnections}
\end{theorem}

\subsection*{FlatDualConnections}

\begin{theorem}[blueprint="infogeometry:canonical:informationtorsion:flatdualconnections"]
  \inputleannode{InfoGeometry.Canonical.InformationTorsion.FlatDualConnections}
\end{theorem}

\subsection*{IsTorsionFree}

\begin{theorem}[blueprint="infogeometry:canonical:informationtorsion:istorsionfree"]
  \inputleannode{InfoGeometry.Canonical.InformationTorsion.IsTorsionFree}
\end{theorem}

\subsection*{TwistedInference}

\begin{theorem}[blueprint="infogeometry:canonical:informationtorsion:twistedinference"]
  \inputleannode{InfoGeometry.Canonical.InformationTorsion.TwistedInference}
\end{theorem}

\subsection*{hessian_is_torsion_free}

\begin{theorem}[blueprint="infogeometry:canonical:informationtorsion:hessian_is_torsion_free"]
  \inputleannode{InfoGeometry.Canonical.InformationTorsion.hessian_is_torsion_free}
\end{theorem}

\subsection*{informationTorsion}

\begin{theorem}[blueprint="infogeometry:canonical:informationtorsion:informationtorsion"]
  \inputleannode{InfoGeometry.Canonical.InformationTorsion.informationTorsion}
\end{theorem}

\section{InfoGeometry.Canonical.InformationTorsion.DualConnections}

\subsection*{casesOn}

\begin{theorem}[blueprint="infogeometry:canonical:informationtorsion:dualconnections:caseson"]
  \inputleannode{InfoGeometry.Canonical.InformationTorsion.DualConnections.casesOn}
\end{theorem}

\subsection*{ctorIdx}

\begin{theorem}[blueprint="infogeometry:canonical:informationtorsion:dualconnections:ctoridx"]
  \inputleannode{InfoGeometry.Canonical.InformationTorsion.DualConnections.ctorIdx}
\end{theorem}

\subsection*{mk}

\begin{theorem}[blueprint="infogeometry:canonical:informationtorsion:dualconnections:mk"]
  \inputleannode{InfoGeometry.Canonical.InformationTorsion.DualConnections.mk}
\end{theorem}

\subsection*{nabla}

\begin{theorem}[blueprint="infogeometry:canonical:informationtorsion:dualconnections:nabla"]
  \inputleannode{InfoGeometry.Canonical.InformationTorsion.DualConnections.nabla}
\end{theorem}

\subsection*{nablaStar}

\begin{theorem}[blueprint="infogeometry:canonical:informationtorsion:dualconnections:nablastar"]
  \inputleannode{InfoGeometry.Canonical.InformationTorsion.DualConnections.nablaStar}
\end{theorem}

\subsection*{noConfusion}

\begin{theorem}[blueprint="infogeometry:canonical:informationtorsion:dualconnections:noconfusion"]
  \inputleannode{InfoGeometry.Canonical.InformationTorsion.DualConnections.noConfusion}
\end{theorem}

\subsection*{noConfusionType}

\begin{theorem}[blueprint="infogeometry:canonical:informationtorsion:dualconnections:noconfusiontype"]
  \inputleannode{InfoGeometry.Canonical.InformationTorsion.DualConnections.noConfusionType}
\end{theorem}

\subsection*{rec}

\begin{theorem}[blueprint="infogeometry:canonical:informationtorsion:dualconnections:rec"]
  \inputleannode{InfoGeometry.Canonical.InformationTorsion.DualConnections.rec}
\end{theorem}

\subsection*{recOn}

\begin{theorem}[blueprint="infogeometry:canonical:informationtorsion:dualconnections:recon"]
  \inputleannode{InfoGeometry.Canonical.InformationTorsion.DualConnections.recOn}
\end{theorem}

\section{InfoGeometry.Canonical.InformationTorsion.FlatDualConnections}

\subsection*{casesOn}

\begin{theorem}[blueprint="infogeometry:canonical:informationtorsion:flatdualconnections:caseson"]
  \inputleannode{InfoGeometry.Canonical.InformationTorsion.FlatDualConnections.casesOn}
\end{theorem}

\subsection*{ctorIdx}

\begin{theorem}[blueprint="infogeometry:canonical:informationtorsion:flatdualconnections:ctoridx"]
  \inputleannode{InfoGeometry.Canonical.InformationTorsion.FlatDualConnections.ctorIdx}
\end{theorem}

\subsection*{dual}

\begin{theorem}[blueprint="infogeometry:canonical:informationtorsion:flatdualconnections:dual"]
  \inputleannode{InfoGeometry.Canonical.InformationTorsion.FlatDualConnections.dual}
\end{theorem}

\subsection*{mk}

\begin{theorem}[blueprint="infogeometry:canonical:informationtorsion:flatdualconnections:mk"]
  \inputleannode{InfoGeometry.Canonical.InformationTorsion.FlatDualConnections.mk}
\end{theorem}

\subsection*{noConfusion}

\begin{theorem}[blueprint="infogeometry:canonical:informationtorsion:flatdualconnections:noconfusion"]
  \inputleannode{InfoGeometry.Canonical.InformationTorsion.FlatDualConnections.noConfusion}
\end{theorem}

\subsection*{noConfusionType}

\begin{theorem}[blueprint="infogeometry:canonical:informationtorsion:flatdualconnections:noconfusiontype"]
  \inputleannode{InfoGeometry.Canonical.InformationTorsion.FlatDualConnections.noConfusionType}
\end{theorem}

\subsection*{rec}

\begin{theorem}[blueprint="infogeometry:canonical:informationtorsion:flatdualconnections:rec"]
  \inputleannode{InfoGeometry.Canonical.InformationTorsion.FlatDualConnections.rec}
\end{theorem}

\subsection*{recOn}

\begin{theorem}[blueprint="infogeometry:canonical:informationtorsion:flatdualconnections:recon"]
  \inputleannode{InfoGeometry.Canonical.InformationTorsion.FlatDualConnections.recOn}
\end{theorem}

\subsection*{torsion_free_nabla}

\begin{theorem}[blueprint="infogeometry:canonical:informationtorsion:flatdualconnections:torsion_free_nabla"]
  \inputleannode{InfoGeometry.Canonical.InformationTorsion.FlatDualConnections.torsion_free_nabla}
\end{theorem}

\subsection*{torsion_free_nablaStar}

\begin{theorem}[blueprint="infogeometry:canonical:informationtorsion:flatdualconnections:torsion_free_nablastar"]
  \inputleannode{InfoGeometry.Canonical.InformationTorsion.FlatDualConnections.torsion_free_nablaStar}
\end{theorem}

\section{InfoGeometry.Canonical.InformationTorsion.TwistedInference}

\subsection*{casesOn}

\begin{theorem}[blueprint="infogeometry:canonical:informationtorsion:twistedinference:caseson"]
  \inputleannode{InfoGeometry.Canonical.InformationTorsion.TwistedInference.casesOn}
\end{theorem}

\subsection*{ctorIdx}

\begin{theorem}[blueprint="infogeometry:canonical:informationtorsion:twistedinference:ctoridx"]
  \inputleannode{InfoGeometry.Canonical.InformationTorsion.TwistedInference.ctorIdx}
\end{theorem}

\subsection*{dual}

\begin{theorem}[blueprint="infogeometry:canonical:informationtorsion:twistedinference:dual"]
  \inputleannode{InfoGeometry.Canonical.InformationTorsion.TwistedInference.dual}
\end{theorem}

\subsection*{has_torsion}

\begin{theorem}[blueprint="infogeometry:canonical:informationtorsion:twistedinference:has_torsion"]
  \inputleannode{InfoGeometry.Canonical.InformationTorsion.TwistedInference.has_torsion}
\end{theorem}

\subsection*{mk}

\begin{theorem}[blueprint="infogeometry:canonical:informationtorsion:twistedinference:mk"]
  \inputleannode{InfoGeometry.Canonical.InformationTorsion.TwistedInference.mk}
\end{theorem}

\subsection*{noConfusion}

\begin{theorem}[blueprint="infogeometry:canonical:informationtorsion:twistedinference:noconfusion"]
  \inputleannode{InfoGeometry.Canonical.InformationTorsion.TwistedInference.noConfusion}
\end{theorem}

\subsection*{noConfusionType}

\begin{theorem}[blueprint="infogeometry:canonical:informationtorsion:twistedinference:noconfusiontype"]
  \inputleannode{InfoGeometry.Canonical.InformationTorsion.TwistedInference.noConfusionType}
\end{theorem}

\subsection*{rec}

\begin{theorem}[blueprint="infogeometry:canonical:informationtorsion:twistedinference:rec"]
  \inputleannode{InfoGeometry.Canonical.InformationTorsion.TwistedInference.rec}
\end{theorem}

\subsection*{recOn}

\begin{theorem}[blueprint="infogeometry:canonical:informationtorsion:twistedinference:recon"]
  \inputleannode{InfoGeometry.Canonical.InformationTorsion.TwistedInference.recOn}
\end{theorem}

\section{InfoGeometry.Canonical.IntegrationLegacy}

\noindent\textbf{Buildable}: yes

\noindent\emph{No exported declarations in the current declaration graph.}

\section{InfoGeometry.Canonical.KMSSinkhornBridge}

\noindent\textbf{Buildable}: yes

\subsection*{RouterAmplitude}

\begin{theorem}[blueprint="infogeometry:canonical:kmssinkhornbridge:routeramplitude"]
  \inputleannode{InfoGeometry.Canonical.KMSSinkhornBridge.RouterAmplitude}
\end{theorem}

\subsection*{SatisfiesApproxKMSLike}

\begin{theorem}[blueprint="infogeometry:canonical:kmssinkhornbridge:satisfiesapproxkmslike"]
  \inputleannode{InfoGeometry.Canonical.KMSSinkhornBridge.SatisfiesApproxKMSLike}
\end{theorem}

\subsection*{SinkhornKMSClosure}

\begin{theorem}[blueprint="infogeometry:canonical:kmssinkhornbridge:sinkhornkmsclosure"]
  \inputleannode{InfoGeometry.Canonical.KMSSinkhornBridge.SinkhornKMSClosure}
\end{theorem}

\subsection*{SinkhornKMSControl}

\begin{theorem}[blueprint="infogeometry:canonical:kmssinkhornbridge:sinkhornkmscontrol"]
  \inputleannode{InfoGeometry.Canonical.KMSSinkhornBridge.SinkhornKMSControl}
\end{theorem}

\subsection*{kmsResidual}

\begin{theorem}[blueprint="infogeometry:canonical:kmssinkhornbridge:kmsresidual"]
  \inputleannode{InfoGeometry.Canonical.KMSSinkhornBridge.kmsResidual}
\end{theorem}

\subsection*{routerMeanEnergy}

\begin{theorem}[blueprint="infogeometry:canonical:kmssinkhornbridge:routermeanenergy"]
  \inputleannode{InfoGeometry.Canonical.KMSSinkhornBridge.routerMeanEnergy}
\end{theorem}

\subsection*{routerModularHamiltonian}

\begin{theorem}[blueprint="infogeometry:canonical:kmssinkhornbridge:routermodularhamiltonian"]
  \inputleannode{InfoGeometry.Canonical.KMSSinkhornBridge.routerModularHamiltonian}
\end{theorem}

\subsection*{routerModularHamiltonian_apply}

\begin{theorem}[blueprint="infogeometry:canonical:kmssinkhornbridge:routermodularhamiltonian_apply"]
  \inputleannode{InfoGeometry.Canonical.KMSSinkhornBridge.routerModularHamiltonian_apply}
\end{theorem}

\subsection*{router_sinkhornIterate_kmsClosure_of_control}

\begin{theorem}[blueprint="infogeometry:canonical:kmssinkhornbridge:router_sinkhorniterate_kmsclosure_of_control"]
  \inputleannode{InfoGeometry.Canonical.KMSSinkhornBridge.router_sinkhornIterate_kmsClosure_of_control}
\end{theorem}

\subsection*{satisfiesApproxKMSLike_of_satisfiesKMSLike}

\begin{theorem}[blueprint="infogeometry:canonical:kmssinkhornbridge:satisfiesapproxkmslike_of_satisfieskmslike"]
  \inputleannode{InfoGeometry.Canonical.KMSSinkhornBridge.satisfiesApproxKMSLike_of_satisfiesKMSLike}
\end{theorem}

\subsection*{satisfiesKMSLike_of_approx_zero}

\begin{theorem}[blueprint="infogeometry:canonical:kmssinkhornbridge:satisfieskmslike_of_approx_zero"]
  \inputleannode{InfoGeometry.Canonical.KMSSinkhornBridge.satisfiesKMSLike_of_approx_zero}
\end{theorem}

\subsection*{sinkhornIterate_kmsClosure_of_control}

\begin{theorem}[blueprint="infogeometry:canonical:kmssinkhornbridge:sinkhorniterate_kmsclosure_of_control"]
  \inputleannode{InfoGeometry.Canonical.KMSSinkhornBridge.sinkhornIterate_kmsClosure_of_control}
\end{theorem}

\subsection*{sinkhorn_control_of_step_kmsClosure}

\begin{theorem}[blueprint="infogeometry:canonical:kmssinkhornbridge:sinkhorn_control_of_step_kmsclosure"]
  \inputleannode{InfoGeometry.Canonical.KMSSinkhornBridge.sinkhorn_control_of_step_kmsClosure}
\end{theorem}

\subsection*{sinkhorn_step_kmsClosure_of_control}

\begin{theorem}[blueprint="infogeometry:canonical:kmssinkhornbridge:sinkhorn_step_kmsclosure_of_control"]
  \inputleannode{InfoGeometry.Canonical.KMSSinkhornBridge.sinkhorn_step_kmsClosure_of_control}
\end{theorem}

\subsection*{sinkhorn_stepwise_kms_bound}

\begin{theorem}[blueprint="infogeometry:canonical:kmssinkhornbridge:sinkhorn_stepwise_kms_bound"]
  \inputleannode{InfoGeometry.Canonical.KMSSinkhornBridge.sinkhorn_stepwise_kms_bound}
\end{theorem}

\section{InfoGeometry.Canonical.KMSSinkhornBridge.kmsResidual}

\subsection*{eq_1}

\begin{theorem}[blueprint="infogeometry:canonical:kmssinkhornbridge:kmsresidual:eq_1"]
  \inputleannode{InfoGeometry.Canonical.KMSSinkhornBridge.kmsResidual.eq_1}
\end{theorem}

\section{InfoGeometry.Canonical.KMSSinkhornBridge.routerModularHamiltonian}

\subsection*{eq_1}

\begin{theorem}[blueprint="infogeometry:canonical:kmssinkhornbridge:routermodularhamiltonian:eq_1"]
  \inputleannode{InfoGeometry.Canonical.KMSSinkhornBridge.routerModularHamiltonian.eq_1}
\end{theorem}

\section{InfoGeometry.Canonical.KaehlerGeometry}

\noindent\textbf{Buildable}: yes

\subsection*{KaehlerInformationGeometry}

\begin{theorem}[blueprint="infogeometry:canonical:kaehlergeometry:kaehlerinformationgeometry"]
  \inputleannode{InfoGeometry.Canonical.KaehlerGeometry.KaehlerInformationGeometry}
\end{theorem}

\subsection*{SymplecticForm}

\begin{theorem}[blueprint="infogeometry:canonical:kaehlergeometry:symplecticform"]
  \inputleannode{InfoGeometry.Canonical.KaehlerGeometry.SymplecticForm}
\end{theorem}

\section{InfoGeometry.Canonical.KaehlerGeometry.KaehlerInformationGeometry}

\subsection*{H}

\begin{theorem}[blueprint="infogeometry:canonical:kaehlergeometry:kaehlerinformationgeometry:h"]
  \inputleannode{InfoGeometry.Canonical.KaehlerGeometry.KaehlerInformationGeometry.H}
\end{theorem}

\subsection*{J}

\begin{theorem}[blueprint="infogeometry:canonical:kaehlergeometry:kaehlerinformationgeometry:j"]
  \inputleannode{InfoGeometry.Canonical.KaehlerGeometry.KaehlerInformationGeometry.J}
\end{theorem}

\subsection*{casesOn}

\begin{theorem}[blueprint="infogeometry:canonical:kaehlergeometry:kaehlerinformationgeometry:caseson"]
  \inputleannode{InfoGeometry.Canonical.KaehlerGeometry.KaehlerInformationGeometry.casesOn}
\end{theorem}

\subsection*{compatibility}

\begin{theorem}[blueprint="infogeometry:canonical:kaehlergeometry:kaehlerinformationgeometry:compatibility"]
  \inputleannode{InfoGeometry.Canonical.KaehlerGeometry.KaehlerInformationGeometry.compatibility}
\end{theorem}

\subsection*{ctorIdx}

\begin{theorem}[blueprint="infogeometry:canonical:kaehlergeometry:kaehlerinformationgeometry:ctoridx"]
  \inputleannode{InfoGeometry.Canonical.KaehlerGeometry.KaehlerInformationGeometry.ctorIdx}
\end{theorem}

\subsection*{j_sq_eq_neg_id}

\begin{theorem}[blueprint="infogeometry:canonical:kaehlergeometry:kaehlerinformationgeometry:j_sq_eq_neg_id"]
  \inputleannode{InfoGeometry.Canonical.KaehlerGeometry.KaehlerInformationGeometry.j_sq_eq_neg_id}
\end{theorem}

\subsection*{logF}

\begin{theorem}[blueprint="infogeometry:canonical:kaehlergeometry:kaehlerinformationgeometry:logf"]
  \inputleannode{InfoGeometry.Canonical.KaehlerGeometry.KaehlerInformationGeometry.logF}
\end{theorem}

\subsection*{mk}

\begin{theorem}[blueprint="infogeometry:canonical:kaehlergeometry:kaehlerinformationgeometry:mk"]
  \inputleannode{InfoGeometry.Canonical.KaehlerGeometry.KaehlerInformationGeometry.mk}
\end{theorem}

\subsection*{noConfusion}

\begin{theorem}[blueprint="infogeometry:canonical:kaehlergeometry:kaehlerinformationgeometry:noconfusion"]
  \inputleannode{InfoGeometry.Canonical.KaehlerGeometry.KaehlerInformationGeometry.noConfusion}
\end{theorem}

\subsection*{noConfusionType}

\begin{theorem}[blueprint="infogeometry:canonical:kaehlergeometry:kaehlerinformationgeometry:noconfusiontype"]
  \inputleannode{InfoGeometry.Canonical.KaehlerGeometry.KaehlerInformationGeometry.noConfusionType}
\end{theorem}

\subsection*{rec}

\begin{theorem}[blueprint="infogeometry:canonical:kaehlergeometry:kaehlerinformationgeometry:rec"]
  \inputleannode{InfoGeometry.Canonical.KaehlerGeometry.KaehlerInformationGeometry.rec}
\end{theorem}

\subsection*{recOn}

\begin{theorem}[blueprint="infogeometry:canonical:kaehlergeometry:kaehlerinformationgeometry:recon"]
  \inputleannode{InfoGeometry.Canonical.KaehlerGeometry.KaehlerInformationGeometry.recOn}
\end{theorem}

\subsection*{symplecticCurvature}

\begin{theorem}[blueprint="infogeometry:canonical:kaehlergeometry:kaehlerinformationgeometry:symplecticcurvature"]
  \inputleannode{InfoGeometry.Canonical.KaehlerGeometry.KaehlerInformationGeometry.symplecticCurvature}
\end{theorem}

\subsection*{ω}

\begin{theorem}[blueprint="infogeometry:canonical:kaehlergeometry:kaehlerinformationgeometry:ω"]
  \inputleannode{InfoGeometry.Canonical.KaehlerGeometry.KaehlerInformationGeometry.ω}
\end{theorem}

\section{InfoGeometry.Canonical.Krein}

\noindent\textbf{Buildable}: yes

\noindent\emph{No exported declarations in the current declaration graph.}

\section{InfoGeometry.Canonical.KreinLadder}

\noindent\textbf{Buildable}: yes

\subsection*{InformationLadder}

\begin{theorem}[blueprint="infogeometry:canonical:kreinladder:informationladder"]
  \inputleannode{InfoGeometry.Canonical.KreinLadder.InformationLadder}
\end{theorem}

\subsection*{drazinProjection}

\begin{theorem}[blueprint="infogeometry:canonical:kreinladder:drazinprojection"]
  \inputleannode{InfoGeometry.Canonical.KreinLadder.drazinProjection}
\end{theorem}

\section{InfoGeometry.Canonical.KreinLadder.InformationLadder}

\subsection*{P}

\begin{theorem}[blueprint="infogeometry:canonical:kreinladder:informationladder:p"]
  \inputleannode{InfoGeometry.Canonical.KreinLadder.InformationLadder.P}
\end{theorem}

\subsection*{RST}

\begin{theorem}[blueprint="infogeometry:canonical:kreinladder:informationladder:rst"]
  \inputleannode{InfoGeometry.Canonical.KreinLadder.InformationLadder.RST}
\end{theorem}

\subsection*{a_op}

\begin{theorem}[blueprint="infogeometry:canonical:kreinladder:informationladder:a_op"]
  \inputleannode{InfoGeometry.Canonical.KreinLadder.InformationLadder.a_op}
\end{theorem}

\subsection*{a_persistent}

\begin{theorem}[blueprint="infogeometry:canonical:kreinladder:informationladder:a_persistent"]
  \inputleannode{InfoGeometry.Canonical.KreinLadder.InformationLadder.a_persistent}
\end{theorem}

\subsection*{annihilation}

\begin{theorem}[blueprint="infogeometry:canonical:kreinladder:informationladder:annihilation"]
  \inputleannode{InfoGeometry.Canonical.KreinLadder.InformationLadder.annihilation}
\end{theorem}

\subsection*{casesOn}

\begin{theorem}[blueprint="infogeometry:canonical:kreinladder:informationladder:caseson"]
  \inputleannode{InfoGeometry.Canonical.KreinLadder.InformationLadder.casesOn}
\end{theorem}

\subsection*{creation}

\begin{theorem}[blueprint="infogeometry:canonical:kreinladder:informationladder:creation"]
  \inputleannode{InfoGeometry.Canonical.KreinLadder.InformationLadder.creation}
\end{theorem}

\subsection*{ctorIdx}

\begin{theorem}[blueprint="infogeometry:canonical:kreinladder:informationladder:ctoridx"]
  \inputleannode{InfoGeometry.Canonical.KreinLadder.InformationLadder.ctorIdx}
\end{theorem}

\subsection*{mk}

\begin{theorem}[blueprint="infogeometry:canonical:kreinladder:informationladder:mk"]
  \inputleannode{InfoGeometry.Canonical.KreinLadder.InformationLadder.mk}
\end{theorem}

\subsection*{noConfusion}

\begin{theorem}[blueprint="infogeometry:canonical:kreinladder:informationladder:noconfusion"]
  \inputleannode{InfoGeometry.Canonical.KreinLadder.InformationLadder.noConfusion}
\end{theorem}

\subsection*{noConfusionType}

\begin{theorem}[blueprint="infogeometry:canonical:kreinladder:informationladder:noconfusiontype"]
  \inputleannode{InfoGeometry.Canonical.KreinLadder.InformationLadder.noConfusionType}
\end{theorem}

\subsection*{rec}

\begin{theorem}[blueprint="infogeometry:canonical:kreinladder:informationladder:rec"]
  \inputleannode{InfoGeometry.Canonical.KreinLadder.InformationLadder.rec}
\end{theorem}

\subsection*{recOn}

\begin{theorem}[blueprint="infogeometry:canonical:kreinladder:informationladder:recon"]
  \inputleannode{InfoGeometry.Canonical.KreinLadder.InformationLadder.recOn}
\end{theorem}

\section{InfoGeometry.Canonical.LLMLegacy}

\noindent\textbf{Buildable}: yes

\noindent\emph{No exported declarations in the current declaration graph.}

\section{InfoGeometry.Canonical.LLN}

\noindent\textbf{Buildable}: yes

\noindent\emph{No exported declarations in the current declaration graph.}

\section{InfoGeometry.Canonical.LorentzianRouting}

\noindent\textbf{Buildable}: yes

\noindent\emph{No exported declarations in the current declaration graph.}

\section{InfoGeometry.Canonical.ManifoldHomology}

\noindent\textbf{Buildable}: yes

\noindent\emph{No exported declarations in the current declaration graph.}

\section{InfoGeometry.Canonical.MixtureOfExperts}

\noindent\textbf{Buildable}: yes

\noindent\emph{No exported declarations in the current declaration graph.}

\section{InfoGeometry.Canonical.MoE}

\subsection*{CliffordLabel}

\begin{theorem}[blueprint="infogeometry:canonical:moe:cliffordlabel"]
  \inputleannode{InfoGeometry.Canonical.MoE.CliffordLabel}
\end{theorem}

\subsection*{CostMatrix}

\begin{theorem}[blueprint="infogeometry:canonical:moe:costmatrix"]
  \inputleannode{InfoGeometry.Canonical.MoE.CostMatrix}
\end{theorem}

\subsection*{Coupling}

\begin{theorem}[blueprint="infogeometry:canonical:moe:coupling"]
  \inputleannode{InfoGeometry.Canonical.MoE.Coupling}
\end{theorem}

\subsection*{Expert}

\begin{theorem}[blueprint="infogeometry:canonical:moe:expert"]
  \inputleannode{InfoGeometry.Canonical.MoE.Expert}
\end{theorem}

\subsection*{ExpertIdx}

\begin{theorem}[blueprint="infogeometry:canonical:moe:expertidx"]
  \inputleannode{InfoGeometry.Canonical.MoE.ExpertIdx}
\end{theorem}

\subsection*{HasMarginals}

\begin{theorem}[blueprint="infogeometry:canonical:moe:hasmarginals"]
  \inputleannode{InfoGeometry.Canonical.MoE.HasMarginals}
\end{theorem}

\subsection*{HasPositiveColSums}

\begin{theorem}[blueprint="infogeometry:canonical:moe:haspositivecolsums"]
  \inputleannode{InfoGeometry.Canonical.MoE.HasPositiveColSums}
\end{theorem}

\subsection*{HasPositiveRowSums}

\begin{theorem}[blueprint="infogeometry:canonical:moe:haspositiverowsums"]
  \inputleannode{InfoGeometry.Canonical.MoE.HasPositiveRowSums}
\end{theorem}

\subsection*{IsBistochasticSwitch}

\begin{theorem}[blueprint="infogeometry:canonical:moe:isbistochasticswitch"]
  \inputleannode{InfoGeometry.Canonical.MoE.IsBistochasticSwitch}
\end{theorem}

\subsection*{Marginal}

\begin{theorem}[blueprint="infogeometry:canonical:moe:marginal"]
  \inputleannode{InfoGeometry.Canonical.MoE.Marginal}
\end{theorem}

\subsection*{MoELayer}

\begin{theorem}[blueprint="infogeometry:canonical:moe:moelayer"]
  \inputleannode{InfoGeometry.Canonical.MoE.MoELayer}
\end{theorem}

\subsection*{ModeDiracMassProfile}

\begin{theorem}[blueprint="infogeometry:canonical:moe:modediracmassprofile"]
  \inputleannode{InfoGeometry.Canonical.MoE.ModeDiracMassProfile}
\end{theorem}

\subsection*{ModeMass}

\begin{theorem}[blueprint="infogeometry:canonical:moe:modemass"]
  \inputleannode{InfoGeometry.Canonical.MoE.ModeMass}
\end{theorem}

\subsection*{ModewiseSplitRep}

\begin{theorem}[blueprint="infogeometry:canonical:moe:modewisesplitrep"]
  \inputleannode{InfoGeometry.Canonical.MoE.ModewiseSplitRep}
\end{theorem}

\subsection*{PermMode}

\begin{theorem}[blueprint="infogeometry:canonical:moe:permmode"]
  \inputleannode{InfoGeometry.Canonical.MoE.PermMode}
\end{theorem}

\subsection*{PermutationMode}

\begin{theorem}[blueprint="infogeometry:canonical:moe:permutationmode"]
  \inputleannode{InfoGeometry.Canonical.MoE.PermutationMode}
\end{theorem}

\subsection*{SinkhornCertificate}

\begin{theorem}[blueprint="infogeometry:canonical:moe:sinkhorncertificate"]
  \inputleannode{InfoGeometry.Canonical.MoE.SinkhornCertificate}
\end{theorem}

\subsection*{SinkhornMatrix}

\begin{theorem}[blueprint="infogeometry:canonical:moe:sinkhornmatrix"]
  \inputleannode{InfoGeometry.Canonical.MoE.SinkhornMatrix}
\end{theorem}

\subsection*{SinkhornPhase}

\begin{theorem}[blueprint="infogeometry:canonical:moe:sinkhornphase"]
  \inputleannode{InfoGeometry.Canonical.MoE.SinkhornPhase}
\end{theorem}

\subsection*{SinkhornStep}

\begin{theorem}[blueprint="infogeometry:canonical:moe:sinkhornstep"]
  \inputleannode{InfoGeometry.Canonical.MoE.SinkhornStep}
\end{theorem}

\subsection*{SinkhornTrajectory}

\begin{theorem}[blueprint="infogeometry:canonical:moe:sinkhorntrajectory"]
  \inputleannode{InfoGeometry.Canonical.MoE.SinkhornTrajectory}
\end{theorem}

\subsection*{SplitCliffordAlg}

\begin{theorem}[blueprint="infogeometry:canonical:moe:splitcliffordalg"]
  \inputleannode{InfoGeometry.Canonical.MoE.SplitCliffordAlg}
\end{theorem}

\subsection*{SplitCliffordSuperData}

\begin{theorem}[blueprint="infogeometry:canonical:moe:splitcliffordsuperdata"]
  \inputleannode{InfoGeometry.Canonical.MoE.SplitCliffordSuperData}
\end{theorem}

\subsection*{bayesianFreeEnergyObjective}

\begin{theorem}[blueprint="infogeometry:canonical:moe:bayesianfreeenergyobjective"]
  \inputleannode{InfoGeometry.Canonical.MoE.bayesianFreeEnergyObjective}
\end{theorem}

\subsection*{bayesianFreeEnergyObjective_eq_regularizedOTObjective}

\begin{theorem}[blueprint="infogeometry:canonical:moe:bayesianfreeenergyobjective_eq_regularizedotobjective"]
  \inputleannode{InfoGeometry.Canonical.MoE.bayesianFreeEnergyObjective_eq_regularizedOTObjective}
\end{theorem}

\subsection*{bayesianLikelihoodKernel}

\begin{theorem}[blueprint="infogeometry:canonical:moe:bayesianlikelihoodkernel"]
  \inputleannode{InfoGeometry.Canonical.MoE.bayesianLikelihoodKernel}
\end{theorem}

\subsection*{bayesianPosteriorCoupling}

\begin{theorem}[blueprint="infogeometry:canonical:moe:bayesianposteriorcoupling"]
  \inputleannode{InfoGeometry.Canonical.MoE.bayesianPosteriorCoupling}
\end{theorem}

\subsection*{canonicalModewiseRep}

\begin{theorem}[blueprint="infogeometry:canonical:moe:canonicalmodewiserep"]
  \inputleannode{InfoGeometry.Canonical.MoE.canonicalModewiseRep}
\end{theorem}

\subsection*{cliffordBasis}

\begin{theorem}[blueprint="infogeometry:canonical:moe:cliffordbasis"]
  \inputleannode{InfoGeometry.Canonical.MoE.cliffordBasis}
\end{theorem}

\subsection*{cliffordBasis_minus}

\begin{theorem}[blueprint="infogeometry:canonical:moe:cliffordbasis_minus"]
  \inputleannode{InfoGeometry.Canonical.MoE.cliffordBasis_minus}
\end{theorem}

\subsection*{cliffordBasis_plus}

\begin{theorem}[blueprint="infogeometry:canonical:moe:cliffordbasis_plus"]
  \inputleannode{InfoGeometry.Canonical.MoE.cliffordBasis_plus}
\end{theorem}

\subsection*{cliffordModeContribution}

\begin{theorem}[blueprint="infogeometry:canonical:moe:cliffordmodecontribution"]
  \inputleannode{InfoGeometry.Canonical.MoE.cliffordModeContribution}
\end{theorem}

\subsection*{cliffordSemanticState}

\begin{theorem}[blueprint="infogeometry:canonical:moe:cliffordsemanticstate"]
  \inputleannode{InfoGeometry.Canonical.MoE.cliffordSemanticState}
\end{theorem}

\subsection*{cliffordSemanticState_coord_sum_eq_weight_sum}

\begin{theorem}[blueprint="infogeometry:canonical:moe:cliffordsemanticstate_coord_sum_eq_weight_sum"]
  \inputleannode{InfoGeometry.Canonical.MoE.cliffordSemanticState_coord_sum_eq_weight_sum}
\end{theorem}

\subsection*{cliffordSemanticState_fst_eq_plusMass}

\begin{theorem}[blueprint="infogeometry:canonical:moe:cliffordsemanticstate_fst_eq_plusmass"]
  \inputleannode{InfoGeometry.Canonical.MoE.cliffordSemanticState_fst_eq_plusMass}
\end{theorem}

\subsection*{cliffordSemanticState_snd_eq_minusMass}

\begin{theorem}[blueprint="infogeometry:canonical:moe:cliffordsemanticstate_snd_eq_minusmass"]
  \inputleannode{InfoGeometry.Canonical.MoE.cliffordSemanticState_snd_eq_minusMass}
\end{theorem}

\subsection*{colBarrierPotential}

\begin{theorem}[blueprint="infogeometry:canonical:moe:colbarrierpotential"]
  \inputleannode{InfoGeometry.Canonical.MoE.colBarrierPotential}
\end{theorem}

\subsection*{colLyapunov}

\begin{theorem}[blueprint="infogeometry:canonical:moe:collyapunov"]
  \inputleannode{InfoGeometry.Canonical.MoE.colLyapunov}
\end{theorem}

\subsection*{colLyapunov_colNormalize_eq_zero}

\begin{theorem}[blueprint="infogeometry:canonical:moe:collyapunov_colnormalize_eq_zero"]
  \inputleannode{InfoGeometry.Canonical.MoE.colLyapunov_colNormalize_eq_zero}
\end{theorem}

\subsection*{colLyapunov_nonneg}

\begin{theorem}[blueprint="infogeometry:canonical:moe:collyapunov_nonneg"]
  \inputleannode{InfoGeometry.Canonical.MoE.colLyapunov_nonneg}
\end{theorem}

\subsection*{colNormalize}

\begin{theorem}[blueprint="infogeometry:canonical:moe:colnormalize"]
  \inputleannode{InfoGeometry.Canonical.MoE.colNormalize}
\end{theorem}

\subsection*{colNormalize_eq_rightDiagonalGauge}

\begin{theorem}[blueprint="infogeometry:canonical:moe:colnormalize_eq_rightdiagonalgauge"]
  \inputleannode{InfoGeometry.Canonical.MoE.colNormalize_eq_rightDiagonalGauge}
\end{theorem}

\subsection*{colNormalize_has_unit_colMarginal}

\begin{theorem}[blueprint="infogeometry:canonical:moe:colnormalize_has_unit_colmarginal"]
  \inputleannode{InfoGeometry.Canonical.MoE.colNormalize_has_unit_colMarginal}
\end{theorem}

\subsection*{colRNBarrier}

\begin{theorem}[blueprint="infogeometry:canonical:moe:colrnbarrier"]
  \inputleannode{InfoGeometry.Canonical.MoE.colRNBarrier}
\end{theorem}

\subsection*{colRNBarrier_colNormalize_eq_zero}

\begin{theorem}[blueprint="infogeometry:canonical:moe:colrnbarrier_colnormalize_eq_zero"]
  \inputleannode{InfoGeometry.Canonical.MoE.colRNBarrier_colNormalize_eq_zero}
\end{theorem}

\subsection*{colRNBarrier_nonneg}

\begin{theorem}[blueprint="infogeometry:canonical:moe:colrnbarrier_nonneg"]
  \inputleannode{InfoGeometry.Canonical.MoE.colRNBarrier_nonneg}
\end{theorem}

\subsection*{colRadonNikodymGenerator}

\begin{theorem}[blueprint="infogeometry:canonical:moe:colradonnikodymgenerator"]
  \inputleannode{InfoGeometry.Canonical.MoE.colRadonNikodymGenerator}
\end{theorem}

\subsection*{colSum}

\begin{theorem}[blueprint="infogeometry:canonical:moe:colsum"]
  \inputleannode{InfoGeometry.Canonical.MoE.colSum}
\end{theorem}

\subsection*{colSum_colNormalize}

\begin{theorem}[blueprint="infogeometry:canonical:moe:colsum_colnormalize"]
  \inputleannode{InfoGeometry.Canonical.MoE.colSum_colNormalize}
\end{theorem}

\subsection*{diracAction}

\begin{theorem}[blueprint="infogeometry:canonical:moe:diracaction"]
  \inputleannode{InfoGeometry.Canonical.MoE.diracAction}
\end{theorem}

\subsection*{diracAction_add_right}

\begin{theorem}[blueprint="infogeometry:canonical:moe:diracaction_add_right"]
  \inputleannode{InfoGeometry.Canonical.MoE.diracAction_add_right}
\end{theorem}

\subsection*{diracAction_def}

\begin{theorem}[blueprint="infogeometry:canonical:moe:diracaction_def"]
  \inputleannode{InfoGeometry.Canonical.MoE.diracAction_def}
\end{theorem}

\subsection*{diracEulerStep}

\begin{theorem}[blueprint="infogeometry:canonical:moe:diraceulerstep"]
  \inputleannode{InfoGeometry.Canonical.MoE.diracEulerStep}
\end{theorem}

\subsection*{diracEulerStep_zero}

\begin{theorem}[blueprint="infogeometry:canonical:moe:diraceulerstep_zero"]
  \inputleannode{InfoGeometry.Canonical.MoE.diracEulerStep_zero}
\end{theorem}

\subsection*{entropicKernel}

\begin{theorem}[blueprint="infogeometry:canonical:moe:entropickernel"]
  \inputleannode{InfoGeometry.Canonical.MoE.entropicKernel}
\end{theorem}

\subsection*{entropicOptimalTransportObjective}

\begin{theorem}[blueprint="infogeometry:canonical:moe:entropicoptimaltransportobjective"]
  \inputleannode{InfoGeometry.Canonical.MoE.entropicOptimalTransportObjective}
\end{theorem}

\subsection*{entropicOptimalTransportObjective_eq_transport_plus_entropy}

\begin{theorem}[blueprint="infogeometry:canonical:moe:entropicoptimaltransportobjective_eq_transport_plus_entropy"]
  \inputleannode{InfoGeometry.Canonical.MoE.entropicOptimalTransportObjective_eq_transport_plus_entropy}
\end{theorem}

\subsection*{exists_clifford_labeled_state_of_bistochastic}

\begin{theorem}[blueprint="infogeometry:canonical:moe:exists_clifford_labeled_state_of_bistochastic"]
  \inputleannode{InfoGeometry.Canonical.MoE.exists_clifford_labeled_state_of_bistochastic}
\end{theorem}

\subsection*{exists_modewiseClifford_rep_of_bistochastic}

\begin{theorem}[blueprint="infogeometry:canonical:moe:exists_modewiseclifford_rep_of_bistochastic"]
  \inputleannode{InfoGeometry.Canonical.MoE.exists_modewiseClifford_rep_of_bistochastic}
\end{theorem}

\subsection*{exists_perm_decomposition_of_bistochastic}

\begin{theorem}[blueprint="infogeometry:canonical:moe:exists_perm_decomposition_of_bistochastic"]
  \inputleannode{InfoGeometry.Canonical.MoE.exists_perm_decomposition_of_bistochastic}
\end{theorem}

\subsection*{exists_perm_decomposition_of_sinkhornBalanced}

\begin{theorem}[blueprint="infogeometry:canonical:moe:exists_perm_decomposition_of_sinkhornbalanced"]
  \inputleannode{InfoGeometry.Canonical.MoE.exists_perm_decomposition_of_sinkhornBalanced}
\end{theorem}

\subsection*{gc_gibbsWeight_eq_normalizedWeights}

\begin{theorem}[blueprint="infogeometry:canonical:moe:gc_gibbsweight_eq_normalizedweights"]
  \inputleannode{InfoGeometry.Canonical.MoE.gc_gibbsWeight_eq_normalizedWeights}
\end{theorem}

\subsection*{gc_partition_eq_routerPartition}

\begin{theorem}[blueprint="infogeometry:canonical:moe:gc_partition_eq_routerpartition"]
  \inputleannode{InfoGeometry.Canonical.MoE.gc_partition_eq_routerPartition}
\end{theorem}

\subsection*{instDecidableEqCliffordLabel}

\begin{theorem}[blueprint="infogeometry:canonical:moe:instdecidableeqcliffordlabel"]
  \inputleannode{InfoGeometry.Canonical.MoE.instDecidableEqCliffordLabel}
\end{theorem}

\subsection*{instDecidableEqSinkhornPhase}

\begin{theorem}[blueprint="infogeometry:canonical:moe:instdecidableeqsinkhornphase"]
  \inputleannode{InfoGeometry.Canonical.MoE.instDecidableEqSinkhornPhase}
\end{theorem}

\subsection*{instReprCliffordLabel}

\begin{theorem}[blueprint="infogeometry:canonical:moe:instreprcliffordlabel"]
  \inputleannode{InfoGeometry.Canonical.MoE.instReprCliffordLabel}
\end{theorem}

\subsection*{instReprSinkhornPhase}

\begin{theorem}[blueprint="infogeometry:canonical:moe:instreprsinkhornphase"]
  \inputleannode{InfoGeometry.Canonical.MoE.instReprSinkhornPhase}
\end{theorem}

\subsection*{labelGenerator}

\begin{theorem}[blueprint="infogeometry:canonical:moe:labelgenerator"]
  \inputleannode{InfoGeometry.Canonical.MoE.labelGenerator}
\end{theorem}

\subsection*{labelGenerator_sq}

\begin{theorem}[blueprint="infogeometry:canonical:moe:labelgenerator_sq"]
  \inputleannode{InfoGeometry.Canonical.MoE.labelGenerator_sq}
\end{theorem}

\subsection*{labelGenerator_sq_minus}

\begin{theorem}[blueprint="infogeometry:canonical:moe:labelgenerator_sq_minus"]
  \inputleannode{InfoGeometry.Canonical.MoE.labelGenerator_sq_minus}
\end{theorem}

\subsection*{labelGenerator_sq_plus}

\begin{theorem}[blueprint="infogeometry:canonical:moe:labelgenerator_sq_plus"]
  \inputleannode{InfoGeometry.Canonical.MoE.labelGenerator_sq_plus}
\end{theorem}

\subsection*{leftWeylScale}

\begin{theorem}[blueprint="infogeometry:canonical:moe:leftweylscale"]
  \inputleannode{InfoGeometry.Canonical.MoE.leftWeylScale}
\end{theorem}

\subsection*{minusMass}

\begin{theorem}[blueprint="infogeometry:canonical:moe:minusmass"]
  \inputleannode{InfoGeometry.Canonical.MoE.minusMass}
\end{theorem}

\subsection*{minusMass_nonneg}

\begin{theorem}[blueprint="infogeometry:canonical:moe:minusmass_nonneg"]
  \inputleannode{InfoGeometry.Canonical.MoE.minusMass_nonneg}
\end{theorem}

\subsection*{mixture_gauge_reduction}

\begin{theorem}[blueprint="infogeometry:canonical:moe:mixture_gauge_reduction"]
  \inputleannode{InfoGeometry.Canonical.MoE.mixture_gauge_reduction}
\end{theorem}

\subsection*{modeDiracAction_respects_grading}

\begin{theorem}[blueprint="infogeometry:canonical:moe:modediracaction_respects_grading"]
  \inputleannode{InfoGeometry.Canonical.MoE.modeDiracAction_respects_grading}
\end{theorem}

\subsection*{modeDiracOperator}

\begin{theorem}[blueprint="infogeometry:canonical:moe:modediracoperator"]
  \inputleannode{InfoGeometry.Canonical.MoE.modeDiracOperator}
\end{theorem}

\subsection*{modeSplitSuperBracket}

\begin{theorem}[blueprint="infogeometry:canonical:moe:modesplitsuperbracket"]
  \inputleannode{InfoGeometry.Canonical.MoE.modeSplitSuperBracket}
\end{theorem}

\subsection*{modeSuperSign}

\begin{theorem}[blueprint="infogeometry:canonical:moe:modesupersign"]
  \inputleannode{InfoGeometry.Canonical.MoE.modeSuperSign}
\end{theorem}

\subsection*{modeWeightEntropy}

\begin{theorem}[blueprint="infogeometry:canonical:moe:modeweightentropy"]
  \inputleannode{InfoGeometry.Canonical.MoE.modeWeightEntropy}
\end{theorem}

\subsection*{modewiseCliffordState}

\begin{theorem}[blueprint="infogeometry:canonical:moe:modewisecliffordstate"]
  \inputleannode{InfoGeometry.Canonical.MoE.modewiseCliffordState}
\end{theorem}

\subsection*{modewiseCliffordStateEven}

\begin{theorem}[blueprint="infogeometry:canonical:moe:modewisecliffordstateeven"]
  \inputleannode{InfoGeometry.Canonical.MoE.modewiseCliffordStateEven}
\end{theorem}

\subsection*{modewiseCliffordStateOdd}

\begin{theorem}[blueprint="infogeometry:canonical:moe:modewisecliffordstateodd"]
  \inputleannode{InfoGeometry.Canonical.MoE.modewiseCliffordStateOdd}
\end{theorem}

\subsection*{modewiseCliffordState_split}

\begin{theorem}[blueprint="infogeometry:canonical:moe:modewisecliffordstate_split"]
  \inputleannode{InfoGeometry.Canonical.MoE.modewiseCliffordState_split}
\end{theorem}

\subsection*{negativeEntropy}

\begin{theorem}[blueprint="infogeometry:canonical:moe:negativeentropy"]
  \inputleannode{InfoGeometry.Canonical.MoE.negativeEntropy}
\end{theorem}

\subsection*{normalizedMixture}

\begin{theorem}[blueprint="infogeometry:canonical:moe:normalizedmixture"]
  \inputleannode{InfoGeometry.Canonical.MoE.normalizedMixture}
\end{theorem}

\subsection*{normalizedWeights}

\begin{theorem}[blueprint="infogeometry:canonical:moe:normalizedweights"]
  \inputleannode{InfoGeometry.Canonical.MoE.normalizedWeights}
\end{theorem}

\subsection*{normalizedWeights_nonneg}

\begin{theorem}[blueprint="infogeometry:canonical:moe:normalizedweights_nonneg"]
  \inputleannode{InfoGeometry.Canonical.MoE.normalizedWeights_nonneg}
\end{theorem}

\subsection*{normalizedWeights_sum_one}

\begin{theorem}[blueprint="infogeometry:canonical:moe:normalizedweights_sum_one"]
  \inputleannode{InfoGeometry.Canonical.MoE.normalizedWeights_sum_one}
\end{theorem}

\subsection*{parity}

\begin{theorem}[blueprint="infogeometry:canonical:moe:parity"]
  \inputleannode{InfoGeometry.Canonical.MoE.parity}
\end{theorem}

\subsection*{parityEntropy}

\begin{theorem}[blueprint="infogeometry:canonical:moe:parityentropy"]
  \inputleannode{InfoGeometry.Canonical.MoE.parityEntropy}
\end{theorem}

\subsection*{parity_minus}

\begin{theorem}[blueprint="infogeometry:canonical:moe:parity_minus"]
  \inputleannode{InfoGeometry.Canonical.MoE.parity_minus}
\end{theorem}

\subsection*{parity_plus}

\begin{theorem}[blueprint="infogeometry:canonical:moe:parity_plus"]
  \inputleannode{InfoGeometry.Canonical.MoE.parity_plus}
\end{theorem}

\subsection*{permutationCliffordSemanticState}

\begin{theorem}[blueprint="infogeometry:canonical:moe:permutationcliffordsemanticstate"]
  \inputleannode{InfoGeometry.Canonical.MoE.permutationCliffordSemanticState}
\end{theorem}

\subsection*{phaseAt}

\begin{theorem}[blueprint="infogeometry:canonical:moe:phaseat"]
  \inputleannode{InfoGeometry.Canonical.MoE.phaseAt}
\end{theorem}

\subsection*{phaseLyapunovAfter}

\begin{theorem}[blueprint="infogeometry:canonical:moe:phaselyapunovafter"]
  \inputleannode{InfoGeometry.Canonical.MoE.phaseLyapunovAfter}
\end{theorem}

\subsection*{phaseLyapunovBefore}

\begin{theorem}[blueprint="infogeometry:canonical:moe:phaselyapunovbefore"]
  \inputleannode{InfoGeometry.Canonical.MoE.phaseLyapunovBefore}
\end{theorem}

\subsection*{phaseLyapunovBefore_nonneg}

\begin{theorem}[blueprint="infogeometry:canonical:moe:phaselyapunovbefore_nonneg"]
  \inputleannode{InfoGeometry.Canonical.MoE.phaseLyapunovBefore_nonneg}
\end{theorem}

\subsection*{phaseRNBarrierAfter}

\begin{theorem}[blueprint="infogeometry:canonical:moe:phasernbarrierafter"]
  \inputleannode{InfoGeometry.Canonical.MoE.phaseRNBarrierAfter}
\end{theorem}

\subsection*{phaseRNBarrierBefore}

\begin{theorem}[blueprint="infogeometry:canonical:moe:phasernbarrierbefore"]
  \inputleannode{InfoGeometry.Canonical.MoE.phaseRNBarrierBefore}
\end{theorem}

\subsection*{phaseRNBarrierBefore_nonneg}

\begin{theorem}[blueprint="infogeometry:canonical:moe:phasernbarrierbefore_nonneg"]
  \inputleannode{InfoGeometry.Canonical.MoE.phaseRNBarrierBefore_nonneg}
\end{theorem}

\subsection*{plusMass}

\begin{theorem}[blueprint="infogeometry:canonical:moe:plusmass"]
  \inputleannode{InfoGeometry.Canonical.MoE.plusMass}
\end{theorem}

\subsection*{plusMass_add_minusMass_eq_sum}

\begin{theorem}[blueprint="infogeometry:canonical:moe:plusmass_add_minusmass_eq_sum"]
  \inputleannode{InfoGeometry.Canonical.MoE.plusMass_add_minusMass_eq_sum}
\end{theorem}

\subsection*{plusMass_nonneg}

\begin{theorem}[blueprint="infogeometry:canonical:moe:plusmass_nonneg"]
  \inputleannode{InfoGeometry.Canonical.MoE.plusMass_nonneg}
\end{theorem}

\subsection*{regularizedOTObjective}

\begin{theorem}[blueprint="infogeometry:canonical:moe:regularizedotobjective"]
  \inputleannode{InfoGeometry.Canonical.MoE.regularizedOTObjective}
\end{theorem}

\subsection*{rightWeylScale}

\begin{theorem}[blueprint="infogeometry:canonical:moe:rightweylscale"]
  \inputleannode{InfoGeometry.Canonical.MoE.rightWeylScale}
\end{theorem}

\subsection*{routerEnergy}

\begin{theorem}[blueprint="infogeometry:canonical:moe:routerenergy"]
  \inputleannode{InfoGeometry.Canonical.MoE.routerEnergy}
\end{theorem}

\subsection*{routerParams}

\begin{theorem}[blueprint="infogeometry:canonical:moe:routerparams"]
  \inputleannode{InfoGeometry.Canonical.MoE.routerParams}
\end{theorem}

\subsection*{routerPartition}

\begin{theorem}[blueprint="infogeometry:canonical:moe:routerpartition"]
  \inputleannode{InfoGeometry.Canonical.MoE.routerPartition}
\end{theorem}

\subsection*{routerPartition_pos}

\begin{theorem}[blueprint="infogeometry:canonical:moe:routerpartition_pos"]
  \inputleannode{InfoGeometry.Canonical.MoE.routerPartition_pos}
\end{theorem}

\subsection*{rowBarrierPotential}

\begin{theorem}[blueprint="infogeometry:canonical:moe:rowbarrierpotential"]
  \inputleannode{InfoGeometry.Canonical.MoE.rowBarrierPotential}
\end{theorem}

\subsection*{rowLyapunov}

\begin{theorem}[blueprint="infogeometry:canonical:moe:rowlyapunov"]
  \inputleannode{InfoGeometry.Canonical.MoE.rowLyapunov}
\end{theorem}

\subsection*{rowLyapunov_nonneg}

\begin{theorem}[blueprint="infogeometry:canonical:moe:rowlyapunov_nonneg"]
  \inputleannode{InfoGeometry.Canonical.MoE.rowLyapunov_nonneg}
\end{theorem}

\subsection*{rowLyapunov_rowNormalize_eq_zero}

\begin{theorem}[blueprint="infogeometry:canonical:moe:rowlyapunov_rownormalize_eq_zero"]
  \inputleannode{InfoGeometry.Canonical.MoE.rowLyapunov_rowNormalize_eq_zero}
\end{theorem}

\subsection*{rowNormalize}

\begin{theorem}[blueprint="infogeometry:canonical:moe:rownormalize"]
  \inputleannode{InfoGeometry.Canonical.MoE.rowNormalize}
\end{theorem}

\subsection*{rowNormalize_eq_leftDiagonalGauge}

\begin{theorem}[blueprint="infogeometry:canonical:moe:rownormalize_eq_leftdiagonalgauge"]
  \inputleannode{InfoGeometry.Canonical.MoE.rowNormalize_eq_leftDiagonalGauge}
\end{theorem}

\subsection*{rowNormalize_has_unit_rowMarginal}

\begin{theorem}[blueprint="infogeometry:canonical:moe:rownormalize_has_unit_rowmarginal"]
  \inputleannode{InfoGeometry.Canonical.MoE.rowNormalize_has_unit_rowMarginal}
\end{theorem}

\subsection*{rowRNBarrier}

\begin{theorem}[blueprint="infogeometry:canonical:moe:rowrnbarrier"]
  \inputleannode{InfoGeometry.Canonical.MoE.rowRNBarrier}
\end{theorem}

\subsection*{rowRNBarrier_nonneg}

\begin{theorem}[blueprint="infogeometry:canonical:moe:rowrnbarrier_nonneg"]
  \inputleannode{InfoGeometry.Canonical.MoE.rowRNBarrier_nonneg}
\end{theorem}

\subsection*{rowRNBarrier_rowNormalize_eq_zero}

\begin{theorem}[blueprint="infogeometry:canonical:moe:rowrnbarrier_rownormalize_eq_zero"]
  \inputleannode{InfoGeometry.Canonical.MoE.rowRNBarrier_rowNormalize_eq_zero}
\end{theorem}

\subsection*{rowRadonNikodymGenerator}

\begin{theorem}[blueprint="infogeometry:canonical:moe:rowradonnikodymgenerator"]
  \inputleannode{InfoGeometry.Canonical.MoE.rowRadonNikodymGenerator}
\end{theorem}

\subsection*{rowSum}

\begin{theorem}[blueprint="infogeometry:canonical:moe:rowsum"]
  \inputleannode{InfoGeometry.Canonical.MoE.rowSum}
\end{theorem}

\subsection*{rowSum_rowNormalize}

\begin{theorem}[blueprint="infogeometry:canonical:moe:rowsum_rownormalize"]
  \inputleannode{InfoGeometry.Canonical.MoE.rowSum_rowNormalize}
\end{theorem}

\subsection*{schroedingerBridgeCoupling}

\begin{theorem}[blueprint="infogeometry:canonical:moe:schroedingerbridgecoupling"]
  \inputleannode{InfoGeometry.Canonical.MoE.schroedingerBridgeCoupling}
\end{theorem}

\subsection*{schroedingerBridgeKernel}

\begin{theorem}[blueprint="infogeometry:canonical:moe:schroedingerbridgekernel"]
  \inputleannode{InfoGeometry.Canonical.MoE.schroedingerBridgeKernel}
\end{theorem}

\subsection*{schroedingerBridgeObjective}

\begin{theorem}[blueprint="infogeometry:canonical:moe:schroedingerbridgeobjective"]
  \inputleannode{InfoGeometry.Canonical.MoE.schroedingerBridgeObjective}
\end{theorem}

\subsection*{schroedingerBridgeStep_monotone}

\begin{theorem}[blueprint="infogeometry:canonical:moe:schroedingerbridgestep_monotone"]
  \inputleannode{InfoGeometry.Canonical.MoE.schroedingerBridgeStep_monotone}
\end{theorem}

\subsection*{schroedingerBridgeStep_radonNikodymBarrier_monotone}

\begin{theorem}[blueprint="infogeometry:canonical:moe:schroedingerbridgestep_radonnikodymbarrier_monotone"]
  \inputleannode{InfoGeometry.Canonical.MoE.schroedingerBridgeStep_radonNikodymBarrier_monotone}
\end{theorem}

\subsection*{sinkhornScaledCoupling}

\begin{theorem}[blueprint="infogeometry:canonical:moe:sinkhornscaledcoupling"]
  \inputleannode{InfoGeometry.Canonical.MoE.sinkhornScaledCoupling}
\end{theorem}

\subsection*{sinkhornScaledCoupling_def}

\begin{theorem}[blueprint="infogeometry:canonical:moe:sinkhornscaledcoupling_def"]
  \inputleannode{InfoGeometry.Canonical.MoE.sinkhornScaledCoupling_def}
\end{theorem}

\subsection*{sinkhornStep_bayesianFreeEnergy_monotone}

\begin{theorem}[blueprint="infogeometry:canonical:moe:sinkhornstep_bayesianfreeenergy_monotone"]
  \inputleannode{InfoGeometry.Canonical.MoE.sinkhornStep_bayesianFreeEnergy_monotone}
\end{theorem}

\subsection*{sinkhornStep_entropicOT_monotone}

\begin{theorem}[blueprint="infogeometry:canonical:moe:sinkhornstep_entropicot_monotone"]
  \inputleannode{InfoGeometry.Canonical.MoE.sinkhornStep_entropicOT_monotone}
\end{theorem}

\subsection*{sinkhornStep_phaseLyapunovAfter_eq_zero}

\begin{theorem}[blueprint="infogeometry:canonical:moe:sinkhornstep_phaselyapunovafter_eq_zero"]
  \inputleannode{InfoGeometry.Canonical.MoE.sinkhornStep_phaseLyapunovAfter_eq_zero}
\end{theorem}

\subsection*{sinkhornStep_phaseLyapunov_monotone}

\begin{theorem}[blueprint="infogeometry:canonical:moe:sinkhornstep_phaselyapunov_monotone"]
  \inputleannode{InfoGeometry.Canonical.MoE.sinkhornStep_phaseLyapunov_monotone}
\end{theorem}

\subsection*{sinkhornStep_phaseRNBarrierAfter_eq_zero}

\begin{theorem}[blueprint="infogeometry:canonical:moe:sinkhornstep_phasernbarrierafter_eq_zero"]
  \inputleannode{InfoGeometry.Canonical.MoE.sinkhornStep_phaseRNBarrierAfter_eq_zero}
\end{theorem}

\subsection*{sinkhornStep_phaseRNBarrier_monotone}

\begin{theorem}[blueprint="infogeometry:canonical:moe:sinkhornstep_phasernbarrier_monotone"]
  \inputleannode{InfoGeometry.Canonical.MoE.sinkhornStep_phaseRNBarrier_monotone}
\end{theorem}

\subsection*{sinkhornStep_radonNikodymBarrier_monotone}

\begin{theorem}[blueprint="infogeometry:canonical:moe:sinkhornstep_radonnikodymbarrier_monotone"]
  \inputleannode{InfoGeometry.Canonical.MoE.sinkhornStep_radonNikodymBarrier_monotone}
\end{theorem}

\subsection*{sinkhornStep_regularizedObjective_monotone}

\begin{theorem}[blueprint="infogeometry:canonical:moe:sinkhornstep_regularizedobjective_monotone"]
  \inputleannode{InfoGeometry.Canonical.MoE.sinkhornStep_regularizedObjective_monotone}
\end{theorem}

\subsection*{sinkhornTwoStep_eq_bayesianPosteriorGauge}

\begin{theorem}[blueprint="infogeometry:canonical:moe:sinkhorntwostep_eq_bayesianposteriorgauge"]
  \inputleannode{InfoGeometry.Canonical.MoE.sinkhornTwoStep_eq_bayesianPosteriorGauge}
\end{theorem}

\subsection*{sinkhornTwoStep_eq_schroedingerBridgeGauge}

\begin{theorem}[blueprint="infogeometry:canonical:moe:sinkhorntwostep_eq_schroedingerbridgegauge"]
  \inputleannode{InfoGeometry.Canonical.MoE.sinkhornTwoStep_eq_schroedingerBridgeGauge}
\end{theorem}

\subsection*{sinkhornTwoStep_eq_twoSidedGauge}

\begin{theorem}[blueprint="infogeometry:canonical:moe:sinkhorntwostep_eq_twosidedgauge"]
  \inputleannode{InfoGeometry.Canonical.MoE.sinkhornTwoStep_eq_twoSidedGauge}
\end{theorem}

\subsection*{sinkhornTwoStep_eq_weylGauge}

\begin{theorem}[blueprint="infogeometry:canonical:moe:sinkhorntwostep_eq_weylgauge"]
  \inputleannode{InfoGeometry.Canonical.MoE.sinkhornTwoStep_eq_weylGauge}
\end{theorem}

\subsection*{splitB11_splitBasis_orthogonal}

\begin{theorem}[blueprint="infogeometry:canonical:moe:splitb11_splitbasis_orthogonal"]
  \inputleannode{InfoGeometry.Canonical.MoE.splitB11_splitBasis_orthogonal}
\end{theorem}

\subsection*{splitB11_symm}

\begin{theorem}[blueprint="infogeometry:canonical:moe:splitb11_symm"]
  \inputleannode{InfoGeometry.Canonical.MoE.splitB11_symm}
\end{theorem}

\subsection*{splitBasisMinus}

\begin{theorem}[blueprint="infogeometry:canonical:moe:splitbasisminus"]
  \inputleannode{InfoGeometry.Canonical.MoE.splitBasisMinus}
\end{theorem}

\subsection*{splitBasisPlus}

\begin{theorem}[blueprint="infogeometry:canonical:moe:splitbasisplus"]
  \inputleannode{InfoGeometry.Canonical.MoE.splitBasisPlus}
\end{theorem}

\subsection*{splitQ11_splitBasisMinus}

\begin{theorem}[blueprint="infogeometry:canonical:moe:splitq11_splitbasisminus"]
  \inputleannode{InfoGeometry.Canonical.MoE.splitQ11_splitBasisMinus}
\end{theorem}

\subsection*{splitQ11_splitBasisPlus}

\begin{theorem}[blueprint="infogeometry:canonical:moe:splitq11_splitbasisplus"]
  \inputleannode{InfoGeometry.Canonical.MoE.splitQ11_splitBasisPlus}
\end{theorem}

\subsection*{splitSuperBracket}

\begin{theorem}[blueprint="infogeometry:canonical:moe:splitsuperbracket"]
  \inputleannode{InfoGeometry.Canonical.MoE.splitSuperBracket}
\end{theorem}

\subsection*{splitSuperBracket_minus_minus}

\begin{theorem}[blueprint="infogeometry:canonical:moe:splitsuperbracket_minus_minus"]
  \inputleannode{InfoGeometry.Canonical.MoE.splitSuperBracket_minus_minus}
\end{theorem}

\subsection*{splitSuperBracket_plus_left}

\begin{theorem}[blueprint="infogeometry:canonical:moe:splitsuperbracket_plus_left"]
  \inputleannode{InfoGeometry.Canonical.MoE.splitSuperBracket_plus_left}
\end{theorem}

\subsection*{splitSuperBracket_plus_right}

\begin{theorem}[blueprint="infogeometry:canonical:moe:splitsuperbracket_plus_right"]
  \inputleannode{InfoGeometry.Canonical.MoE.splitSuperBracket_plus_right}
\end{theorem}

\subsection*{superSign}

\begin{theorem}[blueprint="infogeometry:canonical:moe:supersign"]
  \inputleannode{InfoGeometry.Canonical.MoE.superSign}
\end{theorem}

\subsection*{superSign_minus_minus}

\begin{theorem}[blueprint="infogeometry:canonical:moe:supersign_minus_minus"]
  \inputleannode{InfoGeometry.Canonical.MoE.superSign_minus_minus}
\end{theorem}

\subsection*{superSign_plus_left}

\begin{theorem}[blueprint="infogeometry:canonical:moe:supersign_plus_left"]
  \inputleannode{InfoGeometry.Canonical.MoE.superSign_plus_left}
\end{theorem}

\subsection*{superSign_plus_right}

\begin{theorem}[blueprint="infogeometry:canonical:moe:supersign_plus_right"]
  \inputleannode{InfoGeometry.Canonical.MoE.superSign_plus_right}
\end{theorem}

\subsection*{switchMatrix}

\begin{theorem}[blueprint="infogeometry:canonical:moe:switchmatrix"]
  \inputleannode{InfoGeometry.Canonical.MoE.switchMatrix}
\end{theorem}

\subsection*{switchMatrix_apply}

\begin{theorem}[blueprint="infogeometry:canonical:moe:switchmatrix_apply"]
  \inputleannode{InfoGeometry.Canonical.MoE.switchMatrix_apply}
\end{theorem}

\subsection*{switchMatrix_mem_doublyStochastic}

\begin{theorem}[blueprint="infogeometry:canonical:moe:switchmatrix_mem_doublystochastic"]
  \inputleannode{InfoGeometry.Canonical.MoE.switchMatrix_mem_doublyStochastic}
\end{theorem}

\subsection*{switchMatrix_mem_rowStochastic}

\begin{theorem}[blueprint="infogeometry:canonical:moe:switchmatrix_mem_rowstochastic"]
  \inputleannode{InfoGeometry.Canonical.MoE.switchMatrix_mem_rowStochastic}
\end{theorem}

\subsection*{switchMatrix_row_sum_one}

\begin{theorem}[blueprint="infogeometry:canonical:moe:switchmatrix_row_sum_one"]
  \inputleannode{InfoGeometry.Canonical.MoE.switchMatrix_row_sum_one}
\end{theorem}

\subsection*{trajectoryLyapunov}

\begin{theorem}[blueprint="infogeometry:canonical:moe:trajectorylyapunov"]
  \inputleannode{InfoGeometry.Canonical.MoE.trajectoryLyapunov}
\end{theorem}

\subsection*{trajectoryLyapunovNext}

\begin{theorem}[blueprint="infogeometry:canonical:moe:trajectorylyapunovnext"]
  \inputleannode{InfoGeometry.Canonical.MoE.trajectoryLyapunovNext}
\end{theorem}

\subsection*{trajectoryLyapunovNext_eq_zero}

\begin{theorem}[blueprint="infogeometry:canonical:moe:trajectorylyapunovnext_eq_zero"]
  \inputleannode{InfoGeometry.Canonical.MoE.trajectoryLyapunovNext_eq_zero}
\end{theorem}

\subsection*{trajectoryLyapunov_monotone}

\begin{theorem}[blueprint="infogeometry:canonical:moe:trajectorylyapunov_monotone"]
  \inputleannode{InfoGeometry.Canonical.MoE.trajectoryLyapunov_monotone}
\end{theorem}

\subsection*{trajectoryRNBarrier}

\begin{theorem}[blueprint="infogeometry:canonical:moe:trajectoryrnbarrier"]
  \inputleannode{InfoGeometry.Canonical.MoE.trajectoryRNBarrier}
\end{theorem}

\subsection*{trajectoryRNBarrierNext}

\begin{theorem}[blueprint="infogeometry:canonical:moe:trajectoryrnbarriernext"]
  \inputleannode{InfoGeometry.Canonical.MoE.trajectoryRNBarrierNext}
\end{theorem}

\subsection*{trajectoryRNBarrierNext_eq_zero}

\begin{theorem}[blueprint="infogeometry:canonical:moe:trajectoryrnbarriernext_eq_zero"]
  \inputleannode{InfoGeometry.Canonical.MoE.trajectoryRNBarrierNext_eq_zero}
\end{theorem}

\subsection*{trajectoryRNBarrier_monotone}

\begin{theorem}[blueprint="infogeometry:canonical:moe:trajectoryrnbarrier_monotone"]
  \inputleannode{InfoGeometry.Canonical.MoE.trajectoryRNBarrier_monotone}
\end{theorem}

\subsection*{transportCost}

\begin{theorem}[blueprint="infogeometry:canonical:moe:transportcost"]
  \inputleannode{InfoGeometry.Canonical.MoE.transportCost}
\end{theorem}

\subsection*{unnormalizedMixture}

\begin{theorem}[blueprint="infogeometry:canonical:moe:unnormalizedmixture"]
  \inputleannode{InfoGeometry.Canonical.MoE.unnormalizedMixture}
\end{theorem}

\subsection*{unnormalizedWeights}

\begin{theorem}[blueprint="infogeometry:canonical:moe:unnormalizedweights"]
  \inputleannode{InfoGeometry.Canonical.MoE.unnormalizedWeights}
\end{theorem}

\subsection*{weightedModeCliffordState}

\begin{theorem}[blueprint="infogeometry:canonical:moe:weightedmodecliffordstate"]
  \inputleannode{InfoGeometry.Canonical.MoE.weightedModeCliffordState}
\end{theorem}

\subsection*{weightedModeCliffordStateMinus}

\begin{theorem}[blueprint="infogeometry:canonical:moe:weightedmodecliffordstateminus"]
  \inputleannode{InfoGeometry.Canonical.MoE.weightedModeCliffordStateMinus}
\end{theorem}

\subsection*{weightedModeCliffordStatePlus}

\begin{theorem}[blueprint="infogeometry:canonical:moe:weightedmodecliffordstateplus"]
  \inputleannode{InfoGeometry.Canonical.MoE.weightedModeCliffordStatePlus}
\end{theorem}

\subsection*{weightedModeCliffordState_split}

\begin{theorem}[blueprint="infogeometry:canonical:moe:weightedmodecliffordstate_split"]
  \inputleannode{InfoGeometry.Canonical.MoE.weightedModeCliffordState_split}
\end{theorem}

\section{InfoGeometry.Canonical.MoE.CliffordLabel}

\subsection*{casesOn}

\begin{theorem}[blueprint="infogeometry:canonical:moe:cliffordlabel:caseson"]
  \inputleannode{InfoGeometry.Canonical.MoE.CliffordLabel.casesOn}
\end{theorem}

\subsection*{ctorElim}

\begin{theorem}[blueprint="infogeometry:canonical:moe:cliffordlabel:ctorelim"]
  \inputleannode{InfoGeometry.Canonical.MoE.CliffordLabel.ctorElim}
\end{theorem}

\subsection*{ctorElimType}

\begin{theorem}[blueprint="infogeometry:canonical:moe:cliffordlabel:ctorelimtype"]
  \inputleannode{InfoGeometry.Canonical.MoE.CliffordLabel.ctorElimType}
\end{theorem}

\subsection*{ctorIdx}

\begin{theorem}[blueprint="infogeometry:canonical:moe:cliffordlabel:ctoridx"]
  \inputleannode{InfoGeometry.Canonical.MoE.CliffordLabel.ctorIdx}
\end{theorem}

\subsection*{minus}

\begin{theorem}[blueprint="infogeometry:canonical:moe:cliffordlabel:minus"]
  \inputleannode{InfoGeometry.Canonical.MoE.CliffordLabel.minus}
\end{theorem}

\subsection*{noConfusion}

\begin{theorem}[blueprint="infogeometry:canonical:moe:cliffordlabel:noconfusion"]
  \inputleannode{InfoGeometry.Canonical.MoE.CliffordLabel.noConfusion}
\end{theorem}

\subsection*{noConfusionType}

\begin{theorem}[blueprint="infogeometry:canonical:moe:cliffordlabel:noconfusiontype"]
  \inputleannode{InfoGeometry.Canonical.MoE.CliffordLabel.noConfusionType}
\end{theorem}

\subsection*{ofNat}

\begin{theorem}[blueprint="infogeometry:canonical:moe:cliffordlabel:ofnat"]
  \inputleannode{InfoGeometry.Canonical.MoE.CliffordLabel.ofNat}
\end{theorem}

\subsection*{ofNat_ctorIdx}

\begin{theorem}[blueprint="infogeometry:canonical:moe:cliffordlabel:ofnat_ctoridx"]
  \inputleannode{InfoGeometry.Canonical.MoE.CliffordLabel.ofNat_ctorIdx}
\end{theorem}

\subsection*{plus}

\begin{theorem}[blueprint="infogeometry:canonical:moe:cliffordlabel:plus"]
  \inputleannode{InfoGeometry.Canonical.MoE.CliffordLabel.plus}
\end{theorem}

\subsection*{rec}

\begin{theorem}[blueprint="infogeometry:canonical:moe:cliffordlabel:rec"]
  \inputleannode{InfoGeometry.Canonical.MoE.CliffordLabel.rec}
\end{theorem}

\subsection*{recOn}

\begin{theorem}[blueprint="infogeometry:canonical:moe:cliffordlabel:recon"]
  \inputleannode{InfoGeometry.Canonical.MoE.CliffordLabel.recOn}
\end{theorem}

\subsection*{toCtorIdx}

\begin{theorem}[blueprint="infogeometry:canonical:moe:cliffordlabel:toctoridx"]
  \inputleannode{InfoGeometry.Canonical.MoE.CliffordLabel.toCtorIdx}
\end{theorem}

\section{InfoGeometry.Canonical.MoE.CliffordLabel.minus}

\subsection*{elim}

\begin{theorem}[blueprint="infogeometry:canonical:moe:cliffordlabel:minus:elim"]
  \inputleannode{InfoGeometry.Canonical.MoE.CliffordLabel.minus.elim}
\end{theorem}

\section{InfoGeometry.Canonical.MoE.CliffordLabel.plus}

\subsection*{elim}

\begin{theorem}[blueprint="infogeometry:canonical:moe:cliffordlabel:plus:elim"]
  \inputleannode{InfoGeometry.Canonical.MoE.CliffordLabel.plus.elim}
\end{theorem}

\section{InfoGeometry.Canonical.MoE.Expert}

\subsection*{apply}

\begin{theorem}[blueprint="infogeometry:canonical:moe:expert:apply"]
  \inputleannode{InfoGeometry.Canonical.MoE.Expert.apply}
\end{theorem}

\subsection*{casesOn}

\begin{theorem}[blueprint="infogeometry:canonical:moe:expert:caseson"]
  \inputleannode{InfoGeometry.Canonical.MoE.Expert.casesOn}
\end{theorem}

\subsection*{ctorIdx}

\begin{theorem}[blueprint="infogeometry:canonical:moe:expert:ctoridx"]
  \inputleannode{InfoGeometry.Canonical.MoE.Expert.ctorIdx}
\end{theorem}

\subsection*{mk}

\begin{theorem}[blueprint="infogeometry:canonical:moe:expert:mk"]
  \inputleannode{InfoGeometry.Canonical.MoE.Expert.mk}
\end{theorem}

\subsection*{noConfusion}

\begin{theorem}[blueprint="infogeometry:canonical:moe:expert:noconfusion"]
  \inputleannode{InfoGeometry.Canonical.MoE.Expert.noConfusion}
\end{theorem}

\subsection*{noConfusionType}

\begin{theorem}[blueprint="infogeometry:canonical:moe:expert:noconfusiontype"]
  \inputleannode{InfoGeometry.Canonical.MoE.Expert.noConfusionType}
\end{theorem}

\subsection*{rec}

\begin{theorem}[blueprint="infogeometry:canonical:moe:expert:rec"]
  \inputleannode{InfoGeometry.Canonical.MoE.Expert.rec}
\end{theorem}

\subsection*{recOn}

\begin{theorem}[blueprint="infogeometry:canonical:moe:expert:recon"]
  \inputleannode{InfoGeometry.Canonical.MoE.Expert.recOn}
\end{theorem}

\section{InfoGeometry.Canonical.MoE.MoELayer}

\subsection*{casesOn}

\begin{theorem}[blueprint="infogeometry:canonical:moe:moelayer:caseson"]
  \inputleannode{InfoGeometry.Canonical.MoE.MoELayer.casesOn}
\end{theorem}

\subsection*{ctorIdx}

\begin{theorem}[blueprint="infogeometry:canonical:moe:moelayer:ctoridx"]
  \inputleannode{InfoGeometry.Canonical.MoE.MoELayer.ctorIdx}
\end{theorem}

\subsection*{experts}

\begin{theorem}[blueprint="infogeometry:canonical:moe:moelayer:experts"]
  \inputleannode{InfoGeometry.Canonical.MoE.MoELayer.experts}
\end{theorem}

\subsection*{mk}

\begin{theorem}[blueprint="infogeometry:canonical:moe:moelayer:mk"]
  \inputleannode{InfoGeometry.Canonical.MoE.MoELayer.mk}
\end{theorem}

\subsection*{noConfusion}

\begin{theorem}[blueprint="infogeometry:canonical:moe:moelayer:noconfusion"]
  \inputleannode{InfoGeometry.Canonical.MoE.MoELayer.noConfusion}
\end{theorem}

\subsection*{noConfusionType}

\begin{theorem}[blueprint="infogeometry:canonical:moe:moelayer:noconfusiontype"]
  \inputleannode{InfoGeometry.Canonical.MoE.MoELayer.noConfusionType}
\end{theorem}

\subsection*{rec}

\begin{theorem}[blueprint="infogeometry:canonical:moe:moelayer:rec"]
  \inputleannode{InfoGeometry.Canonical.MoE.MoELayer.rec}
\end{theorem}

\subsection*{recOn}

\begin{theorem}[blueprint="infogeometry:canonical:moe:moelayer:recon"]
  \inputleannode{InfoGeometry.Canonical.MoE.MoELayer.recOn}
\end{theorem}

\section{InfoGeometry.Canonical.MoE.SinkhornCertificate}

\subsection*{balanced_mem_doublyStochastic}

\begin{theorem}[blueprint="infogeometry:canonical:moe:sinkhorncertificate:balanced_mem_doublystochastic"]
  \inputleannode{InfoGeometry.Canonical.MoE.SinkhornCertificate.balanced_mem_doublyStochastic}
\end{theorem}

\subsection*{casesOn}

\begin{theorem}[blueprint="infogeometry:canonical:moe:sinkhorncertificate:caseson"]
  \inputleannode{InfoGeometry.Canonical.MoE.SinkhornCertificate.casesOn}
\end{theorem}

\subsection*{ctorIdx}

\begin{theorem}[blueprint="infogeometry:canonical:moe:sinkhorncertificate:ctoridx"]
  \inputleannode{InfoGeometry.Canonical.MoE.SinkhornCertificate.ctorIdx}
\end{theorem}

\subsection*{leftScale}

\begin{theorem}[blueprint="infogeometry:canonical:moe:sinkhorncertificate:leftscale"]
  \inputleannode{InfoGeometry.Canonical.MoE.SinkhornCertificate.leftScale}
\end{theorem}

\subsection*{leftScale_pos}

\begin{theorem}[blueprint="infogeometry:canonical:moe:sinkhorncertificate:leftscale_pos"]
  \inputleannode{InfoGeometry.Canonical.MoE.SinkhornCertificate.leftScale_pos}
\end{theorem}

\subsection*{mk}

\begin{theorem}[blueprint="infogeometry:canonical:moe:sinkhorncertificate:mk"]
  \inputleannode{InfoGeometry.Canonical.MoE.SinkhornCertificate.mk}
\end{theorem}

\subsection*{noConfusion}

\begin{theorem}[blueprint="infogeometry:canonical:moe:sinkhorncertificate:noconfusion"]
  \inputleannode{InfoGeometry.Canonical.MoE.SinkhornCertificate.noConfusion}
\end{theorem}

\subsection*{noConfusionType}

\begin{theorem}[blueprint="infogeometry:canonical:moe:sinkhorncertificate:noconfusiontype"]
  \inputleannode{InfoGeometry.Canonical.MoE.SinkhornCertificate.noConfusionType}
\end{theorem}

\subsection*{rec}

\begin{theorem}[blueprint="infogeometry:canonical:moe:sinkhorncertificate:rec"]
  \inputleannode{InfoGeometry.Canonical.MoE.SinkhornCertificate.rec}
\end{theorem}

\subsection*{recOn}

\begin{theorem}[blueprint="infogeometry:canonical:moe:sinkhorncertificate:recon"]
  \inputleannode{InfoGeometry.Canonical.MoE.SinkhornCertificate.recOn}
\end{theorem}

\subsection*{rightScale}

\begin{theorem}[blueprint="infogeometry:canonical:moe:sinkhorncertificate:rightscale"]
  \inputleannode{InfoGeometry.Canonical.MoE.SinkhornCertificate.rightScale}
\end{theorem}

\subsection*{rightScale_pos}

\begin{theorem}[blueprint="infogeometry:canonical:moe:sinkhorncertificate:rightscale_pos"]
  \inputleannode{InfoGeometry.Canonical.MoE.SinkhornCertificate.rightScale_pos}
\end{theorem}

\section{InfoGeometry.Canonical.MoE.SinkhornPhase}

\subsection*{casesOn}

\begin{theorem}[blueprint="infogeometry:canonical:moe:sinkhornphase:caseson"]
  \inputleannode{InfoGeometry.Canonical.MoE.SinkhornPhase.casesOn}
\end{theorem}

\subsection*{col}

\begin{theorem}[blueprint="infogeometry:canonical:moe:sinkhornphase:col"]
  \inputleannode{InfoGeometry.Canonical.MoE.SinkhornPhase.col}
\end{theorem}

\subsection*{ctorElim}

\begin{theorem}[blueprint="infogeometry:canonical:moe:sinkhornphase:ctorelim"]
  \inputleannode{InfoGeometry.Canonical.MoE.SinkhornPhase.ctorElim}
\end{theorem}

\subsection*{ctorElimType}

\begin{theorem}[blueprint="infogeometry:canonical:moe:sinkhornphase:ctorelimtype"]
  \inputleannode{InfoGeometry.Canonical.MoE.SinkhornPhase.ctorElimType}
\end{theorem}

\subsection*{ctorIdx}

\begin{theorem}[blueprint="infogeometry:canonical:moe:sinkhornphase:ctoridx"]
  \inputleannode{InfoGeometry.Canonical.MoE.SinkhornPhase.ctorIdx}
\end{theorem}

\subsection*{noConfusion}

\begin{theorem}[blueprint="infogeometry:canonical:moe:sinkhornphase:noconfusion"]
  \inputleannode{InfoGeometry.Canonical.MoE.SinkhornPhase.noConfusion}
\end{theorem}

\subsection*{noConfusionType}

\begin{theorem}[blueprint="infogeometry:canonical:moe:sinkhornphase:noconfusiontype"]
  \inputleannode{InfoGeometry.Canonical.MoE.SinkhornPhase.noConfusionType}
\end{theorem}

\subsection*{ofNat}

\begin{theorem}[blueprint="infogeometry:canonical:moe:sinkhornphase:ofnat"]
  \inputleannode{InfoGeometry.Canonical.MoE.SinkhornPhase.ofNat}
\end{theorem}

\subsection*{ofNat_ctorIdx}

\begin{theorem}[blueprint="infogeometry:canonical:moe:sinkhornphase:ofnat_ctoridx"]
  \inputleannode{InfoGeometry.Canonical.MoE.SinkhornPhase.ofNat_ctorIdx}
\end{theorem}

\subsection*{rec}

\begin{theorem}[blueprint="infogeometry:canonical:moe:sinkhornphase:rec"]
  \inputleannode{InfoGeometry.Canonical.MoE.SinkhornPhase.rec}
\end{theorem}

\subsection*{recOn}

\begin{theorem}[blueprint="infogeometry:canonical:moe:sinkhornphase:recon"]
  \inputleannode{InfoGeometry.Canonical.MoE.SinkhornPhase.recOn}
\end{theorem}

\subsection*{row}

\begin{theorem}[blueprint="infogeometry:canonical:moe:sinkhornphase:row"]
  \inputleannode{InfoGeometry.Canonical.MoE.SinkhornPhase.row}
\end{theorem}

\subsection*{toCtorIdx}

\begin{theorem}[blueprint="infogeometry:canonical:moe:sinkhornphase:toctoridx"]
  \inputleannode{InfoGeometry.Canonical.MoE.SinkhornPhase.toCtorIdx}
\end{theorem}

\section{InfoGeometry.Canonical.MoE.SinkhornPhase.col}

\subsection*{elim}

\begin{theorem}[blueprint="infogeometry:canonical:moe:sinkhornphase:col:elim"]
  \inputleannode{InfoGeometry.Canonical.MoE.SinkhornPhase.col.elim}
\end{theorem}

\section{InfoGeometry.Canonical.MoE.SinkhornPhase.row}

\subsection*{elim}

\begin{theorem}[blueprint="infogeometry:canonical:moe:sinkhornphase:row:elim"]
  \inputleannode{InfoGeometry.Canonical.MoE.SinkhornPhase.row.elim}
\end{theorem}

\section{InfoGeometry.Canonical.MoE.SinkhornTrajectory}

\subsection*{casesOn}

\begin{theorem}[blueprint="infogeometry:canonical:moe:sinkhorntrajectory:caseson"]
  \inputleannode{InfoGeometry.Canonical.MoE.SinkhornTrajectory.casesOn}
\end{theorem}

\subsection*{ctorIdx}

\begin{theorem}[blueprint="infogeometry:canonical:moe:sinkhorntrajectory:ctoridx"]
  \inputleannode{InfoGeometry.Canonical.MoE.SinkhornTrajectory.ctorIdx}
\end{theorem}

\subsection*{mk}

\begin{theorem}[blueprint="infogeometry:canonical:moe:sinkhorntrajectory:mk"]
  \inputleannode{InfoGeometry.Canonical.MoE.SinkhornTrajectory.mk}
\end{theorem}

\subsection*{noConfusion}

\begin{theorem}[blueprint="infogeometry:canonical:moe:sinkhorntrajectory:noconfusion"]
  \inputleannode{InfoGeometry.Canonical.MoE.SinkhornTrajectory.noConfusion}
\end{theorem}

\subsection*{noConfusionType}

\begin{theorem}[blueprint="infogeometry:canonical:moe:sinkhorntrajectory:noconfusiontype"]
  \inputleannode{InfoGeometry.Canonical.MoE.SinkhornTrajectory.noConfusionType}
\end{theorem}

\subsection*{rec}

\begin{theorem}[blueprint="infogeometry:canonical:moe:sinkhorntrajectory:rec"]
  \inputleannode{InfoGeometry.Canonical.MoE.SinkhornTrajectory.rec}
\end{theorem}

\subsection*{recOn}

\begin{theorem}[blueprint="infogeometry:canonical:moe:sinkhorntrajectory:recon"]
  \inputleannode{InfoGeometry.Canonical.MoE.SinkhornTrajectory.recOn}
\end{theorem}

\subsection*{state}

\begin{theorem}[blueprint="infogeometry:canonical:moe:sinkhorntrajectory:state"]
  \inputleannode{InfoGeometry.Canonical.MoE.SinkhornTrajectory.state}
\end{theorem}

\subsection*{step}

\begin{theorem}[blueprint="infogeometry:canonical:moe:sinkhorntrajectory:step"]
  \inputleannode{InfoGeometry.Canonical.MoE.SinkhornTrajectory.step}
\end{theorem}

\section{InfoGeometry.Canonical.MoE.SplitCliffordSuperData}

\subsection*{casesOn}

\begin{theorem}[blueprint="infogeometry:canonical:moe:splitcliffordsuperdata:caseson"]
  \inputleannode{InfoGeometry.Canonical.MoE.SplitCliffordSuperData.casesOn}
\end{theorem}

\subsection*{ctorIdx}

\begin{theorem}[blueprint="infogeometry:canonical:moe:splitcliffordsuperdata:ctoridx"]
  \inputleannode{InfoGeometry.Canonical.MoE.SplitCliffordSuperData.ctorIdx}
\end{theorem}

\subsection*{label}

\begin{theorem}[blueprint="infogeometry:canonical:moe:splitcliffordsuperdata:label"]
  \inputleannode{InfoGeometry.Canonical.MoE.SplitCliffordSuperData.label}
\end{theorem}

\subsection*{mk}

\begin{theorem}[blueprint="infogeometry:canonical:moe:splitcliffordsuperdata:mk"]
  \inputleannode{InfoGeometry.Canonical.MoE.SplitCliffordSuperData.mk}
\end{theorem}

\subsection*{noConfusion}

\begin{theorem}[blueprint="infogeometry:canonical:moe:splitcliffordsuperdata:noconfusion"]
  \inputleannode{InfoGeometry.Canonical.MoE.SplitCliffordSuperData.noConfusion}
\end{theorem}

\subsection*{noConfusionType}

\begin{theorem}[blueprint="infogeometry:canonical:moe:splitcliffordsuperdata:noconfusiontype"]
  \inputleannode{InfoGeometry.Canonical.MoE.SplitCliffordSuperData.noConfusionType}
\end{theorem}

\subsection*{rec}

\begin{theorem}[blueprint="infogeometry:canonical:moe:splitcliffordsuperdata:rec"]
  \inputleannode{InfoGeometry.Canonical.MoE.SplitCliffordSuperData.rec}
\end{theorem}

\subsection*{recOn}

\begin{theorem}[blueprint="infogeometry:canonical:moe:splitcliffordsuperdata:recon"]
  \inputleannode{InfoGeometry.Canonical.MoE.SplitCliffordSuperData.recOn}
\end{theorem}

\section{InfoGeometry.Canonical.MoE.cliffordBasis}

\subsection*{eq_1}

\begin{theorem}[blueprint="infogeometry:canonical:moe:cliffordbasis:eq_1"]
  \inputleannode{InfoGeometry.Canonical.MoE.cliffordBasis.eq_1}
\end{theorem}

\subsection*{eq_2}

\begin{theorem}[blueprint="infogeometry:canonical:moe:cliffordbasis:eq_2"]
  \inputleannode{InfoGeometry.Canonical.MoE.cliffordBasis.eq_2}
\end{theorem}

\section{InfoGeometry.Canonical.MoE.colNormalize}

\subsection*{congr_simp}

\begin{theorem}[blueprint="infogeometry:canonical:moe:colnormalize:congr_simp"]
  \inputleannode{InfoGeometry.Canonical.MoE.colNormalize.congr_simp}
\end{theorem}

\subsection*{eq_1}

\begin{theorem}[blueprint="infogeometry:canonical:moe:colnormalize:eq_1"]
  \inputleannode{InfoGeometry.Canonical.MoE.colNormalize.eq_1}
\end{theorem}

\section{InfoGeometry.Canonical.MoE.colSum}

\subsection*{eq_1}

\begin{theorem}[blueprint="infogeometry:canonical:moe:colsum:eq_1"]
  \inputleannode{InfoGeometry.Canonical.MoE.colSum.eq_1}
\end{theorem}

\section{InfoGeometry.Canonical.MoE.diracAction}

\subsection*{eq_1}

\begin{theorem}[blueprint="infogeometry:canonical:moe:diracaction:eq_1"]
  \inputleannode{InfoGeometry.Canonical.MoE.diracAction.eq_1}
\end{theorem}

\section{InfoGeometry.Canonical.MoE.diracEulerStep}

\subsection*{eq_1}

\begin{theorem}[blueprint="infogeometry:canonical:moe:diraceulerstep:eq_1"]
  \inputleannode{InfoGeometry.Canonical.MoE.diracEulerStep.eq_1}
\end{theorem}

\section{InfoGeometry.Canonical.MoE.instReprCliffordLabel}

\subsection*{repr}

\begin{theorem}[blueprint="infogeometry:canonical:moe:instreprcliffordlabel:repr"]
  \inputleannode{InfoGeometry.Canonical.MoE.instReprCliffordLabel.repr}
\end{theorem}

\section{InfoGeometry.Canonical.MoE.instReprSinkhornPhase}

\subsection*{repr}

\begin{theorem}[blueprint="infogeometry:canonical:moe:instreprsinkhornphase:repr"]
  \inputleannode{InfoGeometry.Canonical.MoE.instReprSinkhornPhase.repr}
\end{theorem}

\section{InfoGeometry.Canonical.MoE.labelGenerator}

\subsection*{eq_1}

\begin{theorem}[blueprint="infogeometry:canonical:moe:labelgenerator:eq_1"]
  \inputleannode{InfoGeometry.Canonical.MoE.labelGenerator.eq_1}
\end{theorem}

\section{InfoGeometry.Canonical.MoE.leftWeylScale}

\subsection*{eq_1}

\begin{theorem}[blueprint="infogeometry:canonical:moe:leftweylscale:eq_1"]
  \inputleannode{InfoGeometry.Canonical.MoE.leftWeylScale.eq_1}
\end{theorem}

\section{InfoGeometry.Canonical.MoE.phaseLyapunovAfter}

\subsection*{eq_1}

\begin{theorem}[blueprint="infogeometry:canonical:moe:phaselyapunovafter:eq_1"]
  \inputleannode{InfoGeometry.Canonical.MoE.phaseLyapunovAfter.eq_1}
\end{theorem}

\subsection*{eq_2}

\begin{theorem}[blueprint="infogeometry:canonical:moe:phaselyapunovafter:eq_2"]
  \inputleannode{InfoGeometry.Canonical.MoE.phaseLyapunovAfter.eq_2}
\end{theorem}

\section{InfoGeometry.Canonical.MoE.phaseLyapunovBefore}

\subsection*{eq_1}

\begin{theorem}[blueprint="infogeometry:canonical:moe:phaselyapunovbefore:eq_1"]
  \inputleannode{InfoGeometry.Canonical.MoE.phaseLyapunovBefore.eq_1}
\end{theorem}

\subsection*{eq_2}

\begin{theorem}[blueprint="infogeometry:canonical:moe:phaselyapunovbefore:eq_2"]
  \inputleannode{InfoGeometry.Canonical.MoE.phaseLyapunovBefore.eq_2}
\end{theorem}

\section{InfoGeometry.Canonical.MoE.phaseRNBarrierAfter}

\subsection*{eq_1}

\begin{theorem}[blueprint="infogeometry:canonical:moe:phasernbarrierafter:eq_1"]
  \inputleannode{InfoGeometry.Canonical.MoE.phaseRNBarrierAfter.eq_1}
\end{theorem}

\subsection*{eq_2}

\begin{theorem}[blueprint="infogeometry:canonical:moe:phasernbarrierafter:eq_2"]
  \inputleannode{InfoGeometry.Canonical.MoE.phaseRNBarrierAfter.eq_2}
\end{theorem}

\section{InfoGeometry.Canonical.MoE.phaseRNBarrierBefore}

\subsection*{eq_1}

\begin{theorem}[blueprint="infogeometry:canonical:moe:phasernbarrierbefore:eq_1"]
  \inputleannode{InfoGeometry.Canonical.MoE.phaseRNBarrierBefore.eq_1}
\end{theorem}

\subsection*{eq_2}

\begin{theorem}[blueprint="infogeometry:canonical:moe:phasernbarrierbefore:eq_2"]
  \inputleannode{InfoGeometry.Canonical.MoE.phaseRNBarrierBefore.eq_2}
\end{theorem}

\section{InfoGeometry.Canonical.MoE.rightWeylScale}

\subsection*{eq_1}

\begin{theorem}[blueprint="infogeometry:canonical:moe:rightweylscale:eq_1"]
  \inputleannode{InfoGeometry.Canonical.MoE.rightWeylScale.eq_1}
\end{theorem}

\section{InfoGeometry.Canonical.MoE.rowNormalize}

\subsection*{congr_simp}

\begin{theorem}[blueprint="infogeometry:canonical:moe:rownormalize:congr_simp"]
  \inputleannode{InfoGeometry.Canonical.MoE.rowNormalize.congr_simp}
\end{theorem}

\subsection*{eq_1}

\begin{theorem}[blueprint="infogeometry:canonical:moe:rownormalize:eq_1"]
  \inputleannode{InfoGeometry.Canonical.MoE.rowNormalize.eq_1}
\end{theorem}

\section{InfoGeometry.Canonical.MoE.rowSum}

\subsection*{eq_1}

\begin{theorem}[blueprint="infogeometry:canonical:moe:rowsum:eq_1"]
  \inputleannode{InfoGeometry.Canonical.MoE.rowSum.eq_1}
\end{theorem}

\section{InfoGeometry.Canonical.MoE.sinkhornScaledCoupling}

\subsection*{eq_1}

\begin{theorem}[blueprint="infogeometry:canonical:moe:sinkhornscaledcoupling:eq_1"]
  \inputleannode{InfoGeometry.Canonical.MoE.sinkhornScaledCoupling.eq_1}
\end{theorem}

\section{InfoGeometry.Canonical.MoE.splitBasisMinus}

\subsection*{eq_1}

\begin{theorem}[blueprint="infogeometry:canonical:moe:splitbasisminus:eq_1"]
  \inputleannode{InfoGeometry.Canonical.MoE.splitBasisMinus.eq_1}
\end{theorem}

\section{InfoGeometry.Canonical.MoE.splitBasisPlus}

\subsection*{eq_1}

\begin{theorem}[blueprint="infogeometry:canonical:moe:splitbasisplus:eq_1"]
  \inputleannode{InfoGeometry.Canonical.MoE.splitBasisPlus.eq_1}
\end{theorem}

\section{InfoGeometry.Canonical.MoE.switchMatrix}

\subsection*{eq_1}

\begin{theorem}[blueprint="infogeometry:canonical:moe:switchmatrix:eq_1"]
  \inputleannode{InfoGeometry.Canonical.MoE.switchMatrix.eq_1}
\end{theorem}

\section{InfoGeometry.Canonical.MoorePenrose}

\noindent\textbf{Buildable}: yes

\subsection*{IsMoorePenroseInverse}

\begin{theorem}[blueprint="infogeometry:canonical:moorepenrose:ismoorepenroseinverse"]
  \inputleannode{InfoGeometry.Canonical.MoorePenrose.IsMoorePenroseInverse}
\end{theorem}

\subsection*{chiralAnomaly}

\begin{theorem}[blueprint="infogeometry:canonical:moorepenrose:chiralanomaly"]
  \inputleannode{InfoGeometry.Canonical.MoorePenrose.chiralAnomaly}
\end{theorem}

\subsection*{chiralScale}

\begin{theorem}[blueprint="infogeometry:canonical:moorepenrose:chiralscale"]
  \inputleannode{InfoGeometry.Canonical.MoorePenrose.chiralScale}
\end{theorem}

\subsection*{epsilon}

\begin{theorem}[blueprint="infogeometry:canonical:moorepenrose:epsilon"]
  \inputleannode{InfoGeometry.Canonical.MoorePenrose.epsilon}
\end{theorem}

\section{InfoGeometry.Canonical.MoorePenrose.IsMoorePenroseInverse}

\subsection*{ab_star}

\begin{theorem}[blueprint="infogeometry:canonical:moorepenrose:ismoorepenroseinverse:ab_star"]
  \inputleannode{InfoGeometry.Canonical.MoorePenrose.IsMoorePenroseInverse.ab_star}
\end{theorem}

\subsection*{aba_eq_a}

\begin{theorem}[blueprint="infogeometry:canonical:moorepenrose:ismoorepenroseinverse:aba_eq_a"]
  \inputleannode{InfoGeometry.Canonical.MoorePenrose.IsMoorePenroseInverse.aba_eq_a}
\end{theorem}

\subsection*{ba_star}

\begin{theorem}[blueprint="infogeometry:canonical:moorepenrose:ismoorepenroseinverse:ba_star"]
  \inputleannode{InfoGeometry.Canonical.MoorePenrose.IsMoorePenroseInverse.ba_star}
\end{theorem}

\subsection*{bab_eq_b}

\begin{theorem}[blueprint="infogeometry:canonical:moorepenrose:ismoorepenroseinverse:bab_eq_b"]
  \inputleannode{InfoGeometry.Canonical.MoorePenrose.IsMoorePenroseInverse.bab_eq_b}
\end{theorem}

\subsection*{casesOn}

\begin{theorem}[blueprint="infogeometry:canonical:moorepenrose:ismoorepenroseinverse:caseson"]
  \inputleannode{InfoGeometry.Canonical.MoorePenrose.IsMoorePenroseInverse.casesOn}
\end{theorem}

\subsection*{leftProjector}

\begin{theorem}[blueprint="infogeometry:canonical:moorepenrose:ismoorepenroseinverse:leftprojector"]
  \inputleannode{InfoGeometry.Canonical.MoorePenrose.IsMoorePenroseInverse.leftProjector}
\end{theorem}

\subsection*{leftProjector_idempotent}

\begin{theorem}[blueprint="infogeometry:canonical:moorepenrose:ismoorepenroseinverse:leftprojector_idempotent"]
  \inputleannode{InfoGeometry.Canonical.MoorePenrose.IsMoorePenroseInverse.leftProjector_idempotent}
\end{theorem}

\subsection*{leftProjector_star}

\begin{theorem}[blueprint="infogeometry:canonical:moorepenrose:ismoorepenroseinverse:leftprojector_star"]
  \inputleannode{InfoGeometry.Canonical.MoorePenrose.IsMoorePenroseInverse.leftProjector_star}
\end{theorem}

\subsection*{mk}

\begin{theorem}[blueprint="infogeometry:canonical:moorepenrose:ismoorepenroseinverse:mk"]
  \inputleannode{InfoGeometry.Canonical.MoorePenrose.IsMoorePenroseInverse.mk}
\end{theorem}

\subsection*{rec}

\begin{theorem}[blueprint="infogeometry:canonical:moorepenrose:ismoorepenroseinverse:rec"]
  \inputleannode{InfoGeometry.Canonical.MoorePenrose.IsMoorePenroseInverse.rec}
\end{theorem}

\subsection*{recOn}

\begin{theorem}[blueprint="infogeometry:canonical:moorepenrose:ismoorepenroseinverse:recon"]
  \inputleannode{InfoGeometry.Canonical.MoorePenrose.IsMoorePenroseInverse.recOn}
\end{theorem}

\subsection*{rightProjector}

\begin{theorem}[blueprint="infogeometry:canonical:moorepenrose:ismoorepenroseinverse:rightprojector"]
  \inputleannode{InfoGeometry.Canonical.MoorePenrose.IsMoorePenroseInverse.rightProjector}
\end{theorem}

\subsection*{rightProjector_idempotent}

\begin{theorem}[blueprint="infogeometry:canonical:moorepenrose:ismoorepenroseinverse:rightprojector_idempotent"]
  \inputleannode{InfoGeometry.Canonical.MoorePenrose.IsMoorePenroseInverse.rightProjector_idempotent}
\end{theorem}

\subsection*{rightProjector_star}

\begin{theorem}[blueprint="infogeometry:canonical:moorepenrose:ismoorepenroseinverse:rightprojector_star"]
  \inputleannode{InfoGeometry.Canonical.MoorePenrose.IsMoorePenroseInverse.rightProjector_star}
\end{theorem}

\section{InfoGeometry.Canonical.PathIntegral}

\noindent\textbf{Buildable}: yes

\subsection*{chiralPathIntegral}

\begin{theorem}[blueprint="infogeometry:canonical:pathintegral:chiralpathintegral"]
  \inputleannode{InfoGeometry.Canonical.PathIntegral.chiralPathIntegral}
\end{theorem}

\subsection*{chiralPathWeight}

\begin{theorem}[blueprint="infogeometry:canonical:pathintegral:chiralpathweight"]
  \inputleannode{InfoGeometry.Canonical.PathIntegral.chiralPathWeight}
\end{theorem}

\subsection*{expectedHolonomy}

\begin{theorem}[blueprint="infogeometry:canonical:pathintegral:expectedholonomy"]
  \inputleannode{InfoGeometry.Canonical.PathIntegral.expectedHolonomy}
\end{theorem}

\section{InfoGeometry.Canonical.PerelmanW}

\noindent\textbf{Buildable}: yes

\subsection*{SatisfiesWLaw}

\begin{theorem}[blueprint="infogeometry:canonical:perelmanw:satisfieswlaw"]
  \inputleannode{InfoGeometry.Canonical.PerelmanW.SatisfiesWLaw}
\end{theorem}

\subsection*{WDissipation}

\begin{theorem}[blueprint="infogeometry:canonical:perelmanw:wdissipation"]
  \inputleannode{InfoGeometry.Canonical.PerelmanW.WDissipation}
\end{theorem}

\subsection*{WFunctional}

\begin{theorem}[blueprint="infogeometry:canonical:perelmanw:wfunctional"]
  \inputleannode{InfoGeometry.Canonical.PerelmanW.WFunctional}
\end{theorem}

\subsection*{WFunctional_monotone_of_nonneg_dissipation}

\begin{theorem}[blueprint="infogeometry:canonical:perelmanw:wfunctional_monotone_of_nonneg_dissipation"]
  \inputleannode{InfoGeometry.Canonical.PerelmanW.WFunctional_monotone_of_nonneg_dissipation}
\end{theorem}

\subsection*{WFunctional_strictMono_of_pos_dissipation}

\begin{theorem}[blueprint="infogeometry:canonical:perelmanw:wfunctional_strictmono_of_pos_dissipation"]
  \inputleannode{InfoGeometry.Canonical.PerelmanW.WFunctional_strictMono_of_pos_dissipation}
\end{theorem}

\subsection*{W_constant_of_spinorial_zero}

\begin{theorem}[blueprint="infogeometry:canonical:perelmanw:w_constant_of_spinorial_zero"]
  \inputleannode{InfoGeometry.Canonical.PerelmanW.W_constant_of_spinorial_zero}
\end{theorem}

\subsection*{W_monotone_of_spinorial_normalized_tracking}

\begin{theorem}[blueprint="infogeometry:canonical:perelmanw:w_monotone_of_spinorial_normalized_tracking"]
  \inputleannode{InfoGeometry.Canonical.PerelmanW.W_monotone_of_spinorial_normalized_tracking}
\end{theorem}

\subsection*{W_strictMono_of_spinorial_nonzero}

\begin{theorem}[blueprint="infogeometry:canonical:perelmanw:w_strictmono_of_spinorial_nonzero"]
  \inputleannode{InfoGeometry.Canonical.PerelmanW.W_strictMono_of_spinorial_nonzero}
\end{theorem}

\subsection*{deriv_WFunctional_eq_of_law}

\begin{theorem}[blueprint="infogeometry:canonical:perelmanw:deriv_wfunctional_eq_of_law"]
  \inputleannode{InfoGeometry.Canonical.PerelmanW.deriv_WFunctional_eq_of_law}
\end{theorem}

\subsection*{deriv_W_eq_abs_spinorial_of_normalized_tracking}

\begin{theorem}[blueprint="infogeometry:canonical:perelmanw:deriv_w_eq_abs_spinorial_of_normalized_tracking"]
  \inputleannode{InfoGeometry.Canonical.PerelmanW.deriv_W_eq_abs_spinorial_of_normalized_tracking}
\end{theorem}

\subsection*{spinorialWDissipation}

\begin{theorem}[blueprint="infogeometry:canonical:perelmanw:spinorialwdissipation"]
  \inputleannode{InfoGeometry.Canonical.PerelmanW.spinorialWDissipation}
\end{theorem}

\subsection*{spinorialWDissipation_eq_abs_spinorial_of_normalized_tracking}

\begin{theorem}[blueprint="infogeometry:canonical:perelmanw:spinorialwdissipation_eq_abs_spinorial_of_normalized_tracking"]
  \inputleannode{InfoGeometry.Canonical.PerelmanW.spinorialWDissipation_eq_abs_spinorial_of_normalized_tracking}
\end{theorem}

\section{InfoGeometry.Canonical.PerelmanW.WDissipation}

\subsection*{eq_1}

\begin{theorem}[blueprint="infogeometry:canonical:perelmanw:wdissipation:eq_1"]
  \inputleannode{InfoGeometry.Canonical.PerelmanW.WDissipation.eq_1}
\end{theorem}

\section{InfoGeometry.Canonical.Prequantum}

\noindent\textbf{Buildable}: yes

\noindent\emph{No exported declarations in the current declaration graph.}

\section{InfoGeometry.Canonical.Projective}

\noindent\textbf{Buildable}: yes

\noindent\emph{No exported declarations in the current declaration graph.}

\section{InfoGeometry.Canonical.QFTTDFTLaunchpad}

\noindent\textbf{Buildable}: yes

\subsection*{HohenbergKohnDualState}

\begin{theorem}[blueprint="infogeometry:canonical:qfttdftlaunchpad:hohenbergkohndualstate"]
  \inputleannode{InfoGeometry.Canonical.QFTTDFTLaunchpad.HohenbergKohnDualState}
\end{theorem}

\subsection*{RungeGrossStationaryDualState}

\begin{theorem}[blueprint="infogeometry:canonical:qfttdftlaunchpad:rungegrossstationarydualstate"]
  \inputleannode{InfoGeometry.Canonical.QFTTDFTLaunchpad.RungeGrossStationaryDualState}
\end{theorem}

\subsection*{aqft_kmsClosure_of_sinkhornControl}

\begin{theorem}[blueprint="infogeometry:canonical:qfttdftlaunchpad:aqft_kmsclosure_of_sinkhorncontrol"]
  \inputleannode{InfoGeometry.Canonical.QFTTDFTLaunchpad.aqft_kmsClosure_of_sinkhornControl}
\end{theorem}

\subsection*{aqft_tdft_constructive_launchpad}

\begin{theorem}[blueprint="infogeometry:canonical:qfttdftlaunchpad:aqft_tdft_constructive_launchpad"]
  \inputleannode{InfoGeometry.Canonical.QFTTDFTLaunchpad.aqft_tdft_constructive_launchpad}
\end{theorem}

\subsection*{fenchelGap_zero_along_grad_of_legendreInvolution}

\begin{theorem}[blueprint="infogeometry:canonical:qfttdftlaunchpad:fenchelgap_zero_along_grad_of_legendreinvolution"]
  \inputleannode{InfoGeometry.Canonical.QFTTDFTLaunchpad.fenchelGap_zero_along_grad_of_legendreInvolution}
\end{theorem}

\subsection*{grandCanonicalFockEulerStep_eq_hamiltonianStep_of_vacuumTransported}

\begin{theorem}[blueprint="infogeometry:canonical:qfttdftlaunchpad:grandcanonicalfockeulerstep_eq_hamiltonianstep_of_vacuumtransported"]
  \inputleannode{InfoGeometry.Canonical.QFTTDFTLaunchpad.grandCanonicalFockEulerStep_eq_hamiltonianStep_of_vacuumTransported}
\end{theorem}

\subsection*{hohenbergKohnDualState_of_legendreInvolution}

\begin{theorem}[blueprint="infogeometry:canonical:qfttdftlaunchpad:hohenbergkohndualstate_of_legendreinvolution"]
  \inputleannode{InfoGeometry.Canonical.QFTTDFTLaunchpad.hohenbergKohnDualState_of_legendreInvolution}
\end{theorem}

\subsection*{rungeGrossStationaryDualState_of_fixedPoint}

\begin{theorem}[blueprint="infogeometry:canonical:qfttdftlaunchpad:rungegrossstationarydualstate_of_fixedpoint"]
  \inputleannode{InfoGeometry.Canonical.QFTTDFTLaunchpad.rungeGrossStationaryDualState_of_fixedPoint}
\end{theorem}

\section{InfoGeometry.Canonical.Quantum}

\noindent\textbf{Buildable}: yes

\noindent\emph{No exported declarations in the current declaration graph.}

\section{InfoGeometry.Canonical.QuantumInference}

\noindent\textbf{Buildable}: yes

\subsection*{CCR}

\begin{theorem}[blueprint="infogeometry:canonical:quantuminference:ccr"]
  \inputleannode{InfoGeometry.Canonical.QuantumInference.CCR}
\end{theorem}

\subsection*{DiracField}

\begin{theorem}[blueprint="infogeometry:canonical:quantuminference:diracfield"]
  \inputleannode{InfoGeometry.Canonical.QuantumInference.DiracField}
\end{theorem}

\subsection*{Op}

\begin{theorem}[blueprint="infogeometry:canonical:quantuminference:op"]
  \inputleannode{InfoGeometry.Canonical.QuantumInference.Op}
\end{theorem}

\subsection*{comm}

\begin{theorem}[blueprint="infogeometry:canonical:quantuminference:comm"]
  \inputleannode{InfoGeometry.Canonical.QuantumInference.comm}
\end{theorem}

\subsection*{continuousWilsonLoop}

\begin{theorem}[blueprint="infogeometry:canonical:quantuminference:continuouswilsonloop"]
  \inputleannode{InfoGeometry.Canonical.QuantumInference.continuousWilsonLoop}
\end{theorem}

\subsection*{diracReg}

\begin{theorem}[blueprint="infogeometry:canonical:quantuminference:diracreg"]
  \inputleannode{InfoGeometry.Canonical.QuantumInference.diracReg}
\end{theorem}

\subsection*{grandCanonicalZ}

\begin{theorem}[blueprint="infogeometry:canonical:quantuminference:grandcanonicalz"]
  \inputleannode{InfoGeometry.Canonical.QuantumInference.grandCanonicalZ}
\end{theorem}

\subsection*{totalNumber}

\begin{theorem}[blueprint="infogeometry:canonical:quantuminference:totalnumber"]
  \inputleannode{InfoGeometry.Canonical.QuantumInference.totalNumber}
\end{theorem}

\subsection*{wilsonLoopDiscrete}

\begin{theorem}[blueprint="infogeometry:canonical:quantuminference:wilsonloopdiscrete"]
  \inputleannode{InfoGeometry.Canonical.QuantumInference.wilsonLoopDiscrete}
\end{theorem}

\subsection*{wilsonPropagatorDiscrete}

\begin{theorem}[blueprint="infogeometry:canonical:quantuminference:wilsonpropagatordiscrete"]
  \inputleannode{InfoGeometry.Canonical.QuantumInference.wilsonPropagatorDiscrete}
\end{theorem}

\subsection*{wilsonStep}

\begin{theorem}[blueprint="infogeometry:canonical:quantuminference:wilsonstep"]
  \inputleannode{InfoGeometry.Canonical.QuantumInference.wilsonStep}
\end{theorem}

\section{InfoGeometry.Canonical.QuantumInference.CCR}

\subsection*{a}

\begin{theorem}[blueprint="infogeometry:canonical:quantuminference:ccr:a"]
  \inputleannode{InfoGeometry.Canonical.QuantumInference.CCR.a}
\end{theorem}

\subsection*{adag}

\begin{theorem}[blueprint="infogeometry:canonical:quantuminference:ccr:adag"]
  \inputleannode{InfoGeometry.Canonical.QuantumInference.CCR.adag}
\end{theorem}

\subsection*{casesOn}

\begin{theorem}[blueprint="infogeometry:canonical:quantuminference:ccr:caseson"]
  \inputleannode{InfoGeometry.Canonical.QuantumInference.CCR.casesOn}
\end{theorem}

\subsection*{comm_a_a}

\begin{theorem}[blueprint="infogeometry:canonical:quantuminference:ccr:comm_a_a"]
  \inputleannode{InfoGeometry.Canonical.QuantumInference.CCR.comm_a_a}
\end{theorem}

\subsection*{comm_a_adag}

\begin{theorem}[blueprint="infogeometry:canonical:quantuminference:ccr:comm_a_adag"]
  \inputleannode{InfoGeometry.Canonical.QuantumInference.CCR.comm_a_adag}
\end{theorem}

\subsection*{comm_adag_adag}

\begin{theorem}[blueprint="infogeometry:canonical:quantuminference:ccr:comm_adag_adag"]
  \inputleannode{InfoGeometry.Canonical.QuantumInference.CCR.comm_adag_adag}
\end{theorem}

\subsection*{ctorIdx}

\begin{theorem}[blueprint="infogeometry:canonical:quantuminference:ccr:ctoridx"]
  \inputleannode{InfoGeometry.Canonical.QuantumInference.CCR.ctorIdx}
\end{theorem}

\subsection*{hessianIndefiniteForm}

\begin{theorem}[blueprint="infogeometry:canonical:quantuminference:ccr:hessianindefiniteform"]
  \inputleannode{InfoGeometry.Canonical.QuantumInference.CCR.hessianIndefiniteForm}
\end{theorem}

\subsection*{mk}

\begin{theorem}[blueprint="infogeometry:canonical:quantuminference:ccr:mk"]
  \inputleannode{InfoGeometry.Canonical.QuantumInference.CCR.mk}
\end{theorem}

\subsection*{noConfusion}

\begin{theorem}[blueprint="infogeometry:canonical:quantuminference:ccr:noconfusion"]
  \inputleannode{InfoGeometry.Canonical.QuantumInference.CCR.noConfusion}
\end{theorem}

\subsection*{noConfusionType}

\begin{theorem}[blueprint="infogeometry:canonical:quantuminference:ccr:noconfusiontype"]
  \inputleannode{InfoGeometry.Canonical.QuantumInference.CCR.noConfusionType}
\end{theorem}

\subsection*{rec}

\begin{theorem}[blueprint="infogeometry:canonical:quantuminference:ccr:rec"]
  \inputleannode{InfoGeometry.Canonical.QuantumInference.CCR.rec}
\end{theorem}

\subsection*{recOn}

\begin{theorem}[blueprint="infogeometry:canonical:quantuminference:ccr:recon"]
  \inputleannode{InfoGeometry.Canonical.QuantumInference.CCR.recOn}
\end{theorem}

\section{InfoGeometry.Canonical.RGFlow}

\noindent\textbf{Buildable}: yes

\subsection*{InformationFlow}

\begin{theorem}[blueprint="infogeometry:canonical:rgflow:informationflow"]
  \inputleannode{InfoGeometry.Canonical.RGFlow.InformationFlow}
\end{theorem}

\subsection*{IsFixedPoint}

\begin{theorem}[blueprint="infogeometry:canonical:rgflow:isfixedpoint"]
  \inputleannode{InfoGeometry.Canonical.RGFlow.IsFixedPoint}
\end{theorem}

\subsection*{betaFunction}

\begin{theorem}[blueprint="infogeometry:canonical:rgflow:betafunction"]
  \inputleannode{InfoGeometry.Canonical.RGFlow.betaFunction}
\end{theorem}

\subsection*{dual_map_invariant_at_fixed_point}

\begin{theorem}[blueprint="infogeometry:canonical:rgflow:dual_map_invariant_at_fixed_point"]
  \inputleannode{InfoGeometry.Canonical.RGFlow.dual_map_invariant_at_fixed_point}
\end{theorem}

\section{InfoGeometry.Canonical.RGFlow.IsFixedPoint}

\subsection*{beta_zero}

\begin{theorem}[blueprint="infogeometry:canonical:rgflow:isfixedpoint:beta_zero"]
  \inputleannode{InfoGeometry.Canonical.RGFlow.IsFixedPoint.beta_zero}
\end{theorem}

\subsection*{casesOn}

\begin{theorem}[blueprint="infogeometry:canonical:rgflow:isfixedpoint:caseson"]
  \inputleannode{InfoGeometry.Canonical.RGFlow.IsFixedPoint.casesOn}
\end{theorem}

\subsection*{dual_zero}

\begin{theorem}[blueprint="infogeometry:canonical:rgflow:isfixedpoint:dual_zero"]
  \inputleannode{InfoGeometry.Canonical.RGFlow.IsFixedPoint.dual_zero}
\end{theorem}

\subsection*{mk}

\begin{theorem}[blueprint="infogeometry:canonical:rgflow:isfixedpoint:mk"]
  \inputleannode{InfoGeometry.Canonical.RGFlow.IsFixedPoint.mk}
\end{theorem}

\subsection*{rec}

\begin{theorem}[blueprint="infogeometry:canonical:rgflow:isfixedpoint:rec"]
  \inputleannode{InfoGeometry.Canonical.RGFlow.IsFixedPoint.rec}
\end{theorem}

\subsection*{recOn}

\begin{theorem}[blueprint="infogeometry:canonical:rgflow:isfixedpoint:recon"]
  \inputleannode{InfoGeometry.Canonical.RGFlow.IsFixedPoint.recOn}
\end{theorem}

\section{InfoGeometry.Canonical.RicciMongeAmpere}

\noindent\textbf{Buildable}: yes

\subsection*{CurvatureTwoFormCompatibleWithEinstein}

\begin{theorem}[blueprint="infogeometry:canonical:riccimongeampere:curvaturetwoformcompatiblewitheinstein"]
  \inputleannode{InfoGeometry.Canonical.RicciMongeAmpere.CurvatureTwoFormCompatibleWithEinstein}
\end{theorem}

\subsection*{EinsteinEquationAt}

\begin{theorem}[blueprint="infogeometry:canonical:riccimongeampere:einsteinequationat"]
  \inputleannode{InfoGeometry.Canonical.RicciMongeAmpere.EinsteinEquationAt}
\end{theorem}

\subsection*{IsEinsteinKaehlerAt}

\begin{theorem}[blueprint="infogeometry:canonical:riccimongeampere:iseinsteinkaehlerat"]
  \inputleannode{InfoGeometry.Canonical.RicciMongeAmpere.IsEinsteinKaehlerAt}
\end{theorem}

\subsection*{IsEinsteinKaehlerAtWith}

\begin{theorem}[blueprint="infogeometry:canonical:riccimongeampere:iseinsteinkaehleratwith"]
  \inputleannode{InfoGeometry.Canonical.RicciMongeAmpere.IsEinsteinKaehlerAtWith}
\end{theorem}

\subsection*{IsRicciFixedPoint}

\begin{theorem}[blueprint="infogeometry:canonical:riccimongeampere:isriccifixedpoint"]
  \inputleannode{InfoGeometry.Canonical.RicciMongeAmpere.IsRicciFixedPoint}
\end{theorem}

\subsection*{IsSkewCurvatureTwoForm}

\begin{theorem}[blueprint="infogeometry:canonical:riccimongeampere:isskewcurvaturetwoform"]
  \inputleannode{InfoGeometry.Canonical.RicciMongeAmpere.IsSkewCurvatureTwoForm}
\end{theorem}

\subsection*{MongeAmpereConsistentWithHeatKernel}

\begin{theorem}[blueprint="infogeometry:canonical:riccimongeampere:mongeampereconsistentwithheatkernel"]
  \inputleannode{InfoGeometry.Canonical.RicciMongeAmpere.MongeAmpereConsistentWithHeatKernel}
\end{theorem}

\subsection*{RicciData}

\begin{theorem}[blueprint="infogeometry:canonical:riccimongeampere:riccidata"]
  \inputleannode{InfoGeometry.Canonical.RicciMongeAmpere.RicciData}
\end{theorem}

\subsection*{RicciFlow}

\begin{theorem}[blueprint="infogeometry:canonical:riccimongeampere:ricciflow"]
  \inputleannode{InfoGeometry.Canonical.RicciMongeAmpere.RicciFlow}
\end{theorem}

\subsection*{RicciTensor}

\begin{theorem}[blueprint="infogeometry:canonical:riccimongeampere:riccitensor"]
  \inputleannode{InfoGeometry.Canonical.RicciMongeAmpere.RicciTensor}
\end{theorem}

\subsection*{SatisfiesKaehlerRicciEvolution}

\begin{theorem}[blueprint="infogeometry:canonical:riccimongeampere:satisfieskaehlerriccievolution"]
  \inputleannode{InfoGeometry.Canonical.RicciMongeAmpere.SatisfiesKaehlerRicciEvolution}
\end{theorem}

\subsection*{SatisfiesMongeAmpere}

\begin{theorem}[blueprint="infogeometry:canonical:riccimongeampere:satisfiesmongeampere"]
  \inputleannode{InfoGeometry.Canonical.RicciMongeAmpere.SatisfiesMongeAmpere}
\end{theorem}

\subsection*{SatisfiesMongeAmperePotential}

\begin{theorem}[blueprint="infogeometry:canonical:riccimongeampere:satisfiesmongeamperepotential"]
  \inputleannode{InfoGeometry.Canonical.RicciMongeAmpere.SatisfiesMongeAmperePotential}
\end{theorem}

\subsection*{SatisfiesNormalizedKaehlerRicciFlow}

\begin{theorem}[blueprint="infogeometry:canonical:riccimongeampere:satisfiesnormalizedkaehlerricciflow"]
  \inputleannode{InfoGeometry.Canonical.RicciMongeAmpere.SatisfiesNormalizedKaehlerRicciFlow}
\end{theorem}

\subsection*{ScalarRicciFlow}

\begin{theorem}[blueprint="infogeometry:canonical:riccimongeampere:scalarricciflow"]
  \inputleannode{InfoGeometry.Canonical.RicciMongeAmpere.ScalarRicciFlow}
\end{theorem}

\subsection*{SpinConnection}

\begin{theorem}[blueprint="infogeometry:canonical:riccimongeampere:spinconnection"]
  \inputleannode{InfoGeometry.Canonical.RicciMongeAmpere.SpinConnection}
\end{theorem}

\subsection*{SplitVielbein}

\begin{theorem}[blueprint="infogeometry:canonical:riccimongeampere:splitvielbein"]
  \inputleannode{InfoGeometry.Canonical.RicciMongeAmpere.SplitVielbein}
\end{theorem}

\subsection*{StressEnergyTensor}

\begin{theorem}[blueprint="infogeometry:canonical:riccimongeampere:stressenergytensor"]
  \inputleannode{InfoGeometry.Canonical.RicciMongeAmpere.StressEnergyTensor}
\end{theorem}

\subsection*{StrongRicciFromHessian}

\begin{theorem}[blueprint="infogeometry:canonical:riccimongeampere:strongriccifromhessian"]
  \inputleannode{InfoGeometry.Canonical.RicciMongeAmpere.StrongRicciFromHessian}
\end{theorem}

\subsection*{VacuumEinsteinEquationAt}

\begin{theorem}[blueprint="infogeometry:canonical:riccimongeampere:vacuumeinsteinequationat"]
  \inputleannode{InfoGeometry.Canonical.RicciMongeAmpere.VacuumEinsteinEquationAt}
\end{theorem}

\subsection*{VacuumEinsteinOnTransportedSplit}

\begin{theorem}[blueprint="infogeometry:canonical:riccimongeampere:vacuumeinsteinontransportedsplit"]
  \inputleannode{InfoGeometry.Canonical.RicciMongeAmpere.VacuumEinsteinOnTransportedSplit}
\end{theorem}

\subsection*{curvature_action_zero_of_vacuum_on_transportedSplit}

\begin{theorem}[blueprint="infogeometry:canonical:riccimongeampere:curvature_action_zero_of_vacuum_on_transportedsplit"]
  \inputleannode{InfoGeometry.Canonical.RicciMongeAmpere.curvature_action_zero_of_vacuum_on_transportedSplit}
\end{theorem}

\subsection*{einsteinEquationAt_of_vacuum}

\begin{theorem}[blueprint="infogeometry:canonical:riccimongeampere:einsteinequationat_of_vacuum"]
  \inputleannode{InfoGeometry.Canonical.RicciMongeAmpere.einsteinEquationAt_of_vacuum}
\end{theorem}

\subsection*{einsteinTensorAt}

\begin{theorem}[blueprint="infogeometry:canonical:riccimongeampere:einsteintensorat"]
  \inputleannode{InfoGeometry.Canonical.RicciMongeAmpere.einsteinTensorAt}
\end{theorem}

\subsection*{einsteinTensor_eq_metric_multiple_of_einsteinKaehlerWith}

\begin{theorem}[blueprint="infogeometry:canonical:riccimongeampere:einsteintensor_eq_metric_multiple_of_einsteinkaehlerwith"]
  \inputleannode{InfoGeometry.Canonical.RicciMongeAmpere.einsteinTensor_eq_metric_multiple_of_einsteinKaehlerWith}
\end{theorem}

\subsection*{einsteinTensor_transport_split_eq_metric_multiple}

\begin{theorem}[blueprint="infogeometry:canonical:riccimongeampere:einsteintensor_transport_split_eq_metric_multiple"]
  \inputleannode{InfoGeometry.Canonical.RicciMongeAmpere.einsteinTensor_transport_split_eq_metric_multiple}
\end{theorem}

\subsection*{einsteinTensor_transport_split_mixed_eq_zero_of_scalar_closure}

\begin{theorem}[blueprint="infogeometry:canonical:riccimongeampere:einsteintensor_transport_split_mixed_eq_zero_of_scalar_closure"]
  \inputleannode{InfoGeometry.Canonical.RicciMongeAmpere.einsteinTensor_transport_split_mixed_eq_zero_of_scalar_closure}
\end{theorem}

\subsection*{isEinsteinKaehlerAtWith_to_exists}

\begin{theorem}[blueprint="infogeometry:canonical:riccimongeampere:iseinsteinkaehleratwith_to_exists"]
  \inputleannode{InfoGeometry.Canonical.RicciMongeAmpere.isEinsteinKaehlerAtWith_to_exists}
\end{theorem}

\subsection*{kaehlerRicciEvolution_beta_eq}

\begin{theorem}[blueprint="infogeometry:canonical:riccimongeampere:kaehlerriccievolution_beta_eq"]
  \inputleannode{InfoGeometry.Canonical.RicciMongeAmpere.kaehlerRicciEvolution_beta_eq}
\end{theorem}

\subsection*{mongeAmpereConsistentWithHeatKernel}

\begin{theorem}[blueprint="infogeometry:canonical:riccimongeampere:mongeampereconsistentwithheatkernel"]
  \inputleannode{InfoGeometry.Canonical.RicciMongeAmpere.mongeAmpereConsistentWithHeatKernel}
\end{theorem}

\subsection*{mongeAmpereDensity}

\begin{theorem}[blueprint="infogeometry:canonical:riccimongeampere:mongeamperedensity"]
  \inputleannode{InfoGeometry.Canonical.RicciMongeAmpere.mongeAmpereDensity}
\end{theorem}

\subsection*{mongeAmpereDensity_pos}

\begin{theorem}[blueprint="infogeometry:canonical:riccimongeampere:mongeamperedensity_pos"]
  \inputleannode{InfoGeometry.Canonical.RicciMongeAmpere.mongeAmpereDensity_pos}
\end{theorem}

\subsection*{normalizedKaehlerRicci_beta_eq_neg}

\begin{theorem}[blueprint="infogeometry:canonical:riccimongeampere:normalizedkaehlerricci_beta_eq_neg"]
  \inputleannode{InfoGeometry.Canonical.RicciMongeAmpere.normalizedKaehlerRicci_beta_eq_neg}
\end{theorem}

\subsection*{normalizedKaehlerRicci_beta_eq_neg_spinorial}

\begin{theorem}[blueprint="infogeometry:canonical:riccimongeampere:normalizedkaehlerricci_beta_eq_neg_spinorial"]
  \inputleannode{InfoGeometry.Canonical.RicciMongeAmpere.normalizedKaehlerRicci_beta_eq_neg_spinorial}
\end{theorem}

\subsection*{normalizedKaehlerRicci_fixedpoint_eq_zero}

\begin{theorem}[blueprint="infogeometry:canonical:riccimongeampere:normalizedkaehlerricci_fixedpoint_eq_zero"]
  \inputleannode{InfoGeometry.Canonical.RicciMongeAmpere.normalizedKaehlerRicci_fixedpoint_eq_zero}
\end{theorem}

\subsection*{normalizedKaehlerRicci_zero_implies_fixedpoint}

\begin{theorem}[blueprint="infogeometry:canonical:riccimongeampere:normalizedkaehlerricci_zero_implies_fixedpoint"]
  \inputleannode{InfoGeometry.Canonical.RicciMongeAmpere.normalizedKaehlerRicci_zero_implies_fixedpoint}
\end{theorem}

\subsection*{ricciBetaFunction}

\begin{theorem}[blueprint="infogeometry:canonical:riccimongeampere:riccibetafunction"]
  \inputleannode{InfoGeometry.Canonical.RicciMongeAmpere.ricciBetaFunction}
\end{theorem}

\subsection*{ricciFromMetricOp}

\begin{theorem}[blueprint="infogeometry:canonical:riccimongeampere:riccifrommetricop"]
  \inputleannode{InfoGeometry.Canonical.RicciMongeAmpere.ricciFromMetricOp}
\end{theorem}

\subsection*{ricci_component_invariant_at_fixed_point}

\begin{theorem}[blueprint="infogeometry:canonical:riccimongeampere:ricci_component_invariant_at_fixed_point"]
  \inputleannode{InfoGeometry.Canonical.RicciMongeAmpere.ricci_component_invariant_at_fixed_point}
\end{theorem}

\subsection*{ricci_tensor_invariant_at_fixed_point}

\begin{theorem}[blueprint="infogeometry:canonical:riccimongeampere:ricci_tensor_invariant_at_fixed_point"]
  \inputleannode{InfoGeometry.Canonical.RicciMongeAmpere.ricci_tensor_invariant_at_fixed_point}
\end{theorem}

\subsection*{satisfiesKaehlerRicciEvolution_iff}

\begin{theorem}[blueprint="infogeometry:canonical:riccimongeampere:satisfieskaehlerriccievolution_iff"]
  \inputleannode{InfoGeometry.Canonical.RicciMongeAmpere.satisfiesKaehlerRicciEvolution_iff}
\end{theorem}

\subsection*{satisfiesMongeAmperePotential_iff}

\begin{theorem}[blueprint="infogeometry:canonical:riccimongeampere:satisfiesmongeamperepotential_iff"]
  \inputleannode{InfoGeometry.Canonical.RicciMongeAmpere.satisfiesMongeAmperePotential_iff}
\end{theorem}

\subsection*{satisfiesMongeAmpere_iff}

\begin{theorem}[blueprint="infogeometry:canonical:riccimongeampere:satisfiesmongeampere_iff"]
  \inputleannode{InfoGeometry.Canonical.RicciMongeAmpere.satisfiesMongeAmpere_iff}
\end{theorem}

\subsection*{scalarCurvatureOnFrame}

\begin{theorem}[blueprint="infogeometry:canonical:riccimongeampere:scalarcurvatureonframe"]
  \inputleannode{InfoGeometry.Canonical.RicciMongeAmpere.scalarCurvatureOnFrame}
\end{theorem}

\subsection*{scalarCurvatureOnFrame_eq_einstein_multiple}

\begin{theorem}[blueprint="infogeometry:canonical:riccimongeampere:scalarcurvatureonframe_eq_einstein_multiple"]
  \inputleannode{InfoGeometry.Canonical.RicciMongeAmpere.scalarCurvatureOnFrame_eq_einstein_multiple}
\end{theorem}

\subsection*{scalarRicciBetaFunction}

\begin{theorem}[blueprint="infogeometry:canonical:riccimongeampere:scalarriccibetafunction"]
  \inputleannode{InfoGeometry.Canonical.RicciMongeAmpere.scalarRicciBetaFunction}
\end{theorem}

\subsection*{scalarRicci_invariant_at_fixed_point}

\begin{theorem}[blueprint="infogeometry:canonical:riccimongeampere:scalarricci_invariant_at_fixed_point"]
  \inputleannode{InfoGeometry.Canonical.RicciMongeAmpere.scalarRicci_invariant_at_fixed_point}
\end{theorem}

\subsection*{spectralMongeAmpereDensity}

\begin{theorem}[blueprint="infogeometry:canonical:riccimongeampere:spectralmongeamperedensity"]
  \inputleannode{InfoGeometry.Canonical.RicciMongeAmpere.spectralMongeAmpereDensity}
\end{theorem}

\subsection*{spectralMongeAmpereDensity_eq_exp_spectralVolume}

\begin{theorem}[blueprint="infogeometry:canonical:riccimongeampere:spectralmongeamperedensity_eq_exp_spectralvolume"]
  \inputleannode{InfoGeometry.Canonical.RicciMongeAmpere.spectralMongeAmpereDensity_eq_exp_spectralVolume}
\end{theorem}

\subsection*{spinorialScalarCurvature}

\begin{theorem}[blueprint="infogeometry:canonical:riccimongeampere:spinorialscalarcurvature"]
  \inputleannode{InfoGeometry.Canonical.RicciMongeAmpere.spinorialScalarCurvature}
\end{theorem}

\subsection*{spinorialScalarCurvature_eq_totalScalarCurvature}

\begin{theorem}[blueprint="infogeometry:canonical:riccimongeampere:spinorialscalarcurvature_eq_totalscalarcurvature"]
  \inputleannode{InfoGeometry.Canonical.RicciMongeAmpere.spinorialScalarCurvature_eq_totalScalarCurvature}
\end{theorem}

\subsection*{spinorialScalarCurvature_eq_zero_of_normalized_fixedpoint}

\begin{theorem}[blueprint="infogeometry:canonical:riccimongeampere:spinorialscalarcurvature_eq_zero_of_normalized_fixedpoint"]
  \inputleannode{InfoGeometry.Canonical.RicciMongeAmpere.spinorialScalarCurvature_eq_zero_of_normalized_fixedpoint}
\end{theorem}

\subsection*{transportedEinsteinResidual}

\begin{theorem}[blueprint="infogeometry:canonical:riccimongeampere:transportedeinsteinresidual"]
  \inputleannode{InfoGeometry.Canonical.RicciMongeAmpere.transportedEinsteinResidual}
\end{theorem}

\subsection*{transportedSplitVielbein}

\begin{theorem}[blueprint="infogeometry:canonical:riccimongeampere:transportedsplitvielbein"]
  \inputleannode{InfoGeometry.Canonical.RicciMongeAmpere.transportedSplitVielbein}
\end{theorem}

\subsection*{vacuumEinsteinAt_implies_on_transportedSplit}

\begin{theorem}[blueprint="infogeometry:canonical:riccimongeampere:vacuumeinsteinat_implies_on_transportedsplit"]
  \inputleannode{InfoGeometry.Canonical.RicciMongeAmpere.vacuumEinsteinAt_implies_on_transportedSplit}
\end{theorem}

\subsection*{vacuumEinsteinEquation_of_scalar_relation}

\begin{theorem}[blueprint="infogeometry:canonical:riccimongeampere:vacuumeinsteinequation_of_scalar_relation"]
  \inputleannode{InfoGeometry.Canonical.RicciMongeAmpere.vacuumEinsteinEquation_of_scalar_relation}
\end{theorem}

\subsection*{vacuumEinsteinEquation_of_spinorial_scalar}

\begin{theorem}[blueprint="infogeometry:canonical:riccimongeampere:vacuumeinsteinequation_of_spinorial_scalar"]
  \inputleannode{InfoGeometry.Canonical.RicciMongeAmpere.vacuumEinsteinEquation_of_spinorial_scalar}
\end{theorem}

\subsection*{vacuumOnTransportedSplit_iff_curvature_action_zero}

\begin{theorem}[blueprint="infogeometry:canonical:riccimongeampere:vacuumontransportedsplit_iff_curvature_action_zero"]
  \inputleannode{InfoGeometry.Canonical.RicciMongeAmpere.vacuumOnTransportedSplit_iff_curvature_action_zero}
\end{theorem}

\subsection*{vacuum_on_transportedSplit_of_curvature_action_zero}

\begin{theorem}[blueprint="infogeometry:canonical:riccimongeampere:vacuum_on_transportedsplit_of_curvature_action_zero"]
  \inputleannode{InfoGeometry.Canonical.RicciMongeAmpere.vacuum_on_transportedSplit_of_curvature_action_zero}
\end{theorem}

\section{InfoGeometry.Canonical.RicciMongeAmpere.RicciData}

\subsection*{K}

\begin{theorem}[blueprint="infogeometry:canonical:riccimongeampere:riccidata:k"]
  \inputleannode{InfoGeometry.Canonical.RicciMongeAmpere.RicciData.K}
\end{theorem}

\subsection*{casesOn}

\begin{theorem}[blueprint="infogeometry:canonical:riccimongeampere:riccidata:caseson"]
  \inputleannode{InfoGeometry.Canonical.RicciMongeAmpere.RicciData.casesOn}
\end{theorem}

\subsection*{ctorIdx}

\begin{theorem}[blueprint="infogeometry:canonical:riccimongeampere:riccidata:ctoridx"]
  \inputleannode{InfoGeometry.Canonical.RicciMongeAmpere.RicciData.ctorIdx}
\end{theorem}

\subsection*{mk}

\begin{theorem}[blueprint="infogeometry:canonical:riccimongeampere:riccidata:mk"]
  \inputleannode{InfoGeometry.Canonical.RicciMongeAmpere.RicciData.mk}
\end{theorem}

\subsection*{noConfusion}

\begin{theorem}[blueprint="infogeometry:canonical:riccimongeampere:riccidata:noconfusion"]
  \inputleannode{InfoGeometry.Canonical.RicciMongeAmpere.RicciData.noConfusion}
\end{theorem}

\subsection*{noConfusionType}

\begin{theorem}[blueprint="infogeometry:canonical:riccimongeampere:riccidata:noconfusiontype"]
  \inputleannode{InfoGeometry.Canonical.RicciMongeAmpere.RicciData.noConfusionType}
\end{theorem}

\subsection*{rec}

\begin{theorem}[blueprint="infogeometry:canonical:riccimongeampere:riccidata:rec"]
  \inputleannode{InfoGeometry.Canonical.RicciMongeAmpere.RicciData.rec}
\end{theorem}

\subsection*{recOn}

\begin{theorem}[blueprint="infogeometry:canonical:riccimongeampere:riccidata:recon"]
  \inputleannode{InfoGeometry.Canonical.RicciMongeAmpere.RicciData.recOn}
\end{theorem}

\subsection*{ricci}

\begin{theorem}[blueprint="infogeometry:canonical:riccimongeampere:riccidata:ricci"]
  \inputleannode{InfoGeometry.Canonical.RicciMongeAmpere.RicciData.ricci}
\end{theorem}

\subsection*{symmetric}

\begin{theorem}[blueprint="infogeometry:canonical:riccimongeampere:riccidata:symmetric"]
  \inputleannode{InfoGeometry.Canonical.RicciMongeAmpere.RicciData.symmetric}
\end{theorem}

\section{InfoGeometry.Canonical.RicciMongeAmpere.SpinConnection}

\subsection*{casesOn}

\begin{theorem}[blueprint="infogeometry:canonical:riccimongeampere:spinconnection:caseson"]
  \inputleannode{InfoGeometry.Canonical.RicciMongeAmpere.SpinConnection.casesOn}
\end{theorem}

\subsection*{ctorIdx}

\begin{theorem}[blueprint="infogeometry:canonical:riccimongeampere:spinconnection:ctoridx"]
  \inputleannode{InfoGeometry.Canonical.RicciMongeAmpere.SpinConnection.ctorIdx}
\end{theorem}

\subsection*{mk}

\begin{theorem}[blueprint="infogeometry:canonical:riccimongeampere:spinconnection:mk"]
  \inputleannode{InfoGeometry.Canonical.RicciMongeAmpere.SpinConnection.mk}
\end{theorem}

\subsection*{noConfusion}

\begin{theorem}[blueprint="infogeometry:canonical:riccimongeampere:spinconnection:noconfusion"]
  \inputleannode{InfoGeometry.Canonical.RicciMongeAmpere.SpinConnection.noConfusion}
\end{theorem}

\subsection*{noConfusionType}

\begin{theorem}[blueprint="infogeometry:canonical:riccimongeampere:spinconnection:noconfusiontype"]
  \inputleannode{InfoGeometry.Canonical.RicciMongeAmpere.SpinConnection.noConfusionType}
\end{theorem}

\subsection*{preserves_minus}

\begin{theorem}[blueprint="infogeometry:canonical:riccimongeampere:spinconnection:preserves_minus"]
  \inputleannode{InfoGeometry.Canonical.RicciMongeAmpere.SpinConnection.preserves_minus}
\end{theorem}

\subsection*{preserves_orthogonal}

\begin{theorem}[blueprint="infogeometry:canonical:riccimongeampere:spinconnection:preserves_orthogonal"]
  \inputleannode{InfoGeometry.Canonical.RicciMongeAmpere.SpinConnection.preserves_orthogonal}
\end{theorem}

\subsection*{preserves_plus}

\begin{theorem}[blueprint="infogeometry:canonical:riccimongeampere:spinconnection:preserves_plus"]
  \inputleannode{InfoGeometry.Canonical.RicciMongeAmpere.SpinConnection.preserves_plus}
\end{theorem}

\subsection*{rec}

\begin{theorem}[blueprint="infogeometry:canonical:riccimongeampere:spinconnection:rec"]
  \inputleannode{InfoGeometry.Canonical.RicciMongeAmpere.SpinConnection.rec}
\end{theorem}

\subsection*{recOn}

\begin{theorem}[blueprint="infogeometry:canonical:riccimongeampere:spinconnection:recon"]
  \inputleannode{InfoGeometry.Canonical.RicciMongeAmpere.SpinConnection.recOn}
\end{theorem}

\subsection*{transport}

\begin{theorem}[blueprint="infogeometry:canonical:riccimongeampere:spinconnection:transport"]
  \inputleannode{InfoGeometry.Canonical.RicciMongeAmpere.SpinConnection.transport}
\end{theorem}

\section{InfoGeometry.Canonical.RicciMongeAmpere.SplitVielbein}

\subsection*{casesOn}

\begin{theorem}[blueprint="infogeometry:canonical:riccimongeampere:splitvielbein:caseson"]
  \inputleannode{InfoGeometry.Canonical.RicciMongeAmpere.SplitVielbein.casesOn}
\end{theorem}

\subsection*{ctorIdx}

\begin{theorem}[blueprint="infogeometry:canonical:riccimongeampere:splitvielbein:ctoridx"]
  \inputleannode{InfoGeometry.Canonical.RicciMongeAmpere.SplitVielbein.ctorIdx}
\end{theorem}

\subsection*{eMinus}

\begin{theorem}[blueprint="infogeometry:canonical:riccimongeampere:splitvielbein:eminus"]
  \inputleannode{InfoGeometry.Canonical.RicciMongeAmpere.SplitVielbein.eMinus}
\end{theorem}

\subsection*{ePlus}

\begin{theorem}[blueprint="infogeometry:canonical:riccimongeampere:splitvielbein:eplus"]
  \inputleannode{InfoGeometry.Canonical.RicciMongeAmpere.SplitVielbein.ePlus}
\end{theorem}

\subsection*{minus_norm}

\begin{theorem}[blueprint="infogeometry:canonical:riccimongeampere:splitvielbein:minus_norm"]
  \inputleannode{InfoGeometry.Canonical.RicciMongeAmpere.SplitVielbein.minus_norm}
\end{theorem}

\subsection*{mk}

\begin{theorem}[blueprint="infogeometry:canonical:riccimongeampere:splitvielbein:mk"]
  \inputleannode{InfoGeometry.Canonical.RicciMongeAmpere.SplitVielbein.mk}
\end{theorem}

\subsection*{noConfusion}

\begin{theorem}[blueprint="infogeometry:canonical:riccimongeampere:splitvielbein:noconfusion"]
  \inputleannode{InfoGeometry.Canonical.RicciMongeAmpere.SplitVielbein.noConfusion}
\end{theorem}

\subsection*{noConfusionType}

\begin{theorem}[blueprint="infogeometry:canonical:riccimongeampere:splitvielbein:noconfusiontype"]
  \inputleannode{InfoGeometry.Canonical.RicciMongeAmpere.SplitVielbein.noConfusionType}
\end{theorem}

\subsection*{orthogonal}

\begin{theorem}[blueprint="infogeometry:canonical:riccimongeampere:splitvielbein:orthogonal"]
  \inputleannode{InfoGeometry.Canonical.RicciMongeAmpere.SplitVielbein.orthogonal}
\end{theorem}

\subsection*{plus_norm}

\begin{theorem}[blueprint="infogeometry:canonical:riccimongeampere:splitvielbein:plus_norm"]
  \inputleannode{InfoGeometry.Canonical.RicciMongeAmpere.SplitVielbein.plus_norm}
\end{theorem}

\subsection*{rec}

\begin{theorem}[blueprint="infogeometry:canonical:riccimongeampere:splitvielbein:rec"]
  \inputleannode{InfoGeometry.Canonical.RicciMongeAmpere.SplitVielbein.rec}
\end{theorem}

\subsection*{recOn}

\begin{theorem}[blueprint="infogeometry:canonical:riccimongeampere:splitvielbein:recon"]
  \inputleannode{InfoGeometry.Canonical.RicciMongeAmpere.SplitVielbein.recOn}
\end{theorem}

\section{InfoGeometry.Canonical.RicciMongeAmpere.StrongRicciFromHessian}

\subsection*{H}

\begin{theorem}[blueprint="infogeometry:canonical:riccimongeampere:strongriccifromhessian:h"]
  \inputleannode{InfoGeometry.Canonical.RicciMongeAmpere.StrongRicciFromHessian.H}
\end{theorem}

\subsection*{casesOn}

\begin{theorem}[blueprint="infogeometry:canonical:riccimongeampere:strongriccifromhessian:caseson"]
  \inputleannode{InfoGeometry.Canonical.RicciMongeAmpere.StrongRicciFromHessian.casesOn}
\end{theorem}

\subsection*{ctorIdx}

\begin{theorem}[blueprint="infogeometry:canonical:riccimongeampere:strongriccifromhessian:ctoridx"]
  \inputleannode{InfoGeometry.Canonical.RicciMongeAmpere.StrongRicciFromHessian.ctorIdx}
\end{theorem}

\subsection*{isEinsteinKaehlerAt}

\begin{theorem}[blueprint="infogeometry:canonical:riccimongeampere:strongriccifromhessian:iseinsteinkaehlerat"]
  \inputleannode{InfoGeometry.Canonical.RicciMongeAmpere.StrongRicciFromHessian.isEinsteinKaehlerAt}
\end{theorem}

\subsection*{mk}

\begin{theorem}[blueprint="infogeometry:canonical:riccimongeampere:strongriccifromhessian:mk"]
  \inputleannode{InfoGeometry.Canonical.RicciMongeAmpere.StrongRicciFromHessian.mk}
\end{theorem}

\subsection*{noConfusion}

\begin{theorem}[blueprint="infogeometry:canonical:riccimongeampere:strongriccifromhessian:noconfusion"]
  \inputleannode{InfoGeometry.Canonical.RicciMongeAmpere.StrongRicciFromHessian.noConfusion}
\end{theorem}

\subsection*{noConfusionType}

\begin{theorem}[blueprint="infogeometry:canonical:riccimongeampere:strongriccifromhessian:noconfusiontype"]
  \inputleannode{InfoGeometry.Canonical.RicciMongeAmpere.StrongRicciFromHessian.noConfusionType}
\end{theorem}

\subsection*{nonneg}

\begin{theorem}[blueprint="infogeometry:canonical:riccimongeampere:strongriccifromhessian:nonneg"]
  \inputleannode{InfoGeometry.Canonical.RicciMongeAmpere.StrongRicciFromHessian.nonneg}
\end{theorem}

\subsection*{rec}

\begin{theorem}[blueprint="infogeometry:canonical:riccimongeampere:strongriccifromhessian:rec"]
  \inputleannode{InfoGeometry.Canonical.RicciMongeAmpere.StrongRicciFromHessian.rec}
\end{theorem}

\subsection*{recOn}

\begin{theorem}[blueprint="infogeometry:canonical:riccimongeampere:strongriccifromhessian:recon"]
  \inputleannode{InfoGeometry.Canonical.RicciMongeAmpere.StrongRicciFromHessian.recOn}
\end{theorem}

\subsection*{ricci}

\begin{theorem}[blueprint="infogeometry:canonical:riccimongeampere:strongriccifromhessian:ricci"]
  \inputleannode{InfoGeometry.Canonical.RicciMongeAmpere.StrongRicciFromHessian.ricci}
\end{theorem}

\subsection*{ricci_apply}

\begin{theorem}[blueprint="infogeometry:canonical:riccimongeampere:strongriccifromhessian:ricci_apply"]
  \inputleannode{InfoGeometry.Canonical.RicciMongeAmpere.StrongRicciFromHessian.ricci_apply}
\end{theorem}

\subsection*{ricci_nonneg_diag}

\begin{theorem}[blueprint="infogeometry:canonical:riccimongeampere:strongriccifromhessian:ricci_nonneg_diag"]
  \inputleannode{InfoGeometry.Canonical.RicciMongeAmpere.StrongRicciFromHessian.ricci_nonneg_diag}
\end{theorem}

\subsection*{ricci_symmetric}

\begin{theorem}[blueprint="infogeometry:canonical:riccimongeampere:strongriccifromhessian:ricci_symmetric"]
  \inputleannode{InfoGeometry.Canonical.RicciMongeAmpere.StrongRicciFromHessian.ricci_symmetric}
\end{theorem}

\subsection*{symmetric}

\begin{theorem}[blueprint="infogeometry:canonical:riccimongeampere:strongriccifromhessian:symmetric"]
  \inputleannode{InfoGeometry.Canonical.RicciMongeAmpere.StrongRicciFromHessian.symmetric}
\end{theorem}

\subsection*{x0}

\begin{theorem}[blueprint="infogeometry:canonical:riccimongeampere:strongriccifromhessian:x0"]
  \inputleannode{InfoGeometry.Canonical.RicciMongeAmpere.StrongRicciFromHessian.x0}
\end{theorem}

\section{InfoGeometry.Canonical.RicciMongeAmpere.StrongRicciFromHessian.ricci}

\subsection*{eq_1}

\begin{theorem}[blueprint="infogeometry:canonical:riccimongeampere:strongriccifromhessian:ricci:eq_1"]
  \inputleannode{InfoGeometry.Canonical.RicciMongeAmpere.StrongRicciFromHessian.ricci.eq_1}
\end{theorem}

\section{InfoGeometry.Canonical.RicciMongeAmpere.ricciBetaFunction}

\subsection*{eq_1}

\begin{theorem}[blueprint="infogeometry:canonical:riccimongeampere:riccibetafunction:eq_1"]
  \inputleannode{InfoGeometry.Canonical.RicciMongeAmpere.ricciBetaFunction.eq_1}
\end{theorem}

\section{InfoGeometry.Canonical.RicciMongeAmpere.ricciFromMetricOp}

\subsection*{eq_1}

\begin{theorem}[blueprint="infogeometry:canonical:riccimongeampere:riccifrommetricop:eq_1"]
  \inputleannode{InfoGeometry.Canonical.RicciMongeAmpere.ricciFromMetricOp.eq_1}
\end{theorem}

\section{InfoGeometry.Canonical.Rosetta}

\noindent\textbf{Buildable}: yes

\noindent\emph{No exported declarations in the current declaration graph.}

\section{InfoGeometry.Canonical.SpectralInference}

\noindent\textbf{Buildable}: yes

\subsection*{ChiralSpectralTriple}

\begin{theorem}[blueprint="infogeometry:canonical:spectralinference:chiralspectraltriple"]
  \inputleannode{InfoGeometry.Canonical.SpectralInference.ChiralSpectralTriple}
\end{theorem}

\subsection*{InfoSpectralTriple}

\begin{theorem}[blueprint="infogeometry:canonical:spectralinference:infospectraltriple"]
  \inputleannode{InfoGeometry.Canonical.SpectralInference.InfoSpectralTriple}
\end{theorem}

\subsection*{RegularizedSpectralTriple}

\begin{theorem}[blueprint="infogeometry:canonical:spectralinference:regularizedspectraltriple"]
  \inputleannode{InfoGeometry.Canonical.SpectralInference.RegularizedSpectralTriple}
\end{theorem}

\subsection*{SpectralTriple}

\begin{theorem}[blueprint="infogeometry:canonical:spectralinference:spectraltriple"]
  \inputleannode{InfoGeometry.Canonical.SpectralInference.SpectralTriple}
\end{theorem}

\subsection*{bayesianAction}

\begin{theorem}[blueprint="infogeometry:canonical:spectralinference:bayesianaction"]
  \inputleannode{InfoGeometry.Canonical.SpectralInference.bayesianAction}
\end{theorem}

\subsection*{bayesianAction_succ}

\begin{theorem}[blueprint="infogeometry:canonical:spectralinference:bayesianaction_succ"]
  \inputleannode{InfoGeometry.Canonical.SpectralInference.bayesianAction_succ}
\end{theorem}

\subsection*{bayesianAction_zero}

\begin{theorem}[blueprint="infogeometry:canonical:spectralinference:bayesianaction_zero"]
  \inputleannode{InfoGeometry.Canonical.SpectralInference.bayesianAction_zero}
\end{theorem}

\section{InfoGeometry.Canonical.SpectralInference.ChiralSpectralTriple}

\subsection*{DD}

\begin{theorem}[blueprint="infogeometry:canonical:spectralinference:chiralspectraltriple:dd"]
  \inputleannode{InfoGeometry.Canonical.SpectralInference.ChiralSpectralTriple.DD}
\end{theorem}

\subsection*{DP}

\begin{theorem}[blueprint="infogeometry:canonical:spectralinference:chiralspectraltriple:dp"]
  \inputleannode{InfoGeometry.Canonical.SpectralInference.ChiralSpectralTriple.DP}
\end{theorem}

\subsection*{casesOn}

\begin{theorem}[blueprint="infogeometry:canonical:spectralinference:chiralspectraltriple:caseson"]
  \inputleannode{InfoGeometry.Canonical.SpectralInference.ChiralSpectralTriple.casesOn}
\end{theorem}

\subsection*{chiralSpectralAction}

\begin{theorem}[blueprint="infogeometry:canonical:spectralinference:chiralspectraltriple:chiralspectralaction"]
  \inputleannode{InfoGeometry.Canonical.SpectralInference.ChiralSpectralTriple.chiralSpectralAction}
\end{theorem}

\subsection*{ctorIdx}

\begin{theorem}[blueprint="infogeometry:canonical:spectralinference:chiralspectraltriple:ctoridx"]
  \inputleannode{InfoGeometry.Canonical.SpectralInference.ChiralSpectralTriple.ctorIdx}
\end{theorem}

\subsection*{epsilon}

\begin{theorem}[blueprint="infogeometry:canonical:spectralinference:chiralspectraltriple:epsilon"]
  \inputleannode{InfoGeometry.Canonical.SpectralInference.ChiralSpectralTriple.epsilon}
\end{theorem}

\subsection*{is_drazin}

\begin{theorem}[blueprint="infogeometry:canonical:spectralinference:chiralspectraltriple:is_drazin"]
  \inputleannode{InfoGeometry.Canonical.SpectralInference.ChiralSpectralTriple.is_drazin}
\end{theorem}

\subsection*{is_penrose}

\begin{theorem}[blueprint="infogeometry:canonical:spectralinference:chiralspectraltriple:is_penrose"]
  \inputleannode{InfoGeometry.Canonical.SpectralInference.ChiralSpectralTriple.is_penrose}
\end{theorem}

\subsection*{mk}

\begin{theorem}[blueprint="infogeometry:canonical:spectralinference:chiralspectraltriple:mk"]
  \inputleannode{InfoGeometry.Canonical.SpectralInference.ChiralSpectralTriple.mk}
\end{theorem}

\subsection*{noConfusion}

\begin{theorem}[blueprint="infogeometry:canonical:spectralinference:chiralspectraltriple:noconfusion"]
  \inputleannode{InfoGeometry.Canonical.SpectralInference.ChiralSpectralTriple.noConfusion}
\end{theorem}

\subsection*{noConfusionType}

\begin{theorem}[blueprint="infogeometry:canonical:spectralinference:chiralspectraltriple:noconfusiontype"]
  \inputleannode{InfoGeometry.Canonical.SpectralInference.ChiralSpectralTriple.noConfusionType}
\end{theorem}

\subsection*{rec}

\begin{theorem}[blueprint="infogeometry:canonical:spectralinference:chiralspectraltriple:rec"]
  \inputleannode{InfoGeometry.Canonical.SpectralInference.ChiralSpectralTriple.rec}
\end{theorem}

\subsection*{recOn}

\begin{theorem}[blueprint="infogeometry:canonical:spectralinference:chiralspectraltriple:recon"]
  \inputleannode{InfoGeometry.Canonical.SpectralInference.ChiralSpectralTriple.recOn}
\end{theorem}

\subsection*{toSpectralTriple}

\begin{theorem}[blueprint="infogeometry:canonical:spectralinference:chiralspectraltriple:tospectraltriple"]
  \inputleannode{InfoGeometry.Canonical.SpectralInference.ChiralSpectralTriple.toSpectralTriple}
\end{theorem}

\section{InfoGeometry.Canonical.SpectralInference.ChiralSpectralTriple.epsilon}

\subsection*{eq_1}

\begin{theorem}[blueprint="infogeometry:canonical:spectralinference:chiralspectraltriple:epsilon:eq_1"]
  \inputleannode{InfoGeometry.Canonical.SpectralInference.ChiralSpectralTriple.epsilon.eq_1}
\end{theorem}

\section{InfoGeometry.Canonical.SpectralInference.InfoSpectralTriple}

\subsection*{H}

\begin{theorem}[blueprint="infogeometry:canonical:spectralinference:infospectraltriple:h"]
  \inputleannode{InfoGeometry.Canonical.SpectralInference.InfoSpectralTriple.H}
\end{theorem}

\subsection*{casesOn}

\begin{theorem}[blueprint="infogeometry:canonical:spectralinference:infospectraltriple:caseson"]
  \inputleannode{InfoGeometry.Canonical.SpectralInference.InfoSpectralTriple.casesOn}
\end{theorem}

\subsection*{ctorIdx}

\begin{theorem}[blueprint="infogeometry:canonical:spectralinference:infospectraltriple:ctoridx"]
  \inputleannode{InfoGeometry.Canonical.SpectralInference.InfoSpectralTriple.ctorIdx}
\end{theorem}

\subsection*{dirac_sq_eq_metric}

\begin{theorem}[blueprint="infogeometry:canonical:spectralinference:infospectraltriple:dirac_sq_eq_metric"]
  \inputleannode{InfoGeometry.Canonical.SpectralInference.InfoSpectralTriple.dirac_sq_eq_metric}
\end{theorem}

\subsection*{inner_dirac_sq}

\begin{theorem}[blueprint="infogeometry:canonical:spectralinference:infospectraltriple:inner_dirac_sq"]
  \inputleannode{InfoGeometry.Canonical.SpectralInference.InfoSpectralTriple.inner_dirac_sq}
\end{theorem}

\subsection*{mk}

\begin{theorem}[blueprint="infogeometry:canonical:spectralinference:infospectraltriple:mk"]
  \inputleannode{InfoGeometry.Canonical.SpectralInference.InfoSpectralTriple.mk}
\end{theorem}

\subsection*{noConfusion}

\begin{theorem}[blueprint="infogeometry:canonical:spectralinference:infospectraltriple:noconfusion"]
  \inputleannode{InfoGeometry.Canonical.SpectralInference.InfoSpectralTriple.noConfusion}
\end{theorem}

\subsection*{noConfusionType}

\begin{theorem}[blueprint="infogeometry:canonical:spectralinference:infospectraltriple:noconfusiontype"]
  \inputleannode{InfoGeometry.Canonical.SpectralInference.InfoSpectralTriple.noConfusionType}
\end{theorem}

\subsection*{rec}

\begin{theorem}[blueprint="infogeometry:canonical:spectralinference:infospectraltriple:rec"]
  \inputleannode{InfoGeometry.Canonical.SpectralInference.InfoSpectralTriple.rec}
\end{theorem}

\subsection*{recOn}

\begin{theorem}[blueprint="infogeometry:canonical:spectralinference:infospectraltriple:recon"]
  \inputleannode{InfoGeometry.Canonical.SpectralInference.InfoSpectralTriple.recOn}
\end{theorem}

\subsection*{toSpectralTriple}

\begin{theorem}[blueprint="infogeometry:canonical:spectralinference:infospectraltriple:tospectraltriple"]
  \inputleannode{InfoGeometry.Canonical.SpectralInference.InfoSpectralTriple.toSpectralTriple}
\end{theorem}

\subsection*{x₀}

\begin{theorem}[blueprint="infogeometry:canonical:spectralinference:infospectraltriple:x₀"]
  \inputleannode{InfoGeometry.Canonical.SpectralInference.InfoSpectralTriple.x₀}
\end{theorem}

\section{InfoGeometry.Canonical.SpectralInference.RegularizedSpectralTriple}

\subsection*{DD}

\begin{theorem}[blueprint="infogeometry:canonical:spectralinference:regularizedspectraltriple:dd"]
  \inputleannode{InfoGeometry.Canonical.SpectralInference.RegularizedSpectralTriple.DD}
\end{theorem}

\subsection*{casesOn}

\begin{theorem}[blueprint="infogeometry:canonical:spectralinference:regularizedspectraltriple:caseson"]
  \inputleannode{InfoGeometry.Canonical.SpectralInference.RegularizedSpectralTriple.casesOn}
\end{theorem}

\subsection*{ctorIdx}

\begin{theorem}[blueprint="infogeometry:canonical:spectralinference:regularizedspectraltriple:ctoridx"]
  \inputleannode{InfoGeometry.Canonical.SpectralInference.RegularizedSpectralTriple.ctorIdx}
\end{theorem}

\subsection*{index}

\begin{theorem}[blueprint="infogeometry:canonical:spectralinference:regularizedspectraltriple:index"]
  \inputleannode{InfoGeometry.Canonical.SpectralInference.RegularizedSpectralTriple.index}
\end{theorem}

\subsection*{is_drazin}

\begin{theorem}[blueprint="infogeometry:canonical:spectralinference:regularizedspectraltriple:is_drazin"]
  \inputleannode{InfoGeometry.Canonical.SpectralInference.RegularizedSpectralTriple.is_drazin}
\end{theorem}

\subsection*{mk}

\begin{theorem}[blueprint="infogeometry:canonical:spectralinference:regularizedspectraltriple:mk"]
  \inputleannode{InfoGeometry.Canonical.SpectralInference.RegularizedSpectralTriple.mk}
\end{theorem}

\subsection*{noConfusion}

\begin{theorem}[blueprint="infogeometry:canonical:spectralinference:regularizedspectraltriple:noconfusion"]
  \inputleannode{InfoGeometry.Canonical.SpectralInference.RegularizedSpectralTriple.noConfusion}
\end{theorem}

\subsection*{noConfusionType}

\begin{theorem}[blueprint="infogeometry:canonical:spectralinference:regularizedspectraltriple:noconfusiontype"]
  \inputleannode{InfoGeometry.Canonical.SpectralInference.RegularizedSpectralTriple.noConfusionType}
\end{theorem}

\subsection*{rec}

\begin{theorem}[blueprint="infogeometry:canonical:spectralinference:regularizedspectraltriple:rec"]
  \inputleannode{InfoGeometry.Canonical.SpectralInference.RegularizedSpectralTriple.rec}
\end{theorem}

\subsection*{recOn}

\begin{theorem}[blueprint="infogeometry:canonical:spectralinference:regularizedspectraltriple:recon"]
  \inputleannode{InfoGeometry.Canonical.SpectralInference.RegularizedSpectralTriple.recOn}
\end{theorem}

\subsection*{spectralAction}

\begin{theorem}[blueprint="infogeometry:canonical:spectralinference:regularizedspectraltriple:spectralaction"]
  \inputleannode{InfoGeometry.Canonical.SpectralInference.RegularizedSpectralTriple.spectralAction}
\end{theorem}

\subsection*{toSpectralTriple}

\begin{theorem}[blueprint="infogeometry:canonical:spectralinference:regularizedspectraltriple:tospectraltriple"]
  \inputleannode{InfoGeometry.Canonical.SpectralInference.RegularizedSpectralTriple.toSpectralTriple}
\end{theorem}

\section{InfoGeometry.Canonical.SpectralInference.SpectralTriple}

\subsection*{D}

\begin{theorem}[blueprint="infogeometry:canonical:spectralinference:spectraltriple:d"]
  \inputleannode{InfoGeometry.Canonical.SpectralInference.SpectralTriple.D}
\end{theorem}

\subsection*{casesOn}

\begin{theorem}[blueprint="infogeometry:canonical:spectralinference:spectraltriple:caseson"]
  \inputleannode{InfoGeometry.Canonical.SpectralInference.SpectralTriple.casesOn}
\end{theorem}

\subsection*{commutator}

\begin{theorem}[blueprint="infogeometry:canonical:spectralinference:spectraltriple:commutator"]
  \inputleannode{InfoGeometry.Canonical.SpectralInference.SpectralTriple.commutator}
\end{theorem}

\subsection*{ctorIdx}

\begin{theorem}[blueprint="infogeometry:canonical:spectralinference:spectraltriple:ctoridx"]
  \inputleannode{InfoGeometry.Canonical.SpectralInference.SpectralTriple.ctorIdx}
\end{theorem}

\subsection*{gaugeField}

\begin{theorem}[blueprint="infogeometry:canonical:spectralinference:spectraltriple:gaugefield"]
  \inputleannode{InfoGeometry.Canonical.SpectralInference.SpectralTriple.gaugeField}
\end{theorem}

\subsection*{is_self_adjoint}

\begin{theorem}[blueprint="infogeometry:canonical:spectralinference:spectraltriple:is_self_adjoint"]
  \inputleannode{InfoGeometry.Canonical.SpectralInference.SpectralTriple.is_self_adjoint}
\end{theorem}

\subsection*{mk}

\begin{theorem}[blueprint="infogeometry:canonical:spectralinference:spectraltriple:mk"]
  \inputleannode{InfoGeometry.Canonical.SpectralInference.SpectralTriple.mk}
\end{theorem}

\subsection*{noConfusion}

\begin{theorem}[blueprint="infogeometry:canonical:spectralinference:spectraltriple:noconfusion"]
  \inputleannode{InfoGeometry.Canonical.SpectralInference.SpectralTriple.noConfusion}
\end{theorem}

\subsection*{noConfusionType}

\begin{theorem}[blueprint="infogeometry:canonical:spectralinference:spectraltriple:noconfusiontype"]
  \inputleannode{InfoGeometry.Canonical.SpectralInference.SpectralTriple.noConfusionType}
\end{theorem}

\subsection*{perturbedDirac}

\begin{theorem}[blueprint="infogeometry:canonical:spectralinference:spectraltriple:perturbeddirac"]
  \inputleannode{InfoGeometry.Canonical.SpectralInference.SpectralTriple.perturbedDirac}
\end{theorem}

\subsection*{rec}

\begin{theorem}[blueprint="infogeometry:canonical:spectralinference:spectraltriple:rec"]
  \inputleannode{InfoGeometry.Canonical.SpectralInference.SpectralTriple.rec}
\end{theorem}

\subsection*{recOn}

\begin{theorem}[blueprint="infogeometry:canonical:spectralinference:spectraltriple:recon"]
  \inputleannode{InfoGeometry.Canonical.SpectralInference.SpectralTriple.recOn}
\end{theorem}

\section{InfoGeometry.Canonical.SpectralInference.bayesianAction}

\subsection*{eq_1}

\begin{theorem}[blueprint="infogeometry:canonical:spectralinference:bayesianaction:eq_1"]
  \inputleannode{InfoGeometry.Canonical.SpectralInference.bayesianAction.eq_1}
\end{theorem}

\section{InfoGeometry.Canonical.Statistics}

\noindent\textbf{Buildable}: yes

\noindent\emph{No exported declarations in the current declaration graph.}

\section{InfoGeometry.Canonical.StatisticsLegacy}

\noindent\textbf{Buildable}: yes

\noindent\emph{No exported declarations in the current declaration graph.}

\section{InfoGeometry.Canonical.SuperInference}

\noindent\textbf{Buildable}: yes

\subsection*{SuperBeliefState}

\begin{theorem}[blueprint="infogeometry:canonical:superinference:superbeliefstate"]
  \inputleannode{InfoGeometry.Canonical.SuperInference.SuperBeliefState}
\end{theorem}

\subsection*{SuperState}

\begin{theorem}[blueprint="infogeometry:canonical:superinference:superstate"]
  \inputleannode{InfoGeometry.Canonical.SuperInference.SuperState}
\end{theorem}

\subsection*{superCharge}

\begin{theorem}[blueprint="infogeometry:canonical:superinference:supercharge"]
  \inputleannode{InfoGeometry.Canonical.SuperInference.superCharge}
\end{theorem}

\subsection*{superCharge_boson_eq_zero}

\begin{theorem}[blueprint="infogeometry:canonical:superinference:supercharge_boson_eq_zero"]
  \inputleannode{InfoGeometry.Canonical.SuperInference.superCharge_boson_eq_zero}
\end{theorem}

\subsection*{superCharge_fermion_eq_dualMap}

\begin{theorem}[blueprint="infogeometry:canonical:superinference:supercharge_fermion_eq_dualmap"]
  \inputleannode{InfoGeometry.Canonical.SuperInference.superCharge_fermion_eq_dualMap}
\end{theorem}

\subsection*{superHamiltonian}

\begin{theorem}[blueprint="infogeometry:canonical:superinference:superhamiltonian"]
  \inputleannode{InfoGeometry.Canonical.SuperInference.superHamiltonian}
\end{theorem}

\subsection*{susyCharge}

\begin{theorem}[blueprint="infogeometry:canonical:superinference:susycharge"]
  \inputleannode{InfoGeometry.Canonical.SuperInference.susyCharge}
\end{theorem}

\subsection*{susyHamiltonian}

\begin{theorem}[blueprint="infogeometry:canonical:superinference:susyhamiltonian"]
  \inputleannode{InfoGeometry.Canonical.SuperInference.susyHamiltonian}
\end{theorem}

\subsection*{susyHamiltonian_eq_self}

\begin{theorem}[blueprint="infogeometry:canonical:superinference:susyhamiltonian_eq_self"]
  \inputleannode{InfoGeometry.Canonical.SuperInference.susyHamiltonian_eq_self}
\end{theorem}

\section{InfoGeometry.Canonical.SuperInference.SuperState}

\subsection*{boson}

\begin{theorem}[blueprint="infogeometry:canonical:superinference:superstate:boson"]
  \inputleannode{InfoGeometry.Canonical.SuperInference.SuperState.boson}
\end{theorem}

\subsection*{casesOn}

\begin{theorem}[blueprint="infogeometry:canonical:superinference:superstate:caseson"]
  \inputleannode{InfoGeometry.Canonical.SuperInference.SuperState.casesOn}
\end{theorem}

\subsection*{ctorIdx}

\begin{theorem}[blueprint="infogeometry:canonical:superinference:superstate:ctoridx"]
  \inputleannode{InfoGeometry.Canonical.SuperInference.SuperState.ctorIdx}
\end{theorem}

\subsection*{fermion}

\begin{theorem}[blueprint="infogeometry:canonical:superinference:superstate:fermion"]
  \inputleannode{InfoGeometry.Canonical.SuperInference.SuperState.fermion}
\end{theorem}

\subsection*{mk}

\begin{theorem}[blueprint="infogeometry:canonical:superinference:superstate:mk"]
  \inputleannode{InfoGeometry.Canonical.SuperInference.SuperState.mk}
\end{theorem}

\subsection*{noConfusion}

\begin{theorem}[blueprint="infogeometry:canonical:superinference:superstate:noconfusion"]
  \inputleannode{InfoGeometry.Canonical.SuperInference.SuperState.noConfusion}
\end{theorem}

\subsection*{noConfusionType}

\begin{theorem}[blueprint="infogeometry:canonical:superinference:superstate:noconfusiontype"]
  \inputleannode{InfoGeometry.Canonical.SuperInference.SuperState.noConfusionType}
\end{theorem}

\subsection*{rec}

\begin{theorem}[blueprint="infogeometry:canonical:superinference:superstate:rec"]
  \inputleannode{InfoGeometry.Canonical.SuperInference.SuperState.rec}
\end{theorem}

\subsection*{recOn}

\begin{theorem}[blueprint="infogeometry:canonical:superinference:superstate:recon"]
  \inputleannode{InfoGeometry.Canonical.SuperInference.SuperState.recOn}
\end{theorem}

\section{InfoGeometry.Canonical.SuperUnified}

\noindent\textbf{Buildable}: yes

\noindent\emph{No exported declarations in the current declaration graph.}

\section{InfoGeometry.Canonical.ThermalLegacy}

\noindent\textbf{Buildable}: yes

\noindent\emph{No exported declarations in the current declaration graph.}

\section{InfoGeometry.Canonical.Thermo}

\noindent\textbf{Buildable}: yes

\noindent\emph{No exported declarations in the current declaration graph.}

\section{InfoGeometry.Canonical.TomitaTakesaki}

\noindent\textbf{Buildable}: yes

\subsection*{RoutingSplitCliffordAlg}

\begin{theorem}[blueprint="infogeometry:canonical:tomitatakesaki:routingsplitcliffordalg"]
  \inputleannode{InfoGeometry.Canonical.TomitaTakesaki.RoutingSplitCliffordAlg}
\end{theorem}

\subsection*{cptAtomSet}

\begin{theorem}[blueprint="infogeometry:canonical:tomitatakesaki:cptatomset"]
  \inputleannode{InfoGeometry.Canonical.TomitaTakesaki.cptAtomSet}
\end{theorem}

\subsection*{cptAtoms_generate_splitCliffordAlg}

\begin{theorem}[blueprint="infogeometry:canonical:tomitatakesaki:cptatoms_generate_splitcliffordalg"]
  \inputleannode{InfoGeometry.Canonical.TomitaTakesaki.cptAtoms_generate_splitCliffordAlg}
\end{theorem}

\subsection*{cptEps}

\begin{theorem}[blueprint="infogeometry:canonical:tomitatakesaki:cpteps"]
  \inputleannode{InfoGeometry.Canonical.TomitaTakesaki.cptEps}
\end{theorem}

\subsection*{cptJ}

\begin{theorem}[blueprint="infogeometry:canonical:tomitatakesaki:cptj"]
  \inputleannode{InfoGeometry.Canonical.TomitaTakesaki.cptJ}
\end{theorem}

\subsection*{cptJ_cptJeps_anticommute}

\begin{theorem}[blueprint="infogeometry:canonical:tomitatakesaki:cptj_cptjeps_anticommute"]
  \inputleannode{InfoGeometry.Canonical.TomitaTakesaki.cptJ_cptJeps_anticommute}
\end{theorem}

\subsection*{cptJ_sq}

\begin{theorem}[blueprint="infogeometry:canonical:tomitatakesaki:cptj_sq"]
  \inputleannode{InfoGeometry.Canonical.TomitaTakesaki.cptJ_sq}
\end{theorem}

\subsection*{cptJeps}

\begin{theorem}[blueprint="infogeometry:canonical:tomitatakesaki:cptjeps"]
  \inputleannode{InfoGeometry.Canonical.TomitaTakesaki.cptJeps}
\end{theorem}

\subsection*{cptJeps_sq}

\begin{theorem}[blueprint="infogeometry:canonical:tomitatakesaki:cptjeps_sq"]
  \inputleannode{InfoGeometry.Canonical.TomitaTakesaki.cptJeps_sq}
\end{theorem}

\subsection*{iota_mem_adjoin_cptAtomSet}

\begin{theorem}[blueprint="infogeometry:canonical:tomitatakesaki:iota_mem_adjoin_cptatomset"]
  \inputleannode{InfoGeometry.Canonical.TomitaTakesaki.iota_mem_adjoin_cptAtomSet}
\end{theorem}

\subsection*{iota_range_subset_adjoin_cptAtomSet}

\begin{theorem}[blueprint="infogeometry:canonical:tomitatakesaki:iota_range_subset_adjoin_cptatomset"]
  \inputleannode{InfoGeometry.Canonical.TomitaTakesaki.iota_range_subset_adjoin_cptAtomSet}
\end{theorem}

\subsection*{modularComplexI}

\begin{theorem}[blueprint="infogeometry:canonical:tomitatakesaki:modularcomplexi"]
  \inputleannode{InfoGeometry.Canonical.TomitaTakesaki.modularComplexI}
\end{theorem}

\subsection*{modularComplexI_sq}

\begin{theorem}[blueprint="infogeometry:canonical:tomitatakesaki:modularcomplexi_sq"]
  \inputleannode{InfoGeometry.Canonical.TomitaTakesaki.modularComplexI_sq}
\end{theorem}

\subsection*{modularConjugationJ}

\begin{theorem}[blueprint="infogeometry:canonical:tomitatakesaki:modularconjugationj"]
  \inputleannode{InfoGeometry.Canonical.TomitaTakesaki.modularConjugationJ}
\end{theorem}

\subsection*{modularConjugationJ_anticommutes_modularSign}

\begin{theorem}[blueprint="infogeometry:canonical:tomitatakesaki:modularconjugationj_anticommutes_modularsign"]
  \inputleannode{InfoGeometry.Canonical.TomitaTakesaki.modularConjugationJ_anticommutes_modularSign}
\end{theorem}

\subsection*{modularConjugationJ_sq}

\begin{theorem}[blueprint="infogeometry:canonical:tomitatakesaki:modularconjugationj_sq"]
  \inputleannode{InfoGeometry.Canonical.TomitaTakesaki.modularConjugationJ_sq}
\end{theorem}

\subsection*{modularSignEpsilon}

\begin{theorem}[blueprint="infogeometry:canonical:tomitatakesaki:modularsignepsilon"]
  \inputleannode{InfoGeometry.Canonical.TomitaTakesaki.modularSignEpsilon}
\end{theorem}

\subsection*{modularSignEpsilon_sq}

\begin{theorem}[blueprint="infogeometry:canonical:tomitatakesaki:modularsignepsilon_sq"]
  \inputleannode{InfoGeometry.Canonical.TomitaTakesaki.modularSignEpsilon_sq}
\end{theorem}

\subsection*{tomitaRepresentation}

\begin{theorem}[blueprint="infogeometry:canonical:tomitatakesaki:tomitarepresentation"]
  \inputleannode{InfoGeometry.Canonical.TomitaTakesaki.tomitaRepresentation}
\end{theorem}

\subsection*{tomitaRepresentation_cptEps}

\begin{theorem}[blueprint="infogeometry:canonical:tomitatakesaki:tomitarepresentation_cpteps"]
  \inputleannode{InfoGeometry.Canonical.TomitaTakesaki.tomitaRepresentation_cptEps}
\end{theorem}

\subsection*{tomitaRepresentation_cptJ}

\begin{theorem}[blueprint="infogeometry:canonical:tomitatakesaki:tomitarepresentation_cptj"]
  \inputleannode{InfoGeometry.Canonical.TomitaTakesaki.tomitaRepresentation_cptJ}
\end{theorem}

\subsection*{tomitaRepresentation_cptJeps}

\begin{theorem}[blueprint="infogeometry:canonical:tomitatakesaki:tomitarepresentation_cptjeps"]
  \inputleannode{InfoGeometry.Canonical.TomitaTakesaki.tomitaRepresentation_cptJeps}
\end{theorem}

\section{InfoGeometry.Canonical.TomitaTakesaki.cptAtomSet}

\subsection*{eq_1}

\begin{theorem}[blueprint="infogeometry:canonical:tomitatakesaki:cptatomset:eq_1"]
  \inputleannode{InfoGeometry.Canonical.TomitaTakesaki.cptAtomSet.eq_1}
\end{theorem}

\section{InfoGeometry.Canonical.TomitaTakesaki.cptEps}

\subsection*{eq_1}

\begin{theorem}[blueprint="infogeometry:canonical:tomitatakesaki:cpteps:eq_1"]
  \inputleannode{InfoGeometry.Canonical.TomitaTakesaki.cptEps.eq_1}
\end{theorem}

\section{InfoGeometry.Canonical.TomitaTakesaki.cptJ}

\subsection*{eq_1}

\begin{theorem}[blueprint="infogeometry:canonical:tomitatakesaki:cptj:eq_1"]
  \inputleannode{InfoGeometry.Canonical.TomitaTakesaki.cptJ.eq_1}
\end{theorem}

\section{InfoGeometry.Canonical.TomitaTakesaki.cptJeps}

\subsection*{eq_1}

\begin{theorem}[blueprint="infogeometry:canonical:tomitatakesaki:cptjeps:eq_1"]
  \inputleannode{InfoGeometry.Canonical.TomitaTakesaki.cptJeps.eq_1}
\end{theorem}

\section{InfoGeometry.Canonical.TopologicalEuler}

\noindent\textbf{Buildable}: yes

\subsection*{HasInformationHoles}

\begin{theorem}[blueprint="infogeometry:canonical:topologicaleuler:hasinformationholes"]
  \inputleannode{InfoGeometry.Canonical.TopologicalEuler.HasInformationHoles}
\end{theorem}

\subsection*{eulerCharacteristic}

\begin{theorem}[blueprint="infogeometry:canonical:topologicaleuler:eulercharacteristic"]
  \inputleannode{InfoGeometry.Canonical.TopologicalEuler.eulerCharacteristic}
\end{theorem}

\subsection*{euler_is_invariant}

\begin{theorem}[blueprint="infogeometry:canonical:topologicaleuler:euler_is_invariant"]
  \inputleannode{InfoGeometry.Canonical.TopologicalEuler.euler_is_invariant}
\end{theorem}

\section{InfoGeometry.Canonical.TopologicalEuler.eulerCharacteristic}

\subsection*{eq_1}

\begin{theorem}[blueprint="infogeometry:canonical:topologicaleuler:eulercharacteristic:eq_1"]
  \inputleannode{InfoGeometry.Canonical.TopologicalEuler.eulerCharacteristic.eq_1}
\end{theorem}

\section{InfoGeometry.Canonical.TopologicalInvariants}

\noindent\textbf{Buildable}: yes

\subsection*{BayesianLoop}

\begin{theorem}[blueprint="infogeometry:canonical:topologicalinvariants:bayesianloop"]
  \inputleannode{InfoGeometry.Canonical.TopologicalInvariants.BayesianLoop}
\end{theorem}

\subsection*{chiralAnomalyIndex}

\begin{theorem}[blueprint="infogeometry:canonical:topologicalinvariants:chiralanomalyindex"]
  \inputleannode{InfoGeometry.Canonical.TopologicalInvariants.chiralAnomalyIndex}
\end{theorem}

\subsection*{cs_invariant_of_flat}

\begin{theorem}[blueprint="infogeometry:canonical:topologicalinvariants:cs_invariant_of_flat"]
  \inputleannode{InfoGeometry.Canonical.TopologicalInvariants.cs_invariant_of_flat}
\end{theorem}

\subsection*{informationChernSimons}

\begin{theorem}[blueprint="infogeometry:canonical:topologicalinvariants:informationchernsimons"]
  \inputleannode{InfoGeometry.Canonical.TopologicalInvariants.informationChernSimons}
\end{theorem}

\section{InfoGeometry.Canonical.TopologicalInvariants.BayesianLoop}

\subsection*{N}

\begin{theorem}[blueprint="infogeometry:canonical:topologicalinvariants:bayesianloop:n"]
  \inputleannode{InfoGeometry.Canonical.TopologicalInvariants.BayesianLoop.N}
\end{theorem}

\subsection*{casesOn}

\begin{theorem}[blueprint="infogeometry:canonical:topologicalinvariants:bayesianloop:caseson"]
  \inputleannode{InfoGeometry.Canonical.TopologicalInvariants.BayesianLoop.casesOn}
\end{theorem}

\subsection*{ctorIdx}

\begin{theorem}[blueprint="infogeometry:canonical:topologicalinvariants:bayesianloop:ctoridx"]
  \inputleannode{InfoGeometry.Canonical.TopologicalInvariants.BayesianLoop.ctorIdx}
\end{theorem}

\subsection*{mk}

\begin{theorem}[blueprint="infogeometry:canonical:topologicalinvariants:bayesianloop:mk"]
  \inputleannode{InfoGeometry.Canonical.TopologicalInvariants.BayesianLoop.mk}
\end{theorem}

\subsection*{noConfusion}

\begin{theorem}[blueprint="infogeometry:canonical:topologicalinvariants:bayesianloop:noconfusion"]
  \inputleannode{InfoGeometry.Canonical.TopologicalInvariants.BayesianLoop.noConfusion}
\end{theorem}

\subsection*{noConfusionType}

\begin{theorem}[blueprint="infogeometry:canonical:topologicalinvariants:bayesianloop:noconfusiontype"]
  \inputleannode{InfoGeometry.Canonical.TopologicalInvariants.BayesianLoop.noConfusionType}
\end{theorem}

\subsection*{periodic}

\begin{theorem}[blueprint="infogeometry:canonical:topologicalinvariants:bayesianloop:periodic"]
  \inputleannode{InfoGeometry.Canonical.TopologicalInvariants.BayesianLoop.periodic}
\end{theorem}

\subsection*{rec}

\begin{theorem}[blueprint="infogeometry:canonical:topologicalinvariants:bayesianloop:rec"]
  \inputleannode{InfoGeometry.Canonical.TopologicalInvariants.BayesianLoop.rec}
\end{theorem}

\subsection*{recOn}

\begin{theorem}[blueprint="infogeometry:canonical:topologicalinvariants:bayesianloop:recon"]
  \inputleannode{InfoGeometry.Canonical.TopologicalInvariants.BayesianLoop.recOn}
\end{theorem}

\subsection*{γ}

\begin{theorem}[blueprint="infogeometry:canonical:topologicalinvariants:bayesianloop:γ"]
  \inputleannode{InfoGeometry.Canonical.TopologicalInvariants.BayesianLoop.γ}
\end{theorem}

\section{InfoGeometry.Canonical.TopologicalInvariants.chiralAnomalyIndex}

\subsection*{eq_1}

\begin{theorem}[blueprint="infogeometry:canonical:topologicalinvariants:chiralanomalyindex:eq_1"]
  \inputleannode{InfoGeometry.Canonical.TopologicalInvariants.chiralAnomalyIndex.eq_1}
\end{theorem}

\section{InfoGeometry.Canonical.TopologicalInvariants.informationChernSimons}

\subsection*{eq_1}

\begin{theorem}[blueprint="infogeometry:canonical:topologicalinvariants:informationchernsimons:eq_1"]
  \inputleannode{InfoGeometry.Canonical.TopologicalInvariants.informationChernSimons.eq_1}
\end{theorem}

\section{InfoGeometry.Canonical.Triality}

\noindent\textbf{Buildable}: yes

\subsection*{GeometricAttentionMap}

\begin{theorem}[blueprint="infogeometry:canonical:triality:geometricattentionmap"]
  \inputleannode{InfoGeometry.Canonical.Triality.GeometricAttentionMap}
\end{theorem}

\subsection*{MetricTriadicCore}

\begin{theorem}[blueprint="infogeometry:canonical:triality:metrictriadiccore"]
  \inputleannode{InfoGeometry.Canonical.Triality.MetricTriadicCore}
\end{theorem}

\subsection*{MultiHeadGeometricAttention}

\begin{theorem}[blueprint="infogeometry:canonical:triality:multiheadgeometricattention"]
  \inputleannode{InfoGeometry.Canonical.Triality.MultiHeadGeometricAttention}
\end{theorem}

\subsection*{TriadicCore}

\begin{theorem}[blueprint="infogeometry:canonical:triality:triadiccore"]
  \inputleannode{InfoGeometry.Canonical.Triality.TriadicCore}
\end{theorem}

\subsection*{softmaxAttention}

\begin{theorem}[blueprint="infogeometry:canonical:triality:softmaxattention"]
  \inputleannode{InfoGeometry.Canonical.Triality.softmaxAttention}
\end{theorem}

\subsection*{splitMetricTriadicInstance}

\begin{theorem}[blueprint="infogeometry:canonical:triality:splitmetrictriadicinstance"]
  \inputleannode{InfoGeometry.Canonical.Triality.splitMetricTriadicInstance}
\end{theorem}

\subsection*{splitMetricTriadicInstance_interact}

\begin{theorem}[blueprint="infogeometry:canonical:triality:splitmetrictriadicinstance_interact"]
  \inputleannode{InfoGeometry.Canonical.Triality.splitMetricTriadicInstance_interact}
\end{theorem}

\subsection*{splitMetricTriadicInstance_quadK}

\begin{theorem}[blueprint="infogeometry:canonical:triality:splitmetrictriadicinstance_quadk"]
  \inputleannode{InfoGeometry.Canonical.Triality.splitMetricTriadicInstance_quadK}
\end{theorem}

\subsection*{splitMetricTriadicInstance_quadQ}

\begin{theorem}[blueprint="infogeometry:canonical:triality:splitmetrictriadicinstance_quadq"]
  \inputleannode{InfoGeometry.Canonical.Triality.splitMetricTriadicInstance_quadQ}
\end{theorem}

\subsection*{splitMetricTriadicInstance_quadV}

\begin{theorem}[blueprint="infogeometry:canonical:triality:splitmetrictriadicinstance_quadv"]
  \inputleannode{InfoGeometry.Canonical.Triality.splitMetricTriadicInstance_quadV}
\end{theorem}

\subsection*{splitMetricTriadicInstance_route}

\begin{theorem}[blueprint="infogeometry:canonical:triality:splitmetrictriadicinstance_route"]
  \inputleannode{InfoGeometry.Canonical.Triality.splitMetricTriadicInstance_route}
\end{theorem}

\subsection*{splitMetricTriadicInstance_route_norm_compat}

\begin{theorem}[blueprint="infogeometry:canonical:triality:splitmetrictriadicinstance_route_norm_compat"]
  \inputleannode{InfoGeometry.Canonical.Triality.splitMetricTriadicInstance_route_norm_compat}
\end{theorem}

\section{InfoGeometry.Canonical.Triality.GeometricAttentionMap}

\subsection*{I}

\begin{theorem}[blueprint="infogeometry:canonical:triality:geometricattentionmap:i"]
  \inputleannode{InfoGeometry.Canonical.Triality.GeometricAttentionMap.I}
\end{theorem}

\subsection*{attention}

\begin{theorem}[blueprint="infogeometry:canonical:triality:geometricattentionmap:attention"]
  \inputleannode{InfoGeometry.Canonical.Triality.GeometricAttentionMap.attention}
\end{theorem}

\subsection*{attention_decomposition_residual}

\begin{theorem}[blueprint="infogeometry:canonical:triality:geometricattentionmap:attention_decomposition_residual"]
  \inputleannode{InfoGeometry.Canonical.Triality.GeometricAttentionMap.attention_decomposition_residual}
\end{theorem}

\subsection*{attention_def}

\begin{theorem}[blueprint="infogeometry:canonical:triality:geometricattentionmap:attention_def"]
  \inputleannode{InfoGeometry.Canonical.Triality.GeometricAttentionMap.attention_def}
\end{theorem}

\subsection*{casesOn}

\begin{theorem}[blueprint="infogeometry:canonical:triality:geometricattentionmap:caseson"]
  \inputleannode{InfoGeometry.Canonical.Triality.GeometricAttentionMap.casesOn}
\end{theorem}

\subsection*{ctorIdx}

\begin{theorem}[blueprint="infogeometry:canonical:triality:geometricattentionmap:ctoridx"]
  \inputleannode{InfoGeometry.Canonical.Triality.GeometricAttentionMap.ctorIdx}
\end{theorem}

\subsection*{keys}

\begin{theorem}[blueprint="infogeometry:canonical:triality:geometricattentionmap:keys"]
  \inputleannode{InfoGeometry.Canonical.Triality.GeometricAttentionMap.keys}
\end{theorem}

\subsection*{mk}

\begin{theorem}[blueprint="infogeometry:canonical:triality:geometricattentionmap:mk"]
  \inputleannode{InfoGeometry.Canonical.Triality.GeometricAttentionMap.mk}
\end{theorem}

\subsection*{noConfusion}

\begin{theorem}[blueprint="infogeometry:canonical:triality:geometricattentionmap:noconfusion"]
  \inputleannode{InfoGeometry.Canonical.Triality.GeometricAttentionMap.noConfusion}
\end{theorem}

\subsection*{noConfusionType}

\begin{theorem}[blueprint="infogeometry:canonical:triality:geometricattentionmap:noconfusiontype"]
  \inputleannode{InfoGeometry.Canonical.Triality.GeometricAttentionMap.noConfusionType}
\end{theorem}

\subsection*{rec}

\begin{theorem}[blueprint="infogeometry:canonical:triality:geometricattentionmap:rec"]
  \inputleannode{InfoGeometry.Canonical.Triality.GeometricAttentionMap.rec}
\end{theorem}

\subsection*{recOn}

\begin{theorem}[blueprint="infogeometry:canonical:triality:geometricattentionmap:recon"]
  \inputleannode{InfoGeometry.Canonical.Triality.GeometricAttentionMap.recOn}
\end{theorem}

\subsection*{weights}

\begin{theorem}[blueprint="infogeometry:canonical:triality:geometricattentionmap:weights"]
  \inputleannode{InfoGeometry.Canonical.Triality.GeometricAttentionMap.weights}
\end{theorem}

\subsection*{weights_sum_one}

\begin{theorem}[blueprint="infogeometry:canonical:triality:geometricattentionmap:weights_sum_one"]
  \inputleannode{InfoGeometry.Canonical.Triality.GeometricAttentionMap.weights_sum_one}
\end{theorem}

\section{InfoGeometry.Canonical.Triality.MetricTriadicCore}

\subsection*{casesOn}

\begin{theorem}[blueprint="infogeometry:canonical:triality:metrictriadiccore:caseson"]
  \inputleannode{InfoGeometry.Canonical.Triality.MetricTriadicCore.casesOn}
\end{theorem}

\subsection*{ctorIdx}

\begin{theorem}[blueprint="infogeometry:canonical:triality:metrictriadiccore:ctoridx"]
  \inputleannode{InfoGeometry.Canonical.Triality.MetricTriadicCore.ctorIdx}
\end{theorem}

\subsection*{mk}

\begin{theorem}[blueprint="infogeometry:canonical:triality:metrictriadiccore:mk"]
  \inputleannode{InfoGeometry.Canonical.Triality.MetricTriadicCore.mk}
\end{theorem}

\subsection*{noConfusion}

\begin{theorem}[blueprint="infogeometry:canonical:triality:metrictriadiccore:noconfusion"]
  \inputleannode{InfoGeometry.Canonical.Triality.MetricTriadicCore.noConfusion}
\end{theorem}

\subsection*{noConfusionType}

\begin{theorem}[blueprint="infogeometry:canonical:triality:metrictriadiccore:noconfusiontype"]
  \inputleannode{InfoGeometry.Canonical.Triality.MetricTriadicCore.noConfusionType}
\end{theorem}

\subsection*{quadK}

\begin{theorem}[blueprint="infogeometry:canonical:triality:metrictriadiccore:quadk"]
  \inputleannode{InfoGeometry.Canonical.Triality.MetricTriadicCore.quadK}
\end{theorem}

\subsection*{quadQ}

\begin{theorem}[blueprint="infogeometry:canonical:triality:metrictriadiccore:quadq"]
  \inputleannode{InfoGeometry.Canonical.Triality.MetricTriadicCore.quadQ}
\end{theorem}

\subsection*{quadV}

\begin{theorem}[blueprint="infogeometry:canonical:triality:metrictriadiccore:quadv"]
  \inputleannode{InfoGeometry.Canonical.Triality.MetricTriadicCore.quadV}
\end{theorem}

\subsection*{rec}

\begin{theorem}[blueprint="infogeometry:canonical:triality:metrictriadiccore:rec"]
  \inputleannode{InfoGeometry.Canonical.Triality.MetricTriadicCore.rec}
\end{theorem}

\subsection*{recOn}

\begin{theorem}[blueprint="infogeometry:canonical:triality:metrictriadiccore:recon"]
  \inputleannode{InfoGeometry.Canonical.Triality.MetricTriadicCore.recOn}
\end{theorem}

\subsection*{route_norm_compat}

\begin{theorem}[blueprint="infogeometry:canonical:triality:metrictriadiccore:route_norm_compat"]
  \inputleannode{InfoGeometry.Canonical.Triality.MetricTriadicCore.route_norm_compat}
\end{theorem}

\subsection*{toTriadicCore}

\begin{theorem}[blueprint="infogeometry:canonical:triality:metrictriadiccore:totriadiccore"]
  \inputleannode{InfoGeometry.Canonical.Triality.MetricTriadicCore.toTriadicCore}
\end{theorem}

\section{InfoGeometry.Canonical.Triality.MultiHeadGeometricAttention}

\subsection*{casesOn}

\begin{theorem}[blueprint="infogeometry:canonical:triality:multiheadgeometricattention:caseson"]
  \inputleannode{InfoGeometry.Canonical.Triality.MultiHeadGeometricAttention.casesOn}
\end{theorem}

\subsection*{cores}

\begin{theorem}[blueprint="infogeometry:canonical:triality:multiheadgeometricattention:cores"]
  \inputleannode{InfoGeometry.Canonical.Triality.MultiHeadGeometricAttention.cores}
\end{theorem}

\subsection*{ctorIdx}

\begin{theorem}[blueprint="infogeometry:canonical:triality:multiheadgeometricattention:ctoridx"]
  \inputleannode{InfoGeometry.Canonical.Triality.MultiHeadGeometricAttention.ctorIdx}
\end{theorem}

\subsection*{heads}

\begin{theorem}[blueprint="infogeometry:canonical:triality:multiheadgeometricattention:heads"]
  \inputleannode{InfoGeometry.Canonical.Triality.MultiHeadGeometricAttention.heads}
\end{theorem}

\subsection*{mk}

\begin{theorem}[blueprint="infogeometry:canonical:triality:multiheadgeometricattention:mk"]
  \inputleannode{InfoGeometry.Canonical.Triality.MultiHeadGeometricAttention.mk}
\end{theorem}

\subsection*{noConfusion}

\begin{theorem}[blueprint="infogeometry:canonical:triality:multiheadgeometricattention:noconfusion"]
  \inputleannode{InfoGeometry.Canonical.Triality.MultiHeadGeometricAttention.noConfusion}
\end{theorem}

\subsection*{noConfusionType}

\begin{theorem}[blueprint="infogeometry:canonical:triality:multiheadgeometricattention:noconfusiontype"]
  \inputleannode{InfoGeometry.Canonical.Triality.MultiHeadGeometricAttention.noConfusionType}
\end{theorem}

\subsection*{outProj}

\begin{theorem}[blueprint="infogeometry:canonical:triality:multiheadgeometricattention:outproj"]
  \inputleannode{InfoGeometry.Canonical.Triality.MultiHeadGeometricAttention.outProj}
\end{theorem}

\subsection*{output}

\begin{theorem}[blueprint="infogeometry:canonical:triality:multiheadgeometricattention:output"]
  \inputleannode{InfoGeometry.Canonical.Triality.MultiHeadGeometricAttention.output}
\end{theorem}

\subsection*{output_decomposition_residual}

\begin{theorem}[blueprint="infogeometry:canonical:triality:multiheadgeometricattention:output_decomposition_residual"]
  \inputleannode{InfoGeometry.Canonical.Triality.MultiHeadGeometricAttention.output_decomposition_residual}
\end{theorem}

\subsection*{output_def}

\begin{theorem}[blueprint="infogeometry:canonical:triality:multiheadgeometricattention:output_def"]
  \inputleannode{InfoGeometry.Canonical.Triality.MultiHeadGeometricAttention.output_def}
\end{theorem}

\subsection*{preOutput}

\begin{theorem}[blueprint="infogeometry:canonical:triality:multiheadgeometricattention:preoutput"]
  \inputleannode{InfoGeometry.Canonical.Triality.MultiHeadGeometricAttention.preOutput}
\end{theorem}

\subsection*{preOutput_def}

\begin{theorem}[blueprint="infogeometry:canonical:triality:multiheadgeometricattention:preoutput_def"]
  \inputleannode{InfoGeometry.Canonical.Triality.MultiHeadGeometricAttention.preOutput_def}
\end{theorem}

\subsection*{rec}

\begin{theorem}[blueprint="infogeometry:canonical:triality:multiheadgeometricattention:rec"]
  \inputleannode{InfoGeometry.Canonical.Triality.MultiHeadGeometricAttention.rec}
\end{theorem}

\subsection*{recOn}

\begin{theorem}[blueprint="infogeometry:canonical:triality:multiheadgeometricattention:recon"]
  \inputleannode{InfoGeometry.Canonical.Triality.MultiHeadGeometricAttention.recOn}
\end{theorem}

\section{InfoGeometry.Canonical.Triality.TriadicCore}

\subsection*{casesOn}

\begin{theorem}[blueprint="infogeometry:canonical:triality:triadiccore:caseson"]
  \inputleannode{InfoGeometry.Canonical.Triality.TriadicCore.casesOn}
\end{theorem}

\subsection*{ctorIdx}

\begin{theorem}[blueprint="infogeometry:canonical:triality:triadiccore:ctoridx"]
  \inputleannode{InfoGeometry.Canonical.Triality.TriadicCore.ctorIdx}
\end{theorem}

\subsection*{interact}

\begin{theorem}[blueprint="infogeometry:canonical:triality:triadiccore:interact"]
  \inputleannode{InfoGeometry.Canonical.Triality.TriadicCore.interact}
\end{theorem}

\subsection*{mk}

\begin{theorem}[blueprint="infogeometry:canonical:triality:triadiccore:mk"]
  \inputleannode{InfoGeometry.Canonical.Triality.TriadicCore.mk}
\end{theorem}

\subsection*{noConfusion}

\begin{theorem}[blueprint="infogeometry:canonical:triality:triadiccore:noconfusion"]
  \inputleannode{InfoGeometry.Canonical.Triality.TriadicCore.noConfusion}
\end{theorem}

\subsection*{noConfusionType}

\begin{theorem}[blueprint="infogeometry:canonical:triality:triadiccore:noconfusiontype"]
  \inputleannode{InfoGeometry.Canonical.Triality.TriadicCore.noConfusionType}
\end{theorem}

\subsection*{rec}

\begin{theorem}[blueprint="infogeometry:canonical:triality:triadiccore:rec"]
  \inputleannode{InfoGeometry.Canonical.Triality.TriadicCore.rec}
\end{theorem}

\subsection*{recOn}

\begin{theorem}[blueprint="infogeometry:canonical:triality:triadiccore:recon"]
  \inputleannode{InfoGeometry.Canonical.Triality.TriadicCore.recOn}
\end{theorem}

\subsection*{route}

\begin{theorem}[blueprint="infogeometry:canonical:triality:triadiccore:route"]
  \inputleannode{InfoGeometry.Canonical.Triality.TriadicCore.route}
\end{theorem}

\section{InfoGeometry.Canonical.Twistor}

\noindent\textbf{Buildable}: yes

\noindent\emph{No exported declarations in the current declaration graph.}

\section{InfoGeometry.Canonical.WeylInformationGauge}

\noindent\textbf{Buildable}: yes

\subsection*{UpdateOrderHysteresis}

\begin{theorem}[blueprint="infogeometry:canonical:weylinformationgauge:updateorderhysteresis"]
  \inputleannode{InfoGeometry.Canonical.WeylInformationGauge.UpdateOrderHysteresis}
\end{theorem}

\subsection*{UpdateOrderPathDependent}

\begin{theorem}[blueprint="infogeometry:canonical:weylinformationgauge:updateorderpathdependent"]
  \inputleannode{InfoGeometry.Canonical.WeylInformationGauge.UpdateOrderPathDependent}
\end{theorem}

\subsection*{chiralInferenceState_of_nonzero_anomaly}

\begin{theorem}[blueprint="infogeometry:canonical:weylinformationgauge:chiralinferencestate_of_nonzero_anomaly"]
  \inputleannode{InfoGeometry.Canonical.WeylInformationGauge.chiralInferenceState_of_nonzero_anomaly}
\end{theorem}

\subsection*{colThenRowUpdate}

\begin{theorem}[blueprint="infogeometry:canonical:weylinformationgauge:colthenrowupdate"]
  \inputleannode{InfoGeometry.Canonical.WeylInformationGauge.colThenRowUpdate}
\end{theorem}

\subsection*{dilationAnomaly_sources_transportedEinsteinResidual}

\begin{theorem}[blueprint="infogeometry:canonical:weylinformationgauge:dilationanomaly_sources_transportedeinsteinresidual"]
  \inputleannode{InfoGeometry.Canonical.WeylInformationGauge.dilationAnomaly_sources_transportedEinsteinResidual}
\end{theorem}

\subsection*{exists_updateOrderHysteresis_n2}

\begin{theorem}[blueprint="infogeometry:canonical:weylinformationgauge:exists_updateorderhysteresis_n2"]
  \inputleannode{InfoGeometry.Canonical.WeylInformationGauge.exists_updateOrderHysteresis_n2}
\end{theorem}

\subsection*{nonzeroAnomaly_sources_transportedEinsteinResidual}

\begin{theorem}[blueprint="infogeometry:canonical:weylinformationgauge:nonzeroanomaly_sources_transportedeinsteinresidual"]
  \inputleannode{InfoGeometry.Canonical.WeylInformationGauge.nonzeroAnomaly_sources_transportedEinsteinResidual}
\end{theorem}

\subsection*{rowThenColUpdate}

\begin{theorem}[blueprint="infogeometry:canonical:weylinformationgauge:rowthencolupdate"]
  \inputleannode{InfoGeometry.Canonical.WeylInformationGauge.rowThenColUpdate}
\end{theorem}

\subsection*{sinkhornGaugeFixing_is_marginal_closure}

\begin{theorem}[blueprint="infogeometry:canonical:weylinformationgauge:sinkhorngaugefixing_is_marginal_closure"]
  \inputleannode{InfoGeometry.Canonical.WeylInformationGauge.sinkhornGaugeFixing_is_marginal_closure}
\end{theorem}

\subsection*{sinkhornTwoStep_eq_informationWeylGauge}

\begin{theorem}[blueprint="infogeometry:canonical:weylinformationgauge:sinkhorntwostep_eq_informationweylgauge"]
  \inputleannode{InfoGeometry.Canonical.WeylInformationGauge.sinkhornTwoStep_eq_informationWeylGauge}
\end{theorem}

\subsection*{twistedInference_not_torsionFree}

\begin{theorem}[blueprint="infogeometry:canonical:weylinformationgauge:twistedinference_not_torsionfree"]
  \inputleannode{InfoGeometry.Canonical.WeylInformationGauge.twistedInference_not_torsionFree}
\end{theorem}

\subsection*{twistedInference_updateOrderPathDependent}

\begin{theorem}[blueprint="infogeometry:canonical:weylinformationgauge:twistedinference_updateorderpathdependent"]
  \inputleannode{InfoGeometry.Canonical.WeylInformationGauge.twistedInference_updateOrderPathDependent}
\end{theorem}

\subsection*{weylOrderHysteresis_on_witnessMatrix2}

\begin{theorem}[blueprint="infogeometry:canonical:weylinformationgauge:weylorderhysteresis_on_witnessmatrix2"]
  \inputleannode{InfoGeometry.Canonical.WeylInformationGauge.weylOrderHysteresis_on_witnessMatrix2}
\end{theorem}

\subsection*{weylOrderWitnessMatrix2}

\begin{theorem}[blueprint="infogeometry:canonical:weylinformationgauge:weylorderwitnessmatrix2"]
  \inputleannode{InfoGeometry.Canonical.WeylInformationGauge.weylOrderWitnessMatrix2}
\end{theorem}

\subsection*{weylOrderWitnessMatrix2_positiveCols}

\begin{theorem}[blueprint="infogeometry:canonical:weylinformationgauge:weylorderwitnessmatrix2_positivecols"]
  \inputleannode{InfoGeometry.Canonical.WeylInformationGauge.weylOrderWitnessMatrix2_positiveCols}
\end{theorem}

\subsection*{weylOrderWitnessMatrix2_positiveCols_afterRow}

\begin{theorem}[blueprint="infogeometry:canonical:weylinformationgauge:weylorderwitnessmatrix2_positivecols_afterrow"]
  \inputleannode{InfoGeometry.Canonical.WeylInformationGauge.weylOrderWitnessMatrix2_positiveCols_afterRow}
\end{theorem}

\subsection*{weylOrderWitnessMatrix2_positiveRows}

\begin{theorem}[blueprint="infogeometry:canonical:weylinformationgauge:weylorderwitnessmatrix2_positiverows"]
  \inputleannode{InfoGeometry.Canonical.WeylInformationGauge.weylOrderWitnessMatrix2_positiveRows}
\end{theorem}

\subsection*{weylOrderWitnessMatrix2_positiveRows_afterCol}

\begin{theorem}[blueprint="infogeometry:canonical:weylinformationgauge:weylorderwitnessmatrix2_positiverows_aftercol"]
  \inputleannode{InfoGeometry.Canonical.WeylInformationGauge.weylOrderWitnessMatrix2_positiveRows_afterCol}
\end{theorem}

\section{InfoGeometry.Canonical.WeylInformationGauge.weylOrderWitnessMatrix2}

\subsection*{eq_1}

\begin{theorem}[blueprint="infogeometry:canonical:weylinformationgauge:weylorderwitnessmatrix2:eq_1"]
  \inputleannode{InfoGeometry.Canonical.WeylInformationGauge.weylOrderWitnessMatrix2.eq_1}
\end{theorem}

\chapter{Root Modules}

\section{InfoGeometry}

\noindent\textbf{Buildable}: yes

\subsection*{EmpiricalCounts}

\begin{theorem}[blueprint="infogeometry:empiricalcounts"]
  \inputleannode{InfoGeometry.EmpiricalCounts}
\end{theorem}

\subsection*{FinProb}

\begin{theorem}[blueprint="infogeometry:finprob"]
  \inputleannode{InfoGeometry.FinProb}
\end{theorem}

\subsection*{LogPotential}

\begin{theorem}[blueprint="infogeometry:logpotential"]
  \inputleannode{InfoGeometry.LogPotential}
\end{theorem}

\subsection*{ProbabilityDist}

\begin{theorem}[blueprint="infogeometry:probabilitydist"]
  \inputleannode{InfoGeometry.ProbabilityDist}
\end{theorem}

\subsection*{RenyiD}

\begin{theorem}[blueprint="infogeometry:renyid"]
  \inputleannode{InfoGeometry.RenyiD}
\end{theorem}

\subsection*{RenyiDShifted}

\begin{theorem}[blueprint="infogeometry:renyidshifted"]
  \inputleannode{InfoGeometry.RenyiDShifted}
\end{theorem}

\subsection*{RenyiDShifted_eq_log_Phi_div}

\begin{theorem}[blueprint="infogeometry:renyidshifted_eq_log_phi_div"]
  \inputleannode{InfoGeometry.RenyiDShifted_eq_log_Phi_div}
\end{theorem}

\subsection*{RenyiD_eq_log_Phi_shift_div}

\begin{theorem}[blueprint="infogeometry:renyid_eq_log_phi_shift_div"]
  \inputleannode{InfoGeometry.RenyiD_eq_log_Phi_shift_div}
\end{theorem}

\subsection*{RenyiD_mul_order_sub_one}

\begin{theorem}[blueprint="infogeometry:renyid_mul_order_sub_one"]
  \inputleannode{InfoGeometry.RenyiD_mul_order_sub_one}
\end{theorem}

\subsection*{RenyiSupportFaithful}

\begin{theorem}[blueprint="infogeometry:renyisupportfaithful"]
  \inputleannode{InfoGeometry.RenyiSupportFaithful}
\end{theorem}

\subsection*{StrictProbabilityDist}

\begin{theorem}[blueprint="infogeometry:strictprobabilitydist"]
  \inputleannode{InfoGeometry.StrictProbabilityDist}
\end{theorem}

\subsection*{bregmanDiv}

\begin{theorem}[blueprint="infogeometry:bregmandiv"]
  \inputleannode{InfoGeometry.bregmanDiv}
\end{theorem}

\subsection*{bregmanPythagoreanIneq}

\begin{theorem}[blueprint="infogeometry:bregmanpythagoreanineq"]
  \inputleannode{InfoGeometry.bregmanPythagoreanIneq}
\end{theorem}

\subsection*{bregmanPythagoreanIneq_of_crossTerm_nonneg}

\begin{theorem}[blueprint="infogeometry:bregmanpythagoreanineq_of_crossterm_nonneg"]
  \inputleannode{InfoGeometry.bregmanPythagoreanIneq_of_crossTerm_nonneg}
\end{theorem}

\subsection*{bregmanThreePoint}

\begin{theorem}[blueprint="infogeometry:bregmanthreepoint"]
  \inputleannode{InfoGeometry.bregmanThreePoint}
\end{theorem}

\subsection*{bregmanThreePoint_sub}

\begin{theorem}[blueprint="infogeometry:bregmanthreepoint_sub"]
  \inputleannode{InfoGeometry.bregmanThreePoint_sub}
\end{theorem}

\subsection*{dirac}

\begin{theorem}[blueprint="infogeometry:dirac"]
  \inputleannode{InfoGeometry.dirac}
\end{theorem}

\subsection*{entropy}

\begin{theorem}[blueprint="infogeometry:entropy"]
  \inputleannode{InfoGeometry.entropy}
\end{theorem}

\subsection*{expectation}

\begin{theorem}[blueprint="infogeometry:expectation"]
  \inputleannode{InfoGeometry.expectation}
\end{theorem}

\subsection*{expectation_coe_FinProb}

\begin{theorem}[blueprint="infogeometry:expectation_coe_finprob"]
  \inputleannode{InfoGeometry.expectation_coe_FinProb}
\end{theorem}

\subsection*{instCoeFinProbProbabilityDist}

\begin{theorem}[blueprint="infogeometry:instcoefinprobprobabilitydist"]
  \inputleannode{InfoGeometry.instCoeFinProbProbabilityDist}
\end{theorem}

\subsection*{instCoeFunEmpiricalCountsForallNat}

\begin{theorem}[blueprint="infogeometry:instcoefunempiricalcountsforallnat"]
  \inputleannode{InfoGeometry.instCoeFunEmpiricalCountsForallNat}
\end{theorem}

\subsection*{instCoeFunFinProbForallReal}

\begin{theorem}[blueprint="infogeometry:instcoefunfinprobforallreal"]
  \inputleannode{InfoGeometry.instCoeFunFinProbForallReal}
\end{theorem}

\subsection*{instCoeFunProbabilityDistForallReal}

\begin{theorem}[blueprint="infogeometry:instcoefunprobabilitydistforallreal"]
  \inputleannode{InfoGeometry.instCoeFunProbabilityDistForallReal}
\end{theorem}

\subsection*{instCoeFunStrictProbabilityDistForallReal}

\begin{theorem}[blueprint="infogeometry:instcoefunstrictprobabilitydistforallreal"]
  \inputleannode{InfoGeometry.instCoeFunStrictProbabilityDistForallReal}
\end{theorem}

\subsection*{klDiv}

\begin{theorem}[blueprint="infogeometry:kldiv"]
  \inputleannode{InfoGeometry.klDiv}
\end{theorem}

\subsection*{logDensity}

\begin{theorem}[blueprint="infogeometry:logdensity"]
  \inputleannode{InfoGeometry.logDensity}
\end{theorem}

\subsection*{normalize}

\begin{theorem}[blueprint="infogeometry:normalize"]
  \inputleannode{InfoGeometry.normalize}
\end{theorem}

\subsection*{normalizeStrict}

\begin{theorem}[blueprint="infogeometry:normalizestrict"]
  \inputleannode{InfoGeometry.normalizeStrict}
\end{theorem}

\subsection*{prob_nonneg}

\begin{theorem}[blueprint="infogeometry:prob_nonneg"]
  \inputleannode{InfoGeometry.prob_nonneg}
\end{theorem}

\subsection*{sum_eq_one}

\begin{theorem}[blueprint="infogeometry:sum_eq_one"]
  \inputleannode{InfoGeometry.sum_eq_one}
\end{theorem}

\subsection*{surprisal}

\begin{theorem}[blueprint="infogeometry:surprisal"]
  \inputleannode{InfoGeometry.surprisal}
\end{theorem}

\chapter{Analysis Modules}

\section{InfoGeometry.Analysis}

\noindent\textbf{Buildable}: yes

\subsection*{Graph}

\begin{theorem}[blueprint="infogeometry:analysis:graph"]
  \inputleannode{InfoGeometry.Analysis.Graph}
\end{theorem}

\section{InfoGeometry.Analysis.Graph}

\noindent\textbf{Buildable}: yes

\subsection*{EdgeKind}

\begin{theorem}[blueprint="infogeometry:analysis:graph:edgekind"]
  \inputleannode{InfoGeometry.Analysis.Graph.EdgeKind}
\end{theorem}

\subsection*{IndexedGraph}

\begin{theorem}[blueprint="infogeometry:analysis:graph:indexedgraph"]
  \inputleannode{InfoGeometry.Analysis.Graph.IndexedGraph}
\end{theorem}

\subsection*{SimpleGraph}

\begin{theorem}[blueprint="infogeometry:analysis:graph:simplegraph"]
  \inputleannode{InfoGeometry.Analysis.Graph.SimpleGraph}
\end{theorem}

\subsection*{arrayReplicate}

\begin{theorem}[blueprint="infogeometry:analysis:graph:arrayreplicate"]
  \inputleannode{InfoGeometry.Analysis.Graph.arrayReplicate}
\end{theorem}

\subsection*{collectConsts}

\begin{theorem}[blueprint="infogeometry:analysis:graph:collectconsts"]
  \inputleannode{InfoGeometry.Analysis.Graph.collectConsts}
\end{theorem}

\subsection*{collectDeps}

\begin{theorem}[blueprint="infogeometry:analysis:graph:collectdeps"]
  \inputleannode{InfoGeometry.Analysis.Graph.collectDeps}
\end{theorem}

\subsection*{envToGraph}

\begin{theorem}[blueprint="infogeometry:analysis:graph:envtograph"]
  \inputleannode{InfoGeometry.Analysis.Graph.envToGraph}
\end{theorem}

\subsection*{envToIndexedGraph}

\begin{theorem}[blueprint="infogeometry:analysis:graph:envtoindexedgraph"]
  \inputleannode{InfoGeometry.Analysis.Graph.envToIndexedGraph}
\end{theorem}

\subsection*{indexedToGraph}

\begin{theorem}[blueprint="infogeometry:analysis:graph:indexedtograph"]
  \inputleannode{InfoGeometry.Analysis.Graph.indexedToGraph}
\end{theorem}

\subsection*{instBEqEdgeKind}

\begin{theorem}[blueprint="infogeometry:analysis:graph:instbeqedgekind"]
  \inputleannode{InfoGeometry.Analysis.Graph.instBEqEdgeKind}
\end{theorem}

\subsection*{instDecidableEqEdgeKind}

\begin{theorem}[blueprint="infogeometry:analysis:graph:instdecidableeqedgekind"]
  \inputleannode{InfoGeometry.Analysis.Graph.instDecidableEqEdgeKind}
\end{theorem}

\subsection*{instHashableEdgeKind}

\begin{theorem}[blueprint="infogeometry:analysis:graph:insthashableedgekind"]
  \inputleannode{InfoGeometry.Analysis.Graph.instHashableEdgeKind}
\end{theorem}

\subsection*{instReprEdgeKind}

\begin{theorem}[blueprint="infogeometry:analysis:graph:instrepredgekind"]
  \inputleannode{InfoGeometry.Analysis.Graph.instReprEdgeKind}
\end{theorem}

\section{InfoGeometry.Analysis.Graph.EdgeKind}

\subsection*{casesOn}

\begin{theorem}[blueprint="infogeometry:analysis:graph:edgekind:caseson"]
  \inputleannode{InfoGeometry.Analysis.Graph.EdgeKind.casesOn}
\end{theorem}

\subsection*{ctorElim}

\begin{theorem}[blueprint="infogeometry:analysis:graph:edgekind:ctorelim"]
  \inputleannode{InfoGeometry.Analysis.Graph.EdgeKind.ctorElim}
\end{theorem}

\subsection*{ctorElimType}

\begin{theorem}[blueprint="infogeometry:analysis:graph:edgekind:ctorelimtype"]
  \inputleannode{InfoGeometry.Analysis.Graph.EdgeKind.ctorElimType}
\end{theorem}

\subsection*{ctorIdx}

\begin{theorem}[blueprint="infogeometry:analysis:graph:edgekind:ctoridx"]
  \inputleannode{InfoGeometry.Analysis.Graph.EdgeKind.ctorIdx}
\end{theorem}

\subsection*{noConfusion}

\begin{theorem}[blueprint="infogeometry:analysis:graph:edgekind:noconfusion"]
  \inputleannode{InfoGeometry.Analysis.Graph.EdgeKind.noConfusion}
\end{theorem}

\subsection*{noConfusionType}

\begin{theorem}[blueprint="infogeometry:analysis:graph:edgekind:noconfusiontype"]
  \inputleannode{InfoGeometry.Analysis.Graph.EdgeKind.noConfusionType}
\end{theorem}

\subsection*{ofNat}

\begin{theorem}[blueprint="infogeometry:analysis:graph:edgekind:ofnat"]
  \inputleannode{InfoGeometry.Analysis.Graph.EdgeKind.ofNat}
\end{theorem}

\subsection*{ofNat_ctorIdx}

\begin{theorem}[blueprint="infogeometry:analysis:graph:edgekind:ofnat_ctoridx"]
  \inputleannode{InfoGeometry.Analysis.Graph.EdgeKind.ofNat_ctorIdx}
\end{theorem}

\subsection*{rec}

\begin{theorem}[blueprint="infogeometry:analysis:graph:edgekind:rec"]
  \inputleannode{InfoGeometry.Analysis.Graph.EdgeKind.rec}
\end{theorem}

\subsection*{recOn}

\begin{theorem}[blueprint="infogeometry:analysis:graph:edgekind:recon"]
  \inputleannode{InfoGeometry.Analysis.Graph.EdgeKind.recOn}
\end{theorem}

\subsection*{toCtorIdx}

\begin{theorem}[blueprint="infogeometry:analysis:graph:edgekind:toctoridx"]
  \inputleannode{InfoGeometry.Analysis.Graph.EdgeKind.toCtorIdx}
\end{theorem}

\subsection*{type}

\begin{theorem}[blueprint="infogeometry:analysis:graph:edgekind:type"]
  \inputleannode{InfoGeometry.Analysis.Graph.EdgeKind.type}
\end{theorem}

\subsection*{value}

\begin{theorem}[blueprint="infogeometry:analysis:graph:edgekind:value"]
  \inputleannode{InfoGeometry.Analysis.Graph.EdgeKind.value}
\end{theorem}

\section{InfoGeometry.Analysis.Graph.EdgeKind.type}

\subsection*{elim}

\begin{theorem}[blueprint="infogeometry:analysis:graph:edgekind:type:elim"]
  \inputleannode{InfoGeometry.Analysis.Graph.EdgeKind.type.elim}
\end{theorem}

\section{InfoGeometry.Analysis.Graph.EdgeKind.value}

\subsection*{elim}

\begin{theorem}[blueprint="infogeometry:analysis:graph:edgekind:value:elim"]
  \inputleannode{InfoGeometry.Analysis.Graph.EdgeKind.value.elim}
\end{theorem}

\section{InfoGeometry.Analysis.Graph.IndexedGraph}

\subsection*{casesOn}

\begin{theorem}[blueprint="infogeometry:analysis:graph:indexedgraph:caseson"]
  \inputleannode{InfoGeometry.Analysis.Graph.IndexedGraph.casesOn}
\end{theorem}

\subsection*{ctorIdx}

\begin{theorem}[blueprint="infogeometry:analysis:graph:indexedgraph:ctoridx"]
  \inputleannode{InfoGeometry.Analysis.Graph.IndexedGraph.ctorIdx}
\end{theorem}

\subsection*{forward}

\begin{theorem}[blueprint="infogeometry:analysis:graph:indexedgraph:forward"]
  \inputleannode{InfoGeometry.Analysis.Graph.IndexedGraph.forward}
\end{theorem}

\subsection*{mk}

\begin{theorem}[blueprint="infogeometry:analysis:graph:indexedgraph:mk"]
  \inputleannode{InfoGeometry.Analysis.Graph.IndexedGraph.mk}
\end{theorem}

\subsection*{noConfusion}

\begin{theorem}[blueprint="infogeometry:analysis:graph:indexedgraph:noconfusion"]
  \inputleannode{InfoGeometry.Analysis.Graph.IndexedGraph.noConfusion}
\end{theorem}

\subsection*{noConfusionType}

\begin{theorem}[blueprint="infogeometry:analysis:graph:indexedgraph:noconfusiontype"]
  \inputleannode{InfoGeometry.Analysis.Graph.IndexedGraph.noConfusionType}
\end{theorem}

\subsection*{nodeToIdx}

\begin{theorem}[blueprint="infogeometry:analysis:graph:indexedgraph:nodetoidx"]
  \inputleannode{InfoGeometry.Analysis.Graph.IndexedGraph.nodeToIdx}
\end{theorem}

\subsection*{nodes}

\begin{theorem}[blueprint="infogeometry:analysis:graph:indexedgraph:nodes"]
  \inputleannode{InfoGeometry.Analysis.Graph.IndexedGraph.nodes}
\end{theorem}

\subsection*{rec}

\begin{theorem}[blueprint="infogeometry:analysis:graph:indexedgraph:rec"]
  \inputleannode{InfoGeometry.Analysis.Graph.IndexedGraph.rec}
\end{theorem}

\subsection*{recOn}

\begin{theorem}[blueprint="infogeometry:analysis:graph:indexedgraph:recon"]
  \inputleannode{InfoGeometry.Analysis.Graph.IndexedGraph.recOn}
\end{theorem}

\section{InfoGeometry.Analysis.Graph.SimpleGraph}

\subsection*{casesOn}

\begin{theorem}[blueprint="infogeometry:analysis:graph:simplegraph:caseson"]
  \inputleannode{InfoGeometry.Analysis.Graph.SimpleGraph.casesOn}
\end{theorem}

\subsection*{ctorIdx}

\begin{theorem}[blueprint="infogeometry:analysis:graph:simplegraph:ctoridx"]
  \inputleannode{InfoGeometry.Analysis.Graph.SimpleGraph.ctorIdx}
\end{theorem}

\subsection*{edges}

\begin{theorem}[blueprint="infogeometry:analysis:graph:simplegraph:edges"]
  \inputleannode{InfoGeometry.Analysis.Graph.SimpleGraph.edges}
\end{theorem}

\subsection*{mk}

\begin{theorem}[blueprint="infogeometry:analysis:graph:simplegraph:mk"]
  \inputleannode{InfoGeometry.Analysis.Graph.SimpleGraph.mk}
\end{theorem}

\subsection*{noConfusion}

\begin{theorem}[blueprint="infogeometry:analysis:graph:simplegraph:noconfusion"]
  \inputleannode{InfoGeometry.Analysis.Graph.SimpleGraph.noConfusion}
\end{theorem}

\subsection*{noConfusionType}

\begin{theorem}[blueprint="infogeometry:analysis:graph:simplegraph:noconfusiontype"]
  \inputleannode{InfoGeometry.Analysis.Graph.SimpleGraph.noConfusionType}
\end{theorem}

\subsection*{nodes}

\begin{theorem}[blueprint="infogeometry:analysis:graph:simplegraph:nodes"]
  \inputleannode{InfoGeometry.Analysis.Graph.SimpleGraph.nodes}
\end{theorem}

\subsection*{rec}

\begin{theorem}[blueprint="infogeometry:analysis:graph:simplegraph:rec"]
  \inputleannode{InfoGeometry.Analysis.Graph.SimpleGraph.rec}
\end{theorem}

\subsection*{recOn}

\begin{theorem}[blueprint="infogeometry:analysis:graph:simplegraph:recon"]
  \inputleannode{InfoGeometry.Analysis.Graph.SimpleGraph.recOn}
\end{theorem}

\section{InfoGeometry.Analysis.Graph.collectConsts}

\subsection*{proj}

\begin{theorem}[blueprint="infogeometry:analysis:graph:collectconsts:proj"]
  \inputleannode{InfoGeometry.Analysis.Graph.collectConsts.proj}
\end{theorem}

\section{InfoGeometry.Analysis.Graph.instBEqEdgeKind}

\subsection*{beq}

\begin{theorem}[blueprint="infogeometry:analysis:graph:instbeqedgekind:beq"]
  \inputleannode{InfoGeometry.Analysis.Graph.instBEqEdgeKind.beq}
\end{theorem}

\section{InfoGeometry.Analysis.Graph.instHashableEdgeKind}

\subsection*{hash}

\begin{theorem}[blueprint="infogeometry:analysis:graph:insthashableedgekind:hash"]
  \inputleannode{InfoGeometry.Analysis.Graph.instHashableEdgeKind.hash}
\end{theorem}

\section{InfoGeometry.Analysis.Graph.instReprEdgeKind}

\subsection*{repr}

\begin{theorem}[blueprint="infogeometry:analysis:graph:instrepredgekind:repr"]
  \inputleannode{InfoGeometry.Analysis.Graph.instReprEdgeKind.repr}
\end{theorem}

\chapter{Analytic Modules}

\section{InfoGeometry.Analytic}

\subsection*{deriv2_logSumExp_eq_softmaxVariance}

\begin{theorem}[blueprint="infogeometry:analytic:deriv2_logsumexp_eq_softmaxvariance"]
  \inputleannode{InfoGeometry.Analytic.deriv2_logSumExp_eq_softmaxVariance}
\end{theorem}

\subsection*{deriv_firstMomentUnnormalized}

\begin{theorem}[blueprint="infogeometry:analytic:deriv_firstmomentunnormalized"]
  \inputleannode{InfoGeometry.Analytic.deriv_firstMomentUnnormalized}
\end{theorem}

\subsection*{deriv_logSumExpMoment1}

\begin{theorem}[blueprint="infogeometry:analytic:deriv_logsumexpmoment1"]
  \inputleannode{InfoGeometry.Analytic.deriv_logSumExpMoment1}
\end{theorem}

\subsection*{deriv_logSumExpPartition}

\begin{theorem}[blueprint="infogeometry:analytic:deriv_logsumexppartition"]
  \inputleannode{InfoGeometry.Analytic.deriv_logSumExpPartition}
\end{theorem}

\subsection*{deriv_logSumExp_eq_firstMoment_div_partition}

\begin{theorem}[blueprint="infogeometry:analytic:deriv_logsumexp_eq_firstmoment_div_partition"]
  \inputleannode{InfoGeometry.Analytic.deriv_logSumExp_eq_firstMoment_div_partition}
\end{theorem}

\subsection*{deriv_logSumExp_eq_softmaxMean}

\begin{theorem}[blueprint="infogeometry:analytic:deriv_logsumexp_eq_softmaxmean"]
  \inputleannode{InfoGeometry.Analytic.deriv_logSumExp_eq_softmaxMean}
\end{theorem}

\subsection*{deriv_softmaxPartition}

\begin{theorem}[blueprint="infogeometry:analytic:deriv_softmaxpartition"]
  \inputleannode{InfoGeometry.Analytic.deriv_softmaxPartition}
\end{theorem}

\subsection*{firstMomentUnnormalized}

\begin{theorem}[blueprint="infogeometry:analytic:firstmomentunnormalized"]
  \inputleannode{InfoGeometry.Analytic.firstMomentUnnormalized}
\end{theorem}

\subsection*{hasDerivAt_firstMomentUnnormalized}

\begin{theorem}[blueprint="infogeometry:analytic:hasderivat_firstmomentunnormalized"]
  \inputleannode{InfoGeometry.Analytic.hasDerivAt_firstMomentUnnormalized}
\end{theorem}

\subsection*{hasDerivAt_logSumExpMoment1}

\begin{theorem}[blueprint="infogeometry:analytic:hasderivat_logsumexpmoment1"]
  \inputleannode{InfoGeometry.Analytic.hasDerivAt_logSumExpMoment1}
\end{theorem}

\subsection*{hasDerivAt_logSumExpPartition}

\begin{theorem}[blueprint="infogeometry:analytic:hasderivat_logsumexppartition"]
  \inputleannode{InfoGeometry.Analytic.hasDerivAt_logSumExpPartition}
\end{theorem}

\subsection*{hasDerivAt_logSumExpPartition_term}

\begin{theorem}[blueprint="infogeometry:analytic:hasderivat_logsumexppartition_term"]
  \inputleannode{InfoGeometry.Analytic.hasDerivAt_logSumExpPartition_term}
\end{theorem}

\subsection*{hasDerivAt_softmaxPartition}

\begin{theorem}[blueprint="infogeometry:analytic:hasderivat_softmaxpartition"]
  \inputleannode{InfoGeometry.Analytic.hasDerivAt_softmaxPartition}
\end{theorem}

\subsection*{logSumExp}

\begin{theorem}[blueprint="infogeometry:analytic:logsumexp"]
  \inputleannode{InfoGeometry.Analytic.logSumExp}
\end{theorem}

\subsection*{logSumExpMean}

\begin{theorem}[blueprint="infogeometry:analytic:logsumexpmean"]
  \inputleannode{InfoGeometry.Analytic.logSumExpMean}
\end{theorem}

\subsection*{logSumExpMean_eq_weighted_sum}

\begin{theorem}[blueprint="infogeometry:analytic:logsumexpmean_eq_weighted_sum"]
  \inputleannode{InfoGeometry.Analytic.logSumExpMean_eq_weighted_sum}
\end{theorem}

\subsection*{logSumExpMoment1}

\begin{theorem}[blueprint="infogeometry:analytic:logsumexpmoment1"]
  \inputleannode{InfoGeometry.Analytic.logSumExpMoment1}
\end{theorem}

\subsection*{logSumExpMoment2}

\begin{theorem}[blueprint="infogeometry:analytic:logsumexpmoment2"]
  \inputleannode{InfoGeometry.Analytic.logSumExpMoment2}
\end{theorem}

\subsection*{logSumExpPartition}

\begin{theorem}[blueprint="infogeometry:analytic:logsumexppartition"]
  \inputleannode{InfoGeometry.Analytic.logSumExpPartition}
\end{theorem}

\subsection*{logSumExpScaled}

\begin{theorem}[blueprint="infogeometry:analytic:logsumexpscaled"]
  \inputleannode{InfoGeometry.Analytic.logSumExpScaled}
\end{theorem}

\subsection*{logSumExpScaledBregman}

\begin{theorem}[blueprint="infogeometry:analytic:logsumexpscaledbregman"]
  \inputleannode{InfoGeometry.Analytic.logSumExpScaledBregman}
\end{theorem}

\subsection*{logSumExpScaledBregman_eq_eps_mul_KL}

\begin{theorem}[blueprint="infogeometry:analytic:logsumexpscaledbregman_eq_eps_mul_kl"]
  \inputleannode{InfoGeometry.Analytic.logSumExpScaledBregman_eq_eps_mul_KL}
\end{theorem}

\subsection*{logSumExpScaledEntropicTransportObjective}

\begin{theorem}[blueprint="infogeometry:analytic:logsumexpscaledentropictransportobjective"]
  \inputleannode{InfoGeometry.Analytic.logSumExpScaledEntropicTransportObjective}
\end{theorem}

\subsection*{logSumExpScaledEntropicTransportObjective_eq}

\begin{theorem}[blueprint="infogeometry:analytic:logsumexpscaledentropictransportobjective_eq"]
  \inputleannode{InfoGeometry.Analytic.logSumExpScaledEntropicTransportObjective_eq}
\end{theorem}

\subsection*{logSumExpScaledEntropicTransportObjective_eq_inv_eps_mul_potentialGap}

\begin{theorem}[blueprint="infogeometry:analytic:logsumexpscaledentropictransportobjective_eq_inv_eps_mul_potentialgap"]
  \inputleannode{InfoGeometry.Analytic.logSumExpScaledEntropicTransportObjective_eq_inv_eps_mul_potentialGap}
\end{theorem}

\subsection*{logSumExpScaledEntropicTransportPotentialGap}

\begin{theorem}[blueprint="infogeometry:analytic:logsumexpscaledentropictransportpotentialgap"]
  \inputleannode{InfoGeometry.Analytic.logSumExpScaledEntropicTransportPotentialGap}
\end{theorem}

\subsection*{logSumExpScaledEntropicTransportPotentialGap_eq}

\begin{theorem}[blueprint="infogeometry:analytic:logsumexpscaledentropictransportpotentialgap_eq"]
  \inputleannode{InfoGeometry.Analytic.logSumExpScaledEntropicTransportPotentialGap_eq}
\end{theorem}

\subsection*{logSumExpScaledEntropicTransportPotentialGap_eq_eps_mul_objective}

\begin{theorem}[blueprint="infogeometry:analytic:logsumexpscaledentropictransportpotentialgap_eq_eps_mul_objective"]
  \inputleannode{InfoGeometry.Analytic.logSumExpScaledEntropicTransportPotentialGap_eq_eps_mul_objective}
\end{theorem}

\subsection*{logSumExpScaledKL}

\begin{theorem}[blueprint="infogeometry:analytic:logsumexpscaledkl"]
  \inputleannode{InfoGeometry.Analytic.logSumExpScaledKL}
\end{theorem}

\subsection*{logSumExpScaledKL_eq}

\begin{theorem}[blueprint="infogeometry:analytic:logsumexpscaledkl_eq"]
  \inputleannode{InfoGeometry.Analytic.logSumExpScaledKL_eq}
\end{theorem}

\subsection*{logSumExpScaledKL_eq_inv_eps_mul_bregman}

\begin{theorem}[blueprint="infogeometry:analytic:logsumexpscaledkl_eq_inv_eps_mul_bregman"]
  \inputleannode{InfoGeometry.Analytic.logSumExpScaledKL_eq_inv_eps_mul_bregman}
\end{theorem}

\subsection*{logSumExpScaledKL_eq_inv_eps_mul_gap}

\begin{theorem}[blueprint="infogeometry:analytic:logsumexpscaledkl_eq_inv_eps_mul_gap"]
  \inputleannode{InfoGeometry.Analytic.logSumExpScaledKL_eq_inv_eps_mul_gap}
\end{theorem}

\subsection*{logSumExpScaledMean}

\begin{theorem}[blueprint="infogeometry:analytic:logsumexpscaledmean"]
  \inputleannode{InfoGeometry.Analytic.logSumExpScaledMean}
\end{theorem}

\subsection*{logSumExpScaledMean_eq_eps_mul_mean}

\begin{theorem}[blueprint="infogeometry:analytic:logsumexpscaledmean_eq_eps_mul_mean"]
  \inputleannode{InfoGeometry.Analytic.logSumExpScaledMean_eq_eps_mul_mean}
\end{theorem}

\subsection*{logSumExpScaledPartition}

\begin{theorem}[blueprint="infogeometry:analytic:logsumexpscaledpartition"]
  \inputleannode{InfoGeometry.Analytic.logSumExpScaledPartition}
\end{theorem}

\subsection*{logSumExpScaledPartition_pos}

\begin{theorem}[blueprint="infogeometry:analytic:logsumexpscaledpartition_pos"]
  \inputleannode{InfoGeometry.Analytic.logSumExpScaledPartition_pos}
\end{theorem}

\subsection*{logSumExpScaledWeight}

\begin{theorem}[blueprint="infogeometry:analytic:logsumexpscaledweight"]
  \inputleannode{InfoGeometry.Analytic.logSumExpScaledWeight}
\end{theorem}

\subsection*{logSumExpScaledWeight_pos}

\begin{theorem}[blueprint="infogeometry:analytic:logsumexpscaledweight_pos"]
  \inputleannode{InfoGeometry.Analytic.logSumExpScaledWeight_pos}
\end{theorem}

\subsection*{logSumExpScaledWeight_sum_one}

\begin{theorem}[blueprint="infogeometry:analytic:logsumexpscaledweight_sum_one"]
  \inputleannode{InfoGeometry.Analytic.logSumExpScaledWeight_sum_one}
\end{theorem}

\subsection*{logSumExpScaled_deriv_eq_mean}

\begin{theorem}[blueprint="infogeometry:analytic:logsumexpscaled_deriv_eq_mean"]
  \inputleannode{InfoGeometry.Analytic.logSumExpScaled_deriv_eq_mean}
\end{theorem}

\subsection*{logSumExpScaled_logWeight_eq}

\begin{theorem}[blueprint="infogeometry:analytic:logsumexpscaled_logweight_eq"]
  \inputleannode{InfoGeometry.Analytic.logSumExpScaled_logWeight_eq}
\end{theorem}

\subsection*{logSumExpSecondMoment}

\begin{theorem}[blueprint="infogeometry:analytic:logsumexpsecondmoment"]
  \inputleannode{InfoGeometry.Analytic.logSumExpSecondMoment}
\end{theorem}

\subsection*{logSumExpSecondMoment_eq_weighted_sum}

\begin{theorem}[blueprint="infogeometry:analytic:logsumexpsecondmoment_eq_weighted_sum"]
  \inputleannode{InfoGeometry.Analytic.logSumExpSecondMoment_eq_weighted_sum}
\end{theorem}

\subsection*{logSumExpVariance}

\begin{theorem}[blueprint="infogeometry:analytic:logsumexpvariance"]
  \inputleannode{InfoGeometry.Analytic.logSumExpVariance}
\end{theorem}

\subsection*{logSumExpWeight}

\begin{theorem}[blueprint="infogeometry:analytic:logsumexpweight"]
  \inputleannode{InfoGeometry.Analytic.logSumExpWeight}
\end{theorem}

\subsection*{logSumExpWeight_pos}

\begin{theorem}[blueprint="infogeometry:analytic:logsumexpweight_pos"]
  \inputleannode{InfoGeometry.Analytic.logSumExpWeight_pos}
\end{theorem}

\subsection*{logSumExpWeight_sum_one}

\begin{theorem}[blueprint="infogeometry:analytic:logsumexpweight_sum_one"]
  \inputleannode{InfoGeometry.Analytic.logSumExpWeight_sum_one}
\end{theorem}

\subsection*{logSumExp_contDiff}

\begin{theorem}[blueprint="infogeometry:analytic:logsumexp_contdiff"]
  \inputleannode{InfoGeometry.Analytic.logSumExp_contDiff}
\end{theorem}

\subsection*{logSumExp_deriv_eq_mean}

\begin{theorem}[blueprint="infogeometry:analytic:logsumexp_deriv_eq_mean"]
  \inputleannode{InfoGeometry.Analytic.logSumExp_deriv_eq_mean}
\end{theorem}

\subsection*{logSumExp_deriv_eq_ratio}

\begin{theorem}[blueprint="infogeometry:analytic:logsumexp_deriv_eq_ratio"]
  \inputleannode{InfoGeometry.Analytic.logSumExp_deriv_eq_ratio}
\end{theorem}

\subsection*{logSumExp_eq_log_partition}

\begin{theorem}[blueprint="infogeometry:analytic:logsumexp_eq_log_partition"]
  \inputleannode{InfoGeometry.Analytic.logSumExp_eq_log_partition}
\end{theorem}

\subsection*{logSumExp_secondDeriv_eq_ratio}

\begin{theorem}[blueprint="infogeometry:analytic:logsumexp_secondderiv_eq_ratio"]
  \inputleannode{InfoGeometry.Analytic.logSumExp_secondDeriv_eq_ratio}
\end{theorem}

\subsection*{logSumExp_secondDeriv_eq_variance}

\begin{theorem}[blueprint="infogeometry:analytic:logsumexp_secondderiv_eq_variance"]
  \inputleannode{InfoGeometry.Analytic.logSumExp_secondDeriv_eq_variance}
\end{theorem}

\subsection*{logSumExp_sum_pos}

\begin{theorem}[blueprint="infogeometry:analytic:logsumexp_sum_pos"]
  \inputleannode{InfoGeometry.Analytic.logSumExp_sum_pos}
\end{theorem}

\subsection*{secondMomentUnnormalized}

\begin{theorem}[blueprint="infogeometry:analytic:secondmomentunnormalized"]
  \inputleannode{InfoGeometry.Analytic.secondMomentUnnormalized}
\end{theorem}

\subsection*{softmaxDist}

\begin{theorem}[blueprint="infogeometry:analytic:softmaxdist"]
  \inputleannode{InfoGeometry.Analytic.softmaxDist}
\end{theorem}

\subsection*{softmaxDist_prob_eq}

\begin{theorem}[blueprint="infogeometry:analytic:softmaxdist_prob_eq"]
  \inputleannode{InfoGeometry.Analytic.softmaxDist_prob_eq}
\end{theorem}

\subsection*{softmaxHessian}

\begin{theorem}[blueprint="infogeometry:analytic:softmaxhessian"]
  \inputleannode{InfoGeometry.Analytic.softmaxHessian}
\end{theorem}

\subsection*{softmaxHessian_eq_secondMoment_sub_mean_sq}

\begin{theorem}[blueprint="infogeometry:analytic:softmaxhessian_eq_secondmoment_sub_mean_sq"]
  \inputleannode{InfoGeometry.Analytic.softmaxHessian_eq_secondMoment_sub_mean_sq}
\end{theorem}

\subsection*{softmaxMean}

\begin{theorem}[blueprint="infogeometry:analytic:softmaxmean"]
  \inputleannode{InfoGeometry.Analytic.softmaxMean}
\end{theorem}

\subsection*{softmaxMean_eq_firstMoment_div_partition}

\begin{theorem}[blueprint="infogeometry:analytic:softmaxmean_eq_firstmoment_div_partition"]
  \inputleannode{InfoGeometry.Analytic.softmaxMean_eq_firstMoment_div_partition}
\end{theorem}

\subsection*{softmaxMean_eq_logSumExpMean}

\begin{theorem}[blueprint="infogeometry:analytic:softmaxmean_eq_logsumexpmean"]
  \inputleannode{InfoGeometry.Analytic.softmaxMean_eq_logSumExpMean}
\end{theorem}

\subsection*{softmaxPartition}

\begin{theorem}[blueprint="infogeometry:analytic:softmaxpartition"]
  \inputleannode{InfoGeometry.Analytic.softmaxPartition}
\end{theorem}

\subsection*{softmaxPartition_pos}

\begin{theorem}[blueprint="infogeometry:analytic:softmaxpartition_pos"]
  \inputleannode{InfoGeometry.Analytic.softmaxPartition_pos}
\end{theorem}

\subsection*{softmaxProb}

\begin{theorem}[blueprint="infogeometry:analytic:softmaxprob"]
  \inputleannode{InfoGeometry.Analytic.softmaxProb}
\end{theorem}

\subsection*{softmaxProb_pos}

\begin{theorem}[blueprint="infogeometry:analytic:softmaxprob_pos"]
  \inputleannode{InfoGeometry.Analytic.softmaxProb_pos}
\end{theorem}

\subsection*{softmaxProb_sum_one}

\begin{theorem}[blueprint="infogeometry:analytic:softmaxprob_sum_one"]
  \inputleannode{InfoGeometry.Analytic.softmaxProb_sum_one}
\end{theorem}

\subsection*{softmaxSecondMoment}

\begin{theorem}[blueprint="infogeometry:analytic:softmaxsecondmoment"]
  \inputleannode{InfoGeometry.Analytic.softmaxSecondMoment}
\end{theorem}

\subsection*{softmaxSecondMoment_eq_logSumExpSecondMoment}

\begin{theorem}[blueprint="infogeometry:analytic:softmaxsecondmoment_eq_logsumexpsecondmoment"]
  \inputleannode{InfoGeometry.Analytic.softmaxSecondMoment_eq_logSumExpSecondMoment}
\end{theorem}

\subsection*{softmaxSecondMoment_eq_secondMomentUnnormalized_div_partition}

\begin{theorem}[blueprint="infogeometry:analytic:softmaxsecondmoment_eq_secondmomentunnormalized_div_partition"]
  \inputleannode{InfoGeometry.Analytic.softmaxSecondMoment_eq_secondMomentUnnormalized_div_partition}
\end{theorem}

\subsection*{softmaxVariance}

\begin{theorem}[blueprint="infogeometry:analytic:softmaxvariance"]
  \inputleannode{InfoGeometry.Analytic.softmaxVariance}
\end{theorem}

\subsection*{softmaxVariance_eq_secondMoment_sub_mean_sq}

\begin{theorem}[blueprint="infogeometry:analytic:softmaxvariance_eq_secondmoment_sub_mean_sq"]
  \inputleannode{InfoGeometry.Analytic.softmaxVariance_eq_secondMoment_sub_mean_sq}
\end{theorem}

\subsection*{softmaxVariance_nonneg}

\begin{theorem}[blueprint="infogeometry:analytic:softmaxvariance_nonneg"]
  \inputleannode{InfoGeometry.Analytic.softmaxVariance_nonneg}
\end{theorem}

\section{InfoGeometry.Analytic.logSumExpMean}

\subsection*{eq_1}

\begin{theorem}[blueprint="infogeometry:analytic:logsumexpmean:eq_1"]
  \inputleannode{InfoGeometry.Analytic.logSumExpMean.eq_1}
\end{theorem}

\section{InfoGeometry.Analytic.logSumExpPartition}

\subsection*{eq_1}

\begin{theorem}[blueprint="infogeometry:analytic:logsumexppartition:eq_1"]
  \inputleannode{InfoGeometry.Analytic.logSumExpPartition.eq_1}
\end{theorem}

\section{InfoGeometry.Analytic.logSumExpScaledPartition}

\subsection*{eq_1}

\begin{theorem}[blueprint="infogeometry:analytic:logsumexpscaledpartition:eq_1"]
  \inputleannode{InfoGeometry.Analytic.logSumExpScaledPartition.eq_1}
\end{theorem}

\section{InfoGeometry.Analytic.logSumExpScaledWeight}

\subsection*{eq_1}

\begin{theorem}[blueprint="infogeometry:analytic:logsumexpscaledweight:eq_1"]
  \inputleannode{InfoGeometry.Analytic.logSumExpScaledWeight.eq_1}
\end{theorem}

\section{InfoGeometry.Analytic.logSumExpSecondMoment}

\subsection*{eq_1}

\begin{theorem}[blueprint="infogeometry:analytic:logsumexpsecondmoment:eq_1"]
  \inputleannode{InfoGeometry.Analytic.logSumExpSecondMoment.eq_1}
\end{theorem}

\section{InfoGeometry.Analytic.logSumExpVariance}

\subsection*{eq_1}

\begin{theorem}[blueprint="infogeometry:analytic:logsumexpvariance:eq_1"]
  \inputleannode{InfoGeometry.Analytic.logSumExpVariance.eq_1}
\end{theorem}

\section{InfoGeometry.Analytic.softmaxDist}

\subsection*{congr_simp}

\begin{theorem}[blueprint="infogeometry:analytic:softmaxdist:congr_simp"]
  \inputleannode{InfoGeometry.Analytic.softmaxDist.congr_simp}
\end{theorem}

\section{InfoGeometry.Analytic.softmaxHessian}

\subsection*{eq_1}

\begin{theorem}[blueprint="infogeometry:analytic:softmaxhessian:eq_1"]
  \inputleannode{InfoGeometry.Analytic.softmaxHessian.eq_1}
\end{theorem}

\section{InfoGeometry.Analytic.softmaxVariance}

\subsection*{congr_simp}

\begin{theorem}[blueprint="infogeometry:analytic:softmaxvariance:congr_simp"]
  \inputleannode{InfoGeometry.Analytic.softmaxVariance.congr_simp}
\end{theorem}

\chapter{Architecture Modules}

\section{InfoGeometry.Architecture}

\subsection*{CartanInvolution}

\begin{theorem}[blueprint="infogeometry:architecture:cartaninvolution"]
  \inputleannode{InfoGeometry.Architecture.CartanInvolution}
\end{theorem}

\subsection*{InvolutiveMulAut}

\begin{theorem}[blueprint="infogeometry:architecture:involutivemulaut"]
  \inputleannode{InfoGeometry.Architecture.InvolutiveMulAut}
\end{theorem}

\subsection*{SymmetricPair}

\begin{theorem}[blueprint="infogeometry:architecture:symmetricpair"]
  \inputleannode{InfoGeometry.Architecture.SymmetricPair}
\end{theorem}

\subsection*{SymmetricSpace}

\begin{theorem}[blueprint="infogeometry:architecture:symmetricspace"]
  \inputleannode{InfoGeometry.Architecture.SymmetricSpace}
\end{theorem}

\subsection*{cartanSymmetry}

\begin{theorem}[blueprint="infogeometry:architecture:cartansymmetry"]
  \inputleannode{InfoGeometry.Architecture.cartanSymmetry}
\end{theorem}

\subsection*{cartanSymmetryOfInvolutiveMulAut}

\begin{theorem}[blueprint="infogeometry:architecture:cartansymmetryofinvolutivemulaut"]
  \inputleannode{InfoGeometry.Architecture.cartanSymmetryOfInvolutiveMulAut}
\end{theorem}

\subsection*{cartanSymmetryOfInvolutiveMulAut_fixpoint}

\begin{theorem}[blueprint="infogeometry:architecture:cartansymmetryofinvolutivemulaut_fixpoint"]
  \inputleannode{InfoGeometry.Architecture.cartanSymmetryOfInvolutiveMulAut_fixpoint}
\end{theorem}

\subsection*{cartanSymmetryOfInvolutiveMulAut_involutive}

\begin{theorem}[blueprint="infogeometry:architecture:cartansymmetryofinvolutivemulaut_involutive"]
  \inputleannode{InfoGeometry.Architecture.cartanSymmetryOfInvolutiveMulAut_involutive}
\end{theorem}

\subsection*{cartanSymmetry_fixpoint}

\begin{theorem}[blueprint="infogeometry:architecture:cartansymmetry_fixpoint"]
  \inputleannode{InfoGeometry.Architecture.cartanSymmetry_fixpoint}
\end{theorem}

\subsection*{cartanSymmetry_involutive}

\begin{theorem}[blueprint="infogeometry:architecture:cartansymmetry_involutive"]
  \inputleannode{InfoGeometry.Architecture.cartanSymmetry_involutive}
\end{theorem}

\subsection*{fixedSubgroup}

\begin{theorem}[blueprint="infogeometry:architecture:fixedsubgroup"]
  \inputleannode{InfoGeometry.Architecture.fixedSubgroup}
\end{theorem}

\subsection*{mem_fixedSubgroup_iff}

\begin{theorem}[blueprint="infogeometry:architecture:mem_fixedsubgroup_iff"]
  \inputleannode{InfoGeometry.Architecture.mem_fixedSubgroup_iff}
\end{theorem}

\subsection*{symmetricPairOfInvolution}

\begin{theorem}[blueprint="infogeometry:architecture:symmetricpairofinvolution"]
  \inputleannode{InfoGeometry.Architecture.symmetricPairOfInvolution}
\end{theorem}

\subsection*{symmetricPairOfInvolutiveMulAut}

\begin{theorem}[blueprint="infogeometry:architecture:symmetricpairofinvolutivemulaut"]
  \inputleannode{InfoGeometry.Architecture.symmetricPairOfInvolutiveMulAut}
\end{theorem}

\subsection*{symmetricSpaceOfCartan}

\begin{theorem}[blueprint="infogeometry:architecture:symmetricspaceofcartan"]
  \inputleannode{InfoGeometry.Architecture.symmetricSpaceOfCartan}
\end{theorem}

\subsection*{symmetricSpaceOfInvolutiveMulAut}

\begin{theorem}[blueprint="infogeometry:architecture:symmetricspaceofinvolutivemulaut"]
  \inputleannode{InfoGeometry.Architecture.symmetricSpaceOfInvolutiveMulAut}
\end{theorem}

\subsection*{symmetry_self}

\begin{theorem}[blueprint="infogeometry:architecture:symmetry_self"]
  \inputleannode{InfoGeometry.Architecture.symmetry_self}
\end{theorem}

\subsection*{symmetry_symmetry}

\begin{theorem}[blueprint="infogeometry:architecture:symmetry_symmetry"]
  \inputleannode{InfoGeometry.Architecture.symmetry_symmetry}
\end{theorem}

\section{InfoGeometry.Architecture.CartanInvolution}

\subsection*{casesOn}

\begin{theorem}[blueprint="infogeometry:architecture:cartaninvolution:caseson"]
  \inputleannode{InfoGeometry.Architecture.CartanInvolution.casesOn}
\end{theorem}

\subsection*{ctorIdx}

\begin{theorem}[blueprint="infogeometry:architecture:cartaninvolution:ctoridx"]
  \inputleannode{InfoGeometry.Architecture.CartanInvolution.ctorIdx}
\end{theorem}

\subsection*{fixedSubgroup}

\begin{theorem}[blueprint="infogeometry:architecture:cartaninvolution:fixedsubgroup"]
  \inputleannode{InfoGeometry.Architecture.CartanInvolution.fixedSubgroup}
\end{theorem}

\subsection*{involutive}

\begin{theorem}[blueprint="infogeometry:architecture:cartaninvolution:involutive"]
  \inputleannode{InfoGeometry.Architecture.CartanInvolution.involutive}
\end{theorem}

\subsection*{mem_fixedSubgroup_iff}

\begin{theorem}[blueprint="infogeometry:architecture:cartaninvolution:mem_fixedsubgroup_iff"]
  \inputleannode{InfoGeometry.Architecture.CartanInvolution.mem_fixedSubgroup_iff}
\end{theorem}

\subsection*{mk}

\begin{theorem}[blueprint="infogeometry:architecture:cartaninvolution:mk"]
  \inputleannode{InfoGeometry.Architecture.CartanInvolution.mk}
\end{theorem}

\subsection*{noConfusion}

\begin{theorem}[blueprint="infogeometry:architecture:cartaninvolution:noconfusion"]
  \inputleannode{InfoGeometry.Architecture.CartanInvolution.noConfusion}
\end{theorem}

\subsection*{noConfusionType}

\begin{theorem}[blueprint="infogeometry:architecture:cartaninvolution:noconfusiontype"]
  \inputleannode{InfoGeometry.Architecture.CartanInvolution.noConfusionType}
\end{theorem}

\subsection*{ofMulAutInvolution}

\begin{theorem}[blueprint="infogeometry:architecture:cartaninvolution:ofmulautinvolution"]
  \inputleannode{InfoGeometry.Architecture.CartanInvolution.ofMulAutInvolution}
\end{theorem}

\subsection*{rec}

\begin{theorem}[blueprint="infogeometry:architecture:cartaninvolution:rec"]
  \inputleannode{InfoGeometry.Architecture.CartanInvolution.rec}
\end{theorem}

\subsection*{recOn}

\begin{theorem}[blueprint="infogeometry:architecture:cartaninvolution:recon"]
  \inputleannode{InfoGeometry.Architecture.CartanInvolution.recOn}
\end{theorem}

\subsection*{toInvolutiveMulAut}

\begin{theorem}[blueprint="infogeometry:architecture:cartaninvolution:toinvolutivemulaut"]
  \inputleannode{InfoGeometry.Architecture.CartanInvolution.toInvolutiveMulAut}
\end{theorem}

\subsection*{toMulAut}

\begin{theorem}[blueprint="infogeometry:architecture:cartaninvolution:tomulaut"]
  \inputleannode{InfoGeometry.Architecture.CartanInvolution.toMulAut}
\end{theorem}

\section{InfoGeometry.Architecture.CartanInvolution.fixedSubgroup}

\subsection*{eq_1}

\begin{theorem}[blueprint="infogeometry:architecture:cartaninvolution:fixedsubgroup:eq_1"]
  \inputleannode{InfoGeometry.Architecture.CartanInvolution.fixedSubgroup.eq_1}
\end{theorem}

\section{InfoGeometry.Architecture.CartanInvolution.ofMulAutInvolution}

\subsection*{congr_simp}

\begin{theorem}[blueprint="infogeometry:architecture:cartaninvolution:ofmulautinvolution:congr_simp"]
  \inputleannode{InfoGeometry.Architecture.CartanInvolution.ofMulAutInvolution.congr_simp}
\end{theorem}

\section{InfoGeometry.Architecture.SymmetricPair}

\subsection*{K}

\begin{theorem}[blueprint="infogeometry:architecture:symmetricpair:k"]
  \inputleannode{InfoGeometry.Architecture.SymmetricPair.K}
\end{theorem}

\subsection*{K_eq_fixedSubgroup}

\begin{theorem}[blueprint="infogeometry:architecture:symmetricpair:k_eq_fixedsubgroup"]
  \inputleannode{InfoGeometry.Architecture.SymmetricPair.K_eq_fixedSubgroup}
\end{theorem}

\subsection*{casesOn}

\begin{theorem}[blueprint="infogeometry:architecture:symmetricpair:caseson"]
  \inputleannode{InfoGeometry.Architecture.SymmetricPair.casesOn}
\end{theorem}

\subsection*{ctorIdx}

\begin{theorem}[blueprint="infogeometry:architecture:symmetricpair:ctoridx"]
  \inputleannode{InfoGeometry.Architecture.SymmetricPair.ctorIdx}
\end{theorem}

\subsection*{fix_eq}

\begin{theorem}[blueprint="infogeometry:architecture:symmetricpair:fix_eq"]
  \inputleannode{InfoGeometry.Architecture.SymmetricPair.fix_eq}
\end{theorem}

\subsection*{mk}

\begin{theorem}[blueprint="infogeometry:architecture:symmetricpair:mk"]
  \inputleannode{InfoGeometry.Architecture.SymmetricPair.mk}
\end{theorem}

\subsection*{noConfusion}

\begin{theorem}[blueprint="infogeometry:architecture:symmetricpair:noconfusion"]
  \inputleannode{InfoGeometry.Architecture.SymmetricPair.noConfusion}
\end{theorem}

\subsection*{noConfusionType}

\begin{theorem}[blueprint="infogeometry:architecture:symmetricpair:noconfusiontype"]
  \inputleannode{InfoGeometry.Architecture.SymmetricPair.noConfusionType}
\end{theorem}

\subsection*{rec}

\begin{theorem}[blueprint="infogeometry:architecture:symmetricpair:rec"]
  \inputleannode{InfoGeometry.Architecture.SymmetricPair.rec}
\end{theorem}

\subsection*{recOn}

\begin{theorem}[blueprint="infogeometry:architecture:symmetricpair:recon"]
  \inputleannode{InfoGeometry.Architecture.SymmetricPair.recOn}
\end{theorem}

\subsection*{θ}

\begin{theorem}[blueprint="infogeometry:architecture:symmetricpair:θ"]
  \inputleannode{InfoGeometry.Architecture.SymmetricPair.θ}
\end{theorem}

\section{InfoGeometry.Architecture.SymmetricSpace}

\noindent\textbf{Buildable}: yes

\subsection*{casesOn}

\begin{theorem}[blueprint="infogeometry:architecture:symmetricspace:caseson"]
  \inputleannode{InfoGeometry.Architecture.SymmetricSpace.casesOn}
\end{theorem}

\subsection*{ctorIdx}

\begin{theorem}[blueprint="infogeometry:architecture:symmetricspace:ctoridx"]
  \inputleannode{InfoGeometry.Architecture.SymmetricSpace.ctorIdx}
\end{theorem}

\subsection*{mk}

\begin{theorem}[blueprint="infogeometry:architecture:symmetricspace:mk"]
  \inputleannode{InfoGeometry.Architecture.SymmetricSpace.mk}
\end{theorem}

\subsection*{noConfusion}

\begin{theorem}[blueprint="infogeometry:architecture:symmetricspace:noconfusion"]
  \inputleannode{InfoGeometry.Architecture.SymmetricSpace.noConfusion}
\end{theorem}

\subsection*{noConfusionType}

\begin{theorem}[blueprint="infogeometry:architecture:symmetricspace:noconfusiontype"]
  \inputleannode{InfoGeometry.Architecture.SymmetricSpace.noConfusionType}
\end{theorem}

\subsection*{rec}

\begin{theorem}[blueprint="infogeometry:architecture:symmetricspace:rec"]
  \inputleannode{InfoGeometry.Architecture.SymmetricSpace.rec}
\end{theorem}

\subsection*{recOn}

\begin{theorem}[blueprint="infogeometry:architecture:symmetricspace:recon"]
  \inputleannode{InfoGeometry.Architecture.SymmetricSpace.recOn}
\end{theorem}

\subsection*{symm_fixpoint}

\begin{theorem}[blueprint="infogeometry:architecture:symmetricspace:symm_fixpoint"]
  \inputleannode{InfoGeometry.Architecture.SymmetricSpace.symm_fixpoint}
\end{theorem}

\subsection*{symm_involutive}

\begin{theorem}[blueprint="infogeometry:architecture:symmetricspace:symm_involutive"]
  \inputleannode{InfoGeometry.Architecture.SymmetricSpace.symm_involutive}
\end{theorem}

\subsection*{symmetry}

\begin{theorem}[blueprint="infogeometry:architecture:symmetricspace:symmetry"]
  \inputleannode{InfoGeometry.Architecture.SymmetricSpace.symmetry}
\end{theorem}

\section{InfoGeometry.Architecture.cartanSymmetryOfInvolutiveMulAut}

\subsection*{eq_1}

\begin{theorem}[blueprint="infogeometry:architecture:cartansymmetryofinvolutivemulaut:eq_1"]
  \inputleannode{InfoGeometry.Architecture.cartanSymmetryOfInvolutiveMulAut.eq_1}
\end{theorem}

\chapter{Archive Modules}

\section{InfoGeometry.Archive.Drafts.Degree_pre_assumptions_extraction}

\noindent\textbf{Buildable}: no

\noindent\emph{No exported declarations in the current declaration graph.}

\section{InfoGeometry.Archive.Drafts.New_pre_assumptions_extraction}

\noindent\textbf{Buildable}: no

\noindent\emph{No exported declarations in the current declaration graph.}

\section{InfoGeometry.Archive.Drafts.Projective_LogSum_pre_assumptions_extraction}

\noindent\textbf{Buildable}: no

\noindent\emph{No exported declarations in the current declaration graph.}

\chapter{Assumptions Modules}

\section{InfoGeometry.Assumptions}

\noindent\textbf{Buildable}: yes

\noindent\emph{No exported declarations in the current declaration graph.}

\section{InfoGeometry.Assumptions.Determinant}

\noindent\textbf{Buildable}: yes

\subsection*{GL}

\begin{theorem}[blueprint="infogeometry:assumptions:determinant:gl"]
  \inputleannode{InfoGeometry.Assumptions.Determinant.GL}
\end{theorem}

\subsection*{SL}

\begin{theorem}[blueprint="infogeometry:assumptions:determinant:sl"]
  \inputleannode{InfoGeometry.Assumptions.Determinant.SL}
\end{theorem}

\subsection*{detHom}

\begin{theorem}[blueprint="infogeometry:assumptions:determinant:dethom"]
  \inputleannode{InfoGeometry.Assumptions.Determinant.detHom}
\end{theorem}

\subsection*{detMonoidHom}

\begin{theorem}[blueprint="infogeometry:assumptions:determinant:detmonoidhom"]
  \inputleannode{InfoGeometry.Assumptions.Determinant.detMonoidHom}
\end{theorem}

\subsection*{jacDet}

\begin{theorem}[blueprint="infogeometry:assumptions:determinant:jacdet"]
  \inputleannode{InfoGeometry.Assumptions.Determinant.jacDet}
\end{theorem}

\subsection*{jacDet_comp}

\begin{theorem}[blueprint="infogeometry:assumptions:determinant:jacdet_comp"]
  \inputleannode{InfoGeometry.Assumptions.Determinant.jacDet_comp}
\end{theorem}

\subsection*{jacobian_functoriality}

\begin{theorem}[blueprint="infogeometry:assumptions:determinant:jacobian_functoriality"]
  \inputleannode{InfoGeometry.Assumptions.Determinant.jacobian_functoriality}
\end{theorem}

\subsection*{ker_det_eq_SL}

\begin{theorem}[blueprint="infogeometry:assumptions:determinant:ker_det_eq_sl"]
  \inputleannode{InfoGeometry.Assumptions.Determinant.ker_det_eq_SL}
\end{theorem}

\subsection*{logAbsDet}

\begin{theorem}[blueprint="infogeometry:assumptions:determinant:logabsdet"]
  \inputleannode{InfoGeometry.Assumptions.Determinant.logAbsDet}
\end{theorem}

\section{InfoGeometry.Assumptions.DualConnections}

\noindent\textbf{Buildable}: yes

\subsection*{alphaConnection}

\begin{theorem}[blueprint="infogeometry:assumptions:dualconnections:alphaconnection"]
  \inputleannode{InfoGeometry.Assumptions.DualConnections.alphaConnection}
\end{theorem}

\subsection*{alpha_duality}

\begin{theorem}[blueprint="infogeometry:assumptions:dualconnections:alpha_duality"]
  \inputleannode{InfoGeometry.Assumptions.DualConnections.alpha_duality}
\end{theorem}

\subsection*{amariChentsovTensor}

\begin{theorem}[blueprint="infogeometry:assumptions:dualconnections:amarichentsovtensor"]
  \inputleannode{InfoGeometry.Assumptions.DualConnections.amariChentsovTensor}
\end{theorem}

\subsection*{amariChentsovTensor_of_finProb}

\begin{theorem}[blueprint="infogeometry:assumptions:dualconnections:amarichentsovtensor_of_finprob"]
  \inputleannode{InfoGeometry.Assumptions.DualConnections.amariChentsovTensor_of_finProb}
\end{theorem}

\subsection*{fisherMetric}

\begin{theorem}[blueprint="infogeometry:assumptions:dualconnections:fishermetric"]
  \inputleannode{InfoGeometry.Assumptions.DualConnections.fisherMetric}
\end{theorem}

\subsection*{fisherMetric_of_finProb}

\begin{theorem}[blueprint="infogeometry:assumptions:dualconnections:fishermetric_of_finprob"]
  \inputleannode{InfoGeometry.Assumptions.DualConnections.fisherMetric_of_finProb}
\end{theorem}

\subsection*{fisher_metric_eq_hessian_KL}

\begin{theorem}[blueprint="infogeometry:assumptions:dualconnections:fisher_metric_eq_hessian_kl"]
  \inputleannode{InfoGeometry.Assumptions.DualConnections.fisher_metric_eq_hessian_KL}
\end{theorem}

\section{InfoGeometry.Assumptions.DualConnections.alphaConnection}

\subsection*{eq_1}

\begin{theorem}[blueprint="infogeometry:assumptions:dualconnections:alphaconnection:eq_1"]
  \inputleannode{InfoGeometry.Assumptions.DualConnections.alphaConnection.eq_1}
\end{theorem}

\section{InfoGeometry.Assumptions.IB}

\noindent\textbf{Buildable}: yes

\subsection*{FinProb}

\begin{theorem}[blueprint="infogeometry:assumptions:ib:finprob"]
  \inputleannode{InfoGeometry.Assumptions.IB.FinProb}
\end{theorem}

\subsection*{IBProblem}

\begin{theorem}[blueprint="infogeometry:assumptions:ib:ibproblem"]
  \inputleannode{InfoGeometry.Assumptions.IB.IBProblem}
\end{theorem}

\subsection*{KLKernel}

\begin{theorem}[blueprint="infogeometry:assumptions:ib:klkernel"]
  \inputleannode{InfoGeometry.Assumptions.IB.KLKernel}
\end{theorem}

\subsection*{argmin_exponentialTilt}

\begin{theorem}[blueprint="infogeometry:assumptions:ib:argmin_exponentialtilt"]
  \inputleannode{InfoGeometry.Assumptions.IB.argmin_exponentialTilt}
\end{theorem}

\subsection*{argmin_exponentialTilt_true}

\begin{theorem}[blueprint="infogeometry:assumptions:ib:argmin_exponentialtilt_true"]
  \inputleannode{InfoGeometry.Assumptions.IB.argmin_exponentialTilt_true}
\end{theorem}

\subsection*{energy}

\begin{theorem}[blueprint="infogeometry:assumptions:ib:energy"]
  \inputleannode{InfoGeometry.Assumptions.IB.energy}
\end{theorem}

\subsection*{exponentialTilt}

\begin{theorem}[blueprint="infogeometry:assumptions:ib:exponentialtilt"]
  \inputleannode{InfoGeometry.Assumptions.IB.exponentialTilt}
\end{theorem}

\subsection*{ibIteration}

\begin{theorem}[blueprint="infogeometry:assumptions:ib:ibiteration"]
  \inputleannode{InfoGeometry.Assumptions.IB.ibIteration}
\end{theorem}

\subsection*{ibLagrangian}

\begin{theorem}[blueprint="infogeometry:assumptions:ib:iblagrangian"]
  \inputleannode{InfoGeometry.Assumptions.IB.ibLagrangian}
\end{theorem}

\subsection*{ib_convergence}

\begin{theorem}[blueprint="infogeometry:assumptions:ib:ib_convergence"]
  \inputleannode{InfoGeometry.Assumptions.IB.ib_convergence}
\end{theorem}

\subsection*{ib_convergence_nonneg}

\begin{theorem}[blueprint="infogeometry:assumptions:ib:ib_convergence_nonneg"]
  \inputleannode{InfoGeometry.Assumptions.IB.ib_convergence_nonneg}
\end{theorem}

\subsection*{ib_convergence_true}

\begin{theorem}[blueprint="infogeometry:assumptions:ib:ib_convergence_true"]
  \inputleannode{InfoGeometry.Assumptions.IB.ib_convergence_true}
\end{theorem}

\subsection*{ib_stationary_point_gibbs}

\begin{theorem}[blueprint="infogeometry:assumptions:ib:ib_stationary_point_gibbs"]
  \inputleannode{InfoGeometry.Assumptions.IB.ib_stationary_point_gibbs}
\end{theorem}

\subsection*{ib_stationary_point_gibbs_true}

\begin{theorem}[blueprint="infogeometry:assumptions:ib:ib_stationary_point_gibbs_true"]
  \inputleannode{InfoGeometry.Assumptions.IB.ib_stationary_point_gibbs_true}
\end{theorem}

\subsection*{inducedMProjection}

\begin{theorem}[blueprint="infogeometry:assumptions:ib:inducedmprojection"]
  \inputleannode{InfoGeometry.Assumptions.IB.inducedMProjection}
\end{theorem}

\subsection*{jointYT}

\begin{theorem}[blueprint="infogeometry:assumptions:ib:jointyt"]
  \inputleannode{InfoGeometry.Assumptions.IB.jointYT}
\end{theorem}

\subsection*{mutualInformation}

\begin{theorem}[blueprint="infogeometry:assumptions:ib:mutualinformation"]
  \inputleannode{InfoGeometry.Assumptions.IB.mutualInformation}
\end{theorem}

\subsection*{partitionFunction}

\begin{theorem}[blueprint="infogeometry:assumptions:ib:partitionfunction"]
  \inputleannode{InfoGeometry.Assumptions.IB.partitionFunction}
\end{theorem}

\section{InfoGeometry.Assumptions.IB.IBProblem}

\subsection*{beta}

\begin{theorem}[blueprint="infogeometry:assumptions:ib:ibproblem:beta"]
  \inputleannode{InfoGeometry.Assumptions.IB.IBProblem.beta}
\end{theorem}

\subsection*{beta_pos}

\begin{theorem}[blueprint="infogeometry:assumptions:ib:ibproblem:beta_pos"]
  \inputleannode{InfoGeometry.Assumptions.IB.IBProblem.beta_pos}
\end{theorem}

\subsection*{casesOn}

\begin{theorem}[blueprint="infogeometry:assumptions:ib:ibproblem:caseson"]
  \inputleannode{InfoGeometry.Assumptions.IB.IBProblem.casesOn}
\end{theorem}

\subsection*{ctorIdx}

\begin{theorem}[blueprint="infogeometry:assumptions:ib:ibproblem:ctoridx"]
  \inputleannode{InfoGeometry.Assumptions.IB.IBProblem.ctorIdx}
\end{theorem}

\subsection*{mk}

\begin{theorem}[blueprint="infogeometry:assumptions:ib:ibproblem:mk"]
  \inputleannode{InfoGeometry.Assumptions.IB.IBProblem.mk}
\end{theorem}

\subsection*{noConfusion}

\begin{theorem}[blueprint="infogeometry:assumptions:ib:ibproblem:noconfusion"]
  \inputleannode{InfoGeometry.Assumptions.IB.IBProblem.noConfusion}
\end{theorem}

\subsection*{noConfusionType}

\begin{theorem}[blueprint="infogeometry:assumptions:ib:ibproblem:noconfusiontype"]
  \inputleannode{InfoGeometry.Assumptions.IB.IBProblem.noConfusionType}
\end{theorem}

\subsection*{pXY}

\begin{theorem}[blueprint="infogeometry:assumptions:ib:ibproblem:pxy"]
  \inputleannode{InfoGeometry.Assumptions.IB.IBProblem.pXY}
\end{theorem}

\subsection*{rec}

\begin{theorem}[blueprint="infogeometry:assumptions:ib:ibproblem:rec"]
  \inputleannode{InfoGeometry.Assumptions.IB.IBProblem.rec}
\end{theorem}

\subsection*{recOn}

\begin{theorem}[blueprint="infogeometry:assumptions:ib:ibproblem:recon"]
  \inputleannode{InfoGeometry.Assumptions.IB.IBProblem.recOn}
\end{theorem}

\section{InfoGeometry.Assumptions.LLN}

\noindent\textbf{Buildable}: yes

\subsection*{ae_tendsto_ratio_to_rnDeriv}

\begin{theorem}[blueprint="infogeometry:assumptions:lln:ae_tendsto_ratio_to_rnderiv"]
  \inputleannode{InfoGeometry.Assumptions.LLN.ae_tendsto_ratio_to_rnDeriv}
\end{theorem}

\subsection*{ae_tendsto_ratio_to_rnDeriv_holds}

\begin{theorem}[blueprint="infogeometry:assumptions:lln:ae_tendsto_ratio_to_rnderiv_holds"]
  \inputleannode{InfoGeometry.Assumptions.LLN.ae_tendsto_ratio_to_rnDeriv_holds}
\end{theorem}

\subsection*{empiricalAverage}

\begin{theorem}[blueprint="infogeometry:assumptions:lln:empiricalaverage"]
  \inputleannode{InfoGeometry.Assumptions.LLN.empiricalAverage}
\end{theorem}

\subsection*{empirical_to_theoretical_slln}

\begin{theorem}[blueprint="infogeometry:assumptions:lln:empirical_to_theoretical_slln"]
  \inputleannode{InfoGeometry.Assumptions.LLN.empirical_to_theoretical_slln}
\end{theorem}

\subsection*{fixed_partition_slln}

\begin{theorem}[blueprint="infogeometry:assumptions:lln:fixed_partition_slln"]
  \inputleannode{InfoGeometry.Assumptions.LLN.fixed_partition_slln}
\end{theorem}

\subsection*{fixed_partition_slln_holds}

\begin{theorem}[blueprint="infogeometry:assumptions:lln:fixed_partition_slln_holds"]
  \inputleannode{InfoGeometry.Assumptions.LLN.fixed_partition_slln_holds}
\end{theorem}

\section{InfoGeometry.Assumptions.LLN.empiricalAverage}

\subsection*{eq_1}

\begin{theorem}[blueprint="infogeometry:assumptions:lln:empiricalaverage:eq_1"]
  \inputleannode{InfoGeometry.Assumptions.LLN.empiricalAverage.eq_1}
\end{theorem}

\section{InfoGeometry.Assumptions.ManifoldDegree}

\noindent\textbf{Buildable}: yes

\subsection*{exists_isolating_nhds_of_nondegenerate}

\begin{theorem}[blueprint="infogeometry:assumptions:manifolddegree:exists_isolating_nhds_of_nondegenerate"]
  \inputleannode{InfoGeometry.Assumptions.ManifoldDegree.exists_isolating_nhds_of_nondegenerate}
\end{theorem}

\subsection*{exists_local_chart_homotopy_to_linear}

\begin{theorem}[blueprint="infogeometry:assumptions:manifolddegree:exists_local_chart_homotopy_to_linear"]
  \inputleannode{InfoGeometry.Assumptions.ManifoldDegree.exists_local_chart_homotopy_to_linear}
\end{theorem}

\subsection*{localDegreeSign}

\begin{theorem}[blueprint="infogeometry:assumptions:manifolddegree:localdegreesign"]
  \inputleannode{InfoGeometry.Assumptions.ManifoldDegree.localDegreeSign}
\end{theorem}

\subsection*{local_degree_eq_sign_jacDet}

\begin{theorem}[blueprint="infogeometry:assumptions:manifolddegree:local_degree_eq_sign_jacdet"]
  \inputleannode{InfoGeometry.Assumptions.ManifoldDegree.local_degree_eq_sign_jacDet}
\end{theorem}

\subsection*{preimage_finite_of_regular_value}

\begin{theorem}[blueprint="infogeometry:assumptions:manifolddegree:preimage_finite_of_regular_value"]
  \inputleannode{InfoGeometry.Assumptions.ManifoldDegree.preimage_finite_of_regular_value}
\end{theorem}

\section{InfoGeometry.Assumptions.ManifoldHomology}

\noindent\textbf{Buildable}: yes

\subsection*{IsRegularValue}

\begin{theorem}[blueprint="infogeometry:assumptions:manifoldhomology:isregularvalue"]
  \inputleannode{InfoGeometry.Assumptions.ManifoldHomology.IsRegularValue}
\end{theorem}

\subsection*{degree_formula_via_jacobian}

\begin{theorem}[blueprint="infogeometry:assumptions:manifoldhomology:degree_formula_via_jacobian"]
  \inputleannode{InfoGeometry.Assumptions.ManifoldHomology.degree_formula_via_jacobian}
\end{theorem}

\subsection*{degree_formula_via_jacobian_true}

\begin{theorem}[blueprint="infogeometry:assumptions:manifoldhomology:degree_formula_via_jacobian_true"]
  \inputleannode{InfoGeometry.Assumptions.ManifoldHomology.degree_formula_via_jacobian_true}
\end{theorem}

\subsection*{mappingDegree}

\begin{theorem}[blueprint="infogeometry:assumptions:manifoldhomology:mappingdegree"]
  \inputleannode{InfoGeometry.Assumptions.ManifoldHomology.mappingDegree}
\end{theorem}

\subsection*{mappingDegree_nonneg}

\begin{theorem}[blueprint="infogeometry:assumptions:manifoldhomology:mappingdegree_nonneg"]
  \inputleannode{InfoGeometry.Assumptions.ManifoldHomology.mappingDegree_nonneg}
\end{theorem}

\subsection*{topHomologyIso}

\begin{theorem}[blueprint="infogeometry:assumptions:manifoldhomology:tophomologyiso"]
  \inputleannode{InfoGeometry.Assumptions.ManifoldHomology.topHomologyIso}
\end{theorem}

\subsection*{top_homology_is_Z}

\begin{theorem}[blueprint="infogeometry:assumptions:manifoldhomology:top_homology_is_z"]
  \inputleannode{InfoGeometry.Assumptions.ManifoldHomology.top_homology_is_Z}
\end{theorem}

\subsection*{top_homology_is_Z_true}

\begin{theorem}[blueprint="infogeometry:assumptions:manifoldhomology:top_homology_is_z_true"]
  \inputleannode{InfoGeometry.Assumptions.ManifoldHomology.top_homology_is_Z_true}
\end{theorem}

\section{InfoGeometry.Assumptions.ManifoldHomology.mappingDegree}

\subsection*{eq_1}

\begin{theorem}[blueprint="infogeometry:assumptions:manifoldhomology:mappingdegree:eq_1"]
  \inputleannode{InfoGeometry.Assumptions.ManifoldHomology.mappingDegree.eq_1}
\end{theorem}

\section{InfoGeometry.Assumptions.Projective}

\subsection*{logSum_inequality}

\begin{theorem}[blueprint="infogeometry:assumptions:projective:logsum_inequality"]
  \inputleannode{InfoGeometry.Assumptions.Projective.logSum_inequality}
\end{theorem}

\chapter{Axioms Modules}

\section{InfoGeometry.Axioms}

\noindent\textbf{Buildable}: yes

\subsection*{GL}

\begin{theorem}[blueprint="infogeometry:axioms:gl"]
  \inputleannode{InfoGeometry.Axioms.GL}
\end{theorem}

\subsection*{SL}

\begin{theorem}[blueprint="infogeometry:axioms:sl"]
  \inputleannode{InfoGeometry.Axioms.SL}
\end{theorem}

\chapter{Basic Modules}

\section{InfoGeometry.Basic}

\noindent\textbf{Buildable}: yes

\noindent\emph{No exported declarations in the current declaration graph.}

\chapter{Clifford Modules}

\section{InfoGeometry.Clifford}

\subsection*{IsLightlike}

\begin{theorem}[blueprint="infogeometry:clifford:islightlike"]
  \inputleannode{InfoGeometry.Clifford.IsLightlike}
\end{theorem}

\subsection*{IsSpacelike}

\begin{theorem}[blueprint="infogeometry:clifford:isspacelike"]
  \inputleannode{InfoGeometry.Clifford.IsSpacelike}
\end{theorem}

\subsection*{IsTimelike}

\begin{theorem}[blueprint="infogeometry:clifford:istimelike"]
  \inputleannode{InfoGeometry.Clifford.IsTimelike}
\end{theorem}

\subsection*{splitB11}

\begin{theorem}[blueprint="infogeometry:clifford:splitb11"]
  \inputleannode{InfoGeometry.Clifford.splitB11}
\end{theorem}

\subsection*{splitB11_apply}

\begin{theorem}[blueprint="infogeometry:clifford:splitb11_apply"]
  \inputleannode{InfoGeometry.Clifford.splitB11_apply}
\end{theorem}

\subsection*{splitB11_expand}

\begin{theorem}[blueprint="infogeometry:clifford:splitb11_expand"]
  \inputleannode{InfoGeometry.Clifford.splitB11_expand}
\end{theorem}

\subsection*{splitQ11}

\begin{theorem}[blueprint="infogeometry:clifford:splitq11"]
  \inputleannode{InfoGeometry.Clifford.splitQ11}
\end{theorem}

\subsection*{splitQ11_add}

\begin{theorem}[blueprint="infogeometry:clifford:splitq11_add"]
  \inputleannode{InfoGeometry.Clifford.splitQ11_add}
\end{theorem}

\subsection*{splitQ11_apply}

\begin{theorem}[blueprint="infogeometry:clifford:splitq11_apply"]
  \inputleannode{InfoGeometry.Clifford.splitQ11_apply}
\end{theorem}

\subsection*{splitQ11_eq_splitB11_diag}

\begin{theorem}[blueprint="infogeometry:clifford:splitq11_eq_splitb11_diag"]
  \inputleannode{InfoGeometry.Clifford.splitQ11_eq_splitB11_diag}
\end{theorem}

\section{InfoGeometry.Clifford.Cl11}

\noindent\textbf{Buildable}: yes

\noindent\emph{No exported declarations in the current declaration graph.}

\section{InfoGeometry.Clifford.Grading}

\noindent\textbf{Buildable}: yes

\noindent\emph{No exported declarations in the current declaration graph.}

\section{InfoGeometry.Clifford.Lift}

\noindent\textbf{Buildable}: yes

\noindent\emph{No exported declarations in the current declaration graph.}

\section{InfoGeometry.Clifford.Relations}

\noindent\textbf{Buildable}: yes

\noindent\emph{No exported declarations in the current declaration graph.}

\section{InfoGeometry.Clifford.SplitQ11}

\noindent\textbf{Buildable}: yes

\noindent\emph{No exported declarations in the current declaration graph.}

\section{InfoGeometry.Clifford.Supercharge}

\noindent\textbf{Buildable}: yes

\noindent\emph{No exported declarations in the current declaration graph.}

\section{InfoGeometry.Clifford.Tower}

\noindent\textbf{Buildable}: yes

\noindent\emph{No exported declarations in the current declaration graph.}

\section{InfoGeometry.Clifford.splitQ11}

\subsection*{eq_1}

\begin{theorem}[blueprint="infogeometry:clifford:splitq11:eq_1"]
  \inputleannode{InfoGeometry.Clifford.splitQ11.eq_1}
\end{theorem}

\chapter{CliffordTower Modules}

\section{InfoGeometry.CliffordTower}

\subsection*{Clsplit}

\begin{theorem}[blueprint="infogeometry:cliffordtower:clsplit"]
  \inputleannode{InfoGeometry.CliffordTower.Clsplit}
\end{theorem}

\subsection*{Q11}

\begin{theorem}[blueprint="infogeometry:cliffordtower:q11"]
  \inputleannode{InfoGeometry.CliffordTower.Q11}
\end{theorem}

\subsection*{Q11_apply}

\begin{theorem}[blueprint="infogeometry:cliffordtower:q11_apply"]
  \inputleannode{InfoGeometry.CliffordTower.Q11_apply}
\end{theorem}

\subsection*{Qsplit}

\begin{theorem}[blueprint="infogeometry:cliffordtower:qsplit"]
  \inputleannode{InfoGeometry.CliffordTower.Qsplit}
\end{theorem}

\subsection*{Qsplit_succ_apply}

\begin{theorem}[blueprint="infogeometry:cliffordtower:qsplit_succ_apply"]
  \inputleannode{InfoGeometry.CliffordTower.Qsplit_succ_apply}
\end{theorem}

\subsection*{Qsplit_zero_apply}

\begin{theorem}[blueprint="infogeometry:cliffordtower:qsplit_zero_apply"]
  \inputleannode{InfoGeometry.CliffordTower.Qsplit_zero_apply}
\end{theorem}

\subsection*{SplitSpace}

\begin{theorem}[blueprint="infogeometry:cliffordtower:splitspace"]
  \inputleannode{InfoGeometry.CliffordTower.SplitSpace}
\end{theorem}

\subsection*{clsplit_succ_equiv}

\begin{theorem}[blueprint="infogeometry:cliffordtower:clsplit_succ_equiv"]
  \inputleannode{InfoGeometry.CliffordTower.clsplit_succ_equiv}
\end{theorem}

\subsection*{splitSpaceAddCommGroup}

\begin{theorem}[blueprint="infogeometry:cliffordtower:splitspaceaddcommgroup"]
  \inputleannode{InfoGeometry.CliffordTower.splitSpaceAddCommGroup}
\end{theorem}

\subsection*{splitSpaceModule}

\begin{theorem}[blueprint="infogeometry:cliffordtower:splitspacemodule"]
  \inputleannode{InfoGeometry.CliffordTower.splitSpaceModule}
\end{theorem}

\section{InfoGeometry.CliffordTower.Qsplit}

\subsection*{eq_1}

\begin{theorem}[blueprint="infogeometry:cliffordtower:qsplit:eq_1"]
  \inputleannode{InfoGeometry.CliffordTower.Qsplit.eq_1}
\end{theorem}

\subsection*{eq_2}

\begin{theorem}[blueprint="infogeometry:cliffordtower:qsplit:eq_2"]
  \inputleannode{InfoGeometry.CliffordTower.Qsplit.eq_2}
\end{theorem}

\subsection*{eq_def}

\begin{theorem}[blueprint="infogeometry:cliffordtower:qsplit:eq_def"]
  \inputleannode{InfoGeometry.CliffordTower.Qsplit.eq_def}
\end{theorem}

\chapter{Convex Modules}

\section{InfoGeometry.Convex}

\noindent\textbf{Buildable}: yes

\subsection*{BregmanDivergence}

\begin{theorem}[blueprint="infogeometry:convex:bregmandivergence"]
  \inputleannode{InfoGeometry.Convex.BregmanDivergence}
\end{theorem}

\subsection*{ConvexFunctional}

\begin{theorem}[blueprint="infogeometry:convex:convexfunctional"]
  \inputleannode{InfoGeometry.Convex.ConvexFunctional}
\end{theorem}

\subsection*{DualSpace}

\begin{theorem}[blueprint="infogeometry:convex:dualspace"]
  \inputleannode{InfoGeometry.Convex.DualSpace}
\end{theorem}

\subsection*{HessianGeometry}

\begin{theorem}[blueprint="infogeometry:convex:hessiangeometry"]
  \inputleannode{InfoGeometry.Convex.HessianGeometry}
\end{theorem}

\subsection*{HessianGeometry1D}

\begin{theorem}[blueprint="infogeometry:convex:hessiangeometry1d"]
  \inputleannode{InfoGeometry.Convex.HessianGeometry1D}
\end{theorem}

\subsection*{LegendrePotential}

\begin{theorem}[blueprint="infogeometry:convex:legendrepotential"]
  \inputleannode{InfoGeometry.Convex.LegendrePotential}
\end{theorem}

\subsection*{NonzeroDoubledState}

\begin{theorem}[blueprint="infogeometry:convex:nonzerodoubledstate"]
  \inputleannode{InfoGeometry.Convex.NonzeroDoubledState}
\end{theorem}

\subsection*{ProjectiveState}

\begin{theorem}[blueprint="infogeometry:convex:projectivestate"]
  \inputleannode{InfoGeometry.Convex.ProjectiveState}
\end{theorem}

\subsection*{bregmanTriadicCore}

\begin{theorem}[blueprint="infogeometry:convex:bregmantriadiccore"]
  \inputleannode{InfoGeometry.Convex.bregmanTriadicCore}
\end{theorem}

\subsection*{dualPair}

\begin{theorem}[blueprint="infogeometry:convex:dualpair"]
  \inputleannode{InfoGeometry.Convex.dualPair}
\end{theorem}

\subsection*{fenchelConj}

\begin{theorem}[blueprint="infogeometry:convex:fenchelconj"]
  \inputleannode{InfoGeometry.Convex.fenchelConj}
\end{theorem}

\subsection*{fenchelConj_eq_of_isGreatest}

\begin{theorem}[blueprint="infogeometry:convex:fenchelconj_eq_of_isgreatest"]
  \inputleannode{InfoGeometry.Convex.fenchelConj_eq_of_isGreatest}
\end{theorem}

\subsection*{fenchelConj_ge_eval_sub}

\begin{theorem}[blueprint="infogeometry:convex:fenchelconj_ge_eval_sub"]
  \inputleannode{InfoGeometry.Convex.fenchelConj_ge_eval_sub}
\end{theorem}

\subsection*{fenchelSet}

\begin{theorem}[blueprint="infogeometry:convex:fenchelset"]
  \inputleannode{InfoGeometry.Convex.fenchelSet}
\end{theorem}

\subsection*{fenchelSet_nonempty}

\begin{theorem}[blueprint="infogeometry:convex:fenchelset_nonempty"]
  \inputleannode{InfoGeometry.Convex.fenchelSet_nonempty}
\end{theorem}

\subsection*{fenchelYoung}

\begin{theorem}[blueprint="infogeometry:convex:fenchelyoung"]
  \inputleannode{InfoGeometry.Convex.fenchelYoung}
\end{theorem}

\subsection*{fenchelYoung_eq_of_conj_eq}

\begin{theorem}[blueprint="infogeometry:convex:fenchelyoung_eq_of_conj_eq"]
  \inputleannode{InfoGeometry.Convex.fenchelYoung_eq_of_conj_eq}
\end{theorem}

\subsection*{helly_theorem}

\begin{theorem}[blueprint="infogeometry:convex:helly_theorem"]
  \inputleannode{InfoGeometry.Convex.helly_theorem}
\end{theorem}

\subsection*{helly_theorem'}

\begin{theorem}[blueprint="infogeometry:convex:helly_theorem'"]
  \inputleannode{InfoGeometry.Convex.helly_theorem'}
\end{theorem}

\subsection*{helly_theorem_compact}

\begin{theorem}[blueprint="infogeometry:convex:helly_theorem_compact"]
  \inputleannode{InfoGeometry.Convex.helly_theorem_compact}
\end{theorem}

\subsection*{helly_theorem_compact'}

\begin{theorem}[blueprint="infogeometry:convex:helly_theorem_compact'"]
  \inputleannode{InfoGeometry.Convex.helly_theorem_compact'}
\end{theorem}

\subsection*{helly_theorem_set}

\begin{theorem}[blueprint="infogeometry:convex:helly_theorem_set"]
  \inputleannode{InfoGeometry.Convex.helly_theorem_set}
\end{theorem}

\subsection*{helly_theorem_set'}

\begin{theorem}[blueprint="infogeometry:convex:helly_theorem_set'"]
  \inputleannode{InfoGeometry.Convex.helly_theorem_set'}
\end{theorem}

\subsection*{helly_theorem_set_compact}

\begin{theorem}[blueprint="infogeometry:convex:helly_theorem_set_compact"]
  \inputleannode{InfoGeometry.Convex.helly_theorem_set_compact}
\end{theorem}

\subsection*{helly_theorem_set_compact'}

\begin{theorem}[blueprint="infogeometry:convex:helly_theorem_set_compact'"]
  \inputleannode{InfoGeometry.Convex.helly_theorem_set_compact'}
\end{theorem}

\subsection*{le_fenchelConj}

\begin{theorem}[blueprint="infogeometry:convex:le_fenchelconj"]
  \inputleannode{InfoGeometry.Convex.le_fenchelConj}
\end{theorem}

\subsection*{projectivize}

\begin{theorem}[blueprint="infogeometry:convex:projectivize"]
  \inputleannode{InfoGeometry.Convex.projectivize}
\end{theorem}

\subsection*{projectivize'}

\begin{theorem}[blueprint="infogeometry:convex:projectivize'"]
  \inputleannode{InfoGeometry.Convex.projectivize'}
\end{theorem}

\subsection*{projectivize'_eq_projectivize}

\begin{theorem}[blueprint="infogeometry:convex:projectivize'_eq_projectivize"]
  \inputleannode{InfoGeometry.Convex.projectivize'_eq_projectivize}
\end{theorem}

\subsection*{projectivize'_smul}

\begin{theorem}[blueprint="infogeometry:convex:projectivize'_smul"]
  \inputleannode{InfoGeometry.Convex.projectivize'_smul}
\end{theorem}

\subsection*{projectivize_smul}

\begin{theorem}[blueprint="infogeometry:convex:projectivize_smul"]
  \inputleannode{InfoGeometry.Convex.projectivize_smul}
\end{theorem}

\subsection*{radon_partition}

\begin{theorem}[blueprint="infogeometry:convex:radon_partition"]
  \inputleannode{InfoGeometry.Convex.radon_partition}
\end{theorem}

\section{InfoGeometry.Convex.Bregman}

\noindent\textbf{Buildable}: yes

\noindent\emph{No exported declarations in the current declaration graph.}

\section{InfoGeometry.Convex.BregmanDivergence}

\subsection*{D}

\begin{theorem}[blueprint="infogeometry:convex:bregmandivergence:d"]
  \inputleannode{InfoGeometry.Convex.BregmanDivergence.D}
\end{theorem}

\subsection*{casesOn}

\begin{theorem}[blueprint="infogeometry:convex:bregmandivergence:caseson"]
  \inputleannode{InfoGeometry.Convex.BregmanDivergence.casesOn}
\end{theorem}

\subsection*{ctorIdx}

\begin{theorem}[blueprint="infogeometry:convex:bregmandivergence:ctoridx"]
  \inputleannode{InfoGeometry.Convex.BregmanDivergence.ctorIdx}
\end{theorem}

\subsection*{mk}

\begin{theorem}[blueprint="infogeometry:convex:bregmandivergence:mk"]
  \inputleannode{InfoGeometry.Convex.BregmanDivergence.mk}
\end{theorem}

\subsection*{noConfusion}

\begin{theorem}[blueprint="infogeometry:convex:bregmandivergence:noconfusion"]
  \inputleannode{InfoGeometry.Convex.BregmanDivergence.noConfusion}
\end{theorem}

\subsection*{noConfusionType}

\begin{theorem}[blueprint="infogeometry:convex:bregmandivergence:noconfusiontype"]
  \inputleannode{InfoGeometry.Convex.BregmanDivergence.noConfusionType}
\end{theorem}

\subsection*{rec}

\begin{theorem}[blueprint="infogeometry:convex:bregmandivergence:rec"]
  \inputleannode{InfoGeometry.Convex.BregmanDivergence.rec}
\end{theorem}

\subsection*{recOn}

\begin{theorem}[blueprint="infogeometry:convex:bregmandivergence:recon"]
  \inputleannode{InfoGeometry.Convex.BregmanDivergence.recOn}
\end{theorem}

\section{InfoGeometry.Convex.ConvexFunctional}

\subsection*{F}

\begin{theorem}[blueprint="infogeometry:convex:convexfunctional:f"]
  \inputleannode{InfoGeometry.Convex.ConvexFunctional.F}
\end{theorem}

\subsection*{StrictConvex}

\begin{theorem}[blueprint="infogeometry:convex:convexfunctional:strictconvex"]
  \inputleannode{InfoGeometry.Convex.ConvexFunctional.StrictConvex}
\end{theorem}

\subsection*{StrictMonotoneGrad}

\begin{theorem}[blueprint="infogeometry:convex:convexfunctional:strictmonotonegrad"]
  \inputleannode{InfoGeometry.Convex.ConvexFunctional.StrictMonotoneGrad}
\end{theorem}

\subsection*{affineSet}

\begin{theorem}[blueprint="infogeometry:convex:convexfunctional:affineset"]
  \inputleannode{InfoGeometry.Convex.ConvexFunctional.affineSet}
\end{theorem}

\subsection*{affineSet_nonempty}

\begin{theorem}[blueprint="infogeometry:convex:convexfunctional:affineset_nonempty"]
  \inputleannode{InfoGeometry.Convex.ConvexFunctional.affineSet_nonempty}
\end{theorem}

\subsection*{casesOn}

\begin{theorem}[blueprint="infogeometry:convex:convexfunctional:caseson"]
  \inputleannode{InfoGeometry.Convex.ConvexFunctional.casesOn}
\end{theorem}

\subsection*{convex}

\begin{theorem}[blueprint="infogeometry:convex:convexfunctional:convex"]
  \inputleannode{InfoGeometry.Convex.ConvexFunctional.convex}
\end{theorem}

\subsection*{ctorIdx}

\begin{theorem}[blueprint="infogeometry:convex:convexfunctional:ctoridx"]
  \inputleannode{InfoGeometry.Convex.ConvexFunctional.ctorIdx}
\end{theorem}

\subsection*{diff}

\begin{theorem}[blueprint="infogeometry:convex:convexfunctional:diff"]
  \inputleannode{InfoGeometry.Convex.ConvexFunctional.diff}
\end{theorem}

\subsection*{fenchel_young}

\begin{theorem}[blueprint="infogeometry:convex:convexfunctional:fenchel_young"]
  \inputleannode{InfoGeometry.Convex.ConvexFunctional.fenchel_young}
\end{theorem}

\subsection*{fenchel_young_eq_of_grad}

\begin{theorem}[blueprint="infogeometry:convex:convexfunctional:fenchel_young_eq_of_grad"]
  \inputleannode{InfoGeometry.Convex.ConvexFunctional.fenchel_young_eq_of_grad}
\end{theorem}

\subsection*{fenchel_young_eq_of_supporting}

\begin{theorem}[blueprint="infogeometry:convex:convexfunctional:fenchel_young_eq_of_supporting"]
  \inputleannode{InfoGeometry.Convex.ConvexFunctional.fenchel_young_eq_of_supporting}
\end{theorem}

\subsection*{grad}

\begin{theorem}[blueprint="infogeometry:convex:convexfunctional:grad"]
  \inputleannode{InfoGeometry.Convex.ConvexFunctional.grad}
\end{theorem}

\subsection*{grad_injective_of_strictMonotone}

\begin{theorem}[blueprint="infogeometry:convex:convexfunctional:grad_injective_of_strictmonotone"]
  \inputleannode{InfoGeometry.Convex.ConvexFunctional.grad_injective_of_strictMonotone}
\end{theorem}

\subsection*{legendre}

\begin{theorem}[blueprint="infogeometry:convex:convexfunctional:legendre"]
  \inputleannode{InfoGeometry.Convex.ConvexFunctional.legendre}
\end{theorem}

\subsection*{mk}

\begin{theorem}[blueprint="infogeometry:convex:convexfunctional:mk"]
  \inputleannode{InfoGeometry.Convex.ConvexFunctional.mk}
\end{theorem}

\subsection*{noConfusion}

\begin{theorem}[blueprint="infogeometry:convex:convexfunctional:noconfusion"]
  \inputleannode{InfoGeometry.Convex.ConvexFunctional.noConfusion}
\end{theorem}

\subsection*{noConfusionType}

\begin{theorem}[blueprint="infogeometry:convex:convexfunctional:noconfusiontype"]
  \inputleannode{InfoGeometry.Convex.ConvexFunctional.noConfusionType}
\end{theorem}

\subsection*{rec}

\begin{theorem}[blueprint="infogeometry:convex:convexfunctional:rec"]
  \inputleannode{InfoGeometry.Convex.ConvexFunctional.rec}
\end{theorem}

\subsection*{recOn}

\begin{theorem}[blueprint="infogeometry:convex:convexfunctional:recon"]
  \inputleannode{InfoGeometry.Convex.ConvexFunctional.recOn}
\end{theorem}

\subsection*{supporting_ineq_of_convex_differentiable}

\begin{theorem}[blueprint="infogeometry:convex:convexfunctional:supporting_ineq_of_convex_differentiable"]
  \inputleannode{InfoGeometry.Convex.ConvexFunctional.supporting_ineq_of_convex_differentiable}
\end{theorem}

\section{InfoGeometry.Convex.ConvexFunctional.grad}

\subsection*{eq_1}

\begin{theorem}[blueprint="infogeometry:convex:convexfunctional:grad:eq_1"]
  \inputleannode{InfoGeometry.Convex.ConvexFunctional.grad.eq_1}
\end{theorem}

\section{InfoGeometry.Convex.Duality}

\noindent\textbf{Buildable}: yes

\noindent\emph{No exported declarations in the current declaration graph.}

\section{InfoGeometry.Convex.FenchelConjugate}

\noindent\textbf{Buildable}: yes

\noindent\emph{No exported declarations in the current declaration graph.}

\section{InfoGeometry.Convex.HessianGeometry}

\noindent\textbf{Buildable}: yes

\subsection*{bregmanDiv}

\begin{theorem}[blueprint="infogeometry:convex:hessiangeometry:bregmandiv"]
  \inputleannode{InfoGeometry.Convex.HessianGeometry.bregmanDiv}
\end{theorem}

\subsection*{bregmanDiv_D}

\begin{theorem}[blueprint="infogeometry:convex:hessiangeometry:bregmandiv_d"]
  \inputleannode{InfoGeometry.Convex.HessianGeometry.bregmanDiv_D}
\end{theorem}

\subsection*{casesOn}

\begin{theorem}[blueprint="infogeometry:convex:hessiangeometry:caseson"]
  \inputleannode{InfoGeometry.Convex.HessianGeometry.casesOn}
\end{theorem}

\subsection*{ctorIdx}

\begin{theorem}[blueprint="infogeometry:convex:hessiangeometry:ctoridx"]
  \inputleannode{InfoGeometry.Convex.HessianGeometry.ctorIdx}
\end{theorem}

\subsection*{divergence}

\begin{theorem}[blueprint="infogeometry:convex:hessiangeometry:divergence"]
  \inputleannode{InfoGeometry.Convex.HessianGeometry.divergence}
\end{theorem}

\subsection*{divergence_def}

\begin{theorem}[blueprint="infogeometry:convex:hessiangeometry:divergence_def"]
  \inputleannode{InfoGeometry.Convex.HessianGeometry.divergence_def}
\end{theorem}

\subsection*{dualMap}

\begin{theorem}[blueprint="infogeometry:convex:hessiangeometry:dualmap"]
  \inputleannode{InfoGeometry.Convex.HessianGeometry.dualMap}
\end{theorem}

\subsection*{dualMap_def}

\begin{theorem}[blueprint="infogeometry:convex:hessiangeometry:dualmap_def"]
  \inputleannode{InfoGeometry.Convex.HessianGeometry.dualMap_def}
\end{theorem}

\subsection*{dualPotential}

\begin{theorem}[blueprint="infogeometry:convex:hessiangeometry:dualpotential"]
  \inputleannode{InfoGeometry.Convex.HessianGeometry.dualPotential}
\end{theorem}

\subsection*{grad}

\begin{theorem}[blueprint="infogeometry:convex:hessiangeometry:grad"]
  \inputleannode{InfoGeometry.Convex.HessianGeometry.grad}
\end{theorem}

\subsection*{has_gradient}

\begin{theorem}[blueprint="infogeometry:convex:hessiangeometry:has_gradient"]
  \inputleannode{InfoGeometry.Convex.HessianGeometry.has_gradient}
\end{theorem}

\subsection*{metric}

\begin{theorem}[blueprint="infogeometry:convex:hessiangeometry:metric"]
  \inputleannode{InfoGeometry.Convex.HessianGeometry.metric}
\end{theorem}

\subsection*{metricOp}

\begin{theorem}[blueprint="infogeometry:convex:hessiangeometry:metricop"]
  \inputleannode{InfoGeometry.Convex.HessianGeometry.metricOp}
\end{theorem}

\subsection*{metric_def}

\begin{theorem}[blueprint="infogeometry:convex:hessiangeometry:metric_def"]
  \inputleannode{InfoGeometry.Convex.HessianGeometry.metric_def}
\end{theorem}

\subsection*{mk}

\begin{theorem}[blueprint="infogeometry:convex:hessiangeometry:mk"]
  \inputleannode{InfoGeometry.Convex.HessianGeometry.mk}
\end{theorem}

\subsection*{noConfusion}

\begin{theorem}[blueprint="infogeometry:convex:hessiangeometry:noconfusion"]
  \inputleannode{InfoGeometry.Convex.HessianGeometry.noConfusion}
\end{theorem}

\subsection*{noConfusionType}

\begin{theorem}[blueprint="infogeometry:convex:hessiangeometry:noconfusiontype"]
  \inputleannode{InfoGeometry.Convex.HessianGeometry.noConfusionType}
\end{theorem}

\subsection*{potential}

\begin{theorem}[blueprint="infogeometry:convex:hessiangeometry:potential"]
  \inputleannode{InfoGeometry.Convex.HessianGeometry.potential}
\end{theorem}

\subsection*{rec}

\begin{theorem}[blueprint="infogeometry:convex:hessiangeometry:rec"]
  \inputleannode{InfoGeometry.Convex.HessianGeometry.rec}
\end{theorem}

\subsection*{recOn}

\begin{theorem}[blueprint="infogeometry:convex:hessiangeometry:recon"]
  \inputleannode{InfoGeometry.Convex.HessianGeometry.recOn}
\end{theorem}

\subsection*{softmaxBregmanAttention}

\begin{theorem}[blueprint="infogeometry:convex:hessiangeometry:softmaxbregmanattention"]
  \inputleannode{InfoGeometry.Convex.HessianGeometry.softmaxBregmanAttention}
\end{theorem}

\section{InfoGeometry.Convex.HessianGeometry1D}

\subsection*{bregmanDiv}

\begin{theorem}[blueprint="infogeometry:convex:hessiangeometry1d:bregmandiv"]
  \inputleannode{InfoGeometry.Convex.HessianGeometry1D.bregmanDiv}
\end{theorem}

\subsection*{bregmanDiv_D}

\begin{theorem}[blueprint="infogeometry:convex:hessiangeometry1d:bregmandiv_d"]
  \inputleannode{InfoGeometry.Convex.HessianGeometry1D.bregmanDiv_D}
\end{theorem}

\subsection*{casesOn}

\begin{theorem}[blueprint="infogeometry:convex:hessiangeometry1d:caseson"]
  \inputleannode{InfoGeometry.Convex.HessianGeometry1D.casesOn}
\end{theorem}

\subsection*{ctorIdx}

\begin{theorem}[blueprint="infogeometry:convex:hessiangeometry1d:ctoridx"]
  \inputleannode{InfoGeometry.Convex.HessianGeometry1D.ctorIdx}
\end{theorem}

\subsection*{divergence}

\begin{theorem}[blueprint="infogeometry:convex:hessiangeometry1d:divergence"]
  \inputleannode{InfoGeometry.Convex.HessianGeometry1D.divergence}
\end{theorem}

\subsection*{divergence_def}

\begin{theorem}[blueprint="infogeometry:convex:hessiangeometry1d:divergence_def"]
  \inputleannode{InfoGeometry.Convex.HessianGeometry1D.divergence_def}
\end{theorem}

\subsection*{dualMap}

\begin{theorem}[blueprint="infogeometry:convex:hessiangeometry1d:dualmap"]
  \inputleannode{InfoGeometry.Convex.HessianGeometry1D.dualMap}
\end{theorem}

\subsection*{dualMap_def}

\begin{theorem}[blueprint="infogeometry:convex:hessiangeometry1d:dualmap_def"]
  \inputleannode{InfoGeometry.Convex.HessianGeometry1D.dualMap_def}
\end{theorem}

\subsection*{dualPotential}

\begin{theorem}[blueprint="infogeometry:convex:hessiangeometry1d:dualpotential"]
  \inputleannode{InfoGeometry.Convex.HessianGeometry1D.dualPotential}
\end{theorem}

\subsection*{metric}

\begin{theorem}[blueprint="infogeometry:convex:hessiangeometry1d:metric"]
  \inputleannode{InfoGeometry.Convex.HessianGeometry1D.metric}
\end{theorem}

\subsection*{metric_def}

\begin{theorem}[blueprint="infogeometry:convex:hessiangeometry1d:metric_def"]
  \inputleannode{InfoGeometry.Convex.HessianGeometry1D.metric_def}
\end{theorem}

\subsection*{mk}

\begin{theorem}[blueprint="infogeometry:convex:hessiangeometry1d:mk"]
  \inputleannode{InfoGeometry.Convex.HessianGeometry1D.mk}
\end{theorem}

\subsection*{noConfusion}

\begin{theorem}[blueprint="infogeometry:convex:hessiangeometry1d:noconfusion"]
  \inputleannode{InfoGeometry.Convex.HessianGeometry1D.noConfusion}
\end{theorem}

\subsection*{noConfusionType}

\begin{theorem}[blueprint="infogeometry:convex:hessiangeometry1d:noconfusiontype"]
  \inputleannode{InfoGeometry.Convex.HessianGeometry1D.noConfusionType}
\end{theorem}

\subsection*{ofLogPotential}

\begin{theorem}[blueprint="infogeometry:convex:hessiangeometry1d:oflogpotential"]
  \inputleannode{InfoGeometry.Convex.HessianGeometry1D.ofLogPotential}
\end{theorem}

\subsection*{potential}

\begin{theorem}[blueprint="infogeometry:convex:hessiangeometry1d:potential"]
  \inputleannode{InfoGeometry.Convex.HessianGeometry1D.potential}
\end{theorem}

\subsection*{rec}

\begin{theorem}[blueprint="infogeometry:convex:hessiangeometry1d:rec"]
  \inputleannode{InfoGeometry.Convex.HessianGeometry1D.rec}
\end{theorem}

\subsection*{recOn}

\begin{theorem}[blueprint="infogeometry:convex:hessiangeometry1d:recon"]
  \inputleannode{InfoGeometry.Convex.HessianGeometry1D.recOn}
\end{theorem}

\subsection*{softmaxBregmanAttention}

\begin{theorem}[blueprint="infogeometry:convex:hessiangeometry1d:softmaxbregmanattention"]
  \inputleannode{InfoGeometry.Convex.HessianGeometry1D.softmaxBregmanAttention}
\end{theorem}

\section{InfoGeometry.Convex.Legendre}

\noindent\textbf{Buildable}: yes

\noindent\emph{No exported declarations in the current declaration graph.}

\section{InfoGeometry.Convex.LegendrePotential}

\subsection*{bregman}

\begin{theorem}[blueprint="infogeometry:convex:legendrepotential:bregman"]
  \inputleannode{InfoGeometry.Convex.LegendrePotential.bregman}
\end{theorem}

\subsection*{bregman_def}

\begin{theorem}[blueprint="infogeometry:convex:legendrepotential:bregman_def"]
  \inputleannode{InfoGeometry.Convex.LegendrePotential.bregman_def}
\end{theorem}

\subsection*{casesOn}

\begin{theorem}[blueprint="infogeometry:convex:legendrepotential:caseson"]
  \inputleannode{InfoGeometry.Convex.LegendrePotential.casesOn}
\end{theorem}

\subsection*{ctorIdx}

\begin{theorem}[blueprint="infogeometry:convex:legendrepotential:ctoridx"]
  \inputleannode{InfoGeometry.Convex.LegendrePotential.ctorIdx}
\end{theorem}

\subsection*{deriv_strictMono}

\begin{theorem}[blueprint="infogeometry:convex:legendrepotential:deriv_strictmono"]
  \inputleannode{InfoGeometry.Convex.LegendrePotential.deriv_strictMono}
\end{theorem}

\subsection*{differentiable}

\begin{theorem}[blueprint="infogeometry:convex:legendrepotential:differentiable"]
  \inputleannode{InfoGeometry.Convex.LegendrePotential.differentiable}
\end{theorem}

\subsection*{differentiableOn}

\begin{theorem}[blueprint="infogeometry:convex:legendrepotential:differentiableon"]
  \inputleannode{InfoGeometry.Convex.LegendrePotential.differentiableOn}
\end{theorem}

\subsection*{divergence}

\begin{theorem}[blueprint="infogeometry:convex:legendrepotential:divergence"]
  \inputleannode{InfoGeometry.Convex.LegendrePotential.divergence}
\end{theorem}

\subsection*{eta}

\begin{theorem}[blueprint="infogeometry:convex:legendrepotential:eta"]
  \inputleannode{InfoGeometry.Convex.LegendrePotential.eta}
\end{theorem}

\subsection*{eta_injective}

\begin{theorem}[blueprint="infogeometry:convex:legendrepotential:eta_injective"]
  \inputleannode{InfoGeometry.Convex.LegendrePotential.eta_injective}
\end{theorem}

\subsection*{eta_strictMono}

\begin{theorem}[blueprint="infogeometry:convex:legendrepotential:eta_strictmono"]
  \inputleannode{InfoGeometry.Convex.LegendrePotential.eta_strictMono}
\end{theorem}

\subsection*{f}

\begin{theorem}[blueprint="infogeometry:convex:legendrepotential:f"]
  \inputleannode{InfoGeometry.Convex.LegendrePotential.f}
\end{theorem}

\subsection*{fenchel_young_eq_of_grad}

\begin{theorem}[blueprint="infogeometry:convex:legendrepotential:fenchel_young_eq_of_grad"]
  \inputleannode{InfoGeometry.Convex.LegendrePotential.fenchel_young_eq_of_grad}
\end{theorem}

\subsection*{fisher}

\begin{theorem}[blueprint="infogeometry:convex:legendrepotential:fisher"]
  \inputleannode{InfoGeometry.Convex.LegendrePotential.fisher}
\end{theorem}

\subsection*{fisher_nonneg}

\begin{theorem}[blueprint="infogeometry:convex:legendrepotential:fisher_nonneg"]
  \inputleannode{InfoGeometry.Convex.LegendrePotential.fisher_nonneg}
\end{theorem}

\subsection*{grad_injective}

\begin{theorem}[blueprint="infogeometry:convex:legendrepotential:grad_injective"]
  \inputleannode{InfoGeometry.Convex.LegendrePotential.grad_injective}
\end{theorem}

\subsection*{legendre}

\begin{theorem}[blueprint="infogeometry:convex:legendrepotential:legendre"]
  \inputleannode{InfoGeometry.Convex.LegendrePotential.legendre}
\end{theorem}

\subsection*{legendreTransform}

\begin{theorem}[blueprint="infogeometry:convex:legendrepotential:legendretransform"]
  \inputleannode{InfoGeometry.Convex.LegendrePotential.legendreTransform}
\end{theorem}

\subsection*{legendreTransform_eta}

\begin{theorem}[blueprint="infogeometry:convex:legendrepotential:legendretransform_eta"]
  \inputleannode{InfoGeometry.Convex.LegendrePotential.legendreTransform_eta}
\end{theorem}

\subsection*{mk}

\begin{theorem}[blueprint="infogeometry:convex:legendrepotential:mk"]
  \inputleannode{InfoGeometry.Convex.LegendrePotential.mk}
\end{theorem}

\subsection*{noConfusion}

\begin{theorem}[blueprint="infogeometry:convex:legendrepotential:noconfusion"]
  \inputleannode{InfoGeometry.Convex.LegendrePotential.noConfusion}
\end{theorem}

\subsection*{noConfusionType}

\begin{theorem}[blueprint="infogeometry:convex:legendrepotential:noconfusiontype"]
  \inputleannode{InfoGeometry.Convex.LegendrePotential.noConfusionType}
\end{theorem}

\subsection*{rec}

\begin{theorem}[blueprint="infogeometry:convex:legendrepotential:rec"]
  \inputleannode{InfoGeometry.Convex.LegendrePotential.rec}
\end{theorem}

\subsection*{recOn}

\begin{theorem}[blueprint="infogeometry:convex:legendrepotential:recon"]
  \inputleannode{InfoGeometry.Convex.LegendrePotential.recOn}
\end{theorem}

\subsection*{smooth}

\begin{theorem}[blueprint="infogeometry:convex:legendrepotential:smooth"]
  \inputleannode{InfoGeometry.Convex.LegendrePotential.smooth}
\end{theorem}

\subsection*{strict_convex}

\begin{theorem}[blueprint="infogeometry:convex:legendrepotential:strict_convex"]
  \inputleannode{InfoGeometry.Convex.LegendrePotential.strict_convex}
\end{theorem}

\subsection*{thetaOfEta}

\begin{theorem}[blueprint="infogeometry:convex:legendrepotential:thetaofeta"]
  \inputleannode{InfoGeometry.Convex.LegendrePotential.thetaOfEta}
\end{theorem}

\subsection*{thetaOfEta_eta}

\begin{theorem}[blueprint="infogeometry:convex:legendrepotential:thetaofeta_eta"]
  \inputleannode{InfoGeometry.Convex.LegendrePotential.thetaOfEta_eta}
\end{theorem}

\subsection*{theta_eta_leftInverse}

\begin{theorem}[blueprint="infogeometry:convex:legendrepotential:theta_eta_leftinverse"]
  \inputleannode{InfoGeometry.Convex.LegendrePotential.theta_eta_leftInverse}
\end{theorem}

\subsection*{toConvexFunctional}

\begin{theorem}[blueprint="infogeometry:convex:legendrepotential:toconvexfunctional"]
  \inputleannode{InfoGeometry.Convex.LegendrePotential.toConvexFunctional}
\end{theorem}

\section{InfoGeometry.Convex.LegendrePotential.legendreTransform}

\subsection*{eq_1}

\begin{theorem}[blueprint="infogeometry:convex:legendrepotential:legendretransform:eq_1"]
  \inputleannode{InfoGeometry.Convex.LegendrePotential.legendreTransform.eq_1}
\end{theorem}

\section{InfoGeometry.Convex.OneD}

\subsection*{DualFlatPotential}

\begin{theorem}[blueprint="infogeometry:convex:oned:dualflatpotential"]
  \inputleannode{InfoGeometry.Convex.OneD.DualFlatPotential}
\end{theorem}

\subsection*{IsFenchelUpperBound}

\begin{theorem}[blueprint="infogeometry:convex:oned:isfenchelupperbound"]
  \inputleannode{InfoGeometry.Convex.OneD.IsFenchelUpperBound}
\end{theorem}

\subsection*{IsGradientDual}

\begin{theorem}[blueprint="infogeometry:convex:oned:isgradientdual"]
  \inputleannode{InfoGeometry.Convex.OneD.IsGradientDual}
\end{theorem}

\subsection*{IsKernelClosed}

\begin{theorem}[blueprint="infogeometry:convex:oned:iskernelclosed"]
  \inputleannode{InfoGeometry.Convex.OneD.IsKernelClosed}
\end{theorem}

\subsection*{IsLegendreConjugate}

\begin{theorem}[blueprint="infogeometry:convex:oned:islegendreconjugate"]
  \inputleannode{InfoGeometry.Convex.OneD.IsLegendreConjugate}
\end{theorem}

\subsection*{KernelClosedFunctions}

\begin{theorem}[blueprint="infogeometry:convex:oned:kernelclosedfunctions"]
  \inputleannode{InfoGeometry.Convex.OneD.KernelClosedFunctions}
\end{theorem}

\subsection*{KernelClosure}

\begin{theorem}[blueprint="infogeometry:convex:oned:kernelclosure"]
  \inputleannode{InfoGeometry.Convex.OneD.KernelClosure}
\end{theorem}

\subsection*{KernelClosureOperator}

\begin{theorem}[blueprint="infogeometry:convex:oned:kernelclosureoperator"]
  \inputleannode{InfoGeometry.Convex.OneD.KernelClosureOperator}
\end{theorem}

\subsection*{KernelOp}

\begin{theorem}[blueprint="infogeometry:convex:oned:kernelop"]
  \inputleannode{InfoGeometry.Convex.OneD.KernelOp}
\end{theorem}

\subsection*{bregmanDiv_pythagorean_ineq}

\begin{theorem}[blueprint="infogeometry:convex:oned:bregmandiv_pythagorean_ineq"]
  \inputleannode{InfoGeometry.Convex.OneD.bregmanDiv_pythagorean_ineq}
\end{theorem}

\subsection*{bregmanDiv_three_point}

\begin{theorem}[blueprint="infogeometry:convex:oned:bregmandiv_three_point"]
  \inputleannode{InfoGeometry.Convex.OneD.bregmanDiv_three_point}
\end{theorem}

\subsection*{bregman_as_fenchelGap}

\begin{theorem}[blueprint="infogeometry:convex:oned:bregman_as_fenchelgap"]
  \inputleannode{InfoGeometry.Convex.OneD.bregman_as_fenchelGap}
\end{theorem}

\subsection*{fenchelGap}

\begin{theorem}[blueprint="infogeometry:convex:oned:fenchelgap"]
  \inputleannode{InfoGeometry.Convex.OneD.fenchelGap}
\end{theorem}

\subsection*{fenchelGap_eq_zero_of_equality}

\begin{theorem}[blueprint="infogeometry:convex:oned:fenchelgap_eq_zero_of_equality"]
  \inputleannode{InfoGeometry.Convex.OneD.fenchelGap_eq_zero_of_equality}
\end{theorem}

\subsection*{fenchelGap_nonneg_of_conjugate}

\begin{theorem}[blueprint="infogeometry:convex:oned:fenchelgap_nonneg_of_conjugate"]
  \inputleannode{InfoGeometry.Convex.OneD.fenchelGap_nonneg_of_conjugate}
\end{theorem}

\subsection*{fenchelGap_zero_iff_eq_deriv_of_dual_value_match}

\begin{theorem}[blueprint="infogeometry:convex:oned:fenchelgap_zero_iff_eq_deriv_of_dual_value_match"]
  \inputleannode{InfoGeometry.Convex.OneD.fenchelGap_zero_iff_eq_deriv_of_dual_value_match}
\end{theorem}

\subsection*{fenchelGap_zero_of_gradient}

\begin{theorem}[blueprint="infogeometry:convex:oned:fenchelgap_zero_of_gradient"]
  \inputleannode{InfoGeometry.Convex.OneD.fenchelGap_zero_of_gradient}
\end{theorem}

\subsection*{fenchelYoung_eq_of_gap_zero}

\begin{theorem}[blueprint="infogeometry:convex:oned:fenchelyoung_eq_of_gap_zero"]
  \inputleannode{InfoGeometry.Convex.OneD.fenchelYoung_eq_of_gap_zero}
\end{theorem}

\subsection*{fenchelYoung_ineq}

\begin{theorem}[blueprint="infogeometry:convex:oned:fenchelyoung_ineq"]
  \inputleannode{InfoGeometry.Convex.OneD.fenchelYoung_ineq}
\end{theorem}

\subsection*{fenchel_legendre_equivalence}

\begin{theorem}[blueprint="infogeometry:convex:oned:fenchel_legendre_equivalence"]
  \inputleannode{InfoGeometry.Convex.OneD.fenchel_legendre_equivalence}
\end{theorem}

\subsection*{fenchel_young_eq_gradient}

\begin{theorem}[blueprint="infogeometry:convex:oned:fenchel_young_eq_gradient"]
  \inputleannode{InfoGeometry.Convex.OneD.fenchel_young_eq_gradient}
\end{theorem}

\subsection*{kernelClosedFunctionsCompleteLattice}

\begin{theorem}[blueprint="infogeometry:convex:oned:kernelclosedfunctionscompletelattice"]
  \inputleannode{InfoGeometry.Convex.OneD.kernelClosedFunctionsCompleteLattice}
\end{theorem}

\subsection*{kernelClosure_monotone}

\begin{theorem}[blueprint="infogeometry:convex:oned:kernelclosure_monotone"]
  \inputleannode{InfoGeometry.Convex.OneD.kernelClosure_monotone}
\end{theorem}

\subsection*{kernelOp_eq_transpose}

\begin{theorem}[blueprint="infogeometry:convex:oned:kernelop_eq_transpose"]
  \inputleannode{InfoGeometry.Convex.OneD.kernelOp_eq_transpose}
\end{theorem}

\subsection*{kernelOp_monotone}

\begin{theorem}[blueprint="infogeometry:convex:oned:kernelop_monotone"]
  \inputleannode{InfoGeometry.Convex.OneD.kernelOp_monotone}
\end{theorem}

\subsection*{kernel_idempotent}

\begin{theorem}[blueprint="infogeometry:convex:oned:kernel_idempotent"]
  \inputleannode{InfoGeometry.Convex.OneD.kernel_idempotent}
\end{theorem}

\subsection*{kernel_le_closure}

\begin{theorem}[blueprint="infogeometry:convex:oned:kernel_le_closure"]
  \inputleannode{InfoGeometry.Convex.OneD.kernel_le_closure}
\end{theorem}

\section{InfoGeometry.Convex.OneD.DualFlatPotential}

\subsection*{bregman_as_gap}

\begin{theorem}[blueprint="infogeometry:convex:oned:dualflatpotential:bregman_as_gap"]
  \inputleannode{InfoGeometry.Convex.OneD.DualFlatPotential.bregman_as_gap}
\end{theorem}

\subsection*{casesOn}

\begin{theorem}[blueprint="infogeometry:convex:oned:dualflatpotential:caseson"]
  \inputleannode{InfoGeometry.Convex.OneD.DualFlatPotential.casesOn}
\end{theorem}

\subsection*{conjugate}

\begin{theorem}[blueprint="infogeometry:convex:oned:dualflatpotential:conjugate"]
  \inputleannode{InfoGeometry.Convex.OneD.DualFlatPotential.conjugate}
\end{theorem}

\subsection*{ctorIdx}

\begin{theorem}[blueprint="infogeometry:convex:oned:dualflatpotential:ctoridx"]
  \inputleannode{InfoGeometry.Convex.OneD.DualFlatPotential.ctorIdx}
\end{theorem}

\subsection*{dual}

\begin{theorem}[blueprint="infogeometry:convex:oned:dualflatpotential:dual"]
  \inputleannode{InfoGeometry.Convex.OneD.DualFlatPotential.dual}
\end{theorem}

\subsection*{fenchelGap_nonneg}

\begin{theorem}[blueprint="infogeometry:convex:oned:dualflatpotential:fenchelgap_nonneg"]
  \inputleannode{InfoGeometry.Convex.OneD.DualFlatPotential.fenchelGap_nonneg}
\end{theorem}

\subsection*{fenchelGap_zero}

\begin{theorem}[blueprint="infogeometry:convex:oned:dualflatpotential:fenchelgap_zero"]
  \inputleannode{InfoGeometry.Convex.OneD.DualFlatPotential.fenchelGap_zero}
\end{theorem}

\subsection*{gradientDual}

\begin{theorem}[blueprint="infogeometry:convex:oned:dualflatpotential:gradientdual"]
  \inputleannode{InfoGeometry.Convex.OneD.DualFlatPotential.gradientDual}
\end{theorem}

\subsection*{mk}

\begin{theorem}[blueprint="infogeometry:convex:oned:dualflatpotential:mk"]
  \inputleannode{InfoGeometry.Convex.OneD.DualFlatPotential.mk}
\end{theorem}

\subsection*{noConfusion}

\begin{theorem}[blueprint="infogeometry:convex:oned:dualflatpotential:noconfusion"]
  \inputleannode{InfoGeometry.Convex.OneD.DualFlatPotential.noConfusion}
\end{theorem}

\subsection*{noConfusionType}

\begin{theorem}[blueprint="infogeometry:convex:oned:dualflatpotential:noconfusiontype"]
  \inputleannode{InfoGeometry.Convex.OneD.DualFlatPotential.noConfusionType}
\end{theorem}

\subsection*{primal}

\begin{theorem}[blueprint="infogeometry:convex:oned:dualflatpotential:primal"]
  \inputleannode{InfoGeometry.Convex.OneD.DualFlatPotential.primal}
\end{theorem}

\subsection*{rec}

\begin{theorem}[blueprint="infogeometry:convex:oned:dualflatpotential:rec"]
  \inputleannode{InfoGeometry.Convex.OneD.DualFlatPotential.rec}
\end{theorem}

\subsection*{recOn}

\begin{theorem}[blueprint="infogeometry:convex:oned:dualflatpotential:recon"]
  \inputleannode{InfoGeometry.Convex.OneD.DualFlatPotential.recOn}
\end{theorem}

\section{InfoGeometry.Convex.OneD.IsKernelClosed}

\subsection*{eq_1}

\begin{theorem}[blueprint="infogeometry:convex:oned:iskernelclosed:eq_1"]
  \inputleannode{InfoGeometry.Convex.OneD.IsKernelClosed.eq_1}
\end{theorem}

\section{InfoGeometry.Convex.OneD.KernelClosure}

\subsection*{eq_1}

\begin{theorem}[blueprint="infogeometry:convex:oned:kernelclosure:eq_1"]
  \inputleannode{InfoGeometry.Convex.OneD.KernelClosure.eq_1}
\end{theorem}

\section{InfoGeometry.Convex.OneD.KernelClosureOperator}

\subsection*{eq_1}

\begin{theorem}[blueprint="infogeometry:convex:oned:kernelclosureoperator:eq_1"]
  \inputleannode{InfoGeometry.Convex.OneD.KernelClosureOperator.eq_1}
\end{theorem}

\section{InfoGeometry.Convex.OneD.KernelOp}

\subsection*{eq_1}

\begin{theorem}[blueprint="infogeometry:convex:oned:kernelop:eq_1"]
  \inputleannode{InfoGeometry.Convex.OneD.KernelOp.eq_1}
\end{theorem}

\section{InfoGeometry.Convex.ProjectiveRays}

\noindent\textbf{Buildable}: yes

\noindent\emph{No exported declarations in the current declaration graph.}

\section{InfoGeometry.Convex.RadonHelly}

\noindent\textbf{Buildable}: yes

\noindent\emph{No exported declarations in the current declaration graph.}

\section{InfoGeometry.Convex.fenchelConj}

\subsection*{eq_1}

\begin{theorem}[blueprint="infogeometry:convex:fenchelconj:eq_1"]
  \inputleannode{InfoGeometry.Convex.fenchelConj.eq_1}
\end{theorem}

\section{InfoGeometry.Convex.projectivize}

\subsection*{congr_simp}

\begin{theorem}[blueprint="infogeometry:convex:projectivize:congr_simp"]
  \inputleannode{InfoGeometry.Convex.projectivize.congr_simp}
\end{theorem}

\subsection*{eq_1}

\begin{theorem}[blueprint="infogeometry:convex:projectivize:eq_1"]
  \inputleannode{InfoGeometry.Convex.projectivize.eq_1}
\end{theorem}

\section{InfoGeometry.Convex.projectivize'}

\subsection*{eq_1}

\begin{theorem}[blueprint="infogeometry:convex:projectivize':eq_1"]
  \inputleannode{InfoGeometry.Convex.projectivize'.eq_1}
\end{theorem}

\chapter{ConvexDuality Modules}

\section{InfoGeometry.ConvexDuality}

\subsection*{KL_param}

\begin{theorem}[blueprint="infogeometry:convexduality:kl_param"]
  \inputleannode{InfoGeometry.ConvexDuality.KL_param}
\end{theorem}

\subsection*{KL_param_eq_bregman_swap}

\begin{theorem}[blueprint="infogeometry:convexduality:kl_param_eq_bregman_swap"]
  \inputleannode{InfoGeometry.ConvexDuality.KL_param_eq_bregman_swap}
\end{theorem}

\subsection*{KL_param_expanded}

\begin{theorem}[blueprint="infogeometry:convexduality:kl_param_expanded"]
  \inputleannode{InfoGeometry.ConvexDuality.KL_param_expanded}
\end{theorem}

\subsection*{KL_param_nonneg_of_convex}

\begin{theorem}[blueprint="infogeometry:convexduality:kl_param_nonneg_of_convex"]
  \inputleannode{InfoGeometry.ConvexDuality.KL_param_nonneg_of_convex}
\end{theorem}

\subsection*{KL_param_nonneg_of_convex_at}

\begin{theorem}[blueprint="infogeometry:convexduality:kl_param_nonneg_of_convex_at"]
  \inputleannode{InfoGeometry.ConvexDuality.KL_param_nonneg_of_convex_at}
\end{theorem}

\subsection*{bregman}

\begin{theorem}[blueprint="infogeometry:convexduality:bregman"]
  \inputleannode{InfoGeometry.ConvexDuality.bregman}
\end{theorem}

\subsection*{bregman_def}

\begin{theorem}[blueprint="infogeometry:convexduality:bregman_def"]
  \inputleannode{InfoGeometry.ConvexDuality.bregman_def}
\end{theorem}

\subsection*{bregman_nonneg_of_convex}

\begin{theorem}[blueprint="infogeometry:convexduality:bregman_nonneg_of_convex"]
  \inputleannode{InfoGeometry.ConvexDuality.bregman_nonneg_of_convex}
\end{theorem}

\subsection*{bregman_nonneg_of_convex_at}

\begin{theorem}[blueprint="infogeometry:convexduality:bregman_nonneg_of_convex_at"]
  \inputleannode{InfoGeometry.ConvexDuality.bregman_nonneg_of_convex_at}
\end{theorem}

\subsection*{fenchelGap}

\begin{theorem}[blueprint="infogeometry:convexduality:fenchelgap"]
  \inputleannode{InfoGeometry.ConvexDuality.fenchelGap}
\end{theorem}

\section{InfoGeometry.ConvexDuality.KL_param}

\subsection*{eq_1}

\begin{theorem}[blueprint="infogeometry:convexduality:kl_param:eq_1"]
  \inputleannode{InfoGeometry.ConvexDuality.KL_param.eq_1}
\end{theorem}

\chapter{Core Modules}

\section{InfoGeometry.Core}

\noindent\textbf{Buildable}: yes

\subsection*{CartanDecomposition}

\begin{theorem}[blueprint="infogeometry:core:cartandecomposition"]
  \inputleannode{InfoGeometry.Core.CartanDecomposition}
\end{theorem}

\subsection*{CartanInvolution}

\begin{theorem}[blueprint="infogeometry:core:cartaninvolution"]
  \inputleannode{InfoGeometry.Core.CartanInvolution}
\end{theorem}

\subsection*{CartanLieAlgebra}

\begin{theorem}[blueprint="infogeometry:core:cartanliealgebra"]
  \inputleannode{InfoGeometry.Core.CartanLieAlgebra}
\end{theorem}

\subsection*{CartanLieStructure}

\begin{theorem}[blueprint="infogeometry:core:cartanliestructure"]
  \inputleannode{InfoGeometry.Core.CartanLieStructure}
\end{theorem}

\subsection*{GrandCanonicalParams}

\begin{theorem}[blueprint="infogeometry:core:grandcanonicalparams"]
  \inputleannode{InfoGeometry.Core.GrandCanonicalParams}
\end{theorem}

\subsection*{GrandCanonicalTwoParam}

\begin{theorem}[blueprint="infogeometry:core:grandcanonicaltwoparam"]
  \inputleannode{InfoGeometry.Core.GrandCanonicalTwoParam}
\end{theorem}

\subsection*{Involution}

\begin{theorem}[blueprint="infogeometry:core:involution"]
  \inputleannode{InfoGeometry.Core.Involution}
\end{theorem}

\subsection*{InvolutiveAutomorphism}

\begin{theorem}[blueprint="infogeometry:core:involutiveautomorphism"]
  \inputleannode{InfoGeometry.Core.InvolutiveAutomorphism}
\end{theorem}

\subsection*{InvolutiveLieAut}

\begin{theorem}[blueprint="infogeometry:core:involutivelieaut"]
  \inputleannode{InfoGeometry.Core.InvolutiveLieAut}
\end{theorem}

\subsection*{InvolutiveMulAut}

\begin{theorem}[blueprint="infogeometry:core:involutivemulaut"]
  \inputleannode{InfoGeometry.Core.InvolutiveMulAut}
\end{theorem}

\subsection*{IsEven}

\begin{theorem}[blueprint="infogeometry:core:iseven"]
  \inputleannode{InfoGeometry.Core.IsEven}
\end{theorem}

\subsection*{IsOdd}

\begin{theorem}[blueprint="infogeometry:core:isodd"]
  \inputleannode{InfoGeometry.Core.IsOdd}
\end{theorem}

\subsection*{JordanAlgebra}

\begin{theorem}[blueprint="infogeometry:core:jordanalgebra"]
  \inputleannode{InfoGeometry.Core.JordanAlgebra}
\end{theorem}

\subsection*{LieAut}

\begin{theorem}[blueprint="infogeometry:core:lieaut"]
  \inputleannode{InfoGeometry.Core.LieAut}
\end{theorem}

\subsection*{LieTripleSystem}

\begin{theorem}[blueprint="infogeometry:core:lietriplesystem"]
  \inputleannode{InfoGeometry.Core.LieTripleSystem}
\end{theorem}

\subsection*{LinearInvolutiveAutomorphism}

\begin{theorem}[blueprint="infogeometry:core:linearinvolutiveautomorphism"]
  \inputleannode{InfoGeometry.Core.LinearInvolutiveAutomorphism}
\end{theorem}

\subsection*{MulInvolutiveAutomorphism}

\begin{theorem}[blueprint="infogeometry:core:mulinvolutiveautomorphism"]
  \inputleannode{InfoGeometry.Core.MulInvolutiveAutomorphism}
\end{theorem}

\subsection*{OddLocalModel}

\begin{theorem}[blueprint="infogeometry:core:oddlocalmodel"]
  \inputleannode{InfoGeometry.Core.OddLocalModel}
\end{theorem}

\subsection*{PreservesLie}

\begin{theorem}[blueprint="infogeometry:core:preserveslie"]
  \inputleannode{InfoGeometry.Core.PreservesLie}
\end{theorem}

\subsection*{PreservesLieBracket}

\begin{theorem}[blueprint="infogeometry:core:preservesliebracket"]
  \inputleannode{InfoGeometry.Core.PreservesLieBracket}
\end{theorem}

\subsection*{PreservesLinear}

\begin{theorem}[blueprint="infogeometry:core:preserveslinear"]
  \inputleannode{InfoGeometry.Core.PreservesLinear}
\end{theorem}

\subsection*{PreservesMul}

\begin{theorem}[blueprint="infogeometry:core:preservesmul"]
  \inputleannode{InfoGeometry.Core.PreservesMul}
\end{theorem}

\subsection*{SPD}

\begin{theorem}[blueprint="infogeometry:core:spd"]
  \inputleannode{InfoGeometry.Core.SPD}
\end{theorem}

\subsection*{SPD_posDef}

\begin{theorem}[blueprint="infogeometry:core:spd_posdef"]
  \inputleannode{InfoGeometry.Core.SPD_posDef}
\end{theorem}

\subsection*{SPD_transpose_eq_self}

\begin{theorem}[blueprint="infogeometry:core:spd_transpose_eq_self"]
  \inputleannode{InfoGeometry.Core.SPD_transpose_eq_self}
\end{theorem}

\subsection*{Spinodal}

\begin{theorem}[blueprint="infogeometry:core:spinodal"]
  \inputleannode{InfoGeometry.Core.Spinodal}
\end{theorem}

\subsection*{SymmetricLieAlgebra}

\begin{theorem}[blueprint="infogeometry:core:symmetricliealgebra"]
  \inputleannode{InfoGeometry.Core.SymmetricLieAlgebra}
\end{theorem}

\subsection*{SymmetricPair}

\begin{theorem}[blueprint="infogeometry:core:symmetricpair"]
  \inputleannode{InfoGeometry.Core.SymmetricPair}
\end{theorem}

\subsection*{SymmetricSpace}

\begin{theorem}[blueprint="infogeometry:core:symmetricspace"]
  \inputleannode{InfoGeometry.Core.SymmetricSpace}
\end{theorem}

\subsection*{cartanMinus}

\begin{theorem}[blueprint="infogeometry:core:cartanminus"]
  \inputleannode{InfoGeometry.Core.cartanMinus}
\end{theorem}

\subsection*{cartanMinus_neg_fixed}

\begin{theorem}[blueprint="infogeometry:core:cartanminus_neg_fixed"]
  \inputleannode{InfoGeometry.Core.cartanMinus_neg_fixed}
\end{theorem}

\subsection*{cartanMinus_odd}

\begin{theorem}[blueprint="infogeometry:core:cartanminus_odd"]
  \inputleannode{InfoGeometry.Core.cartanMinus_odd}
\end{theorem}

\subsection*{cartanOfMulInvolutive}

\begin{theorem}[blueprint="infogeometry:core:cartanofmulinvolutive"]
  \inputleannode{InfoGeometry.Core.cartanOfMulInvolutive}
\end{theorem}

\subsection*{cartanOfMulInvolutive_toMulInvolutive}

\begin{theorem}[blueprint="infogeometry:core:cartanofmulinvolutive_tomulinvolutive"]
  \inputleannode{InfoGeometry.Core.cartanOfMulInvolutive_toMulInvolutive}
\end{theorem}

\subsection*{cartanPlus}

\begin{theorem}[blueprint="infogeometry:core:cartanplus"]
  \inputleannode{InfoGeometry.Core.cartanPlus}
\end{theorem}

\subsection*{cartanPlus_even}

\begin{theorem}[blueprint="infogeometry:core:cartanplus_even"]
  \inputleannode{InfoGeometry.Core.cartanPlus_even}
\end{theorem}

\subsection*{cartanPlus_fixed}

\begin{theorem}[blueprint="infogeometry:core:cartanplus_fixed"]
  \inputleannode{InfoGeometry.Core.cartanPlus_fixed}
\end{theorem}

\subsection*{cartanSymmetry}

\begin{theorem}[blueprint="infogeometry:core:cartansymmetry"]
  \inputleannode{InfoGeometry.Core.cartanSymmetry}
\end{theorem}

\subsection*{cartanSymmetryOfInvolutiveMulAut}

\begin{theorem}[blueprint="infogeometry:core:cartansymmetryofinvolutivemulaut"]
  \inputleannode{InfoGeometry.Core.cartanSymmetryOfInvolutiveMulAut}
\end{theorem}

\subsection*{cartanSymmetry_fixpoint_ofInvolutiveMulAut}

\begin{theorem}[blueprint="infogeometry:core:cartansymmetry_fixpoint_ofinvolutivemulaut"]
  \inputleannode{InfoGeometry.Core.cartanSymmetry_fixpoint_ofInvolutiveMulAut}
\end{theorem}

\subsection*{cartanSymmetry_involutive_ofInvolutiveMulAut}

\begin{theorem}[blueprint="infogeometry:core:cartansymmetry_involutive_ofinvolutivemulaut"]
  \inputleannode{InfoGeometry.Core.cartanSymmetry_involutive_ofInvolutiveMulAut}
\end{theorem}

\subsection*{cartan_decomposition}

\begin{theorem}[blueprint="infogeometry:core:cartan_decomposition"]
  \inputleannode{InfoGeometry.Core.cartan_decomposition}
\end{theorem}

\subsection*{cartan_split}

\begin{theorem}[blueprint="infogeometry:core:cartan_split"]
  \inputleannode{InfoGeometry.Core.cartan_split}
\end{theorem}

\subsection*{conjugationByInvolution}

\begin{theorem}[blueprint="infogeometry:core:conjugationbyinvolution"]
  \inputleannode{InfoGeometry.Core.conjugationByInvolution}
\end{theorem}

\subsection*{conjugationEnd}

\begin{theorem}[blueprint="infogeometry:core:conjugationend"]
  \inputleannode{InfoGeometry.Core.conjugationEnd}
\end{theorem}

\subsection*{conjugationEvenLieSubalgebra}

\begin{theorem}[blueprint="infogeometry:core:conjugationevenliesubalgebra"]
  \inputleannode{InfoGeometry.Core.conjugationEvenLieSubalgebra}
\end{theorem}

\subsection*{conjugationEvenSubmodule}

\begin{theorem}[blueprint="infogeometry:core:conjugationevensubmodule"]
  \inputleannode{InfoGeometry.Core.conjugationEvenSubmodule}
\end{theorem}

\subsection*{conjugationLieEquiv}

\begin{theorem}[blueprint="infogeometry:core:conjugationlieequiv"]
  \inputleannode{InfoGeometry.Core.conjugationLieEquiv}
\end{theorem}

\subsection*{conjugationMap}

\begin{theorem}[blueprint="infogeometry:core:conjugationmap"]
  \inputleannode{InfoGeometry.Core.conjugationMap}
\end{theorem}

\subsection*{conjugationMap_bracket}

\begin{theorem}[blueprint="infogeometry:core:conjugationmap_bracket"]
  \inputleannode{InfoGeometry.Core.conjugationMap_bracket}
\end{theorem}

\subsection*{conjugationMap_involutive}

\begin{theorem}[blueprint="infogeometry:core:conjugationmap_involutive"]
  \inputleannode{InfoGeometry.Core.conjugationMap_involutive}
\end{theorem}

\subsection*{conjugationOddSubmodule}

\begin{theorem}[blueprint="infogeometry:core:conjugationoddsubmodule"]
  \inputleannode{InfoGeometry.Core.conjugationOddSubmodule}
\end{theorem}

\subsection*{conjugationSymmetricLieAlgebra}

\begin{theorem}[blueprint="infogeometry:core:conjugationsymmetricliealgebra"]
  \inputleannode{InfoGeometry.Core.conjugationSymmetricLieAlgebra}
\end{theorem}

\subsection*{conjugation_bracket_even_even_mem_even}

\begin{theorem}[blueprint="infogeometry:core:conjugation_bracket_even_even_mem_even"]
  \inputleannode{InfoGeometry.Core.conjugation_bracket_even_even_mem_even}
\end{theorem}

\subsection*{conjugation_bracket_even_odd_mem_odd}

\begin{theorem}[blueprint="infogeometry:core:conjugation_bracket_even_odd_mem_odd"]
  \inputleannode{InfoGeometry.Core.conjugation_bracket_even_odd_mem_odd}
\end{theorem}

\subsection*{conjugation_bracket_odd_odd_mem_even}

\begin{theorem}[blueprint="infogeometry:core:conjugation_bracket_odd_odd_mem_even"]
  \inputleannode{InfoGeometry.Core.conjugation_bracket_odd_odd_mem_even}
\end{theorem}

\subsection*{endoSymmetricLieAlgebra}

\begin{theorem}[blueprint="infogeometry:core:endosymmetricliealgebra"]
  \inputleannode{InfoGeometry.Core.endoSymmetricLieAlgebra}
\end{theorem}

\subsection*{endomorphismCartanDecomposition}

\begin{theorem}[blueprint="infogeometry:core:endomorphismcartandecomposition"]
  \inputleannode{InfoGeometry.Core.endomorphismCartanDecomposition}
\end{theorem}

\subsection*{entropy_def}

\begin{theorem}[blueprint="infogeometry:core:entropy_def"]
  \inputleannode{InfoGeometry.Core.entropy_def}
\end{theorem}

\subsection*{expectation_const}

\begin{theorem}[blueprint="infogeometry:core:expectation_const"]
  \inputleannode{InfoGeometry.Core.expectation_const}
\end{theorem}

\subsection*{expectation_one}

\begin{theorem}[blueprint="infogeometry:core:expectation_one"]
  \inputleannode{InfoGeometry.Core.expectation_one}
\end{theorem}

\subsection*{expectation_zero}

\begin{theorem}[blueprint="infogeometry:core:expectation_zero"]
  \inputleannode{InfoGeometry.Core.expectation_zero}
\end{theorem}

\subsection*{fixedSubgroup}

\begin{theorem}[blueprint="infogeometry:core:fixedsubgroup"]
  \inputleannode{InfoGeometry.Core.fixedSubgroup}
\end{theorem}

\subsection*{gc2_gibbsWeight_sum_one}

\begin{theorem}[blueprint="infogeometry:core:gc2_gibbsweight_sum_one"]
  \inputleannode{InfoGeometry.Core.gc2_gibbsWeight_sum_one}
\end{theorem}

\subsection*{gc2_partition_pos}

\begin{theorem}[blueprint="infogeometry:core:gc2_partition_pos"]
  \inputleannode{InfoGeometry.Core.gc2_partition_pos}
\end{theorem}

\subsection*{gc2_potential_deriv_beta_eq_neg_meanShift}

\begin{theorem}[blueprint="infogeometry:core:gc2_potential_deriv_beta_eq_neg_meanshift"]
  \inputleannode{InfoGeometry.Core.gc2_potential_deriv_beta_eq_neg_meanShift}
\end{theorem}

\subsection*{gc2_potential_deriv_mu_eq_beta_meanNumber}

\begin{theorem}[blueprint="infogeometry:core:gc2_potential_deriv_mu_eq_beta_meannumber"]
  \inputleannode{InfoGeometry.Core.gc2_potential_deriv_mu_eq_beta_meanNumber}
\end{theorem}

\subsection*{gc_gibbsWeight_sum_one}

\begin{theorem}[blueprint="infogeometry:core:gc_gibbsweight_sum_one"]
  \inputleannode{InfoGeometry.Core.gc_gibbsWeight_sum_one}
\end{theorem}

\subsection*{gc_hessian_eq_variance}

\begin{theorem}[blueprint="infogeometry:core:gc_hessian_eq_variance"]
  \inputleannode{InfoGeometry.Core.gc_hessian_eq_variance}
\end{theorem}

\subsection*{gc_hessian_nonneg}

\begin{theorem}[blueprint="infogeometry:core:gc_hessian_nonneg"]
  \inputleannode{InfoGeometry.Core.gc_hessian_nonneg}
\end{theorem}

\subsection*{gc_partition_pos}

\begin{theorem}[blueprint="infogeometry:core:gc_partition_pos"]
  \inputleannode{InfoGeometry.Core.gc_partition_pos}
\end{theorem}

\subsection*{gc_potential_deriv_eq_neg_mean}

\begin{theorem}[blueprint="infogeometry:core:gc_potential_deriv_eq_neg_mean"]
  \inputleannode{InfoGeometry.Core.gc_potential_deriv_eq_neg_mean}
\end{theorem}

\subsection*{gc_potential_second_derivative_eq_variance}

\begin{theorem}[blueprint="infogeometry:core:gc_potential_second_derivative_eq_variance"]
  \inputleannode{InfoGeometry.Core.gc_potential_second_derivative_eq_variance}
\end{theorem}

\subsection*{gc_spinodal_iff_variance_eq_zero}

\begin{theorem}[blueprint="infogeometry:core:gc_spinodal_iff_variance_eq_zero"]
  \inputleannode{InfoGeometry.Core.gc_spinodal_iff_variance_eq_zero}
\end{theorem}

\subsection*{gibbsWeight}

\begin{theorem}[blueprint="infogeometry:core:gibbsweight"]
  \inputleannode{InfoGeometry.Core.gibbsWeight}
\end{theorem}

\subsection*{gibbsWeightGC}

\begin{theorem}[blueprint="infogeometry:core:gibbsweightgc"]
  \inputleannode{InfoGeometry.Core.gibbsWeightGC}
\end{theorem}

\subsection*{hessian}

\begin{theorem}[blueprint="infogeometry:core:hessian"]
  \inputleannode{InfoGeometry.Core.hessian}
\end{theorem}

\subsection*{instCoeFunInvolutiveAutomorphismForall}

\begin{theorem}[blueprint="infogeometry:core:instcoefuninvolutiveautomorphismforall"]
  \inputleannode{InfoGeometry.Core.instCoeFunInvolutiveAutomorphismForall}
\end{theorem}

\subsection*{instCoeFunLinearInvolutiveAutomorphismForall}

\begin{theorem}[blueprint="infogeometry:core:instcoefunlinearinvolutiveautomorphismforall"]
  \inputleannode{InfoGeometry.Core.instCoeFunLinearInvolutiveAutomorphismForall}
\end{theorem}

\subsection*{instCoeFunMulInvolutiveAutomorphismForall}

\begin{theorem}[blueprint="infogeometry:core:instcoefunmulinvolutiveautomorphismforall"]
  \inputleannode{InfoGeometry.Core.instCoeFunMulInvolutiveAutomorphismForall}
\end{theorem}

\subsection*{instPreservesLieBracketOfPreservesLie}

\begin{theorem}[blueprint="infogeometry:core:instpreservesliebracketofpreserveslie"]
  \inputleannode{InfoGeometry.Core.instPreservesLieBracketOfPreservesLie}
\end{theorem}

\subsection*{instPreservesLinearOfPreservesLie}

\begin{theorem}[blueprint="infogeometry:core:instpreserveslinearofpreserveslie"]
  \inputleannode{InfoGeometry.Core.instPreservesLinearOfPreservesLie}
\end{theorem}

\subsection*{klDiv_def}

\begin{theorem}[blueprint="infogeometry:core:kldiv_def"]
  \inputleannode{InfoGeometry.Core.klDiv_def}
\end{theorem}

\subsection*{logDensity_def}

\begin{theorem}[blueprint="infogeometry:core:logdensity_def"]
  \inputleannode{InfoGeometry.Core.logDensity_def}
\end{theorem}

\subsection*{mean}

\begin{theorem}[blueprint="infogeometry:core:mean"]
  \inputleannode{InfoGeometry.Core.mean}
\end{theorem}

\subsection*{meanNumber}

\begin{theorem}[blueprint="infogeometry:core:meannumber"]
  \inputleannode{InfoGeometry.Core.meanNumber}
\end{theorem}

\subsection*{meanShift}

\begin{theorem}[blueprint="infogeometry:core:meanshift"]
  \inputleannode{InfoGeometry.Core.meanShift}
\end{theorem}

\subsection*{minusPart}

\begin{theorem}[blueprint="infogeometry:core:minuspart"]
  \inputleannode{InfoGeometry.Core.minusPart}
\end{theorem}

\subsection*{mulInvolutiveOfCartan}

\begin{theorem}[blueprint="infogeometry:core:mulinvolutiveofcartan"]
  \inputleannode{InfoGeometry.Core.mulInvolutiveOfCartan}
\end{theorem}

\subsection*{mulInvolutiveOfCartan_toCartan}

\begin{theorem}[blueprint="infogeometry:core:mulinvolutiveofcartan_tocartan"]
  \inputleannode{InfoGeometry.Core.mulInvolutiveOfCartan_toCartan}
\end{theorem}

\subsection*{partition}

\begin{theorem}[blueprint="infogeometry:core:partition"]
  \inputleannode{InfoGeometry.Core.partition}
\end{theorem}

\subsection*{partitionGC}

\begin{theorem}[blueprint="infogeometry:core:partitiongc"]
  \inputleannode{InfoGeometry.Core.partitionGC}
\end{theorem}

\subsection*{plusPart}

\begin{theorem}[blueprint="infogeometry:core:pluspart"]
  \inputleannode{InfoGeometry.Core.plusPart}
\end{theorem}

\subsection*{potential}

\begin{theorem}[blueprint="infogeometry:core:potential"]
  \inputleannode{InfoGeometry.Core.potential}
\end{theorem}

\subsection*{potentialGC}

\begin{theorem}[blueprint="infogeometry:core:potentialgc"]
  \inputleannode{InfoGeometry.Core.potentialGC}
\end{theorem}

\subsection*{reflection_involutive}

\begin{theorem}[blueprint="infogeometry:core:reflection_involutive"]
  \inputleannode{InfoGeometry.Core.reflection_involutive}
\end{theorem}

\subsection*{reflection_self}

\begin{theorem}[blueprint="infogeometry:core:reflection_self"]
  \inputleannode{InfoGeometry.Core.reflection_self}
\end{theorem}

\subsection*{shiftedEnergy}

\begin{theorem}[blueprint="infogeometry:core:shiftedenergy"]
  \inputleannode{InfoGeometry.Core.shiftedEnergy}
\end{theorem}

\subsection*{spd_pos_def}

\begin{theorem}[blueprint="infogeometry:core:spd_pos_def"]
  \inputleannode{InfoGeometry.Core.spd_pos_def}
\end{theorem}

\subsection*{spd_transpose_eq_self}

\begin{theorem}[blueprint="infogeometry:core:spd_transpose_eq_self"]
  \inputleannode{InfoGeometry.Core.spd_transpose_eq_self}
\end{theorem}

\subsection*{surprisal_def}

\begin{theorem}[blueprint="infogeometry:core:surprisal_def"]
  \inputleannode{InfoGeometry.Core.surprisal_def}
\end{theorem}

\subsection*{symmetricPairOfInvolutiveMulAut}

\begin{theorem}[blueprint="infogeometry:core:symmetricpairofinvolutivemulaut"]
  \inputleannode{InfoGeometry.Core.symmetricPairOfInvolutiveMulAut}
\end{theorem}

\subsection*{symmetricPairOfInvolutiveMulAut_K_eq_fixed}

\begin{theorem}[blueprint="infogeometry:core:symmetricpairofinvolutivemulaut_k_eq_fixed"]
  \inputleannode{InfoGeometry.Core.symmetricPairOfInvolutiveMulAut_K_eq_fixed}
\end{theorem}

\subsection*{symmetricPairOfInvolutiveMulAut_K_eq_fixedSubgroup}

\begin{theorem}[blueprint="infogeometry:core:symmetricpairofinvolutivemulaut_k_eq_fixedsubgroup"]
  \inputleannode{InfoGeometry.Core.symmetricPairOfInvolutiveMulAut_K_eq_fixedSubgroup}
\end{theorem}

\subsection*{symmetricPair_K_eq_fixedSubgroup}

\begin{theorem}[blueprint="infogeometry:core:symmetricpair_k_eq_fixedsubgroup"]
  \inputleannode{InfoGeometry.Core.symmetricPair_K_eq_fixedSubgroup}
\end{theorem}

\subsection*{variance}

\begin{theorem}[blueprint="infogeometry:core:variance"]
  \inputleannode{InfoGeometry.Core.variance}
\end{theorem}

\section{InfoGeometry.Core.CartanDecomposition}

\subsection*{casesOn}

\begin{theorem}[blueprint="infogeometry:core:cartandecomposition:caseson"]
  \inputleannode{InfoGeometry.Core.CartanDecomposition.casesOn}
\end{theorem}

\subsection*{ctorIdx}

\begin{theorem}[blueprint="infogeometry:core:cartandecomposition:ctoridx"]
  \inputleannode{InfoGeometry.Core.CartanDecomposition.ctorIdx}
\end{theorem}

\subsection*{ext}

\begin{theorem}[blueprint="infogeometry:core:cartandecomposition:ext"]
  \inputleannode{InfoGeometry.Core.CartanDecomposition.ext}
\end{theorem}

\subsection*{ext_iff}

\begin{theorem}[blueprint="infogeometry:core:cartandecomposition:ext_iff"]
  \inputleannode{InfoGeometry.Core.CartanDecomposition.ext_iff}
\end{theorem}

\subsection*{mk}

\begin{theorem}[blueprint="infogeometry:core:cartandecomposition:mk"]
  \inputleannode{InfoGeometry.Core.CartanDecomposition.mk}
\end{theorem}

\subsection*{noConfusion}

\begin{theorem}[blueprint="infogeometry:core:cartandecomposition:noconfusion"]
  \inputleannode{InfoGeometry.Core.CartanDecomposition.noConfusion}
\end{theorem}

\subsection*{noConfusionType}

\begin{theorem}[blueprint="infogeometry:core:cartandecomposition:noconfusiontype"]
  \inputleannode{InfoGeometry.Core.CartanDecomposition.noConfusionType}
\end{theorem}

\subsection*{rec}

\begin{theorem}[blueprint="infogeometry:core:cartandecomposition:rec"]
  \inputleannode{InfoGeometry.Core.CartanDecomposition.rec}
\end{theorem}

\subsection*{recOn}

\begin{theorem}[blueprint="infogeometry:core:cartandecomposition:recon"]
  \inputleannode{InfoGeometry.Core.CartanDecomposition.recOn}
\end{theorem}

\subsection*{θ}

\begin{theorem}[blueprint="infogeometry:core:cartandecomposition:θ"]
  \inputleannode{InfoGeometry.Core.CartanDecomposition.θ}
\end{theorem}

\section{InfoGeometry.Core.CartanLieAlgebra}

\subsection*{casesOn}

\begin{theorem}[blueprint="infogeometry:core:cartanliealgebra:caseson"]
  \inputleannode{InfoGeometry.Core.CartanLieAlgebra.casesOn}
\end{theorem}

\subsection*{ctorIdx}

\begin{theorem}[blueprint="infogeometry:core:cartanliealgebra:ctoridx"]
  \inputleannode{InfoGeometry.Core.CartanLieAlgebra.ctorIdx}
\end{theorem}

\subsection*{killing_nondegenerate}

\begin{theorem}[blueprint="infogeometry:core:cartanliealgebra:killing_nondegenerate"]
  \inputleannode{InfoGeometry.Core.CartanLieAlgebra.killing_nondegenerate}
\end{theorem}

\subsection*{mk}

\begin{theorem}[blueprint="infogeometry:core:cartanliealgebra:mk"]
  \inputleannode{InfoGeometry.Core.CartanLieAlgebra.mk}
\end{theorem}

\subsection*{noConfusion}

\begin{theorem}[blueprint="infogeometry:core:cartanliealgebra:noconfusion"]
  \inputleannode{InfoGeometry.Core.CartanLieAlgebra.noConfusion}
\end{theorem}

\subsection*{noConfusionType}

\begin{theorem}[blueprint="infogeometry:core:cartanliealgebra:noconfusiontype"]
  \inputleannode{InfoGeometry.Core.CartanLieAlgebra.noConfusionType}
\end{theorem}

\subsection*{rec}

\begin{theorem}[blueprint="infogeometry:core:cartanliealgebra:rec"]
  \inputleannode{InfoGeometry.Core.CartanLieAlgebra.rec}
\end{theorem}

\subsection*{recOn}

\begin{theorem}[blueprint="infogeometry:core:cartanliealgebra:recon"]
  \inputleannode{InfoGeometry.Core.CartanLieAlgebra.recOn}
\end{theorem}

\subsection*{toSymmetricLieAlgebra}

\begin{theorem}[blueprint="infogeometry:core:cartanliealgebra:tosymmetricliealgebra"]
  \inputleannode{InfoGeometry.Core.CartanLieAlgebra.toSymmetricLieAlgebra}
\end{theorem}

\section{InfoGeometry.Core.CartanLieStructure}

\subsection*{bracket}

\begin{theorem}[blueprint="infogeometry:core:cartanliestructure:bracket"]
  \inputleannode{InfoGeometry.Core.CartanLieStructure.bracket}
\end{theorem}

\subsection*{bracket_compat}

\begin{theorem}[blueprint="infogeometry:core:cartanliestructure:bracket_compat"]
  \inputleannode{InfoGeometry.Core.CartanLieStructure.bracket_compat}
\end{theorem}

\subsection*{cartan}

\begin{theorem}[blueprint="infogeometry:core:cartanliestructure:cartan"]
  \inputleannode{InfoGeometry.Core.CartanLieStructure.cartan}
\end{theorem}

\subsection*{casesOn}

\begin{theorem}[blueprint="infogeometry:core:cartanliestructure:caseson"]
  \inputleannode{InfoGeometry.Core.CartanLieStructure.casesOn}
\end{theorem}

\subsection*{ctorIdx}

\begin{theorem}[blueprint="infogeometry:core:cartanliestructure:ctoridx"]
  \inputleannode{InfoGeometry.Core.CartanLieStructure.ctorIdx}
\end{theorem}

\subsection*{mk}

\begin{theorem}[blueprint="infogeometry:core:cartanliestructure:mk"]
  \inputleannode{InfoGeometry.Core.CartanLieStructure.mk}
\end{theorem}

\subsection*{noConfusion}

\begin{theorem}[blueprint="infogeometry:core:cartanliestructure:noconfusion"]
  \inputleannode{InfoGeometry.Core.CartanLieStructure.noConfusion}
\end{theorem}

\subsection*{noConfusionType}

\begin{theorem}[blueprint="infogeometry:core:cartanliestructure:noconfusiontype"]
  \inputleannode{InfoGeometry.Core.CartanLieStructure.noConfusionType}
\end{theorem}

\subsection*{rec}

\begin{theorem}[blueprint="infogeometry:core:cartanliestructure:rec"]
  \inputleannode{InfoGeometry.Core.CartanLieStructure.rec}
\end{theorem}

\subsection*{recOn}

\begin{theorem}[blueprint="infogeometry:core:cartanliestructure:recon"]
  \inputleannode{InfoGeometry.Core.CartanLieStructure.recOn}
\end{theorem}

\section{InfoGeometry.Core.Derivatives}

\noindent\textbf{Buildable}: yes

\noindent\emph{No exported declarations in the current declaration graph.}

\section{InfoGeometry.Core.DerivativesSmoke}

\noindent\textbf{Buildable}: yes

\noindent\emph{No exported declarations in the current declaration graph.}

\section{InfoGeometry.Core.Entropy}

\noindent\textbf{Buildable}: yes

\noindent\emph{No exported declarations in the current declaration graph.}

\section{InfoGeometry.Core.Generic}

\subsection*{SymmetricLieAlgebra}

\begin{theorem}[blueprint="infogeometry:core:generic:symmetricliealgebra"]
  \inputleannode{InfoGeometry.Core.Generic.SymmetricLieAlgebra}
\end{theorem}

\section{InfoGeometry.Core.Generic.SymmetricLieAlgebra}

\subsection*{B}

\begin{theorem}[blueprint="infogeometry:core:generic:symmetricliealgebra:b"]
  \inputleannode{InfoGeometry.Core.Generic.SymmetricLieAlgebra.B}
\end{theorem}

\subsection*{CartanLieAlgebra}

\begin{theorem}[blueprint="infogeometry:core:generic:symmetricliealgebra:cartanliealgebra"]
  \inputleannode{InfoGeometry.Core.Generic.SymmetricLieAlgebra.CartanLieAlgebra}
\end{theorem}

\subsection*{CartanSignature}

\begin{theorem}[blueprint="infogeometry:core:generic:symmetricliealgebra:cartansignature"]
  \inputleannode{InfoGeometry.Core.Generic.SymmetricLieAlgebra.CartanSignature}
\end{theorem}

\subsection*{P_minus}

\begin{theorem}[blueprint="infogeometry:core:generic:symmetricliealgebra:p_minus"]
  \inputleannode{InfoGeometry.Core.Generic.SymmetricLieAlgebra.P_minus}
\end{theorem}

\subsection*{P_minus_apply}

\begin{theorem}[blueprint="infogeometry:core:generic:symmetricliealgebra:p_minus_apply"]
  \inputleannode{InfoGeometry.Core.Generic.SymmetricLieAlgebra.P_minus_apply}
\end{theorem}

\subsection*{P_minus_comp_P_plus}

\begin{theorem}[blueprint="infogeometry:core:generic:symmetricliealgebra:p_minus_comp_p_plus"]
  \inputleannode{InfoGeometry.Core.Generic.SymmetricLieAlgebra.P_minus_comp_P_plus}
\end{theorem}

\subsection*{P_minus_eq_self_of_mem_𝔭}

\begin{theorem}[blueprint="infogeometry:core:generic:symmetricliealgebra:p_minus_eq_self_of_mem_𝔭"]
  \inputleannode{InfoGeometry.Core.Generic.SymmetricLieAlgebra.P_minus_eq_self_of_mem_𝔭}
\end{theorem}

\subsection*{P_minus_eq_zero_of_mem_𝔨}

\begin{theorem}[blueprint="infogeometry:core:generic:symmetricliealgebra:p_minus_eq_zero_of_mem_𝔨"]
  \inputleannode{InfoGeometry.Core.Generic.SymmetricLieAlgebra.P_minus_eq_zero_of_mem_𝔨}
\end{theorem}

\subsection*{P_minus_idempotent}

\begin{theorem}[blueprint="infogeometry:core:generic:symmetricliealgebra:p_minus_idempotent"]
  \inputleannode{InfoGeometry.Core.Generic.SymmetricLieAlgebra.P_minus_idempotent}
\end{theorem}

\subsection*{P_minus_mem_𝔭}

\begin{theorem}[blueprint="infogeometry:core:generic:symmetricliealgebra:p_minus_mem_𝔭"]
  \inputleannode{InfoGeometry.Core.Generic.SymmetricLieAlgebra.P_minus_mem_𝔭}
\end{theorem}

\subsection*{P_minus_neg_fixed}

\begin{theorem}[blueprint="infogeometry:core:generic:symmetricliealgebra:p_minus_neg_fixed"]
  \inputleannode{InfoGeometry.Core.Generic.SymmetricLieAlgebra.P_minus_neg_fixed}
\end{theorem}

\subsection*{P_plus}

\begin{theorem}[blueprint="infogeometry:core:generic:symmetricliealgebra:p_plus"]
  \inputleannode{InfoGeometry.Core.Generic.SymmetricLieAlgebra.P_plus}
\end{theorem}

\subsection*{P_plus_apply}

\begin{theorem}[blueprint="infogeometry:core:generic:symmetricliealgebra:p_plus_apply"]
  \inputleannode{InfoGeometry.Core.Generic.SymmetricLieAlgebra.P_plus_apply}
\end{theorem}

\subsection*{P_plus_comp_P_minus}

\begin{theorem}[blueprint="infogeometry:core:generic:symmetricliealgebra:p_plus_comp_p_minus"]
  \inputleannode{InfoGeometry.Core.Generic.SymmetricLieAlgebra.P_plus_comp_P_minus}
\end{theorem}

\subsection*{P_plus_eq_self_of_mem_𝔨}

\begin{theorem}[blueprint="infogeometry:core:generic:symmetricliealgebra:p_plus_eq_self_of_mem_𝔨"]
  \inputleannode{InfoGeometry.Core.Generic.SymmetricLieAlgebra.P_plus_eq_self_of_mem_𝔨}
\end{theorem}

\subsection*{P_plus_eq_zero_of_mem_𝔭}

\begin{theorem}[blueprint="infogeometry:core:generic:symmetricliealgebra:p_plus_eq_zero_of_mem_𝔭"]
  \inputleannode{InfoGeometry.Core.Generic.SymmetricLieAlgebra.P_plus_eq_zero_of_mem_𝔭}
\end{theorem}

\subsection*{P_plus_fixed}

\begin{theorem}[blueprint="infogeometry:core:generic:symmetricliealgebra:p_plus_fixed"]
  \inputleannode{InfoGeometry.Core.Generic.SymmetricLieAlgebra.P_plus_fixed}
\end{theorem}

\subsection*{P_plus_idempotent}

\begin{theorem}[blueprint="infogeometry:core:generic:symmetricliealgebra:p_plus_idempotent"]
  \inputleannode{InfoGeometry.Core.Generic.SymmetricLieAlgebra.P_plus_idempotent}
\end{theorem}

\subsection*{P_plus_mem_𝔨}

\begin{theorem}[blueprint="infogeometry:core:generic:symmetricliealgebra:p_plus_mem_𝔨"]
  \inputleannode{InfoGeometry.Core.Generic.SymmetricLieAlgebra.P_plus_mem_𝔨}
\end{theorem}

\subsection*{bracket_k_k}

\begin{theorem}[blueprint="infogeometry:core:generic:symmetricliealgebra:bracket_k_k"]
  \inputleannode{InfoGeometry.Core.Generic.SymmetricLieAlgebra.bracket_k_k}
\end{theorem}

\subsection*{bracket_k_p}

\begin{theorem}[blueprint="infogeometry:core:generic:symmetricliealgebra:bracket_k_p"]
  \inputleannode{InfoGeometry.Core.Generic.SymmetricLieAlgebra.bracket_k_p}
\end{theorem}

\subsection*{bracket_p_p}

\begin{theorem}[blueprint="infogeometry:core:generic:symmetricliealgebra:bracket_p_p"]
  \inputleannode{InfoGeometry.Core.Generic.SymmetricLieAlgebra.bracket_p_p}
\end{theorem}

\subsection*{cartanForm}

\begin{theorem}[blueprint="infogeometry:core:generic:symmetricliealgebra:cartanform"]
  \inputleannode{InfoGeometry.Core.Generic.SymmetricLieAlgebra.cartanForm}
\end{theorem}

\subsection*{cartan_decomposition}

\begin{theorem}[blueprint="infogeometry:core:generic:symmetricliealgebra:cartan_decomposition"]
  \inputleannode{InfoGeometry.Core.Generic.SymmetricLieAlgebra.cartan_decomposition}
\end{theorem}

\subsection*{casesOn}

\begin{theorem}[blueprint="infogeometry:core:generic:symmetricliealgebra:caseson"]
  \inputleannode{InfoGeometry.Core.Generic.SymmetricLieAlgebra.casesOn}
\end{theorem}

\subsection*{ctorIdx}

\begin{theorem}[blueprint="infogeometry:core:generic:symmetricliealgebra:ctoridx"]
  \inputleannode{InfoGeometry.Core.Generic.SymmetricLieAlgebra.ctorIdx}
\end{theorem}

\subsection*{curvature}

\begin{theorem}[blueprint="infogeometry:core:generic:symmetricliealgebra:curvature"]
  \inputleannode{InfoGeometry.Core.Generic.SymmetricLieAlgebra.curvature}
\end{theorem}

\subsection*{involution}

\begin{theorem}[blueprint="infogeometry:core:generic:symmetricliealgebra:involution"]
  \inputleannode{InfoGeometry.Core.Generic.SymmetricLieAlgebra.involution}
\end{theorem}

\subsection*{involution_apply}

\begin{theorem}[blueprint="infogeometry:core:generic:symmetricliealgebra:involution_apply"]
  \inputleannode{InfoGeometry.Core.Generic.SymmetricLieAlgebra.involution_apply}
\end{theorem}

\subsection*{killing_invariant}

\begin{theorem}[blueprint="infogeometry:core:generic:symmetricliealgebra:killing_invariant"]
  \inputleannode{InfoGeometry.Core.Generic.SymmetricLieAlgebra.killing_invariant}
\end{theorem}

\subsection*{killing_orthogonal}

\begin{theorem}[blueprint="infogeometry:core:generic:symmetricliealgebra:killing_orthogonal"]
  \inputleannode{InfoGeometry.Core.Generic.SymmetricLieAlgebra.killing_orthogonal}
\end{theorem}

\subsection*{mem_𝔨_iff}

\begin{theorem}[blueprint="infogeometry:core:generic:symmetricliealgebra:mem_𝔨_iff"]
  \inputleannode{InfoGeometry.Core.Generic.SymmetricLieAlgebra.mem_𝔨_iff}
\end{theorem}

\subsection*{mem_𝔭_iff}

\begin{theorem}[blueprint="infogeometry:core:generic:symmetricliealgebra:mem_𝔭_iff"]
  \inputleannode{InfoGeometry.Core.Generic.SymmetricLieAlgebra.mem_𝔭_iff}
\end{theorem}

\subsection*{mk}

\begin{theorem}[blueprint="infogeometry:core:generic:symmetricliealgebra:mk"]
  \inputleannode{InfoGeometry.Core.Generic.SymmetricLieAlgebra.mk}
\end{theorem}

\subsection*{noConfusion}

\begin{theorem}[blueprint="infogeometry:core:generic:symmetricliealgebra:noconfusion"]
  \inputleannode{InfoGeometry.Core.Generic.SymmetricLieAlgebra.noConfusion}
\end{theorem}

\subsection*{noConfusionType}

\begin{theorem}[blueprint="infogeometry:core:generic:symmetricliealgebra:noconfusiontype"]
  \inputleannode{InfoGeometry.Core.Generic.SymmetricLieAlgebra.noConfusionType}
\end{theorem}

\subsection*{rec}

\begin{theorem}[blueprint="infogeometry:core:generic:symmetricliealgebra:rec"]
  \inputleannode{InfoGeometry.Core.Generic.SymmetricLieAlgebra.rec}
\end{theorem}

\subsection*{recOn}

\begin{theorem}[blueprint="infogeometry:core:generic:symmetricliealgebra:recon"]
  \inputleannode{InfoGeometry.Core.Generic.SymmetricLieAlgebra.recOn}
\end{theorem}

\subsection*{triple}

\begin{theorem}[blueprint="infogeometry:core:generic:symmetricliealgebra:triple"]
  \inputleannode{InfoGeometry.Core.Generic.SymmetricLieAlgebra.triple}
\end{theorem}

\subsection*{θ}

\begin{theorem}[blueprint="infogeometry:core:generic:symmetricliealgebra:θ"]
  \inputleannode{InfoGeometry.Core.Generic.SymmetricLieAlgebra.θ}
\end{theorem}

\subsection*{𝔨}

\begin{theorem}[blueprint="infogeometry:core:generic:symmetricliealgebra:𝔨"]
  \inputleannode{InfoGeometry.Core.Generic.SymmetricLieAlgebra.𝔨}
\end{theorem}

\subsection*{𝔭}

\begin{theorem}[blueprint="infogeometry:core:generic:symmetricliealgebra:𝔭"]
  \inputleannode{InfoGeometry.Core.Generic.SymmetricLieAlgebra.𝔭}
\end{theorem}

\section{InfoGeometry.Core.Generic.SymmetricLieAlgebra.CartanLieAlgebra}

\subsection*{casesOn}

\begin{theorem}[blueprint="infogeometry:core:generic:symmetricliealgebra:cartanliealgebra:caseson"]
  \inputleannode{InfoGeometry.Core.Generic.SymmetricLieAlgebra.CartanLieAlgebra.casesOn}
\end{theorem}

\subsection*{ctorIdx}

\begin{theorem}[blueprint="infogeometry:core:generic:symmetricliealgebra:cartanliealgebra:ctoridx"]
  \inputleannode{InfoGeometry.Core.Generic.SymmetricLieAlgebra.CartanLieAlgebra.ctorIdx}
\end{theorem}

\subsection*{killing_nondegenerate}

\begin{theorem}[blueprint="infogeometry:core:generic:symmetricliealgebra:cartanliealgebra:killing_nondegenerate"]
  \inputleannode{InfoGeometry.Core.Generic.SymmetricLieAlgebra.CartanLieAlgebra.killing_nondegenerate}
\end{theorem}

\subsection*{mk}

\begin{theorem}[blueprint="infogeometry:core:generic:symmetricliealgebra:cartanliealgebra:mk"]
  \inputleannode{InfoGeometry.Core.Generic.SymmetricLieAlgebra.CartanLieAlgebra.mk}
\end{theorem}

\subsection*{noConfusion}

\begin{theorem}[blueprint="infogeometry:core:generic:symmetricliealgebra:cartanliealgebra:noconfusion"]
  \inputleannode{InfoGeometry.Core.Generic.SymmetricLieAlgebra.CartanLieAlgebra.noConfusion}
\end{theorem}

\subsection*{noConfusionType}

\begin{theorem}[blueprint="infogeometry:core:generic:symmetricliealgebra:cartanliealgebra:noconfusiontype"]
  \inputleannode{InfoGeometry.Core.Generic.SymmetricLieAlgebra.CartanLieAlgebra.noConfusionType}
\end{theorem}

\subsection*{rec}

\begin{theorem}[blueprint="infogeometry:core:generic:symmetricliealgebra:cartanliealgebra:rec"]
  \inputleannode{InfoGeometry.Core.Generic.SymmetricLieAlgebra.CartanLieAlgebra.rec}
\end{theorem}

\subsection*{recOn}

\begin{theorem}[blueprint="infogeometry:core:generic:symmetricliealgebra:cartanliealgebra:recon"]
  \inputleannode{InfoGeometry.Core.Generic.SymmetricLieAlgebra.CartanLieAlgebra.recOn}
\end{theorem}

\subsection*{toSymmetricLieAlgebra}

\begin{theorem}[blueprint="infogeometry:core:generic:symmetricliealgebra:cartanliealgebra:tosymmetricliealgebra"]
  \inputleannode{InfoGeometry.Core.Generic.SymmetricLieAlgebra.CartanLieAlgebra.toSymmetricLieAlgebra}
\end{theorem}

\section{InfoGeometry.Core.Generic.SymmetricLieAlgebra.CartanSignature}

\subsection*{casesOn}

\begin{theorem}[blueprint="infogeometry:core:generic:symmetricliealgebra:cartansignature:caseson"]
  \inputleannode{InfoGeometry.Core.Generic.SymmetricLieAlgebra.CartanSignature.casesOn}
\end{theorem}

\subsection*{mk}

\begin{theorem}[blueprint="infogeometry:core:generic:symmetricliealgebra:cartansignature:mk"]
  \inputleannode{InfoGeometry.Core.Generic.SymmetricLieAlgebra.CartanSignature.mk}
\end{theorem}

\subsection*{neg_on_k}

\begin{theorem}[blueprint="infogeometry:core:generic:symmetricliealgebra:cartansignature:neg_on_k"]
  \inputleannode{InfoGeometry.Core.Generic.SymmetricLieAlgebra.CartanSignature.neg_on_k}
\end{theorem}

\subsection*{pos_on_p}

\begin{theorem}[blueprint="infogeometry:core:generic:symmetricliealgebra:cartansignature:pos_on_p"]
  \inputleannode{InfoGeometry.Core.Generic.SymmetricLieAlgebra.CartanSignature.pos_on_p}
\end{theorem}

\subsection*{rec}

\begin{theorem}[blueprint="infogeometry:core:generic:symmetricliealgebra:cartansignature:rec"]
  \inputleannode{InfoGeometry.Core.Generic.SymmetricLieAlgebra.CartanSignature.rec}
\end{theorem}

\subsection*{recOn}

\begin{theorem}[blueprint="infogeometry:core:generic:symmetricliealgebra:cartansignature:recon"]
  \inputleannode{InfoGeometry.Core.Generic.SymmetricLieAlgebra.CartanSignature.recOn}
\end{theorem}

\section{InfoGeometry.Core.Generic.SymmetricLieAlgebra.P_minus}

\subsection*{eq_1}

\begin{theorem}[blueprint="infogeometry:core:generic:symmetricliealgebra:p_minus:eq_1"]
  \inputleannode{InfoGeometry.Core.Generic.SymmetricLieAlgebra.P_minus.eq_1}
\end{theorem}

\section{InfoGeometry.Core.Generic.SymmetricLieAlgebra.P_plus}

\subsection*{eq_1}

\begin{theorem}[blueprint="infogeometry:core:generic:symmetricliealgebra:p_plus:eq_1"]
  \inputleannode{InfoGeometry.Core.Generic.SymmetricLieAlgebra.P_plus.eq_1}
\end{theorem}

\section{InfoGeometry.Core.Generic.SymmetricLieAlgebra.𝔨}

\subsection*{eq_1}

\begin{theorem}[blueprint="infogeometry:core:generic:symmetricliealgebra:𝔨:eq_1"]
  \inputleannode{InfoGeometry.Core.Generic.SymmetricLieAlgebra.𝔨.eq_1}
\end{theorem}

\section{InfoGeometry.Core.Generic.SymmetricLieAlgebra.𝔭}

\subsection*{eq_1}

\begin{theorem}[blueprint="infogeometry:core:generic:symmetricliealgebra:𝔭:eq_1"]
  \inputleannode{InfoGeometry.Core.Generic.SymmetricLieAlgebra.𝔭.eq_1}
\end{theorem}

\section{InfoGeometry.Core.GrandCanonical}

\noindent\textbf{Buildable}: yes

\noindent\emph{No exported declarations in the current declaration graph.}

\section{InfoGeometry.Core.Involution}

\noindent\textbf{Buildable}: yes

\noindent\emph{No exported declarations in the current declaration graph.}

\section{InfoGeometry.Core.InvolutiveAutomorphism}

\subsection*{casesOn}

\begin{theorem}[blueprint="infogeometry:core:involutiveautomorphism:caseson"]
  \inputleannode{InfoGeometry.Core.InvolutiveAutomorphism.casesOn}
\end{theorem}

\subsection*{ctorIdx}

\begin{theorem}[blueprint="infogeometry:core:involutiveautomorphism:ctoridx"]
  \inputleannode{InfoGeometry.Core.InvolutiveAutomorphism.ctorIdx}
\end{theorem}

\subsection*{ext}

\begin{theorem}[blueprint="infogeometry:core:involutiveautomorphism:ext"]
  \inputleannode{InfoGeometry.Core.InvolutiveAutomorphism.ext}
\end{theorem}

\subsection*{ext_iff}

\begin{theorem}[blueprint="infogeometry:core:involutiveautomorphism:ext_iff"]
  \inputleannode{InfoGeometry.Core.InvolutiveAutomorphism.ext_iff}
\end{theorem}

\subsection*{involutive}

\begin{theorem}[blueprint="infogeometry:core:involutiveautomorphism:involutive"]
  \inputleannode{InfoGeometry.Core.InvolutiveAutomorphism.involutive}
\end{theorem}

\subsection*{map_add}

\begin{theorem}[blueprint="infogeometry:core:involutiveautomorphism:map_add"]
  \inputleannode{InfoGeometry.Core.InvolutiveAutomorphism.map_add}
\end{theorem}

\subsection*{map_div}

\begin{theorem}[blueprint="infogeometry:core:involutiveautomorphism:map_div"]
  \inputleannode{InfoGeometry.Core.InvolutiveAutomorphism.map_div}
\end{theorem}

\subsection*{map_inv}

\begin{theorem}[blueprint="infogeometry:core:involutiveautomorphism:map_inv"]
  \inputleannode{InfoGeometry.Core.InvolutiveAutomorphism.map_inv}
\end{theorem}

\subsection*{map_lie}

\begin{theorem}[blueprint="infogeometry:core:involutiveautomorphism:map_lie"]
  \inputleannode{InfoGeometry.Core.InvolutiveAutomorphism.map_lie}
\end{theorem}

\subsection*{map_mul}

\begin{theorem}[blueprint="infogeometry:core:involutiveautomorphism:map_mul"]
  \inputleannode{InfoGeometry.Core.InvolutiveAutomorphism.map_mul}
\end{theorem}

\subsection*{map_neg}

\begin{theorem}[blueprint="infogeometry:core:involutiveautomorphism:map_neg"]
  \inputleannode{InfoGeometry.Core.InvolutiveAutomorphism.map_neg}
\end{theorem}

\subsection*{map_one}

\begin{theorem}[blueprint="infogeometry:core:involutiveautomorphism:map_one"]
  \inputleannode{InfoGeometry.Core.InvolutiveAutomorphism.map_one}
\end{theorem}

\subsection*{map_smul}

\begin{theorem}[blueprint="infogeometry:core:involutiveautomorphism:map_smul"]
  \inputleannode{InfoGeometry.Core.InvolutiveAutomorphism.map_smul}
\end{theorem}

\subsection*{map_sub}

\begin{theorem}[blueprint="infogeometry:core:involutiveautomorphism:map_sub"]
  \inputleannode{InfoGeometry.Core.InvolutiveAutomorphism.map_sub}
\end{theorem}

\subsection*{mk}

\begin{theorem}[blueprint="infogeometry:core:involutiveautomorphism:mk"]
  \inputleannode{InfoGeometry.Core.InvolutiveAutomorphism.mk}
\end{theorem}

\subsection*{noConfusion}

\begin{theorem}[blueprint="infogeometry:core:involutiveautomorphism:noconfusion"]
  \inputleannode{InfoGeometry.Core.InvolutiveAutomorphism.noConfusion}
\end{theorem}

\subsection*{noConfusionType}

\begin{theorem}[blueprint="infogeometry:core:involutiveautomorphism:noconfusiontype"]
  \inputleannode{InfoGeometry.Core.InvolutiveAutomorphism.noConfusionType}
\end{theorem}

\subsection*{rec}

\begin{theorem}[blueprint="infogeometry:core:involutiveautomorphism:rec"]
  \inputleannode{InfoGeometry.Core.InvolutiveAutomorphism.rec}
\end{theorem}

\subsection*{recOn}

\begin{theorem}[blueprint="infogeometry:core:involutiveautomorphism:recon"]
  \inputleannode{InfoGeometry.Core.InvolutiveAutomorphism.recOn}
\end{theorem}

\subsection*{toFun}

\begin{theorem}[blueprint="infogeometry:core:involutiveautomorphism:tofun"]
  \inputleannode{InfoGeometry.Core.InvolutiveAutomorphism.toFun}
\end{theorem}

\section{InfoGeometry.Core.Jordan}

\noindent\textbf{Buildable}: yes

\noindent\emph{No exported declarations in the current declaration graph.}

\section{InfoGeometry.Core.LieTripleSystem}

\subsection*{casesOn}

\begin{theorem}[blueprint="infogeometry:core:lietriplesystem:caseson"]
  \inputleannode{InfoGeometry.Core.LieTripleSystem.casesOn}
\end{theorem}

\subsection*{ctorIdx}

\begin{theorem}[blueprint="infogeometry:core:lietriplesystem:ctoridx"]
  \inputleannode{InfoGeometry.Core.LieTripleSystem.ctorIdx}
\end{theorem}

\subsection*{jacobi_like}

\begin{theorem}[blueprint="infogeometry:core:lietriplesystem:jacobi_like"]
  \inputleannode{InfoGeometry.Core.LieTripleSystem.jacobi_like}
\end{theorem}

\subsection*{mk}

\begin{theorem}[blueprint="infogeometry:core:lietriplesystem:mk"]
  \inputleannode{InfoGeometry.Core.LieTripleSystem.mk}
\end{theorem}

\subsection*{noConfusion}

\begin{theorem}[blueprint="infogeometry:core:lietriplesystem:noconfusion"]
  \inputleannode{InfoGeometry.Core.LieTripleSystem.noConfusion}
\end{theorem}

\subsection*{noConfusionType}

\begin{theorem}[blueprint="infogeometry:core:lietriplesystem:noconfusiontype"]
  \inputleannode{InfoGeometry.Core.LieTripleSystem.noConfusionType}
\end{theorem}

\subsection*{rec}

\begin{theorem}[blueprint="infogeometry:core:lietriplesystem:rec"]
  \inputleannode{InfoGeometry.Core.LieTripleSystem.rec}
\end{theorem}

\subsection*{recOn}

\begin{theorem}[blueprint="infogeometry:core:lietriplesystem:recon"]
  \inputleannode{InfoGeometry.Core.LieTripleSystem.recOn}
\end{theorem}

\subsection*{skew₁}

\begin{theorem}[blueprint="infogeometry:core:lietriplesystem:skew₁"]
  \inputleannode{InfoGeometry.Core.LieTripleSystem.skew₁}
\end{theorem}

\subsection*{triple}

\begin{theorem}[blueprint="infogeometry:core:lietriplesystem:triple"]
  \inputleannode{InfoGeometry.Core.LieTripleSystem.triple}
\end{theorem}

\section{InfoGeometry.Core.LinearInvolutiveAutomorphism}

\subsection*{instPreservesLinear}

\begin{theorem}[blueprint="infogeometry:core:linearinvolutiveautomorphism:instpreserveslinear"]
  \inputleannode{InfoGeometry.Core.LinearInvolutiveAutomorphism.instPreservesLinear}
\end{theorem}

\subsection*{involutive}

\begin{theorem}[blueprint="infogeometry:core:linearinvolutiveautomorphism:involutive"]
  \inputleannode{InfoGeometry.Core.LinearInvolutiveAutomorphism.involutive}
\end{theorem}

\subsection*{map_add}

\begin{theorem}[blueprint="infogeometry:core:linearinvolutiveautomorphism:map_add"]
  \inputleannode{InfoGeometry.Core.LinearInvolutiveAutomorphism.map_add}
\end{theorem}

\subsection*{map_neg}

\begin{theorem}[blueprint="infogeometry:core:linearinvolutiveautomorphism:map_neg"]
  \inputleannode{InfoGeometry.Core.LinearInvolutiveAutomorphism.map_neg}
\end{theorem}

\subsection*{map_smul}

\begin{theorem}[blueprint="infogeometry:core:linearinvolutiveautomorphism:map_smul"]
  \inputleannode{InfoGeometry.Core.LinearInvolutiveAutomorphism.map_smul}
\end{theorem}

\subsection*{map_sub}

\begin{theorem}[blueprint="infogeometry:core:linearinvolutiveautomorphism:map_sub"]
  \inputleannode{InfoGeometry.Core.LinearInvolutiveAutomorphism.map_sub}
\end{theorem}

\section{InfoGeometry.Core.MulInvolutiveAutomorphism}

\subsection*{instPreservesMul}

\begin{theorem}[blueprint="infogeometry:core:mulinvolutiveautomorphism:instpreservesmul"]
  \inputleannode{InfoGeometry.Core.MulInvolutiveAutomorphism.instPreservesMul}
\end{theorem}

\subsection*{involutive}

\begin{theorem}[blueprint="infogeometry:core:mulinvolutiveautomorphism:involutive"]
  \inputleannode{InfoGeometry.Core.MulInvolutiveAutomorphism.involutive}
\end{theorem}

\subsection*{map_div}

\begin{theorem}[blueprint="infogeometry:core:mulinvolutiveautomorphism:map_div"]
  \inputleannode{InfoGeometry.Core.MulInvolutiveAutomorphism.map_div}
\end{theorem}

\subsection*{map_inv}

\begin{theorem}[blueprint="infogeometry:core:mulinvolutiveautomorphism:map_inv"]
  \inputleannode{InfoGeometry.Core.MulInvolutiveAutomorphism.map_inv}
\end{theorem}

\subsection*{map_mul}

\begin{theorem}[blueprint="infogeometry:core:mulinvolutiveautomorphism:map_mul"]
  \inputleannode{InfoGeometry.Core.MulInvolutiveAutomorphism.map_mul}
\end{theorem}

\subsection*{map_one}

\begin{theorem}[blueprint="infogeometry:core:mulinvolutiveautomorphism:map_one"]
  \inputleannode{InfoGeometry.Core.MulInvolutiveAutomorphism.map_one}
\end{theorem}

\section{InfoGeometry.Core.OddLocalModel}

\subsection*{casesOn}

\begin{theorem}[blueprint="infogeometry:core:oddlocalmodel:caseson"]
  \inputleannode{InfoGeometry.Core.OddLocalModel.casesOn}
\end{theorem}

\subsection*{convex}

\begin{theorem}[blueprint="infogeometry:core:oddlocalmodel:convex"]
  \inputleannode{InfoGeometry.Core.OddLocalModel.convex}
\end{theorem}

\subsection*{ctorIdx}

\begin{theorem}[blueprint="infogeometry:core:oddlocalmodel:ctoridx"]
  \inputleannode{InfoGeometry.Core.OddLocalModel.ctorIdx}
\end{theorem}

\subsection*{mk}

\begin{theorem}[blueprint="infogeometry:core:oddlocalmodel:mk"]
  \inputleannode{InfoGeometry.Core.OddLocalModel.mk}
\end{theorem}

\subsection*{noConfusion}

\begin{theorem}[blueprint="infogeometry:core:oddlocalmodel:noconfusion"]
  \inputleannode{InfoGeometry.Core.OddLocalModel.noConfusion}
\end{theorem}

\subsection*{noConfusionType}

\begin{theorem}[blueprint="infogeometry:core:oddlocalmodel:noconfusiontype"]
  \inputleannode{InfoGeometry.Core.OddLocalModel.noConfusionType}
\end{theorem}

\subsection*{odd}

\begin{theorem}[blueprint="infogeometry:core:oddlocalmodel:odd"]
  \inputleannode{InfoGeometry.Core.OddLocalModel.odd}
\end{theorem}

\subsection*{pathConnected}

\begin{theorem}[blueprint="infogeometry:core:oddlocalmodel:pathconnected"]
  \inputleannode{InfoGeometry.Core.OddLocalModel.pathConnected}
\end{theorem}

\subsection*{rec}

\begin{theorem}[blueprint="infogeometry:core:oddlocalmodel:rec"]
  \inputleannode{InfoGeometry.Core.OddLocalModel.rec}
\end{theorem}

\subsection*{recOn}

\begin{theorem}[blueprint="infogeometry:core:oddlocalmodel:recon"]
  \inputleannode{InfoGeometry.Core.OddLocalModel.recOn}
\end{theorem}

\subsection*{tripleSystem}

\begin{theorem}[blueprint="infogeometry:core:oddlocalmodel:triplesystem"]
  \inputleannode{InfoGeometry.Core.OddLocalModel.tripleSystem}
\end{theorem}

\section{InfoGeometry.Core.PreservesLie}

\subsection*{casesOn}

\begin{theorem}[blueprint="infogeometry:core:preserveslie:caseson"]
  \inputleannode{InfoGeometry.Core.PreservesLie.casesOn}
\end{theorem}

\subsection*{mk}

\begin{theorem}[blueprint="infogeometry:core:preserveslie:mk"]
  \inputleannode{InfoGeometry.Core.PreservesLie.mk}
\end{theorem}

\subsection*{rec}

\begin{theorem}[blueprint="infogeometry:core:preserveslie:rec"]
  \inputleannode{InfoGeometry.Core.PreservesLie.rec}
\end{theorem}

\subsection*{recOn}

\begin{theorem}[blueprint="infogeometry:core:preserveslie:recon"]
  \inputleannode{InfoGeometry.Core.PreservesLie.recOn}
\end{theorem}

\subsection*{toPreservesLieBracket}

\begin{theorem}[blueprint="infogeometry:core:preserveslie:topreservesliebracket"]
  \inputleannode{InfoGeometry.Core.PreservesLie.toPreservesLieBracket}
\end{theorem}

\subsection*{toPreservesLinear}

\begin{theorem}[blueprint="infogeometry:core:preserveslie:topreserveslinear"]
  \inputleannode{InfoGeometry.Core.PreservesLie.toPreservesLinear}
\end{theorem}

\section{InfoGeometry.Core.PreservesLieBracket}

\subsection*{casesOn}

\begin{theorem}[blueprint="infogeometry:core:preservesliebracket:caseson"]
  \inputleannode{InfoGeometry.Core.PreservesLieBracket.casesOn}
\end{theorem}

\subsection*{map_lie}

\begin{theorem}[blueprint="infogeometry:core:preservesliebracket:map_lie"]
  \inputleannode{InfoGeometry.Core.PreservesLieBracket.map_lie}
\end{theorem}

\subsection*{mk}

\begin{theorem}[blueprint="infogeometry:core:preservesliebracket:mk"]
  \inputleannode{InfoGeometry.Core.PreservesLieBracket.mk}
\end{theorem}

\subsection*{rec}

\begin{theorem}[blueprint="infogeometry:core:preservesliebracket:rec"]
  \inputleannode{InfoGeometry.Core.PreservesLieBracket.rec}
\end{theorem}

\subsection*{recOn}

\begin{theorem}[blueprint="infogeometry:core:preservesliebracket:recon"]
  \inputleannode{InfoGeometry.Core.PreservesLieBracket.recOn}
\end{theorem}

\section{InfoGeometry.Core.PreservesLinear}

\subsection*{casesOn}

\begin{theorem}[blueprint="infogeometry:core:preserveslinear:caseson"]
  \inputleannode{InfoGeometry.Core.PreservesLinear.casesOn}
\end{theorem}

\subsection*{map_add}

\begin{theorem}[blueprint="infogeometry:core:preserveslinear:map_add"]
  \inputleannode{InfoGeometry.Core.PreservesLinear.map_add}
\end{theorem}

\subsection*{map_smul}

\begin{theorem}[blueprint="infogeometry:core:preserveslinear:map_smul"]
  \inputleannode{InfoGeometry.Core.PreservesLinear.map_smul}
\end{theorem}

\subsection*{mk}

\begin{theorem}[blueprint="infogeometry:core:preserveslinear:mk"]
  \inputleannode{InfoGeometry.Core.PreservesLinear.mk}
\end{theorem}

\subsection*{rec}

\begin{theorem}[blueprint="infogeometry:core:preserveslinear:rec"]
  \inputleannode{InfoGeometry.Core.PreservesLinear.rec}
\end{theorem}

\subsection*{recOn}

\begin{theorem}[blueprint="infogeometry:core:preserveslinear:recon"]
  \inputleannode{InfoGeometry.Core.PreservesLinear.recOn}
\end{theorem}

\section{InfoGeometry.Core.PreservesMul}

\subsection*{casesOn}

\begin{theorem}[blueprint="infogeometry:core:preservesmul:caseson"]
  \inputleannode{InfoGeometry.Core.PreservesMul.casesOn}
\end{theorem}

\subsection*{map_mul}

\begin{theorem}[blueprint="infogeometry:core:preservesmul:map_mul"]
  \inputleannode{InfoGeometry.Core.PreservesMul.map_mul}
\end{theorem}

\subsection*{map_one}

\begin{theorem}[blueprint="infogeometry:core:preservesmul:map_one"]
  \inputleannode{InfoGeometry.Core.PreservesMul.map_one}
\end{theorem}

\subsection*{mk}

\begin{theorem}[blueprint="infogeometry:core:preservesmul:mk"]
  \inputleannode{InfoGeometry.Core.PreservesMul.mk}
\end{theorem}

\subsection*{rec}

\begin{theorem}[blueprint="infogeometry:core:preservesmul:rec"]
  \inputleannode{InfoGeometry.Core.PreservesMul.rec}
\end{theorem}

\subsection*{recOn}

\begin{theorem}[blueprint="infogeometry:core:preservesmul:recon"]
  \inputleannode{InfoGeometry.Core.PreservesMul.recOn}
\end{theorem}

\section{InfoGeometry.Core.Projector}

\subsection*{decomposition}

\begin{theorem}[blueprint="infogeometry:core:projector:decomposition"]
  \inputleannode{InfoGeometry.Core.Projector.decomposition}
\end{theorem}

\subsection*{fixed_iff_minus_eq_zero}

\begin{theorem}[blueprint="infogeometry:core:projector:fixed_iff_minus_eq_zero"]
  \inputleannode{InfoGeometry.Core.Projector.fixed_iff_minus_eq_zero}
\end{theorem}

\subsection*{minus}

\begin{theorem}[blueprint="infogeometry:core:projector:minus"]
  \inputleannode{InfoGeometry.Core.Projector.minus}
\end{theorem}

\subsection*{minus_add}

\begin{theorem}[blueprint="infogeometry:core:projector:minus_add"]
  \inputleannode{InfoGeometry.Core.Projector.minus_add}
\end{theorem}

\subsection*{minus_idempotent}

\begin{theorem}[blueprint="infogeometry:core:projector:minus_idempotent"]
  \inputleannode{InfoGeometry.Core.Projector.minus_idempotent}
\end{theorem}

\subsection*{minus_neg_fixed}

\begin{theorem}[blueprint="infogeometry:core:projector:minus_neg_fixed"]
  \inputleannode{InfoGeometry.Core.Projector.minus_neg_fixed}
\end{theorem}

\subsection*{minus_plus}

\begin{theorem}[blueprint="infogeometry:core:projector:minus_plus"]
  \inputleannode{InfoGeometry.Core.Projector.minus_plus}
\end{theorem}

\subsection*{minus_smul}

\begin{theorem}[blueprint="infogeometry:core:projector:minus_smul"]
  \inputleannode{InfoGeometry.Core.Projector.minus_smul}
\end{theorem}

\subsection*{neg_fixed_iff_plus_eq_zero}

\begin{theorem}[blueprint="infogeometry:core:projector:neg_fixed_iff_plus_eq_zero"]
  \inputleannode{InfoGeometry.Core.Projector.neg_fixed_iff_plus_eq_zero}
\end{theorem}

\subsection*{plus}

\begin{theorem}[blueprint="infogeometry:core:projector:plus"]
  \inputleannode{InfoGeometry.Core.Projector.plus}
\end{theorem}

\subsection*{plus_add}

\begin{theorem}[blueprint="infogeometry:core:projector:plus_add"]
  \inputleannode{InfoGeometry.Core.Projector.plus_add}
\end{theorem}

\subsection*{plus_fixed}

\begin{theorem}[blueprint="infogeometry:core:projector:plus_fixed"]
  \inputleannode{InfoGeometry.Core.Projector.plus_fixed}
\end{theorem}

\subsection*{plus_idempotent}

\begin{theorem}[blueprint="infogeometry:core:projector:plus_idempotent"]
  \inputleannode{InfoGeometry.Core.Projector.plus_idempotent}
\end{theorem}

\subsection*{plus_minus}

\begin{theorem}[blueprint="infogeometry:core:projector:plus_minus"]
  \inputleannode{InfoGeometry.Core.Projector.plus_minus}
\end{theorem}

\subsection*{plus_smul}

\begin{theorem}[blueprint="infogeometry:core:projector:plus_smul"]
  \inputleannode{InfoGeometry.Core.Projector.plus_smul}
\end{theorem}

\section{InfoGeometry.Core.Projector.minus}

\subsection*{eq_1}

\begin{theorem}[blueprint="infogeometry:core:projector:minus:eq_1"]
  \inputleannode{InfoGeometry.Core.Projector.minus.eq_1}
\end{theorem}

\section{InfoGeometry.Core.Projector.plus}

\subsection*{eq_1}

\begin{theorem}[blueprint="infogeometry:core:projector:plus:eq_1"]
  \inputleannode{InfoGeometry.Core.Projector.plus.eq_1}
\end{theorem}

\section{InfoGeometry.Core.SymmetricLie}

\noindent\textbf{Buildable}: yes

\noindent\emph{No exported declarations in the current declaration graph.}

\section{InfoGeometry.Core.SymmetricLieAlgebra}

\subsection*{CartanSignature}

\begin{theorem}[blueprint="infogeometry:core:symmetricliealgebra:cartansignature"]
  \inputleannode{InfoGeometry.Core.SymmetricLieAlgebra.CartanSignature}
\end{theorem}

\subsection*{P_minus}

\begin{theorem}[blueprint="infogeometry:core:symmetricliealgebra:p_minus"]
  \inputleannode{InfoGeometry.Core.SymmetricLieAlgebra.P_minus}
\end{theorem}

\subsection*{P_plus}

\begin{theorem}[blueprint="infogeometry:core:symmetricliealgebra:p_plus"]
  \inputleannode{InfoGeometry.Core.SymmetricLieAlgebra.P_plus}
\end{theorem}

\subsection*{add_smul_sub_mem_even}

\begin{theorem}[blueprint="infogeometry:core:symmetricliealgebra:add_smul_sub_mem_even"]
  \inputleannode{InfoGeometry.Core.SymmetricLieAlgebra.add_smul_sub_mem_even}
\end{theorem}

\subsection*{add_smul_sub_mem_odd}

\begin{theorem}[blueprint="infogeometry:core:symmetricliealgebra:add_smul_sub_mem_odd"]
  \inputleannode{InfoGeometry.Core.SymmetricLieAlgebra.add_smul_sub_mem_odd}
\end{theorem}

\subsection*{bracket_k_k}

\begin{theorem}[blueprint="infogeometry:core:symmetricliealgebra:bracket_k_k"]
  \inputleannode{InfoGeometry.Core.SymmetricLieAlgebra.bracket_k_k}
\end{theorem}

\subsection*{bracket_k_p}

\begin{theorem}[blueprint="infogeometry:core:symmetricliealgebra:bracket_k_p"]
  \inputleannode{InfoGeometry.Core.SymmetricLieAlgebra.bracket_k_p}
\end{theorem}

\subsection*{bracket_p_p}

\begin{theorem}[blueprint="infogeometry:core:symmetricliealgebra:bracket_p_p"]
  \inputleannode{InfoGeometry.Core.SymmetricLieAlgebra.bracket_p_p}
\end{theorem}

\subsection*{cartanAssembleLinear}

\begin{theorem}[blueprint="infogeometry:core:symmetricliealgebra:cartanassemblelinear"]
  \inputleannode{InfoGeometry.Core.SymmetricLieAlgebra.cartanAssembleLinear}
\end{theorem}

\subsection*{cartanDecomposeLinear}

\begin{theorem}[blueprint="infogeometry:core:symmetricliealgebra:cartandecomposelinear"]
  \inputleannode{InfoGeometry.Core.SymmetricLieAlgebra.cartanDecomposeLinear}
\end{theorem}

\subsection*{cartanForm}

\begin{theorem}[blueprint="infogeometry:core:symmetricliealgebra:cartanform"]
  \inputleannode{InfoGeometry.Core.SymmetricLieAlgebra.cartanForm}
\end{theorem}

\subsection*{cartanForm_def}

\begin{theorem}[blueprint="infogeometry:core:symmetricliealgebra:cartanform_def"]
  \inputleannode{InfoGeometry.Core.SymmetricLieAlgebra.cartanForm_def}
\end{theorem}

\subsection*{cartanForm_orthogonal}

\begin{theorem}[blueprint="infogeometry:core:symmetricliealgebra:cartanform_orthogonal"]
  \inputleannode{InfoGeometry.Core.SymmetricLieAlgebra.cartanForm_orthogonal}
\end{theorem}

\subsection*{cartanLinearEquiv}

\begin{theorem}[blueprint="infogeometry:core:symmetricliealgebra:cartanlinearequiv"]
  \inputleannode{InfoGeometry.Core.SymmetricLieAlgebra.cartanLinearEquiv}
\end{theorem}

\subsection*{cartan_decomposition}

\begin{theorem}[blueprint="infogeometry:core:symmetricliealgebra:cartan_decomposition"]
  \inputleannode{InfoGeometry.Core.SymmetricLieAlgebra.cartan_decomposition}
\end{theorem}

\subsection*{cartan_direct_sum}

\begin{theorem}[blueprint="infogeometry:core:symmetricliealgebra:cartan_direct_sum"]
  \inputleannode{InfoGeometry.Core.SymmetricLieAlgebra.cartan_direct_sum}
\end{theorem}

\subsection*{cartan_left_inverse}

\begin{theorem}[blueprint="infogeometry:core:symmetricliealgebra:cartan_left_inverse"]
  \inputleannode{InfoGeometry.Core.SymmetricLieAlgebra.cartan_left_inverse}
\end{theorem}

\subsection*{cartan_right_inverse}

\begin{theorem}[blueprint="infogeometry:core:symmetricliealgebra:cartan_right_inverse"]
  \inputleannode{InfoGeometry.Core.SymmetricLieAlgebra.cartan_right_inverse}
\end{theorem}

\subsection*{casesOn}

\begin{theorem}[blueprint="infogeometry:core:symmetricliealgebra:caseson"]
  \inputleannode{InfoGeometry.Core.SymmetricLieAlgebra.casesOn}
\end{theorem}

\subsection*{convex_minusPart_image}

\begin{theorem}[blueprint="infogeometry:core:symmetricliealgebra:convex_minuspart_image"]
  \inputleannode{InfoGeometry.Core.SymmetricLieAlgebra.convex_minusPart_image}
\end{theorem}

\subsection*{convex_minusPart_preimage}

\begin{theorem}[blueprint="infogeometry:core:symmetricliealgebra:convex_minuspart_preimage"]
  \inputleannode{InfoGeometry.Core.SymmetricLieAlgebra.convex_minusPart_preimage}
\end{theorem}

\subsection*{convex_plusPart_image}

\begin{theorem}[blueprint="infogeometry:core:symmetricliealgebra:convex_pluspart_image"]
  \inputleannode{InfoGeometry.Core.SymmetricLieAlgebra.convex_plusPart_image}
\end{theorem}

\subsection*{convex_plusPart_preimage}

\begin{theorem}[blueprint="infogeometry:core:symmetricliealgebra:convex_pluspart_preimage"]
  \inputleannode{InfoGeometry.Core.SymmetricLieAlgebra.convex_plusPart_preimage}
\end{theorem}

\subsection*{ctorIdx}

\begin{theorem}[blueprint="infogeometry:core:symmetricliealgebra:ctoridx"]
  \inputleannode{InfoGeometry.Core.SymmetricLieAlgebra.ctorIdx}
\end{theorem}

\subsection*{curvature}

\begin{theorem}[blueprint="infogeometry:core:symmetricliealgebra:curvature"]
  \inputleannode{InfoGeometry.Core.SymmetricLieAlgebra.curvature}
\end{theorem}

\subsection*{decomposition}

\begin{theorem}[blueprint="infogeometry:core:symmetricliealgebra:decomposition"]
  \inputleannode{InfoGeometry.Core.SymmetricLieAlgebra.decomposition}
\end{theorem}

\subsection*{equivInvolutiveLieAut}

\begin{theorem}[blueprint="infogeometry:core:symmetricliealgebra:equivinvolutivelieaut"]
  \inputleannode{InfoGeometry.Core.SymmetricLieAlgebra.equivInvolutiveLieAut}
\end{theorem}

\subsection*{evenLieSubalgebra}

\begin{theorem}[blueprint="infogeometry:core:symmetricliealgebra:evenliesubalgebra"]
  \inputleannode{InfoGeometry.Core.SymmetricLieAlgebra.evenLieSubalgebra}
\end{theorem}

\subsection*{evenSubmodule}

\begin{theorem}[blueprint="infogeometry:core:symmetricliealgebra:evensubmodule"]
  \inputleannode{InfoGeometry.Core.SymmetricLieAlgebra.evenSubmodule}
\end{theorem}

\subsection*{even_convex}

\begin{theorem}[blueprint="infogeometry:core:symmetricliealgebra:even_convex"]
  \inputleannode{InfoGeometry.Core.SymmetricLieAlgebra.even_convex}
\end{theorem}

\subsection*{even_isConnected}

\begin{theorem}[blueprint="infogeometry:core:symmetricliealgebra:even_isconnected"]
  \inputleannode{InfoGeometry.Core.SymmetricLieAlgebra.even_isConnected}
\end{theorem}

\subsection*{even_isPathConnected}

\begin{theorem}[blueprint="infogeometry:core:symmetricliealgebra:even_ispathconnected"]
  \inputleannode{InfoGeometry.Core.SymmetricLieAlgebra.even_isPathConnected}
\end{theorem}

\subsection*{instPreservesLie}

\begin{theorem}[blueprint="infogeometry:core:symmetricliealgebra:instpreserveslie"]
  \inputleannode{InfoGeometry.Core.SymmetricLieAlgebra.instPreservesLie}
\end{theorem}

\subsection*{instPreservesLieBracket}

\begin{theorem}[blueprint="infogeometry:core:symmetricliealgebra:instpreservesliebracket"]
  \inputleannode{InfoGeometry.Core.SymmetricLieAlgebra.instPreservesLieBracket}
\end{theorem}

\subsection*{instPreservesLinear}

\begin{theorem}[blueprint="infogeometry:core:symmetricliealgebra:instpreserveslinear"]
  \inputleannode{InfoGeometry.Core.SymmetricLieAlgebra.instPreservesLinear}
\end{theorem}

\subsection*{involutive}

\begin{theorem}[blueprint="infogeometry:core:symmetricliealgebra:involutive"]
  \inputleannode{InfoGeometry.Core.SymmetricLieAlgebra.involutive}
\end{theorem}

\subsection*{killing_invariant}

\begin{theorem}[blueprint="infogeometry:core:symmetricliealgebra:killing_invariant"]
  \inputleannode{InfoGeometry.Core.SymmetricLieAlgebra.killing_invariant}
\end{theorem}

\subsection*{killing_orthogonal}

\begin{theorem}[blueprint="infogeometry:core:symmetricliealgebra:killing_orthogonal"]
  \inputleannode{InfoGeometry.Core.SymmetricLieAlgebra.killing_orthogonal}
\end{theorem}

\subsection*{lie}

\begin{theorem}[blueprint="infogeometry:core:symmetricliealgebra:lie"]
  \inputleannode{InfoGeometry.Core.SymmetricLieAlgebra.lie}
\end{theorem}

\subsection*{lin}

\begin{theorem}[blueprint="infogeometry:core:symmetricliealgebra:lin"]
  \inputleannode{InfoGeometry.Core.SymmetricLieAlgebra.lin}
\end{theorem}

\subsection*{mem_evenSubmodule_iff}

\begin{theorem}[blueprint="infogeometry:core:symmetricliealgebra:mem_evensubmodule_iff"]
  \inputleannode{InfoGeometry.Core.SymmetricLieAlgebra.mem_evenSubmodule_iff}
\end{theorem}

\subsection*{mem_even_iff}

\begin{theorem}[blueprint="infogeometry:core:symmetricliealgebra:mem_even_iff"]
  \inputleannode{InfoGeometry.Core.SymmetricLieAlgebra.mem_even_iff}
\end{theorem}

\subsection*{mem_oddSubmodule_iff}

\begin{theorem}[blueprint="infogeometry:core:symmetricliealgebra:mem_oddsubmodule_iff"]
  \inputleannode{InfoGeometry.Core.SymmetricLieAlgebra.mem_oddSubmodule_iff}
\end{theorem}

\subsection*{minusPart}

\begin{theorem}[blueprint="infogeometry:core:symmetricliealgebra:minuspart"]
  \inputleannode{InfoGeometry.Core.SymmetricLieAlgebra.minusPart}
\end{theorem}

\subsection*{minusPartLinear}

\begin{theorem}[blueprint="infogeometry:core:symmetricliealgebra:minuspartlinear"]
  \inputleannode{InfoGeometry.Core.SymmetricLieAlgebra.minusPartLinear}
\end{theorem}

\subsection*{minusPart_add}

\begin{theorem}[blueprint="infogeometry:core:symmetricliealgebra:minuspart_add"]
  \inputleannode{InfoGeometry.Core.SymmetricLieAlgebra.minusPart_add}
\end{theorem}

\subsection*{minusPart_mem_odd}

\begin{theorem}[blueprint="infogeometry:core:symmetricliealgebra:minuspart_mem_odd"]
  \inputleannode{InfoGeometry.Core.SymmetricLieAlgebra.minusPart_mem_odd}
\end{theorem}

\subsection*{minusPart_of_mem_even}

\begin{theorem}[blueprint="infogeometry:core:symmetricliealgebra:minuspart_of_mem_even"]
  \inputleannode{InfoGeometry.Core.SymmetricLieAlgebra.minusPart_of_mem_even}
\end{theorem}

\subsection*{minusPart_of_mem_odd}

\begin{theorem}[blueprint="infogeometry:core:symmetricliealgebra:minuspart_of_mem_odd"]
  \inputleannode{InfoGeometry.Core.SymmetricLieAlgebra.minusPart_of_mem_odd}
\end{theorem}

\subsection*{minusPart_smul}

\begin{theorem}[blueprint="infogeometry:core:symmetricliealgebra:minuspart_smul"]
  \inputleannode{InfoGeometry.Core.SymmetricLieAlgebra.minusPart_smul}
\end{theorem}

\subsection*{mk}

\begin{theorem}[blueprint="infogeometry:core:symmetricliealgebra:mk"]
  \inputleannode{InfoGeometry.Core.SymmetricLieAlgebra.mk}
\end{theorem}

\subsection*{noConfusion}

\begin{theorem}[blueprint="infogeometry:core:symmetricliealgebra:noconfusion"]
  \inputleannode{InfoGeometry.Core.SymmetricLieAlgebra.noConfusion}
\end{theorem}

\subsection*{noConfusionType}

\begin{theorem}[blueprint="infogeometry:core:symmetricliealgebra:noconfusiontype"]
  \inputleannode{InfoGeometry.Core.SymmetricLieAlgebra.noConfusionType}
\end{theorem}

\subsection*{oddLieTripleSystem}

\begin{theorem}[blueprint="infogeometry:core:symmetricliealgebra:oddlietriplesystem"]
  \inputleannode{InfoGeometry.Core.SymmetricLieAlgebra.oddLieTripleSystem}
\end{theorem}

\subsection*{oddLocalModel}

\begin{theorem}[blueprint="infogeometry:core:symmetricliealgebra:oddlocalmodel"]
  \inputleannode{InfoGeometry.Core.SymmetricLieAlgebra.oddLocalModel}
\end{theorem}

\subsection*{oddLocalModel_convex}

\begin{theorem}[blueprint="infogeometry:core:symmetricliealgebra:oddlocalmodel_convex"]
  \inputleannode{InfoGeometry.Core.SymmetricLieAlgebra.oddLocalModel_convex}
\end{theorem}

\subsection*{oddLocalModel_odd}

\begin{theorem}[blueprint="infogeometry:core:symmetricliealgebra:oddlocalmodel_odd"]
  \inputleannode{InfoGeometry.Core.SymmetricLieAlgebra.oddLocalModel_odd}
\end{theorem}

\subsection*{oddLocalModel_pathConnected}

\begin{theorem}[blueprint="infogeometry:core:symmetricliealgebra:oddlocalmodel_pathconnected"]
  \inputleannode{InfoGeometry.Core.SymmetricLieAlgebra.oddLocalModel_pathConnected}
\end{theorem}

\subsection*{oddLocalModel_tripleSystem}

\begin{theorem}[blueprint="infogeometry:core:symmetricliealgebra:oddlocalmodel_triplesystem"]
  \inputleannode{InfoGeometry.Core.SymmetricLieAlgebra.oddLocalModel_tripleSystem}
\end{theorem}

\subsection*{oddSubmodule}

\begin{theorem}[blueprint="infogeometry:core:symmetricliealgebra:oddsubmodule"]
  \inputleannode{InfoGeometry.Core.SymmetricLieAlgebra.oddSubmodule}
\end{theorem}

\subsection*{oddTriple}

\begin{theorem}[blueprint="infogeometry:core:symmetricliealgebra:oddtriple"]
  \inputleannode{InfoGeometry.Core.SymmetricLieAlgebra.oddTriple}
\end{theorem}

\subsection*{oddTriple_jacobi}

\begin{theorem}[blueprint="infogeometry:core:symmetricliealgebra:oddtriple_jacobi"]
  \inputleannode{InfoGeometry.Core.SymmetricLieAlgebra.oddTriple_jacobi}
\end{theorem}

\subsection*{oddTriple_skew₁}

\begin{theorem}[blueprint="infogeometry:core:symmetricliealgebra:oddtriple_skew₁"]
  \inputleannode{InfoGeometry.Core.SymmetricLieAlgebra.oddTriple_skew₁}
\end{theorem}

\subsection*{odd_convex}

\begin{theorem}[blueprint="infogeometry:core:symmetricliealgebra:odd_convex"]
  \inputleannode{InfoGeometry.Core.SymmetricLieAlgebra.odd_convex}
\end{theorem}

\subsection*{odd_isConnected}

\begin{theorem}[blueprint="infogeometry:core:symmetricliealgebra:odd_isconnected"]
  \inputleannode{InfoGeometry.Core.SymmetricLieAlgebra.odd_isConnected}
\end{theorem}

\subsection*{odd_isPathConnected}

\begin{theorem}[blueprint="infogeometry:core:symmetricliealgebra:odd_ispathconnected"]
  \inputleannode{InfoGeometry.Core.SymmetricLieAlgebra.odd_isPathConnected}
\end{theorem}

\subsection*{ofInvolutiveLieAut}

\begin{theorem}[blueprint="infogeometry:core:symmetricliealgebra:ofinvolutivelieaut"]
  \inputleannode{InfoGeometry.Core.SymmetricLieAlgebra.ofInvolutiveLieAut}
\end{theorem}

\subsection*{of_toInvolutiveLieAut}

\begin{theorem}[blueprint="infogeometry:core:symmetricliealgebra:of_toinvolutivelieaut"]
  \inputleannode{InfoGeometry.Core.SymmetricLieAlgebra.of_toInvolutiveLieAut}
\end{theorem}

\subsection*{plusPart}

\begin{theorem}[blueprint="infogeometry:core:symmetricliealgebra:pluspart"]
  \inputleannode{InfoGeometry.Core.SymmetricLieAlgebra.plusPart}
\end{theorem}

\subsection*{plusPartLinear}

\begin{theorem}[blueprint="infogeometry:core:symmetricliealgebra:pluspartlinear"]
  \inputleannode{InfoGeometry.Core.SymmetricLieAlgebra.plusPartLinear}
\end{theorem}

\subsection*{plusPart_add}

\begin{theorem}[blueprint="infogeometry:core:symmetricliealgebra:pluspart_add"]
  \inputleannode{InfoGeometry.Core.SymmetricLieAlgebra.plusPart_add}
\end{theorem}

\subsection*{plusPart_mem_even}

\begin{theorem}[blueprint="infogeometry:core:symmetricliealgebra:pluspart_mem_even"]
  \inputleannode{InfoGeometry.Core.SymmetricLieAlgebra.plusPart_mem_even}
\end{theorem}

\subsection*{plusPart_of_mem_even}

\begin{theorem}[blueprint="infogeometry:core:symmetricliealgebra:pluspart_of_mem_even"]
  \inputleannode{InfoGeometry.Core.SymmetricLieAlgebra.plusPart_of_mem_even}
\end{theorem}

\subsection*{plusPart_of_mem_odd}

\begin{theorem}[blueprint="infogeometry:core:symmetricliealgebra:pluspart_of_mem_odd"]
  \inputleannode{InfoGeometry.Core.SymmetricLieAlgebra.plusPart_of_mem_odd}
\end{theorem}

\subsection*{plusPart_smul}

\begin{theorem}[blueprint="infogeometry:core:symmetricliealgebra:pluspart_smul"]
  \inputleannode{InfoGeometry.Core.SymmetricLieAlgebra.plusPart_smul}
\end{theorem}

\subsection*{rec}

\begin{theorem}[blueprint="infogeometry:core:symmetricliealgebra:rec"]
  \inputleannode{InfoGeometry.Core.SymmetricLieAlgebra.rec}
\end{theorem}

\subsection*{recOn}

\begin{theorem}[blueprint="infogeometry:core:symmetricliealgebra:recon"]
  \inputleannode{InfoGeometry.Core.SymmetricLieAlgebra.recOn}
\end{theorem}

\subsection*{segment_subset_even}

\begin{theorem}[blueprint="infogeometry:core:symmetricliealgebra:segment_subset_even"]
  \inputleannode{InfoGeometry.Core.SymmetricLieAlgebra.segment_subset_even}
\end{theorem}

\subsection*{segment_subset_odd}

\begin{theorem}[blueprint="infogeometry:core:symmetricliealgebra:segment_subset_odd"]
  \inputleannode{InfoGeometry.Core.SymmetricLieAlgebra.segment_subset_odd}
\end{theorem}

\subsection*{symmetric_pair_axioms}

\begin{theorem}[blueprint="infogeometry:core:symmetricliealgebra:symmetric_pair_axioms"]
  \inputleannode{InfoGeometry.Core.SymmetricLieAlgebra.symmetric_pair_axioms}
\end{theorem}

\subsection*{toInvolutiveLieAut}

\begin{theorem}[blueprint="infogeometry:core:symmetricliealgebra:toinvolutivelieaut"]
  \inputleannode{InfoGeometry.Core.SymmetricLieAlgebra.toInvolutiveLieAut}
\end{theorem}

\subsection*{toLieAut}

\begin{theorem}[blueprint="infogeometry:core:symmetricliealgebra:tolieaut"]
  \inputleannode{InfoGeometry.Core.SymmetricLieAlgebra.toLieAut}
\end{theorem}

\subsection*{to_ofInvolutiveLieAut}

\begin{theorem}[blueprint="infogeometry:core:symmetricliealgebra:to_ofinvolutivelieaut"]
  \inputleannode{InfoGeometry.Core.SymmetricLieAlgebra.to_ofInvolutiveLieAut}
\end{theorem}

\subsection*{triple}

\begin{theorem}[blueprint="infogeometry:core:symmetricliealgebra:triple"]
  \inputleannode{InfoGeometry.Core.SymmetricLieAlgebra.triple}
\end{theorem}

\subsection*{tripleOnP}

\begin{theorem}[blueprint="infogeometry:core:symmetricliealgebra:tripleonp"]
  \inputleannode{InfoGeometry.Core.SymmetricLieAlgebra.tripleOnP}
\end{theorem}

\subsection*{triple_closed}

\begin{theorem}[blueprint="infogeometry:core:symmetricliealgebra:triple_closed"]
  \inputleannode{InfoGeometry.Core.SymmetricLieAlgebra.triple_closed}
\end{theorem}

\subsection*{triple_jacobi}

\begin{theorem}[blueprint="infogeometry:core:symmetricliealgebra:triple_jacobi"]
  \inputleannode{InfoGeometry.Core.SymmetricLieAlgebra.triple_jacobi}
\end{theorem}

\subsection*{θ}

\begin{theorem}[blueprint="infogeometry:core:symmetricliealgebra:θ"]
  \inputleannode{InfoGeometry.Core.SymmetricLieAlgebra.θ}
\end{theorem}

\subsection*{𝔨}

\begin{theorem}[blueprint="infogeometry:core:symmetricliealgebra:𝔨"]
  \inputleannode{InfoGeometry.Core.SymmetricLieAlgebra.𝔨}
\end{theorem}

\subsection*{𝔭}

\begin{theorem}[blueprint="infogeometry:core:symmetricliealgebra:𝔭"]
  \inputleannode{InfoGeometry.Core.SymmetricLieAlgebra.𝔭}
\end{theorem}

\section{InfoGeometry.Core.SymmetricLieAlgebra.CartanSignature}

\subsection*{casesOn}

\begin{theorem}[blueprint="infogeometry:core:symmetricliealgebra:cartansignature:caseson"]
  \inputleannode{InfoGeometry.Core.SymmetricLieAlgebra.CartanSignature.casesOn}
\end{theorem}

\subsection*{mk}

\begin{theorem}[blueprint="infogeometry:core:symmetricliealgebra:cartansignature:mk"]
  \inputleannode{InfoGeometry.Core.SymmetricLieAlgebra.CartanSignature.mk}
\end{theorem}

\subsection*{nonneg_on_odd}

\begin{theorem}[blueprint="infogeometry:core:symmetricliealgebra:cartansignature:nonneg_on_odd"]
  \inputleannode{InfoGeometry.Core.SymmetricLieAlgebra.CartanSignature.nonneg_on_odd}
\end{theorem}

\subsection*{nonpos_on_even}

\begin{theorem}[blueprint="infogeometry:core:symmetricliealgebra:cartansignature:nonpos_on_even"]
  \inputleannode{InfoGeometry.Core.SymmetricLieAlgebra.CartanSignature.nonpos_on_even}
\end{theorem}

\subsection*{rec}

\begin{theorem}[blueprint="infogeometry:core:symmetricliealgebra:cartansignature:rec"]
  \inputleannode{InfoGeometry.Core.SymmetricLieAlgebra.CartanSignature.rec}
\end{theorem}

\subsection*{recOn}

\begin{theorem}[blueprint="infogeometry:core:symmetricliealgebra:cartansignature:recon"]
  \inputleannode{InfoGeometry.Core.SymmetricLieAlgebra.CartanSignature.recOn}
\end{theorem}

\section{InfoGeometry.Core.SymmetricLieAlgebra.P_minus}

\subsection*{eq_1}

\begin{theorem}[blueprint="infogeometry:core:symmetricliealgebra:p_minus:eq_1"]
  \inputleannode{InfoGeometry.Core.SymmetricLieAlgebra.P_minus.eq_1}
\end{theorem}

\section{InfoGeometry.Core.SymmetricLieAlgebra.P_plus}

\subsection*{eq_1}

\begin{theorem}[blueprint="infogeometry:core:symmetricliealgebra:p_plus:eq_1"]
  \inputleannode{InfoGeometry.Core.SymmetricLieAlgebra.P_plus.eq_1}
\end{theorem}

\section{InfoGeometry.Core.SymmetricLieAlgebra.cartanAssembleLinear}

\subsection*{eq_1}

\begin{theorem}[blueprint="infogeometry:core:symmetricliealgebra:cartanassemblelinear:eq_1"]
  \inputleannode{InfoGeometry.Core.SymmetricLieAlgebra.cartanAssembleLinear.eq_1}
\end{theorem}

\section{InfoGeometry.Core.SymmetricLieAlgebra.cartanDecomposeLinear}

\subsection*{eq_1}

\begin{theorem}[blueprint="infogeometry:core:symmetricliealgebra:cartandecomposelinear:eq_1"]
  \inputleannode{InfoGeometry.Core.SymmetricLieAlgebra.cartanDecomposeLinear.eq_1}
\end{theorem}

\section{InfoGeometry.Core.SymmetricLieAlgebra.minusPart}

\subsection*{eq_1}

\begin{theorem}[blueprint="infogeometry:core:symmetricliealgebra:minuspart:eq_1"]
  \inputleannode{InfoGeometry.Core.SymmetricLieAlgebra.minusPart.eq_1}
\end{theorem}

\section{InfoGeometry.Core.SymmetricLieAlgebra.minusPartLinear}

\subsection*{eq_1}

\begin{theorem}[blueprint="infogeometry:core:symmetricliealgebra:minuspartlinear:eq_1"]
  \inputleannode{InfoGeometry.Core.SymmetricLieAlgebra.minusPartLinear.eq_1}
\end{theorem}

\section{InfoGeometry.Core.SymmetricLieAlgebra.oddTriple}

\subsection*{eq_1}

\begin{theorem}[blueprint="infogeometry:core:symmetricliealgebra:oddtriple:eq_1"]
  \inputleannode{InfoGeometry.Core.SymmetricLieAlgebra.oddTriple.eq_1}
\end{theorem}

\section{InfoGeometry.Core.SymmetricLieAlgebra.plusPart}

\subsection*{eq_1}

\begin{theorem}[blueprint="infogeometry:core:symmetricliealgebra:pluspart:eq_1"]
  \inputleannode{InfoGeometry.Core.SymmetricLieAlgebra.plusPart.eq_1}
\end{theorem}

\section{InfoGeometry.Core.SymmetricLieAlgebra.plusPartLinear}

\subsection*{eq_1}

\begin{theorem}[blueprint="infogeometry:core:symmetricliealgebra:pluspartlinear:eq_1"]
  \inputleannode{InfoGeometry.Core.SymmetricLieAlgebra.plusPartLinear.eq_1}
\end{theorem}

\section{InfoGeometry.Core.SymmetricLieAlgebra.triple}

\subsection*{eq_1}

\begin{theorem}[blueprint="infogeometry:core:symmetricliealgebra:triple:eq_1"]
  \inputleannode{InfoGeometry.Core.SymmetricLieAlgebra.triple.eq_1}
\end{theorem}

\section{InfoGeometry.Core.SymmetricLieGeneric}

\noindent\textbf{Buildable}: yes

\noindent\emph{No exported declarations in the current declaration graph.}

\section{InfoGeometry.Core.SymmetricSpace}

\subsection*{casesOn}

\begin{theorem}[blueprint="infogeometry:core:symmetricspace:caseson"]
  \inputleannode{InfoGeometry.Core.SymmetricSpace.casesOn}
\end{theorem}

\subsection*{ctorIdx}

\begin{theorem}[blueprint="infogeometry:core:symmetricspace:ctoridx"]
  \inputleannode{InfoGeometry.Core.SymmetricSpace.ctorIdx}
\end{theorem}

\subsection*{ext}

\begin{theorem}[blueprint="infogeometry:core:symmetricspace:ext"]
  \inputleannode{InfoGeometry.Core.SymmetricSpace.ext}
\end{theorem}

\subsection*{ext_iff}

\begin{theorem}[blueprint="infogeometry:core:symmetricspace:ext_iff"]
  \inputleannode{InfoGeometry.Core.SymmetricSpace.ext_iff}
\end{theorem}

\subsection*{mk}

\begin{theorem}[blueprint="infogeometry:core:symmetricspace:mk"]
  \inputleannode{InfoGeometry.Core.SymmetricSpace.mk}
\end{theorem}

\subsection*{noConfusion}

\begin{theorem}[blueprint="infogeometry:core:symmetricspace:noconfusion"]
  \inputleannode{InfoGeometry.Core.SymmetricSpace.noConfusion}
\end{theorem}

\subsection*{noConfusionType}

\begin{theorem}[blueprint="infogeometry:core:symmetricspace:noconfusiontype"]
  \inputleannode{InfoGeometry.Core.SymmetricSpace.noConfusionType}
\end{theorem}

\subsection*{ofLegacy}

\begin{theorem}[blueprint="infogeometry:core:symmetricspace:oflegacy"]
  \inputleannode{InfoGeometry.Core.SymmetricSpace.ofLegacy}
\end{theorem}

\subsection*{ofLegacy_toLegacy}

\begin{theorem}[blueprint="infogeometry:core:symmetricspace:oflegacy_tolegacy"]
  \inputleannode{InfoGeometry.Core.SymmetricSpace.ofLegacy_toLegacy}
\end{theorem}

\subsection*{rec}

\begin{theorem}[blueprint="infogeometry:core:symmetricspace:rec"]
  \inputleannode{InfoGeometry.Core.SymmetricSpace.rec}
\end{theorem}

\subsection*{recOn}

\begin{theorem}[blueprint="infogeometry:core:symmetricspace:recon"]
  \inputleannode{InfoGeometry.Core.SymmetricSpace.recOn}
\end{theorem}

\subsection*{reflection}

\begin{theorem}[blueprint="infogeometry:core:symmetricspace:reflection"]
  \inputleannode{InfoGeometry.Core.SymmetricSpace.reflection}
\end{theorem}

\subsection*{reflection_fix}

\begin{theorem}[blueprint="infogeometry:core:symmetricspace:reflection_fix"]
  \inputleannode{InfoGeometry.Core.SymmetricSpace.reflection_fix}
\end{theorem}

\subsection*{toLegacy}

\begin{theorem}[blueprint="infogeometry:core:symmetricspace:tolegacy"]
  \inputleannode{InfoGeometry.Core.SymmetricSpace.toLegacy}
\end{theorem}

\subsection*{toLegacy_ofLegacy}

\begin{theorem}[blueprint="infogeometry:core:symmetricspace:tolegacy_oflegacy"]
  \inputleannode{InfoGeometry.Core.SymmetricSpace.toLegacy_ofLegacy}
\end{theorem}

\section{InfoGeometry.Core.SymmetricSpaces}

\noindent\textbf{Buildable}: yes

\noindent\emph{No exported declarations in the current declaration graph.}

\section{InfoGeometry.Core.UnifiedGeometry}

\noindent\textbf{Buildable}: yes

\noindent\emph{No exported declarations in the current declaration graph.}

\section{InfoGeometry.Core.cartanMinus}

\subsection*{eq_1}

\begin{theorem}[blueprint="infogeometry:core:cartanminus:eq_1"]
  \inputleannode{InfoGeometry.Core.cartanMinus.eq_1}
\end{theorem}

\section{InfoGeometry.Core.cartanPlus}

\subsection*{eq_1}

\begin{theorem}[blueprint="infogeometry:core:cartanplus:eq_1"]
  \inputleannode{InfoGeometry.Core.cartanPlus.eq_1}
\end{theorem}

\section{InfoGeometry.Core.conjugationEvenSubmodule}

\subsection*{eq_1}

\begin{theorem}[blueprint="infogeometry:core:conjugationevensubmodule:eq_1"]
  \inputleannode{InfoGeometry.Core.conjugationEvenSubmodule.eq_1}
\end{theorem}

\section{InfoGeometry.Core.conjugationLieEquiv}

\subsection*{eq_1}

\begin{theorem}[blueprint="infogeometry:core:conjugationlieequiv:eq_1"]
  \inputleannode{InfoGeometry.Core.conjugationLieEquiv.eq_1}
\end{theorem}

\section{InfoGeometry.Core.conjugationMap}

\subsection*{eq_1}

\begin{theorem}[blueprint="infogeometry:core:conjugationmap:eq_1"]
  \inputleannode{InfoGeometry.Core.conjugationMap.eq_1}
\end{theorem}

\chapter{Cramer Modules}

\section{InfoGeometry.Cramer}

\noindent\textbf{Buildable}: yes

\noindent\emph{No exported declarations in the current declaration graph.}

\chapter{Degree Modules}

\section{InfoGeometry.Degree}

\noindent\textbf{Buildable}: yes

\noindent\emph{No exported declarations in the current declaration graph.}

\chapter{EmpiricalCounts Modules}

\section{InfoGeometry.EmpiricalCounts}

\subsection*{casesOn}

\begin{theorem}[blueprint="infogeometry:empiricalcounts:caseson"]
  \inputleannode{InfoGeometry.EmpiricalCounts.casesOn}
\end{theorem}

\subsection*{count}

\begin{theorem}[blueprint="infogeometry:empiricalcounts:count"]
  \inputleannode{InfoGeometry.EmpiricalCounts.count}
\end{theorem}

\subsection*{ctorIdx}

\begin{theorem}[blueprint="infogeometry:empiricalcounts:ctoridx"]
  \inputleannode{InfoGeometry.EmpiricalCounts.ctorIdx}
\end{theorem}

\subsection*{mk}

\begin{theorem}[blueprint="infogeometry:empiricalcounts:mk"]
  \inputleannode{InfoGeometry.EmpiricalCounts.mk}
\end{theorem}

\subsection*{noConfusion}

\begin{theorem}[blueprint="infogeometry:empiricalcounts:noconfusion"]
  \inputleannode{InfoGeometry.EmpiricalCounts.noConfusion}
\end{theorem}

\subsection*{noConfusionType}

\begin{theorem}[blueprint="infogeometry:empiricalcounts:noconfusiontype"]
  \inputleannode{InfoGeometry.EmpiricalCounts.noConfusionType}
\end{theorem}

\subsection*{rec}

\begin{theorem}[blueprint="infogeometry:empiricalcounts:rec"]
  \inputleannode{InfoGeometry.EmpiricalCounts.rec}
\end{theorem}

\subsection*{recOn}

\begin{theorem}[blueprint="infogeometry:empiricalcounts:recon"]
  \inputleannode{InfoGeometry.EmpiricalCounts.recOn}
\end{theorem}

\chapter{EntropicInference Modules}

\section{InfoGeometry.EntropicInference}

\noindent\textbf{Buildable}: yes

\subsection*{Entropy}

\begin{theorem}[blueprint="infogeometry:entropicinference:entropy"]
  \inputleannode{InfoGeometry.EntropicInference.Entropy}
\end{theorem}

\subsection*{Entropy_le_Entropy_of_eq}

\begin{theorem}[blueprint="infogeometry:entropicinference:entropy_le_entropy_of_eq"]
  \inputleannode{InfoGeometry.EntropicInference.Entropy_le_Entropy_of_eq}
\end{theorem}

\subsection*{FinProb}

\begin{theorem}[blueprint="infogeometry:entropicinference:finprob"]
  \inputleannode{InfoGeometry.EntropicInference.FinProb}
\end{theorem}

\subsection*{Joint}

\begin{theorem}[blueprint="infogeometry:entropicinference:joint"]
  \inputleannode{InfoGeometry.EntropicInference.Joint}
\end{theorem}

\subsection*{KL}

\begin{theorem}[blueprint="infogeometry:entropicinference:kl"]
  \inputleannode{InfoGeometry.EntropicInference.KL}
\end{theorem}

\subsection*{KL_chain_rule}

\begin{theorem}[blueprint="infogeometry:entropicinference:kl_chain_rule"]
  \inputleannode{InfoGeometry.EntropicInference.KL_chain_rule}
\end{theorem}

\subsection*{KL_eq_klDiv}

\begin{theorem}[blueprint="infogeometry:entropicinference:kl_eq_kldiv"]
  \inputleannode{InfoGeometry.EntropicInference.KL_eq_klDiv}
\end{theorem}

\subsection*{KL_nonneg}

\begin{theorem}[blueprint="infogeometry:entropicinference:kl_nonneg"]
  \inputleannode{InfoGeometry.EntropicInference.KL_nonneg}
\end{theorem}

\subsection*{KL_self}

\begin{theorem}[blueprint="infogeometry:entropicinference:kl_self"]
  \inputleannode{InfoGeometry.EntropicInference.KL_self}
\end{theorem}

\subsection*{assemble}

\begin{theorem}[blueprint="infogeometry:entropicinference:assemble"]
  \inputleannode{InfoGeometry.EntropicInference.assemble}
\end{theorem}

\subsection*{assemble_toFun}

\begin{theorem}[blueprint="infogeometry:entropicinference:assemble_tofun"]
  \inputleannode{InfoGeometry.EntropicInference.assemble_toFun}
\end{theorem}

\subsection*{bayesJoint}

\begin{theorem}[blueprint="infogeometry:entropicinference:bayesjoint"]
  \inputleannode{InfoGeometry.EntropicInference.bayesJoint}
\end{theorem}

\subsection*{bayesPosterior}

\begin{theorem}[blueprint="infogeometry:entropicinference:bayesposterior"]
  \inputleannode{InfoGeometry.EntropicInference.bayesPosterior}
\end{theorem}

\subsection*{bayes_is_ME}

\begin{theorem}[blueprint="infogeometry:entropicinference:bayes_is_me"]
  \inputleannode{InfoGeometry.EntropicInference.bayes_is_ME}
\end{theorem}

\subsection*{condΘGivenX}

\begin{theorem}[blueprint="infogeometry:entropicinference:condθgivenx"]
  \inputleannode{InfoGeometry.EntropicInference.condΘGivenX}
\end{theorem}

\subsection*{condΘGivenX_toFun}

\begin{theorem}[blueprint="infogeometry:entropicinference:condθgivenx_tofun"]
  \inputleannode{InfoGeometry.EntropicInference.condΘGivenX_toFun}
\end{theorem}

\subsection*{dirac}

\begin{theorem}[blueprint="infogeometry:entropicinference:dirac"]
  \inputleannode{InfoGeometry.EntropicInference.dirac}
\end{theorem}

\subsection*{factorizedJoint}

\begin{theorem}[blueprint="infogeometry:entropicinference:factorizedjoint"]
  \inputleannode{InfoGeometry.EntropicInference.factorizedJoint}
\end{theorem}

\subsection*{jeffreyJoint}

\begin{theorem}[blueprint="infogeometry:entropicinference:jeffreyjoint"]
  \inputleannode{InfoGeometry.EntropicInference.jeffreyJoint}
\end{theorem}

\subsection*{jeffrey_is_ME}

\begin{theorem}[blueprint="infogeometry:entropicinference:jeffrey_is_me"]
  \inputleannode{InfoGeometry.EntropicInference.jeffrey_is_ME}
\end{theorem}

\subsection*{klTerm}

\begin{theorem}[blueprint="infogeometry:entropicinference:klterm"]
  \inputleannode{InfoGeometry.EntropicInference.klTerm}
\end{theorem}

\subsection*{marginalX}

\begin{theorem}[blueprint="infogeometry:entropicinference:marginalx"]
  \inputleannode{InfoGeometry.EntropicInference.marginalX}
\end{theorem}

\subsection*{marginalX_assemble}

\begin{theorem}[blueprint="infogeometry:entropicinference:marginalx_assemble"]
  \inputleannode{InfoGeometry.EntropicInference.marginalX_assemble}
\end{theorem}

\subsection*{marginalΘ}

\begin{theorem}[blueprint="infogeometry:entropicinference:marginalθ"]
  \inputleannode{InfoGeometry.EntropicInference.marginalΘ}
\end{theorem}

\subsection*{normalize}

\begin{theorem}[blueprint="infogeometry:entropicinference:normalize"]
  \inputleannode{InfoGeometry.EntropicInference.normalize}
\end{theorem}

\subsection*{toProbabilityDist}

\begin{theorem}[blueprint="infogeometry:entropicinference:toprobabilitydist"]
  \inputleannode{InfoGeometry.EntropicInference.toProbabilityDist}
\end{theorem}

\section{InfoGeometry.EntropicInference.Entropy}

\subsection*{eq_1}

\begin{theorem}[blueprint="infogeometry:entropicinference:entropy:eq_1"]
  \inputleannode{InfoGeometry.EntropicInference.Entropy.eq_1}
\end{theorem}

\section{InfoGeometry.EntropicInference.klTerm}

\subsection*{eq_1}

\begin{theorem}[blueprint="infogeometry:entropicinference:klterm:eq_1"]
  \inputleannode{InfoGeometry.EntropicInference.klTerm.eq_1}
\end{theorem}

\section{InfoGeometry.EntropicInference.marginalΘ}

\subsection*{eq_1}

\begin{theorem}[blueprint="infogeometry:entropicinference:marginalθ:eq_1"]
  \inputleannode{InfoGeometry.EntropicInference.marginalΘ.eq_1}
\end{theorem}

\section{InfoGeometry.EntropicInference.toProbabilityDist}

\subsection*{eq_1}

\begin{theorem}[blueprint="infogeometry:entropicinference:toprobabilitydist:eq_1"]
  \inputleannode{InfoGeometry.EntropicInference.toProbabilityDist.eq_1}
\end{theorem}

\chapter{EntropicInferenceTest Modules}

\section{InfoGeometry.EntropicInferenceTest}

\noindent\textbf{Buildable}: yes

\noindent\emph{No exported declarations in the current declaration graph.}

\chapter{ExponentialFamily Modules}

\section{InfoGeometry.ExponentialFamily}

\noindent\textbf{Buildable}: yes

\subsection*{FiniteExponentialFamilyData}

\begin{theorem}[blueprint="infogeometry:exponentialfamily:finiteexponentialfamilydata"]
  \inputleannode{InfoGeometry.ExponentialFamily.FiniteExponentialFamilyData}
\end{theorem}

\subsection*{KLParam}

\begin{theorem}[blueprint="infogeometry:exponentialfamily:klparam"]
  \inputleannode{InfoGeometry.ExponentialFamily.KLParam}
\end{theorem}

\subsection*{KLParam_eq}

\begin{theorem}[blueprint="infogeometry:exponentialfamily:klparam_eq"]
  \inputleannode{InfoGeometry.ExponentialFamily.KLParam_eq}
\end{theorem}

\subsection*{KLParam_eq_LogPotential_bregman_of_statMean_eq_deriv}

\begin{theorem}[blueprint="infogeometry:exponentialfamily:klparam_eq_logpotential_bregman_of_statmean_eq_deriv"]
  \inputleannode{InfoGeometry.ExponentialFamily.KLParam_eq_LogPotential_bregman_of_statMean_eq_deriv}
\end{theorem}

\subsection*{KLParam_eq_bregman_if_deriv_mean}

\begin{theorem}[blueprint="infogeometry:exponentialfamily:klparam_eq_bregman_if_deriv_mean"]
  \inputleannode{InfoGeometry.ExponentialFamily.KLParam_eq_bregman_if_deriv_mean}
\end{theorem}

\subsection*{KLParam_self}

\begin{theorem}[blueprint="infogeometry:exponentialfamily:klparam_self"]
  \inputleannode{InfoGeometry.ExponentialFamily.KLParam_self}
\end{theorem}

\subsection*{entropicTransportObjective}

\begin{theorem}[blueprint="infogeometry:exponentialfamily:entropictransportobjective"]
  \inputleannode{InfoGeometry.ExponentialFamily.entropicTransportObjective}
\end{theorem}

\subsection*{entropicTransportObjective_eq}

\begin{theorem}[blueprint="infogeometry:exponentialfamily:entropictransportobjective_eq"]
  \inputleannode{InfoGeometry.ExponentialFamily.entropicTransportObjective_eq}
\end{theorem}

\subsection*{entropicTransportObjective_eq_potentialGap_if_deriv_mean}

\begin{theorem}[blueprint="infogeometry:exponentialfamily:entropictransportobjective_eq_potentialgap_if_deriv_mean"]
  \inputleannode{InfoGeometry.ExponentialFamily.entropicTransportObjective_eq_potentialGap_if_deriv_mean}
\end{theorem}

\subsection*{entropicTransportPotentialGap}

\begin{theorem}[blueprint="infogeometry:exponentialfamily:entropictransportpotentialgap"]
  \inputleannode{InfoGeometry.ExponentialFamily.entropicTransportPotentialGap}
\end{theorem}

\subsection*{entropicTransportPotentialGap_eq}

\begin{theorem}[blueprint="infogeometry:exponentialfamily:entropictransportpotentialgap_eq"]
  \inputleannode{InfoGeometry.ExponentialFamily.entropicTransportPotentialGap_eq}
\end{theorem}

\subsection*{familyDensity}

\begin{theorem}[blueprint="infogeometry:exponentialfamily:familydensity"]
  \inputleannode{InfoGeometry.ExponentialFamily.familyDensity}
\end{theorem}

\subsection*{familyDensity_eq}

\begin{theorem}[blueprint="infogeometry:exponentialfamily:familydensity_eq"]
  \inputleannode{InfoGeometry.ExponentialFamily.familyDensity_eq}
\end{theorem}

\subsection*{familyDensity_pos}

\begin{theorem}[blueprint="infogeometry:exponentialfamily:familydensity_pos"]
  \inputleannode{InfoGeometry.ExponentialFamily.familyDensity_pos}
\end{theorem}

\subsection*{familyLogPartition}

\begin{theorem}[blueprint="infogeometry:exponentialfamily:familylogpartition"]
  \inputleannode{InfoGeometry.ExponentialFamily.familyLogPartition}
\end{theorem}

\subsection*{familyNormalization}

\begin{theorem}[blueprint="infogeometry:exponentialfamily:familynormalization"]
  \inputleannode{InfoGeometry.ExponentialFamily.familyNormalization}
\end{theorem}

\subsection*{familyPartition}

\begin{theorem}[blueprint="infogeometry:exponentialfamily:familypartition"]
  \inputleannode{InfoGeometry.ExponentialFamily.familyPartition}
\end{theorem}

\subsection*{familyPartition_pos}

\begin{theorem}[blueprint="infogeometry:exponentialfamily:familypartition_pos"]
  \inputleannode{InfoGeometry.ExponentialFamily.familyPartition_pos}
\end{theorem}

\subsection*{familyStatistic}

\begin{theorem}[blueprint="infogeometry:exponentialfamily:familystatistic"]
  \inputleannode{InfoGeometry.ExponentialFamily.familyStatistic}
\end{theorem}

\subsection*{logPotential}

\begin{theorem}[blueprint="infogeometry:exponentialfamily:logpotential"]
  \inputleannode{InfoGeometry.ExponentialFamily.logPotential}
\end{theorem}

\subsection*{logSumExpLegendre}

\begin{theorem}[blueprint="infogeometry:exponentialfamily:logsumexplegendre"]
  \inputleannode{InfoGeometry.ExponentialFamily.logSumExpLegendre}
\end{theorem}

\subsection*{logSumExpLegendre_f}

\begin{theorem}[blueprint="infogeometry:exponentialfamily:logsumexplegendre_f"]
  \inputleannode{InfoGeometry.ExponentialFamily.logSumExpLegendre_f}
\end{theorem}

\subsection*{logSumExpLegendre_smooth}

\begin{theorem}[blueprint="infogeometry:exponentialfamily:logsumexplegendre_smooth"]
  \inputleannode{InfoGeometry.ExponentialFamily.logSumExpLegendre_smooth}
\end{theorem}

\subsection*{logSumExpLegendre_strictConvex}

\begin{theorem}[blueprint="infogeometry:exponentialfamily:logsumexplegendre_strictconvex"]
  \inputleannode{InfoGeometry.ExponentialFamily.logSumExpLegendre_strictConvex}
\end{theorem}

\subsection*{log_density_eq}

\begin{theorem}[blueprint="infogeometry:exponentialfamily:log_density_eq"]
  \inputleannode{InfoGeometry.ExponentialFamily.log_density_eq}
\end{theorem}

\subsection*{statMean}

\begin{theorem}[blueprint="infogeometry:exponentialfamily:statmean"]
  \inputleannode{InfoGeometry.ExponentialFamily.statMean}
\end{theorem}

\subsection*{toFiniteExponentialFamily}

\begin{theorem}[blueprint="infogeometry:exponentialfamily:tofiniteexponentialfamily"]
  \inputleannode{InfoGeometry.ExponentialFamily.toFiniteExponentialFamily}
\end{theorem}

\section{InfoGeometry.ExponentialFamily.Analytic.LogSumExp}

\noindent\textbf{Buildable}: yes

\noindent\emph{No exported declarations in the current declaration graph.}

\section{InfoGeometry.ExponentialFamily.Analytic.Softmax}

\noindent\textbf{Buildable}: yes

\noindent\emph{No exported declarations in the current declaration graph.}

\section{InfoGeometry.ExponentialFamily.Bernoulli}

\noindent\textbf{Buildable}: yes

\subsection*{bernoulliHessianGeometry}

\begin{theorem}[blueprint="infogeometry:exponentialfamily:bernoulli:bernoullihessiangeometry"]
  \inputleannode{InfoGeometry.ExponentialFamily.Bernoulli.bernoulliHessianGeometry}
\end{theorem}

\subsection*{bernoulliTriadicCore}

\begin{theorem}[blueprint="infogeometry:exponentialfamily:bernoulli:bernoullitriadiccore"]
  \inputleannode{InfoGeometry.ExponentialFamily.Bernoulli.bernoulliTriadicCore}
\end{theorem}

\subsection*{bernoulli_divergence_eq_kl}

\begin{theorem}[blueprint="infogeometry:exponentialfamily:bernoulli:bernoulli_divergence_eq_kl"]
  \inputleannode{InfoGeometry.ExponentialFamily.Bernoulli.bernoulli_divergence_eq_kl}
\end{theorem}

\subsection*{bernoulli_metric}

\begin{theorem}[blueprint="infogeometry:exponentialfamily:bernoulli:bernoulli_metric"]
  \inputleannode{InfoGeometry.ExponentialFamily.Bernoulli.bernoulli_metric}
\end{theorem}

\subsection*{deriv_logPartition}

\begin{theorem}[blueprint="infogeometry:exponentialfamily:bernoulli:deriv_logpartition"]
  \inputleannode{InfoGeometry.ExponentialFamily.Bernoulli.deriv_logPartition}
\end{theorem}

\subsection*{instBregmanDivergenceReal}

\begin{theorem}[blueprint="infogeometry:exponentialfamily:bernoulli:instbregmandivergencereal"]
  \inputleannode{InfoGeometry.ExponentialFamily.Bernoulli.instBregmanDivergenceReal}
\end{theorem}

\subsection*{logPartition}

\begin{theorem}[blueprint="infogeometry:exponentialfamily:bernoulli:logpartition"]
  \inputleannode{InfoGeometry.ExponentialFamily.Bernoulli.logPartition}
\end{theorem}

\section{InfoGeometry.ExponentialFamily.Bernoulli.logPartition}

\subsection*{eq_1}

\begin{theorem}[blueprint="infogeometry:exponentialfamily:bernoulli:logpartition:eq_1"]
  \inputleannode{InfoGeometry.ExponentialFamily.Bernoulli.logPartition.eq_1}
\end{theorem}

\section{InfoGeometry.ExponentialFamily.Class}

\noindent\textbf{Buildable}: yes

\noindent\emph{No exported declarations in the current declaration graph.}

\section{InfoGeometry.ExponentialFamily.Finite}

\noindent\textbf{Buildable}: yes

\noindent\emph{No exported declarations in the current declaration graph.}

\section{InfoGeometry.ExponentialFamily.FiniteExponentialFamilyData}

\subsection*{base}

\begin{theorem}[blueprint="infogeometry:exponentialfamily:finiteexponentialfamilydata:base"]
  \inputleannode{InfoGeometry.ExponentialFamily.FiniteExponentialFamilyData.base}
\end{theorem}

\subsection*{base_pos}

\begin{theorem}[blueprint="infogeometry:exponentialfamily:finiteexponentialfamilydata:base_pos"]
  \inputleannode{InfoGeometry.ExponentialFamily.FiniteExponentialFamilyData.base_pos}
\end{theorem}

\subsection*{casesOn}

\begin{theorem}[blueprint="infogeometry:exponentialfamily:finiteexponentialfamilydata:caseson"]
  \inputleannode{InfoGeometry.ExponentialFamily.FiniteExponentialFamilyData.casesOn}
\end{theorem}

\subsection*{ctorIdx}

\begin{theorem}[blueprint="infogeometry:exponentialfamily:finiteexponentialfamilydata:ctoridx"]
  \inputleannode{InfoGeometry.ExponentialFamily.FiniteExponentialFamilyData.ctorIdx}
\end{theorem}

\subsection*{mk}

\begin{theorem}[blueprint="infogeometry:exponentialfamily:finiteexponentialfamilydata:mk"]
  \inputleannode{InfoGeometry.ExponentialFamily.FiniteExponentialFamilyData.mk}
\end{theorem}

\subsection*{noConfusion}

\begin{theorem}[blueprint="infogeometry:exponentialfamily:finiteexponentialfamilydata:noconfusion"]
  \inputleannode{InfoGeometry.ExponentialFamily.FiniteExponentialFamilyData.noConfusion}
\end{theorem}

\subsection*{noConfusionType}

\begin{theorem}[blueprint="infogeometry:exponentialfamily:finiteexponentialfamilydata:noconfusiontype"]
  \inputleannode{InfoGeometry.ExponentialFamily.FiniteExponentialFamilyData.noConfusionType}
\end{theorem}

\subsection*{rec}

\begin{theorem}[blueprint="infogeometry:exponentialfamily:finiteexponentialfamilydata:rec"]
  \inputleannode{InfoGeometry.ExponentialFamily.FiniteExponentialFamilyData.rec}
\end{theorem}

\subsection*{recOn}

\begin{theorem}[blueprint="infogeometry:exponentialfamily:finiteexponentialfamilydata:recon"]
  \inputleannode{InfoGeometry.ExponentialFamily.FiniteExponentialFamilyData.recOn}
\end{theorem}

\subsection*{stat}

\begin{theorem}[blueprint="infogeometry:exponentialfamily:finiteexponentialfamilydata:stat"]
  \inputleannode{InfoGeometry.ExponentialFamily.FiniteExponentialFamilyData.stat}
\end{theorem}

\section{InfoGeometry.ExponentialFamily.Gaussian}

\noindent\textbf{Buildable}: yes

\subsection*{GaussianFamily}

\begin{theorem}[blueprint="infogeometry:exponentialfamily:gaussian:gaussianfamily"]
  \inputleannode{InfoGeometry.ExponentialFamily.Gaussian.GaussianFamily}
\end{theorem}

\section{InfoGeometry.ExponentialFamily.Gaussian.GaussianFamily}

\subsection*{casesOn}

\begin{theorem}[blueprint="infogeometry:exponentialfamily:gaussian:gaussianfamily:caseson"]
  \inputleannode{InfoGeometry.ExponentialFamily.Gaussian.GaussianFamily.casesOn}
\end{theorem}

\subsection*{ctorIdx}

\begin{theorem}[blueprint="infogeometry:exponentialfamily:gaussian:gaussianfamily:ctoridx"]
  \inputleannode{InfoGeometry.ExponentialFamily.Gaussian.GaussianFamily.ctorIdx}
\end{theorem}

\subsection*{divergence_eq_mahalanobis}

\begin{theorem}[blueprint="infogeometry:exponentialfamily:gaussian:gaussianfamily:divergence_eq_mahalanobis"]
  \inputleannode{InfoGeometry.ExponentialFamily.Gaussian.GaussianFamily.divergence_eq_mahalanobis}
\end{theorem}

\subsection*{hasFDerivAt_logPartition}

\begin{theorem}[blueprint="infogeometry:exponentialfamily:gaussian:gaussianfamily:hasfderivat_logpartition"]
  \inputleannode{InfoGeometry.ExponentialFamily.Gaussian.GaussianFamily.hasFDerivAt_logPartition}
\end{theorem}

\subsection*{hessianGeometry}

\begin{theorem}[blueprint="infogeometry:exponentialfamily:gaussian:gaussianfamily:hessiangeometry"]
  \inputleannode{InfoGeometry.ExponentialFamily.Gaussian.GaussianFamily.hessianGeometry}
\end{theorem}

\subsection*{logPartition}

\begin{theorem}[blueprint="infogeometry:exponentialfamily:gaussian:gaussianfamily:logpartition"]
  \inputleannode{InfoGeometry.ExponentialFamily.Gaussian.GaussianFamily.logPartition}
\end{theorem}

\subsection*{mk}

\begin{theorem}[blueprint="infogeometry:exponentialfamily:gaussian:gaussianfamily:mk"]
  \inputleannode{InfoGeometry.ExponentialFamily.Gaussian.GaussianFamily.mk}
\end{theorem}

\subsection*{noConfusion}

\begin{theorem}[blueprint="infogeometry:exponentialfamily:gaussian:gaussianfamily:noconfusion"]
  \inputleannode{InfoGeometry.ExponentialFamily.Gaussian.GaussianFamily.noConfusion}
\end{theorem}

\subsection*{noConfusionType}

\begin{theorem}[blueprint="infogeometry:exponentialfamily:gaussian:gaussianfamily:noconfusiontype"]
  \inputleannode{InfoGeometry.ExponentialFamily.Gaussian.GaussianFamily.noConfusionType}
\end{theorem}

\subsection*{rec}

\begin{theorem}[blueprint="infogeometry:exponentialfamily:gaussian:gaussianfamily:rec"]
  \inputleannode{InfoGeometry.ExponentialFamily.Gaussian.GaussianFamily.rec}
\end{theorem}

\subsection*{recOn}

\begin{theorem}[blueprint="infogeometry:exponentialfamily:gaussian:gaussianfamily:recon"]
  \inputleannode{InfoGeometry.ExponentialFamily.Gaussian.GaussianFamily.recOn}
\end{theorem}

\subsection*{sigma}

\begin{theorem}[blueprint="infogeometry:exponentialfamily:gaussian:gaussianfamily:sigma"]
  \inputleannode{InfoGeometry.ExponentialFamily.Gaussian.GaussianFamily.sigma}
\end{theorem}

\subsection*{sigma_pos}

\begin{theorem}[blueprint="infogeometry:exponentialfamily:gaussian:gaussianfamily:sigma_pos"]
  \inputleannode{InfoGeometry.ExponentialFamily.Gaussian.GaussianFamily.sigma_pos}
\end{theorem}

\subsection*{sigma_symm}

\begin{theorem}[blueprint="infogeometry:exponentialfamily:gaussian:gaussianfamily:sigma_symm"]
  \inputleannode{InfoGeometry.ExponentialFamily.Gaussian.GaussianFamily.sigma_symm}
\end{theorem}

\subsection*{softmaxGaussianAttention}

\begin{theorem}[blueprint="infogeometry:exponentialfamily:gaussian:gaussianfamily:softmaxgaussianattention"]
  \inputleannode{InfoGeometry.ExponentialFamily.Gaussian.GaussianFamily.softmaxGaussianAttention}
\end{theorem}

\section{InfoGeometry.ExponentialFamily.Gaussian.GaussianFamily.hessianGeometry}

\subsection*{eq_1}

\begin{theorem}[blueprint="infogeometry:exponentialfamily:gaussian:gaussianfamily:hessiangeometry:eq_1"]
  \inputleannode{InfoGeometry.ExponentialFamily.Gaussian.GaussianFamily.hessianGeometry.eq_1}
\end{theorem}

\section{InfoGeometry.ExponentialFamily.Gaussian.GaussianFamily.logPartition}

\subsection*{eq_1}

\begin{theorem}[blueprint="infogeometry:exponentialfamily:gaussian:gaussianfamily:logpartition:eq_1"]
  \inputleannode{InfoGeometry.ExponentialFamily.Gaussian.GaussianFamily.logPartition.eq_1}
\end{theorem}

\section{InfoGeometry.ExponentialFamily.GaussianHolonomy}

\noindent\textbf{Buildable}: yes

\subsection*{gaussianDiracField}

\begin{theorem}[blueprint="infogeometry:exponentialfamily:gaussianholonomy:gaussiandiracfield"]
  \inputleannode{InfoGeometry.ExponentialFamily.GaussianHolonomy.gaussianDiracField}
\end{theorem}

\subsection*{gaussianWilsonLoopDiscrete}

\begin{theorem}[blueprint="infogeometry:exponentialfamily:gaussianholonomy:gaussianwilsonloopdiscrete"]
  \inputleannode{InfoGeometry.ExponentialFamily.GaussianHolonomy.gaussianWilsonLoopDiscrete}
\end{theorem}

\subsection*{gaussian_holonomy_flat}

\begin{theorem}[blueprint="infogeometry:exponentialfamily:gaussianholonomy:gaussian_holonomy_flat"]
  \inputleannode{InfoGeometry.ExponentialFamily.GaussianHolonomy.gaussian_holonomy_flat}
\end{theorem}

\section{InfoGeometry.ExponentialFamily.GaussianLadder}

\noindent\textbf{Buildable}: yes

\subsection*{gaussianAnnihilation}

\begin{theorem}[blueprint="infogeometry:exponentialfamily:gaussianladder:gaussianannihilation"]
  \inputleannode{InfoGeometry.ExponentialFamily.GaussianLadder.gaussianAnnihilation}
\end{theorem}

\subsection*{gaussianCreation}

\begin{theorem}[blueprint="infogeometry:exponentialfamily:gaussianladder:gaussiancreation"]
  \inputleannode{InfoGeometry.ExponentialFamily.GaussianLadder.gaussianCreation}
\end{theorem}

\subsection*{gaussian_ccr}

\begin{theorem}[blueprint="infogeometry:exponentialfamily:gaussianladder:gaussian_ccr"]
  \inputleannode{InfoGeometry.ExponentialFamily.GaussianLadder.gaussian_ccr}
\end{theorem}

\section{InfoGeometry.ExponentialFamily.GaussianLadder.gaussianAnnihilation}

\subsection*{eq_1}

\begin{theorem}[blueprint="infogeometry:exponentialfamily:gaussianladder:gaussianannihilation:eq_1"]
  \inputleannode{InfoGeometry.ExponentialFamily.GaussianLadder.gaussianAnnihilation.eq_1}
\end{theorem}

\section{InfoGeometry.ExponentialFamily.GaussianLadder.gaussianCreation}

\subsection*{eq_1}

\begin{theorem}[blueprint="infogeometry:exponentialfamily:gaussianladder:gaussiancreation:eq_1"]
  \inputleannode{InfoGeometry.ExponentialFamily.GaussianLadder.gaussianCreation.eq_1}
\end{theorem}

\section{InfoGeometry.ExponentialFamily.KLBregman}

\noindent\textbf{Buildable}: yes

\noindent\emph{No exported declarations in the current declaration graph.}

\section{InfoGeometry.ExponentialFamily.Legendre}

\noindent\textbf{Buildable}: yes

\noindent\emph{No exported declarations in the current declaration graph.}

\section{InfoGeometry.ExponentialFamily.TwistedGaussian}

\noindent\textbf{Buildable}: yes

\subsection*{TwistedGaussianFamily}

\begin{theorem}[blueprint="infogeometry:exponentialfamily:twistedgaussian:twistedgaussianfamily"]
  \inputleannode{InfoGeometry.ExponentialFamily.TwistedGaussian.TwistedGaussianFamily}
\end{theorem}

\section{InfoGeometry.ExponentialFamily.TwistedGaussian.TwistedGaussianFamily}

\subsection*{T}

\begin{theorem}[blueprint="infogeometry:exponentialfamily:twistedgaussian:twistedgaussianfamily:t"]
  \inputleannode{InfoGeometry.ExponentialFamily.TwistedGaussian.TwistedGaussianFamily.T}
\end{theorem}

\subsection*{T_antisymm}

\begin{theorem}[blueprint="infogeometry:exponentialfamily:twistedgaussian:twistedgaussianfamily:t_antisymm"]
  \inputleannode{InfoGeometry.ExponentialFamily.TwistedGaussian.TwistedGaussianFamily.T_antisymm}
\end{theorem}

\subsection*{adjoint_T_eq_neg}

\begin{theorem}[blueprint="infogeometry:exponentialfamily:twistedgaussian:twistedgaussianfamily:adjoint_t_eq_neg"]
  \inputleannode{InfoGeometry.ExponentialFamily.TwistedGaussian.TwistedGaussianFamily.adjoint_T_eq_neg}
\end{theorem}

\subsection*{adjoint_infoOperator}

\begin{theorem}[blueprint="infogeometry:exponentialfamily:twistedgaussian:twistedgaussianfamily:adjoint_infooperator"]
  \inputleannode{InfoGeometry.ExponentialFamily.TwistedGaussian.TwistedGaussianFamily.adjoint_infoOperator}
\end{theorem}

\subsection*{casesOn}

\begin{theorem}[blueprint="infogeometry:exponentialfamily:twistedgaussian:twistedgaussianfamily:caseson"]
  \inputleannode{InfoGeometry.ExponentialFamily.TwistedGaussian.TwistedGaussianFamily.casesOn}
\end{theorem}

\subsection*{ctorIdx}

\begin{theorem}[blueprint="infogeometry:exponentialfamily:twistedgaussian:twistedgaussianfamily:ctoridx"]
  \inputleannode{InfoGeometry.ExponentialFamily.TwistedGaussian.TwistedGaussianFamily.ctorIdx}
\end{theorem}

\subsection*{infoOperator}

\begin{theorem}[blueprint="infogeometry:exponentialfamily:twistedgaussian:twistedgaussianfamily:infooperator"]
  \inputleannode{InfoGeometry.ExponentialFamily.TwistedGaussian.TwistedGaussianFamily.infoOperator}
\end{theorem}

\subsection*{mk}

\begin{theorem}[blueprint="infogeometry:exponentialfamily:twistedgaussian:twistedgaussianfamily:mk"]
  \inputleannode{InfoGeometry.ExponentialFamily.TwistedGaussian.TwistedGaussianFamily.mk}
\end{theorem}

\subsection*{noConfusion}

\begin{theorem}[blueprint="infogeometry:exponentialfamily:twistedgaussian:twistedgaussianfamily:noconfusion"]
  \inputleannode{InfoGeometry.ExponentialFamily.TwistedGaussian.TwistedGaussianFamily.noConfusion}
\end{theorem}

\subsection*{noConfusionType}

\begin{theorem}[blueprint="infogeometry:exponentialfamily:twistedgaussian:twistedgaussianfamily:noconfusiontype"]
  \inputleannode{InfoGeometry.ExponentialFamily.TwistedGaussian.TwistedGaussianFamily.noConfusionType}
\end{theorem}

\subsection*{normal_commutes}

\begin{theorem}[blueprint="infogeometry:exponentialfamily:twistedgaussian:twistedgaussianfamily:normal_commutes"]
  \inputleannode{InfoGeometry.ExponentialFamily.TwistedGaussian.TwistedGaussianFamily.normal_commutes}
\end{theorem}

\subsection*{normal_iff_commutes}

\begin{theorem}[blueprint="infogeometry:exponentialfamily:twistedgaussian:twistedgaussianfamily:normal_iff_commutes"]
  \inputleannode{InfoGeometry.ExponentialFamily.TwistedGaussian.TwistedGaussianFamily.normal_iff_commutes}
\end{theorem}

\subsection*{rec}

\begin{theorem}[blueprint="infogeometry:exponentialfamily:twistedgaussian:twistedgaussianfamily:rec"]
  \inputleannode{InfoGeometry.ExponentialFamily.TwistedGaussian.TwistedGaussianFamily.rec}
\end{theorem}

\subsection*{recOn}

\begin{theorem}[blueprint="infogeometry:exponentialfamily:twistedgaussian:twistedgaussianfamily:recon"]
  \inputleannode{InfoGeometry.ExponentialFamily.TwistedGaussian.TwistedGaussianFamily.recOn}
\end{theorem}

\subsection*{toGaussianFamily}

\begin{theorem}[blueprint="infogeometry:exponentialfamily:twistedgaussian:twistedgaussianfamily:togaussianfamily"]
  \inputleannode{InfoGeometry.ExponentialFamily.TwistedGaussian.TwistedGaussianFamily.toGaussianFamily}
\end{theorem}

\subsection*{twistedChiralScale}

\begin{theorem}[blueprint="infogeometry:exponentialfamily:twistedgaussian:twistedgaussianfamily:twistedchiralscale"]
  \inputleannode{InfoGeometry.ExponentialFamily.TwistedGaussian.TwistedGaussianFamily.twistedChiralScale}
\end{theorem}

\section{InfoGeometry.ExponentialFamily.TwistedGaussian.TwistedGaussianFamily.infoOperator}

\subsection*{eq_1}

\begin{theorem}[blueprint="infogeometry:exponentialfamily:twistedgaussian:twistedgaussianfamily:infooperator:eq_1"]
  \inputleannode{InfoGeometry.ExponentialFamily.TwistedGaussian.TwistedGaussianFamily.infoOperator.eq_1}
\end{theorem}

\section{InfoGeometry.ExponentialFamily.familyPartition}

\subsection*{eq_1}

\begin{theorem}[blueprint="infogeometry:exponentialfamily:familypartition:eq_1"]
  \inputleannode{InfoGeometry.ExponentialFamily.familyPartition.eq_1}
\end{theorem}

\chapter{Fenchel Modules}

\section{InfoGeometry.Fenchel}

\noindent\textbf{Buildable}: yes

\noindent\emph{No exported declarations in the current declaration graph.}

\chapter{FinProb Modules}

\section{InfoGeometry.FinProb}

\subsection*{casesOn}

\begin{theorem}[blueprint="infogeometry:finprob:caseson"]
  \inputleannode{InfoGeometry.FinProb.casesOn}
\end{theorem}

\subsection*{ctorIdx}

\begin{theorem}[blueprint="infogeometry:finprob:ctoridx"]
  \inputleannode{InfoGeometry.FinProb.ctorIdx}
\end{theorem}

\subsection*{exists_pos}

\begin{theorem}[blueprint="infogeometry:finprob:exists_pos"]
  \inputleannode{InfoGeometry.FinProb.exists_pos}
\end{theorem}

\subsection*{ext}

\begin{theorem}[blueprint="infogeometry:finprob:ext"]
  \inputleannode{InfoGeometry.FinProb.ext}
\end{theorem}

\subsection*{ext_iff}

\begin{theorem}[blueprint="infogeometry:finprob:ext_iff"]
  \inputleannode{InfoGeometry.FinProb.ext_iff}
\end{theorem}

\subsection*{mk}

\begin{theorem}[blueprint="infogeometry:finprob:mk"]
  \inputleannode{InfoGeometry.FinProb.mk}
\end{theorem}

\subsection*{noConfusion}

\begin{theorem}[blueprint="infogeometry:finprob:noconfusion"]
  \inputleannode{InfoGeometry.FinProb.noConfusion}
\end{theorem}

\subsection*{noConfusionType}

\begin{theorem}[blueprint="infogeometry:finprob:noconfusiontype"]
  \inputleannode{InfoGeometry.FinProb.noConfusionType}
\end{theorem}

\subsection*{nonneg}

\begin{theorem}[blueprint="infogeometry:finprob:nonneg"]
  \inputleannode{InfoGeometry.FinProb.nonneg}
\end{theorem}

\subsection*{rec}

\begin{theorem}[blueprint="infogeometry:finprob:rec"]
  \inputleannode{InfoGeometry.FinProb.rec}
\end{theorem}

\subsection*{recOn}

\begin{theorem}[blueprint="infogeometry:finprob:recon"]
  \inputleannode{InfoGeometry.FinProb.recOn}
\end{theorem}

\subsection*{sum_one}

\begin{theorem}[blueprint="infogeometry:finprob:sum_one"]
  \inputleannode{InfoGeometry.FinProb.sum_one}
\end{theorem}

\subsection*{toFun}

\begin{theorem}[blueprint="infogeometry:finprob:tofun"]
  \inputleannode{InfoGeometry.FinProb.toFun}
\end{theorem}

\subsection*{toProbabilityDist}

\begin{theorem}[blueprint="infogeometry:finprob:toprobabilitydist"]
  \inputleannode{InfoGeometry.FinProb.toProbabilityDist}
\end{theorem}

\subsection*{toProbabilityDist_prob}

\begin{theorem}[blueprint="infogeometry:finprob:toprobabilitydist_prob"]
  \inputleannode{InfoGeometry.FinProb.toProbabilityDist_prob}
\end{theorem}

\subsection*{toProbabilityDist_toFinProb}

\begin{theorem}[blueprint="infogeometry:finprob:toprobabilitydist_tofinprob"]
  \inputleannode{InfoGeometry.FinProb.toProbabilityDist_toFinProb}
\end{theorem}

\chapter{Geometry Modules}

\section{InfoGeometry.Geometry}

\subsection*{Doubled}

\begin{theorem}[blueprint="infogeometry:geometry:doubled"]
  \inputleannode{InfoGeometry.Geometry.Doubled}
\end{theorem}

\subsection*{FenchelYoungEquality}

\begin{theorem}[blueprint="infogeometry:geometry:fenchelyoungequality"]
  \inputleannode{InfoGeometry.Geometry.FenchelYoungEquality}
\end{theorem}

\subsection*{GradientBijection}

\begin{theorem}[blueprint="infogeometry:geometry:gradientbijection"]
  \inputleannode{InfoGeometry.Geometry.GradientBijection}
\end{theorem}

\subsection*{HessianGeometry}

\begin{theorem}[blueprint="infogeometry:geometry:hessiangeometry"]
  \inputleannode{InfoGeometry.Geometry.HessianGeometry}
\end{theorem}

\subsection*{HessianManifold}

\begin{theorem}[blueprint="infogeometry:geometry:hessianmanifold"]
  \inputleannode{InfoGeometry.Geometry.HessianManifold}
\end{theorem}

\subsection*{HilbertHessian}

\begin{theorem}[blueprint="infogeometry:geometry:hilberthessian"]
  \inputleannode{InfoGeometry.Geometry.HilbertHessian}
\end{theorem}

\subsection*{IsFenchelConjugate}

\begin{theorem}[blueprint="infogeometry:geometry:isfenchelconjugate"]
  \inputleannode{InfoGeometry.Geometry.IsFenchelConjugate}
\end{theorem}

\subsection*{IsFenchelMajorized}

\begin{theorem}[blueprint="infogeometry:geometry:isfenchelmajorized"]
  \inputleannode{InfoGeometry.Geometry.IsFenchelMajorized}
\end{theorem}

\subsection*{IsLegendreConjugate}

\begin{theorem}[blueprint="infogeometry:geometry:islegendreconjugate"]
  \inputleannode{InfoGeometry.Geometry.IsLegendreConjugate}
\end{theorem}

\subsection*{LegendreInvolutionAssumptions}

\begin{theorem}[blueprint="infogeometry:geometry:legendreinvolutionassumptions"]
  \inputleannode{InfoGeometry.Geometry.LegendreInvolutionAssumptions}
\end{theorem}

\subsection*{LegendreWellPosed}

\begin{theorem}[blueprint="infogeometry:geometry:legendrewellposed"]
  \inputleannode{InfoGeometry.Geometry.LegendreWellPosed}
\end{theorem}

\subsection*{convexOn_univ_iff_eGeodesicConvex}

\begin{theorem}[blueprint="infogeometry:geometry:convexon_univ_iff_egeodesicconvex"]
  \inputleannode{InfoGeometry.Geometry.convexOn_univ_iff_eGeodesicConvex}
\end{theorem}

\subsection*{convex_iff_eGeodesicConvex}

\begin{theorem}[blueprint="infogeometry:geometry:convex_iff_egeodesicconvex"]
  \inputleannode{InfoGeometry.Geometry.convex_iff_eGeodesicConvex}
\end{theorem}

\subsection*{divergence}

\begin{theorem}[blueprint="infogeometry:geometry:divergence"]
  \inputleannode{InfoGeometry.Geometry.divergence}
\end{theorem}

\subsection*{divergenceVec}

\begin{theorem}[blueprint="infogeometry:geometry:divergencevec"]
  \inputleannode{InfoGeometry.Geometry.divergenceVec}
\end{theorem}

\subsection*{divergenceVec_pythagorean}

\begin{theorem}[blueprint="infogeometry:geometry:divergencevec_pythagorean"]
  \inputleannode{InfoGeometry.Geometry.divergenceVec_pythagorean}
\end{theorem}

\subsection*{divergenceVec_self}

\begin{theorem}[blueprint="infogeometry:geometry:divergencevec_self"]
  \inputleannode{InfoGeometry.Geometry.divergenceVec_self}
\end{theorem}

\subsection*{divergenceVec_three_point}

\begin{theorem}[blueprint="infogeometry:geometry:divergencevec_three_point"]
  \inputleannode{InfoGeometry.Geometry.divergenceVec_three_point}
\end{theorem}

\subsection*{divergence_eq_divergenceVec}

\begin{theorem}[blueprint="infogeometry:geometry:divergence_eq_divergencevec"]
  \inputleannode{InfoGeometry.Geometry.divergence_eq_divergenceVec}
\end{theorem}

\subsection*{divergence_pythagorean}

\begin{theorem}[blueprint="infogeometry:geometry:divergence_pythagorean"]
  \inputleannode{InfoGeometry.Geometry.divergence_pythagorean}
\end{theorem}

\subsection*{divergence_self}

\begin{theorem}[blueprint="infogeometry:geometry:divergence_self"]
  \inputleannode{InfoGeometry.Geometry.divergence_self}
\end{theorem}

\subsection*{divergence_three_point}

\begin{theorem}[blueprint="infogeometry:geometry:divergence_three_point"]
  \inputleannode{InfoGeometry.Geometry.divergence_three_point}
\end{theorem}

\subsection*{dualCoord}

\begin{theorem}[blueprint="infogeometry:geometry:dualcoord"]
  \inputleannode{InfoGeometry.Geometry.dualCoord}
\end{theorem}

\subsection*{dualCoordVec}

\begin{theorem}[blueprint="infogeometry:geometry:dualcoordvec"]
  \inputleannode{InfoGeometry.Geometry.dualCoordVec}
\end{theorem}

\subsection*{dualCoordVec_spec}

\begin{theorem}[blueprint="infogeometry:geometry:dualcoordvec_spec"]
  \inputleannode{InfoGeometry.Geometry.dualCoordVec_spec}
\end{theorem}

\subsection*{dualCoord_mGeodesic}

\begin{theorem}[blueprint="infogeometry:geometry:dualcoord_mgeodesic"]
  \inputleannode{InfoGeometry.Geometry.dualCoord_mGeodesic}
\end{theorem}

\subsection*{dual_value_of_fenchelYoungEquality}

\begin{theorem}[blueprint="infogeometry:geometry:dual_value_of_fenchelyoungequality"]
  \inputleannode{InfoGeometry.Geometry.dual_value_of_fenchelYoungEquality}
\end{theorem}

\subsection*{eGeodesic}

\begin{theorem}[blueprint="infogeometry:geometry:egeodesic"]
  \inputleannode{InfoGeometry.Geometry.eGeodesic}
\end{theorem}

\subsection*{eGeodesicConvex}

\begin{theorem}[blueprint="infogeometry:geometry:egeodesicconvex"]
  \inputleannode{InfoGeometry.Geometry.eGeodesicConvex}
\end{theorem}

\subsection*{eGeodesic_one}

\begin{theorem}[blueprint="infogeometry:geometry:egeodesic_one"]
  \inputleannode{InfoGeometry.Geometry.eGeodesic_one}
\end{theorem}

\subsection*{eGeodesic_zero}

\begin{theorem}[blueprint="infogeometry:geometry:egeodesic_zero"]
  \inputleannode{InfoGeometry.Geometry.eGeodesic_zero}
\end{theorem}

\subsection*{fenchelGap}

\begin{theorem}[blueprint="infogeometry:geometry:fenchelgap"]
  \inputleannode{InfoGeometry.Geometry.fenchelGap}
\end{theorem}

\subsection*{fenchelGap_nonneg}

\begin{theorem}[blueprint="infogeometry:geometry:fenchelgap_nonneg"]
  \inputleannode{InfoGeometry.Geometry.fenchelGap_nonneg}
\end{theorem}

\subsection*{fenchelYoungEquality_iff_gap_eq_zero}

\begin{theorem}[blueprint="infogeometry:geometry:fenchelyoungequality_iff_gap_eq_zero"]
  \inputleannode{InfoGeometry.Geometry.fenchelYoungEquality_iff_gap_eq_zero}
\end{theorem}

\subsection*{fenchelYoungEquality_of_dual_value}

\begin{theorem}[blueprint="infogeometry:geometry:fenchelyoungequality_of_dual_value"]
  \inputleannode{InfoGeometry.Geometry.fenchelYoungEquality_of_dual_value}
\end{theorem}

\subsection*{fenchelYoung_ineq}

\begin{theorem}[blueprint="infogeometry:geometry:fenchelyoung_ineq"]
  \inputleannode{InfoGeometry.Geometry.fenchelYoung_ineq}
\end{theorem}

\subsection*{hessian}

\begin{theorem}[blueprint="infogeometry:geometry:hessian"]
  \inputleannode{InfoGeometry.Geometry.hessian}
\end{theorem}

\subsection*{kreinForm}

\begin{theorem}[blueprint="infogeometry:geometry:kreinform"]
  \inputleannode{InfoGeometry.Geometry.kreinForm}
\end{theorem}

\subsection*{kreinForm_explicit}

\begin{theorem}[blueprint="infogeometry:geometry:kreinform_explicit"]
  \inputleannode{InfoGeometry.Geometry.kreinForm_explicit}
\end{theorem}

\subsection*{kreinForm_self}

\begin{theorem}[blueprint="infogeometry:geometry:kreinform_self"]
  \inputleannode{InfoGeometry.Geometry.kreinForm_self}
\end{theorem}

\subsection*{kreinForm_symm}

\begin{theorem}[blueprint="infogeometry:geometry:kreinform_symm"]
  \inputleannode{InfoGeometry.Geometry.kreinForm_symm}
\end{theorem}

\subsection*{kreinGrad}

\begin{theorem}[blueprint="infogeometry:geometry:kreingrad"]
  \inputleannode{InfoGeometry.Geometry.kreinGrad}
\end{theorem}

\subsection*{kreinGrad_eq_hessian_apply}

\begin{theorem}[blueprint="infogeometry:geometry:kreingrad_eq_hessian_apply"]
  \inputleannode{InfoGeometry.Geometry.kreinGrad_eq_hessian_apply}
\end{theorem}

\subsection*{kreinHessian}

\begin{theorem}[blueprint="infogeometry:geometry:kreinhessian"]
  \inputleannode{InfoGeometry.Geometry.kreinHessian}
\end{theorem}

\subsection*{kreinHessian_apply}

\begin{theorem}[blueprint="infogeometry:geometry:kreinhessian_apply"]
  \inputleannode{InfoGeometry.Geometry.kreinHessian_apply}
\end{theorem}

\subsection*{kreinHessian_eq_spectralEpsilon}

\begin{theorem}[blueprint="infogeometry:geometry:kreinhessian_eq_spectralepsilon"]
  \inputleannode{InfoGeometry.Geometry.kreinHessian_eq_spectralEpsilon}
\end{theorem}

\subsection*{kreinHessian_sq}

\begin{theorem}[blueprint="infogeometry:geometry:kreinhessian_sq"]
  \inputleannode{InfoGeometry.Geometry.kreinHessian_sq}
\end{theorem}

\subsection*{kreinPotential}

\begin{theorem}[blueprint="infogeometry:geometry:kreinpotential"]
  \inputleannode{InfoGeometry.Geometry.kreinPotential}
\end{theorem}

\subsection*{le_legendre_of_bddAbove}

\begin{theorem}[blueprint="infogeometry:geometry:le_legendre_of_bddabove"]
  \inputleannode{InfoGeometry.Geometry.le_legendre_of_bddAbove}
\end{theorem}

\subsection*{legendre}

\begin{theorem}[blueprint="infogeometry:geometry:legendre"]
  \inputleannode{InfoGeometry.Geometry.legendre}
\end{theorem}

\subsection*{legendreSupport}

\begin{theorem}[blueprint="infogeometry:geometry:legendresupport"]
  \inputleannode{InfoGeometry.Geometry.legendreSupport}
\end{theorem}

\subsection*{legendre_involution_assumption_theorem}

\begin{theorem}[blueprint="infogeometry:geometry:legendre_involution_assumption_theorem"]
  \inputleannode{InfoGeometry.Geometry.legendre_involution_assumption_theorem}
\end{theorem}

\subsection*{legendre_involution_gap_theorem}

\begin{theorem}[blueprint="infogeometry:geometry:legendre_involution_gap_theorem"]
  \inputleannode{InfoGeometry.Geometry.legendre_involution_gap_theorem}
\end{theorem}

\subsection*{mGeodesic}

\begin{theorem}[blueprint="infogeometry:geometry:mgeodesic"]
  \inputleannode{InfoGeometry.Geometry.mGeodesic}
\end{theorem}

\subsection*{mGeodesic_one}

\begin{theorem}[blueprint="infogeometry:geometry:mgeodesic_one"]
  \inputleannode{InfoGeometry.Geometry.mGeodesic_one}
\end{theorem}

\subsection*{mGeodesic_zero}

\begin{theorem}[blueprint="infogeometry:geometry:mgeodesic_zero"]
  \inputleannode{InfoGeometry.Geometry.mGeodesic_zero}
\end{theorem}

\subsection*{metric}

\begin{theorem}[blueprint="infogeometry:geometry:metric"]
  \inputleannode{InfoGeometry.Geometry.metric}
\end{theorem}

\subsection*{two_mul_kreinPotential}

\begin{theorem}[blueprint="infogeometry:geometry:two_mul_kreinpotential"]
  \inputleannode{InfoGeometry.Geometry.two_mul_kreinPotential}
\end{theorem}

\section{InfoGeometry.Geometry.ConvexOn}

\subsection*{eGeodesicConvex}

\begin{theorem}[blueprint="infogeometry:geometry:convexon:egeodesicconvex"]
  \inputleannode{InfoGeometry.Geometry.ConvexOn.eGeodesicConvex}
\end{theorem}

\section{InfoGeometry.Geometry.DualFlat}

\noindent\textbf{Buildable}: yes

\subsection*{DualFlatStructure}

\begin{theorem}[blueprint="infogeometry:geometry:dualflat:dualflatstructure"]
  \inputleannode{InfoGeometry.Geometry.DualFlat.DualFlatStructure}
\end{theorem}

\subsection*{bregman_pythagorean}

\begin{theorem}[blueprint="infogeometry:geometry:dualflat:bregman_pythagorean"]
  \inputleannode{InfoGeometry.Geometry.DualFlat.bregman_pythagorean}
\end{theorem}

\subsection*{divergence}

\begin{theorem}[blueprint="infogeometry:geometry:dualflat:divergence"]
  \inputleannode{InfoGeometry.Geometry.DualFlat.divergence}
\end{theorem}

\subsection*{divergence_eq_bregmanDiv_real}

\begin{theorem}[blueprint="infogeometry:geometry:dualflat:divergence_eq_bregmandiv_real"]
  \inputleannode{InfoGeometry.Geometry.DualFlat.divergence_eq_bregmanDiv_real}
\end{theorem}

\subsection*{divergence_self}

\begin{theorem}[blueprint="infogeometry:geometry:dualflat:divergence_self"]
  \inputleannode{InfoGeometry.Geometry.DualFlat.divergence_self}
\end{theorem}

\subsection*{dualCoord_apply_sub_eq_deriv_mul}

\begin{theorem}[blueprint="infogeometry:geometry:dualflat:dualcoord_apply_sub_eq_deriv_mul"]
  \inputleannode{InfoGeometry.Geometry.DualFlat.dualCoord_apply_sub_eq_deriv_mul}
\end{theorem}

\subsection*{nabla}

\begin{theorem}[blueprint="infogeometry:geometry:dualflat:nabla"]
  \inputleannode{InfoGeometry.Geometry.DualFlat.nabla}
\end{theorem}

\subsection*{three_point_identity}

\begin{theorem}[blueprint="infogeometry:geometry:dualflat:three_point_identity"]
  \inputleannode{InfoGeometry.Geometry.DualFlat.three_point_identity}
\end{theorem}

\subsection*{toHessianGeometry}

\begin{theorem}[blueprint="infogeometry:geometry:dualflat:tohessiangeometry"]
  \inputleannode{InfoGeometry.Geometry.DualFlat.toHessianGeometry}
\end{theorem}

\section{InfoGeometry.Geometry.DualFlat.DualFlatStructure}

\subsection*{casesOn}

\begin{theorem}[blueprint="infogeometry:geometry:dualflat:dualflatstructure:caseson"]
  \inputleannode{InfoGeometry.Geometry.DualFlat.DualFlatStructure.casesOn}
\end{theorem}

\subsection*{ctorIdx}

\begin{theorem}[blueprint="infogeometry:geometry:dualflat:dualflatstructure:ctoridx"]
  \inputleannode{InfoGeometry.Geometry.DualFlat.DualFlatStructure.ctorIdx}
\end{theorem}

\subsection*{h_convex}

\begin{theorem}[blueprint="infogeometry:geometry:dualflat:dualflatstructure:h_convex"]
  \inputleannode{InfoGeometry.Geometry.DualFlat.DualFlatStructure.h_convex}
\end{theorem}

\subsection*{h_diff}

\begin{theorem}[blueprint="infogeometry:geometry:dualflat:dualflatstructure:h_diff"]
  \inputleannode{InfoGeometry.Geometry.DualFlat.DualFlatStructure.h_diff}
\end{theorem}

\subsection*{mk}

\begin{theorem}[blueprint="infogeometry:geometry:dualflat:dualflatstructure:mk"]
  \inputleannode{InfoGeometry.Geometry.DualFlat.DualFlatStructure.mk}
\end{theorem}

\subsection*{noConfusion}

\begin{theorem}[blueprint="infogeometry:geometry:dualflat:dualflatstructure:noconfusion"]
  \inputleannode{InfoGeometry.Geometry.DualFlat.DualFlatStructure.noConfusion}
\end{theorem}

\subsection*{noConfusionType}

\begin{theorem}[blueprint="infogeometry:geometry:dualflat:dualflatstructure:noconfusiontype"]
  \inputleannode{InfoGeometry.Geometry.DualFlat.DualFlatStructure.noConfusionType}
\end{theorem}

\subsection*{rec}

\begin{theorem}[blueprint="infogeometry:geometry:dualflat:dualflatstructure:rec"]
  \inputleannode{InfoGeometry.Geometry.DualFlat.DualFlatStructure.rec}
\end{theorem}

\subsection*{recOn}

\begin{theorem}[blueprint="infogeometry:geometry:dualflat:dualflatstructure:recon"]
  \inputleannode{InfoGeometry.Geometry.DualFlat.DualFlatStructure.recOn}
\end{theorem}

\subsection*{ψ}

\begin{theorem}[blueprint="infogeometry:geometry:dualflat:dualflatstructure:ψ"]
  \inputleannode{InfoGeometry.Geometry.DualFlat.DualFlatStructure.ψ}
\end{theorem}

\section{InfoGeometry.Geometry.DualFlat.ProjectiveBridge}

\subsection*{ProjectiveProbability}

\begin{theorem}[blueprint="infogeometry:geometry:dualflat:projectivebridge:projectiveprobability"]
  \inputleannode{InfoGeometry.Geometry.DualFlat.ProjectiveBridge.ProjectiveProbability}
\end{theorem}

\subsection*{UnnormalizedMeasure}

\begin{theorem}[blueprint="infogeometry:geometry:dualflat:projectivebridge:unnormalizedmeasure"]
  \inputleannode{InfoGeometry.Geometry.DualFlat.ProjectiveBridge.UnnormalizedMeasure}
\end{theorem}

\subsection*{bayesian_update_pythagorean}

\begin{theorem}[blueprint="infogeometry:geometry:dualflat:projectivebridge:bayesian_update_pythagorean"]
  \inputleannode{InfoGeometry.Geometry.DualFlat.ProjectiveBridge.bayesian_update_pythagorean}
\end{theorem}

\subsection*{bayesian_update_pythagorean_bregman}

\begin{theorem}[blueprint="infogeometry:geometry:dualflat:projectivebridge:bayesian_update_pythagorean_bregman"]
  \inputleannode{InfoGeometry.Geometry.DualFlat.ProjectiveBridge.bayesian_update_pythagorean_bregman}
\end{theorem}

\subsection*{bayesian_update_pythagorean_point}

\begin{theorem}[blueprint="infogeometry:geometry:dualflat:projectivebridge:bayesian_update_pythagorean_point"]
  \inputleannode{InfoGeometry.Geometry.DualFlat.ProjectiveBridge.bayesian_update_pythagorean_point}
\end{theorem}

\subsection*{projectiveDivergence}

\begin{theorem}[blueprint="infogeometry:geometry:dualflat:projectivebridge:projectivedivergence"]
  \inputleannode{InfoGeometry.Geometry.DualFlat.ProjectiveBridge.projectiveDivergence}
\end{theorem}

\subsection*{projectiveDivergence_eq_kl}

\begin{theorem}[blueprint="infogeometry:geometry:dualflat:projectivebridge:projectivedivergence_eq_kl"]
  \inputleannode{InfoGeometry.Geometry.DualFlat.ProjectiveBridge.projectiveDivergence_eq_kl}
\end{theorem}

\subsection*{projectiveDivergence_eq_of_representative_eq}

\begin{theorem}[blueprint="infogeometry:geometry:dualflat:projectivebridge:projectivedivergence_eq_of_representative_eq"]
  \inputleannode{InfoGeometry.Geometry.DualFlat.ProjectiveBridge.projectiveDivergence_eq_of_representative_eq}
\end{theorem}

\section{InfoGeometry.Geometry.DualFlat.ProjectiveBridge.UnnormalizedMeasure}

\subsection*{casesOn}

\begin{theorem}[blueprint="infogeometry:geometry:dualflat:projectivebridge:unnormalizedmeasure:caseson"]
  \inputleannode{InfoGeometry.Geometry.DualFlat.ProjectiveBridge.UnnormalizedMeasure.casesOn}
\end{theorem}

\subsection*{ctorIdx}

\begin{theorem}[blueprint="infogeometry:geometry:dualflat:projectivebridge:unnormalizedmeasure:ctoridx"]
  \inputleannode{InfoGeometry.Geometry.DualFlat.ProjectiveBridge.UnnormalizedMeasure.ctorIdx}
\end{theorem}

\subsection*{mass_pos}

\begin{theorem}[blueprint="infogeometry:geometry:dualflat:projectivebridge:unnormalizedmeasure:mass_pos"]
  \inputleannode{InfoGeometry.Geometry.DualFlat.ProjectiveBridge.UnnormalizedMeasure.mass_pos}
\end{theorem}

\subsection*{mk}

\begin{theorem}[blueprint="infogeometry:geometry:dualflat:projectivebridge:unnormalizedmeasure:mk"]
  \inputleannode{InfoGeometry.Geometry.DualFlat.ProjectiveBridge.UnnormalizedMeasure.mk}
\end{theorem}

\subsection*{noConfusion}

\begin{theorem}[blueprint="infogeometry:geometry:dualflat:projectivebridge:unnormalizedmeasure:noconfusion"]
  \inputleannode{InfoGeometry.Geometry.DualFlat.ProjectiveBridge.UnnormalizedMeasure.noConfusion}
\end{theorem}

\subsection*{noConfusionType}

\begin{theorem}[blueprint="infogeometry:geometry:dualflat:projectivebridge:unnormalizedmeasure:noconfusiontype"]
  \inputleannode{InfoGeometry.Geometry.DualFlat.ProjectiveBridge.UnnormalizedMeasure.noConfusionType}
\end{theorem}

\subsection*{rec}

\begin{theorem}[blueprint="infogeometry:geometry:dualflat:projectivebridge:unnormalizedmeasure:rec"]
  \inputleannode{InfoGeometry.Geometry.DualFlat.ProjectiveBridge.UnnormalizedMeasure.rec}
\end{theorem}

\subsection*{recOn}

\begin{theorem}[blueprint="infogeometry:geometry:dualflat:projectivebridge:unnormalizedmeasure:recon"]
  \inputleannode{InfoGeometry.Geometry.DualFlat.ProjectiveBridge.UnnormalizedMeasure.recOn}
\end{theorem}

\subsection*{representative}

\begin{theorem}[blueprint="infogeometry:geometry:dualflat:projectivebridge:unnormalizedmeasure:representative"]
  \inputleannode{InfoGeometry.Geometry.DualFlat.ProjectiveBridge.UnnormalizedMeasure.representative}
\end{theorem}

\subsection*{totalMass}

\begin{theorem}[blueprint="infogeometry:geometry:dualflat:projectivebridge:unnormalizedmeasure:totalmass"]
  \inputleannode{InfoGeometry.Geometry.DualFlat.ProjectiveBridge.UnnormalizedMeasure.totalMass}
\end{theorem}

\section{InfoGeometry.Geometry.DualFlat.divergence}

\subsection*{congr_simp}

\begin{theorem}[blueprint="infogeometry:geometry:dualflat:divergence:congr_simp"]
  \inputleannode{InfoGeometry.Geometry.DualFlat.divergence.congr_simp}
\end{theorem}

\subsection*{eq_1}

\begin{theorem}[blueprint="infogeometry:geometry:dualflat:divergence:eq_1"]
  \inputleannode{InfoGeometry.Geometry.DualFlat.divergence.eq_1}
\end{theorem}

\section{InfoGeometry.Geometry.DualFlat.nabla}

\subsection*{congr_simp}

\begin{theorem}[blueprint="infogeometry:geometry:dualflat:nabla:congr_simp"]
  \inputleannode{InfoGeometry.Geometry.DualFlat.nabla.congr_simp}
\end{theorem}

\subsection*{eq_1}

\begin{theorem}[blueprint="infogeometry:geometry:dualflat:nabla:eq_1"]
  \inputleannode{InfoGeometry.Geometry.DualFlat.nabla.eq_1}
\end{theorem}

\section{InfoGeometry.Geometry.DualFlat.toHessianGeometry}

\subsection*{eq_1}

\begin{theorem}[blueprint="infogeometry:geometry:dualflat:tohessiangeometry:eq_1"]
  \inputleannode{InfoGeometry.Geometry.DualFlat.toHessianGeometry.eq_1}
\end{theorem}

\section{InfoGeometry.Geometry.GradientBijection}

\subsection*{casesOn}

\begin{theorem}[blueprint="infogeometry:geometry:gradientbijection:caseson"]
  \inputleannode{InfoGeometry.Geometry.GradientBijection.casesOn}
\end{theorem}

\subsection*{ctorIdx}

\begin{theorem}[blueprint="infogeometry:geometry:gradientbijection:ctoridx"]
  \inputleannode{InfoGeometry.Geometry.GradientBijection.ctorIdx}
\end{theorem}

\subsection*{invDual}

\begin{theorem}[blueprint="infogeometry:geometry:gradientbijection:invdual"]
  \inputleannode{InfoGeometry.Geometry.GradientBijection.invDual}
\end{theorem}

\subsection*{left_inv}

\begin{theorem}[blueprint="infogeometry:geometry:gradientbijection:left_inv"]
  \inputleannode{InfoGeometry.Geometry.GradientBijection.left_inv}
\end{theorem}

\subsection*{mk}

\begin{theorem}[blueprint="infogeometry:geometry:gradientbijection:mk"]
  \inputleannode{InfoGeometry.Geometry.GradientBijection.mk}
\end{theorem}

\subsection*{noConfusion}

\begin{theorem}[blueprint="infogeometry:geometry:gradientbijection:noconfusion"]
  \inputleannode{InfoGeometry.Geometry.GradientBijection.noConfusion}
\end{theorem}

\subsection*{noConfusionType}

\begin{theorem}[blueprint="infogeometry:geometry:gradientbijection:noconfusiontype"]
  \inputleannode{InfoGeometry.Geometry.GradientBijection.noConfusionType}
\end{theorem}

\subsection*{rec}

\begin{theorem}[blueprint="infogeometry:geometry:gradientbijection:rec"]
  \inputleannode{InfoGeometry.Geometry.GradientBijection.rec}
\end{theorem}

\subsection*{recOn}

\begin{theorem}[blueprint="infogeometry:geometry:gradientbijection:recon"]
  \inputleannode{InfoGeometry.Geometry.GradientBijection.recOn}
\end{theorem}

\subsection*{right_inv}

\begin{theorem}[blueprint="infogeometry:geometry:gradientbijection:right_inv"]
  \inputleannode{InfoGeometry.Geometry.GradientBijection.right_inv}
\end{theorem}

\section{InfoGeometry.Geometry.HessianGeometry}

\subsection*{casesOn}

\begin{theorem}[blueprint="infogeometry:geometry:hessiangeometry:caseson"]
  \inputleannode{InfoGeometry.Geometry.HessianGeometry.casesOn}
\end{theorem}

\subsection*{ctorIdx}

\begin{theorem}[blueprint="infogeometry:geometry:hessiangeometry:ctoridx"]
  \inputleannode{InfoGeometry.Geometry.HessianGeometry.ctorIdx}
\end{theorem}

\subsection*{mk}

\begin{theorem}[blueprint="infogeometry:geometry:hessiangeometry:mk"]
  \inputleannode{InfoGeometry.Geometry.HessianGeometry.mk}
\end{theorem}

\subsection*{noConfusion}

\begin{theorem}[blueprint="infogeometry:geometry:hessiangeometry:noconfusion"]
  \inputleannode{InfoGeometry.Geometry.HessianGeometry.noConfusion}
\end{theorem}

\subsection*{noConfusionType}

\begin{theorem}[blueprint="infogeometry:geometry:hessiangeometry:noconfusiontype"]
  \inputleannode{InfoGeometry.Geometry.HessianGeometry.noConfusionType}
\end{theorem}

\subsection*{rec}

\begin{theorem}[blueprint="infogeometry:geometry:hessiangeometry:rec"]
  \inputleannode{InfoGeometry.Geometry.HessianGeometry.rec}
\end{theorem}

\subsection*{recOn}

\begin{theorem}[blueprint="infogeometry:geometry:hessiangeometry:recon"]
  \inputleannode{InfoGeometry.Geometry.HessianGeometry.recOn}
\end{theorem}

\subsection*{ψ}

\begin{theorem}[blueprint="infogeometry:geometry:hessiangeometry:ψ"]
  \inputleannode{InfoGeometry.Geometry.HessianGeometry.ψ}
\end{theorem}

\section{InfoGeometry.Geometry.HessianManifold}

\subsection*{casesOn}

\begin{theorem}[blueprint="infogeometry:geometry:hessianmanifold:caseson"]
  \inputleannode{InfoGeometry.Geometry.HessianManifold.casesOn}
\end{theorem}

\subsection*{ctorIdx}

\begin{theorem}[blueprint="infogeometry:geometry:hessianmanifold:ctoridx"]
  \inputleannode{InfoGeometry.Geometry.HessianManifold.ctorIdx}
\end{theorem}

\subsection*{metric}

\begin{theorem}[blueprint="infogeometry:geometry:hessianmanifold:metric"]
  \inputleannode{InfoGeometry.Geometry.HessianManifold.metric}
\end{theorem}

\subsection*{mk}

\begin{theorem}[blueprint="infogeometry:geometry:hessianmanifold:mk"]
  \inputleannode{InfoGeometry.Geometry.HessianManifold.mk}
\end{theorem}

\subsection*{noConfusion}

\begin{theorem}[blueprint="infogeometry:geometry:hessianmanifold:noconfusion"]
  \inputleannode{InfoGeometry.Geometry.HessianManifold.noConfusion}
\end{theorem}

\subsection*{noConfusionType}

\begin{theorem}[blueprint="infogeometry:geometry:hessianmanifold:noconfusiontype"]
  \inputleannode{InfoGeometry.Geometry.HessianManifold.noConfusionType}
\end{theorem}

\subsection*{rec}

\begin{theorem}[blueprint="infogeometry:geometry:hessianmanifold:rec"]
  \inputleannode{InfoGeometry.Geometry.HessianManifold.rec}
\end{theorem}

\subsection*{recOn}

\begin{theorem}[blueprint="infogeometry:geometry:hessianmanifold:recon"]
  \inputleannode{InfoGeometry.Geometry.HessianManifold.recOn}
\end{theorem}

\subsection*{twiceDiff}

\begin{theorem}[blueprint="infogeometry:geometry:hessianmanifold:twicediff"]
  \inputleannode{InfoGeometry.Geometry.HessianManifold.twiceDiff}
\end{theorem}

\subsection*{ψ}

\begin{theorem}[blueprint="infogeometry:geometry:hessianmanifold:ψ"]
  \inputleannode{InfoGeometry.Geometry.HessianManifold.ψ}
\end{theorem}

\section{InfoGeometry.Geometry.HilbertHessian}

\subsection*{casesOn}

\begin{theorem}[blueprint="infogeometry:geometry:hilberthessian:caseson"]
  \inputleannode{InfoGeometry.Geometry.HilbertHessian.casesOn}
\end{theorem}

\subsection*{ctorIdx}

\begin{theorem}[blueprint="infogeometry:geometry:hilberthessian:ctoridx"]
  \inputleannode{InfoGeometry.Geometry.HilbertHessian.ctorIdx}
\end{theorem}

\subsection*{metric}

\begin{theorem}[blueprint="infogeometry:geometry:hilberthessian:metric"]
  \inputleannode{InfoGeometry.Geometry.HilbertHessian.metric}
\end{theorem}

\subsection*{mk}

\begin{theorem}[blueprint="infogeometry:geometry:hilberthessian:mk"]
  \inputleannode{InfoGeometry.Geometry.HilbertHessian.mk}
\end{theorem}

\subsection*{noConfusion}

\begin{theorem}[blueprint="infogeometry:geometry:hilberthessian:noconfusion"]
  \inputleannode{InfoGeometry.Geometry.HilbertHessian.noConfusion}
\end{theorem}

\subsection*{noConfusionType}

\begin{theorem}[blueprint="infogeometry:geometry:hilberthessian:noconfusiontype"]
  \inputleannode{InfoGeometry.Geometry.HilbertHessian.noConfusionType}
\end{theorem}

\subsection*{rec}

\begin{theorem}[blueprint="infogeometry:geometry:hilberthessian:rec"]
  \inputleannode{InfoGeometry.Geometry.HilbertHessian.rec}
\end{theorem}

\subsection*{recOn}

\begin{theorem}[blueprint="infogeometry:geometry:hilberthessian:recon"]
  \inputleannode{InfoGeometry.Geometry.HilbertHessian.recOn}
\end{theorem}

\subsection*{twiceDiff}

\begin{theorem}[blueprint="infogeometry:geometry:hilberthessian:twicediff"]
  \inputleannode{InfoGeometry.Geometry.HilbertHessian.twiceDiff}
\end{theorem}

\subsection*{ψ}

\begin{theorem}[blueprint="infogeometry:geometry:hilberthessian:ψ"]
  \inputleannode{InfoGeometry.Geometry.HilbertHessian.ψ}
\end{theorem}

\section{InfoGeometry.Geometry.IsFenchelConjugate}

\subsection*{isFenchelMajorized}

\begin{theorem}[blueprint="infogeometry:geometry:isfenchelconjugate:isfenchelmajorized"]
  \inputleannode{InfoGeometry.Geometry.IsFenchelConjugate.isFenchelMajorized}
\end{theorem}

\section{InfoGeometry.Geometry.KreinAsHessian}

\noindent\textbf{Buildable}: yes

\noindent\emph{No exported declarations in the current declaration graph.}

\section{InfoGeometry.Geometry.LegendreDuality}

\noindent\textbf{Buildable}: yes

\noindent\emph{No exported declarations in the current declaration graph.}

\section{InfoGeometry.Geometry.LegendreInvolutionAssumptions}

\subsection*{casesOn}

\begin{theorem}[blueprint="infogeometry:geometry:legendreinvolutionassumptions:caseson"]
  \inputleannode{InfoGeometry.Geometry.LegendreInvolutionAssumptions.casesOn}
\end{theorem}

\subsection*{conjugate}

\begin{theorem}[blueprint="infogeometry:geometry:legendreinvolutionassumptions:conjugate"]
  \inputleannode{InfoGeometry.Geometry.LegendreInvolutionAssumptions.conjugate}
\end{theorem}

\subsection*{ctorIdx}

\begin{theorem}[blueprint="infogeometry:geometry:legendreinvolutionassumptions:ctoridx"]
  \inputleannode{InfoGeometry.Geometry.LegendreInvolutionAssumptions.ctorIdx}
\end{theorem}

\subsection*{diff}

\begin{theorem}[blueprint="infogeometry:geometry:legendreinvolutionassumptions:diff"]
  \inputleannode{InfoGeometry.Geometry.LegendreInvolutionAssumptions.diff}
\end{theorem}

\subsection*{dual_value_along_grad}

\begin{theorem}[blueprint="infogeometry:geometry:legendreinvolutionassumptions:dual_value_along_grad"]
  \inputleannode{InfoGeometry.Geometry.LegendreInvolutionAssumptions.dual_value_along_grad}
\end{theorem}

\subsection*{fenchelGap_eq_zero_along_grad}

\begin{theorem}[blueprint="infogeometry:geometry:legendreinvolutionassumptions:fenchelgap_eq_zero_along_grad"]
  \inputleannode{InfoGeometry.Geometry.LegendreInvolutionAssumptions.fenchelGap_eq_zero_along_grad}
\end{theorem}

\subsection*{fenchelYoung_along_grad}

\begin{theorem}[blueprint="infogeometry:geometry:legendreinvolutionassumptions:fenchelyoung_along_grad"]
  \inputleannode{InfoGeometry.Geometry.LegendreInvolutionAssumptions.fenchelYoung_along_grad}
\end{theorem}

\subsection*{fenchelYoung_eq}

\begin{theorem}[blueprint="infogeometry:geometry:legendreinvolutionassumptions:fenchelyoung_eq"]
  \inputleannode{InfoGeometry.Geometry.LegendreInvolutionAssumptions.fenchelYoung_eq}
\end{theorem}

\subsection*{grad}

\begin{theorem}[blueprint="infogeometry:geometry:legendreinvolutionassumptions:grad"]
  \inputleannode{InfoGeometry.Geometry.LegendreInvolutionAssumptions.grad}
\end{theorem}

\subsection*{gradStar}

\begin{theorem}[blueprint="infogeometry:geometry:legendreinvolutionassumptions:gradstar"]
  \inputleannode{InfoGeometry.Geometry.LegendreInvolutionAssumptions.gradStar}
\end{theorem}

\subsection*{grad_eq_fderiv}

\begin{theorem}[blueprint="infogeometry:geometry:legendreinvolutionassumptions:grad_eq_fderiv"]
  \inputleannode{InfoGeometry.Geometry.LegendreInvolutionAssumptions.grad_eq_fderiv}
\end{theorem}

\subsection*{left_inv}

\begin{theorem}[blueprint="infogeometry:geometry:legendreinvolutionassumptions:left_inv"]
  \inputleannode{InfoGeometry.Geometry.LegendreInvolutionAssumptions.left_inv}
\end{theorem}

\subsection*{mk}

\begin{theorem}[blueprint="infogeometry:geometry:legendreinvolutionassumptions:mk"]
  \inputleannode{InfoGeometry.Geometry.LegendreInvolutionAssumptions.mk}
\end{theorem}

\subsection*{noConfusion}

\begin{theorem}[blueprint="infogeometry:geometry:legendreinvolutionassumptions:noconfusion"]
  \inputleannode{InfoGeometry.Geometry.LegendreInvolutionAssumptions.noConfusion}
\end{theorem}

\subsection*{noConfusionType}

\begin{theorem}[blueprint="infogeometry:geometry:legendreinvolutionassumptions:noconfusiontype"]
  \inputleannode{InfoGeometry.Geometry.LegendreInvolutionAssumptions.noConfusionType}
\end{theorem}

\subsection*{rec}

\begin{theorem}[blueprint="infogeometry:geometry:legendreinvolutionassumptions:rec"]
  \inputleannode{InfoGeometry.Geometry.LegendreInvolutionAssumptions.rec}
\end{theorem}

\subsection*{recOn}

\begin{theorem}[blueprint="infogeometry:geometry:legendreinvolutionassumptions:recon"]
  \inputleannode{InfoGeometry.Geometry.LegendreInvolutionAssumptions.recOn}
\end{theorem}

\subsection*{right_inv}

\begin{theorem}[blueprint="infogeometry:geometry:legendreinvolutionassumptions:right_inv"]
  \inputleannode{InfoGeometry.Geometry.LegendreInvolutionAssumptions.right_inv}
\end{theorem}

\section{InfoGeometry.Geometry.LegendreWellPosed}

\subsection*{bddAbove}

\begin{theorem}[blueprint="infogeometry:geometry:legendrewellposed:bddabove"]
  \inputleannode{InfoGeometry.Geometry.LegendreWellPosed.bddAbove}
\end{theorem}

\subsection*{casesOn}

\begin{theorem}[blueprint="infogeometry:geometry:legendrewellposed:caseson"]
  \inputleannode{InfoGeometry.Geometry.LegendreWellPosed.casesOn}
\end{theorem}

\subsection*{mk}

\begin{theorem}[blueprint="infogeometry:geometry:legendrewellposed:mk"]
  \inputleannode{InfoGeometry.Geometry.LegendreWellPosed.mk}
\end{theorem}

\subsection*{nonempty}

\begin{theorem}[blueprint="infogeometry:geometry:legendrewellposed:nonempty"]
  \inputleannode{InfoGeometry.Geometry.LegendreWellPosed.nonempty}
\end{theorem}

\subsection*{rec}

\begin{theorem}[blueprint="infogeometry:geometry:legendrewellposed:rec"]
  \inputleannode{InfoGeometry.Geometry.LegendreWellPosed.rec}
\end{theorem}

\subsection*{recOn}

\begin{theorem}[blueprint="infogeometry:geometry:legendrewellposed:recon"]
  \inputleannode{InfoGeometry.Geometry.LegendreWellPosed.recOn}
\end{theorem}

\section{InfoGeometry.Geometry.divergenceVec}

\subsection*{congr_simp}

\begin{theorem}[blueprint="infogeometry:geometry:divergencevec:congr_simp"]
  \inputleannode{InfoGeometry.Geometry.divergenceVec.congr_simp}
\end{theorem}

\section{InfoGeometry.Geometry.dualCoordVec}

\subsection*{congr_simp}

\begin{theorem}[blueprint="infogeometry:geometry:dualcoordvec:congr_simp"]
  \inputleannode{InfoGeometry.Geometry.dualCoordVec.congr_simp}
\end{theorem}

\section{InfoGeometry.Geometry.eGeodesic}

\subsection*{eq_1}

\begin{theorem}[blueprint="infogeometry:geometry:egeodesic:eq_1"]
  \inputleannode{InfoGeometry.Geometry.eGeodesic.eq_1}
\end{theorem}

\section{InfoGeometry.Geometry.eGeodesicConvex}

\subsection*{convexOn_univ}

\begin{theorem}[blueprint="infogeometry:geometry:egeodesicconvex:convexon_univ"]
  \inputleannode{InfoGeometry.Geometry.eGeodesicConvex.convexOn_univ}
\end{theorem}

\section{InfoGeometry.Geometry.kreinHessian}

\subsection*{eq_1}

\begin{theorem}[blueprint="infogeometry:geometry:kreinhessian:eq_1"]
  \inputleannode{InfoGeometry.Geometry.kreinHessian.eq_1}
\end{theorem}

\section{InfoGeometry.Geometry.mGeodesic}

\subsection*{eq_1}

\begin{theorem}[blueprint="infogeometry:geometry:mgeodesic:eq_1"]
  \inputleannode{InfoGeometry.Geometry.mGeodesic.eq_1}
\end{theorem}

\chapter{GrandCanonical Modules}

\section{InfoGeometry.GrandCanonical}

\subsection*{GrandCanonicalParams}

\begin{theorem}[blueprint="infogeometry:grandcanonical:grandcanonicalparams"]
  \inputleannode{InfoGeometry.GrandCanonical.GrandCanonicalParams}
\end{theorem}

\subsection*{GrandCanonicalTwoParam}

\begin{theorem}[blueprint="infogeometry:grandcanonical:grandcanonicaltwoparam"]
  \inputleannode{InfoGeometry.GrandCanonical.GrandCanonicalTwoParam}
\end{theorem}

\subsection*{Spinodal}

\begin{theorem}[blueprint="infogeometry:grandcanonical:spinodal"]
  \inputleannode{InfoGeometry.GrandCanonical.Spinodal}
\end{theorem}

\subsection*{covariance}

\begin{theorem}[blueprint="infogeometry:grandcanonical:covariance"]
  \inputleannode{InfoGeometry.GrandCanonical.covariance}
\end{theorem}

\subsection*{deriv_partition}

\begin{theorem}[blueprint="infogeometry:grandcanonical:deriv_partition"]
  \inputleannode{InfoGeometry.GrandCanonical.deriv_partition}
\end{theorem}

\subsection*{deriv_partitionDerivFun}

\begin{theorem}[blueprint="infogeometry:grandcanonical:deriv_partitionderivfun"]
  \inputleannode{InfoGeometry.GrandCanonical.deriv_partitionDerivFun}
\end{theorem}

\subsection*{firstMomentUnnormalized}

\begin{theorem}[blueprint="infogeometry:grandcanonical:firstmomentunnormalized"]
  \inputleannode{InfoGeometry.GrandCanonical.firstMomentUnnormalized}
\end{theorem}

\subsection*{firstNumberUnnormalized}

\begin{theorem}[blueprint="infogeometry:grandcanonical:firstnumberunnormalized"]
  \inputleannode{InfoGeometry.GrandCanonical.firstNumberUnnormalized}
\end{theorem}

\subsection*{firstShiftUnnormalized}

\begin{theorem}[blueprint="infogeometry:grandcanonical:firstshiftunnormalized"]
  \inputleannode{InfoGeometry.GrandCanonical.firstShiftUnnormalized}
\end{theorem}

\subsection*{gibbsWeight}

\begin{theorem}[blueprint="infogeometry:grandcanonical:gibbsweight"]
  \inputleannode{InfoGeometry.GrandCanonical.gibbsWeight}
\end{theorem}

\subsection*{gibbsWeightGC}

\begin{theorem}[blueprint="infogeometry:grandcanonical:gibbsweightgc"]
  \inputleannode{InfoGeometry.GrandCanonical.gibbsWeightGC}
\end{theorem}

\subsection*{gibbsWeightGC_nonneg}

\begin{theorem}[blueprint="infogeometry:grandcanonical:gibbsweightgc_nonneg"]
  \inputleannode{InfoGeometry.GrandCanonical.gibbsWeightGC_nonneg}
\end{theorem}

\subsection*{gibbsWeightGC_pos}

\begin{theorem}[blueprint="infogeometry:grandcanonical:gibbsweightgc_pos"]
  \inputleannode{InfoGeometry.GrandCanonical.gibbsWeightGC_pos}
\end{theorem}

\subsection*{gibbsWeightGC_sum_one}

\begin{theorem}[blueprint="infogeometry:grandcanonical:gibbsweightgc_sum_one"]
  \inputleannode{InfoGeometry.GrandCanonical.gibbsWeightGC_sum_one}
\end{theorem}

\subsection*{gibbsWeight_nonneg}

\begin{theorem}[blueprint="infogeometry:grandcanonical:gibbsweight_nonneg"]
  \inputleannode{InfoGeometry.GrandCanonical.gibbsWeight_nonneg}
\end{theorem}

\subsection*{gibbsWeight_pos}

\begin{theorem}[blueprint="infogeometry:grandcanonical:gibbsweight_pos"]
  \inputleannode{InfoGeometry.GrandCanonical.gibbsWeight_pos}
\end{theorem}

\subsection*{gibbsWeight_sum_one}

\begin{theorem}[blueprint="infogeometry:grandcanonical:gibbsweight_sum_one"]
  \inputleannode{InfoGeometry.GrandCanonical.gibbsWeight_sum_one}
\end{theorem}

\subsection*{hasDerivAt_exp_neg_mul_energy}

\begin{theorem}[blueprint="infogeometry:grandcanonical:hasderivat_exp_neg_mul_energy"]
  \inputleannode{InfoGeometry.GrandCanonical.hasDerivAt_exp_neg_mul_energy}
\end{theorem}

\subsection*{hasDerivAt_exp_neg_mul_shiftedEnergy_beta}

\begin{theorem}[blueprint="infogeometry:grandcanonical:hasderivat_exp_neg_mul_shiftedenergy_beta"]
  \inputleannode{InfoGeometry.GrandCanonical.hasDerivAt_exp_neg_mul_shiftedEnergy_beta}
\end{theorem}

\subsection*{hasDerivAt_exp_neg_mul_shiftedEnergy_mu}

\begin{theorem}[blueprint="infogeometry:grandcanonical:hasderivat_exp_neg_mul_shiftedenergy_mu"]
  \inputleannode{InfoGeometry.GrandCanonical.hasDerivAt_exp_neg_mul_shiftedEnergy_mu}
\end{theorem}

\subsection*{hasDerivAt_partition}

\begin{theorem}[blueprint="infogeometry:grandcanonical:hasderivat_partition"]
  \inputleannode{InfoGeometry.GrandCanonical.hasDerivAt_partition}
\end{theorem}

\subsection*{hasDerivAt_partitionDerivFun}

\begin{theorem}[blueprint="infogeometry:grandcanonical:hasderivat_partitionderivfun"]
  \inputleannode{InfoGeometry.GrandCanonical.hasDerivAt_partitionDerivFun}
\end{theorem}

\subsection*{hasDerivAt_partitionDerivTerm}

\begin{theorem}[blueprint="infogeometry:grandcanonical:hasderivat_partitionderivterm"]
  \inputleannode{InfoGeometry.GrandCanonical.hasDerivAt_partitionDerivTerm}
\end{theorem}

\subsection*{hasDerivAt_partitionGC_beta}

\begin{theorem}[blueprint="infogeometry:grandcanonical:hasderivat_partitiongc_beta"]
  \inputleannode{InfoGeometry.GrandCanonical.hasDerivAt_partitionGC_beta}
\end{theorem}

\subsection*{hasDerivAt_partitionGC_mu}

\begin{theorem}[blueprint="infogeometry:grandcanonical:hasderivat_partitiongc_mu"]
  \inputleannode{InfoGeometry.GrandCanonical.hasDerivAt_partitionGC_mu}
\end{theorem}

\subsection*{hessian}

\begin{theorem}[blueprint="infogeometry:grandcanonical:hessian"]
  \inputleannode{InfoGeometry.GrandCanonical.hessian}
\end{theorem}

\subsection*{hessian_eq_secondMoment_sub_mean_sq}

\begin{theorem}[blueprint="infogeometry:grandcanonical:hessian_eq_secondmoment_sub_mean_sq"]
  \inputleannode{InfoGeometry.GrandCanonical.hessian_eq_secondMoment_sub_mean_sq}
\end{theorem}

\subsection*{hessian_nonneg}

\begin{theorem}[blueprint="infogeometry:grandcanonical:hessian_nonneg"]
  \inputleannode{InfoGeometry.GrandCanonical.hessian_nonneg}
\end{theorem}

\subsection*{mean}

\begin{theorem}[blueprint="infogeometry:grandcanonical:mean"]
  \inputleannode{InfoGeometry.GrandCanonical.mean}
\end{theorem}

\subsection*{meanNumber}

\begin{theorem}[blueprint="infogeometry:grandcanonical:meannumber"]
  \inputleannode{InfoGeometry.GrandCanonical.meanNumber}
\end{theorem}

\subsection*{meanNumber_eq_firstNumber_div_partition}

\begin{theorem}[blueprint="infogeometry:grandcanonical:meannumber_eq_firstnumber_div_partition"]
  \inputleannode{InfoGeometry.GrandCanonical.meanNumber_eq_firstNumber_div_partition}
\end{theorem}

\subsection*{meanShift}

\begin{theorem}[blueprint="infogeometry:grandcanonical:meanshift"]
  \inputleannode{InfoGeometry.GrandCanonical.meanShift}
\end{theorem}

\subsection*{meanShift_eq_firstShift_div_partition}

\begin{theorem}[blueprint="infogeometry:grandcanonical:meanshift_eq_firstshift_div_partition"]
  \inputleannode{InfoGeometry.GrandCanonical.meanShift_eq_firstShift_div_partition}
\end{theorem}

\subsection*{mean_eq_firstMoment_div_partition}

\begin{theorem}[blueprint="infogeometry:grandcanonical:mean_eq_firstmoment_div_partition"]
  \inputleannode{InfoGeometry.GrandCanonical.mean_eq_firstMoment_div_partition}
\end{theorem}

\subsection*{partition}

\begin{theorem}[blueprint="infogeometry:grandcanonical:partition"]
  \inputleannode{InfoGeometry.GrandCanonical.partition}
\end{theorem}

\subsection*{partitionDerivFun}

\begin{theorem}[blueprint="infogeometry:grandcanonical:partitionderivfun"]
  \inputleannode{InfoGeometry.GrandCanonical.partitionDerivFun}
\end{theorem}

\subsection*{partitionDerivFun_eq_neg_firstMoment}

\begin{theorem}[blueprint="infogeometry:grandcanonical:partitionderivfun_eq_neg_firstmoment"]
  \inputleannode{InfoGeometry.GrandCanonical.partitionDerivFun_eq_neg_firstMoment}
\end{theorem}

\subsection*{partitionGC}

\begin{theorem}[blueprint="infogeometry:grandcanonical:partitiongc"]
  \inputleannode{InfoGeometry.GrandCanonical.partitionGC}
\end{theorem}

\subsection*{partitionGCDerivBetaFun}

\begin{theorem}[blueprint="infogeometry:grandcanonical:partitiongcderivbetafun"]
  \inputleannode{InfoGeometry.GrandCanonical.partitionGCDerivBetaFun}
\end{theorem}

\subsection*{partitionGCDerivBetaFun_eq_neg_firstShift}

\begin{theorem}[blueprint="infogeometry:grandcanonical:partitiongcderivbetafun_eq_neg_firstshift"]
  \inputleannode{InfoGeometry.GrandCanonical.partitionGCDerivBetaFun_eq_neg_firstShift}
\end{theorem}

\subsection*{partitionGCDerivMuFun}

\begin{theorem}[blueprint="infogeometry:grandcanonical:partitiongcderivmufun"]
  \inputleannode{InfoGeometry.GrandCanonical.partitionGCDerivMuFun}
\end{theorem}

\subsection*{partitionGCDerivMuFun_eq_beta_mul_firstNumber}

\begin{theorem}[blueprint="infogeometry:grandcanonical:partitiongcderivmufun_eq_beta_mul_firstnumber"]
  \inputleannode{InfoGeometry.GrandCanonical.partitionGCDerivMuFun_eq_beta_mul_firstNumber}
\end{theorem}

\subsection*{partitionGC_pos}

\begin{theorem}[blueprint="infogeometry:grandcanonical:partitiongc_pos"]
  \inputleannode{InfoGeometry.GrandCanonical.partitionGC_pos}
\end{theorem}

\subsection*{partition_pos}

\begin{theorem}[blueprint="infogeometry:grandcanonical:partition_pos"]
  \inputleannode{InfoGeometry.GrandCanonical.partition_pos}
\end{theorem}

\subsection*{potential}

\begin{theorem}[blueprint="infogeometry:grandcanonical:potential"]
  \inputleannode{InfoGeometry.GrandCanonical.potential}
\end{theorem}

\subsection*{potentialGC}

\begin{theorem}[blueprint="infogeometry:grandcanonical:potentialgc"]
  \inputleannode{InfoGeometry.GrandCanonical.potentialGC}
\end{theorem}

\subsection*{potentialGC_deriv_beta_eq_neg_meanShift}

\begin{theorem}[blueprint="infogeometry:grandcanonical:potentialgc_deriv_beta_eq_neg_meanshift"]
  \inputleannode{InfoGeometry.GrandCanonical.potentialGC_deriv_beta_eq_neg_meanShift}
\end{theorem}

\subsection*{potentialGC_deriv_beta_eq_partitionDeriv_div_partition}

\begin{theorem}[blueprint="infogeometry:grandcanonical:potentialgc_deriv_beta_eq_partitionderiv_div_partition"]
  \inputleannode{InfoGeometry.GrandCanonical.potentialGC_deriv_beta_eq_partitionDeriv_div_partition}
\end{theorem}

\subsection*{potentialGC_deriv_mu_eq_beta_meanNumber}

\begin{theorem}[blueprint="infogeometry:grandcanonical:potentialgc_deriv_mu_eq_beta_meannumber"]
  \inputleannode{InfoGeometry.GrandCanonical.potentialGC_deriv_mu_eq_beta_meanNumber}
\end{theorem}

\subsection*{potentialGC_deriv_mu_eq_partitionDeriv_div_partition}

\begin{theorem}[blueprint="infogeometry:grandcanonical:potentialgc_deriv_mu_eq_partitionderiv_div_partition"]
  \inputleannode{InfoGeometry.GrandCanonical.potentialGC_deriv_mu_eq_partitionDeriv_div_partition}
\end{theorem}

\subsection*{potential_deriv_eq_neg_mean}

\begin{theorem}[blueprint="infogeometry:grandcanonical:potential_deriv_eq_neg_mean"]
  \inputleannode{InfoGeometry.GrandCanonical.potential_deriv_eq_neg_mean}
\end{theorem}

\subsection*{potential_deriv_eq_partitionDeriv_div_partition}

\begin{theorem}[blueprint="infogeometry:grandcanonical:potential_deriv_eq_partitionderiv_div_partition"]
  \inputleannode{InfoGeometry.GrandCanonical.potential_deriv_eq_partitionDeriv_div_partition}
\end{theorem}

\subsection*{potential_second_derivative_eq_variance}

\begin{theorem}[blueprint="infogeometry:grandcanonical:potential_second_derivative_eq_variance"]
  \inputleannode{InfoGeometry.GrandCanonical.potential_second_derivative_eq_variance}
\end{theorem}

\subsection*{secondMoment}

\begin{theorem}[blueprint="infogeometry:grandcanonical:secondmoment"]
  \inputleannode{InfoGeometry.GrandCanonical.secondMoment}
\end{theorem}

\subsection*{secondMomentUnnormalized}

\begin{theorem}[blueprint="infogeometry:grandcanonical:secondmomentunnormalized"]
  \inputleannode{InfoGeometry.GrandCanonical.secondMomentUnnormalized}
\end{theorem}

\subsection*{secondMoment_eq_secondMomentUnnormalized_div_partition}

\begin{theorem}[blueprint="infogeometry:grandcanonical:secondmoment_eq_secondmomentunnormalized_div_partition"]
  \inputleannode{InfoGeometry.GrandCanonical.secondMoment_eq_secondMomentUnnormalized_div_partition}
\end{theorem}

\subsection*{shiftedEnergy}

\begin{theorem}[blueprint="infogeometry:grandcanonical:shiftedenergy"]
  \inputleannode{InfoGeometry.GrandCanonical.shiftedEnergy}
\end{theorem}

\subsection*{spinodal_iff_energy_eq_mean}

\begin{theorem}[blueprint="infogeometry:grandcanonical:spinodal_iff_energy_eq_mean"]
  \inputleannode{InfoGeometry.GrandCanonical.spinodal_iff_energy_eq_mean}
\end{theorem}

\subsection*{spinodal_iff_variance_eq_zero}

\begin{theorem}[blueprint="infogeometry:grandcanonical:spinodal_iff_variance_eq_zero"]
  \inputleannode{InfoGeometry.GrandCanonical.spinodal_iff_variance_eq_zero}
\end{theorem}

\subsection*{variance}

\begin{theorem}[blueprint="infogeometry:grandcanonical:variance"]
  \inputleannode{InfoGeometry.GrandCanonical.variance}
\end{theorem}

\subsection*{variance_eq_secondMoment_sub_mean_sq}

\begin{theorem}[blueprint="infogeometry:grandcanonical:variance_eq_secondmoment_sub_mean_sq"]
  \inputleannode{InfoGeometry.GrandCanonical.variance_eq_secondMoment_sub_mean_sq}
\end{theorem}

\subsection*{variance_eq_zero_iff_energy_eq_mean}

\begin{theorem}[blueprint="infogeometry:grandcanonical:variance_eq_zero_iff_energy_eq_mean"]
  \inputleannode{InfoGeometry.GrandCanonical.variance_eq_zero_iff_energy_eq_mean}
\end{theorem}

\subsection*{variance_nonneg}

\begin{theorem}[blueprint="infogeometry:grandcanonical:variance_nonneg"]
  \inputleannode{InfoGeometry.GrandCanonical.variance_nonneg}
\end{theorem}

\section{InfoGeometry.GrandCanonical.Core}

\noindent\textbf{Buildable}: yes

\noindent\emph{No exported declarations in the current declaration graph.}

\section{InfoGeometry.GrandCanonical.GrandCanonicalParams}

\subsection*{casesOn}

\begin{theorem}[blueprint="infogeometry:grandcanonical:grandcanonicalparams:caseson"]
  \inputleannode{InfoGeometry.GrandCanonical.GrandCanonicalParams.casesOn}
\end{theorem}

\subsection*{ctorIdx}

\begin{theorem}[blueprint="infogeometry:grandcanonical:grandcanonicalparams:ctoridx"]
  \inputleannode{InfoGeometry.GrandCanonical.GrandCanonicalParams.ctorIdx}
\end{theorem}

\subsection*{energy}

\begin{theorem}[blueprint="infogeometry:grandcanonical:grandcanonicalparams:energy"]
  \inputleannode{InfoGeometry.GrandCanonical.GrandCanonicalParams.energy}
\end{theorem}

\subsection*{mk}

\begin{theorem}[blueprint="infogeometry:grandcanonical:grandcanonicalparams:mk"]
  \inputleannode{InfoGeometry.GrandCanonical.GrandCanonicalParams.mk}
\end{theorem}

\subsection*{noConfusion}

\begin{theorem}[blueprint="infogeometry:grandcanonical:grandcanonicalparams:noconfusion"]
  \inputleannode{InfoGeometry.GrandCanonical.GrandCanonicalParams.noConfusion}
\end{theorem}

\subsection*{noConfusionType}

\begin{theorem}[blueprint="infogeometry:grandcanonical:grandcanonicalparams:noconfusiontype"]
  \inputleannode{InfoGeometry.GrandCanonical.GrandCanonicalParams.noConfusionType}
\end{theorem}

\subsection*{rec}

\begin{theorem}[blueprint="infogeometry:grandcanonical:grandcanonicalparams:rec"]
  \inputleannode{InfoGeometry.GrandCanonical.GrandCanonicalParams.rec}
\end{theorem}

\subsection*{recOn}

\begin{theorem}[blueprint="infogeometry:grandcanonical:grandcanonicalparams:recon"]
  \inputleannode{InfoGeometry.GrandCanonical.GrandCanonicalParams.recOn}
\end{theorem}

\section{InfoGeometry.GrandCanonical.GrandCanonicalTwoParam}

\subsection*{casesOn}

\begin{theorem}[blueprint="infogeometry:grandcanonical:grandcanonicaltwoparam:caseson"]
  \inputleannode{InfoGeometry.GrandCanonical.GrandCanonicalTwoParam.casesOn}
\end{theorem}

\subsection*{ctorIdx}

\begin{theorem}[blueprint="infogeometry:grandcanonical:grandcanonicaltwoparam:ctoridx"]
  \inputleannode{InfoGeometry.GrandCanonical.GrandCanonicalTwoParam.ctorIdx}
\end{theorem}

\subsection*{energy}

\begin{theorem}[blueprint="infogeometry:grandcanonical:grandcanonicaltwoparam:energy"]
  \inputleannode{InfoGeometry.GrandCanonical.GrandCanonicalTwoParam.energy}
\end{theorem}

\subsection*{mk}

\begin{theorem}[blueprint="infogeometry:grandcanonical:grandcanonicaltwoparam:mk"]
  \inputleannode{InfoGeometry.GrandCanonical.GrandCanonicalTwoParam.mk}
\end{theorem}

\subsection*{noConfusion}

\begin{theorem}[blueprint="infogeometry:grandcanonical:grandcanonicaltwoparam:noconfusion"]
  \inputleannode{InfoGeometry.GrandCanonical.GrandCanonicalTwoParam.noConfusion}
\end{theorem}

\subsection*{noConfusionType}

\begin{theorem}[blueprint="infogeometry:grandcanonical:grandcanonicaltwoparam:noconfusiontype"]
  \inputleannode{InfoGeometry.GrandCanonical.GrandCanonicalTwoParam.noConfusionType}
\end{theorem}

\subsection*{number}

\begin{theorem}[blueprint="infogeometry:grandcanonical:grandcanonicaltwoparam:number"]
  \inputleannode{InfoGeometry.GrandCanonical.GrandCanonicalTwoParam.number}
\end{theorem}

\subsection*{rec}

\begin{theorem}[blueprint="infogeometry:grandcanonical:grandcanonicaltwoparam:rec"]
  \inputleannode{InfoGeometry.GrandCanonical.GrandCanonicalTwoParam.rec}
\end{theorem}

\subsection*{recOn}

\begin{theorem}[blueprint="infogeometry:grandcanonical:grandcanonicaltwoparam:recon"]
  \inputleannode{InfoGeometry.GrandCanonical.GrandCanonicalTwoParam.recOn}
\end{theorem}

\section{InfoGeometry.GrandCanonical.variance}

\subsection*{eq_1}

\begin{theorem}[blueprint="infogeometry:grandcanonical:variance:eq_1"]
  \inputleannode{InfoGeometry.GrandCanonical.variance.eq_1}
\end{theorem}

\chapter{Information Modules}

\section{InfoGeometry.Information}

\subsection*{FiniteExpFamily}

\begin{theorem}[blueprint="infogeometry:information:finiteexpfamily"]
  \inputleannode{InfoGeometry.Information.FiniteExpFamily}
\end{theorem}

\subsection*{centeredLinearStat}

\begin{theorem}[blueprint="infogeometry:information:centeredlinearstat"]
  \inputleannode{InfoGeometry.Information.centeredLinearStat}
\end{theorem}

\subsection*{centeredLinearStat_eq_sum_centeredStat}

\begin{theorem}[blueprint="infogeometry:information:centeredlinearstat_eq_sum_centeredstat"]
  \inputleannode{InfoGeometry.Information.centeredLinearStat_eq_sum_centeredStat}
\end{theorem}

\subsection*{centeredStat}

\begin{theorem}[blueprint="infogeometry:information:centeredstat"]
  \inputleannode{InfoGeometry.Information.centeredStat}
\end{theorem}

\subsection*{covariance}

\begin{theorem}[blueprint="infogeometry:information:covariance"]
  \inputleannode{InfoGeometry.Information.covariance}
\end{theorem}

\subsection*{covariance_symm}

\begin{theorem}[blueprint="infogeometry:information:covariance_symm"]
  \inputleannode{InfoGeometry.Information.covariance_symm}
\end{theorem}

\subsection*{density}

\begin{theorem}[blueprint="infogeometry:information:density"]
  \inputleannode{InfoGeometry.Information.density}
\end{theorem}

\subsection*{density_nonneg}

\begin{theorem}[blueprint="infogeometry:information:density_nonneg"]
  \inputleannode{InfoGeometry.Information.density_nonneg}
\end{theorem}

\subsection*{density_pos}

\begin{theorem}[blueprint="infogeometry:information:density_pos"]
  \inputleannode{InfoGeometry.Information.density_pos}
\end{theorem}

\subsection*{density_sum_one}

\begin{theorem}[blueprint="infogeometry:information:density_sum_one"]
  \inputleannode{InfoGeometry.Information.density_sum_one}
\end{theorem}

\subsection*{expLogPotential}

\begin{theorem}[blueprint="infogeometry:information:explogpotential"]
  \inputleannode{InfoGeometry.Information.expLogPotential}
\end{theorem}

\subsection*{expectation}

\begin{theorem}[blueprint="infogeometry:information:expectation"]
  \inputleannode{InfoGeometry.Information.expectation}
\end{theorem}

\subsection*{expectation_linearStat}

\begin{theorem}[blueprint="infogeometry:information:expectation_linearstat"]
  \inputleannode{InfoGeometry.Information.expectation_linearStat}
\end{theorem}

\subsection*{fisherMetric}

\begin{theorem}[blueprint="infogeometry:information:fishermetric"]
  \inputleannode{InfoGeometry.Information.fisherMetric}
\end{theorem}

\subsection*{fisherMetric_symm}

\begin{theorem}[blueprint="infogeometry:information:fishermetric_symm"]
  \inputleannode{InfoGeometry.Information.fisherMetric_symm}
\end{theorem}

\subsection*{fisherQuadratic}

\begin{theorem}[blueprint="infogeometry:information:fisherquadratic"]
  \inputleannode{InfoGeometry.Information.fisherQuadratic}
\end{theorem}

\subsection*{fisher_as_variance}

\begin{theorem}[blueprint="infogeometry:information:fisher_as_variance"]
  \inputleannode{InfoGeometry.Information.fisher_as_variance}
\end{theorem}

\subsection*{fisher_positive_definite}

\begin{theorem}[blueprint="infogeometry:information:fisher_positive_definite"]
  \inputleannode{InfoGeometry.Information.fisher_positive_definite}
\end{theorem}

\subsection*{fisher_positive_definite_of_nonconstant_linearStats}

\begin{theorem}[blueprint="infogeometry:information:fisher_positive_definite_of_nonconstant_linearstats"]
  \inputleannode{InfoGeometry.Information.fisher_positive_definite_of_nonconstant_linearStats}
\end{theorem}

\subsection*{fisher_positive_semidefinite}

\begin{theorem}[blueprint="infogeometry:information:fisher_positive_semidefinite"]
  \inputleannode{InfoGeometry.Information.fisher_positive_semidefinite}
\end{theorem}

\subsection*{linearStat}

\begin{theorem}[blueprint="infogeometry:information:linearstat"]
  \inputleannode{InfoGeometry.Information.linearStat}
\end{theorem}

\subsection*{logPartition}

\begin{theorem}[blueprint="infogeometry:information:logpartition"]
  \inputleannode{InfoGeometry.Information.logPartition}
\end{theorem}

\subsection*{naturalParam}

\begin{theorem}[blueprint="infogeometry:information:naturalparam"]
  \inputleannode{InfoGeometry.Information.naturalParam}
\end{theorem}

\subsection*{partition}

\begin{theorem}[blueprint="infogeometry:information:partition"]
  \inputleannode{InfoGeometry.Information.partition}
\end{theorem}

\subsection*{partition_pos}

\begin{theorem}[blueprint="infogeometry:information:partition_pos"]
  \inputleannode{InfoGeometry.Information.partition_pos}
\end{theorem}

\subsection*{statMean}

\begin{theorem}[blueprint="infogeometry:information:statmean"]
  \inputleannode{InfoGeometry.Information.statMean}
\end{theorem}

\section{InfoGeometry.Information.FiniteExpFamily}

\subsection*{base}

\begin{theorem}[blueprint="infogeometry:information:finiteexpfamily:base"]
  \inputleannode{InfoGeometry.Information.FiniteExpFamily.base}
\end{theorem}

\subsection*{base_pos}

\begin{theorem}[blueprint="infogeometry:information:finiteexpfamily:base_pos"]
  \inputleannode{InfoGeometry.Information.FiniteExpFamily.base_pos}
\end{theorem}

\subsection*{casesOn}

\begin{theorem}[blueprint="infogeometry:information:finiteexpfamily:caseson"]
  \inputleannode{InfoGeometry.Information.FiniteExpFamily.casesOn}
\end{theorem}

\subsection*{ctorIdx}

\begin{theorem}[blueprint="infogeometry:information:finiteexpfamily:ctoridx"]
  \inputleannode{InfoGeometry.Information.FiniteExpFamily.ctorIdx}
\end{theorem}

\subsection*{mk}

\begin{theorem}[blueprint="infogeometry:information:finiteexpfamily:mk"]
  \inputleannode{InfoGeometry.Information.FiniteExpFamily.mk}
\end{theorem}

\subsection*{noConfusion}

\begin{theorem}[blueprint="infogeometry:information:finiteexpfamily:noconfusion"]
  \inputleannode{InfoGeometry.Information.FiniteExpFamily.noConfusion}
\end{theorem}

\subsection*{noConfusionType}

\begin{theorem}[blueprint="infogeometry:information:finiteexpfamily:noconfusiontype"]
  \inputleannode{InfoGeometry.Information.FiniteExpFamily.noConfusionType}
\end{theorem}

\subsection*{rec}

\begin{theorem}[blueprint="infogeometry:information:finiteexpfamily:rec"]
  \inputleannode{InfoGeometry.Information.FiniteExpFamily.rec}
\end{theorem}

\subsection*{recOn}

\begin{theorem}[blueprint="infogeometry:information:finiteexpfamily:recon"]
  \inputleannode{InfoGeometry.Information.FiniteExpFamily.recOn}
\end{theorem}

\subsection*{stat}

\begin{theorem}[blueprint="infogeometry:information:finiteexpfamily:stat"]
  \inputleannode{InfoGeometry.Information.FiniteExpFamily.stat}
\end{theorem}

\section{InfoGeometry.Information.MultiLogPotential}

\noindent\textbf{Buildable}: yes

\noindent\emph{No exported declarations in the current declaration graph.}

\section{InfoGeometry.Information.centeredLinearStat}

\subsection*{eq_1}

\begin{theorem}[blueprint="infogeometry:information:centeredlinearstat:eq_1"]
  \inputleannode{InfoGeometry.Information.centeredLinearStat.eq_1}
\end{theorem}

\section{InfoGeometry.Information.expectation}

\subsection*{eq_1}

\begin{theorem}[blueprint="infogeometry:information:expectation:eq_1"]
  \inputleannode{InfoGeometry.Information.expectation.eq_1}
\end{theorem}

\section{InfoGeometry.Information.linearStat}

\subsection*{eq_1}

\begin{theorem}[blueprint="infogeometry:information:linearstat:eq_1"]
  \inputleannode{InfoGeometry.Information.linearStat.eq_1}
\end{theorem}

\section{InfoGeometry.Information.statMean}

\subsection*{eq_1}

\begin{theorem}[blueprint="infogeometry:information:statmean:eq_1"]
  \inputleannode{InfoGeometry.Information.statMean.eq_1}
\end{theorem}

\chapter{Jordan Modules}

\section{InfoGeometry.Jordan}

\subsection*{JordanAlgebra}

\begin{theorem}[blueprint="infogeometry:jordan:jordanalgebra"]
  \inputleannode{InfoGeometry.Jordan.JordanAlgebra}
\end{theorem}

\subsection*{SPD}

\begin{theorem}[blueprint="infogeometry:jordan:spd"]
  \inputleannode{InfoGeometry.Jordan.SPD}
\end{theorem}

\subsection*{instCoeSPDMatrixFinReal}

\begin{theorem}[blueprint="infogeometry:jordan:instcoespdmatrixfinreal"]
  \inputleannode{InfoGeometry.Jordan.instCoeSPDMatrixFinReal}
\end{theorem}

\subsection*{jordanProd_add_left}

\begin{theorem}[blueprint="infogeometry:jordan:jordanprod_add_left"]
  \inputleannode{InfoGeometry.Jordan.jordanProd_add_left}
\end{theorem}

\subsection*{jordanProd_add_right}

\begin{theorem}[blueprint="infogeometry:jordan:jordanprod_add_right"]
  \inputleannode{InfoGeometry.Jordan.jordanProd_add_right}
\end{theorem}

\subsection*{jordanProd_comm}

\begin{theorem}[blueprint="infogeometry:jordan:jordanprod_comm"]
  \inputleannode{InfoGeometry.Jordan.jordanProd_comm}
\end{theorem}

\subsection*{jordanProd_identity}

\begin{theorem}[blueprint="infogeometry:jordan:jordanprod_identity"]
  \inputleannode{InfoGeometry.Jordan.jordanProd_identity}
\end{theorem}

\subsection*{jordanProd_smul_left}

\begin{theorem}[blueprint="infogeometry:jordan:jordanprod_smul_left"]
  \inputleannode{InfoGeometry.Jordan.jordanProd_smul_left}
\end{theorem}

\subsection*{jordanProd_smul_right}

\begin{theorem}[blueprint="infogeometry:jordan:jordanprod_smul_right"]
  \inputleannode{InfoGeometry.Jordan.jordanProd_smul_right}
\end{theorem}

\subsection*{jordanProd_zero_left}

\begin{theorem}[blueprint="infogeometry:jordan:jordanprod_zero_left"]
  \inputleannode{InfoGeometry.Jordan.jordanProd_zero_left}
\end{theorem}

\subsection*{jordanProd_zero_right}

\begin{theorem}[blueprint="infogeometry:jordan:jordanprod_zero_right"]
  \inputleannode{InfoGeometry.Jordan.jordanProd_zero_right}
\end{theorem}

\subsection*{logDetBarrier}

\begin{theorem}[blueprint="infogeometry:jordan:logdetbarrier"]
  \inputleannode{InfoGeometry.Jordan.logDetBarrier}
\end{theorem}

\subsection*{logDetBregman}

\begin{theorem}[blueprint="infogeometry:jordan:logdetbregman"]
  \inputleannode{InfoGeometry.Jordan.logDetBregman}
\end{theorem}

\subsection*{logDetBregman_eq_burg_form}

\begin{theorem}[blueprint="infogeometry:jordan:logdetbregman_eq_burg_form"]
  \inputleannode{InfoGeometry.Jordan.logDetBregman_eq_burg_form}
\end{theorem}

\subsection*{logDetBregman_nonneg}

\begin{theorem}[blueprint="infogeometry:jordan:logdetbregman_nonneg"]
  \inputleannode{InfoGeometry.Jordan.logDetBregman_nonneg}
\end{theorem}

\subsection*{logDetBregman_nonneg_of_commute}

\begin{theorem}[blueprint="infogeometry:jordan:logdetbregman_nonneg_of_commute"]
  \inputleannode{InfoGeometry.Jordan.logDetBregman_nonneg_of_commute}
\end{theorem}

\subsection*{logDetBregman_self}

\begin{theorem}[blueprint="infogeometry:jordan:logdetbregman_self"]
  \inputleannode{InfoGeometry.Jordan.logDetBregman_self}
\end{theorem}

\subsection*{normalizedDistortion}

\begin{theorem}[blueprint="infogeometry:jordan:normalizeddistortion"]
  \inputleannode{InfoGeometry.Jordan.normalizedDistortion}
\end{theorem}

\subsection*{normalizedDistortion_det}

\begin{theorem}[blueprint="infogeometry:jordan:normalizeddistortion_det"]
  \inputleannode{InfoGeometry.Jordan.normalizedDistortion_det}
\end{theorem}

\subsection*{normalizedDistortion_det_pos}

\begin{theorem}[blueprint="infogeometry:jordan:normalizeddistortion_det_pos"]
  \inputleannode{InfoGeometry.Jordan.normalizedDistortion_det_pos}
\end{theorem}

\subsection*{trace_sub_logdet_sub_dim_nonneg_of_posDef}

\begin{theorem}[blueprint="infogeometry:jordan:trace_sub_logdet_sub_dim_nonneg_of_posdef"]
  \inputleannode{InfoGeometry.Jordan.trace_sub_logdet_sub_dim_nonneg_of_posDef}
\end{theorem}

\section{InfoGeometry.Jordan.Core}

\noindent\textbf{Buildable}: yes

\noindent\emph{No exported declarations in the current declaration graph.}

\section{InfoGeometry.Jordan.JordanAlgebra}

\subsection*{add_left}

\begin{theorem}[blueprint="infogeometry:jordan:jordanalgebra:add_left"]
  \inputleannode{InfoGeometry.Jordan.JordanAlgebra.add_left}
\end{theorem}

\subsection*{casesOn}

\begin{theorem}[blueprint="infogeometry:jordan:jordanalgebra:caseson"]
  \inputleannode{InfoGeometry.Jordan.JordanAlgebra.casesOn}
\end{theorem}

\subsection*{comm}

\begin{theorem}[blueprint="infogeometry:jordan:jordanalgebra:comm"]
  \inputleannode{InfoGeometry.Jordan.JordanAlgebra.comm}
\end{theorem}

\subsection*{ctorIdx}

\begin{theorem}[blueprint="infogeometry:jordan:jordanalgebra:ctoridx"]
  \inputleannode{InfoGeometry.Jordan.JordanAlgebra.ctorIdx}
\end{theorem}

\subsection*{jordanProd}

\begin{theorem}[blueprint="infogeometry:jordan:jordanalgebra:jordanprod"]
  \inputleannode{InfoGeometry.Jordan.JordanAlgebra.jordanProd}
\end{theorem}

\subsection*{jordan_identity}

\begin{theorem}[blueprint="infogeometry:jordan:jordanalgebra:jordan_identity"]
  \inputleannode{InfoGeometry.Jordan.JordanAlgebra.jordan_identity}
\end{theorem}

\subsection*{mk}

\begin{theorem}[blueprint="infogeometry:jordan:jordanalgebra:mk"]
  \inputleannode{InfoGeometry.Jordan.JordanAlgebra.mk}
\end{theorem}

\subsection*{noConfusion}

\begin{theorem}[blueprint="infogeometry:jordan:jordanalgebra:noconfusion"]
  \inputleannode{InfoGeometry.Jordan.JordanAlgebra.noConfusion}
\end{theorem}

\subsection*{noConfusionType}

\begin{theorem}[blueprint="infogeometry:jordan:jordanalgebra:noconfusiontype"]
  \inputleannode{InfoGeometry.Jordan.JordanAlgebra.noConfusionType}
\end{theorem}

\subsection*{rec}

\begin{theorem}[blueprint="infogeometry:jordan:jordanalgebra:rec"]
  \inputleannode{InfoGeometry.Jordan.JordanAlgebra.rec}
\end{theorem}

\subsection*{recOn}

\begin{theorem}[blueprint="infogeometry:jordan:jordanalgebra:recon"]
  \inputleannode{InfoGeometry.Jordan.JordanAlgebra.recOn}
\end{theorem}

\subsection*{smul_left}

\begin{theorem}[blueprint="infogeometry:jordan:jordanalgebra:smul_left"]
  \inputleannode{InfoGeometry.Jordan.JordanAlgebra.smul_left}
\end{theorem}

\section{InfoGeometry.Jordan.LogDet}

\noindent\textbf{Buildable}: yes

\noindent\emph{No exported declarations in the current declaration graph.}

\section{InfoGeometry.Jordan.SPD}

\noindent\textbf{Buildable}: yes

\subsection*{casesOn}

\begin{theorem}[blueprint="infogeometry:jordan:spd:caseson"]
  \inputleannode{InfoGeometry.Jordan.SPD.casesOn}
\end{theorem}

\subsection*{ctorIdx}

\begin{theorem}[blueprint="infogeometry:jordan:spd:ctoridx"]
  \inputleannode{InfoGeometry.Jordan.SPD.ctorIdx}
\end{theorem}

\subsection*{det_ne_zero}

\begin{theorem}[blueprint="infogeometry:jordan:spd:det_ne_zero"]
  \inputleannode{InfoGeometry.Jordan.SPD.det_ne_zero}
\end{theorem}

\subsection*{det_pos}

\begin{theorem}[blueprint="infogeometry:jordan:spd:det_pos"]
  \inputleannode{InfoGeometry.Jordan.SPD.det_pos}
\end{theorem}

\subsection*{mat}

\begin{theorem}[blueprint="infogeometry:jordan:spd:mat"]
  \inputleannode{InfoGeometry.Jordan.SPD.mat}
\end{theorem}

\subsection*{mk}

\begin{theorem}[blueprint="infogeometry:jordan:spd:mk"]
  \inputleannode{InfoGeometry.Jordan.SPD.mk}
\end{theorem}

\subsection*{noConfusion}

\begin{theorem}[blueprint="infogeometry:jordan:spd:noconfusion"]
  \inputleannode{InfoGeometry.Jordan.SPD.noConfusion}
\end{theorem}

\subsection*{noConfusionType}

\begin{theorem}[blueprint="infogeometry:jordan:spd:noconfusiontype"]
  \inputleannode{InfoGeometry.Jordan.SPD.noConfusionType}
\end{theorem}

\subsection*{pos}

\begin{theorem}[blueprint="infogeometry:jordan:spd:pos"]
  \inputleannode{InfoGeometry.Jordan.SPD.pos}
\end{theorem}

\subsection*{posDef}

\begin{theorem}[blueprint="infogeometry:jordan:spd:posdef"]
  \inputleannode{InfoGeometry.Jordan.SPD.posDef}
\end{theorem}

\subsection*{rec}

\begin{theorem}[blueprint="infogeometry:jordan:spd:rec"]
  \inputleannode{InfoGeometry.Jordan.SPD.rec}
\end{theorem}

\subsection*{recOn}

\begin{theorem}[blueprint="infogeometry:jordan:spd:recon"]
  \inputleannode{InfoGeometry.Jordan.SPD.recOn}
\end{theorem}

\subsection*{symm}

\begin{theorem}[blueprint="infogeometry:jordan:spd:symm"]
  \inputleannode{InfoGeometry.Jordan.SPD.symm}
\end{theorem}

\subsection*{transpose_eq_self}

\begin{theorem}[blueprint="infogeometry:jordan:spd:transpose_eq_self"]
  \inputleannode{InfoGeometry.Jordan.SPD.transpose_eq_self}
\end{theorem}

\section{InfoGeometry.Jordan.logDetBregman}

\subsection*{eq_1}

\begin{theorem}[blueprint="infogeometry:jordan:logdetbregman:eq_1"]
  \inputleannode{InfoGeometry.Jordan.logDetBregman.eq_1}
\end{theorem}

\section{InfoGeometry.Jordan.normalizedDistortion}

\subsection*{eq_1}

\begin{theorem}[blueprint="infogeometry:jordan:normalizeddistortion:eq_1"]
  \inputleannode{InfoGeometry.Jordan.normalizedDistortion.eq_1}
\end{theorem}

\chapter{KL Modules}

\section{InfoGeometry.KL}

\noindent\textbf{Buildable}: yes

\subsection*{AbsolutelyContinuous}

\begin{theorem}[blueprint="infogeometry:kl:absolutelycontinuous"]
  \inputleannode{InfoGeometry.KL.AbsolutelyContinuous}
\end{theorem}

\subsection*{EmpiricalNontrivial}

\begin{theorem}[blueprint="infogeometry:kl:empiricalnontrivial"]
  \inputleannode{InfoGeometry.KL.EmpiricalNontrivial}
\end{theorem}

\subsection*{FullSupport}

\begin{theorem}[blueprint="infogeometry:kl:fullsupport"]
  \inputleannode{InfoGeometry.KL.FullSupport}
\end{theorem}

\subsection*{KL_eq_zero_iff}

\begin{theorem}[blueprint="infogeometry:kl:kl_eq_zero_iff"]
  \inputleannode{InfoGeometry.KL.KL_eq_zero_iff}
\end{theorem}

\subsection*{KL_nonneg}

\begin{theorem}[blueprint="infogeometry:kl:kl_nonneg"]
  \inputleannode{InfoGeometry.KL.KL_nonneg}
\end{theorem}

\subsection*{KLdivergence}

\begin{theorem}[blueprint="infogeometry:kl:kldivergence"]
  \inputleannode{InfoGeometry.KL.KLdivergence}
\end{theorem}

\subsection*{KLdivergence_eq_klDiv}

\begin{theorem}[blueprint="infogeometry:kl:kldivergence_eq_kldiv"]
  \inputleannode{InfoGeometry.KL.KLdivergence_eq_klDiv}
\end{theorem}

\subsection*{KLdivergence_eq_neg_entropyExpectation}

\begin{theorem}[blueprint="infogeometry:kl:kldivergence_eq_neg_entropyexpectation"]
  \inputleannode{InfoGeometry.KL.KLdivergence_eq_neg_entropyExpectation}
\end{theorem}

\subsection*{KLdivergence_eq_sum_mul_log_densityRatio}

\begin{theorem}[blueprint="infogeometry:kl:kldivergence_eq_sum_mul_log_densityratio"]
  \inputleannode{InfoGeometry.KL.KLdivergence_eq_sum_mul_log_densityRatio}
\end{theorem}

\subsection*{Phi}

\begin{theorem}[blueprint="infogeometry:kl:phi"]
  \inputleannode{InfoGeometry.KL.Phi}
\end{theorem}

\subsection*{PhiSupportFaithful}

\begin{theorem}[blueprint="infogeometry:kl:phisupportfaithful"]
  \inputleannode{InfoGeometry.KL.PhiSupportFaithful}
\end{theorem}

\subsection*{SupportMatches}

\begin{theorem}[blueprint="infogeometry:kl:supportmatches"]
  \inputleannode{InfoGeometry.KL.SupportMatches}
\end{theorem}

\subsection*{absolutelyContinuous_of_fullSupport}

\begin{theorem}[blueprint="infogeometry:kl:absolutelycontinuous_of_fullsupport"]
  \inputleannode{InfoGeometry.KL.absolutelyContinuous_of_fullSupport}
\end{theorem}

\subsection*{densityRatio}

\begin{theorem}[blueprint="infogeometry:kl:densityratio"]
  \inputleannode{InfoGeometry.KL.densityRatio}
\end{theorem}

\subsection*{empiricalDistribution}

\begin{theorem}[blueprint="infogeometry:kl:empiricaldistribution"]
  \inputleannode{InfoGeometry.KL.empiricalDistribution}
\end{theorem}

\subsection*{empiricalProbDist}

\begin{theorem}[blueprint="infogeometry:kl:empiricalprobdist"]
  \inputleannode{InfoGeometry.KL.empiricalProbDist}
\end{theorem}

\subsection*{empirical_nonneg}

\begin{theorem}[blueprint="infogeometry:kl:empirical_nonneg"]
  \inputleannode{InfoGeometry.KL.empirical_nonneg}
\end{theorem}

\subsection*{empirical_sum_one}

\begin{theorem}[blueprint="infogeometry:kl:empirical_sum_one"]
  \inputleannode{InfoGeometry.KL.empirical_sum_one}
\end{theorem}

\subsection*{entropyExpectation}

\begin{theorem}[blueprint="infogeometry:kl:entropyexpectation"]
  \inputleannode{InfoGeometry.KL.entropyExpectation}
\end{theorem}

\subsection*{klDiv}

\begin{theorem}[blueprint="infogeometry:kl:kldiv"]
  \inputleannode{InfoGeometry.KL.klDiv}
\end{theorem}

\subsection*{klDivMeasure}

\begin{theorem}[blueprint="infogeometry:kl:kldivmeasure"]
  \inputleannode{InfoGeometry.KL.klDivMeasure}
\end{theorem}

\subsection*{klDiv_eq_zero_iff}

\begin{theorem}[blueprint="infogeometry:kl:kldiv_eq_zero_iff"]
  \inputleannode{InfoGeometry.KL.klDiv_eq_zero_iff}
\end{theorem}

\subsection*{klDiv_nonneg_of_fullSupport}

\begin{theorem}[blueprint="infogeometry:kl:kldiv_nonneg_of_fullsupport"]
  \inputleannode{InfoGeometry.KL.klDiv_nonneg_of_fullSupport}
\end{theorem}

\subsection*{klDiv_of_not_ac}

\begin{theorem}[blueprint="infogeometry:kl:kldiv_of_not_ac"]
  \inputleannode{InfoGeometry.KL.klDiv_of_not_ac}
\end{theorem}

\subsection*{klDiv_self}

\begin{theorem}[blueprint="infogeometry:kl:kldiv_self"]
  \inputleannode{InfoGeometry.KL.klDiv_self}
\end{theorem}

\subsection*{klDivergence}

\begin{theorem}[blueprint="infogeometry:kl:kldivergence"]
  \inputleannode{InfoGeometry.KL.klDivergence}
\end{theorem}

\subsection*{kl_pointwise_eq_iff}

\begin{theorem}[blueprint="infogeometry:kl:kl_pointwise_eq_iff"]
  \inputleannode{InfoGeometry.KL.kl_pointwise_eq_iff}
\end{theorem}

\subsection*{kl_pointwise_ge_sub}

\begin{theorem}[blueprint="infogeometry:kl:kl_pointwise_ge_sub"]
  \inputleannode{InfoGeometry.KL.kl_pointwise_ge_sub}
\end{theorem}

\subsection*{kl_pointwise_gt_sub_of_ne}

\begin{theorem}[blueprint="infogeometry:kl:kl_pointwise_gt_sub_of_ne"]
  \inputleannode{InfoGeometry.KL.kl_pointwise_gt_sub_of_ne}
\end{theorem}

\subsection*{measureKlDiv}

\begin{theorem}[blueprint="infogeometry:kl:measurekldiv"]
  \inputleannode{InfoGeometry.KL.measureKlDiv}
\end{theorem}

\subsection*{renyiBridge}

\begin{theorem}[blueprint="infogeometry:kl:renyibridge"]
  \inputleannode{InfoGeometry.KL.renyiBridge}
\end{theorem}

\subsection*{renyiBridge_mul}

\begin{theorem}[blueprint="infogeometry:kl:renyibridge_mul"]
  \inputleannode{InfoGeometry.KL.renyiBridge_mul}
\end{theorem}

\subsection*{surprisal}

\begin{theorem}[blueprint="infogeometry:kl:surprisal"]
  \inputleannode{InfoGeometry.KL.surprisal}
\end{theorem}

\subsection*{totalMass}

\begin{theorem}[blueprint="infogeometry:kl:totalmass"]
  \inputleannode{InfoGeometry.KL.totalMass}
\end{theorem}

\section{InfoGeometry.KL.Finite}

\noindent\textbf{Buildable}: yes

\noindent\emph{No exported declarations in the current declaration graph.}

\section{InfoGeometry.KL.KLdivergence}

\subsection*{congr_simp}

\begin{theorem}[blueprint="infogeometry:kl:kldivergence:congr_simp"]
  \inputleannode{InfoGeometry.KL.KLdivergence.congr_simp}
\end{theorem}

\section{InfoGeometry.KL.Measure}

\noindent\textbf{Buildable}: yes

\noindent\emph{No exported declarations in the current declaration graph.}

\section{InfoGeometry.KL.renyiBridge}

\subsection*{congr_simp}

\begin{theorem}[blueprint="infogeometry:kl:renyibridge:congr_simp"]
  \inputleannode{InfoGeometry.KL.renyiBridge.congr_simp}
\end{theorem}

\section{InfoGeometry.KL.totalMass}

\subsection*{eq_1}

\begin{theorem}[blueprint="infogeometry:kl:totalmass:eq_1"]
  \inputleannode{InfoGeometry.KL.totalMass.eq_1}
\end{theorem}

\chapter{Krein Modules}

\section{InfoGeometry.Krein.Automorphisms}

\noindent\textbf{Buildable}: yes

\noindent\emph{No exported declarations in the current declaration graph.}

\section{InfoGeometry.Krein.CartanDecomposition}

\noindent\textbf{Buildable}: yes

\noindent\emph{No exported declarations in the current declaration graph.}

\section{InfoGeometry.Krein.HilbertBridge}

\noindent\textbf{Buildable}: yes

\noindent\emph{No exported declarations in the current declaration graph.}

\section{InfoGeometry.Krein.Metric}

\noindent\textbf{Buildable}: yes

\noindent\emph{No exported declarations in the current declaration graph.}

\section{InfoGeometry.Krein.Modular}

\noindent\textbf{Buildable}: yes

\noindent\emph{No exported declarations in the current declaration graph.}

\section{InfoGeometry.Krein.Prelude}

\noindent\textbf{Buildable}: yes

\noindent\emph{No exported declarations in the current declaration graph.}

\section{InfoGeometry.Krein.Thermal}

\noindent\textbf{Buildable}: yes

\noindent\emph{No exported declarations in the current declaration graph.}

\chapter{LLM Modules}

\section{InfoGeometry.LLM}

\noindent\textbf{Buildable}: yes

\subsection*{CausalMask}

\begin{theorem}[blueprint="infogeometry:llm:causalmask"]
  \inputleannode{InfoGeometry.LLM.CausalMask}
\end{theorem}

\subsection*{IsIdentityMap}

\begin{theorem}[blueprint="infogeometry:llm:isidentitymap"]
  \inputleannode{InfoGeometry.LLM.IsIdentityMap}
\end{theorem}

\subsection*{MaskedTransformerBlock}

\begin{theorem}[blueprint="infogeometry:llm:maskedtransformerblock"]
  \inputleannode{InfoGeometry.LLM.MaskedTransformerBlock}
\end{theorem}

\subsection*{PositionalEncoding}

\begin{theorem}[blueprint="infogeometry:llm:positionalencoding"]
  \inputleannode{InfoGeometry.LLM.PositionalEncoding}
\end{theorem}

\subsection*{TokenContext}

\begin{theorem}[blueprint="infogeometry:llm:tokencontext"]
  \inputleannode{InfoGeometry.LLM.TokenContext}
\end{theorem}

\subsection*{TransformerBlock}

\begin{theorem}[blueprint="infogeometry:llm:transformerblock"]
  \inputleannode{InfoGeometry.LLM.TransformerBlock}
\end{theorem}

\section{InfoGeometry.LLM.CausalMask}

\subsection*{allow}

\begin{theorem}[blueprint="infogeometry:llm:causalmask:allow"]
  \inputleannode{InfoGeometry.LLM.CausalMask.allow}
\end{theorem}

\subsection*{casesOn}

\begin{theorem}[blueprint="infogeometry:llm:causalmask:caseson"]
  \inputleannode{InfoGeometry.LLM.CausalMask.casesOn}
\end{theorem}

\subsection*{ctorIdx}

\begin{theorem}[blueprint="infogeometry:llm:causalmask:ctoridx"]
  \inputleannode{InfoGeometry.LLM.CausalMask.ctorIdx}
\end{theorem}

\subsection*{mk}

\begin{theorem}[blueprint="infogeometry:llm:causalmask:mk"]
  \inputleannode{InfoGeometry.LLM.CausalMask.mk}
\end{theorem}

\subsection*{noConfusion}

\begin{theorem}[blueprint="infogeometry:llm:causalmask:noconfusion"]
  \inputleannode{InfoGeometry.LLM.CausalMask.noConfusion}
\end{theorem}

\subsection*{noConfusionType}

\begin{theorem}[blueprint="infogeometry:llm:causalmask:noconfusiontype"]
  \inputleannode{InfoGeometry.LLM.CausalMask.noConfusionType}
\end{theorem}

\subsection*{rec}

\begin{theorem}[blueprint="infogeometry:llm:causalmask:rec"]
  \inputleannode{InfoGeometry.LLM.CausalMask.rec}
\end{theorem}

\subsection*{recOn}

\begin{theorem}[blueprint="infogeometry:llm:causalmask:recon"]
  \inputleannode{InfoGeometry.LLM.CausalMask.recOn}
\end{theorem}

\section{InfoGeometry.LLM.MaskedTransformerBlock}

\noindent\textbf{Buildable}: yes

\subsection*{base}

\begin{theorem}[blueprint="infogeometry:llm:maskedtransformerblock:base"]
  \inputleannode{InfoGeometry.LLM.MaskedTransformerBlock.base}
\end{theorem}

\subsection*{casesOn}

\begin{theorem}[blueprint="infogeometry:llm:maskedtransformerblock:caseson"]
  \inputleannode{InfoGeometry.LLM.MaskedTransformerBlock.casesOn}
\end{theorem}

\subsection*{causalMask}

\begin{theorem}[blueprint="infogeometry:llm:maskedtransformerblock:causalmask"]
  \inputleannode{InfoGeometry.LLM.MaskedTransformerBlock.causalMask}
\end{theorem}

\subsection*{ctorIdx}

\begin{theorem}[blueprint="infogeometry:llm:maskedtransformerblock:ctoridx"]
  \inputleannode{InfoGeometry.LLM.MaskedTransformerBlock.ctorIdx}
\end{theorem}

\subsection*{maskedAfterAttention}

\begin{theorem}[blueprint="infogeometry:llm:maskedtransformerblock:maskedafterattention"]
  \inputleannode{InfoGeometry.LLM.MaskedTransformerBlock.maskedAfterAttention}
\end{theorem}

\subsection*{maskedAfterMLP}

\begin{theorem}[blueprint="infogeometry:llm:maskedtransformerblock:maskedaftermlp"]
  \inputleannode{InfoGeometry.LLM.MaskedTransformerBlock.maskedAfterMLP}
\end{theorem}

\subsection*{maskedAttentionOut}

\begin{theorem}[blueprint="infogeometry:llm:maskedtransformerblock:maskedattentionout"]
  \inputleannode{InfoGeometry.LLM.MaskedTransformerBlock.maskedAttentionOut}
\end{theorem}

\subsection*{maskedAttentionOut_eq_base_attentionOut_of_allows_all}

\begin{theorem}[blueprint="infogeometry:llm:maskedtransformerblock:maskedattentionout_eq_base_attentionout_of_allows_all"]
  \inputleannode{InfoGeometry.LLM.MaskedTransformerBlock.maskedAttentionOut_eq_base_attentionOut_of_allows_all}
\end{theorem}

\subsection*{maskedHeadWeight}

\begin{theorem}[blueprint="infogeometry:llm:maskedtransformerblock:maskedheadweight"]
  \inputleannode{InfoGeometry.LLM.MaskedTransformerBlock.maskedHeadWeight}
\end{theorem}

\subsection*{maskedHeadWeight_of_allow}

\begin{theorem}[blueprint="infogeometry:llm:maskedtransformerblock:maskedheadweight_of_allow"]
  \inputleannode{InfoGeometry.LLM.MaskedTransformerBlock.maskedHeadWeight_of_allow}
\end{theorem}

\subsection*{maskedHeadWeight_of_not_allow}

\begin{theorem}[blueprint="infogeometry:llm:maskedtransformerblock:maskedheadweight_of_not_allow"]
  \inputleannode{InfoGeometry.LLM.MaskedTransformerBlock.maskedHeadWeight_of_not_allow}
\end{theorem}

\subsection*{maskedPreOutput}

\begin{theorem}[blueprint="infogeometry:llm:maskedtransformerblock:maskedpreoutput"]
  \inputleannode{InfoGeometry.LLM.MaskedTransformerBlock.maskedPreOutput}
\end{theorem}

\subsection*{maskedPreOutput_eq_base_preOutput_of_allows_all}

\begin{theorem}[blueprint="infogeometry:llm:maskedtransformerblock:maskedpreoutput_eq_base_preoutput_of_allows_all"]
  \inputleannode{InfoGeometry.LLM.MaskedTransformerBlock.maskedPreOutput_eq_base_preOutput_of_allows_all}
\end{theorem}

\subsection*{mk}

\begin{theorem}[blueprint="infogeometry:llm:maskedtransformerblock:mk"]
  \inputleannode{InfoGeometry.LLM.MaskedTransformerBlock.mk}
\end{theorem}

\subsection*{noConfusion}

\begin{theorem}[blueprint="infogeometry:llm:maskedtransformerblock:noconfusion"]
  \inputleannode{InfoGeometry.LLM.MaskedTransformerBlock.noConfusion}
\end{theorem}

\subsection*{noConfusionType}

\begin{theorem}[blueprint="infogeometry:llm:maskedtransformerblock:noconfusiontype"]
  \inputleannode{InfoGeometry.LLM.MaskedTransformerBlock.noConfusionType}
\end{theorem}

\subsection*{rec}

\begin{theorem}[blueprint="infogeometry:llm:maskedtransformerblock:rec"]
  \inputleannode{InfoGeometry.LLM.MaskedTransformerBlock.rec}
\end{theorem}

\subsection*{recOn}

\begin{theorem}[blueprint="infogeometry:llm:maskedtransformerblock:recon"]
  \inputleannode{InfoGeometry.LLM.MaskedTransformerBlock.recOn}
\end{theorem}

\subsection*{run}

\begin{theorem}[blueprint="infogeometry:llm:maskedtransformerblock:run"]
  \inputleannode{InfoGeometry.LLM.MaskedTransformerBlock.run}
\end{theorem}

\subsection*{runWithPositionalEncoding}

\begin{theorem}[blueprint="infogeometry:llm:maskedtransformerblock:runwithpositionalencoding"]
  \inputleannode{InfoGeometry.LLM.MaskedTransformerBlock.runWithPositionalEncoding}
\end{theorem}

\subsection*{runWithPositionalEncoding_comp}

\begin{theorem}[blueprint="infogeometry:llm:maskedtransformerblock:runwithpositionalencoding_comp"]
  \inputleannode{InfoGeometry.LLM.MaskedTransformerBlock.runWithPositionalEncoding_comp}
\end{theorem}

\subsection*{runWithPositionalEncoding_eq_comp}

\begin{theorem}[blueprint="infogeometry:llm:maskedtransformerblock:runwithpositionalencoding_eq_comp"]
  \inputleannode{InfoGeometry.LLM.MaskedTransformerBlock.runWithPositionalEncoding_eq_comp}
\end{theorem}

\subsection*{runWithPositionalEncoding_id}

\begin{theorem}[blueprint="infogeometry:llm:maskedtransformerblock:runwithpositionalencoding_id"]
  \inputleannode{InfoGeometry.LLM.MaskedTransformerBlock.runWithPositionalEncoding_id}
\end{theorem}

\subsection*{run_eq_base_run_of_allows_all}

\begin{theorem}[blueprint="infogeometry:llm:maskedtransformerblock:run_eq_base_run_of_allows_all"]
  \inputleannode{InfoGeometry.LLM.MaskedTransformerBlock.run_eq_base_run_of_allows_all}
\end{theorem}

\section{InfoGeometry.LLM.MaskedTransformerBlock.maskedHeadWeight}

\subsection*{eq_1}

\begin{theorem}[blueprint="infogeometry:llm:maskedtransformerblock:maskedheadweight:eq_1"]
  \inputleannode{InfoGeometry.LLM.MaskedTransformerBlock.maskedHeadWeight.eq_1}
\end{theorem}

\section{InfoGeometry.LLM.PositionalEncoding}

\noindent\textbf{Buildable}: yes

\subsection*{casesOn}

\begin{theorem}[blueprint="infogeometry:llm:positionalencoding:caseson"]
  \inputleannode{InfoGeometry.LLM.PositionalEncoding.casesOn}
\end{theorem}

\subsection*{comp}

\begin{theorem}[blueprint="infogeometry:llm:positionalencoding:comp"]
  \inputleannode{InfoGeometry.LLM.PositionalEncoding.comp}
\end{theorem}

\subsection*{comp_apply}

\begin{theorem}[blueprint="infogeometry:llm:positionalencoding:comp_apply"]
  \inputleannode{InfoGeometry.LLM.PositionalEncoding.comp_apply}
\end{theorem}

\subsection*{ctorIdx}

\begin{theorem}[blueprint="infogeometry:llm:positionalencoding:ctoridx"]
  \inputleannode{InfoGeometry.LLM.PositionalEncoding.ctorIdx}
\end{theorem}

\subsection*{encode}

\begin{theorem}[blueprint="infogeometry:llm:positionalencoding:encode"]
  \inputleannode{InfoGeometry.LLM.PositionalEncoding.encode}
\end{theorem}

\subsection*{id}

\begin{theorem}[blueprint="infogeometry:llm:positionalencoding:id"]
  \inputleannode{InfoGeometry.LLM.PositionalEncoding.id}
\end{theorem}

\subsection*{id_apply}

\begin{theorem}[blueprint="infogeometry:llm:positionalencoding:id_apply"]
  \inputleannode{InfoGeometry.LLM.PositionalEncoding.id_apply}
\end{theorem}

\subsection*{mk}

\begin{theorem}[blueprint="infogeometry:llm:positionalencoding:mk"]
  \inputleannode{InfoGeometry.LLM.PositionalEncoding.mk}
\end{theorem}

\subsection*{noConfusion}

\begin{theorem}[blueprint="infogeometry:llm:positionalencoding:noconfusion"]
  \inputleannode{InfoGeometry.LLM.PositionalEncoding.noConfusion}
\end{theorem}

\subsection*{noConfusionType}

\begin{theorem}[blueprint="infogeometry:llm:positionalencoding:noconfusiontype"]
  \inputleannode{InfoGeometry.LLM.PositionalEncoding.noConfusionType}
\end{theorem}

\subsection*{rec}

\begin{theorem}[blueprint="infogeometry:llm:positionalencoding:rec"]
  \inputleannode{InfoGeometry.LLM.PositionalEncoding.rec}
\end{theorem}

\subsection*{recOn}

\begin{theorem}[blueprint="infogeometry:llm:positionalencoding:recon"]
  \inputleannode{InfoGeometry.LLM.PositionalEncoding.recOn}
\end{theorem}

\section{InfoGeometry.LLM.TransformerBlock}

\noindent\textbf{Buildable}: yes

\subsection*{afterAttention}

\begin{theorem}[blueprint="infogeometry:llm:transformerblock:afterattention"]
  \inputleannode{InfoGeometry.LLM.TransformerBlock.afterAttention}
\end{theorem}

\subsection*{afterAttention_def}

\begin{theorem}[blueprint="infogeometry:llm:transformerblock:afterattention_def"]
  \inputleannode{InfoGeometry.LLM.TransformerBlock.afterAttention_def}
\end{theorem}

\subsection*{afterAttention_eq_residual_add_attention}

\begin{theorem}[blueprint="infogeometry:llm:transformerblock:afterattention_eq_residual_add_attention"]
  \inputleannode{InfoGeometry.LLM.TransformerBlock.afterAttention_eq_residual_add_attention}
\end{theorem}

\subsection*{afterMLP}

\begin{theorem}[blueprint="infogeometry:llm:transformerblock:aftermlp"]
  \inputleannode{InfoGeometry.LLM.TransformerBlock.afterMLP}
\end{theorem}

\subsection*{afterMLP_def}

\begin{theorem}[blueprint="infogeometry:llm:transformerblock:aftermlp_def"]
  \inputleannode{InfoGeometry.LLM.TransformerBlock.afterMLP_def}
\end{theorem}

\subsection*{afterMLP_eq_afterAttention_of_zero_mlp}

\begin{theorem}[blueprint="infogeometry:llm:transformerblock:aftermlp_eq_afterattention_of_zero_mlp"]
  \inputleannode{InfoGeometry.LLM.TransformerBlock.afterMLP_eq_afterAttention_of_zero_mlp}
\end{theorem}

\subsection*{attentionOut}

\begin{theorem}[blueprint="infogeometry:llm:transformerblock:attentionout"]
  \inputleannode{InfoGeometry.LLM.TransformerBlock.attentionOut}
\end{theorem}

\subsection*{attentionOut_decomposition_residual}

\begin{theorem}[blueprint="infogeometry:llm:transformerblock:attentionout_decomposition_residual"]
  \inputleannode{InfoGeometry.LLM.TransformerBlock.attentionOut_decomposition_residual}
\end{theorem}

\subsection*{attentionOut_def}

\begin{theorem}[blueprint="infogeometry:llm:transformerblock:attentionout_def"]
  \inputleannode{InfoGeometry.LLM.TransformerBlock.attentionOut_def}
\end{theorem}

\subsection*{attn}

\begin{theorem}[blueprint="infogeometry:llm:transformerblock:attn"]
  \inputleannode{InfoGeometry.LLM.TransformerBlock.attn}
\end{theorem}

\subsection*{casesOn}

\begin{theorem}[blueprint="infogeometry:llm:transformerblock:caseson"]
  \inputleannode{InfoGeometry.LLM.TransformerBlock.casesOn}
\end{theorem}

\subsection*{ctorIdx}

\begin{theorem}[blueprint="infogeometry:llm:transformerblock:ctoridx"]
  \inputleannode{InfoGeometry.LLM.TransformerBlock.ctorIdx}
\end{theorem}

\subsection*{mk}

\begin{theorem}[blueprint="infogeometry:llm:transformerblock:mk"]
  \inputleannode{InfoGeometry.LLM.TransformerBlock.mk}
\end{theorem}

\subsection*{mlp}

\begin{theorem}[blueprint="infogeometry:llm:transformerblock:mlp"]
  \inputleannode{InfoGeometry.LLM.TransformerBlock.mlp}
\end{theorem}

\subsection*{noConfusion}

\begin{theorem}[blueprint="infogeometry:llm:transformerblock:noconfusion"]
  \inputleannode{InfoGeometry.LLM.TransformerBlock.noConfusion}
\end{theorem}

\subsection*{noConfusionType}

\begin{theorem}[blueprint="infogeometry:llm:transformerblock:noconfusiontype"]
  \inputleannode{InfoGeometry.LLM.TransformerBlock.noConfusionType}
\end{theorem}

\subsection*{norm1}

\begin{theorem}[blueprint="infogeometry:llm:transformerblock:norm1"]
  \inputleannode{InfoGeometry.LLM.TransformerBlock.norm1}
\end{theorem}

\subsection*{norm2}

\begin{theorem}[blueprint="infogeometry:llm:transformerblock:norm2"]
  \inputleannode{InfoGeometry.LLM.TransformerBlock.norm2}
\end{theorem}

\subsection*{rec}

\begin{theorem}[blueprint="infogeometry:llm:transformerblock:rec"]
  \inputleannode{InfoGeometry.LLM.TransformerBlock.rec}
\end{theorem}

\subsection*{recOn}

\begin{theorem}[blueprint="infogeometry:llm:transformerblock:recon"]
  \inputleannode{InfoGeometry.LLM.TransformerBlock.recOn}
\end{theorem}

\subsection*{residualLift}

\begin{theorem}[blueprint="infogeometry:llm:transformerblock:residuallift"]
  \inputleannode{InfoGeometry.LLM.TransformerBlock.residualLift}
\end{theorem}

\subsection*{run}

\begin{theorem}[blueprint="infogeometry:llm:transformerblock:run"]
  \inputleannode{InfoGeometry.LLM.TransformerBlock.run}
\end{theorem}

\subsection*{runWithPositionalEncoding}

\begin{theorem}[blueprint="infogeometry:llm:transformerblock:runwithpositionalencoding"]
  \inputleannode{InfoGeometry.LLM.TransformerBlock.runWithPositionalEncoding}
\end{theorem}

\subsection*{runWithPositionalEncoding_comp}

\begin{theorem}[blueprint="infogeometry:llm:transformerblock:runwithpositionalencoding_comp"]
  \inputleannode{InfoGeometry.LLM.TransformerBlock.runWithPositionalEncoding_comp}
\end{theorem}

\subsection*{runWithPositionalEncoding_eq_comp}

\begin{theorem}[blueprint="infogeometry:llm:transformerblock:runwithpositionalencoding_eq_comp"]
  \inputleannode{InfoGeometry.LLM.TransformerBlock.runWithPositionalEncoding_eq_comp}
\end{theorem}

\subsection*{runWithPositionalEncoding_id}

\begin{theorem}[blueprint="infogeometry:llm:transformerblock:runwithpositionalencoding_id"]
  \inputleannode{InfoGeometry.LLM.TransformerBlock.runWithPositionalEncoding_id}
\end{theorem}

\subsection*{run_def}

\begin{theorem}[blueprint="infogeometry:llm:transformerblock:run_def"]
  \inputleannode{InfoGeometry.LLM.TransformerBlock.run_def}
\end{theorem}

\subsection*{run_eq_residual_add_attention_of_trivial_post}

\begin{theorem}[blueprint="infogeometry:llm:transformerblock:run_eq_residual_add_attention_of_trivial_post"]
  \inputleannode{InfoGeometry.LLM.TransformerBlock.run_eq_residual_add_attention_of_trivial_post}
\end{theorem}

\section{InfoGeometry.LLM.TransformerBlock.attentionOut}

\subsection*{eq_1}

\begin{theorem}[blueprint="infogeometry:llm:transformerblock:attentionout:eq_1"]
  \inputleannode{InfoGeometry.LLM.TransformerBlock.attentionOut.eq_1}
\end{theorem}

\chapter{Library Modules}

\section{InfoGeometry.Library}

\noindent\textbf{Buildable}: yes

\noindent\emph{No exported declarations in the current declaration graph.}

\chapter{LogPotential Modules}

\section{InfoGeometry.LogPotential}

\subsection*{LegendreModel}

\begin{theorem}[blueprint="infogeometry:logpotential:legendremodel"]
  \inputleannode{InfoGeometry.LogPotential.LegendreModel}
\end{theorem}

\subsection*{bregman}

\begin{theorem}[blueprint="infogeometry:logpotential:bregman"]
  \inputleannode{InfoGeometry.LogPotential.bregman}
\end{theorem}

\subsection*{bregman_def}

\begin{theorem}[blueprint="infogeometry:logpotential:bregman_def"]
  \inputleannode{InfoGeometry.LogPotential.bregman_def}
\end{theorem}

\subsection*{casesOn}

\begin{theorem}[blueprint="infogeometry:logpotential:caseson"]
  \inputleannode{InfoGeometry.LogPotential.casesOn}
\end{theorem}

\subsection*{ctorIdx}

\begin{theorem}[blueprint="infogeometry:logpotential:ctoridx"]
  \inputleannode{InfoGeometry.LogPotential.ctorIdx}
\end{theorem}

\subsection*{mk}

\begin{theorem}[blueprint="infogeometry:logpotential:mk"]
  \inputleannode{InfoGeometry.LogPotential.mk}
\end{theorem}

\subsection*{noConfusion}

\begin{theorem}[blueprint="infogeometry:logpotential:noconfusion"]
  \inputleannode{InfoGeometry.LogPotential.noConfusion}
\end{theorem}

\subsection*{noConfusionType}

\begin{theorem}[blueprint="infogeometry:logpotential:noconfusiontype"]
  \inputleannode{InfoGeometry.LogPotential.noConfusionType}
\end{theorem}

\subsection*{rec}

\begin{theorem}[blueprint="infogeometry:logpotential:rec"]
  \inputleannode{InfoGeometry.LogPotential.rec}
\end{theorem}

\subsection*{recOn}

\begin{theorem}[blueprint="infogeometry:logpotential:recon"]
  \inputleannode{InfoGeometry.LogPotential.recOn}
\end{theorem}

\subsection*{ψ}

\begin{theorem}[blueprint="infogeometry:logpotential:ψ"]
  \inputleannode{InfoGeometry.LogPotential.ψ}
\end{theorem}

\section{InfoGeometry.LogPotential.LegendreModel}

\subsection*{L}

\begin{theorem}[blueprint="infogeometry:logpotential:legendremodel:l"]
  \inputleannode{InfoGeometry.LogPotential.LegendreModel.L}
\end{theorem}

\subsection*{canonicalEnergy}

\begin{theorem}[blueprint="infogeometry:logpotential:legendremodel:canonicalenergy"]
  \inputleannode{InfoGeometry.LogPotential.LegendreModel.canonicalEnergy}
\end{theorem}

\subsection*{canonicalEnergy_def}

\begin{theorem}[blueprint="infogeometry:logpotential:legendremodel:canonicalenergy_def"]
  \inputleannode{InfoGeometry.LogPotential.LegendreModel.canonicalEnergy_def}
\end{theorem}

\subsection*{canonicalEntropy}

\begin{theorem}[blueprint="infogeometry:logpotential:legendremodel:canonicalentropy"]
  \inputleannode{InfoGeometry.LogPotential.LegendreModel.canonicalEntropy}
\end{theorem}

\subsection*{canonicalEntropy_def}

\begin{theorem}[blueprint="infogeometry:logpotential:legendremodel:canonicalentropy_def"]
  \inputleannode{InfoGeometry.LogPotential.LegendreModel.canonicalEntropy_def}
\end{theorem}

\subsection*{canonicalEntropy_eq_neg_dual_at_contact}

\begin{theorem}[blueprint="infogeometry:logpotential:legendremodel:canonicalentropy_eq_neg_dual_at_contact"]
  \inputleannode{InfoGeometry.LogPotential.LegendreModel.canonicalEntropy_eq_neg_dual_at_contact}
\end{theorem}

\subsection*{canonicalEntropy_eq_neg_entropy}

\begin{theorem}[blueprint="infogeometry:logpotential:legendremodel:canonicalentropy_eq_neg_entropy"]
  \inputleannode{InfoGeometry.LogPotential.LegendreModel.canonicalEntropy_eq_neg_entropy}
\end{theorem}

\subsection*{canonicalFreeEnergy}

\begin{theorem}[blueprint="infogeometry:logpotential:legendremodel:canonicalfreeenergy"]
  \inputleannode{InfoGeometry.LogPotential.LegendreModel.canonicalFreeEnergy}
\end{theorem}

\subsection*{canonicalFreeEnergy_def}

\begin{theorem}[blueprint="infogeometry:logpotential:legendremodel:canonicalfreeenergy_def"]
  \inputleannode{InfoGeometry.LogPotential.LegendreModel.canonicalFreeEnergy_def}
\end{theorem}

\subsection*{canonicalFreeEnergy_eq_scaled_entropy_energy}

\begin{theorem}[blueprint="infogeometry:logpotential:legendremodel:canonicalfreeenergy_eq_scaled_entropy_energy"]
  \inputleannode{InfoGeometry.LogPotential.LegendreModel.canonicalFreeEnergy_eq_scaled_entropy_energy}
\end{theorem}

\subsection*{casesOn}

\begin{theorem}[blueprint="infogeometry:logpotential:legendremodel:caseson"]
  \inputleannode{InfoGeometry.LogPotential.LegendreModel.casesOn}
\end{theorem}

\subsection*{contact}

\begin{theorem}[blueprint="infogeometry:logpotential:legendremodel:contact"]
  \inputleannode{InfoGeometry.LogPotential.LegendreModel.contact}
\end{theorem}

\subsection*{contact_balance}

\begin{theorem}[blueprint="infogeometry:logpotential:legendremodel:contact_balance"]
  \inputleannode{InfoGeometry.LogPotential.LegendreModel.contact_balance}
\end{theorem}

\subsection*{ctorIdx}

\begin{theorem}[blueprint="infogeometry:logpotential:legendremodel:ctoridx"]
  \inputleannode{InfoGeometry.LogPotential.LegendreModel.ctorIdx}
\end{theorem}

\subsection*{dualBregman}

\begin{theorem}[blueprint="infogeometry:logpotential:legendremodel:dualbregman"]
  \inputleannode{InfoGeometry.LogPotential.LegendreModel.dualBregman}
\end{theorem}

\subsection*{dualCoord}

\begin{theorem}[blueprint="infogeometry:logpotential:legendremodel:dualcoord"]
  \inputleannode{InfoGeometry.LogPotential.LegendreModel.dualCoord}
\end{theorem}

\subsection*{dualCoord_def}

\begin{theorem}[blueprint="infogeometry:logpotential:legendremodel:dualcoord_def"]
  \inputleannode{InfoGeometry.LogPotential.LegendreModel.dualCoord_def}
\end{theorem}

\subsection*{entropy}

\begin{theorem}[blueprint="infogeometry:logpotential:legendremodel:entropy"]
  \inputleannode{InfoGeometry.LogPotential.LegendreModel.entropy}
\end{theorem}

\subsection*{entropy_eq_dual_at_contact}

\begin{theorem}[blueprint="infogeometry:logpotential:legendremodel:entropy_eq_dual_at_contact"]
  \inputleannode{InfoGeometry.LogPotential.LegendreModel.entropy_eq_dual_at_contact}
\end{theorem}

\subsection*{fenchelGap}

\begin{theorem}[blueprint="infogeometry:logpotential:legendremodel:fenchelgap"]
  \inputleannode{InfoGeometry.LogPotential.LegendreModel.fenchelGap}
\end{theorem}

\subsection*{fenchelGap_def}

\begin{theorem}[blueprint="infogeometry:logpotential:legendremodel:fenchelgap_def"]
  \inputleannode{InfoGeometry.LogPotential.LegendreModel.fenchelGap_def}
\end{theorem}

\subsection*{fenchelGap_eq_zero_at_contact}

\begin{theorem}[blueprint="infogeometry:logpotential:legendremodel:fenchelgap_eq_zero_at_contact"]
  \inputleannode{InfoGeometry.LogPotential.LegendreModel.fenchelGap_eq_zero_at_contact}
\end{theorem}

\subsection*{fenchelGap_nonneg}

\begin{theorem}[blueprint="infogeometry:logpotential:legendremodel:fenchelgap_nonneg"]
  \inputleannode{InfoGeometry.LogPotential.LegendreModel.fenchelGap_nonneg}
\end{theorem}

\subsection*{fenchel_ineq}

\begin{theorem}[blueprint="infogeometry:logpotential:legendremodel:fenchel_ineq"]
  \inputleannode{InfoGeometry.LogPotential.LegendreModel.fenchel_ineq}
\end{theorem}

\subsection*{freeEnergy}

\begin{theorem}[blueprint="infogeometry:logpotential:legendremodel:freeenergy"]
  \inputleannode{InfoGeometry.LogPotential.LegendreModel.freeEnergy}
\end{theorem}

\subsection*{freeEnergy_def}

\begin{theorem}[blueprint="infogeometry:logpotential:legendremodel:freeenergy_def"]
  \inputleannode{InfoGeometry.LogPotential.LegendreModel.freeEnergy_def}
\end{theorem}

\subsection*{grad}

\begin{theorem}[blueprint="infogeometry:logpotential:legendremodel:grad"]
  \inputleannode{InfoGeometry.LogPotential.LegendreModel.grad}
\end{theorem}

\subsection*{massieu}

\begin{theorem}[blueprint="infogeometry:logpotential:legendremodel:massieu"]
  \inputleannode{InfoGeometry.LogPotential.LegendreModel.massieu}
\end{theorem}

\subsection*{massieu_def}

\begin{theorem}[blueprint="infogeometry:logpotential:legendremodel:massieu_def"]
  \inputleannode{InfoGeometry.LogPotential.LegendreModel.massieu_def}
\end{theorem}

\subsection*{massieu_eq_canonicalEntropy_minus_theta_energy}

\begin{theorem}[blueprint="infogeometry:logpotential:legendremodel:massieu_eq_canonicalentropy_minus_theta_energy"]
  \inputleannode{InfoGeometry.LogPotential.LegendreModel.massieu_eq_canonicalEntropy_minus_theta_energy}
\end{theorem}

\subsection*{mk}

\begin{theorem}[blueprint="infogeometry:logpotential:legendremodel:mk"]
  \inputleannode{InfoGeometry.LogPotential.LegendreModel.mk}
\end{theorem}

\subsection*{noConfusion}

\begin{theorem}[blueprint="infogeometry:logpotential:legendremodel:noconfusion"]
  \inputleannode{InfoGeometry.LogPotential.LegendreModel.noConfusion}
\end{theorem}

\subsection*{noConfusionType}

\begin{theorem}[blueprint="infogeometry:logpotential:legendremodel:noconfusiontype"]
  \inputleannode{InfoGeometry.LogPotential.LegendreModel.noConfusionType}
\end{theorem}

\subsection*{partition}

\begin{theorem}[blueprint="infogeometry:logpotential:legendremodel:partition"]
  \inputleannode{InfoGeometry.LogPotential.LegendreModel.partition}
\end{theorem}

\subsection*{partition_def}

\begin{theorem}[blueprint="infogeometry:logpotential:legendremodel:partition_def"]
  \inputleannode{InfoGeometry.LogPotential.LegendreModel.partition_def}
\end{theorem}

\subsection*{primalBregman}

\begin{theorem}[blueprint="infogeometry:logpotential:legendremodel:primalbregman"]
  \inputleannode{InfoGeometry.LogPotential.LegendreModel.primalBregman}
\end{theorem}

\subsection*{rec}

\begin{theorem}[blueprint="infogeometry:logpotential:legendremodel:rec"]
  \inputleannode{InfoGeometry.LogPotential.LegendreModel.rec}
\end{theorem}

\subsection*{recOn}

\begin{theorem}[blueprint="infogeometry:logpotential:legendremodel:recon"]
  \inputleannode{InfoGeometry.LogPotential.LegendreModel.recOn}
\end{theorem}

\subsection*{scaledFenchelGap}

\begin{theorem}[blueprint="infogeometry:logpotential:legendremodel:scaledfenchelgap"]
  \inputleannode{InfoGeometry.LogPotential.LegendreModel.scaledFenchelGap}
\end{theorem}

\subsection*{scaledFenchelGap_def}

\begin{theorem}[blueprint="infogeometry:logpotential:legendremodel:scaledfenchelgap_def"]
  \inputleannode{InfoGeometry.LogPotential.LegendreModel.scaledFenchelGap_def}
\end{theorem}

\subsection*{scaledFenchelGap_eq_zero_at_contact}

\begin{theorem}[blueprint="infogeometry:logpotential:legendremodel:scaledfenchelgap_eq_zero_at_contact"]
  \inputleannode{InfoGeometry.LogPotential.LegendreModel.scaledFenchelGap_eq_zero_at_contact}
\end{theorem}

\subsection*{scaledFenchelGap_nonneg}

\begin{theorem}[blueprint="infogeometry:logpotential:legendremodel:scaledfenchelgap_nonneg"]
  \inputleannode{InfoGeometry.LogPotential.LegendreModel.scaledFenchelGap_nonneg}
\end{theorem}

\subsection*{φ}

\begin{theorem}[blueprint="infogeometry:logpotential:legendremodel:φ"]
  \inputleannode{InfoGeometry.LogPotential.LegendreModel.φ}
\end{theorem}

\section{InfoGeometry.LogPotential.LegendreModel.fenchelGap}

\subsection*{eq_1}

\begin{theorem}[blueprint="infogeometry:logpotential:legendremodel:fenchelgap:eq_1"]
  \inputleannode{InfoGeometry.LogPotential.LegendreModel.fenchelGap.eq_1}
\end{theorem}

\section{InfoGeometry.LogPotential.LegendreModel.scaledFenchelGap}

\subsection*{eq_1}

\begin{theorem}[blueprint="infogeometry:logpotential:legendremodel:scaledfenchelgap:eq_1"]
  \inputleannode{InfoGeometry.LogPotential.LegendreModel.scaledFenchelGap.eq_1}
\end{theorem}

\chapter{ManifoldTopology Modules}

\section{InfoGeometry.ManifoldTopology}

\subsection*{exists_isolating_nhds_of_nondegenerate}

\begin{theorem}[blueprint="infogeometry:manifoldtopology:exists_isolating_nhds_of_nondegenerate"]
  \inputleannode{InfoGeometry.ManifoldTopology.exists_isolating_nhds_of_nondegenerate}
\end{theorem}

\subsection*{exists_local_chart_homotopy_to_linear}

\begin{theorem}[blueprint="infogeometry:manifoldtopology:exists_local_chart_homotopy_to_linear"]
  \inputleannode{InfoGeometry.ManifoldTopology.exists_local_chart_homotopy_to_linear}
\end{theorem}

\subsection*{local_degree_eq_sign_jacDet}

\begin{theorem}[blueprint="infogeometry:manifoldtopology:local_degree_eq_sign_jacdet"]
  \inputleannode{InfoGeometry.ManifoldTopology.local_degree_eq_sign_jacDet}
\end{theorem}

\subsection*{preimage_finite_of_regular_value}

\begin{theorem}[blueprint="infogeometry:manifoldtopology:preimage_finite_of_regular_value"]
  \inputleannode{InfoGeometry.ManifoldTopology.preimage_finite_of_regular_value}
\end{theorem}

\chapter{MaxEnt Modules}

\section{InfoGeometry.MaxEnt}

\noindent\textbf{Buildable}: yes

\subsection*{ACProbMeasure}

\begin{theorem}[blueprint="infogeometry:maxent:acprobmeasure"]
  \inputleannode{InfoGeometry.MaxEnt.ACProbMeasure}
\end{theorem}

\subsection*{AsymptoticEquipartitionWitness}

\begin{theorem}[blueprint="infogeometry:maxent:asymptoticequipartitionwitness"]
  \inputleannode{InfoGeometry.MaxEnt.AsymptoticEquipartitionWitness}
\end{theorem}

\subsection*{BurgModel}

\begin{theorem}[blueprint="infogeometry:maxent:burgmodel"]
  \inputleannode{InfoGeometry.MaxEnt.BurgModel}
\end{theorem}

\subsection*{CanonicalProblem}

\begin{theorem}[blueprint="infogeometry:maxent:canonicalproblem"]
  \inputleannode{InfoGeometry.MaxEnt.CanonicalProblem}
\end{theorem}

\subsection*{ConstraintsIntegrableOn}

\begin{theorem}[blueprint="infogeometry:maxent:constraintsintegrableon"]
  \inputleannode{InfoGeometry.MaxEnt.ConstraintsIntegrableOn}
\end{theorem}

\subsection*{ConstraintsIntegrableOn_of_memFeasibleIntegrable}

\begin{theorem}[blueprint="infogeometry:maxent:constraintsintegrableon_of_memfeasibleintegrable"]
  \inputleannode{InfoGeometry.MaxEnt.ConstraintsIntegrableOn_of_memFeasibleIntegrable}
\end{theorem}

\subsection*{EntropyBand}

\begin{theorem}[blueprint="infogeometry:maxent:entropyband"]
  \inputleannode{InfoGeometry.MaxEnt.EntropyBand}
\end{theorem}

\subsection*{EntropyConcentrationCertificate}

\begin{theorem}[blueprint="infogeometry:maxent:entropyconcentrationcertificate"]
  \inputleannode{InfoGeometry.MaxEnt.EntropyConcentrationCertificate}
\end{theorem}

\subsection*{Feasible}

\begin{theorem}[blueprint="infogeometry:maxent:feasible"]
  \inputleannode{InfoGeometry.MaxEnt.Feasible}
\end{theorem}

\subsection*{FeasibleClass}

\begin{theorem}[blueprint="infogeometry:maxent:feasibleclass"]
  \inputleannode{InfoGeometry.MaxEnt.FeasibleClass}
\end{theorem}

\subsection*{FeasibleIntegrable}

\begin{theorem}[blueprint="infogeometry:maxent:feasibleintegrable"]
  \inputleannode{InfoGeometry.MaxEnt.FeasibleIntegrable}
\end{theorem}

\subsection*{FeasibleIntegrable_subset_Feasible}

\begin{theorem}[blueprint="infogeometry:maxent:feasibleintegrable_subset_feasible"]
  \inputleannode{InfoGeometry.MaxEnt.FeasibleIntegrable_subset_Feasible}
\end{theorem}

\subsection*{FenchelDualAssumptions}

\begin{theorem}[blueprint="infogeometry:maxent:fencheldualassumptions"]
  \inputleannode{InfoGeometry.MaxEnt.FenchelDualAssumptions}
\end{theorem}

\subsection*{FiniteParamSpace}

\begin{theorem}[blueprint="infogeometry:maxent:finiteparamspace"]
  \inputleannode{InfoGeometry.MaxEnt.FiniteParamSpace}
\end{theorem}

\subsection*{HasExponentialRNForm}

\begin{theorem}[blueprint="infogeometry:maxent:hasexponentialrnform"]
  \inputleannode{InfoGeometry.MaxEnt.HasExponentialRNForm}
\end{theorem}

\subsection*{HasExponentialRNFormSigmaFinite}

\begin{theorem}[blueprint="infogeometry:maxent:hasexponentialrnformsigmafinite"]
  \inputleannode{InfoGeometry.MaxEnt.HasExponentialRNFormSigmaFinite}
\end{theorem}

\subsection*{IsMaxEntSolution}

\begin{theorem}[blueprint="infogeometry:maxent:ismaxentsolution"]
  \inputleannode{InfoGeometry.MaxEnt.IsMaxEntSolution}
\end{theorem}

\subsection*{IsMaxEntSolutionIntegrable}

\begin{theorem}[blueprint="infogeometry:maxent:ismaxentsolutionintegrable"]
  \inputleannode{InfoGeometry.MaxEnt.IsMaxEntSolutionIntegrable}
\end{theorem}

\subsection*{LinearConstraint}

\begin{theorem}[blueprint="infogeometry:maxent:linearconstraint"]
  \inputleannode{InfoGeometry.MaxEnt.LinearConstraint}
\end{theorem}

\subsection*{MaxEntConstraint}

\begin{theorem}[blueprint="infogeometry:maxent:maxentconstraint"]
  \inputleannode{InfoGeometry.MaxEnt.MaxEntConstraint}
\end{theorem}

\subsection*{MaxEntProblem}

\begin{theorem}[blueprint="infogeometry:maxent:maxentproblem"]
  \inputleannode{InfoGeometry.MaxEnt.MaxEntProblem}
\end{theorem}

\subsection*{Normalized}

\begin{theorem}[blueprint="infogeometry:maxent:normalized"]
  \inputleannode{InfoGeometry.MaxEnt.Normalized}
\end{theorem}

\subsection*{Objective}

\begin{theorem}[blueprint="infogeometry:maxent:objective"]
  \inputleannode{InfoGeometry.MaxEnt.Objective}
\end{theorem}

\subsection*{PointwiseNonneg}

\begin{theorem}[blueprint="infogeometry:maxent:pointwisenonneg"]
  \inputleannode{InfoGeometry.MaxEnt.PointwiseNonneg}
\end{theorem}

\subsection*{ProbabilitySimplex}

\begin{theorem}[blueprint="infogeometry:maxent:probabilitysimplex"]
  \inputleannode{InfoGeometry.MaxEnt.ProbabilitySimplex}
\end{theorem}

\subsection*{SatisfiesLinearConstraints}

\begin{theorem}[blueprint="infogeometry:maxent:satisfieslinearconstraints"]
  \inputleannode{InfoGeometry.MaxEnt.SatisfiesLinearConstraints}
\end{theorem}

\subsection*{ShannonEntropy}

\begin{theorem}[blueprint="infogeometry:maxent:shannonentropy"]
  \inputleannode{InfoGeometry.MaxEnt.ShannonEntropy}
\end{theorem}

\subsection*{ShannonEntropyXlogx}

\begin{theorem}[blueprint="infogeometry:maxent:shannonentropyxlogx"]
  \inputleannode{InfoGeometry.MaxEnt.ShannonEntropyXlogx}
\end{theorem}

\subsection*{ShannonEntropy_eq_entropy}

\begin{theorem}[blueprint="infogeometry:maxent:shannonentropy_eq_entropy"]
  \inputleannode{InfoGeometry.MaxEnt.ShannonEntropy_eq_entropy}
\end{theorem}

\subsection*{ShannonEntropy_eq_xlogx}

\begin{theorem}[blueprint="infogeometry:maxent:shannonentropy_eq_xlogx"]
  \inputleannode{InfoGeometry.MaxEnt.ShannonEntropy_eq_xlogx}
\end{theorem}

\subsection*{Simplex}

\begin{theorem}[blueprint="infogeometry:maxent:simplex"]
  \inputleannode{InfoGeometry.MaxEnt.Simplex}
\end{theorem}

\subsection*{arSpectralDensity}

\begin{theorem}[blueprint="infogeometry:maxent:arspectraldensity"]
  \inputleannode{InfoGeometry.MaxEnt.arSpectralDensity}
\end{theorem}

\subsection*{arSpectralDensity_nonneg}

\begin{theorem}[blueprint="infogeometry:maxent:arspectraldensity_nonneg"]
  \inputleannode{InfoGeometry.MaxEnt.arSpectralDensity_nonneg}
\end{theorem}

\subsection*{autocovarianceData}

\begin{theorem}[blueprint="infogeometry:maxent:autocovariancedata"]
  \inputleannode{InfoGeometry.MaxEnt.autocovarianceData}
\end{theorem}

\subsection*{constraintFamily}

\begin{theorem}[blueprint="infogeometry:maxent:constraintfamily"]
  \inputleannode{InfoGeometry.MaxEnt.constraintFamily}
\end{theorem}

\subsection*{crossEntropyToGibbs_eq}

\begin{theorem}[blueprint="infogeometry:maxent:crossentropytogibbs_eq"]
  \inputleannode{InfoGeometry.MaxEnt.crossEntropyToGibbs_eq}
\end{theorem}

\subsection*{crossEntropyToGibbs_eq_dualObjective}

\begin{theorem}[blueprint="infogeometry:maxent:crossentropytogibbs_eq_dualobjective"]
  \inputleannode{InfoGeometry.MaxEnt.crossEntropyToGibbs_eq_dualObjective}
\end{theorem}

\subsection*{dualObjective}

\begin{theorem}[blueprint="infogeometry:maxent:dualobjective"]
  \inputleannode{InfoGeometry.MaxEnt.dualObjective}
\end{theorem}

\subsection*{empiricalAutocovariance}

\begin{theorem}[blueprint="infogeometry:maxent:empiricalautocovariance"]
  \inputleannode{InfoGeometry.MaxEnt.empiricalAutocovariance}
\end{theorem}

\subsection*{empiricalFreq}

\begin{theorem}[blueprint="infogeometry:maxent:empiricalfreq"]
  \inputleannode{InfoGeometry.MaxEnt.empiricalFreq}
\end{theorem}

\subsection*{empiricalFreq_nonneg}

\begin{theorem}[blueprint="infogeometry:maxent:empiricalfreq_nonneg"]
  \inputleannode{InfoGeometry.MaxEnt.empiricalFreq_nonneg}
\end{theorem}

\subsection*{empiricalFreq_sum_one}

\begin{theorem}[blueprint="infogeometry:maxent:empiricalfreq_sum_one"]
  \inputleannode{InfoGeometry.MaxEnt.empiricalFreq_sum_one}
\end{theorem}

\subsection*{empiricalShannonEntropy}

\begin{theorem}[blueprint="infogeometry:maxent:empiricalshannonentropy"]
  \inputleannode{InfoGeometry.MaxEnt.empiricalShannonEntropy}
\end{theorem}

\subsection*{entropy}

\begin{theorem}[blueprint="infogeometry:maxent:entropy"]
  \inputleannode{InfoGeometry.MaxEnt.entropy}
\end{theorem}

\subsection*{entropyBand_of_mem_confidenceInterval}

\begin{theorem}[blueprint="infogeometry:maxent:entropyband_of_mem_confidenceinterval"]
  \inputleannode{InfoGeometry.MaxEnt.entropyBand_of_mem_confidenceInterval}
\end{theorem}

\subsection*{entropyConfidenceInterval}

\begin{theorem}[blueprint="infogeometry:maxent:entropyconfidenceinterval"]
  \inputleannode{InfoGeometry.MaxEnt.entropyConfidenceInterval}
\end{theorem}

\subsection*{entropyObjective}

\begin{theorem}[blueprint="infogeometry:maxent:entropyobjective"]
  \inputleannode{InfoGeometry.MaxEnt.entropyObjective}
\end{theorem}

\subsection*{entropy_gibbs_eq}

\begin{theorem}[blueprint="infogeometry:maxent:entropy_gibbs_eq"]
  \inputleannode{InfoGeometry.MaxEnt.entropy_gibbs_eq}
\end{theorem}

\subsection*{entropy_gibbs_eq_dualObjective}

\begin{theorem}[blueprint="infogeometry:maxent:entropy_gibbs_eq_dualobjective"]
  \inputleannode{InfoGeometry.MaxEnt.entropy_gibbs_eq_dualObjective}
\end{theorem}

\subsection*{entropy_le_dualObjective_on_constraint}

\begin{theorem}[blueprint="infogeometry:maxent:entropy_le_dualobjective_on_constraint"]
  \inputleannode{InfoGeometry.MaxEnt.entropy_le_dualObjective_on_constraint}
\end{theorem}

\subsection*{expectedAffineDimension}

\begin{theorem}[blueprint="infogeometry:maxent:expectedaffinedimension"]
  \inputleannode{InfoGeometry.MaxEnt.expectedAffineDimension}
\end{theorem}

\subsection*{gibbs}

\begin{theorem}[blueprint="infogeometry:maxent:gibbs"]
  \inputleannode{InfoGeometry.MaxEnt.gibbs}
\end{theorem}

\subsection*{gibbsDist}

\begin{theorem}[blueprint="infogeometry:maxent:gibbsdist"]
  \inputleannode{InfoGeometry.MaxEnt.gibbsDist}
\end{theorem}

\subsection*{gibbsExpectation}

\begin{theorem}[blueprint="infogeometry:maxent:gibbsexpectation"]
  \inputleannode{InfoGeometry.MaxEnt.gibbsExpectation}
\end{theorem}

\subsection*{gibbsExpectationWithPrior}

\begin{theorem}[blueprint="infogeometry:maxent:gibbsexpectationwithprior"]
  \inputleannode{InfoGeometry.MaxEnt.gibbsExpectationWithPrior}
\end{theorem}

\subsection*{gibbsMaxEntProblemOfExpectation}

\begin{theorem}[blueprint="infogeometry:maxent:gibbsmaxentproblemofexpectation"]
  \inputleannode{InfoGeometry.MaxEnt.gibbsMaxEntProblemOfExpectation}
\end{theorem}

\subsection*{gibbsWithPrior}

\begin{theorem}[blueprint="infogeometry:maxent:gibbswithprior"]
  \inputleannode{InfoGeometry.MaxEnt.gibbsWithPrior}
\end{theorem}

\subsection*{gibbsWithPriorMaxEntProblemOfExpectation}

\begin{theorem}[blueprint="infogeometry:maxent:gibbswithpriormaxentproblemofexpectation"]
  \inputleannode{InfoGeometry.MaxEnt.gibbsWithPriorMaxEntProblemOfExpectation}
\end{theorem}

\subsection*{gibbsWithPrior_eq_exp_sub_logPartition}

\begin{theorem}[blueprint="infogeometry:maxent:gibbswithprior_eq_exp_sub_logpartition"]
  \inputleannode{InfoGeometry.MaxEnt.gibbsWithPrior_eq_exp_sub_logPartition}
\end{theorem}

\subsection*{gibbsWithPrior_nonneg}

\begin{theorem}[blueprint="infogeometry:maxent:gibbswithprior_nonneg"]
  \inputleannode{InfoGeometry.MaxEnt.gibbsWithPrior_nonneg}
\end{theorem}

\subsection*{gibbsWithPrior_sum_one}

\begin{theorem}[blueprint="infogeometry:maxent:gibbswithprior_sum_one"]
  \inputleannode{InfoGeometry.MaxEnt.gibbsWithPrior_sum_one}
\end{theorem}

\subsection*{gibbs_attains_dualObjective}

\begin{theorem}[blueprint="infogeometry:maxent:gibbs_attains_dualobjective"]
  \inputleannode{InfoGeometry.MaxEnt.gibbs_attains_dualObjective}
\end{theorem}

\subsection*{gibbs_eq_exp_sub_logPartition}

\begin{theorem}[blueprint="infogeometry:maxent:gibbs_eq_exp_sub_logpartition"]
  \inputleannode{InfoGeometry.MaxEnt.gibbs_eq_exp_sub_logPartition}
\end{theorem}

\subsection*{gibbs_form_candidate_maximizes_entropy}

\begin{theorem}[blueprint="infogeometry:maxent:gibbs_form_candidate_maximizes_entropy"]
  \inputleannode{InfoGeometry.MaxEnt.gibbs_form_candidate_maximizes_entropy}
\end{theorem}

\subsection*{gibbs_maximizes_entropy_on_constraint}

\begin{theorem}[blueprint="infogeometry:maxent:gibbs_maximizes_entropy_on_constraint"]
  \inputleannode{InfoGeometry.MaxEnt.gibbs_maximizes_entropy_on_constraint}
\end{theorem}

\subsection*{gibbs_maximizes_shannon_under_moment}

\begin{theorem}[blueprint="infogeometry:maxent:gibbs_maximizes_shannon_under_moment"]
  \inputleannode{InfoGeometry.MaxEnt.gibbs_maximizes_shannon_under_moment}
\end{theorem}

\subsection*{gibbs_mem_maxEntConstraint_iff}

\begin{theorem}[blueprint="infogeometry:maxent:gibbs_mem_maxentconstraint_iff"]
  \inputleannode{InfoGeometry.MaxEnt.gibbs_mem_maxEntConstraint_iff}
\end{theorem}

\subsection*{gibbs_nonneg}

\begin{theorem}[blueprint="infogeometry:maxent:gibbs_nonneg"]
  \inputleannode{InfoGeometry.MaxEnt.gibbs_nonneg}
\end{theorem}

\subsection*{gibbs_pos}

\begin{theorem}[blueprint="infogeometry:maxent:gibbs_pos"]
  \inputleannode{InfoGeometry.MaxEnt.gibbs_pos}
\end{theorem}

\subsection*{gibbs_sum_one}

\begin{theorem}[blueprint="infogeometry:maxent:gibbs_sum_one"]
  \inputleannode{InfoGeometry.MaxEnt.gibbs_sum_one}
\end{theorem}

\subsection*{goodFraction}

\begin{theorem}[blueprint="infogeometry:maxent:goodfraction"]
  \inputleannode{InfoGeometry.MaxEnt.goodFraction}
\end{theorem}

\subsection*{goodFraction_le_one}

\begin{theorem}[blueprint="infogeometry:maxent:goodfraction_le_one"]
  \inputleannode{InfoGeometry.MaxEnt.goodFraction_le_one}
\end{theorem}

\subsection*{goodFraction_nonneg}

\begin{theorem}[blueprint="infogeometry:maxent:goodfraction_nonneg"]
  \inputleannode{InfoGeometry.MaxEnt.goodFraction_nonneg}
\end{theorem}

\subsection*{goodProfiles}

\begin{theorem}[blueprint="infogeometry:maxent:goodprofiles"]
  \inputleannode{InfoGeometry.MaxEnt.goodProfiles}
\end{theorem}

\subsection*{logMultiplicity}

\begin{theorem}[blueprint="infogeometry:maxent:logmultiplicity"]
  \inputleannode{InfoGeometry.MaxEnt.logMultiplicity}
\end{theorem}

\subsection*{logMultiplicity_nonneg}

\begin{theorem}[blueprint="infogeometry:maxent:logmultiplicity_nonneg"]
  \inputleannode{InfoGeometry.MaxEnt.logMultiplicity_nonneg}
\end{theorem}

\subsection*{logPartition}

\begin{theorem}[blueprint="infogeometry:maxent:logpartition"]
  \inputleannode{InfoGeometry.MaxEnt.logPartition}
\end{theorem}

\subsection*{logPartitionWithPrior}

\begin{theorem}[blueprint="infogeometry:maxent:logpartitionwithprior"]
  \inputleannode{InfoGeometry.MaxEnt.logPartitionWithPrior}
\end{theorem}

\subsection*{maxEnt_stationary}

\begin{theorem}[blueprint="infogeometry:maxent:maxent_stationary"]
  \inputleannode{InfoGeometry.MaxEnt.maxEnt_stationary}
\end{theorem}

\subsection*{max_ent_lagrange_multiplier}

\begin{theorem}[blueprint="infogeometry:maxent:max_ent_lagrange_multiplier"]
  \inputleannode{InfoGeometry.MaxEnt.max_ent_lagrange_multiplier}
\end{theorem}

\subsection*{max_ent_lagrange_multiplier_gibbs}

\begin{theorem}[blueprint="infogeometry:maxent:max_ent_lagrange_multiplier_gibbs"]
  \inputleannode{InfoGeometry.MaxEnt.max_ent_lagrange_multiplier_gibbs}
\end{theorem}

\subsection*{mem_confidenceInterval_of_entropyBand}

\begin{theorem}[blueprint="infogeometry:maxent:mem_confidenceinterval_of_entropyband"]
  \inputleannode{InfoGeometry.MaxEnt.mem_confidenceInterval_of_entropyBand}
\end{theorem}

\subsection*{mem_feasibleClass_iff}

\begin{theorem}[blueprint="infogeometry:maxent:mem_feasibleclass_iff"]
  \inputleannode{InfoGeometry.MaxEnt.mem_feasibleClass_iff}
\end{theorem}

\subsection*{mem_simplex_iff}

\begin{theorem}[blueprint="infogeometry:maxent:mem_simplex_iff"]
  \inputleannode{InfoGeometry.MaxEnt.mem_simplex_iff}
\end{theorem}

\subsection*{momentResidual}

\begin{theorem}[blueprint="infogeometry:maxent:momentresidual"]
  \inputleannode{InfoGeometry.MaxEnt.momentResidual}
\end{theorem}

\subsection*{momentResidual_eq_zero_iff}

\begin{theorem}[blueprint="infogeometry:maxent:momentresidual_eq_zero_iff"]
  \inputleannode{InfoGeometry.MaxEnt.momentResidual_eq_zero_iff}
\end{theorem}

\subsection*{multiplicity}

\begin{theorem}[blueprint="infogeometry:maxent:multiplicity"]
  \inputleannode{InfoGeometry.MaxEnt.multiplicity}
\end{theorem}

\subsection*{multiplicity_one_le}

\begin{theorem}[blueprint="infogeometry:maxent:multiplicity_one_le"]
  \inputleannode{InfoGeometry.MaxEnt.multiplicity_one_le}
\end{theorem}

\subsection*{multiplicity_pos}

\begin{theorem}[blueprint="infogeometry:maxent:multiplicity_pos"]
  \inputleannode{InfoGeometry.MaxEnt.multiplicity_pos}
\end{theorem}

\subsection*{multiplicity_spec}

\begin{theorem}[blueprint="infogeometry:maxent:multiplicity_spec"]
  \inputleannode{InfoGeometry.MaxEnt.multiplicity_spec}
\end{theorem}

\subsection*{normalizedLogMultiplicity}

\begin{theorem}[blueprint="infogeometry:maxent:normalizedlogmultiplicity"]
  \inputleannode{InfoGeometry.MaxEnt.normalizedLogMultiplicity}
\end{theorem}

\subsection*{normalized_of_mem_feasibleClass}

\begin{theorem}[blueprint="infogeometry:maxent:normalized_of_mem_feasibleclass"]
  \inputleannode{InfoGeometry.MaxEnt.normalized_of_mem_feasibleClass}
\end{theorem}

\subsection*{normalized_of_mem_simplex}

\begin{theorem}[blueprint="infogeometry:maxent:normalized_of_mem_simplex"]
  \inputleannode{InfoGeometry.MaxEnt.normalized_of_mem_simplex}
\end{theorem}

\subsection*{partition}

\begin{theorem}[blueprint="infogeometry:maxent:partition"]
  \inputleannode{InfoGeometry.MaxEnt.partition}
\end{theorem}

\subsection*{partitionWithPrior}

\begin{theorem}[blueprint="infogeometry:maxent:partitionwithprior"]
  \inputleannode{InfoGeometry.MaxEnt.partitionWithPrior}
\end{theorem}

\subsection*{partitionWithPrior_ne_zero}

\begin{theorem}[blueprint="infogeometry:maxent:partitionwithprior_ne_zero"]
  \inputleannode{InfoGeometry.MaxEnt.partitionWithPrior_ne_zero}
\end{theorem}

\subsection*{partitionWithPrior_pos}

\begin{theorem}[blueprint="infogeometry:maxent:partitionwithprior_pos"]
  \inputleannode{InfoGeometry.MaxEnt.partitionWithPrior_pos}
\end{theorem}

\subsection*{partition_ne_zero}

\begin{theorem}[blueprint="infogeometry:maxent:partition_ne_zero"]
  \inputleannode{InfoGeometry.MaxEnt.partition_ne_zero}
\end{theorem}

\subsection*{partition_pos}

\begin{theorem}[blueprint="infogeometry:maxent:partition_pos"]
  \inputleannode{InfoGeometry.MaxEnt.partition_pos}
\end{theorem}

\subsection*{probDistOfSimplex}

\begin{theorem}[blueprint="infogeometry:maxent:probdistofsimplex"]
  \inputleannode{InfoGeometry.MaxEnt.probDistOfSimplex}
\end{theorem}

\subsection*{sampleMean}

\begin{theorem}[blueprint="infogeometry:maxent:samplemean"]
  \inputleannode{InfoGeometry.MaxEnt.sampleMean}
\end{theorem}

\subsection*{thermodynamicPreference}

\begin{theorem}[blueprint="infogeometry:maxent:thermodynamicpreference"]
  \inputleannode{InfoGeometry.MaxEnt.thermodynamicPreference}
\end{theorem}

\subsection*{thermodynamicPreference_pos}

\begin{theorem}[blueprint="infogeometry:maxent:thermodynamicpreference_pos"]
  \inputleannode{InfoGeometry.MaxEnt.thermodynamicPreference_pos}
\end{theorem}

\subsection*{totalCount}

\begin{theorem}[blueprint="infogeometry:maxent:totalcount"]
  \inputleannode{InfoGeometry.MaxEnt.totalCount}
\end{theorem}

\subsection*{xlogx}

\begin{theorem}[blueprint="infogeometry:maxent:xlogx"]
  \inputleannode{InfoGeometry.MaxEnt.xlogx}
\end{theorem}

\subsection*{xlogx_eq_mul_log}

\begin{theorem}[blueprint="infogeometry:maxent:xlogx_eq_mul_log"]
  \inputleannode{InfoGeometry.MaxEnt.xlogx_eq_mul_log}
\end{theorem}

\subsection*{xlogx_zero}

\begin{theorem}[blueprint="infogeometry:maxent:xlogx_zero"]
  \inputleannode{InfoGeometry.MaxEnt.xlogx_zero}
\end{theorem}

\section{InfoGeometry.MaxEnt.ACProbMeasure}

\subsection*{ac}

\begin{theorem}[blueprint="infogeometry:maxent:acprobmeasure:ac"]
  \inputleannode{InfoGeometry.MaxEnt.ACProbMeasure.ac}
\end{theorem}

\subsection*{casesOn}

\begin{theorem}[blueprint="infogeometry:maxent:acprobmeasure:caseson"]
  \inputleannode{InfoGeometry.MaxEnt.ACProbMeasure.casesOn}
\end{theorem}

\subsection*{coe_eq_of_measure_eq}

\begin{theorem}[blueprint="infogeometry:maxent:acprobmeasure:coe_eq_of_measure_eq"]
  \inputleannode{InfoGeometry.MaxEnt.ACProbMeasure.coe_eq_of_measure_eq}
\end{theorem}

\subsection*{ctorIdx}

\begin{theorem}[blueprint="infogeometry:maxent:acprobmeasure:ctoridx"]
  \inputleannode{InfoGeometry.MaxEnt.ACProbMeasure.ctorIdx}
\end{theorem}

\subsection*{instCoeTCMeasure}

\begin{theorem}[blueprint="infogeometry:maxent:acprobmeasure:instcoetcmeasure"]
  \inputleannode{InfoGeometry.MaxEnt.ACProbMeasure.instCoeTCMeasure}
\end{theorem}

\subsection*{instIsProbabilityMeasureμ}

\begin{theorem}[blueprint="infogeometry:maxent:acprobmeasure:instisprobabilitymeasureμ"]
  \inputleannode{InfoGeometry.MaxEnt.ACProbMeasure.instIsProbabilityMeasureμ}
\end{theorem}

\subsection*{mk}

\begin{theorem}[blueprint="infogeometry:maxent:acprobmeasure:mk"]
  \inputleannode{InfoGeometry.MaxEnt.ACProbMeasure.mk}
\end{theorem}

\subsection*{noConfusion}

\begin{theorem}[blueprint="infogeometry:maxent:acprobmeasure:noconfusion"]
  \inputleannode{InfoGeometry.MaxEnt.ACProbMeasure.noConfusion}
\end{theorem}

\subsection*{noConfusionType}

\begin{theorem}[blueprint="infogeometry:maxent:acprobmeasure:noconfusiontype"]
  \inputleannode{InfoGeometry.MaxEnt.ACProbMeasure.noConfusionType}
\end{theorem}

\subsection*{prob}

\begin{theorem}[blueprint="infogeometry:maxent:acprobmeasure:prob"]
  \inputleannode{InfoGeometry.MaxEnt.ACProbMeasure.prob}
\end{theorem}

\subsection*{rec}

\begin{theorem}[blueprint="infogeometry:maxent:acprobmeasure:rec"]
  \inputleannode{InfoGeometry.MaxEnt.ACProbMeasure.rec}
\end{theorem}

\subsection*{recOn}

\begin{theorem}[blueprint="infogeometry:maxent:acprobmeasure:recon"]
  \inputleannode{InfoGeometry.MaxEnt.ACProbMeasure.recOn}
\end{theorem}

\subsection*{μ}

\begin{theorem}[blueprint="infogeometry:maxent:acprobmeasure:μ"]
  \inputleannode{InfoGeometry.MaxEnt.ACProbMeasure.μ}
\end{theorem}

\section{InfoGeometry.MaxEnt.AsymptoticEquipartitionWitness}

\subsection*{casesOn}

\begin{theorem}[blueprint="infogeometry:maxent:asymptoticequipartitionwitness:caseson"]
  \inputleannode{InfoGeometry.MaxEnt.AsymptoticEquipartitionWitness.casesOn}
\end{theorem}

\subsection*{countsSeq}

\begin{theorem}[blueprint="infogeometry:maxent:asymptoticequipartitionwitness:countsseq"]
  \inputleannode{InfoGeometry.MaxEnt.AsymptoticEquipartitionWitness.countsSeq}
\end{theorem}

\subsection*{ctorIdx}

\begin{theorem}[blueprint="infogeometry:maxent:asymptoticequipartitionwitness:ctoridx"]
  \inputleannode{InfoGeometry.MaxEnt.AsymptoticEquipartitionWitness.ctorIdx}
\end{theorem}

\subsection*{mk}

\begin{theorem}[blueprint="infogeometry:maxent:asymptoticequipartitionwitness:mk"]
  \inputleannode{InfoGeometry.MaxEnt.AsymptoticEquipartitionWitness.mk}
\end{theorem}

\subsection*{noConfusion}

\begin{theorem}[blueprint="infogeometry:maxent:asymptoticequipartitionwitness:noconfusion"]
  \inputleannode{InfoGeometry.MaxEnt.AsymptoticEquipartitionWitness.noConfusion}
\end{theorem}

\subsection*{noConfusionType}

\begin{theorem}[blueprint="infogeometry:maxent:asymptoticequipartitionwitness:noconfusiontype"]
  \inputleannode{InfoGeometry.MaxEnt.AsymptoticEquipartitionWitness.noConfusionType}
\end{theorem}

\subsection*{rec}

\begin{theorem}[blueprint="infogeometry:maxent:asymptoticequipartitionwitness:rec"]
  \inputleannode{InfoGeometry.MaxEnt.AsymptoticEquipartitionWitness.rec}
\end{theorem}

\subsection*{recOn}

\begin{theorem}[blueprint="infogeometry:maxent:asymptoticequipartitionwitness:recon"]
  \inputleannode{InfoGeometry.MaxEnt.AsymptoticEquipartitionWitness.recOn}
\end{theorem}

\subsection*{tendsToEntropy}

\begin{theorem}[blueprint="infogeometry:maxent:asymptoticequipartitionwitness:tendstoentropy"]
  \inputleannode{InfoGeometry.MaxEnt.AsymptoticEquipartitionWitness.tendsToEntropy}
\end{theorem}

\subsection*{total_pos}

\begin{theorem}[blueprint="infogeometry:maxent:asymptoticequipartitionwitness:total_pos"]
  \inputleannode{InfoGeometry.MaxEnt.AsymptoticEquipartitionWitness.total_pos}
\end{theorem}

\section{InfoGeometry.MaxEnt.BurgModel}

\subsection*{casesOn}

\begin{theorem}[blueprint="infogeometry:maxent:burgmodel:caseson"]
  \inputleannode{InfoGeometry.MaxEnt.BurgModel.casesOn}
\end{theorem}

\subsection*{coeff}

\begin{theorem}[blueprint="infogeometry:maxent:burgmodel:coeff"]
  \inputleannode{InfoGeometry.MaxEnt.BurgModel.coeff}
\end{theorem}

\subsection*{ctorIdx}

\begin{theorem}[blueprint="infogeometry:maxent:burgmodel:ctoridx"]
  \inputleannode{InfoGeometry.MaxEnt.BurgModel.ctorIdx}
\end{theorem}

\subsection*{mk}

\begin{theorem}[blueprint="infogeometry:maxent:burgmodel:mk"]
  \inputleannode{InfoGeometry.MaxEnt.BurgModel.mk}
\end{theorem}

\subsection*{noConfusion}

\begin{theorem}[blueprint="infogeometry:maxent:burgmodel:noconfusion"]
  \inputleannode{InfoGeometry.MaxEnt.BurgModel.noConfusion}
\end{theorem}

\subsection*{noConfusionType}

\begin{theorem}[blueprint="infogeometry:maxent:burgmodel:noconfusiontype"]
  \inputleannode{InfoGeometry.MaxEnt.BurgModel.noConfusionType}
\end{theorem}

\subsection*{noiseVar}

\begin{theorem}[blueprint="infogeometry:maxent:burgmodel:noisevar"]
  \inputleannode{InfoGeometry.MaxEnt.BurgModel.noiseVar}
\end{theorem}

\subsection*{noiseVar_nonneg}

\begin{theorem}[blueprint="infogeometry:maxent:burgmodel:noisevar_nonneg"]
  \inputleannode{InfoGeometry.MaxEnt.BurgModel.noiseVar_nonneg}
\end{theorem}

\subsection*{rec}

\begin{theorem}[blueprint="infogeometry:maxent:burgmodel:rec"]
  \inputleannode{InfoGeometry.MaxEnt.BurgModel.rec}
\end{theorem}

\subsection*{recOn}

\begin{theorem}[blueprint="infogeometry:maxent:burgmodel:recon"]
  \inputleannode{InfoGeometry.MaxEnt.BurgModel.recOn}
\end{theorem}

\subsection*{spectrum}

\begin{theorem}[blueprint="infogeometry:maxent:burgmodel:spectrum"]
  \inputleannode{InfoGeometry.MaxEnt.BurgModel.spectrum}
\end{theorem}

\subsection*{spectrum_nonneg}

\begin{theorem}[blueprint="infogeometry:maxent:burgmodel:spectrum_nonneg"]
  \inputleannode{InfoGeometry.MaxEnt.BurgModel.spectrum_nonneg}
\end{theorem}

\section{InfoGeometry.MaxEnt.CanonicalProblem}

\subsection*{A}

\begin{theorem}[blueprint="infogeometry:maxent:canonicalproblem:a"]
  \inputleannode{InfoGeometry.MaxEnt.CanonicalProblem.A}
\end{theorem}

\subsection*{b}

\begin{theorem}[blueprint="infogeometry:maxent:canonicalproblem:b"]
  \inputleannode{InfoGeometry.MaxEnt.CanonicalProblem.b}
\end{theorem}

\subsection*{casesOn}

\begin{theorem}[blueprint="infogeometry:maxent:canonicalproblem:caseson"]
  \inputleannode{InfoGeometry.MaxEnt.CanonicalProblem.casesOn}
\end{theorem}

\subsection*{ctorIdx}

\begin{theorem}[blueprint="infogeometry:maxent:canonicalproblem:ctoridx"]
  \inputleannode{InfoGeometry.MaxEnt.CanonicalProblem.ctorIdx}
\end{theorem}

\subsection*{expected}

\begin{theorem}[blueprint="infogeometry:maxent:canonicalproblem:expected"]
  \inputleannode{InfoGeometry.MaxEnt.CanonicalProblem.expected}
\end{theorem}

\subsection*{feasible}

\begin{theorem}[blueprint="infogeometry:maxent:canonicalproblem:feasible"]
  \inputleannode{InfoGeometry.MaxEnt.CanonicalProblem.feasible}
\end{theorem}

\subsection*{gibbsCandidate}

\begin{theorem}[blueprint="infogeometry:maxent:canonicalproblem:gibbscandidate"]
  \inputleannode{InfoGeometry.MaxEnt.CanonicalProblem.gibbsCandidate}
\end{theorem}

\subsection*{gibbsCandidate_maximizes_entropy_on_feasible}

\begin{theorem}[blueprint="infogeometry:maxent:canonicalproblem:gibbscandidate_maximizes_entropy_on_feasible"]
  \inputleannode{InfoGeometry.MaxEnt.CanonicalProblem.gibbsCandidate_maximizes_entropy_on_feasible}
\end{theorem}

\subsection*{gibbsCandidate_mem_feasible}

\begin{theorem}[blueprint="infogeometry:maxent:canonicalproblem:gibbscandidate_mem_feasible"]
  \inputleannode{InfoGeometry.MaxEnt.CanonicalProblem.gibbsCandidate_mem_feasible}
\end{theorem}

\subsection*{mk}

\begin{theorem}[blueprint="infogeometry:maxent:canonicalproblem:mk"]
  \inputleannode{InfoGeometry.MaxEnt.CanonicalProblem.mk}
\end{theorem}

\subsection*{noConfusion}

\begin{theorem}[blueprint="infogeometry:maxent:canonicalproblem:noconfusion"]
  \inputleannode{InfoGeometry.MaxEnt.CanonicalProblem.noConfusion}
\end{theorem}

\subsection*{noConfusionType}

\begin{theorem}[blueprint="infogeometry:maxent:canonicalproblem:noconfusiontype"]
  \inputleannode{InfoGeometry.MaxEnt.CanonicalProblem.noConfusionType}
\end{theorem}

\subsection*{obs}

\begin{theorem}[blueprint="infogeometry:maxent:canonicalproblem:obs"]
  \inputleannode{InfoGeometry.MaxEnt.CanonicalProblem.obs}
\end{theorem}

\subsection*{rec}

\begin{theorem}[blueprint="infogeometry:maxent:canonicalproblem:rec"]
  \inputleannode{InfoGeometry.MaxEnt.CanonicalProblem.rec}
\end{theorem}

\subsection*{recOn}

\begin{theorem}[blueprint="infogeometry:maxent:canonicalproblem:recon"]
  \inputleannode{InfoGeometry.MaxEnt.CanonicalProblem.recOn}
\end{theorem}

\section{InfoGeometry.MaxEnt.CanonicalProblem.gibbsCandidate}

\subsection*{eq_1}

\begin{theorem}[blueprint="infogeometry:maxent:canonicalproblem:gibbscandidate:eq_1"]
  \inputleannode{InfoGeometry.MaxEnt.CanonicalProblem.gibbsCandidate.eq_1}
\end{theorem}

\section{InfoGeometry.MaxEnt.Core}

\noindent\textbf{Buildable}: yes

\noindent\emph{No exported declarations in the current declaration graph.}

\section{InfoGeometry.MaxEnt.DualBridge}

\noindent\textbf{Buildable}: yes

\noindent\emph{No exported declarations in the current declaration graph.}

\section{InfoGeometry.MaxEnt.EntropyConcentrationCertificate}

\subsection*{casesOn}

\begin{theorem}[blueprint="infogeometry:maxent:entropyconcentrationcertificate:caseson"]
  \inputleannode{InfoGeometry.MaxEnt.EntropyConcentrationCertificate.casesOn}
\end{theorem}

\subsection*{eps_nonneg}

\begin{theorem}[blueprint="infogeometry:maxent:entropyconcentrationcertificate:eps_nonneg"]
  \inputleannode{InfoGeometry.MaxEnt.EntropyConcentrationCertificate.eps_nonneg}
\end{theorem}

\subsection*{eta_bounds}

\begin{theorem}[blueprint="infogeometry:maxent:entropyconcentrationcertificate:eta_bounds"]
  \inputleannode{InfoGeometry.MaxEnt.EntropyConcentrationCertificate.eta_bounds}
\end{theorem}

\subsection*{lower_fraction_bound}

\begin{theorem}[blueprint="infogeometry:maxent:entropyconcentrationcertificate:lower_fraction_bound"]
  \inputleannode{InfoGeometry.MaxEnt.EntropyConcentrationCertificate.lower_fraction_bound}
\end{theorem}

\subsection*{mk}

\begin{theorem}[blueprint="infogeometry:maxent:entropyconcentrationcertificate:mk"]
  \inputleannode{InfoGeometry.MaxEnt.EntropyConcentrationCertificate.mk}
\end{theorem}

\subsection*{nonempty}

\begin{theorem}[blueprint="infogeometry:maxent:entropyconcentrationcertificate:nonempty"]
  \inputleannode{InfoGeometry.MaxEnt.EntropyConcentrationCertificate.nonempty}
\end{theorem}

\subsection*{rec}

\begin{theorem}[blueprint="infogeometry:maxent:entropyconcentrationcertificate:rec"]
  \inputleannode{InfoGeometry.MaxEnt.EntropyConcentrationCertificate.rec}
\end{theorem}

\subsection*{recOn}

\begin{theorem}[blueprint="infogeometry:maxent:entropyconcentrationcertificate:recon"]
  \inputleannode{InfoGeometry.MaxEnt.EntropyConcentrationCertificate.recOn}
\end{theorem}

\section{InfoGeometry.MaxEnt.FenchelDualAssumptions}

\subsection*{casesOn}

\begin{theorem}[blueprint="infogeometry:maxent:fencheldualassumptions:caseson"]
  \inputleannode{InfoGeometry.MaxEnt.FenchelDualAssumptions.casesOn}
\end{theorem}

\subsection*{ctorIdx}

\begin{theorem}[blueprint="infogeometry:maxent:fencheldualassumptions:ctoridx"]
  \inputleannode{InfoGeometry.MaxEnt.FenchelDualAssumptions.ctorIdx}
\end{theorem}

\subsection*{dualObjective}

\begin{theorem}[blueprint="infogeometry:maxent:fencheldualassumptions:dualobjective"]
  \inputleannode{InfoGeometry.MaxEnt.FenchelDualAssumptions.dualObjective}
\end{theorem}

\subsection*{equalityImpliesExponential}

\begin{theorem}[blueprint="infogeometry:maxent:fencheldualassumptions:equalityimpliesexponential"]
  \inputleannode{InfoGeometry.MaxEnt.FenchelDualAssumptions.equalityImpliesExponential}
\end{theorem}

\subsection*{index_exact}

\begin{theorem}[blueprint="infogeometry:maxent:fencheldualassumptions:index_exact"]
  \inputleannode{InfoGeometry.MaxEnt.FenchelDualAssumptions.index_exact}
\end{theorem}

\subsection*{mk}

\begin{theorem}[blueprint="infogeometry:maxent:fencheldualassumptions:mk"]
  \inputleannode{InfoGeometry.MaxEnt.FenchelDualAssumptions.mk}
\end{theorem}

\subsection*{noConfusion}

\begin{theorem}[blueprint="infogeometry:maxent:fencheldualassumptions:noconfusion"]
  \inputleannode{InfoGeometry.MaxEnt.FenchelDualAssumptions.noConfusion}
\end{theorem}

\subsection*{noConfusionType}

\begin{theorem}[blueprint="infogeometry:maxent:fencheldualassumptions:noconfusiontype"]
  \inputleannode{InfoGeometry.MaxEnt.FenchelDualAssumptions.noConfusionType}
\end{theorem}

\subsection*{rec}

\begin{theorem}[blueprint="infogeometry:maxent:fencheldualassumptions:rec"]
  \inputleannode{InfoGeometry.MaxEnt.FenchelDualAssumptions.rec}
\end{theorem}

\subsection*{recOn}

\begin{theorem}[blueprint="infogeometry:maxent:fencheldualassumptions:recon"]
  \inputleannode{InfoGeometry.MaxEnt.FenchelDualAssumptions.recOn}
\end{theorem}

\subsection*{strongDualityAtOpt}

\begin{theorem}[blueprint="infogeometry:maxent:fencheldualassumptions:strongdualityatopt"]
  \inputleannode{InfoGeometry.MaxEnt.FenchelDualAssumptions.strongDualityAtOpt}
\end{theorem}

\subsection*{weakDuality}

\begin{theorem}[blueprint="infogeometry:maxent:fencheldualassumptions:weakduality"]
  \inputleannode{InfoGeometry.MaxEnt.FenchelDualAssumptions.weakDuality}
\end{theorem}

\section{InfoGeometry.MaxEnt.Finite}

\noindent\textbf{Buildable}: yes

\subsection*{FiniteJaynesProblem}

\begin{theorem}[blueprint="infogeometry:maxent:finite:finitejaynesproblem"]
  \inputleannode{InfoGeometry.MaxEnt.Finite.FiniteJaynesProblem}
\end{theorem}

\section{InfoGeometry.MaxEnt.Finite.FiniteJaynesProblem}

\subsection*{FullSupportPrior}

\begin{theorem}[blueprint="infogeometry:maxent:finite:finitejaynesproblem:fullsupportprior"]
  \inputleannode{InfoGeometry.MaxEnt.Finite.FiniteJaynesProblem.FullSupportPrior}
\end{theorem}

\subsection*{SatisfiesTargetMoments}

\begin{theorem}[blueprint="infogeometry:maxent:finite:finitejaynesproblem:satisfiestargetmoments"]
  \inputleannode{InfoGeometry.MaxEnt.Finite.FiniteJaynesProblem.SatisfiesTargetMoments}
\end{theorem}

\subsection*{casesOn}

\begin{theorem}[blueprint="infogeometry:maxent:finite:finitejaynesproblem:caseson"]
  \inputleannode{InfoGeometry.MaxEnt.Finite.FiniteJaynesProblem.casesOn}
\end{theorem}

\subsection*{ctorIdx}

\begin{theorem}[blueprint="infogeometry:maxent:finite:finitejaynesproblem:ctoridx"]
  \inputleannode{InfoGeometry.MaxEnt.Finite.FiniteJaynesProblem.ctorIdx}
\end{theorem}

\subsection*{energy}

\begin{theorem}[blueprint="infogeometry:maxent:finite:finitejaynesproblem:energy"]
  \inputleannode{InfoGeometry.MaxEnt.Finite.FiniteJaynesProblem.energy}
\end{theorem}

\subsection*{energy_zero}

\begin{theorem}[blueprint="infogeometry:maxent:finite:finitejaynesproblem:energy_zero"]
  \inputleannode{InfoGeometry.MaxEnt.Finite.FiniteJaynesProblem.energy_zero}
\end{theorem}

\subsection*{feature}

\begin{theorem}[blueprint="infogeometry:maxent:finite:finitejaynesproblem:feature"]
  \inputleannode{InfoGeometry.MaxEnt.Finite.FiniteJaynesProblem.feature}
\end{theorem}

\subsection*{gibbsDist}

\begin{theorem}[blueprint="infogeometry:maxent:finite:finitejaynesproblem:gibbsdist"]
  \inputleannode{InfoGeometry.MaxEnt.Finite.FiniteJaynesProblem.gibbsDist}
\end{theorem}

\subsection*{gibbsProb}

\begin{theorem}[blueprint="infogeometry:maxent:finite:finitejaynesproblem:gibbsprob"]
  \inputleannode{InfoGeometry.MaxEnt.Finite.FiniteJaynesProblem.gibbsProb}
\end{theorem}

\subsection*{gibbsProb_eq_prior_mul_exp_tilt}

\begin{theorem}[blueprint="infogeometry:maxent:finite:finitejaynesproblem:gibbsprob_eq_prior_mul_exp_tilt"]
  \inputleannode{InfoGeometry.MaxEnt.Finite.FiniteJaynesProblem.gibbsProb_eq_prior_mul_exp_tilt}
\end{theorem}

\subsection*{gibbsProb_nonneg}

\begin{theorem}[blueprint="infogeometry:maxent:finite:finitejaynesproblem:gibbsprob_nonneg"]
  \inputleannode{InfoGeometry.MaxEnt.Finite.FiniteJaynesProblem.gibbsProb_nonneg}
\end{theorem}

\subsection*{gibbsWeight}

\begin{theorem}[blueprint="infogeometry:maxent:finite:finitejaynesproblem:gibbsweight"]
  \inputleannode{InfoGeometry.MaxEnt.Finite.FiniteJaynesProblem.gibbsWeight}
\end{theorem}

\subsection*{gibbsWeight_nonneg}

\begin{theorem}[blueprint="infogeometry:maxent:finite:finitejaynesproblem:gibbsweight_nonneg"]
  \inputleannode{InfoGeometry.MaxEnt.Finite.FiniteJaynesProblem.gibbsWeight_nonneg}
\end{theorem}

\subsection*{index}

\begin{theorem}[blueprint="infogeometry:maxent:finite:finitejaynesproblem:index"]
  \inputleannode{InfoGeometry.MaxEnt.Finite.FiniteJaynesProblem.index}
\end{theorem}

\subsection*{logPartition}

\begin{theorem}[blueprint="infogeometry:maxent:finite:finitejaynesproblem:logpartition"]
  \inputleannode{InfoGeometry.MaxEnt.Finite.FiniteJaynesProblem.logPartition}
\end{theorem}

\subsection*{logPartition_zero}

\begin{theorem}[blueprint="infogeometry:maxent:finite:finitejaynesproblem:logpartition_zero"]
  \inputleannode{InfoGeometry.MaxEnt.Finite.FiniteJaynesProblem.logPartition_zero}
\end{theorem}

\subsection*{log_gibbsRatio_eq_energy_sub_logPartition}

\begin{theorem}[blueprint="infogeometry:maxent:finite:finitejaynesproblem:log_gibbsratio_eq_energy_sub_logpartition"]
  \inputleannode{InfoGeometry.MaxEnt.Finite.FiniteJaynesProblem.log_gibbsRatio_eq_energy_sub_logPartition}
\end{theorem}

\subsection*{mk}

\begin{theorem}[blueprint="infogeometry:maxent:finite:finitejaynesproblem:mk"]
  \inputleannode{InfoGeometry.MaxEnt.Finite.FiniteJaynesProblem.mk}
\end{theorem}

\subsection*{moment}

\begin{theorem}[blueprint="infogeometry:maxent:finite:finitejaynesproblem:moment"]
  \inputleannode{InfoGeometry.MaxEnt.Finite.FiniteJaynesProblem.moment}
\end{theorem}

\subsection*{noConfusion}

\begin{theorem}[blueprint="infogeometry:maxent:finite:finitejaynesproblem:noconfusion"]
  \inputleannode{InfoGeometry.MaxEnt.Finite.FiniteJaynesProblem.noConfusion}
\end{theorem}

\subsection*{noConfusionType}

\begin{theorem}[blueprint="infogeometry:maxent:finite:finitejaynesproblem:noconfusiontype"]
  \inputleannode{InfoGeometry.MaxEnt.Finite.FiniteJaynesProblem.noConfusionType}
\end{theorem}

\subsection*{ofLogLikelihood}

\begin{theorem}[blueprint="infogeometry:maxent:finite:finitejaynesproblem:ofloglikelihood"]
  \inputleannode{InfoGeometry.MaxEnt.Finite.FiniteJaynesProblem.ofLogLikelihood}
\end{theorem}

\subsection*{ofLogLikelihood_energy}

\begin{theorem}[blueprint="infogeometry:maxent:finite:finitejaynesproblem:ofloglikelihood_energy"]
  \inputleannode{InfoGeometry.MaxEnt.Finite.FiniteJaynesProblem.ofLogLikelihood_energy}
\end{theorem}

\subsection*{partition}

\begin{theorem}[blueprint="infogeometry:maxent:finite:finitejaynesproblem:partition"]
  \inputleannode{InfoGeometry.MaxEnt.Finite.FiniteJaynesProblem.partition}
\end{theorem}

\subsection*{partition_ne_zero_of_fullSupport}

\begin{theorem}[blueprint="infogeometry:maxent:finite:finitejaynesproblem:partition_ne_zero_of_fullsupport"]
  \inputleannode{InfoGeometry.MaxEnt.Finite.FiniteJaynesProblem.partition_ne_zero_of_fullSupport}
\end{theorem}

\subsection*{partition_nonneg}

\begin{theorem}[blueprint="infogeometry:maxent:finite:finitejaynesproblem:partition_nonneg"]
  \inputleannode{InfoGeometry.MaxEnt.Finite.FiniteJaynesProblem.partition_nonneg}
\end{theorem}

\subsection*{partition_pos_of_fullSupport}

\begin{theorem}[blueprint="infogeometry:maxent:finite:finitejaynesproblem:partition_pos_of_fullsupport"]
  \inputleannode{InfoGeometry.MaxEnt.Finite.FiniteJaynesProblem.partition_pos_of_fullSupport}
\end{theorem}

\subsection*{partition_zero}

\begin{theorem}[blueprint="infogeometry:maxent:finite:finitejaynesproblem:partition_zero"]
  \inputleannode{InfoGeometry.MaxEnt.Finite.FiniteJaynesProblem.partition_zero}
\end{theorem}

\subsection*{prior}

\begin{theorem}[blueprint="infogeometry:maxent:finite:finitejaynesproblem:prior"]
  \inputleannode{InfoGeometry.MaxEnt.Finite.FiniteJaynesProblem.prior}
\end{theorem}

\subsection*{rec}

\begin{theorem}[blueprint="infogeometry:maxent:finite:finitejaynesproblem:rec"]
  \inputleannode{InfoGeometry.MaxEnt.Finite.FiniteJaynesProblem.rec}
\end{theorem}

\subsection*{recOn}

\begin{theorem}[blueprint="infogeometry:maxent:finite:finitejaynesproblem:recon"]
  \inputleannode{InfoGeometry.MaxEnt.Finite.FiniteJaynesProblem.recOn}
\end{theorem}

\subsection*{target}

\begin{theorem}[blueprint="infogeometry:maxent:finite:finitejaynesproblem:target"]
  \inputleannode{InfoGeometry.MaxEnt.Finite.FiniteJaynesProblem.target}
\end{theorem}

\section{InfoGeometry.MaxEnt.Finite.FiniteJaynesProblem.energy}

\subsection*{eq_1}

\begin{theorem}[blueprint="infogeometry:maxent:finite:finitejaynesproblem:energy:eq_1"]
  \inputleannode{InfoGeometry.MaxEnt.Finite.FiniteJaynesProblem.energy.eq_1}
\end{theorem}

\section{InfoGeometry.MaxEnt.Finite.FiniteJaynesProblem.ofLogLikelihood}

\subsection*{eq_1}

\begin{theorem}[blueprint="infogeometry:maxent:finite:finitejaynesproblem:ofloglikelihood:eq_1"]
  \inputleannode{InfoGeometry.MaxEnt.Finite.FiniteJaynesProblem.ofLogLikelihood.eq_1}
\end{theorem}

\section{InfoGeometry.MaxEnt.IsMaxEntSolution}

\subsection*{hasExponentialRNForm_of_fenchel}

\begin{theorem}[blueprint="infogeometry:maxent:ismaxentsolution:hasexponentialrnform_of_fenchel"]
  \inputleannode{InfoGeometry.MaxEnt.IsMaxEntSolution.hasExponentialRNForm_of_fenchel}
\end{theorem}

\subsection*{toIntegrable}

\begin{theorem}[blueprint="infogeometry:maxent:ismaxentsolution:tointegrable"]
  \inputleannode{InfoGeometry.MaxEnt.IsMaxEntSolution.toIntegrable}
\end{theorem}

\section{InfoGeometry.MaxEnt.Jaynes}

\noindent\textbf{Buildable}: yes

\noindent\emph{No exported declarations in the current declaration graph.}

\section{InfoGeometry.MaxEnt.JaynesCanonical}

\noindent\textbf{Buildable}: yes

\subsection*{AsymptoticEquipartition}

\begin{theorem}[blueprint="infogeometry:maxent:jaynescanonical:asymptoticequipartition"]
  \inputleannode{InfoGeometry.MaxEnt.JaynesCanonical.AsymptoticEquipartition}
\end{theorem}

\subsection*{FeasibleSet}

\begin{theorem}[blueprint="infogeometry:maxent:jaynescanonical:feasibleset"]
  \inputleannode{InfoGeometry.MaxEnt.JaynesCanonical.FeasibleSet}
\end{theorem}

\subsection*{LinearConstraint}

\begin{theorem}[blueprint="infogeometry:maxent:jaynescanonical:linearconstraint"]
  \inputleannode{InfoGeometry.MaxEnt.JaynesCanonical.LinearConstraint}
\end{theorem}

\subsection*{ProbDist}

\begin{theorem}[blueprint="infogeometry:maxent:jaynescanonical:probdist"]
  \inputleannode{InfoGeometry.MaxEnt.JaynesCanonical.ProbDist}
\end{theorem}

\subsection*{asymptotic_equipartition}

\begin{theorem}[blueprint="infogeometry:maxent:jaynescanonical:asymptotic_equipartition"]
  \inputleannode{InfoGeometry.MaxEnt.JaynesCanonical.asymptotic_equipartition}
\end{theorem}

\subsection*{crossEntropyToGibbs_eq}

\begin{theorem}[blueprint="infogeometry:maxent:jaynescanonical:crossentropytogibbs_eq"]
  \inputleannode{InfoGeometry.MaxEnt.JaynesCanonical.crossEntropyToGibbs_eq}
\end{theorem}

\subsection*{empiricalAutocovariance}

\begin{theorem}[blueprint="infogeometry:maxent:jaynescanonical:empiricalautocovariance"]
  \inputleannode{InfoGeometry.MaxEnt.JaynesCanonical.empiricalAutocovariance}
\end{theorem}

\subsection*{entropy_concentration_bound}

\begin{theorem}[blueprint="infogeometry:maxent:jaynescanonical:entropy_concentration_bound"]
  \inputleannode{InfoGeometry.MaxEnt.JaynesCanonical.entropy_concentration_bound}
\end{theorem}

\subsection*{entropy_gibbs_eq}

\begin{theorem}[blueprint="infogeometry:maxent:jaynescanonical:entropy_gibbs_eq"]
  \inputleannode{InfoGeometry.MaxEnt.JaynesCanonical.entropy_gibbs_eq}
\end{theorem}

\subsection*{gibbsDist}

\begin{theorem}[blueprint="infogeometry:maxent:jaynescanonical:gibbsdist"]
  \inputleannode{InfoGeometry.MaxEnt.JaynesCanonical.gibbsDist}
\end{theorem}

\subsection*{gibbsDist_f}

\begin{theorem}[blueprint="infogeometry:maxent:jaynescanonical:gibbsdist_f"]
  \inputleannode{InfoGeometry.MaxEnt.JaynesCanonical.gibbsDist_f}
\end{theorem}

\subsection*{gibbsDist_pos}

\begin{theorem}[blueprint="infogeometry:maxent:jaynescanonical:gibbsdist_pos"]
  \inputleannode{InfoGeometry.MaxEnt.JaynesCanonical.gibbsDist_pos}
\end{theorem}

\subsection*{gibbs_is_maximum_entropy}

\begin{theorem}[blueprint="infogeometry:maxent:jaynescanonical:gibbs_is_maximum_entropy"]
  \inputleannode{InfoGeometry.MaxEnt.JaynesCanonical.gibbs_is_maximum_entropy}
\end{theorem}

\subsection*{isAutoregressiveModel}

\begin{theorem}[blueprint="infogeometry:maxent:jaynescanonical:isautoregressivemodel"]
  \inputleannode{InfoGeometry.MaxEnt.JaynesCanonical.isAutoregressiveModel}
\end{theorem}

\subsection*{logMultiplicity}

\begin{theorem}[blueprint="infogeometry:maxent:jaynescanonical:logmultiplicity"]
  \inputleannode{InfoGeometry.MaxEnt.JaynesCanonical.logMultiplicity}
\end{theorem}

\subsection*{logMultiplicity_nonneg}

\begin{theorem}[blueprint="infogeometry:maxent:jaynescanonical:logmultiplicity_nonneg"]
  \inputleannode{InfoGeometry.MaxEnt.JaynesCanonical.logMultiplicity_nonneg}
\end{theorem}

\subsection*{log_gibbsDist_f}

\begin{theorem}[blueprint="infogeometry:maxent:jaynescanonical:log_gibbsdist_f"]
  \inputleannode{InfoGeometry.MaxEnt.JaynesCanonical.log_gibbsDist_f}
\end{theorem}

\subsection*{multiplicity}

\begin{theorem}[blueprint="infogeometry:maxent:jaynescanonical:multiplicity"]
  \inputleannode{InfoGeometry.MaxEnt.JaynesCanonical.multiplicity}
\end{theorem}

\subsection*{multiplicity_one_le}

\begin{theorem}[blueprint="infogeometry:maxent:jaynescanonical:multiplicity_one_le"]
  \inputleannode{InfoGeometry.MaxEnt.JaynesCanonical.multiplicity_one_le}
\end{theorem}

\subsection*{multiplicity_pos}

\begin{theorem}[blueprint="infogeometry:maxent:jaynescanonical:multiplicity_pos"]
  \inputleannode{InfoGeometry.MaxEnt.JaynesCanonical.multiplicity_pos}
\end{theorem}

\subsection*{multiplicity_spec}

\begin{theorem}[blueprint="infogeometry:maxent:jaynescanonical:multiplicity_spec"]
  \inputleannode{InfoGeometry.MaxEnt.JaynesCanonical.multiplicity_spec}
\end{theorem}

\subsection*{partitionFunction}

\begin{theorem}[blueprint="infogeometry:maxent:jaynescanonical:partitionfunction"]
  \inputleannode{InfoGeometry.MaxEnt.JaynesCanonical.partitionFunction}
\end{theorem}

\subsection*{partitionFunction_pos}

\begin{theorem}[blueprint="infogeometry:maxent:jaynescanonical:partitionfunction_pos"]
  \inputleannode{InfoGeometry.MaxEnt.JaynesCanonical.partitionFunction_pos}
\end{theorem}

\subsection*{satisfiesConstraint}

\begin{theorem}[blueprint="infogeometry:maxent:jaynescanonical:satisfiesconstraint"]
  \inputleannode{InfoGeometry.MaxEnt.JaynesCanonical.satisfiesConstraint}
\end{theorem}

\subsection*{shannonEntropy}

\begin{theorem}[blueprint="infogeometry:maxent:jaynescanonical:shannonentropy"]
  \inputleannode{InfoGeometry.MaxEnt.JaynesCanonical.shannonEntropy}
\end{theorem}

\subsection*{totalCount}

\begin{theorem}[blueprint="infogeometry:maxent:jaynescanonical:totalcount"]
  \inputleannode{InfoGeometry.MaxEnt.JaynesCanonical.totalCount}
\end{theorem}

\section{InfoGeometry.MaxEnt.JaynesCanonical.LinearConstraint}

\subsection*{A}

\begin{theorem}[blueprint="infogeometry:maxent:jaynescanonical:linearconstraint:a"]
  \inputleannode{InfoGeometry.MaxEnt.JaynesCanonical.LinearConstraint.A}
\end{theorem}

\subsection*{casesOn}

\begin{theorem}[blueprint="infogeometry:maxent:jaynescanonical:linearconstraint:caseson"]
  \inputleannode{InfoGeometry.MaxEnt.JaynesCanonical.LinearConstraint.casesOn}
\end{theorem}

\subsection*{ctorIdx}

\begin{theorem}[blueprint="infogeometry:maxent:jaynescanonical:linearconstraint:ctoridx"]
  \inputleannode{InfoGeometry.MaxEnt.JaynesCanonical.LinearConstraint.ctorIdx}
\end{theorem}

\subsection*{d}

\begin{theorem}[blueprint="infogeometry:maxent:jaynescanonical:linearconstraint:d"]
  \inputleannode{InfoGeometry.MaxEnt.JaynesCanonical.LinearConstraint.d}
\end{theorem}

\subsection*{mk}

\begin{theorem}[blueprint="infogeometry:maxent:jaynescanonical:linearconstraint:mk"]
  \inputleannode{InfoGeometry.MaxEnt.JaynesCanonical.LinearConstraint.mk}
\end{theorem}

\subsection*{noConfusion}

\begin{theorem}[blueprint="infogeometry:maxent:jaynescanonical:linearconstraint:noconfusion"]
  \inputleannode{InfoGeometry.MaxEnt.JaynesCanonical.LinearConstraint.noConfusion}
\end{theorem}

\subsection*{noConfusionType}

\begin{theorem}[blueprint="infogeometry:maxent:jaynescanonical:linearconstraint:noconfusiontype"]
  \inputleannode{InfoGeometry.MaxEnt.JaynesCanonical.LinearConstraint.noConfusionType}
\end{theorem}

\subsection*{rec}

\begin{theorem}[blueprint="infogeometry:maxent:jaynescanonical:linearconstraint:rec"]
  \inputleannode{InfoGeometry.MaxEnt.JaynesCanonical.LinearConstraint.rec}
\end{theorem}

\subsection*{recOn}

\begin{theorem}[blueprint="infogeometry:maxent:jaynescanonical:linearconstraint:recon"]
  \inputleannode{InfoGeometry.MaxEnt.JaynesCanonical.LinearConstraint.recOn}
\end{theorem}

\section{InfoGeometry.MaxEnt.JaynesCanonical.ProbDist}

\subsection*{casesOn}

\begin{theorem}[blueprint="infogeometry:maxent:jaynescanonical:probdist:caseson"]
  \inputleannode{InfoGeometry.MaxEnt.JaynesCanonical.ProbDist.casesOn}
\end{theorem}

\subsection*{ctorIdx}

\begin{theorem}[blueprint="infogeometry:maxent:jaynescanonical:probdist:ctoridx"]
  \inputleannode{InfoGeometry.MaxEnt.JaynesCanonical.ProbDist.ctorIdx}
\end{theorem}

\subsection*{f}

\begin{theorem}[blueprint="infogeometry:maxent:jaynescanonical:probdist:f"]
  \inputleannode{InfoGeometry.MaxEnt.JaynesCanonical.ProbDist.f}
\end{theorem}

\subsection*{mk}

\begin{theorem}[blueprint="infogeometry:maxent:jaynescanonical:probdist:mk"]
  \inputleannode{InfoGeometry.MaxEnt.JaynesCanonical.ProbDist.mk}
\end{theorem}

\subsection*{noConfusion}

\begin{theorem}[blueprint="infogeometry:maxent:jaynescanonical:probdist:noconfusion"]
  \inputleannode{InfoGeometry.MaxEnt.JaynesCanonical.ProbDist.noConfusion}
\end{theorem}

\subsection*{noConfusionType}

\begin{theorem}[blueprint="infogeometry:maxent:jaynescanonical:probdist:noconfusiontype"]
  \inputleannode{InfoGeometry.MaxEnt.JaynesCanonical.ProbDist.noConfusionType}
\end{theorem}

\subsection*{nonneg}

\begin{theorem}[blueprint="infogeometry:maxent:jaynescanonical:probdist:nonneg"]
  \inputleannode{InfoGeometry.MaxEnt.JaynesCanonical.ProbDist.nonneg}
\end{theorem}

\subsection*{rec}

\begin{theorem}[blueprint="infogeometry:maxent:jaynescanonical:probdist:rec"]
  \inputleannode{InfoGeometry.MaxEnt.JaynesCanonical.ProbDist.rec}
\end{theorem}

\subsection*{recOn}

\begin{theorem}[blueprint="infogeometry:maxent:jaynescanonical:probdist:recon"]
  \inputleannode{InfoGeometry.MaxEnt.JaynesCanonical.ProbDist.recOn}
\end{theorem}

\subsection*{sum_one}

\begin{theorem}[blueprint="infogeometry:maxent:jaynescanonical:probdist:sum_one"]
  \inputleannode{InfoGeometry.MaxEnt.JaynesCanonical.ProbDist.sum_one}
\end{theorem}

\subsection*{toInfoProbabilityDist}

\begin{theorem}[blueprint="infogeometry:maxent:jaynescanonical:probdist:toinfoprobabilitydist"]
  \inputleannode{InfoGeometry.MaxEnt.JaynesCanonical.ProbDist.toInfoProbabilityDist}
\end{theorem}

\subsection*{toInfoProbabilityDist_prob}

\begin{theorem}[blueprint="infogeometry:maxent:jaynescanonical:probdist:toinfoprobabilitydist_prob"]
  \inputleannode{InfoGeometry.MaxEnt.JaynesCanonical.ProbDist.toInfoProbabilityDist_prob}
\end{theorem}

\section{InfoGeometry.MaxEnt.JaynesCanonical.ProbDist.toInfoProbabilityDist}

\subsection*{eq_1}

\begin{theorem}[blueprint="infogeometry:maxent:jaynescanonical:probdist:toinfoprobabilitydist:eq_1"]
  \inputleannode{InfoGeometry.MaxEnt.JaynesCanonical.ProbDist.toInfoProbabilityDist.eq_1}
\end{theorem}

\section{InfoGeometry.MaxEnt.JaynesCanonical.gibbsDist}

\subsection*{congr_simp}

\begin{theorem}[blueprint="infogeometry:maxent:jaynescanonical:gibbsdist:congr_simp"]
  \inputleannode{InfoGeometry.MaxEnt.JaynesCanonical.gibbsDist.congr_simp}
\end{theorem}

\section{InfoGeometry.MaxEnt.JaynesCanonical.multiplicity}

\subsection*{eq_1}

\begin{theorem}[blueprint="infogeometry:maxent:jaynescanonical:multiplicity:eq_1"]
  \inputleannode{InfoGeometry.MaxEnt.JaynesCanonical.multiplicity.eq_1}
\end{theorem}

\section{InfoGeometry.MaxEnt.JaynesCanonical.partitionFunction}

\subsection*{eq_1}

\begin{theorem}[blueprint="infogeometry:maxent:jaynescanonical:partitionfunction:eq_1"]
  \inputleannode{InfoGeometry.MaxEnt.JaynesCanonical.partitionFunction.eq_1}
\end{theorem}

\section{InfoGeometry.MaxEnt.JaynesCanonical.shannonEntropy}

\subsection*{eq_1}

\begin{theorem}[blueprint="infogeometry:maxent:jaynescanonical:shannonentropy:eq_1"]
  \inputleannode{InfoGeometry.MaxEnt.JaynesCanonical.shannonEntropy.eq_1}
\end{theorem}

\section{InfoGeometry.MaxEnt.JaynesCanonical.totalCount}

\subsection*{eq_1}

\begin{theorem}[blueprint="infogeometry:maxent:jaynescanonical:totalcount:eq_1"]
  \inputleannode{InfoGeometry.MaxEnt.JaynesCanonical.totalCount.eq_1}
\end{theorem}

\section{InfoGeometry.MaxEnt.JaynesInfoStatMech}

\noindent\textbf{Buildable}: yes

\subsection*{ConstraintFamily}

\begin{theorem}[blueprint="infogeometry:maxent:jaynesinfostatmech:constraintfamily"]
  \inputleannode{InfoGeometry.MaxEnt.JaynesInfoStatMech.ConstraintFamily}
\end{theorem}

\subsection*{IsFeasible}

\begin{theorem}[blueprint="infogeometry:maxent:jaynesinfostatmech:isfeasible"]
  \inputleannode{InfoGeometry.MaxEnt.JaynesInfoStatMech.IsFeasible}
\end{theorem}

\subsection*{Observable}

\begin{theorem}[blueprint="infogeometry:maxent:jaynesinfostatmech:observable"]
  \inputleannode{InfoGeometry.MaxEnt.JaynesInfoStatMech.Observable}
\end{theorem}

\subsection*{ProbDist}

\begin{theorem}[blueprint="infogeometry:maxent:jaynesinfostatmech:probdist"]
  \inputleannode{InfoGeometry.MaxEnt.JaynesInfoStatMech.ProbDist}
\end{theorem}

\subsection*{ThermodynamicRegime}

\begin{theorem}[blueprint="infogeometry:maxent:jaynesinfostatmech:thermodynamicregime"]
  \inputleannode{InfoGeometry.MaxEnt.JaynesInfoStatMech.ThermodynamicRegime}
\end{theorem}

\subsection*{beta}

\begin{theorem}[blueprint="infogeometry:maxent:jaynesinfostatmech:beta"]
  \inputleannode{InfoGeometry.MaxEnt.JaynesInfoStatMech.beta}
\end{theorem}

\subsection*{betaInRegime}

\begin{theorem}[blueprint="infogeometry:maxent:jaynesinfostatmech:betainregime"]
  \inputleannode{InfoGeometry.MaxEnt.JaynesInfoStatMech.betaInRegime}
\end{theorem}

\subsection*{boltzmannDist}

\begin{theorem}[blueprint="infogeometry:maxent:jaynesinfostatmech:boltzmanndist"]
  \inputleannode{InfoGeometry.MaxEnt.JaynesInfoStatMech.boltzmannDist}
\end{theorem}

\subsection*{energyFamily}

\begin{theorem}[blueprint="infogeometry:maxent:jaynesinfostatmech:energyfamily"]
  \inputleannode{InfoGeometry.MaxEnt.JaynesInfoStatMech.energyFamily}
\end{theorem}

\subsection*{expectedValue}

\begin{theorem}[blueprint="infogeometry:maxent:jaynesinfostatmech:expectedvalue"]
  \inputleannode{InfoGeometry.MaxEnt.JaynesInfoStatMech.expectedValue}
\end{theorem}

\subsection*{helmholtzFreeEnergy}

\begin{theorem}[blueprint="infogeometry:maxent:jaynesinfostatmech:helmholtzfreeenergy"]
  \inputleannode{InfoGeometry.MaxEnt.JaynesInfoStatMech.helmholtzFreeEnergy}
\end{theorem}

\subsection*{logZ}

\begin{theorem}[blueprint="infogeometry:maxent:jaynesinfostatmech:logz"]
  \inputleannode{InfoGeometry.MaxEnt.JaynesInfoStatMech.logZ}
\end{theorem}

\subsection*{maxEntDist}

\begin{theorem}[blueprint="infogeometry:maxent:jaynesinfostatmech:maxentdist"]
  \inputleannode{InfoGeometry.MaxEnt.JaynesInfoStatMech.maxEntDist}
\end{theorem}

\subsection*{partitionFunction}

\begin{theorem}[blueprint="infogeometry:maxent:jaynesinfostatmech:partitionfunction"]
  \inputleannode{InfoGeometry.MaxEnt.JaynesInfoStatMech.partitionFunction}
\end{theorem}

\subsection*{partitionFunction_pos}

\begin{theorem}[blueprint="infogeometry:maxent:jaynesinfostatmech:partitionfunction_pos"]
  \inputleannode{InfoGeometry.MaxEnt.JaynesInfoStatMech.partitionFunction_pos}
\end{theorem}

\subsection*{shannonEntropy}

\begin{theorem}[blueprint="infogeometry:maxent:jaynesinfostatmech:shannonentropy"]
  \inputleannode{InfoGeometry.MaxEnt.JaynesInfoStatMech.shannonEntropy}
\end{theorem}

\subsection*{weight}

\begin{theorem}[blueprint="infogeometry:maxent:jaynesinfostatmech:weight"]
  \inputleannode{InfoGeometry.MaxEnt.JaynesInfoStatMech.weight}
\end{theorem}

\section{InfoGeometry.MaxEnt.JaynesInfoStatMech.ConstraintFamily}

\subsection*{casesOn}

\begin{theorem}[blueprint="infogeometry:maxent:jaynesinfostatmech:constraintfamily:caseson"]
  \inputleannode{InfoGeometry.MaxEnt.JaynesInfoStatMech.ConstraintFamily.casesOn}
\end{theorem}

\subsection*{ctorIdx}

\begin{theorem}[blueprint="infogeometry:maxent:jaynesinfostatmech:constraintfamily:ctoridx"]
  \inputleannode{InfoGeometry.MaxEnt.JaynesInfoStatMech.ConstraintFamily.ctorIdx}
\end{theorem}

\subsection*{d}

\begin{theorem}[blueprint="infogeometry:maxent:jaynesinfostatmech:constraintfamily:d"]
  \inputleannode{InfoGeometry.MaxEnt.JaynesInfoStatMech.ConstraintFamily.d}
\end{theorem}

\subsection*{f}

\begin{theorem}[blueprint="infogeometry:maxent:jaynesinfostatmech:constraintfamily:f"]
  \inputleannode{InfoGeometry.MaxEnt.JaynesInfoStatMech.ConstraintFamily.f}
\end{theorem}

\subsection*{mk}

\begin{theorem}[blueprint="infogeometry:maxent:jaynesinfostatmech:constraintfamily:mk"]
  \inputleannode{InfoGeometry.MaxEnt.JaynesInfoStatMech.ConstraintFamily.mk}
\end{theorem}

\subsection*{noConfusion}

\begin{theorem}[blueprint="infogeometry:maxent:jaynesinfostatmech:constraintfamily:noconfusion"]
  \inputleannode{InfoGeometry.MaxEnt.JaynesInfoStatMech.ConstraintFamily.noConfusion}
\end{theorem}

\subsection*{noConfusionType}

\begin{theorem}[blueprint="infogeometry:maxent:jaynesinfostatmech:constraintfamily:noconfusiontype"]
  \inputleannode{InfoGeometry.MaxEnt.JaynesInfoStatMech.ConstraintFamily.noConfusionType}
\end{theorem}

\subsection*{rec}

\begin{theorem}[blueprint="infogeometry:maxent:jaynesinfostatmech:constraintfamily:rec"]
  \inputleannode{InfoGeometry.MaxEnt.JaynesInfoStatMech.ConstraintFamily.rec}
\end{theorem}

\subsection*{recOn}

\begin{theorem}[blueprint="infogeometry:maxent:jaynesinfostatmech:constraintfamily:recon"]
  \inputleannode{InfoGeometry.MaxEnt.JaynesInfoStatMech.ConstraintFamily.recOn}
\end{theorem}

\section{InfoGeometry.MaxEnt.JaynesInfoStatMech.ProbDist}

\subsection*{casesOn}

\begin{theorem}[blueprint="infogeometry:maxent:jaynesinfostatmech:probdist:caseson"]
  \inputleannode{InfoGeometry.MaxEnt.JaynesInfoStatMech.ProbDist.casesOn}
\end{theorem}

\subsection*{ctorIdx}

\begin{theorem}[blueprint="infogeometry:maxent:jaynesinfostatmech:probdist:ctoridx"]
  \inputleannode{InfoGeometry.MaxEnt.JaynesInfoStatMech.ProbDist.ctorIdx}
\end{theorem}

\subsection*{mk}

\begin{theorem}[blueprint="infogeometry:maxent:jaynesinfostatmech:probdist:mk"]
  \inputleannode{InfoGeometry.MaxEnt.JaynesInfoStatMech.ProbDist.mk}
\end{theorem}

\subsection*{noConfusion}

\begin{theorem}[blueprint="infogeometry:maxent:jaynesinfostatmech:probdist:noconfusion"]
  \inputleannode{InfoGeometry.MaxEnt.JaynesInfoStatMech.ProbDist.noConfusion}
\end{theorem}

\subsection*{noConfusionType}

\begin{theorem}[blueprint="infogeometry:maxent:jaynesinfostatmech:probdist:noconfusiontype"]
  \inputleannode{InfoGeometry.MaxEnt.JaynesInfoStatMech.ProbDist.noConfusionType}
\end{theorem}

\subsection*{nonneg}

\begin{theorem}[blueprint="infogeometry:maxent:jaynesinfostatmech:probdist:nonneg"]
  \inputleannode{InfoGeometry.MaxEnt.JaynesInfoStatMech.ProbDist.nonneg}
\end{theorem}

\subsection*{p}

\begin{theorem}[blueprint="infogeometry:maxent:jaynesinfostatmech:probdist:p"]
  \inputleannode{InfoGeometry.MaxEnt.JaynesInfoStatMech.ProbDist.p}
\end{theorem}

\subsection*{rec}

\begin{theorem}[blueprint="infogeometry:maxent:jaynesinfostatmech:probdist:rec"]
  \inputleannode{InfoGeometry.MaxEnt.JaynesInfoStatMech.ProbDist.rec}
\end{theorem}

\subsection*{recOn}

\begin{theorem}[blueprint="infogeometry:maxent:jaynesinfostatmech:probdist:recon"]
  \inputleannode{InfoGeometry.MaxEnt.JaynesInfoStatMech.ProbDist.recOn}
\end{theorem}

\subsection*{sum_one}

\begin{theorem}[blueprint="infogeometry:maxent:jaynesinfostatmech:probdist:sum_one"]
  \inputleannode{InfoGeometry.MaxEnt.JaynesInfoStatMech.ProbDist.sum_one}
\end{theorem}

\section{InfoGeometry.MaxEnt.JaynesInfoStatMech.ThermalDiagonal}

\subsection*{DiagObservable}

\begin{theorem}[blueprint="infogeometry:maxent:jaynesinfostatmech:thermaldiagonal:diagobservable"]
  \inputleannode{InfoGeometry.MaxEnt.JaynesInfoStatMech.ThermalDiagonal.DiagObservable}
\end{theorem}

\subsection*{FinMat}

\begin{theorem}[blueprint="infogeometry:maxent:jaynesinfostatmech:thermaldiagonal:finmat"]
  \inputleannode{InfoGeometry.MaxEnt.JaynesInfoStatMech.ThermalDiagonal.FinMat}
\end{theorem}

\subsection*{densityMatrix}

\begin{theorem}[blueprint="infogeometry:maxent:jaynesinfostatmech:thermaldiagonal:densitymatrix"]
  \inputleannode{InfoGeometry.MaxEnt.JaynesInfoStatMech.ThermalDiagonal.densityMatrix}
\end{theorem}

\subsection*{densityMatrix_diag}

\begin{theorem}[blueprint="infogeometry:maxent:jaynesinfostatmech:thermaldiagonal:densitymatrix_diag"]
  \inputleannode{InfoGeometry.MaxEnt.JaynesInfoStatMech.ThermalDiagonal.densityMatrix_diag}
\end{theorem}

\subsection*{densityMatrix_offdiag}

\begin{theorem}[blueprint="infogeometry:maxent:jaynesinfostatmech:thermaldiagonal:densitymatrix_offdiag"]
  \inputleannode{InfoGeometry.MaxEnt.JaynesInfoStatMech.ThermalDiagonal.densityMatrix_offdiag}
\end{theorem}

\subsection*{densityMatrix_trace_one}

\begin{theorem}[blueprint="infogeometry:maxent:jaynesinfostatmech:thermaldiagonal:densitymatrix_trace_one"]
  \inputleannode{InfoGeometry.MaxEnt.JaynesInfoStatMech.ThermalDiagonal.densityMatrix_trace_one}
\end{theorem}

\subsection*{diagMatrix}

\begin{theorem}[blueprint="infogeometry:maxent:jaynesinfostatmech:thermaldiagonal:diagmatrix"]
  \inputleannode{InfoGeometry.MaxEnt.JaynesInfoStatMech.ThermalDiagonal.diagMatrix}
\end{theorem}

\subsection*{diagMul}

\begin{theorem}[blueprint="infogeometry:maxent:jaynesinfostatmech:thermaldiagonal:diagmul"]
  \inputleannode{InfoGeometry.MaxEnt.JaynesInfoStatMech.ThermalDiagonal.diagMul}
\end{theorem}

\subsection*{gibbsEntropy}

\begin{theorem}[blueprint="infogeometry:maxent:jaynesinfostatmech:thermaldiagonal:gibbsentropy"]
  \inputleannode{InfoGeometry.MaxEnt.JaynesInfoStatMech.ThermalDiagonal.gibbsEntropy}
\end{theorem}

\subsection*{gibbsEntropy_eq_beta_internal_plus_logPartition}

\begin{theorem}[blueprint="infogeometry:maxent:jaynesinfostatmech:thermaldiagonal:gibbsentropy_eq_beta_internal_plus_logpartition"]
  \inputleannode{InfoGeometry.MaxEnt.JaynesInfoStatMech.ThermalDiagonal.gibbsEntropy_eq_beta_internal_plus_logPartition}
\end{theorem}

\subsection*{gibbsProb}

\begin{theorem}[blueprint="infogeometry:maxent:jaynesinfostatmech:thermaldiagonal:gibbsprob"]
  \inputleannode{InfoGeometry.MaxEnt.JaynesInfoStatMech.ThermalDiagonal.gibbsProb}
\end{theorem}

\subsection*{gibbsProb_nonneg}

\begin{theorem}[blueprint="infogeometry:maxent:jaynesinfostatmech:thermaldiagonal:gibbsprob_nonneg"]
  \inputleannode{InfoGeometry.MaxEnt.JaynesInfoStatMech.ThermalDiagonal.gibbsProb_nonneg}
\end{theorem}

\subsection*{gibbsProb_pos}

\begin{theorem}[blueprint="infogeometry:maxent:jaynesinfostatmech:thermaldiagonal:gibbsprob_pos"]
  \inputleannode{InfoGeometry.MaxEnt.JaynesInfoStatMech.ThermalDiagonal.gibbsProb_pos}
\end{theorem}

\subsection*{gibbsProb_sum_one}

\begin{theorem}[blueprint="infogeometry:maxent:jaynesinfostatmech:thermaldiagonal:gibbsprob_sum_one"]
  \inputleannode{InfoGeometry.MaxEnt.JaynesInfoStatMech.ThermalDiagonal.gibbsProb_sum_one}
\end{theorem}

\subsection*{gibbsStateDiag}

\begin{theorem}[blueprint="infogeometry:maxent:jaynesinfostatmech:thermaldiagonal:gibbsstatediag"]
  \inputleannode{InfoGeometry.MaxEnt.JaynesInfoStatMech.ThermalDiagonal.gibbsStateDiag}
\end{theorem}

\subsection*{gibbsStateMatrix}

\begin{theorem}[blueprint="infogeometry:maxent:jaynesinfostatmech:thermaldiagonal:gibbsstatematrix"]
  \inputleannode{InfoGeometry.MaxEnt.JaynesInfoStatMech.ThermalDiagonal.gibbsStateMatrix}
\end{theorem}

\subsection*{gibbsStateMatrix_on_diag}

\begin{theorem}[blueprint="infogeometry:maxent:jaynesinfostatmech:thermaldiagonal:gibbsstatematrix_on_diag"]
  \inputleannode{InfoGeometry.MaxEnt.JaynesInfoStatMech.ThermalDiagonal.gibbsStateMatrix_on_diag}
\end{theorem}

\subsection*{gibbsState_kmsLike_diag}

\begin{theorem}[blueprint="infogeometry:maxent:jaynesinfostatmech:thermaldiagonal:gibbsstate_kmslike_diag"]
  \inputleannode{InfoGeometry.MaxEnt.JaynesInfoStatMech.ThermalDiagonal.gibbsState_kmsLike_diag}
\end{theorem}

\subsection*{gibbsWeight}

\begin{theorem}[blueprint="infogeometry:maxent:jaynesinfostatmech:thermaldiagonal:gibbsweight"]
  \inputleannode{InfoGeometry.MaxEnt.JaynesInfoStatMech.ThermalDiagonal.gibbsWeight}
\end{theorem}

\subsection*{hamiltonianMatrix}

\begin{theorem}[blueprint="infogeometry:maxent:jaynesinfostatmech:thermaldiagonal:hamiltonianmatrix"]
  \inputleannode{InfoGeometry.MaxEnt.JaynesInfoStatMech.ThermalDiagonal.hamiltonianMatrix}
\end{theorem}

\subsection*{helmholtzFreeEnergy}

\begin{theorem}[blueprint="infogeometry:maxent:jaynesinfostatmech:thermaldiagonal:helmholtzfreeenergy"]
  \inputleannode{InfoGeometry.MaxEnt.JaynesInfoStatMech.ThermalDiagonal.helmholtzFreeEnergy}
\end{theorem}

\subsection*{internalEnergy}

\begin{theorem}[blueprint="infogeometry:maxent:jaynesinfostatmech:thermaldiagonal:internalenergy"]
  \inputleannode{InfoGeometry.MaxEnt.JaynesInfoStatMech.ThermalDiagonal.internalEnergy}
\end{theorem}

\subsection*{logDensityMatrix}

\begin{theorem}[blueprint="infogeometry:maxent:jaynesinfostatmech:thermaldiagonal:logdensitymatrix"]
  \inputleannode{InfoGeometry.MaxEnt.JaynesInfoStatMech.ThermalDiagonal.logDensityMatrix}
\end{theorem}

\subsection*{logDensityMatrix_diag}

\begin{theorem}[blueprint="infogeometry:maxent:jaynesinfostatmech:thermaldiagonal:logdensitymatrix_diag"]
  \inputleannode{InfoGeometry.MaxEnt.JaynesInfoStatMech.ThermalDiagonal.logDensityMatrix_diag}
\end{theorem}

\subsection*{logDensityMatrix_diag_affine}

\begin{theorem}[blueprint="infogeometry:maxent:jaynesinfostatmech:thermaldiagonal:logdensitymatrix_diag_affine"]
  \inputleannode{InfoGeometry.MaxEnt.JaynesInfoStatMech.ThermalDiagonal.logDensityMatrix_diag_affine}
\end{theorem}

\subsection*{logPartition}

\begin{theorem}[blueprint="infogeometry:maxent:jaynesinfostatmech:thermaldiagonal:logpartition"]
  \inputleannode{InfoGeometry.MaxEnt.JaynesInfoStatMech.ThermalDiagonal.logPartition}
\end{theorem}

\subsection*{log_gibbsProb}

\begin{theorem}[blueprint="infogeometry:maxent:jaynesinfostatmech:thermaldiagonal:log_gibbsprob"]
  \inputleannode{InfoGeometry.MaxEnt.JaynesInfoStatMech.ThermalDiagonal.log_gibbsProb}
\end{theorem}

\subsection*{matrixTrace}

\begin{theorem}[blueprint="infogeometry:maxent:jaynesinfostatmech:thermaldiagonal:matrixtrace"]
  \inputleannode{InfoGeometry.MaxEnt.JaynesInfoStatMech.ThermalDiagonal.matrixTrace}
\end{theorem}

\subsection*{modularConj}

\begin{theorem}[blueprint="infogeometry:maxent:jaynesinfostatmech:thermaldiagonal:modularconj"]
  \inputleannode{InfoGeometry.MaxEnt.JaynesInfoStatMech.ThermalDiagonal.modularConj}
\end{theorem}

\subsection*{modularConj_diag_entry}

\begin{theorem}[blueprint="infogeometry:maxent:jaynesinfostatmech:thermaldiagonal:modularconj_diag_entry"]
  \inputleannode{InfoGeometry.MaxEnt.JaynesInfoStatMech.ThermalDiagonal.modularConj_diag_entry}
\end{theorem}

\subsection*{modularShiftDiag}

\begin{theorem}[blueprint="infogeometry:maxent:jaynesinfostatmech:thermaldiagonal:modularshiftdiag"]
  \inputleannode{InfoGeometry.MaxEnt.JaynesInfoStatMech.ThermalDiagonal.modularShiftDiag}
\end{theorem}

\subsection*{modularShiftDiag_eq}

\begin{theorem}[blueprint="infogeometry:maxent:jaynesinfostatmech:thermaldiagonal:modularshiftdiag_eq"]
  \inputleannode{InfoGeometry.MaxEnt.JaynesInfoStatMech.ThermalDiagonal.modularShiftDiag_eq}
\end{theorem}

\subsection*{partitionFunction}

\begin{theorem}[blueprint="infogeometry:maxent:jaynesinfostatmech:thermaldiagonal:partitionfunction"]
  \inputleannode{InfoGeometry.MaxEnt.JaynesInfoStatMech.ThermalDiagonal.partitionFunction}
\end{theorem}

\subsection*{partitionFunction_ne_zero}

\begin{theorem}[blueprint="infogeometry:maxent:jaynesinfostatmech:thermaldiagonal:partitionfunction_ne_zero"]
  \inputleannode{InfoGeometry.MaxEnt.JaynesInfoStatMech.ThermalDiagonal.partitionFunction_ne_zero}
\end{theorem}

\subsection*{partitionFunction_pos}

\begin{theorem}[blueprint="infogeometry:maxent:jaynesinfostatmech:thermaldiagonal:partitionfunction_pos"]
  \inputleannode{InfoGeometry.MaxEnt.JaynesInfoStatMech.ThermalDiagonal.partitionFunction_pos}
\end{theorem}

\section{InfoGeometry.MaxEnt.JaynesInfoStatMech.ThermalDiagonal.densityMatrix}

\subsection*{eq_1}

\begin{theorem}[blueprint="infogeometry:maxent:jaynesinfostatmech:thermaldiagonal:densitymatrix:eq_1"]
  \inputleannode{InfoGeometry.MaxEnt.JaynesInfoStatMech.ThermalDiagonal.densityMatrix.eq_1}
\end{theorem}

\section{InfoGeometry.MaxEnt.JaynesInfoStatMech.ThermalDiagonal.diagMatrix}

\subsection*{eq_1}

\begin{theorem}[blueprint="infogeometry:maxent:jaynesinfostatmech:thermaldiagonal:diagmatrix:eq_1"]
  \inputleannode{InfoGeometry.MaxEnt.JaynesInfoStatMech.ThermalDiagonal.diagMatrix.eq_1}
\end{theorem}

\section{InfoGeometry.MaxEnt.JaynesInfoStatMech.ThermalDiagonal.logDensityMatrix}

\subsection*{eq_1}

\begin{theorem}[blueprint="infogeometry:maxent:jaynesinfostatmech:thermaldiagonal:logdensitymatrix:eq_1"]
  \inputleannode{InfoGeometry.MaxEnt.JaynesInfoStatMech.ThermalDiagonal.logDensityMatrix.eq_1}
\end{theorem}

\section{InfoGeometry.MaxEnt.JaynesInfoStatMech.ThermalDiagonal.modularConj}

\subsection*{eq_1}

\begin{theorem}[blueprint="infogeometry:maxent:jaynesinfostatmech:thermaldiagonal:modularconj:eq_1"]
  \inputleannode{InfoGeometry.MaxEnt.JaynesInfoStatMech.ThermalDiagonal.modularConj.eq_1}
\end{theorem}

\section{InfoGeometry.MaxEnt.JaynesInfoStatMech.partitionFunction}

\subsection*{eq_1}

\begin{theorem}[blueprint="infogeometry:maxent:jaynesinfostatmech:partitionfunction:eq_1"]
  \inputleannode{InfoGeometry.MaxEnt.JaynesInfoStatMech.partitionFunction.eq_1}
\end{theorem}

\section{InfoGeometry.MaxEnt.JaynesInfoStatMech.weight}

\subsection*{eq_1}

\begin{theorem}[blueprint="infogeometry:maxent:jaynesinfostatmech:weight:eq_1"]
  \inputleannode{InfoGeometry.MaxEnt.JaynesInfoStatMech.weight.eq_1}
\end{theorem}

\section{InfoGeometry.MaxEnt.JaynesInfoStatMechTest}

\noindent\textbf{Buildable}: yes

\noindent\emph{No exported declarations in the current declaration graph.}

\section{InfoGeometry.MaxEnt.JaynesRNMaxEnt}

\noindent\textbf{Buildable}: yes

\subsection*{FeasibleSet}

\begin{theorem}[blueprint="infogeometry:maxent:jaynesrnmaxent:feasibleset"]
  \inputleannode{InfoGeometry.MaxEnt.JaynesRNMaxEnt.FeasibleSet}
\end{theorem}

\subsection*{FeasibleSetIntegrable}

\begin{theorem}[blueprint="infogeometry:maxent:jaynesrnmaxent:feasiblesetintegrable"]
  \inputleannode{InfoGeometry.MaxEnt.JaynesRNMaxEnt.FeasibleSetIntegrable}
\end{theorem}

\subsection*{FeasibleSetIntegrable_subset_FeasibleSet}

\begin{theorem}[blueprint="infogeometry:maxent:jaynesrnmaxent:feasiblesetintegrable_subset_feasibleset"]
  \inputleannode{InfoGeometry.MaxEnt.JaynesRNMaxEnt.FeasibleSetIntegrable_subset_FeasibleSet}
\end{theorem}

\subsection*{GibbsMinimizesKL}

\begin{theorem}[blueprint="infogeometry:maxent:jaynesrnmaxent:gibbsminimizeskl"]
  \inputleannode{InfoGeometry.MaxEnt.JaynesRNMaxEnt.GibbsMinimizesKL}
\end{theorem}

\subsection*{MomentFamily}

\begin{theorem}[blueprint="infogeometry:maxent:jaynesrnmaxent:momentfamily"]
  \inputleannode{InfoGeometry.MaxEnt.JaynesRNMaxEnt.MomentFamily}
\end{theorem}

\subsection*{PartitionIntegrable}

\begin{theorem}[blueprint="infogeometry:maxent:jaynesrnmaxent:partitionintegrable"]
  \inputleannode{InfoGeometry.MaxEnt.JaynesRNMaxEnt.PartitionIntegrable}
\end{theorem}

\subsection*{Satisfies}

\begin{theorem}[blueprint="infogeometry:maxent:jaynesrnmaxent:satisfies"]
  \inputleannode{InfoGeometry.MaxEnt.JaynesRNMaxEnt.Satisfies}
\end{theorem}

\subsection*{SatisfiesIntegrable}

\begin{theorem}[blueprint="infogeometry:maxent:jaynesrnmaxent:satisfiesintegrable"]
  \inputleannode{InfoGeometry.MaxEnt.JaynesRNMaxEnt.SatisfiesIntegrable}
\end{theorem}

\subsection*{aemeasurable_potential}

\begin{theorem}[blueprint="infogeometry:maxent:jaynesrnmaxent:aemeasurable_potential"]
  \inputleannode{InfoGeometry.MaxEnt.JaynesRNMaxEnt.aemeasurable_potential}
\end{theorem}

\subsection*{gibbs}

\begin{theorem}[blueprint="infogeometry:maxent:jaynesrnmaxent:gibbs"]
  \inputleannode{InfoGeometry.MaxEnt.JaynesRNMaxEnt.gibbs}
\end{theorem}

\subsection*{gibbsMeasure}

\begin{theorem}[blueprint="infogeometry:maxent:jaynesrnmaxent:gibbsmeasure"]
  \inputleannode{InfoGeometry.MaxEnt.JaynesRNMaxEnt.gibbsMeasure}
\end{theorem}

\subsection*{gibbsMeasure_ac}

\begin{theorem}[blueprint="infogeometry:maxent:jaynesrnmaxent:gibbsmeasure_ac"]
  \inputleannode{InfoGeometry.MaxEnt.JaynesRNMaxEnt.gibbsMeasure_ac}
\end{theorem}

\subsection*{gibbs_minimizes_kl}

\begin{theorem}[blueprint="infogeometry:maxent:jaynesrnmaxent:gibbs_minimizes_kl"]
  \inputleannode{InfoGeometry.MaxEnt.JaynesRNMaxEnt.gibbs_minimizes_kl}
\end{theorem}

\subsection*{objectiveKL}

\begin{theorem}[blueprint="infogeometry:maxent:jaynesrnmaxent:objectivekl"]
  \inputleannode{InfoGeometry.MaxEnt.JaynesRNMaxEnt.objectiveKL}
\end{theorem}

\subsection*{partitionFunction}

\begin{theorem}[blueprint="infogeometry:maxent:jaynesrnmaxent:partitionfunction"]
  \inputleannode{InfoGeometry.MaxEnt.JaynesRNMaxEnt.partitionFunction}
\end{theorem}

\subsection*{partitionFunction_nonneg}

\begin{theorem}[blueprint="infogeometry:maxent:jaynesrnmaxent:partitionfunction_nonneg"]
  \inputleannode{InfoGeometry.MaxEnt.JaynesRNMaxEnt.partitionFunction_nonneg}
\end{theorem}

\subsection*{potential}

\begin{theorem}[blueprint="infogeometry:maxent:jaynesrnmaxent:potential"]
  \inputleannode{InfoGeometry.MaxEnt.JaynesRNMaxEnt.potential}
\end{theorem}

\subsection*{rnDeriv_gibbsMeasure_eq}

\begin{theorem}[blueprint="infogeometry:maxent:jaynesrnmaxent:rnderiv_gibbsmeasure_eq"]
  \inputleannode{InfoGeometry.MaxEnt.JaynesRNMaxEnt.rnDeriv_gibbsMeasure_eq}
\end{theorem}

\subsection*{rnDeriv_gibbsMeasure_toReal_eq}

\begin{theorem}[blueprint="infogeometry:maxent:jaynesrnmaxent:rnderiv_gibbsmeasure_toreal_eq"]
  \inputleannode{InfoGeometry.MaxEnt.JaynesRNMaxEnt.rnDeriv_gibbsMeasure_toReal_eq}
\end{theorem}

\section{InfoGeometry.MaxEnt.JaynesRNMaxEnt.MomentFamily}

\subsection*{casesOn}

\begin{theorem}[blueprint="infogeometry:maxent:jaynesrnmaxent:momentfamily:caseson"]
  \inputleannode{InfoGeometry.MaxEnt.JaynesRNMaxEnt.MomentFamily.casesOn}
\end{theorem}

\subsection*{ctorIdx}

\begin{theorem}[blueprint="infogeometry:maxent:jaynesrnmaxent:momentfamily:ctoridx"]
  \inputleannode{InfoGeometry.MaxEnt.JaynesRNMaxEnt.MomentFamily.ctorIdx}
\end{theorem}

\subsection*{d}

\begin{theorem}[blueprint="infogeometry:maxent:jaynesrnmaxent:momentfamily:d"]
  \inputleannode{InfoGeometry.MaxEnt.JaynesRNMaxEnt.MomentFamily.d}
\end{theorem}

\subsection*{f}

\begin{theorem}[blueprint="infogeometry:maxent:jaynesrnmaxent:momentfamily:f"]
  \inputleannode{InfoGeometry.MaxEnt.JaynesRNMaxEnt.MomentFamily.f}
\end{theorem}

\subsection*{measurable_f}

\begin{theorem}[blueprint="infogeometry:maxent:jaynesrnmaxent:momentfamily:measurable_f"]
  \inputleannode{InfoGeometry.MaxEnt.JaynesRNMaxEnt.MomentFamily.measurable_f}
\end{theorem}

\subsection*{mk}

\begin{theorem}[blueprint="infogeometry:maxent:jaynesrnmaxent:momentfamily:mk"]
  \inputleannode{InfoGeometry.MaxEnt.JaynesRNMaxEnt.MomentFamily.mk}
\end{theorem}

\subsection*{noConfusion}

\begin{theorem}[blueprint="infogeometry:maxent:jaynesrnmaxent:momentfamily:noconfusion"]
  \inputleannode{InfoGeometry.MaxEnt.JaynesRNMaxEnt.MomentFamily.noConfusion}
\end{theorem}

\subsection*{noConfusionType}

\begin{theorem}[blueprint="infogeometry:maxent:jaynesrnmaxent:momentfamily:noconfusiontype"]
  \inputleannode{InfoGeometry.MaxEnt.JaynesRNMaxEnt.MomentFamily.noConfusionType}
\end{theorem}

\subsection*{rec}

\begin{theorem}[blueprint="infogeometry:maxent:jaynesrnmaxent:momentfamily:rec"]
  \inputleannode{InfoGeometry.MaxEnt.JaynesRNMaxEnt.MomentFamily.rec}
\end{theorem}

\subsection*{recOn}

\begin{theorem}[blueprint="infogeometry:maxent:jaynesrnmaxent:momentfamily:recon"]
  \inputleannode{InfoGeometry.MaxEnt.JaynesRNMaxEnt.MomentFamily.recOn}
\end{theorem}

\section{InfoGeometry.MaxEnt.JaynesRNMaxEnt.SatisfiesIntegrable}

\subsection*{satisfies}

\begin{theorem}[blueprint="infogeometry:maxent:jaynesrnmaxent:satisfiesintegrable:satisfies"]
  \inputleannode{InfoGeometry.MaxEnt.JaynesRNMaxEnt.SatisfiesIntegrable.satisfies}
\end{theorem}

\section{InfoGeometry.MaxEnt.JaynesRNMaxEnt.gibbsMeasure}

\subsection*{eq_1}

\begin{theorem}[blueprint="infogeometry:maxent:jaynesrnmaxent:gibbsmeasure:eq_1"]
  \inputleannode{InfoGeometry.MaxEnt.JaynesRNMaxEnt.gibbsMeasure.eq_1}
\end{theorem}

\section{InfoGeometry.MaxEnt.JaynesRNMaxEnt.partitionFunction}

\subsection*{eq_1}

\begin{theorem}[blueprint="infogeometry:maxent:jaynesrnmaxent:partitionfunction:eq_1"]
  \inputleannode{InfoGeometry.MaxEnt.JaynesRNMaxEnt.partitionFunction.eq_1}
\end{theorem}

\section{InfoGeometry.MaxEnt.Lagrange}

\noindent\textbf{Buildable}: yes

\noindent\emph{No exported declarations in the current declaration graph.}

\section{InfoGeometry.MaxEnt.LinearConstraint}

\subsection*{IntegrableUnder}

\begin{theorem}[blueprint="infogeometry:maxent:linearconstraint:integrableunder"]
  \inputleannode{InfoGeometry.MaxEnt.LinearConstraint.IntegrableUnder}
\end{theorem}

\subsection*{Satisfies}

\begin{theorem}[blueprint="infogeometry:maxent:linearconstraint:satisfies"]
  \inputleannode{InfoGeometry.MaxEnt.LinearConstraint.Satisfies}
\end{theorem}

\subsection*{SatisfiesIntegrable}

\begin{theorem}[blueprint="infogeometry:maxent:linearconstraint:satisfiesintegrable"]
  \inputleannode{InfoGeometry.MaxEnt.LinearConstraint.SatisfiesIntegrable}
\end{theorem}

\subsection*{c}

\begin{theorem}[blueprint="infogeometry:maxent:linearconstraint:c"]
  \inputleannode{InfoGeometry.MaxEnt.LinearConstraint.c}
\end{theorem}

\subsection*{casesOn}

\begin{theorem}[blueprint="infogeometry:maxent:linearconstraint:caseson"]
  \inputleannode{InfoGeometry.MaxEnt.LinearConstraint.casesOn}
\end{theorem}

\subsection*{ctorIdx}

\begin{theorem}[blueprint="infogeometry:maxent:linearconstraint:ctoridx"]
  \inputleannode{InfoGeometry.MaxEnt.LinearConstraint.ctorIdx}
\end{theorem}

\subsection*{f}

\begin{theorem}[blueprint="infogeometry:maxent:linearconstraint:f"]
  \inputleannode{InfoGeometry.MaxEnt.LinearConstraint.f}
\end{theorem}

\subsection*{measurable_f}

\begin{theorem}[blueprint="infogeometry:maxent:linearconstraint:measurable_f"]
  \inputleannode{InfoGeometry.MaxEnt.LinearConstraint.measurable_f}
\end{theorem}

\subsection*{mk}

\begin{theorem}[blueprint="infogeometry:maxent:linearconstraint:mk"]
  \inputleannode{InfoGeometry.MaxEnt.LinearConstraint.mk}
\end{theorem}

\subsection*{noConfusion}

\begin{theorem}[blueprint="infogeometry:maxent:linearconstraint:noconfusion"]
  \inputleannode{InfoGeometry.MaxEnt.LinearConstraint.noConfusion}
\end{theorem}

\subsection*{noConfusionType}

\begin{theorem}[blueprint="infogeometry:maxent:linearconstraint:noconfusiontype"]
  \inputleannode{InfoGeometry.MaxEnt.LinearConstraint.noConfusionType}
\end{theorem}

\subsection*{rec}

\begin{theorem}[blueprint="infogeometry:maxent:linearconstraint:rec"]
  \inputleannode{InfoGeometry.MaxEnt.LinearConstraint.rec}
\end{theorem}

\subsection*{recOn}

\begin{theorem}[blueprint="infogeometry:maxent:linearconstraint:recon"]
  \inputleannode{InfoGeometry.MaxEnt.LinearConstraint.recOn}
\end{theorem}

\section{InfoGeometry.MaxEnt.LinearConstraint.SatisfiesIntegrable}

\subsection*{integrable}

\begin{theorem}[blueprint="infogeometry:maxent:linearconstraint:satisfiesintegrable:integrable"]
  \inputleannode{InfoGeometry.MaxEnt.LinearConstraint.SatisfiesIntegrable.integrable}
\end{theorem}

\subsection*{satisfies}

\begin{theorem}[blueprint="infogeometry:maxent:linearconstraint:satisfiesintegrable:satisfies"]
  \inputleannode{InfoGeometry.MaxEnt.LinearConstraint.SatisfiesIntegrable.satisfies}
\end{theorem}

\section{InfoGeometry.MaxEnt.MaxEntProblem}

\subsection*{casesOn}

\begin{theorem}[blueprint="infogeometry:maxent:maxentproblem:caseson"]
  \inputleannode{InfoGeometry.MaxEnt.MaxEntProblem.casesOn}
\end{theorem}

\subsection*{ctorIdx}

\begin{theorem}[blueprint="infogeometry:maxent:maxentproblem:ctoridx"]
  \inputleannode{InfoGeometry.MaxEnt.MaxEntProblem.ctorIdx}
\end{theorem}

\subsection*{expectation}

\begin{theorem}[blueprint="infogeometry:maxent:maxentproblem:expectation"]
  \inputleannode{InfoGeometry.MaxEnt.MaxEntProblem.expectation}
\end{theorem}

\subsection*{mk}

\begin{theorem}[blueprint="infogeometry:maxent:maxentproblem:mk"]
  \inputleannode{InfoGeometry.MaxEnt.MaxEntProblem.mk}
\end{theorem}

\subsection*{noConfusion}

\begin{theorem}[blueprint="infogeometry:maxent:maxentproblem:noconfusion"]
  \inputleannode{InfoGeometry.MaxEnt.MaxEntProblem.noConfusion}
\end{theorem}

\subsection*{noConfusionType}

\begin{theorem}[blueprint="infogeometry:maxent:maxentproblem:noconfusiontype"]
  \inputleannode{InfoGeometry.MaxEnt.MaxEntProblem.noConfusionType}
\end{theorem}

\subsection*{nonneg}

\begin{theorem}[blueprint="infogeometry:maxent:maxentproblem:nonneg"]
  \inputleannode{InfoGeometry.MaxEnt.MaxEntProblem.nonneg}
\end{theorem}

\subsection*{norm}

\begin{theorem}[blueprint="infogeometry:maxent:maxentproblem:norm"]
  \inputleannode{InfoGeometry.MaxEnt.MaxEntProblem.norm}
\end{theorem}

\subsection*{p}

\begin{theorem}[blueprint="infogeometry:maxent:maxentproblem:p"]
  \inputleannode{InfoGeometry.MaxEnt.MaxEntProblem.p}
\end{theorem}

\subsection*{rec}

\begin{theorem}[blueprint="infogeometry:maxent:maxentproblem:rec"]
  \inputleannode{InfoGeometry.MaxEnt.MaxEntProblem.rec}
\end{theorem}

\subsection*{recOn}

\begin{theorem}[blueprint="infogeometry:maxent:maxentproblem:recon"]
  \inputleannode{InfoGeometry.MaxEnt.MaxEntProblem.recOn}
\end{theorem}

\section{InfoGeometry.MaxEnt.Optimality}

\noindent\textbf{Buildable}: yes

\noindent\emph{No exported declarations in the current declaration graph.}

\section{InfoGeometry.MaxEnt.ShannonEntropy}

\subsection*{eq_1}

\begin{theorem}[blueprint="infogeometry:maxent:shannonentropy:eq_1"]
  \inputleannode{InfoGeometry.MaxEnt.ShannonEntropy.eq_1}
\end{theorem}

\section{InfoGeometry.MaxEnt.dualObjective}

\subsection*{eq_1}

\begin{theorem}[blueprint="infogeometry:maxent:dualobjective:eq_1"]
  \inputleannode{InfoGeometry.MaxEnt.dualObjective.eq_1}
\end{theorem}

\section{InfoGeometry.MaxEnt.entropy}

\subsection*{eq_1}

\begin{theorem}[blueprint="infogeometry:maxent:entropy:eq_1"]
  \inputleannode{InfoGeometry.MaxEnt.entropy.eq_1}
\end{theorem}

\section{InfoGeometry.MaxEnt.gibbs}

\subsection*{eq_1}

\begin{theorem}[blueprint="infogeometry:maxent:gibbs:eq_1"]
  \inputleannode{InfoGeometry.MaxEnt.gibbs.eq_1}
\end{theorem}

\section{InfoGeometry.MaxEnt.gibbsDist}

\subsection*{eq_1}

\begin{theorem}[blueprint="infogeometry:maxent:gibbsdist:eq_1"]
  \inputleannode{InfoGeometry.MaxEnt.gibbsDist.eq_1}
\end{theorem}

\section{InfoGeometry.MaxEnt.gibbsExpectation}

\subsection*{eq_1}

\begin{theorem}[blueprint="infogeometry:maxent:gibbsexpectation:eq_1"]
  \inputleannode{InfoGeometry.MaxEnt.gibbsExpectation.eq_1}
\end{theorem}

\section{InfoGeometry.MaxEnt.gibbsExpectationWithPrior}

\subsection*{eq_1}

\begin{theorem}[blueprint="infogeometry:maxent:gibbsexpectationwithprior:eq_1"]
  \inputleannode{InfoGeometry.MaxEnt.gibbsExpectationWithPrior.eq_1}
\end{theorem}

\section{InfoGeometry.MaxEnt.goodProfiles}

\subsection*{eq_1}

\begin{theorem}[blueprint="infogeometry:maxent:goodprofiles:eq_1"]
  \inputleannode{InfoGeometry.MaxEnt.goodProfiles.eq_1}
\end{theorem}

\section{InfoGeometry.MaxEnt.multiplicity}

\subsection*{eq_1}

\begin{theorem}[blueprint="infogeometry:maxent:multiplicity:eq_1"]
  \inputleannode{InfoGeometry.MaxEnt.multiplicity.eq_1}
\end{theorem}

\section{InfoGeometry.MaxEnt.probDistOfSimplex}

\subsection*{eq_1}

\begin{theorem}[blueprint="infogeometry:maxent:probdistofsimplex:eq_1"]
  \inputleannode{InfoGeometry.MaxEnt.probDistOfSimplex.eq_1}
\end{theorem}

\section{InfoGeometry.MaxEnt.totalCount}

\subsection*{eq_1}

\begin{theorem}[blueprint="infogeometry:maxent:totalcount:eq_1"]
  \inputleannode{InfoGeometry.MaxEnt.totalCount.eq_1}
\end{theorem}

\section{InfoGeometry.MaxEnt.xlogx}

\subsection*{eq_1}

\begin{theorem}[blueprint="infogeometry:maxent:xlogx:eq_1"]
  \inputleannode{InfoGeometry.MaxEnt.xlogx.eq_1}
\end{theorem}

\chapter{OptimalTransport Modules}

\section{InfoGeometry.OptimalTransport}

\noindent\textbf{Buildable}: yes

\noindent\emph{No exported declarations in the current declaration graph.}

\chapter{Other Modules}

\section{InfoGeometryJordan}

\subsection*{«_aux_InfoGeometry_Jordan_Core___macroRules_InfoGeometryJordan_term_⊙__1»}

\begin{theorem}[blueprint="infogeometryjordan:«_aux_infogeometry_jordan_core___macrorules_infogeometryjordan_term_⊙__1»"]
  \inputleannode{InfoGeometryJordan.«_aux_InfoGeometry_Jordan_Core___macroRules_InfoGeometryJordan_term_⊙__1»}
\end{theorem}

\subsection*{«term_⊙_»}

\begin{theorem}[blueprint="infogeometryjordan:«term_⊙_»"]
  \inputleannode{InfoGeometryJordan.«term_⊙_»}
\end{theorem}

\chapter{PositiveMeasure Modules}

\section{InfoGeometry.PositiveMeasure}

\noindent\textbf{Buildable}: yes

\noindent\emph{No exported declarations in the current declaration graph.}

\chapter{Potential Modules}

\section{InfoGeometry.Potential}

\noindent\textbf{Buildable}: yes

\noindent\emph{No exported declarations in the current declaration graph.}

\section{InfoGeometry.Potential.LogPotential}

\noindent\textbf{Buildable}: yes

\noindent\emph{No exported declarations in the current declaration graph.}

\section{InfoGeometry.Potential.Thermo}

\noindent\textbf{Buildable}: yes

\noindent\emph{No exported declarations in the current declaration graph.}

\chapter{Prequantum Modules}

\section{InfoGeometry.Prequantum.Bundle}

\noindent\textbf{Buildable}: yes

\noindent\emph{No exported declarations in the current declaration graph.}

\section{InfoGeometry.Prequantum.Connection}

\noindent\textbf{Buildable}: yes

\noindent\emph{No exported declarations in the current declaration graph.}

\section{InfoGeometry.Prequantum.Quotient}

\noindent\textbf{Buildable}: yes

\noindent\emph{No exported declarations in the current declaration graph.}

\section{InfoGeometry.Prequantum.Scaling}

\noindent\textbf{Buildable}: yes

\noindent\emph{No exported declarations in the current declaration graph.}

\chapter{ProbabilityDist Modules}

\section{InfoGeometry.ProbabilityDist}

\subsection*{casesOn}

\begin{theorem}[blueprint="infogeometry:probabilitydist:caseson"]
  \inputleannode{InfoGeometry.ProbabilityDist.casesOn}
\end{theorem}

\subsection*{coe_prob}

\begin{theorem}[blueprint="infogeometry:probabilitydist:coe_prob"]
  \inputleannode{InfoGeometry.ProbabilityDist.coe_prob}
\end{theorem}

\subsection*{ctorIdx}

\begin{theorem}[blueprint="infogeometry:probabilitydist:ctoridx"]
  \inputleannode{InfoGeometry.ProbabilityDist.ctorIdx}
\end{theorem}

\subsection*{ext}

\begin{theorem}[blueprint="infogeometry:probabilitydist:ext"]
  \inputleannode{InfoGeometry.ProbabilityDist.ext}
\end{theorem}

\subsection*{ext_iff}

\begin{theorem}[blueprint="infogeometry:probabilitydist:ext_iff"]
  \inputleannode{InfoGeometry.ProbabilityDist.ext_iff}
\end{theorem}

\subsection*{mk}

\begin{theorem}[blueprint="infogeometry:probabilitydist:mk"]
  \inputleannode{InfoGeometry.ProbabilityDist.mk}
\end{theorem}

\subsection*{noConfusion}

\begin{theorem}[blueprint="infogeometry:probabilitydist:noconfusion"]
  \inputleannode{InfoGeometry.ProbabilityDist.noConfusion}
\end{theorem}

\subsection*{noConfusionType}

\begin{theorem}[blueprint="infogeometry:probabilitydist:noconfusiontype"]
  \inputleannode{InfoGeometry.ProbabilityDist.noConfusionType}
\end{theorem}

\subsection*{nonneg}

\begin{theorem}[blueprint="infogeometry:probabilitydist:nonneg"]
  \inputleannode{InfoGeometry.ProbabilityDist.nonneg}
\end{theorem}

\subsection*{prob}

\begin{theorem}[blueprint="infogeometry:probabilitydist:prob"]
  \inputleannode{InfoGeometry.ProbabilityDist.prob}
\end{theorem}

\subsection*{rec}

\begin{theorem}[blueprint="infogeometry:probabilitydist:rec"]
  \inputleannode{InfoGeometry.ProbabilityDist.rec}
\end{theorem}

\subsection*{recOn}

\begin{theorem}[blueprint="infogeometry:probabilitydist:recon"]
  \inputleannode{InfoGeometry.ProbabilityDist.recOn}
\end{theorem}

\subsection*{sum_one}

\begin{theorem}[blueprint="infogeometry:probabilitydist:sum_one"]
  \inputleannode{InfoGeometry.ProbabilityDist.sum_one}
\end{theorem}

\subsection*{toFinProb}

\begin{theorem}[blueprint="infogeometry:probabilitydist:tofinprob"]
  \inputleannode{InfoGeometry.ProbabilityDist.toFinProb}
\end{theorem}

\subsection*{toFinProb_toFun}

\begin{theorem}[blueprint="infogeometry:probabilitydist:tofinprob_tofun"]
  \inputleannode{InfoGeometry.ProbabilityDist.toFinProb_toFun}
\end{theorem}

\subsection*{toFinProb_toProbabilityDist}

\begin{theorem}[blueprint="infogeometry:probabilitydist:tofinprob_toprobabilitydist"]
  \inputleannode{InfoGeometry.ProbabilityDist.toFinProb_toProbabilityDist}
\end{theorem}

\chapter{Projective Modules}

\section{InfoGeometry.Projective}

\subsection*{GaugeInvariant}

\begin{theorem}[blueprint="infogeometry:projective:gaugeinvariant"]
  \inputleannode{InfoGeometry.Projective.GaugeInvariant}
\end{theorem}

\subsection*{PhysicalKinematics}

\begin{theorem}[blueprint="infogeometry:projective:physicalkinematics"]
  \inputleannode{InfoGeometry.Projective.PhysicalKinematics}
\end{theorem}

\subsection*{Unnormalized}

\begin{theorem}[blueprint="infogeometry:projective:unnormalized"]
  \inputleannode{InfoGeometry.Projective.Unnormalized}
\end{theorem}

\subsection*{descend}

\begin{theorem}[blueprint="infogeometry:projective:descend"]
  \inputleannode{InfoGeometry.Projective.descend}
\end{theorem}

\subsection*{descend_projectivize}

\begin{theorem}[blueprint="infogeometry:projective:descend_projectivize"]
  \inputleannode{InfoGeometry.Projective.descend_projectivize}
\end{theorem}

\subsection*{doubledProjectiveKinematics}

\begin{theorem}[blueprint="infogeometry:projective:doubledprojectivekinematics"]
  \inputleannode{InfoGeometry.Projective.doubledProjectiveKinematics}
\end{theorem}

\subsection*{liftGaugeInvariant}

\begin{theorem}[blueprint="infogeometry:projective:liftgaugeinvariant"]
  \inputleannode{InfoGeometry.Projective.liftGaugeInvariant}
\end{theorem}

\subsection*{liftGaugeInvariant_projectivize}

\begin{theorem}[blueprint="infogeometry:projective:liftgaugeinvariant_projectivize"]
  \inputleannode{InfoGeometry.Projective.liftGaugeInvariant_projectivize}
\end{theorem}

\subsection*{logSumInequality}

\begin{theorem}[blueprint="infogeometry:projective:logsuminequality"]
  \inputleannode{InfoGeometry.Projective.logSumInequality}
\end{theorem}

\subsection*{logSum_inequality}

\begin{theorem}[blueprint="infogeometry:projective:logsum_inequality"]
  \inputleannode{InfoGeometry.Projective.logSum_inequality}
\end{theorem}

\subsection*{occupancy}

\begin{theorem}[blueprint="infogeometry:projective:occupancy"]
  \inputleannode{InfoGeometry.Projective.occupancy}
\end{theorem}

\subsection*{occupancyOnProjective}

\begin{theorem}[blueprint="infogeometry:projective:occupancyonprojective"]
  \inputleannode{InfoGeometry.Projective.occupancyOnProjective}
\end{theorem}

\subsection*{occupancyOnProjective_binary}

\begin{theorem}[blueprint="infogeometry:projective:occupancyonprojective_binary"]
  \inputleannode{InfoGeometry.Projective.occupancyOnProjective_binary}
\end{theorem}

\subsection*{occupancyOnProjective_of_ne_zero}

\begin{theorem}[blueprint="infogeometry:projective:occupancyonprojective_of_ne_zero"]
  \inputleannode{InfoGeometry.Projective.occupancyOnProjective_of_ne_zero}
\end{theorem}

\subsection*{occupancyOnProjective_projectivize}

\begin{theorem}[blueprint="infogeometry:projective:occupancyonprojective_projectivize"]
  \inputleannode{InfoGeometry.Projective.occupancyOnProjective_projectivize}
\end{theorem}

\subsection*{occupancyOnProjective_zero}

\begin{theorem}[blueprint="infogeometry:projective:occupancyonprojective_zero"]
  \inputleannode{InfoGeometry.Projective.occupancyOnProjective_zero}
\end{theorem}

\subsection*{occupancy_eq_one_iff}

\begin{theorem}[blueprint="infogeometry:projective:occupancy_eq_one_iff"]
  \inputleannode{InfoGeometry.Projective.occupancy_eq_one_iff}
\end{theorem}

\subsection*{occupancy_eq_zero_iff}

\begin{theorem}[blueprint="infogeometry:projective:occupancy_eq_zero_iff"]
  \inputleannode{InfoGeometry.Projective.occupancy_eq_zero_iff}
\end{theorem}

\subsection*{occupancy_gaugeInvariant}

\begin{theorem}[blueprint="infogeometry:projective:occupancy_gaugeinvariant"]
  \inputleannode{InfoGeometry.Projective.occupancy_gaugeInvariant}
\end{theorem}

\subsection*{occupancy_of_ne_zero}

\begin{theorem}[blueprint="infogeometry:projective:occupancy_of_ne_zero"]
  \inputleannode{InfoGeometry.Projective.occupancy_of_ne_zero}
\end{theorem}

\subsection*{occupancy_zero}

\begin{theorem}[blueprint="infogeometry:projective:occupancy_zero"]
  \inputleannode{InfoGeometry.Projective.occupancy_zero}
\end{theorem}

\subsection*{ray}

\begin{theorem}[blueprint="infogeometry:projective:ray"]
  \inputleannode{InfoGeometry.Projective.ray}
\end{theorem}

\subsection*{ray_smul}

\begin{theorem}[blueprint="infogeometry:projective:ray_smul"]
  \inputleannode{InfoGeometry.Projective.ray_smul}
\end{theorem}

\subsection*{ray_smul₀}

\begin{theorem}[blueprint="infogeometry:projective:ray_smul₀"]
  \inputleannode{InfoGeometry.Projective.ray_smul₀}
\end{theorem}

\subsection*{smul_eq_zero_iff_of_ne_zero}

\begin{theorem}[blueprint="infogeometry:projective:smul_eq_zero_iff_of_ne_zero"]
  \inputleannode{InfoGeometry.Projective.smul_eq_zero_iff_of_ne_zero}
\end{theorem}

\section{InfoGeometry.Projective.Dynamics}

\noindent\textbf{Buildable}: yes

\noindent\emph{No exported declarations in the current declaration graph.}

\section{InfoGeometry.Projective.FaithfulKL}

\noindent\textbf{Buildable}: yes

\noindent\emph{No exported declarations in the current declaration graph.}

\section{InfoGeometry.Projective.GaugeInvariant}

\subsection*{compat}

\begin{theorem}[blueprint="infogeometry:projective:gaugeinvariant:compat"]
  \inputleannode{InfoGeometry.Projective.GaugeInvariant.compat}
\end{theorem}

\section{InfoGeometry.Projective.GaugeQuotient}

\noindent\textbf{Buildable}: yes

\noindent\emph{No exported declarations in the current declaration graph.}

\section{InfoGeometry.Projective.GaugeReduction}

\noindent\textbf{Buildable}: yes

\noindent\emph{No exported declarations in the current declaration graph.}

\section{InfoGeometry.Projective.LogSum}

\noindent\textbf{Buildable}: yes

\noindent\emph{No exported declarations in the current declaration graph.}

\section{InfoGeometry.Projective.LogSumIneq}

\noindent\textbf{Buildable}: yes

\noindent\emph{No exported declarations in the current declaration graph.}

\section{InfoGeometry.Projective.Normalize}

\noindent\textbf{Buildable}: yes

\noindent\emph{No exported declarations in the current declaration graph.}

\section{InfoGeometry.Projective.Null}

\noindent\textbf{Buildable}: yes

\noindent\emph{No exported declarations in the current declaration graph.}

\section{InfoGeometry.Projective.PhysicalKinematics}

\noindent\textbf{Buildable}: yes

\subsection*{Gauge}

\begin{theorem}[blueprint="infogeometry:projective:physicalkinematics:gauge"]
  \inputleannode{InfoGeometry.Projective.PhysicalKinematics.Gauge}
\end{theorem}

\subsection*{State}

\begin{theorem}[blueprint="infogeometry:projective:physicalkinematics:state"]
  \inputleannode{InfoGeometry.Projective.PhysicalKinematics.State}
\end{theorem}

\subsection*{casesOn}

\begin{theorem}[blueprint="infogeometry:projective:physicalkinematics:caseson"]
  \inputleannode{InfoGeometry.Projective.PhysicalKinematics.casesOn}
\end{theorem}

\subsection*{ctorIdx}

\begin{theorem}[blueprint="infogeometry:projective:physicalkinematics:ctoridx"]
  \inputleannode{InfoGeometry.Projective.PhysicalKinematics.ctorIdx}
\end{theorem}

\subsection*{instAct}

\begin{theorem}[blueprint="infogeometry:projective:physicalkinematics:instact"]
  \inputleannode{InfoGeometry.Projective.PhysicalKinematics.instAct}
\end{theorem}

\subsection*{instMonoid}

\begin{theorem}[blueprint="infogeometry:projective:physicalkinematics:instmonoid"]
  \inputleannode{InfoGeometry.Projective.PhysicalKinematics.instMonoid}
\end{theorem}

\subsection*{mk}

\begin{theorem}[blueprint="infogeometry:projective:physicalkinematics:mk"]
  \inputleannode{InfoGeometry.Projective.PhysicalKinematics.mk}
\end{theorem}

\subsection*{noConfusion}

\begin{theorem}[blueprint="infogeometry:projective:physicalkinematics:noconfusion"]
  \inputleannode{InfoGeometry.Projective.PhysicalKinematics.noConfusion}
\end{theorem}

\subsection*{noConfusionType}

\begin{theorem}[blueprint="infogeometry:projective:physicalkinematics:noconfusiontype"]
  \inputleannode{InfoGeometry.Projective.PhysicalKinematics.noConfusionType}
\end{theorem}

\subsection*{ray}

\begin{theorem}[blueprint="infogeometry:projective:physicalkinematics:ray"]
  \inputleannode{InfoGeometry.Projective.PhysicalKinematics.ray}
\end{theorem}

\subsection*{ray_gauge}

\begin{theorem}[blueprint="infogeometry:projective:physicalkinematics:ray_gauge"]
  \inputleannode{InfoGeometry.Projective.PhysicalKinematics.ray_gauge}
\end{theorem}

\subsection*{rec}

\begin{theorem}[blueprint="infogeometry:projective:physicalkinematics:rec"]
  \inputleannode{InfoGeometry.Projective.PhysicalKinematics.rec}
\end{theorem}

\subsection*{recOn}

\begin{theorem}[blueprint="infogeometry:projective:physicalkinematics:recon"]
  \inputleannode{InfoGeometry.Projective.PhysicalKinematics.recOn}
\end{theorem}

\section{InfoGeometry.Projective.Projective}

\noindent\textbf{Buildable}: yes

\noindent\emph{No exported declarations in the current declaration graph.}

\section{InfoGeometry.Projective.ProjectiveMap}

\noindent\textbf{Buildable}: yes

\noindent\emph{No exported declarations in the current declaration graph.}

\section{InfoGeometry.Projective.Rays}

\noindent\textbf{Buildable}: yes

\noindent\emph{No exported declarations in the current declaration graph.}

\section{InfoGeometry.Projective.occupancy}

\subsection*{eq_1}

\begin{theorem}[blueprint="infogeometry:projective:occupancy:eq_1"]
  \inputleannode{InfoGeometry.Projective.occupancy.eq_1}
\end{theorem}

\chapter{Quantum Modules}

\section{InfoGeometry.Quantum}

\subsection*{annihilationOp}

\begin{theorem}[blueprint="infogeometry:quantum:annihilationop"]
  \inputleannode{InfoGeometry.Quantum.annihilationOp}
\end{theorem}

\subsection*{annihilation_eq_minus_projector}

\begin{theorem}[blueprint="infogeometry:quantum:annihilation_eq_minus_projector"]
  \inputleannode{InfoGeometry.Quantum.annihilation_eq_minus_projector}
\end{theorem}

\subsection*{annihilation_kills_vacuum_vector}

\begin{theorem}[blueprint="infogeometry:quantum:annihilation_kills_vacuum_vector"]
  \inputleannode{InfoGeometry.Quantum.annihilation_kills_vacuum_vector}
\end{theorem}

\subsection*{bayesianAddData}

\begin{theorem}[blueprint="infogeometry:quantum:bayesianadddata"]
  \inputleannode{InfoGeometry.Quantum.bayesianAddData}
\end{theorem}

\subsection*{bayesianAddData_zero}

\begin{theorem}[blueprint="infogeometry:quantum:bayesianadddata_zero"]
  \inputleannode{InfoGeometry.Quantum.bayesianAddData_zero}
\end{theorem}

\subsection*{bayesianUpdate}

\begin{theorem}[blueprint="infogeometry:quantum:bayesianupdate"]
  \inputleannode{InfoGeometry.Quantum.bayesianUpdate}
\end{theorem}

\subsection*{bayesian_update_preserves_data_independence}

\begin{theorem}[blueprint="infogeometry:quantum:bayesian_update_preserves_data_independence"]
  \inputleannode{InfoGeometry.Quantum.bayesian_update_preserves_data_independence}
\end{theorem}

\subsection*{cliffordGenE}

\begin{theorem}[blueprint="infogeometry:quantum:cliffordgene"]
  \inputleannode{InfoGeometry.Quantum.cliffordGenE}
\end{theorem}

\subsection*{cliffordGenF}

\begin{theorem}[blueprint="infogeometry:quantum:cliffordgenf"]
  \inputleannode{InfoGeometry.Quantum.cliffordGenF}
\end{theorem}

\subsection*{commutator}

\begin{theorem}[blueprint="infogeometry:quantum:commutator"]
  \inputleannode{InfoGeometry.Quantum.commutator}
\end{theorem}

\subsection*{commutator_J_epsilon_eq_two_I}

\begin{theorem}[blueprint="infogeometry:quantum:commutator_j_epsilon_eq_two_i"]
  \inputleannode{InfoGeometry.Quantum.commutator_J_epsilon_eq_two_I}
\end{theorem}

\subsection*{creationOp}

\begin{theorem}[blueprint="infogeometry:quantum:creationop"]
  \inputleannode{InfoGeometry.Quantum.creationOp}
\end{theorem}

\subsection*{creation_add_annihilation}

\begin{theorem}[blueprint="infogeometry:quantum:creation_add_annihilation"]
  \inputleannode{InfoGeometry.Quantum.creation_add_annihilation}
\end{theorem}

\subsection*{creation_annihilation_orthogonal}

\begin{theorem}[blueprint="infogeometry:quantum:creation_annihilation_orthogonal"]
  \inputleannode{InfoGeometry.Quantum.creation_annihilation_orthogonal}
\end{theorem}

\subsection*{creation_eq_plus_projector}

\begin{theorem}[blueprint="infogeometry:quantum:creation_eq_plus_projector"]
  \inputleannode{InfoGeometry.Quantum.creation_eq_plus_projector}
\end{theorem}

\subsection*{dataPart}

\begin{theorem}[blueprint="infogeometry:quantum:datapart"]
  \inputleannode{InfoGeometry.Quantum.dataPart}
\end{theorem}

\subsection*{data_model_decomposition}

\begin{theorem}[blueprint="infogeometry:quantum:data_model_decomposition"]
  \inputleannode{InfoGeometry.Quantum.data_model_decomposition}
\end{theorem}

\subsection*{inducedSymplecticForm}

\begin{theorem}[blueprint="infogeometry:quantum:inducedsymplecticform"]
  \inputleannode{InfoGeometry.Quantum.inducedSymplecticForm}
\end{theorem}

\subsection*{inducedSymplecticForm_eq_complex_pairing}

\begin{theorem}[blueprint="infogeometry:quantum:inducedsymplecticform_eq_complex_pairing"]
  \inputleannode{InfoGeometry.Quantum.inducedSymplecticForm_eq_complex_pairing}
\end{theorem}

\subsection*{modelPart}

\begin{theorem}[blueprint="infogeometry:quantum:modelpart"]
  \inputleannode{InfoGeometry.Quantum.modelPart}
\end{theorem}

\subsection*{vacuum_is_zero_ray}

\begin{theorem}[blueprint="infogeometry:quantum:vacuum_is_zero_ray"]
  \inputleannode{InfoGeometry.Quantum.vacuum_is_zero_ray}
\end{theorem}

\section{InfoGeometry.Quantum.Fock}

\noindent\textbf{Buildable}: yes

\noindent\emph{No exported declarations in the current declaration graph.}

\section{InfoGeometry.Quantum.annihilationOp}

\subsection*{eq_1}

\begin{theorem}[blueprint="infogeometry:quantum:annihilationop:eq_1"]
  \inputleannode{InfoGeometry.Quantum.annihilationOp.eq_1}
\end{theorem}

\section{InfoGeometry.Quantum.bayesianAddData}

\subsection*{eq_1}

\begin{theorem}[blueprint="infogeometry:quantum:bayesianadddata:eq_1"]
  \inputleannode{InfoGeometry.Quantum.bayesianAddData.eq_1}
\end{theorem}

\section{InfoGeometry.Quantum.commutator}

\subsection*{eq_1}

\begin{theorem}[blueprint="infogeometry:quantum:commutator:eq_1"]
  \inputleannode{InfoGeometry.Quantum.commutator.eq_1}
\end{theorem}

\section{InfoGeometry.Quantum.creationOp}

\subsection*{eq_1}

\begin{theorem}[blueprint="infogeometry:quantum:creationop:eq_1"]
  \inputleannode{InfoGeometry.Quantum.creationOp.eq_1}
\end{theorem}

\section{InfoGeometry.Quantum.dataPart}

\subsection*{eq_1}

\begin{theorem}[blueprint="infogeometry:quantum:datapart:eq_1"]
  \inputleannode{InfoGeometry.Quantum.dataPart.eq_1}
\end{theorem}

\section{InfoGeometry.Quantum.modelPart}

\subsection*{eq_1}

\begin{theorem}[blueprint="infogeometry:quantum:modelpart:eq_1"]
  \inputleannode{InfoGeometry.Quantum.modelPart.eq_1}
\end{theorem}

\chapter{RegularizedKL Modules}

\section{InfoGeometry.RegularizedKL}

\noindent\textbf{Buildable}: yes

\noindent\emph{No exported declarations in the current declaration graph.}

\chapter{RegularizedKLTest Modules}

\section{InfoGeometry.RegularizedKLTest}

\noindent\textbf{Buildable}: yes

\noindent\emph{No exported declarations in the current declaration graph.}

\chapter{Renyi Modules}

\section{InfoGeometry.Renyi}

\noindent\textbf{Buildable}: yes

\noindent\emph{No exported declarations in the current declaration graph.}

\chapter{RenyiD Modules}

\section{InfoGeometry.RenyiD}

\subsection*{eq_1}

\begin{theorem}[blueprint="infogeometry:renyid:eq_1"]
  \inputleannode{InfoGeometry.RenyiD.eq_1}
\end{theorem}

\chapter{RenyiDShifted Modules}

\section{InfoGeometry.RenyiDShifted}

\subsection*{eq_1}

\begin{theorem}[blueprint="infogeometry:renyidshifted:eq_1"]
  \inputleannode{InfoGeometry.RenyiDShifted.eq_1}
\end{theorem}

\chapter{SLT Modules}

\section{InfoGeometry.SLT}

\subsection*{condExpExcept}

\begin{theorem}[blueprint="infogeometry:slt:condexpexcept"]
  \inputleannode{InfoGeometry.SLT.condExpExcept}
\end{theorem}

\subsection*{condVarExcept}

\begin{theorem}[blueprint="infogeometry:slt:condvarexcept"]
  \inputleannode{InfoGeometry.SLT.condVarExcept}
\end{theorem}

\subsection*{measurableSpaceExcept}

\begin{theorem}[blueprint="infogeometry:slt:measurablespaceexcept"]
  \inputleannode{InfoGeometry.SLT.measurableSpaceExcept}
\end{theorem}

\subsection*{setIntegral_condVarExcept_eq}

\begin{theorem}[blueprint="infogeometry:slt:setintegral_condvarexcept_eq"]
  \inputleannode{InfoGeometry.SLT.setIntegral_condVarExcept_eq}
\end{theorem}

\subsection*{setIntegral_condVar_except_eq}

\begin{theorem}[blueprint="infogeometry:slt:setintegral_condvar_except_eq"]
  \inputleannode{InfoGeometry.SLT.setIntegral_condVar_except_eq}
\end{theorem}

\section{InfoGeometry.SLT.ConditionalExpectation}

\noindent\textbf{Buildable}: yes

\noindent\emph{No exported declarations in the current declaration graph.}

\section{InfoGeometry.SLT.condExpExcept}

\subsection*{eq_1}

\begin{theorem}[blueprint="infogeometry:slt:condexpexcept:eq_1"]
  \inputleannode{InfoGeometry.SLT.condExpExcept.eq_1}
\end{theorem}

\section{InfoGeometry.SLT.condVarExcept}

\subsection*{eq_1}

\begin{theorem}[blueprint="infogeometry:slt:condvarexcept:eq_1"]
  \inputleannode{InfoGeometry.SLT.condVarExcept.eq_1}
\end{theorem}

\chapter{StrictProbabilityDist Modules}

\section{InfoGeometry.StrictProbabilityDist}

\subsection*{casesOn}

\begin{theorem}[blueprint="infogeometry:strictprobabilitydist:caseson"]
  \inputleannode{InfoGeometry.StrictProbabilityDist.casesOn}
\end{theorem}

\subsection*{coe_prob}

\begin{theorem}[blueprint="infogeometry:strictprobabilitydist:coe_prob"]
  \inputleannode{InfoGeometry.StrictProbabilityDist.coe_prob}
\end{theorem}

\subsection*{ctorIdx}

\begin{theorem}[blueprint="infogeometry:strictprobabilitydist:ctoridx"]
  \inputleannode{InfoGeometry.StrictProbabilityDist.ctorIdx}
\end{theorem}

\subsection*{ext}

\begin{theorem}[blueprint="infogeometry:strictprobabilitydist:ext"]
  \inputleannode{InfoGeometry.StrictProbabilityDist.ext}
\end{theorem}

\subsection*{ext_iff}

\begin{theorem}[blueprint="infogeometry:strictprobabilitydist:ext_iff"]
  \inputleannode{InfoGeometry.StrictProbabilityDist.ext_iff}
\end{theorem}

\subsection*{instCoeProbabilityDist}

\begin{theorem}[blueprint="infogeometry:strictprobabilitydist:instcoeprobabilitydist"]
  \inputleannode{InfoGeometry.StrictProbabilityDist.instCoeProbabilityDist}
\end{theorem}

\subsection*{mk}

\begin{theorem}[blueprint="infogeometry:strictprobabilitydist:mk"]
  \inputleannode{InfoGeometry.StrictProbabilityDist.mk}
\end{theorem}

\subsection*{noConfusion}

\begin{theorem}[blueprint="infogeometry:strictprobabilitydist:noconfusion"]
  \inputleannode{InfoGeometry.StrictProbabilityDist.noConfusion}
\end{theorem}

\subsection*{noConfusionType}

\begin{theorem}[blueprint="infogeometry:strictprobabilitydist:noconfusiontype"]
  \inputleannode{InfoGeometry.StrictProbabilityDist.noConfusionType}
\end{theorem}

\subsection*{pos}

\begin{theorem}[blueprint="infogeometry:strictprobabilitydist:pos"]
  \inputleannode{InfoGeometry.StrictProbabilityDist.pos}
\end{theorem}

\subsection*{prob}

\begin{theorem}[blueprint="infogeometry:strictprobabilitydist:prob"]
  \inputleannode{InfoGeometry.StrictProbabilityDist.prob}
\end{theorem}

\subsection*{prob_pos}

\begin{theorem}[blueprint="infogeometry:strictprobabilitydist:prob_pos"]
  \inputleannode{InfoGeometry.StrictProbabilityDist.prob_pos}
\end{theorem}

\subsection*{rec}

\begin{theorem}[blueprint="infogeometry:strictprobabilitydist:rec"]
  \inputleannode{InfoGeometry.StrictProbabilityDist.rec}
\end{theorem}

\subsection*{recOn}

\begin{theorem}[blueprint="infogeometry:strictprobabilitydist:recon"]
  \inputleannode{InfoGeometry.StrictProbabilityDist.recOn}
\end{theorem}

\subsection*{sum_one}

\begin{theorem}[blueprint="infogeometry:strictprobabilitydist:sum_one"]
  \inputleannode{InfoGeometry.StrictProbabilityDist.sum_one}
\end{theorem}

\subsection*{toProbabilityDist}

\begin{theorem}[blueprint="infogeometry:strictprobabilitydist:toprobabilitydist"]
  \inputleannode{InfoGeometry.StrictProbabilityDist.toProbabilityDist}
\end{theorem}

\subsection*{toProbabilityDist_prob}

\begin{theorem}[blueprint="infogeometry:strictprobabilitydist:toprobabilitydist_prob"]
  \inputleannode{InfoGeometry.StrictProbabilityDist.toProbabilityDist_prob}
\end{theorem}

\chapter{SuperUnified Modules}

\section{InfoGeometry.SuperUnified}

\noindent\textbf{Buildable}: yes

\subsection*{End}

\begin{theorem}[blueprint="infogeometry:superunified:end"]
  \inputleannode{InfoGeometry.SuperUnified.End}
\end{theorem}

\subsection*{IsBosonic}

\begin{theorem}[blueprint="infogeometry:superunified:isbosonic"]
  \inputleannode{InfoGeometry.SuperUnified.IsBosonic}
\end{theorem}

\subsection*{IsFermionic}

\begin{theorem}[blueprint="infogeometry:superunified:isfermionic"]
  \inputleannode{InfoGeometry.SuperUnified.IsFermionic}
\end{theorem}

\subsection*{StructureGroupLieAlgebra}

\begin{theorem}[blueprint="infogeometry:superunified:structuregroupliealgebra"]
  \inputleannode{InfoGeometry.SuperUnified.StructureGroupLieAlgebra}
\end{theorem}

\subsection*{SuperKaehlerStructure}

\begin{theorem}[blueprint="infogeometry:superunified:superkaehlerstructure"]
  \inputleannode{InfoGeometry.SuperUnified.SuperKaehlerStructure}
\end{theorem}

\subsection*{SuperchargeSpace}

\begin{theorem}[blueprint="infogeometry:superunified:superchargespace"]
  \inputleannode{InfoGeometry.SuperUnified.SuperchargeSpace}
\end{theorem}

\subsection*{bracket_boson_boson_is_bosonic}

\begin{theorem}[blueprint="infogeometry:superunified:bracket_boson_boson_is_bosonic"]
  \inputleannode{InfoGeometry.SuperUnified.bracket_boson_boson_is_bosonic}
\end{theorem}

\subsection*{bracket_boson_fermion_is_fermionic}

\begin{theorem}[blueprint="infogeometry:superunified:bracket_boson_fermion_is_fermionic"]
  \inputleannode{InfoGeometry.SuperUnified.bracket_boson_fermion_is_fermionic}
\end{theorem}

\subsection*{bracket_fermion_fermion_is_bosonic}

\begin{theorem}[blueprint="infogeometry:superunified:bracket_fermion_fermion_is_bosonic"]
  \inputleannode{InfoGeometry.SuperUnified.bracket_fermion_fermion_is_bosonic}
\end{theorem}

\subsection*{clifford_decomposition}

\begin{theorem}[blueprint="infogeometry:superunified:clifford_decomposition"]
  \inputleannode{InfoGeometry.SuperUnified.clifford_decomposition}
\end{theorem}

\subsection*{epsilon_is_fermionic}

\begin{theorem}[blueprint="infogeometry:superunified:epsilon_is_fermionic"]
  \inputleannode{InfoGeometry.SuperUnified.epsilon_is_fermionic}
\end{theorem}

\subsection*{hamiltonian_is_bosonic}

\begin{theorem}[blueprint="infogeometry:superunified:hamiltonian_is_bosonic"]
  \inputleannode{InfoGeometry.SuperUnified.hamiltonian_is_bosonic}
\end{theorem}

\subsection*{jordanProduct}

\begin{theorem}[blueprint="infogeometry:superunified:jordanproduct"]
  \inputleannode{InfoGeometry.SuperUnified.jordanProduct}
\end{theorem}

\subsection*{lieBracket}

\begin{theorem}[blueprint="infogeometry:superunified:liebracket"]
  \inputleannode{InfoGeometry.SuperUnified.lieBracket}
\end{theorem}

\subsection*{metricOp}

\begin{theorem}[blueprint="infogeometry:superunified:metricop"]
  \inputleannode{InfoGeometry.SuperUnified.metricOp}
\end{theorem}

\subsection*{supercharge_squared_is_hamiltonian}

\begin{theorem}[blueprint="infogeometry:superunified:supercharge_squared_is_hamiltonian"]
  \inputleannode{InfoGeometry.SuperUnified.supercharge_squared_is_hamiltonian}
\end{theorem}

\subsection*{symplecticFormOp}

\begin{theorem}[blueprint="infogeometry:superunified:symplecticformop"]
  \inputleannode{InfoGeometry.SuperUnified.symplecticFormOp}
\end{theorem}

\subsection*{symplectic_is_complex_structure}

\begin{theorem}[blueprint="infogeometry:superunified:symplectic_is_complex_structure"]
  \inputleannode{InfoGeometry.SuperUnified.symplectic_is_complex_structure}
\end{theorem}

\section{InfoGeometry.SuperUnified.IsBosonic}

\subsection*{eq_1}

\begin{theorem}[blueprint="infogeometry:superunified:isbosonic:eq_1"]
  \inputleannode{InfoGeometry.SuperUnified.IsBosonic.eq_1}
\end{theorem}

\section{InfoGeometry.SuperUnified.IsFermionic}

\subsection*{eq_1}

\begin{theorem}[blueprint="infogeometry:superunified:isfermionic:eq_1"]
  \inputleannode{InfoGeometry.SuperUnified.IsFermionic.eq_1}
\end{theorem}

\section{InfoGeometry.SuperUnified.SuperKaehlerStructure}

\subsection*{J}

\begin{theorem}[blueprint="infogeometry:superunified:superkaehlerstructure:j"]
  \inputleannode{InfoGeometry.SuperUnified.SuperKaehlerStructure.J}
\end{theorem}

\subsection*{J_sq}

\begin{theorem}[blueprint="infogeometry:superunified:superkaehlerstructure:j_sq"]
  \inputleannode{InfoGeometry.SuperUnified.SuperKaehlerStructure.J_sq}
\end{theorem}

\subsection*{anticomm}

\begin{theorem}[blueprint="infogeometry:superunified:superkaehlerstructure:anticomm"]
  \inputleannode{InfoGeometry.SuperUnified.SuperKaehlerStructure.anticomm}
\end{theorem}

\subsection*{casesOn}

\begin{theorem}[blueprint="infogeometry:superunified:superkaehlerstructure:caseson"]
  \inputleannode{InfoGeometry.SuperUnified.SuperKaehlerStructure.casesOn}
\end{theorem}

\subsection*{ctorIdx}

\begin{theorem}[blueprint="infogeometry:superunified:superkaehlerstructure:ctoridx"]
  \inputleannode{InfoGeometry.SuperUnified.SuperKaehlerStructure.ctorIdx}
\end{theorem}

\subsection*{eps_sq}

\begin{theorem}[blueprint="infogeometry:superunified:superkaehlerstructure:eps_sq"]
  \inputleannode{InfoGeometry.SuperUnified.SuperKaehlerStructure.eps_sq}
\end{theorem}

\subsection*{epsilon}

\begin{theorem}[blueprint="infogeometry:superunified:superkaehlerstructure:epsilon"]
  \inputleannode{InfoGeometry.SuperUnified.SuperKaehlerStructure.epsilon}
\end{theorem}

\subsection*{mk}

\begin{theorem}[blueprint="infogeometry:superunified:superkaehlerstructure:mk"]
  \inputleannode{InfoGeometry.SuperUnified.SuperKaehlerStructure.mk}
\end{theorem}

\subsection*{noConfusion}

\begin{theorem}[blueprint="infogeometry:superunified:superkaehlerstructure:noconfusion"]
  \inputleannode{InfoGeometry.SuperUnified.SuperKaehlerStructure.noConfusion}
\end{theorem}

\subsection*{noConfusionType}

\begin{theorem}[blueprint="infogeometry:superunified:superkaehlerstructure:noconfusiontype"]
  \inputleannode{InfoGeometry.SuperUnified.SuperKaehlerStructure.noConfusionType}
\end{theorem}

\subsection*{rec}

\begin{theorem}[blueprint="infogeometry:superunified:superkaehlerstructure:rec"]
  \inputleannode{InfoGeometry.SuperUnified.SuperKaehlerStructure.rec}
\end{theorem}

\subsection*{recOn}

\begin{theorem}[blueprint="infogeometry:superunified:superkaehlerstructure:recon"]
  \inputleannode{InfoGeometry.SuperUnified.SuperKaehlerStructure.recOn}
\end{theorem}

\section{InfoGeometry.SuperUnified.jordanProduct}

\subsection*{eq_1}

\begin{theorem}[blueprint="infogeometry:superunified:jordanproduct:eq_1"]
  \inputleannode{InfoGeometry.SuperUnified.jordanProduct.eq_1}
\end{theorem}

\section{InfoGeometry.SuperUnified.lieBracket}

\subsection*{eq_1}

\begin{theorem}[blueprint="infogeometry:superunified:liebracket:eq_1"]
  \inputleannode{InfoGeometry.SuperUnified.lieBracket.eq_1}
\end{theorem}

\chapter{Thermal Modules}

\section{InfoGeometry.Thermal}

\subsection*{DiagonalObservable}

\begin{theorem}[blueprint="infogeometry:thermal:diagonalobservable"]
  \inputleannode{InfoGeometry.Thermal.DiagonalObservable}
\end{theorem}

\subsection*{Hamiltonian}

\begin{theorem}[blueprint="infogeometry:thermal:hamiltonian"]
  \inputleannode{InfoGeometry.Thermal.Hamiltonian}
\end{theorem}

\section{InfoGeometry.Thermal.DiagonalObservable}

\subsection*{add_apply}

\begin{theorem}[blueprint="infogeometry:thermal:diagonalobservable:add_apply"]
  \inputleannode{InfoGeometry.Thermal.DiagonalObservable.add_apply}
\end{theorem}

\subsection*{casesOn}

\begin{theorem}[blueprint="infogeometry:thermal:diagonalobservable:caseson"]
  \inputleannode{InfoGeometry.Thermal.DiagonalObservable.casesOn}
\end{theorem}

\subsection*{coeff}

\begin{theorem}[blueprint="infogeometry:thermal:diagonalobservable:coeff"]
  \inputleannode{InfoGeometry.Thermal.DiagonalObservable.coeff}
\end{theorem}

\subsection*{ctorIdx}

\begin{theorem}[blueprint="infogeometry:thermal:diagonalobservable:ctoridx"]
  \inputleannode{InfoGeometry.Thermal.DiagonalObservable.ctorIdx}
\end{theorem}

\subsection*{instAdd}

\begin{theorem}[blueprint="infogeometry:thermal:diagonalobservable:instadd"]
  \inputleannode{InfoGeometry.Thermal.DiagonalObservable.instAdd}
\end{theorem}

\subsection*{instCoeFunForallFinReal}

\begin{theorem}[blueprint="infogeometry:thermal:diagonalobservable:instcoefunforallfinreal"]
  \inputleannode{InfoGeometry.Thermal.DiagonalObservable.instCoeFunForallFinReal}
\end{theorem}

\subsection*{instMul}

\begin{theorem}[blueprint="infogeometry:thermal:diagonalobservable:instmul"]
  \inputleannode{InfoGeometry.Thermal.DiagonalObservable.instMul}
\end{theorem}

\subsection*{instNeg}

\begin{theorem}[blueprint="infogeometry:thermal:diagonalobservable:instneg"]
  \inputleannode{InfoGeometry.Thermal.DiagonalObservable.instNeg}
\end{theorem}

\subsection*{instOne}

\begin{theorem}[blueprint="infogeometry:thermal:diagonalobservable:instone"]
  \inputleannode{InfoGeometry.Thermal.DiagonalObservable.instOne}
\end{theorem}

\subsection*{instSMulReal}

\begin{theorem}[blueprint="infogeometry:thermal:diagonalobservable:instsmulreal"]
  \inputleannode{InfoGeometry.Thermal.DiagonalObservable.instSMulReal}
\end{theorem}

\subsection*{instSub}

\begin{theorem}[blueprint="infogeometry:thermal:diagonalobservable:instsub"]
  \inputleannode{InfoGeometry.Thermal.DiagonalObservable.instSub}
\end{theorem}

\subsection*{instZero}

\begin{theorem}[blueprint="infogeometry:thermal:diagonalobservable:instzero"]
  \inputleannode{InfoGeometry.Thermal.DiagonalObservable.instZero}
\end{theorem}

\subsection*{mk}

\begin{theorem}[blueprint="infogeometry:thermal:diagonalobservable:mk"]
  \inputleannode{InfoGeometry.Thermal.DiagonalObservable.mk}
\end{theorem}

\subsection*{mul_apply}

\begin{theorem}[blueprint="infogeometry:thermal:diagonalobservable:mul_apply"]
  \inputleannode{InfoGeometry.Thermal.DiagonalObservable.mul_apply}
\end{theorem}

\subsection*{neg_apply}

\begin{theorem}[blueprint="infogeometry:thermal:diagonalobservable:neg_apply"]
  \inputleannode{InfoGeometry.Thermal.DiagonalObservable.neg_apply}
\end{theorem}

\subsection*{noConfusion}

\begin{theorem}[blueprint="infogeometry:thermal:diagonalobservable:noconfusion"]
  \inputleannode{InfoGeometry.Thermal.DiagonalObservable.noConfusion}
\end{theorem}

\subsection*{noConfusionType}

\begin{theorem}[blueprint="infogeometry:thermal:diagonalobservable:noconfusiontype"]
  \inputleannode{InfoGeometry.Thermal.DiagonalObservable.noConfusionType}
\end{theorem}

\subsection*{one_apply}

\begin{theorem}[blueprint="infogeometry:thermal:diagonalobservable:one_apply"]
  \inputleannode{InfoGeometry.Thermal.DiagonalObservable.one_apply}
\end{theorem}

\subsection*{rec}

\begin{theorem}[blueprint="infogeometry:thermal:diagonalobservable:rec"]
  \inputleannode{InfoGeometry.Thermal.DiagonalObservable.rec}
\end{theorem}

\subsection*{recOn}

\begin{theorem}[blueprint="infogeometry:thermal:diagonalobservable:recon"]
  \inputleannode{InfoGeometry.Thermal.DiagonalObservable.recOn}
\end{theorem}

\subsection*{smul_apply}

\begin{theorem}[blueprint="infogeometry:thermal:diagonalobservable:smul_apply"]
  \inputleannode{InfoGeometry.Thermal.DiagonalObservable.smul_apply}
\end{theorem}

\subsection*{sub_apply}

\begin{theorem}[blueprint="infogeometry:thermal:diagonalobservable:sub_apply"]
  \inputleannode{InfoGeometry.Thermal.DiagonalObservable.sub_apply}
\end{theorem}

\subsection*{toMatrix}

\begin{theorem}[blueprint="infogeometry:thermal:diagonalobservable:tomatrix"]
  \inputleannode{InfoGeometry.Thermal.DiagonalObservable.toMatrix}
\end{theorem}

\subsection*{toMatrix_apply_diag}

\begin{theorem}[blueprint="infogeometry:thermal:diagonalobservable:tomatrix_apply_diag"]
  \inputleannode{InfoGeometry.Thermal.DiagonalObservable.toMatrix_apply_diag}
\end{theorem}

\subsection*{toMatrix_apply_offdiag}

\begin{theorem}[blueprint="infogeometry:thermal:diagonalobservable:tomatrix_apply_offdiag"]
  \inputleannode{InfoGeometry.Thermal.DiagonalObservable.toMatrix_apply_offdiag}
\end{theorem}

\subsection*{zero_apply}

\begin{theorem}[blueprint="infogeometry:thermal:diagonalobservable:zero_apply"]
  \inputleannode{InfoGeometry.Thermal.DiagonalObservable.zero_apply}
\end{theorem}

\section{InfoGeometry.Thermal.DiagonalObservable.toMatrix}

\subsection*{eq_1}

\begin{theorem}[blueprint="infogeometry:thermal:diagonalobservable:tomatrix:eq_1"]
  \inputleannode{InfoGeometry.Thermal.DiagonalObservable.toMatrix.eq_1}
\end{theorem}

\section{InfoGeometry.Thermal.FiniteMatrix}

\noindent\textbf{Buildable}: yes

\noindent\emph{No exported declarations in the current declaration graph.}

\section{InfoGeometry.Thermal.Hamiltonian}

\subsection*{SatisfiesDiagonalKMS}

\begin{theorem}[blueprint="infogeometry:thermal:hamiltonian:satisfiesdiagonalkms"]
  \inputleannode{InfoGeometry.Thermal.Hamiltonian.SatisfiesDiagonalKMS}
\end{theorem}

\subsection*{casesOn}

\begin{theorem}[blueprint="infogeometry:thermal:hamiltonian:caseson"]
  \inputleannode{InfoGeometry.Thermal.Hamiltonian.casesOn}
\end{theorem}

\subsection*{ctorIdx}

\begin{theorem}[blueprint="infogeometry:thermal:hamiltonian:ctoridx"]
  \inputleannode{InfoGeometry.Thermal.Hamiltonian.ctorIdx}
\end{theorem}

\subsection*{densityMatrix}

\begin{theorem}[blueprint="infogeometry:thermal:hamiltonian:densitymatrix"]
  \inputleannode{InfoGeometry.Thermal.Hamiltonian.densityMatrix}
\end{theorem}

\subsection*{densityMatrix_apply_diag}

\begin{theorem}[blueprint="infogeometry:thermal:hamiltonian:densitymatrix_apply_diag"]
  \inputleannode{InfoGeometry.Thermal.Hamiltonian.densityMatrix_apply_diag}
\end{theorem}

\subsection*{densityMatrix_apply_offdiag}

\begin{theorem}[blueprint="infogeometry:thermal:hamiltonian:densitymatrix_apply_offdiag"]
  \inputleannode{InfoGeometry.Thermal.Hamiltonian.densityMatrix_apply_offdiag}
\end{theorem}

\subsection*{energy}

\begin{theorem}[blueprint="infogeometry:thermal:hamiltonian:energy"]
  \inputleannode{InfoGeometry.Thermal.Hamiltonian.energy}
\end{theorem}

\subsection*{energyObs}

\begin{theorem}[blueprint="infogeometry:thermal:hamiltonian:energyobs"]
  \inputleannode{InfoGeometry.Thermal.Hamiltonian.energyObs}
\end{theorem}

\subsection*{gibbsWeight}

\begin{theorem}[blueprint="infogeometry:thermal:hamiltonian:gibbsweight"]
  \inputleannode{InfoGeometry.Thermal.Hamiltonian.gibbsWeight}
\end{theorem}

\subsection*{gibbsWeight_eq_exp_logDensity}

\begin{theorem}[blueprint="infogeometry:thermal:hamiltonian:gibbsweight_eq_exp_logdensity"]
  \inputleannode{InfoGeometry.Thermal.Hamiltonian.gibbsWeight_eq_exp_logDensity}
\end{theorem}

\subsection*{gibbsWeight_nonneg}

\begin{theorem}[blueprint="infogeometry:thermal:hamiltonian:gibbsweight_nonneg"]
  \inputleannode{InfoGeometry.Thermal.Hamiltonian.gibbsWeight_nonneg}
\end{theorem}

\subsection*{gibbsWeight_pos}

\begin{theorem}[blueprint="infogeometry:thermal:hamiltonian:gibbsweight_pos"]
  \inputleannode{InfoGeometry.Thermal.Hamiltonian.gibbsWeight_pos}
\end{theorem}

\subsection*{gibbsWeight_sum_one}

\begin{theorem}[blueprint="infogeometry:thermal:hamiltonian:gibbsweight_sum_one"]
  \inputleannode{InfoGeometry.Thermal.Hamiltonian.gibbsWeight_sum_one}
\end{theorem}

\subsection*{internalEnergy}

\begin{theorem}[blueprint="infogeometry:thermal:hamiltonian:internalenergy"]
  \inputleannode{InfoGeometry.Thermal.Hamiltonian.internalEnergy}
\end{theorem}

\subsection*{internalEnergy_eq_gibbsExpectation}

\begin{theorem}[blueprint="infogeometry:thermal:hamiltonian:internalenergy_eq_gibbsexpectation"]
  \inputleannode{InfoGeometry.Thermal.Hamiltonian.internalEnergy_eq_gibbsExpectation}
\end{theorem}

\subsection*{logDensityCoeff}

\begin{theorem}[blueprint="infogeometry:thermal:hamiltonian:logdensitycoeff"]
  \inputleannode{InfoGeometry.Thermal.Hamiltonian.logDensityCoeff}
\end{theorem}

\subsection*{logDensityObs}

\begin{theorem}[blueprint="infogeometry:thermal:hamiltonian:logdensityobs"]
  \inputleannode{InfoGeometry.Thermal.Hamiltonian.logDensityObs}
\end{theorem}

\subsection*{logPartition}

\begin{theorem}[blueprint="infogeometry:thermal:hamiltonian:logpartition"]
  \inputleannode{InfoGeometry.Thermal.Hamiltonian.logPartition}
\end{theorem}

\subsection*{mk}

\begin{theorem}[blueprint="infogeometry:thermal:hamiltonian:mk"]
  \inputleannode{InfoGeometry.Thermal.Hamiltonian.mk}
\end{theorem}

\subsection*{modularShift}

\begin{theorem}[blueprint="infogeometry:thermal:hamiltonian:modularshift"]
  \inputleannode{InfoGeometry.Thermal.Hamiltonian.modularShift}
\end{theorem}

\subsection*{modularShift_add}

\begin{theorem}[blueprint="infogeometry:thermal:hamiltonian:modularshift_add"]
  \inputleannode{InfoGeometry.Thermal.Hamiltonian.modularShift_add}
\end{theorem}

\subsection*{modularShift_apply}

\begin{theorem}[blueprint="infogeometry:thermal:hamiltonian:modularshift_apply"]
  \inputleannode{InfoGeometry.Thermal.Hamiltonian.modularShift_apply}
\end{theorem}

\subsection*{modularShift_mul}

\begin{theorem}[blueprint="infogeometry:thermal:hamiltonian:modularshift_mul"]
  \inputleannode{InfoGeometry.Thermal.Hamiltonian.modularShift_mul}
\end{theorem}

\subsection*{modularShift_zero}

\begin{theorem}[blueprint="infogeometry:thermal:hamiltonian:modularshift_zero"]
  \inputleannode{InfoGeometry.Thermal.Hamiltonian.modularShift_zero}
\end{theorem}

\subsection*{noConfusion}

\begin{theorem}[blueprint="infogeometry:thermal:hamiltonian:noconfusion"]
  \inputleannode{InfoGeometry.Thermal.Hamiltonian.noConfusion}
\end{theorem}

\subsection*{noConfusionType}

\begin{theorem}[blueprint="infogeometry:thermal:hamiltonian:noconfusiontype"]
  \inputleannode{InfoGeometry.Thermal.Hamiltonian.noConfusionType}
\end{theorem}

\subsection*{partition}

\begin{theorem}[blueprint="infogeometry:thermal:hamiltonian:partition"]
  \inputleannode{InfoGeometry.Thermal.Hamiltonian.partition}
\end{theorem}

\subsection*{partition_ne_zero}

\begin{theorem}[blueprint="infogeometry:thermal:hamiltonian:partition_ne_zero"]
  \inputleannode{InfoGeometry.Thermal.Hamiltonian.partition_ne_zero}
\end{theorem}

\subsection*{partition_pos}

\begin{theorem}[blueprint="infogeometry:thermal:hamiltonian:partition_pos"]
  \inputleannode{InfoGeometry.Thermal.Hamiltonian.partition_pos}
\end{theorem}

\subsection*{rec}

\begin{theorem}[blueprint="infogeometry:thermal:hamiltonian:rec"]
  \inputleannode{InfoGeometry.Thermal.Hamiltonian.rec}
\end{theorem}

\subsection*{recOn}

\begin{theorem}[blueprint="infogeometry:thermal:hamiltonian:recon"]
  \inputleannode{InfoGeometry.Thermal.Hamiltonian.recOn}
\end{theorem}

\subsection*{satisfiesDiagonalKMS}

\begin{theorem}[blueprint="infogeometry:thermal:hamiltonian:satisfiesdiagonalkms"]
  \inputleannode{InfoGeometry.Thermal.Hamiltonian.satisfiesDiagonalKMS}
\end{theorem}

\subsection*{thermalState}

\begin{theorem}[blueprint="infogeometry:thermal:hamiltonian:thermalstate"]
  \inputleannode{InfoGeometry.Thermal.Hamiltonian.thermalState}
\end{theorem}

\subsection*{thermalState_eq_trace_density_mul_diag}

\begin{theorem}[blueprint="infogeometry:thermal:hamiltonian:thermalstate_eq_trace_density_mul_diag"]
  \inputleannode{InfoGeometry.Thermal.Hamiltonian.thermalState_eq_trace_density_mul_diag}
\end{theorem}

\subsection*{thermalState_modularShift_invariant}

\begin{theorem}[blueprint="infogeometry:thermal:hamiltonian:thermalstate_modularshift_invariant"]
  \inputleannode{InfoGeometry.Thermal.Hamiltonian.thermalState_modularShift_invariant}
\end{theorem}

\subsection*{thermalState_nonneg}

\begin{theorem}[blueprint="infogeometry:thermal:hamiltonian:thermalstate_nonneg"]
  \inputleannode{InfoGeometry.Thermal.Hamiltonian.thermalState_nonneg}
\end{theorem}

\subsection*{thermalState_one}

\begin{theorem}[blueprint="infogeometry:thermal:hamiltonian:thermalstate_one"]
  \inputleannode{InfoGeometry.Thermal.Hamiltonian.thermalState_one}
\end{theorem}

\subsection*{toMatrix}

\begin{theorem}[blueprint="infogeometry:thermal:hamiltonian:tomatrix"]
  \inputleannode{InfoGeometry.Thermal.Hamiltonian.toMatrix}
\end{theorem}

\subsection*{toMatrix_apply_diag}

\begin{theorem}[blueprint="infogeometry:thermal:hamiltonian:tomatrix_apply_diag"]
  \inputleannode{InfoGeometry.Thermal.Hamiltonian.toMatrix_apply_diag}
\end{theorem}

\subsection*{trace}

\begin{theorem}[blueprint="infogeometry:thermal:hamiltonian:trace"]
  \inputleannode{InfoGeometry.Thermal.Hamiltonian.trace}
\end{theorem}

\subsection*{trace_densityMatrix}

\begin{theorem}[blueprint="infogeometry:thermal:hamiltonian:trace_densitymatrix"]
  \inputleannode{InfoGeometry.Thermal.Hamiltonian.trace_densityMatrix}
\end{theorem}

\section{InfoGeometry.Thermal.Hamiltonian.densityMatrix}

\subsection*{eq_1}

\begin{theorem}[blueprint="infogeometry:thermal:hamiltonian:densitymatrix:eq_1"]
  \inputleannode{InfoGeometry.Thermal.Hamiltonian.densityMatrix.eq_1}
\end{theorem}

\section{InfoGeometry.Thermal.Hamiltonian.gibbsWeight}

\subsection*{eq_1}

\begin{theorem}[blueprint="infogeometry:thermal:hamiltonian:gibbsweight:eq_1"]
  \inputleannode{InfoGeometry.Thermal.Hamiltonian.gibbsWeight.eq_1}
\end{theorem}

\section{InfoGeometry.Thermal.Hamiltonian.modularShift}

\subsection*{eq_1}

\begin{theorem}[blueprint="infogeometry:thermal:hamiltonian:modularshift:eq_1"]
  \inputleannode{InfoGeometry.Thermal.Hamiltonian.modularShift.eq_1}
\end{theorem}

\section{InfoGeometry.Thermal.Hamiltonian.partition}

\subsection*{eq_1}

\begin{theorem}[blueprint="infogeometry:thermal:hamiltonian:partition:eq_1"]
  \inputleannode{InfoGeometry.Thermal.Hamiltonian.partition.eq_1}
\end{theorem}

\section{InfoGeometry.Thermal.Hamiltonian.thermalState}

\subsection*{eq_1}

\begin{theorem}[blueprint="infogeometry:thermal:hamiltonian:thermalstate:eq_1"]
  \inputleannode{InfoGeometry.Thermal.Hamiltonian.thermalState.eq_1}
\end{theorem}

\section{InfoGeometry.Thermal.Hamiltonian.toMatrix}

\subsection*{eq_1}

\begin{theorem}[blueprint="infogeometry:thermal:hamiltonian:tomatrix:eq_1"]
  \inputleannode{InfoGeometry.Thermal.Hamiltonian.toMatrix.eq_1}
\end{theorem}

\section{InfoGeometry.Thermal.Hamiltonian.trace}

\subsection*{eq_1}

\begin{theorem}[blueprint="infogeometry:thermal:hamiltonian:trace:eq_1"]
  \inputleannode{InfoGeometry.Thermal.Hamiltonian.trace.eq_1}
\end{theorem}

\chapter{Thermo Modules}

\section{InfoGeometry.Thermo}

\subsection*{KL_param_eq_bregman_energy}

\begin{theorem}[blueprint="infogeometry:thermo:kl_param_eq_bregman_energy"]
  \inputleannode{InfoGeometry.Thermo.KL_param_eq_bregman_energy}
\end{theorem}

\subsection*{Op}

\begin{theorem}[blueprint="infogeometry:thermo:op"]
  \inputleannode{InfoGeometry.Thermo.Op}
\end{theorem}

\subsection*{ThermalModel}

\begin{theorem}[blueprint="infogeometry:thermo:thermalmodel"]
  \inputleannode{InfoGeometry.Thermo.ThermalModel}
\end{theorem}

\subsection*{Z}

\begin{theorem}[blueprint="infogeometry:thermo:z"]
  \inputleannode{InfoGeometry.Thermo.Z}
\end{theorem}

\subsection*{Z_ne_zero}

\begin{theorem}[blueprint="infogeometry:thermo:z_ne_zero"]
  \inputleannode{InfoGeometry.Thermo.Z_ne_zero}
\end{theorem}

\subsection*{Z_pos}

\begin{theorem}[blueprint="infogeometry:thermo:z_pos"]
  \inputleannode{InfoGeometry.Thermo.Z_pos}
\end{theorem}

\subsection*{bayesianFreeEnergyFromBregman}

\begin{theorem}[blueprint="infogeometry:thermo:bayesianfreeenergyfrombregman"]
  \inputleannode{InfoGeometry.Thermo.bayesianFreeEnergyFromBregman}
\end{theorem}

\subsection*{bayesianFreeEnergyFromBregman_eq}

\begin{theorem}[blueprint="infogeometry:thermo:bayesianfreeenergyfrombregman_eq"]
  \inputleannode{InfoGeometry.Thermo.bayesianFreeEnergyFromBregman_eq}
\end{theorem}

\subsection*{bayesianFreeEnergyFromBregman_eq_internal_sub_scale_entropy}

\begin{theorem}[blueprint="infogeometry:thermo:bayesianfreeenergyfrombregman_eq_internal_sub_scale_entropy"]
  \inputleannode{InfoGeometry.Thermo.bayesianFreeEnergyFromBregman_eq_internal_sub_scale_entropy}
\end{theorem}

\subsection*{energyFromBregman}

\begin{theorem}[blueprint="infogeometry:thermo:energyfrombregman"]
  \inputleannode{InfoGeometry.Thermo.energyFromBregman}
\end{theorem}

\subsection*{energyFromDivergence}

\begin{theorem}[blueprint="infogeometry:thermo:energyfromdivergence"]
  \inputleannode{InfoGeometry.Thermo.energyFromDivergence}
\end{theorem}

\subsection*{energyFromLogDet}

\begin{theorem}[blueprint="infogeometry:thermo:energyfromlogdet"]
  \inputleannode{InfoGeometry.Thermo.energyFromLogDet}
\end{theorem}

\subsection*{entropicTransportCost_eq_bregman_gap}

\begin{theorem}[blueprint="infogeometry:thermo:entropictransportcost_eq_bregman_gap"]
  \inputleannode{InfoGeometry.Thermo.entropicTransportCost_eq_bregman_gap}
\end{theorem}

\subsection*{entropicTransportEnergyFromBregman}

\begin{theorem}[blueprint="infogeometry:thermo:entropictransportenergyfrombregman"]
  \inputleannode{InfoGeometry.Thermo.entropicTransportEnergyFromBregman}
\end{theorem}

\subsection*{entropicTransportEnergyFromBregman_eq}

\begin{theorem}[blueprint="infogeometry:thermo:entropictransportenergyfrombregman_eq"]
  \inputleannode{InfoGeometry.Thermo.entropicTransportEnergyFromBregman_eq}
\end{theorem}

\subsection*{entropicTransportPlanFromBregman}

\begin{theorem}[blueprint="infogeometry:thermo:entropictransportplanfrombregman"]
  \inputleannode{InfoGeometry.Thermo.entropicTransportPlanFromBregman}
\end{theorem}

\subsection*{entropicTransportPlanFromBregman_eq}

\begin{theorem}[blueprint="infogeometry:thermo:entropictransportplanfrombregman_eq"]
  \inputleannode{InfoGeometry.Thermo.entropicTransportPlanFromBregman_eq}
\end{theorem}

\subsection*{epsilon_mul_shannonEntropy}

\begin{theorem}[blueprint="infogeometry:thermo:epsilon_mul_shannonentropy"]
  \inputleannode{InfoGeometry.Thermo.epsilon_mul_shannonEntropy}
\end{theorem}

\subsection*{expectation}

\begin{theorem}[blueprint="infogeometry:thermo:expectation"]
  \inputleannode{InfoGeometry.Thermo.expectation}
\end{theorem}

\subsection*{freeEnergy}

\begin{theorem}[blueprint="infogeometry:thermo:freeenergy"]
  \inputleannode{InfoGeometry.Thermo.freeEnergy}
\end{theorem}

\subsection*{freeEnergyDivergence}

\begin{theorem}[blueprint="infogeometry:thermo:freeenergydivergence"]
  \inputleannode{InfoGeometry.Thermo.freeEnergyDivergence}
\end{theorem}

\subsection*{freeEnergyDivergence_def}

\begin{theorem}[blueprint="infogeometry:thermo:freeenergydivergence_def"]
  \inputleannode{InfoGeometry.Thermo.freeEnergyDivergence_def}
\end{theorem}

\subsection*{freeEnergyDivergence_eq_freeEnergyFromBregman}

\begin{theorem}[blueprint="infogeometry:thermo:freeenergydivergence_eq_freeenergyfrombregman"]
  \inputleannode{InfoGeometry.Thermo.freeEnergyDivergence_eq_freeEnergyFromBregman}
\end{theorem}

\subsection*{freeEnergyDivergence_eq_internal_sub_scale_entropy}

\begin{theorem}[blueprint="infogeometry:thermo:freeenergydivergence_eq_internal_sub_scale_entropy"]
  \inputleannode{InfoGeometry.Thermo.freeEnergyDivergence_eq_internal_sub_scale_entropy}
\end{theorem}

\subsection*{freeEnergyFromBregman}

\begin{theorem}[blueprint="infogeometry:thermo:freeenergyfrombregman"]
  \inputleannode{InfoGeometry.Thermo.freeEnergyFromBregman}
\end{theorem}

\subsection*{freeEnergyFromBregman_eq_internal_sub_scale_entropy}

\begin{theorem}[blueprint="infogeometry:thermo:freeenergyfrombregman_eq_internal_sub_scale_entropy"]
  \inputleannode{InfoGeometry.Thermo.freeEnergyFromBregman_eq_internal_sub_scale_entropy}
\end{theorem}

\subsection*{freeEnergyFromLogDet}

\begin{theorem}[blueprint="infogeometry:thermo:freeenergyfromlogdet"]
  \inputleannode{InfoGeometry.Thermo.freeEnergyFromLogDet}
\end{theorem}

\subsection*{freeEnergyFromLogDet_eq_internal_sub_scale_entropy}

\begin{theorem}[blueprint="infogeometry:thermo:freeenergyfromlogdet_eq_internal_sub_scale_entropy"]
  \inputleannode{InfoGeometry.Thermo.freeEnergyFromLogDet_eq_internal_sub_scale_entropy}
\end{theorem}

\subsection*{freeEnergyFromLogDet_eq_internal_sub_scale_entropy_of_pos}

\begin{theorem}[blueprint="infogeometry:thermo:freeenergyfromlogdet_eq_internal_sub_scale_entropy_of_pos"]
  \inputleannode{InfoGeometry.Thermo.freeEnergyFromLogDet_eq_internal_sub_scale_entropy_of_pos}
\end{theorem}

\subsection*{freeEnergyFromLogDet_eq_neg_scale_log_partition}

\begin{theorem}[blueprint="infogeometry:thermo:freeenergyfromlogdet_eq_neg_scale_log_partition"]
  \inputleannode{InfoGeometry.Thermo.freeEnergyFromLogDet_eq_neg_scale_log_partition}
\end{theorem}

\subsection*{freeEnergy_def'}

\begin{theorem}[blueprint="infogeometry:thermo:freeenergy_def'"]
  \inputleannode{InfoGeometry.Thermo.freeEnergy_def'}
\end{theorem}

\subsection*{freeEnergy_eq_internal_sub_scale_entropy}

\begin{theorem}[blueprint="infogeometry:thermo:freeenergy_eq_internal_sub_scale_entropy"]
  \inputleannode{InfoGeometry.Thermo.freeEnergy_eq_internal_sub_scale_entropy}
\end{theorem}

\subsection*{gibbsProb}

\begin{theorem}[blueprint="infogeometry:thermo:gibbsprob"]
  \inputleannode{InfoGeometry.Thermo.gibbsProb}
\end{theorem}

\subsection*{gibbsProbFromBregman}

\begin{theorem}[blueprint="infogeometry:thermo:gibbsprobfrombregman"]
  \inputleannode{InfoGeometry.Thermo.gibbsProbFromBregman}
\end{theorem}

\subsection*{gibbsProbFromBregman_eq_weight_over_partition}

\begin{theorem}[blueprint="infogeometry:thermo:gibbsprobfrombregman_eq_weight_over_partition"]
  \inputleannode{InfoGeometry.Thermo.gibbsProbFromBregman_eq_weight_over_partition}
\end{theorem}

\subsection*{gibbsProbFromBregman_nonneg}

\begin{theorem}[blueprint="infogeometry:thermo:gibbsprobfrombregman_nonneg"]
  \inputleannode{InfoGeometry.Thermo.gibbsProbFromBregman_nonneg}
\end{theorem}

\subsection*{gibbsProbFromBregman_sum_one}

\begin{theorem}[blueprint="infogeometry:thermo:gibbsprobfrombregman_sum_one"]
  \inputleannode{InfoGeometry.Thermo.gibbsProbFromBregman_sum_one}
\end{theorem}

\subsection*{gibbsProbFromLogDet}

\begin{theorem}[blueprint="infogeometry:thermo:gibbsprobfromlogdet"]
  \inputleannode{InfoGeometry.Thermo.gibbsProbFromLogDet}
\end{theorem}

\subsection*{gibbsProbFromLogDet_nonneg}

\begin{theorem}[blueprint="infogeometry:thermo:gibbsprobfromlogdet_nonneg"]
  \inputleannode{InfoGeometry.Thermo.gibbsProbFromLogDet_nonneg}
\end{theorem}

\subsection*{gibbsProbFromLogDet_sum_one}

\begin{theorem}[blueprint="infogeometry:thermo:gibbsprobfromlogdet_sum_one"]
  \inputleannode{InfoGeometry.Thermo.gibbsProbFromLogDet_sum_one}
\end{theorem}

\subsection*{gibbsProb_nonneg}

\begin{theorem}[blueprint="infogeometry:thermo:gibbsprob_nonneg"]
  \inputleannode{InfoGeometry.Thermo.gibbsProb_nonneg}
\end{theorem}

\subsection*{gibbsProb_pos}

\begin{theorem}[blueprint="infogeometry:thermo:gibbsprob_pos"]
  \inputleannode{InfoGeometry.Thermo.gibbsProb_pos}
\end{theorem}

\subsection*{gibbsProb_sum_one}

\begin{theorem}[blueprint="infogeometry:thermo:gibbsprob_sum_one"]
  \inputleannode{InfoGeometry.Thermo.gibbsProb_sum_one}
\end{theorem}

\subsection*{gibbsWeight}

\begin{theorem}[blueprint="infogeometry:thermo:gibbsweight"]
  \inputleannode{InfoGeometry.Thermo.gibbsWeight}
\end{theorem}

\subsection*{gibbsWeightDivergence}

\begin{theorem}[blueprint="infogeometry:thermo:gibbsweightdivergence"]
  \inputleannode{InfoGeometry.Thermo.gibbsWeightDivergence}
\end{theorem}

\subsection*{internalEnergy}

\begin{theorem}[blueprint="infogeometry:thermo:internalenergy"]
  \inputleannode{InfoGeometry.Thermo.internalEnergy}
\end{theorem}

\subsection*{internalEnergy_sub_freeEnergy}

\begin{theorem}[blueprint="infogeometry:thermo:internalenergy_sub_freeenergy"]
  \inputleannode{InfoGeometry.Thermo.internalEnergy_sub_freeEnergy}
\end{theorem}

\subsection*{logZ}

\begin{theorem}[blueprint="infogeometry:thermo:logz"]
  \inputleannode{InfoGeometry.Thermo.logZ}
\end{theorem}

\subsection*{log_gibbsProb}

\begin{theorem}[blueprint="infogeometry:thermo:log_gibbsprob"]
  \inputleannode{InfoGeometry.Thermo.log_gibbsProb}
\end{theorem}

\subsection*{partition}

\begin{theorem}[blueprint="infogeometry:thermo:partition"]
  \inputleannode{InfoGeometry.Thermo.partition}
\end{theorem}

\subsection*{partitionDivergence}

\begin{theorem}[blueprint="infogeometry:thermo:partitiondivergence"]
  \inputleannode{InfoGeometry.Thermo.partitionDivergence}
\end{theorem}

\subsection*{partitionDivergence_eq_Z}

\begin{theorem}[blueprint="infogeometry:thermo:partitiondivergence_eq_z"]
  \inputleannode{InfoGeometry.Thermo.partitionDivergence_eq_Z}
\end{theorem}

\subsection*{partitionDivergence_ne_zero}

\begin{theorem}[blueprint="infogeometry:thermo:partitiondivergence_ne_zero"]
  \inputleannode{InfoGeometry.Thermo.partitionDivergence_ne_zero}
\end{theorem}

\subsection*{partitionDivergence_pos}

\begin{theorem}[blueprint="infogeometry:thermo:partitiondivergence_pos"]
  \inputleannode{InfoGeometry.Thermo.partitionDivergence_pos}
\end{theorem}

\subsection*{partitionFromLogDet}

\begin{theorem}[blueprint="infogeometry:thermo:partitionfromlogdet"]
  \inputleannode{InfoGeometry.Thermo.partitionFromLogDet}
\end{theorem}

\subsection*{partitionFromLogDet_def}

\begin{theorem}[blueprint="infogeometry:thermo:partitionfromlogdet_def"]
  \inputleannode{InfoGeometry.Thermo.partitionFromLogDet_def}
\end{theorem}

\subsection*{partitionFromLogDet_ne_zero}

\begin{theorem}[blueprint="infogeometry:thermo:partitionfromlogdet_ne_zero"]
  \inputleannode{InfoGeometry.Thermo.partitionFromLogDet_ne_zero}
\end{theorem}

\subsection*{partitionFromLogDet_pos}

\begin{theorem}[blueprint="infogeometry:thermo:partitionfromlogdet_pos"]
  \inputleannode{InfoGeometry.Thermo.partitionFromLogDet_pos}
\end{theorem}

\subsection*{partitionFromLogDet_pos'}

\begin{theorem}[blueprint="infogeometry:thermo:partitionfromlogdet_pos'"]
  \inputleannode{InfoGeometry.Thermo.partitionFromLogDet_pos'}
\end{theorem}

\subsection*{shannonEntropy}

\begin{theorem}[blueprint="infogeometry:thermo:shannonentropy"]
  \inputleannode{InfoGeometry.Thermo.shannonEntropy}
\end{theorem}

\subsection*{shannonEntropy_eq_internal_div_add_logZ}

\begin{theorem}[blueprint="infogeometry:thermo:shannonentropy_eq_internal_div_add_logz"]
  \inputleannode{InfoGeometry.Thermo.shannonEntropy_eq_internal_div_add_logZ}
\end{theorem}

\subsection*{softMin}

\begin{theorem}[blueprint="infogeometry:thermo:softmin"]
  \inputleannode{InfoGeometry.Thermo.softMin}
\end{theorem}

\subsection*{softMin_def}

\begin{theorem}[blueprint="infogeometry:thermo:softmin_def"]
  \inputleannode{InfoGeometry.Thermo.softMin_def}
\end{theorem}

\subsection*{weight}

\begin{theorem}[blueprint="infogeometry:thermo:weight"]
  \inputleannode{InfoGeometry.Thermo.weight}
\end{theorem}

\subsection*{weight_pos}

\begin{theorem}[blueprint="infogeometry:thermo:weight_pos"]
  \inputleannode{InfoGeometry.Thermo.weight_pos}
\end{theorem}

\section{InfoGeometry.Thermo.FiniteDiagonal}

\noindent\textbf{Buildable}: yes

\subsection*{Mat}

\begin{theorem}[blueprint="infogeometry:thermo:finitediagonal:mat"]
  \inputleannode{InfoGeometry.Thermo.FiniteDiagonal.Mat}
\end{theorem}

\subsection*{beta_mul_freeEnergy}

\begin{theorem}[blueprint="infogeometry:thermo:finitediagonal:beta_mul_freeenergy"]
  \inputleannode{InfoGeometry.Thermo.FiniteDiagonal.beta_mul_freeEnergy}
\end{theorem}

\subsection*{entropy}

\begin{theorem}[blueprint="infogeometry:thermo:finitediagonal:entropy"]
  \inputleannode{InfoGeometry.Thermo.FiniteDiagonal.entropy}
\end{theorem}

\subsection*{entropy_eq_beta_internal_plus_massieu}

\begin{theorem}[blueprint="infogeometry:thermo:finitediagonal:entropy_eq_beta_internal_plus_massieu"]
  \inputleannode{InfoGeometry.Thermo.FiniteDiagonal.entropy_eq_beta_internal_plus_massieu}
\end{theorem}

\subsection*{freeEnergy}

\begin{theorem}[blueprint="infogeometry:thermo:finitediagonal:freeenergy"]
  \inputleannode{InfoGeometry.Thermo.FiniteDiagonal.freeEnergy}
\end{theorem}

\subsection*{gibbsDensity}

\begin{theorem}[blueprint="infogeometry:thermo:finitediagonal:gibbsdensity"]
  \inputleannode{InfoGeometry.Thermo.FiniteDiagonal.gibbsDensity}
\end{theorem}

\subsection*{gibbsDensity_diag_pos}

\begin{theorem}[blueprint="infogeometry:thermo:finitediagonal:gibbsdensity_diag_pos"]
  \inputleannode{InfoGeometry.Thermo.FiniteDiagonal.gibbsDensity_diag_pos}
\end{theorem}

\subsection*{gibbsState}

\begin{theorem}[blueprint="infogeometry:thermo:finitediagonal:gibbsstate"]
  \inputleannode{InfoGeometry.Thermo.FiniteDiagonal.gibbsState}
\end{theorem}

\subsection*{gibbsState_modularShift_invariant}

\begin{theorem}[blueprint="infogeometry:thermo:finitediagonal:gibbsstate_modularshift_invariant"]
  \inputleannode{InfoGeometry.Thermo.FiniteDiagonal.gibbsState_modularShift_invariant}
\end{theorem}

\subsection*{gibbsState_one}

\begin{theorem}[blueprint="infogeometry:thermo:finitediagonal:gibbsstate_one"]
  \inputleannode{InfoGeometry.Thermo.FiniteDiagonal.gibbsState_one}
\end{theorem}

\subsection*{gibbsWeight}

\begin{theorem}[blueprint="infogeometry:thermo:finitediagonal:gibbsweight"]
  \inputleannode{InfoGeometry.Thermo.FiniteDiagonal.gibbsWeight}
\end{theorem}

\subsection*{gibbsWeight_eq_exp_logDensity}

\begin{theorem}[blueprint="infogeometry:thermo:finitediagonal:gibbsweight_eq_exp_logdensity"]
  \inputleannode{InfoGeometry.Thermo.FiniteDiagonal.gibbsWeight_eq_exp_logDensity}
\end{theorem}

\subsection*{gibbsWeight_nonneg}

\begin{theorem}[blueprint="infogeometry:thermo:finitediagonal:gibbsweight_nonneg"]
  \inputleannode{InfoGeometry.Thermo.FiniteDiagonal.gibbsWeight_nonneg}
\end{theorem}

\subsection*{gibbsWeight_pos}

\begin{theorem}[blueprint="infogeometry:thermo:finitediagonal:gibbsweight_pos"]
  \inputleannode{InfoGeometry.Thermo.FiniteDiagonal.gibbsWeight_pos}
\end{theorem}

\subsection*{gibbsWeight_sum_one}

\begin{theorem}[blueprint="infogeometry:thermo:finitediagonal:gibbsweight_sum_one"]
  \inputleannode{InfoGeometry.Thermo.FiniteDiagonal.gibbsWeight_sum_one}
\end{theorem}

\subsection*{gibbsWeight_transport}

\begin{theorem}[blueprint="infogeometry:thermo:finitediagonal:gibbsweight_transport"]
  \inputleannode{InfoGeometry.Thermo.FiniteDiagonal.gibbsWeight_transport}
\end{theorem}

\subsection*{gibbs_detailedBalance_entry}

\begin{theorem}[blueprint="infogeometry:thermo:finitediagonal:gibbs_detailedbalance_entry"]
  \inputleannode{InfoGeometry.Thermo.FiniteDiagonal.gibbs_detailedBalance_entry}
\end{theorem}

\subsection*{hamiltonianOp}

\begin{theorem}[blueprint="infogeometry:thermo:finitediagonal:hamiltonianop"]
  \inputleannode{InfoGeometry.Thermo.FiniteDiagonal.hamiltonianOp}
\end{theorem}

\subsection*{internalEnergy}

\begin{theorem}[blueprint="infogeometry:thermo:finitediagonal:internalenergy"]
  \inputleannode{InfoGeometry.Thermo.FiniteDiagonal.internalEnergy}
\end{theorem}

\subsection*{internalEnergy_eq_gibbsState_hamiltonian}

\begin{theorem}[blueprint="infogeometry:thermo:finitediagonal:internalenergy_eq_gibbsstate_hamiltonian"]
  \inputleannode{InfoGeometry.Thermo.FiniteDiagonal.internalEnergy_eq_gibbsState_hamiltonian}
\end{theorem}

\subsection*{logDensityEntry}

\begin{theorem}[blueprint="infogeometry:thermo:finitediagonal:logdensityentry"]
  \inputleannode{InfoGeometry.Thermo.FiniteDiagonal.logDensityEntry}
\end{theorem}

\subsection*{logDensityOp}

\begin{theorem}[blueprint="infogeometry:thermo:finitediagonal:logdensityop"]
  \inputleannode{InfoGeometry.Thermo.FiniteDiagonal.logDensityOp}
\end{theorem}

\subsection*{logDensityOp_exp_entry}

\begin{theorem}[blueprint="infogeometry:thermo:finitediagonal:logdensityop_exp_entry"]
  \inputleannode{InfoGeometry.Thermo.FiniteDiagonal.logDensityOp_exp_entry}
\end{theorem}

\subsection*{log_gibbsWeight}

\begin{theorem}[blueprint="infogeometry:thermo:finitediagonal:log_gibbsweight"]
  \inputleannode{InfoGeometry.Thermo.FiniteDiagonal.log_gibbsWeight}
\end{theorem}

\subsection*{massieu}

\begin{theorem}[blueprint="infogeometry:thermo:finitediagonal:massieu"]
  \inputleannode{InfoGeometry.Thermo.FiniteDiagonal.massieu}
\end{theorem}

\subsection*{modularShift}

\begin{theorem}[blueprint="infogeometry:thermo:finitediagonal:modularshift"]
  \inputleannode{InfoGeometry.Thermo.FiniteDiagonal.modularShift}
\end{theorem}

\subsection*{modularShift_add}

\begin{theorem}[blueprint="infogeometry:thermo:finitediagonal:modularshift_add"]
  \inputleannode{InfoGeometry.Thermo.FiniteDiagonal.modularShift_add}
\end{theorem}

\subsection*{modularShift_apply}

\begin{theorem}[blueprint="infogeometry:thermo:finitediagonal:modularshift_apply"]
  \inputleannode{InfoGeometry.Thermo.FiniteDiagonal.modularShift_apply}
\end{theorem}

\subsection*{modularShift_diag_fixed}

\begin{theorem}[blueprint="infogeometry:thermo:finitediagonal:modularshift_diag_fixed"]
  \inputleannode{InfoGeometry.Thermo.FiniteDiagonal.modularShift_diag_fixed}
\end{theorem}

\subsection*{modularShift_zero}

\begin{theorem}[blueprint="infogeometry:thermo:finitediagonal:modularshift_zero"]
  \inputleannode{InfoGeometry.Thermo.FiniteDiagonal.modularShift_zero}
\end{theorem}

\subsection*{partition}

\begin{theorem}[blueprint="infogeometry:thermo:finitediagonal:partition"]
  \inputleannode{InfoGeometry.Thermo.FiniteDiagonal.partition}
\end{theorem}

\subsection*{partition_ne_zero}

\begin{theorem}[blueprint="infogeometry:thermo:finitediagonal:partition_ne_zero"]
  \inputleannode{InfoGeometry.Thermo.FiniteDiagonal.partition_ne_zero}
\end{theorem}

\subsection*{partition_pos}

\begin{theorem}[blueprint="infogeometry:thermo:finitediagonal:partition_pos"]
  \inputleannode{InfoGeometry.Thermo.FiniteDiagonal.partition_pos}
\end{theorem}

\section{InfoGeometry.Thermo.FiniteDiagonal.gibbsDensity}

\subsection*{eq_1}

\begin{theorem}[blueprint="infogeometry:thermo:finitediagonal:gibbsdensity:eq_1"]
  \inputleannode{InfoGeometry.Thermo.FiniteDiagonal.gibbsDensity.eq_1}
\end{theorem}

\section{InfoGeometry.Thermo.FiniteDiagonal.logDensityOp}

\subsection*{eq_1}

\begin{theorem}[blueprint="infogeometry:thermo:finitediagonal:logdensityop:eq_1"]
  \inputleannode{InfoGeometry.Thermo.FiniteDiagonal.logDensityOp.eq_1}
\end{theorem}

\section{InfoGeometry.Thermo.FiniteDiagonal.modularShift}

\subsection*{eq_1}

\begin{theorem}[blueprint="infogeometry:thermo:finitediagonal:modularshift:eq_1"]
  \inputleannode{InfoGeometry.Thermo.FiniteDiagonal.modularShift.eq_1}
\end{theorem}

\section{InfoGeometry.Thermo.FiniteMatrix}

\noindent\textbf{Buildable}: yes

\noindent\emph{No exported declarations in the current declaration graph.}

\section{InfoGeometry.Thermo.FromBregman}

\noindent\textbf{Buildable}: yes

\noindent\emph{No exported declarations in the current declaration graph.}

\section{InfoGeometry.Thermo.FromLogDet}

\noindent\textbf{Buildable}: yes

\noindent\emph{No exported declarations in the current declaration graph.}

\section{InfoGeometry.Thermo.Gibbs}

\noindent\textbf{Buildable}: yes

\noindent\emph{No exported declarations in the current declaration graph.}

\section{InfoGeometry.Thermo.ThermalModel}

\subsection*{H}

\begin{theorem}[blueprint="infogeometry:thermo:thermalmodel:h"]
  \inputleannode{InfoGeometry.Thermo.ThermalModel.H}
\end{theorem}

\subsection*{SatisfiesKMSLike}

\begin{theorem}[blueprint="infogeometry:thermo:thermalmodel:satisfieskmslike"]
  \inputleannode{InfoGeometry.Thermo.ThermalModel.SatisfiesKMSLike}
\end{theorem}

\subsection*{SatisfiesKMSLikeOn}

\begin{theorem}[blueprint="infogeometry:thermo:thermalmodel:satisfieskmslikeon"]
  \inputleannode{InfoGeometry.Thermo.ThermalModel.SatisfiesKMSLikeOn}
\end{theorem}

\subsection*{casesOn}

\begin{theorem}[blueprint="infogeometry:thermo:thermalmodel:caseson"]
  \inputleannode{InfoGeometry.Thermo.ThermalModel.casesOn}
\end{theorem}

\subsection*{ctorIdx}

\begin{theorem}[blueprint="infogeometry:thermo:thermalmodel:ctoridx"]
  \inputleannode{InfoGeometry.Thermo.ThermalModel.ctorIdx}
\end{theorem}

\subsection*{densityMatrix}

\begin{theorem}[blueprint="infogeometry:thermo:thermalmodel:densitymatrix"]
  \inputleannode{InfoGeometry.Thermo.ThermalModel.densityMatrix}
\end{theorem}

\subsection*{densityMatrix_trace_one}

\begin{theorem}[blueprint="infogeometry:thermo:thermalmodel:densitymatrix_trace_one"]
  \inputleannode{InfoGeometry.Thermo.ThermalModel.densityMatrix_trace_one}
\end{theorem}

\subsection*{gibbsExpectation}

\begin{theorem}[blueprint="infogeometry:thermo:thermalmodel:gibbsexpectation"]
  \inputleannode{InfoGeometry.Thermo.ThermalModel.gibbsExpectation}
\end{theorem}

\subsection*{gibbsExpectation_add}

\begin{theorem}[blueprint="infogeometry:thermo:thermalmodel:gibbsexpectation_add"]
  \inputleannode{InfoGeometry.Thermo.ThermalModel.gibbsExpectation_add}
\end{theorem}

\subsection*{gibbsExpectation_cyclic_of_modularFixed_commute}

\begin{theorem}[blueprint="infogeometry:thermo:thermalmodel:gibbsexpectation_cyclic_of_modularfixed_commute"]
  \inputleannode{InfoGeometry.Thermo.ThermalModel.gibbsExpectation_cyclic_of_modularFixed_commute}
\end{theorem}

\subsection*{gibbsExpectation_kms_like_of_fixedpoint}

\begin{theorem}[blueprint="infogeometry:thermo:thermalmodel:gibbsexpectation_kms_like_of_fixedpoint"]
  \inputleannode{InfoGeometry.Thermo.ThermalModel.gibbsExpectation_kms_like_of_fixedpoint}
\end{theorem}

\subsection*{gibbsExpectation_mul_swap_of_commute_density}

\begin{theorem}[blueprint="infogeometry:thermo:thermalmodel:gibbsexpectation_mul_swap_of_commute_density"]
  \inputleannode{InfoGeometry.Thermo.ThermalModel.gibbsExpectation_mul_swap_of_commute_density}
\end{theorem}

\subsection*{gibbsExpectation_one}

\begin{theorem}[blueprint="infogeometry:thermo:thermalmodel:gibbsexpectation_one"]
  \inputleannode{InfoGeometry.Thermo.ThermalModel.gibbsExpectation_one}
\end{theorem}

\subsection*{gibbsState_satisfiesKMSLike_on}

\begin{theorem}[blueprint="infogeometry:thermo:thermalmodel:gibbsstate_satisfieskmslike_on"]
  \inputleannode{InfoGeometry.Thermo.ThermalModel.gibbsState_satisfiesKMSLike_on}
\end{theorem}

\subsection*{gibbsState_satisfiesKMSLike_on_fixed_commuting_set}

\begin{theorem}[blueprint="infogeometry:thermo:thermalmodel:gibbsstate_satisfieskmslike_on_fixed_commuting_set"]
  \inputleannode{InfoGeometry.Thermo.ThermalModel.gibbsState_satisfiesKMSLike_on_fixed_commuting_set}
\end{theorem}

\subsection*{gibbsWeight}

\begin{theorem}[blueprint="infogeometry:thermo:thermalmodel:gibbsweight"]
  \inputleannode{InfoGeometry.Thermo.ThermalModel.gibbsWeight}
\end{theorem}

\subsection*{logUnnormalizedDensity}

\begin{theorem}[blueprint="infogeometry:thermo:thermalmodel:logunnormalizeddensity"]
  \inputleannode{InfoGeometry.Thermo.ThermalModel.logUnnormalizedDensity}
\end{theorem}

\subsection*{mk}

\begin{theorem}[blueprint="infogeometry:thermo:thermalmodel:mk"]
  \inputleannode{InfoGeometry.Thermo.ThermalModel.mk}
\end{theorem}

\subsection*{modularShift}

\begin{theorem}[blueprint="infogeometry:thermo:thermalmodel:modularshift"]
  \inputleannode{InfoGeometry.Thermo.ThermalModel.modularShift}
\end{theorem}

\subsection*{modularShiftFromLogDensity}

\begin{theorem}[blueprint="infogeometry:thermo:thermalmodel:modularshiftfromlogdensity"]
  \inputleannode{InfoGeometry.Thermo.ThermalModel.modularShiftFromLogDensity}
\end{theorem}

\subsection*{modularShiftFromLogDensity_eq_modularShift}

\begin{theorem}[blueprint="infogeometry:thermo:thermalmodel:modularshiftfromlogdensity_eq_modularshift"]
  \inputleannode{InfoGeometry.Thermo.ThermalModel.modularShiftFromLogDensity_eq_modularShift}
\end{theorem}

\subsection*{modularShift_zero}

\begin{theorem}[blueprint="infogeometry:thermo:thermalmodel:modularshift_zero"]
  \inputleannode{InfoGeometry.Thermo.ThermalModel.modularShift_zero}
\end{theorem}

\subsection*{noConfusion}

\begin{theorem}[blueprint="infogeometry:thermo:thermalmodel:noconfusion"]
  \inputleannode{InfoGeometry.Thermo.ThermalModel.noConfusion}
\end{theorem}

\subsection*{noConfusionType}

\begin{theorem}[blueprint="infogeometry:thermo:thermalmodel:noconfusiontype"]
  \inputleannode{InfoGeometry.Thermo.ThermalModel.noConfusionType}
\end{theorem}

\subsection*{partitionFunction}

\begin{theorem}[blueprint="infogeometry:thermo:thermalmodel:partitionfunction"]
  \inputleannode{InfoGeometry.Thermo.ThermalModel.partitionFunction}
\end{theorem}

\subsection*{partitionFunction_ne_zero}

\begin{theorem}[blueprint="infogeometry:thermo:thermalmodel:partitionfunction_ne_zero"]
  \inputleannode{InfoGeometry.Thermo.ThermalModel.partitionFunction_ne_zero}
\end{theorem}

\subsection*{partition_ne_zero}

\begin{theorem}[blueprint="infogeometry:thermo:thermalmodel:partition_ne_zero"]
  \inputleannode{InfoGeometry.Thermo.ThermalModel.partition_ne_zero}
\end{theorem}

\subsection*{rec}

\begin{theorem}[blueprint="infogeometry:thermo:thermalmodel:rec"]
  \inputleannode{InfoGeometry.Thermo.ThermalModel.rec}
\end{theorem}

\subsection*{recOn}

\begin{theorem}[blueprint="infogeometry:thermo:thermalmodel:recon"]
  \inputleannode{InfoGeometry.Thermo.ThermalModel.recOn}
\end{theorem}

\subsection*{β}

\begin{theorem}[blueprint="infogeometry:thermo:thermalmodel:β"]
  \inputleannode{InfoGeometry.Thermo.ThermalModel.β}
\end{theorem}

\section{InfoGeometry.Thermo.ThermalModel.gibbsWeight}

\subsection*{eq_1}

\begin{theorem}[blueprint="infogeometry:thermo:thermalmodel:gibbsweight:eq_1"]
  \inputleannode{InfoGeometry.Thermo.ThermalModel.gibbsWeight.eq_1}
\end{theorem}

\section{InfoGeometry.Thermo.ThermalModel.logUnnormalizedDensity}

\subsection*{eq_1}

\begin{theorem}[blueprint="infogeometry:thermo:thermalmodel:logunnormalizeddensity:eq_1"]
  \inputleannode{InfoGeometry.Thermo.ThermalModel.logUnnormalizedDensity.eq_1}
\end{theorem}

\section{InfoGeometry.Thermo.ThermalModel.partitionFunction}

\subsection*{eq_1}

\begin{theorem}[blueprint="infogeometry:thermo:thermalmodel:partitionfunction:eq_1"]
  \inputleannode{InfoGeometry.Thermo.ThermalModel.partitionFunction.eq_1}
\end{theorem}

\section{InfoGeometry.Thermo.ThermodynamicIdentities}

\noindent\textbf{Buildable}: yes

\noindent\emph{No exported declarations in the current declaration graph.}

\section{InfoGeometry.Thermo.Z}

\subsection*{eq_1}

\begin{theorem}[blueprint="infogeometry:thermo:z:eq_1"]
  \inputleannode{InfoGeometry.Thermo.Z.eq_1}
\end{theorem}

\section{InfoGeometry.Thermo.energyFromBregman}

\subsection*{eq_1}

\begin{theorem}[blueprint="infogeometry:thermo:energyfrombregman:eq_1"]
  \inputleannode{InfoGeometry.Thermo.energyFromBregman.eq_1}
\end{theorem}

\section{InfoGeometry.Thermo.expectation}

\subsection*{eq_1}

\begin{theorem}[blueprint="infogeometry:thermo:expectation:eq_1"]
  \inputleannode{InfoGeometry.Thermo.expectation.eq_1}
\end{theorem}

\section{InfoGeometry.Thermo.freeEnergy}

\subsection*{eq_1}

\begin{theorem}[blueprint="infogeometry:thermo:freeenergy:eq_1"]
  \inputleannode{InfoGeometry.Thermo.freeEnergy.eq_1}
\end{theorem}

\section{InfoGeometry.Thermo.freeEnergyDivergence}

\subsection*{eq_1}

\begin{theorem}[blueprint="infogeometry:thermo:freeenergydivergence:eq_1"]
  \inputleannode{InfoGeometry.Thermo.freeEnergyDivergence.eq_1}
\end{theorem}

\section{InfoGeometry.Thermo.freeEnergyFromBregman}

\subsection*{eq_1}

\begin{theorem}[blueprint="infogeometry:thermo:freeenergyfrombregman:eq_1"]
  \inputleannode{InfoGeometry.Thermo.freeEnergyFromBregman.eq_1}
\end{theorem}

\section{InfoGeometry.Thermo.freeEnergyFromLogDet}

\subsection*{eq_1}

\begin{theorem}[blueprint="infogeometry:thermo:freeenergyfromlogdet:eq_1"]
  \inputleannode{InfoGeometry.Thermo.freeEnergyFromLogDet.eq_1}
\end{theorem}

\section{InfoGeometry.Thermo.gibbsProb}

\subsection*{eq_1}

\begin{theorem}[blueprint="infogeometry:thermo:gibbsprob:eq_1"]
  \inputleannode{InfoGeometry.Thermo.gibbsProb.eq_1}
\end{theorem}

\section{InfoGeometry.Thermo.gibbsProbFromBregman}

\subsection*{eq_1}

\begin{theorem}[blueprint="infogeometry:thermo:gibbsprobfrombregman:eq_1"]
  \inputleannode{InfoGeometry.Thermo.gibbsProbFromBregman.eq_1}
\end{theorem}

\section{InfoGeometry.Thermo.gibbsProbFromLogDet}

\subsection*{eq_1}

\begin{theorem}[blueprint="infogeometry:thermo:gibbsprobfromlogdet:eq_1"]
  \inputleannode{InfoGeometry.Thermo.gibbsProbFromLogDet.eq_1}
\end{theorem}

\section{InfoGeometry.Thermo.gibbsWeightDivergence}

\subsection*{eq_1}

\begin{theorem}[blueprint="infogeometry:thermo:gibbsweightdivergence:eq_1"]
  \inputleannode{InfoGeometry.Thermo.gibbsWeightDivergence.eq_1}
\end{theorem}

\section{InfoGeometry.Thermo.internalEnergy}

\subsection*{eq_1}

\begin{theorem}[blueprint="infogeometry:thermo:internalenergy:eq_1"]
  \inputleannode{InfoGeometry.Thermo.internalEnergy.eq_1}
\end{theorem}

\section{InfoGeometry.Thermo.logZ}

\subsection*{eq_1}

\begin{theorem}[blueprint="infogeometry:thermo:logz:eq_1"]
  \inputleannode{InfoGeometry.Thermo.logZ.eq_1}
\end{theorem}

\section{InfoGeometry.Thermo.partitionDivergence}

\subsection*{eq_1}

\begin{theorem}[blueprint="infogeometry:thermo:partitiondivergence:eq_1"]
  \inputleannode{InfoGeometry.Thermo.partitionDivergence.eq_1}
\end{theorem}

\section{InfoGeometry.Thermo.partitionFromLogDet}

\subsection*{eq_1}

\begin{theorem}[blueprint="infogeometry:thermo:partitionfromlogdet:eq_1"]
  \inputleannode{InfoGeometry.Thermo.partitionFromLogDet.eq_1}
\end{theorem}

\section{InfoGeometry.Thermo.softMin}

\subsection*{eq_1}

\begin{theorem}[blueprint="infogeometry:thermo:softmin:eq_1"]
  \inputleannode{InfoGeometry.Thermo.softMin.eq_1}
\end{theorem}

\section{InfoGeometry.Thermo.weight}

\subsection*{eq_1}

\begin{theorem}[blueprint="infogeometry:thermo:weight:eq_1"]
  \inputleannode{InfoGeometry.Thermo.weight.eq_1}
\end{theorem}

\chapter{TransformationGroups Modules}

\section{InfoGeometry.TransformationGroups}

\noindent\textbf{Buildable}: yes

\subsection*{FinProb}

\begin{theorem}[blueprint="infogeometry:transformationgroups:finprob"]
  \inputleannode{InfoGeometry.TransformationGroups.FinProb}
\end{theorem}

\subsection*{InvariantUnder}

\begin{theorem}[blueprint="infogeometry:transformationgroups:invariantunder"]
  \inputleannode{InfoGeometry.TransformationGroups.InvariantUnder}
\end{theorem}

\subsection*{IsTransitive}

\begin{theorem}[blueprint="infogeometry:transformationgroups:istransitive"]
  \inputleannode{InfoGeometry.TransformationGroups.IsTransitive}
\end{theorem}

\subsection*{LocationInvariant}

\begin{theorem}[blueprint="infogeometry:transformationgroups:locationinvariant"]
  \inputleannode{InfoGeometry.TransformationGroups.LocationInvariant}
\end{theorem}

\subsection*{ScaleInvariant}

\begin{theorem}[blueprint="infogeometry:transformationgroups:scaleinvariant"]
  \inputleannode{InfoGeometry.TransformationGroups.ScaleInvariant}
\end{theorem}

\subsection*{eq_of_invariant_transitive}

\begin{theorem}[blueprint="infogeometry:transformationgroups:eq_of_invariant_transitive"]
  \inputleannode{InfoGeometry.TransformationGroups.eq_of_invariant_transitive}
\end{theorem}

\subsection*{locationInvariant_const}

\begin{theorem}[blueprint="infogeometry:transformationgroups:locationinvariant_const"]
  \inputleannode{InfoGeometry.TransformationGroups.locationInvariant_const}
\end{theorem}

\subsection*{logCoord_flat}

\begin{theorem}[blueprint="infogeometry:transformationgroups:logcoord_flat"]
  \inputleannode{InfoGeometry.TransformationGroups.logCoord_flat}
\end{theorem}

\subsection*{scaleInvariant_inv}

\begin{theorem}[blueprint="infogeometry:transformationgroups:scaleinvariant_inv"]
  \inputleannode{InfoGeometry.TransformationGroups.scaleInvariant_inv}
\end{theorem}

\subsection*{uniform_of_all_eq}

\begin{theorem}[blueprint="infogeometry:transformationgroups:uniform_of_all_eq"]
  \inputleannode{InfoGeometry.TransformationGroups.uniform_of_all_eq}
\end{theorem}

\subsection*{uniform_of_transformation_group}

\begin{theorem}[blueprint="infogeometry:transformationgroups:uniform_of_transformation_group"]
  \inputleannode{InfoGeometry.TransformationGroups.uniform_of_transformation_group}
\end{theorem}

\chapter{Twistor Modules}

\section{InfoGeometry.Twistor}

\subsection*{ChiralFactorization}

\begin{theorem}[blueprint="infogeometry:twistor:chiralfactorization"]
  \inputleannode{InfoGeometry.Twistor.ChiralFactorization}
\end{theorem}

\subsection*{DoubledProjectiveState}

\begin{theorem}[blueprint="infogeometry:twistor:doubledprojectivestate"]
  \inputleannode{InfoGeometry.Twistor.DoubledProjectiveState}
\end{theorem}

\subsection*{DoubledTwistorSpace}

\begin{theorem}[blueprint="infogeometry:twistor:doubledtwistorspace"]
  \inputleannode{InfoGeometry.Twistor.DoubledTwistorSpace}
\end{theorem}

\subsection*{IsNull}

\begin{theorem}[blueprint="infogeometry:twistor:isnull"]
  \inputleannode{InfoGeometry.Twistor.IsNull}
\end{theorem}

\subsection*{TwistorSpace}

\begin{theorem}[blueprint="infogeometry:twistor:twistorspace"]
  \inputleannode{InfoGeometry.Twistor.TwistorSpace}
\end{theorem}

\subsection*{doubledTwistorMk}

\begin{theorem}[blueprint="infogeometry:twistor:doubledtwistormk"]
  \inputleannode{InfoGeometry.Twistor.doubledTwistorMk}
\end{theorem}

\subsection*{isNull_mk_iff}

\begin{theorem}[blueprint="infogeometry:twistor:isnull_mk_iff"]
  \inputleannode{InfoGeometry.Twistor.isNull_mk_iff}
\end{theorem}

\subsection*{isNull_projectivize_iff}

\begin{theorem}[blueprint="infogeometry:twistor:isnull_projectivize_iff"]
  \inputleannode{InfoGeometry.Twistor.isNull_projectivize_iff}
\end{theorem}

\subsection*{twistorMk}

\begin{theorem}[blueprint="infogeometry:twistor:twistormk"]
  \inputleannode{InfoGeometry.Twistor.twistorMk}
\end{theorem}

\section{InfoGeometry.Twistor.ChiralFactorization}

\subsection*{casesOn}

\begin{theorem}[blueprint="infogeometry:twistor:chiralfactorization:caseson"]
  \inputleannode{InfoGeometry.Twistor.ChiralFactorization.casesOn}
\end{theorem}

\subsection*{ctorIdx}

\begin{theorem}[blueprint="infogeometry:twistor:chiralfactorization:ctoridx"]
  \inputleannode{InfoGeometry.Twistor.ChiralFactorization.ctorIdx}
\end{theorem}

\subsection*{factor}

\begin{theorem}[blueprint="infogeometry:twistor:chiralfactorization:factor"]
  \inputleannode{InfoGeometry.Twistor.ChiralFactorization.factor}
\end{theorem}

\subsection*{factor_of_null}

\begin{theorem}[blueprint="infogeometry:twistor:chiralfactorization:factor_of_null"]
  \inputleannode{InfoGeometry.Twistor.ChiralFactorization.factor_of_null}
\end{theorem}

\subsection*{mk}

\begin{theorem}[blueprint="infogeometry:twistor:chiralfactorization:mk"]
  \inputleannode{InfoGeometry.Twistor.ChiralFactorization.mk}
\end{theorem}

\subsection*{noConfusion}

\begin{theorem}[blueprint="infogeometry:twistor:chiralfactorization:noconfusion"]
  \inputleannode{InfoGeometry.Twistor.ChiralFactorization.noConfusion}
\end{theorem}

\subsection*{noConfusionType}

\begin{theorem}[blueprint="infogeometry:twistor:chiralfactorization:noconfusiontype"]
  \inputleannode{InfoGeometry.Twistor.ChiralFactorization.noConfusionType}
\end{theorem}

\subsection*{null_of_factor}

\begin{theorem}[blueprint="infogeometry:twistor:chiralfactorization:null_of_factor"]
  \inputleannode{InfoGeometry.Twistor.ChiralFactorization.null_of_factor}
\end{theorem}

\subsection*{rec}

\begin{theorem}[blueprint="infogeometry:twistor:chiralfactorization:rec"]
  \inputleannode{InfoGeometry.Twistor.ChiralFactorization.rec}
\end{theorem}

\subsection*{recOn}

\begin{theorem}[blueprint="infogeometry:twistor:chiralfactorization:recon"]
  \inputleannode{InfoGeometry.Twistor.ChiralFactorization.recOn}
\end{theorem}

\section{InfoGeometry.Twistor.IsNull}

\subsection*{eq_1}

\begin{theorem}[blueprint="infogeometry:twistor:isnull:eq_1"]
  \inputleannode{InfoGeometry.Twistor.IsNull.eq_1}
\end{theorem}

\section{InfoGeometry.Twistor.NullProjective}

\noindent\textbf{Buildable}: yes

\noindent\emph{No exported declarations in the current declaration graph.}

\chapter{auto_blueprints Modules}

\section{InfoGeometry.auto_blueprints}

\noindent\textbf{Buildable}: yes

\noindent\emph{No exported declarations in the current declaration graph.}

\chapter{bregmanDiv Modules}

\section{InfoGeometry.bregmanDiv}

\subsection*{eq_1}

\begin{theorem}[blueprint="infogeometry:bregmandiv:eq_1"]
  \inputleannode{InfoGeometry.bregmanDiv.eq_1}
\end{theorem}

\chapter{generalizedKL Modules}

\section{InfoGeometry.generalizedKL}

\noindent\textbf{Buildable}: yes

\noindent\emph{No exported declarations in the current declaration graph.}

